% 本文件由 LaTeX 排版 Agent 根据有效 translation.json 自动生成。
% author=Gregory of Nyssa; work_id=2901; title=驳斥尤诺米乌斯

\chapter{驳斥尤诺米乌斯}

% structure: p0001 source=h1 target=omitted reason=page-title-already-rendered-by-parent

第一卷\newline{}第二卷\newline{}第三卷\newline{}第四卷\newline{}第五卷\newline{}第六卷\newline{}第七卷\newline{}第八卷\newline{}第九卷\newline{}第十卷\newline{}第十一卷\newline{}第十二卷\par

\SourceRecord{Translated by W. Moore, H.A. Wilson and H.C. Ogle.；Nicene and Post-Nicene Fathers, Second Series；Vol. 5.；Edited by Philip Schaff and Henry Wace.；Buffalo, NY: Christian Literature Publishing Co.,；1893.；<http://www.newadvent.org/fathers/2901.htm>.}

% structure: p0001 source=h1 target=section reason=page-title

\section{\WorkTitle{驳斥犹诺米（卷一）}}

% structure: p0002 source=h2 target=subsection reason=relative-heading-rank; base=h2

\subsection{第一封序函}

\Person{额我略}致其兄弟、\Place{塞巴斯特}主教\Person{伯多禄}。\par

我好不容易得到了一点闲暇，从\Place{亚美尼亚}返回后得以恢复身体疲劳，并收集了我对\Person{犹诺米}的答复的稿子，这是出于您明智的建议；因此我的作品现在被整理成一部完整的论著，可以装订阅读。然而，我并没有针对他的两本小册子都写了反驳；甚至没有时间那样做；因为借给我那本异端卷册的人非常不客气地又把它要回去了，没有给我时间抄写或研究。在短短的十七天里，不可能准备好回应他的两篇攻击。\par

不知怎的，我们努力回应这篇亵渎宣言的事情变得众所周知，许多对真理有些热忱的人都催促我处理它；但我认为应当首先以您的智慧向您请教，是应该将这篇作品公之于众，还是采取其他方式。我犹豫的原因是这样的：当我们的圣\Person{巴西略}安息后，我接过了与\Person{犹诺米}争论的遗产，当时我心中因失去他而炽热，而且，除了这一共同损失带给教会的深切悲伤外，\Person{犹诺米}并不局限于各种可作为其观点辩护的话题，而是将其精力主要花费在辛苦编写的对我们神父的辱骂上。对此我感到愤怒，有些段落中我发自内心的义愤之火爆发出来，针对这位作者。公众在其他很多事情上已经原谅了我们，因为我们善于耐心应对无法无天的攻击，并尽可能实践了圣人所教导的感情节制；但我担心，从现在所写反对这位对手的文字中，读者会认为我们是非常生涩的辩论者，一遇到傲慢的辱骂就立刻发怒。然而，或许人们会因记得我们的愤怒并非为自己，而是因为对我们的神父的侮辱，从而打消对我们的怀疑；在这种情况下，温和比愤怒更加不可原谅。\par

因此，如果我这篇论著的第一部分似乎有点超出了争论的范围，我想对它的以下解释会得到能够公正判断的读者的接受。不为我们高贵圣人的名誉辩护是不对的，它被对手的亵渎所损害，正如让为他而战的争斗蔓延在整个讨论线索中也是不妥的；此外，如果有人思考，这些篇幅确实构成了争论的一部分。我们对手的论著有两条独立的武装，即辱骂我们和反驳正统教义；因此我们的论著也必须展示双面。但为了清晰起见，并且为了使关于信仰事务的讨论线索不被回应人身攻击的插入语打断，我们将作品分为两部分，首先致力于驳斥这些指控；然后尽全力应对他们对信仰提出的反对。我们论著还包含除了驳斥他们异端观点之外，对我们自己教导的教义阐述；因为当敌人毫不掩饰他们的亵渎时，我们不敢大胆陈述真理，那是极为可耻的态度缺乏。\par

% structure: p0007 source=h2 target=subsection reason=relative-heading-rank; base=h2

\subsection{第二封序函}

致最虔诚的兄弟\Person{额我略}。\Person{伯多禄}在主内问安。\par

我看到了您圣座的作品，并在您反对这一异端的短论中察觉了您对真理和对我们圣教父的热忱，我判断这部作品不仅单纯出于您的能力，更是源于您钻研如何使真理发言，即使是在发表自己观点的时候。我将这为真理的辩护归于真理的圣神；正如这敌对正统信仰的敌意应归于谎言之父，而非\Person{犹诺米}。确实，那从起初就是谋杀者、在\Person{犹诺米}内说话的那位，小心翼翼地磨快了剑刃指向自己；因为他若不是如此大胆反对真理，就没有人会激发您承担我们宗教的事业。但为要使他们道理的腐朽和脆弱暴露出来，祂\BibleQuote{祂使智者中了自己的诡计}，容许他们既顽固反对真理，又徒劳在这虚妄的言语上劳苦。\par

既然那开始了善工的必要完成它，所以不要懈怠推进圣神的能力，也不要半途放弃战胜攻击基督光荣之人的胜利；而要效法您的真父亲，他像热忱的\Person{丕乃哈斯}一样，用\WorkTitle{答辩}一击刺穿了师徒二人。用您智慧的臂膀将圣神的剑刺穿这两本异端小册子，免得那蛇虽断了头，尾巴仍颤抖着惊吓较为单纯的人。当最初的论点得到回答后，如果最后的论点仍未被注意，许多人会怀疑它们仍对真理持有一些力量。\par

您的短论中所表现的情感，将像盐一样给灵魂的味觉带来愉悦。正如根据\Person{约伯}所说，面包没有盐就不能吃，同样，没有浸透天主的圣言的深情的讲论，永远不会唤醒或引发渴望。\par

那么，要坚强，因为你是一个美好的榜样，给后来的时代展示了善良的子女应如何对待有德的父亲。\par

% structure: p0013 source=h2 target=subsection reason=relative-heading-rank; base=h2

\subsection{一、序言——试图帮助那些不愿接受帮助的人是徒劳的。}

看来，愿意造福众人，毫无区别地将自己的恩赐滥施于初来乍到者，在许多人眼中并非完全值得称赞，甚至不免指责；因为许多已预备好的药物白白浪费在无法治愈的病人身上，无论对于接受者还是对于那位有意行善者，都不会产生值得关心的结果：或使接受者获益，或使行善者获得声誉。相反，这样的尝试在许多情况下反而成为恶化的缘由。那些绝望患病、奄奄一息的患者，只会因更猛烈的药物而更快地终结；那凶猛而无理性的性情，只会在滥施珍珠的善意中变得更糟，正如福音告诉我们的。因此，我认为，按照天主的命令，最好在施与时将宝贵的与无价值的分开，从而避免慷慨施与者因那\Quote{践踏其珍珠}且因完全缺乏对美的感受而侮辱他的人所受的痛苦。\par

当我想到那位自由地将自己灵魂的美善通传给他人的人，就是那位天主的人、虔敬的口舌，\Person{巴西略}；他因自己属灵财富的富足，将智慧的恩宠倾注于他从未考验过的邪恶灵魂之中，并倾注于其中之一，\Person{欧诺弥}，此人对自己得救的一切努力完全无感。对于教会所有真正的成员而言，这个可怜人的状况确实可悲，因为他灵魂在信德方面极为软弱；因为谁如此缺乏感觉，以至于至少不怜悯一个沉沦的灵魂呢？但唯有\Person{巴西略}，出于他永恒炽热的爱，被感动去承担对他的医治，从而尝试不可能的事；唯有他如此关心这人绝望的状况，以至于创作了驳斥此异端的著作，作为致命毒药的解毒剂，旨在拯救其作者，并使他回归教会。\par

相反，他如同狂暴之人，抗拒他的医生；他争斗挣扎；他将那位竭尽全力把他从谬误深渊中拉出来的人视为死敌；他不仅偶尔在偶然听众面前发泄这愚蠢的愤怒；他还为自己树立了一座文学纪念碑，以记载他胆汁的黑暗；当经过多年获得必要的闲暇时，他在整个期间孕育他的作品，其痛苦比最大最笨拙的野兽的阵痛更为剧烈；当他还在秘密塑造他的构思时，他对即将到来的东西的威胁是可怕的：但当他最后艰难地将其带到光明中时，却是一个可怜的小流产儿，过早地出生了。然而，那些与他一同毁灭的人却哺育它、溺爱它；而我们，寻求先知所说的祝福（\BibleQuote{凡拿起你的幼童摔在磐石上的，便为有福}\BibleRef{咏 137:9}），只渴望，既然它已落入我们手中，就拿起这个哭泣的宣言，把它摔在磐石上，如同它是巴比伦的婴孩之一；而磐石必须是基督；换言之，是真理的宣告。但愿那在软弱中使力量成全之人的祈祷，将那加强软弱的能力降临在我们身上。\par

% structure: p0017 source=h2 target=subsection reason=relative-heading-rank; base=h2

\subsection{二、我们被\Person{欧诺弥}对我们兄弟的指控所刺痛，因而被正义地激发作出此答复。}

如果那神圣而圣善的灵魂仍在肉体中，关注人间事务，如果那些高亢的音调仍以其特有的优雅和不可抗拒的表达被听到，那么谁会有如此大的胆子，试图就这个问题说一句话呢？那神圣的号角之声会淹没任何可言说的话语。但现在，他的一切都已飞回天主；起初，在他身体那轻微模糊的幻影中，他仍留在地上；但现在，他甚至完全脱去了那无形的形状，并将其遗赠给这个世界。与此同时，黄蜂在圣言的蜂房周围嗡嗡作响，掠夺着蜜；所以，请任何人不要指责我只是出于胆大妄为，而站起来替那沉默的双唇说话。我接受这劳苦的工作，并非因为意识到自己在论辩能力上优于其他可能提到的人；我，如果有人知道，则深知在教会中有成千上万的人，他们得赐了哲学技艺的恩赐。尽管如此，我断言，无论根据成文法还是自然法，这逝者的遗产更特别地属于我；因此，我本人，优先于他人，承受这争论的遗赠。我可能被算作在天主教会中服役者中最小的一个，但我仍然不是太软弱，以至于不能作为她的捍卫者站出来反对一个与她决裂的人。一个强壮身体中最小的肢体，由于与整个身体的生命合一，会比一个被剪除而濒死的肢体更强壮，无论后者有多大，前者有多小。\par

% structure: p0019 source=h2 target=subsection reason=relative-heading-rank; base=h2

\subsection{三、我们在\Person{欧诺弥}的论文中看不出逻辑力量有任何非凡之处，因此我们以正义的信心开始我们的答复。}

不要让任何人认为，我说这话是在夸张，并空夸自己做能力之外的事。我不会被任何幼稚的野心驱使，在仅仅论证和措辞的较量中下降到他的粗俗水平。如果胜利是无用且无益的，我们乐意把它让给那些渴望胜利的人；此外，我们只需看看此人在争论中的长期实践，就可以断定他是一个十足的言辞匠人，再加上他在撰写这篇论文上花费了生命中的不少时间，以及他的密友们对这些劳动的极大喜悦，就可以断定他对这部作品特别费心。一个为此劳作这么多奥林匹亚周期的人，很可能比临时涂鸦者做得更好。就连他在炮制作品时所用的那些庸俗的比喻花样，也进一步表明这种辛劳的谨慎写作。他收集了大量新近搭配的术语，为此他利用了某些其他书籍，并在这非常单薄的思想核心上堆积了大量词语；因此他精心制作了这部矫揉造作的作品，他的错误学生们对此赞叹不已——无疑是因为他们在关键点上的麻木使他们无法感受美丑的区别——但在那些内心洞察力不被任何不信的污秽所蒙蔽的人看来，这是可笑且毫无价值的。他采用这些荒谬的言谈方式、这种新奇而独特的安排、这种忙碌的自负和自负的忙碌，对于任何先前的模式都没有热情，这怎么可能如他所希望的那样有助于证明他所说的以及确立他思辨的真理性呢？因为确实很难发现，在所有那些以口才著称的人中，他参考了谁，才使自己达到这种高度；因为他就像那些在舞台上制造效果的人，将自己的论证调整到节奏性短语的曲调，就像他们用平行的等长、相似声音和相似结尾的句子，将自己的歌与响板配合。在他的序言中，除其他许多缺陷外，还有这些无力的颤抖和矫揉造作的伎俩；人们可以想象他全部展现出来，不是以平静的动作，而是以跺脚和尖锐的响指，按照这样敲出的节拍朗诵，然后评论说不需要其他论证和之后第二次表演了。\par

% structure: p0021 source=h2 target=subsection reason=relative-heading-rank; base=h2

\subsection{四、\Person{欧诺弥}展现了许多愚昧和华丽的文笔，但在关键点上却很少认真。}

在这些和类似的花招中，我承认他占了上风；他尽可以陶醉于他的胜利。我非常愿意放弃这样的竞争，它只能吸引那些追求名声的人；确实，如果有什么名声来自沉迷于这种论证方式的话，考虑到\Person{保罗}\BibleRef{格前 2:1}，那位真正的圣言侍者，他的唯一装饰是真理，他既不屑于将自己的风格降低到这种华丽，也用一个高尚而恰当的劝诫，教导我们只关注真理。的确，一个拥有真理之美的人，何需引入装饰者的行头来制造虚假的人工美？或许对于那些不拥有真理的人来说，用吸引人的风格粉饰他们的谎言，并给他们的论证纹理磨上一层奇异的抛光，是有利的。当他们用牵强的语言和所有矫揉造作的风格来教导其错误时，他们有可能被听众接受而显得合理。但那些唯一目标是简单真理、不被任何误导性陪衬所玷污的人，从他们的话语中散发出自然美的光芒。\par

但现在我即将开始检查他所提出的一切，我感到与农夫在风平浪静时同样的困难；我不知道如何将他的麦子与糠秕分开；事实上，在这堆话语中，废物和糠秕如此之多，以至于让人认为他所说的一切中，事实和真实思想的残留几乎为零。详细审查他所有的言论，将不利于速度且非常烦人，甚至会有悖于我们的目标；我们没有条件获得如此多的闲暇，以至于可以任意将其花在这样无聊的事情上；我认为，一个精明工人的职责不是把力气浪费在琐事上，而是用在显然会回报劳苦的事情上。\par

至于他序言中的所有事情，他如何自封为真理的捍卫者，并将不信的指控加于反对者，并宣称一种持久且不可磨灭的仇恨已渗入他的灵魂，他如何在他的\Quote{新发现}中趾高气扬，虽然他没有告诉我们这些发现是什么，而只是说对其中争议点进行了审查，一场\Quote{合法}的审判，迫使那些胆敢违法的人保持安静，或者用他那种吕底亚式的歌唱风格（他有的）的话说：\Quote{公然违法的——在公开法庭上——被迫保持安静}（他称之为针对他的阴谋的\Quote{放逐令}，无论这个词可能是什么意思）——所有这些令人厌烦的事情，我都略过，认为很不重要。另一方面，他所有为他异端妄想所作的特别辩护，可能值得我们密切关注。我们自己对神学原理的解释者在他的论文中遵循了这条路线；因为他虽然有能力拓宽他的论证，但他只处理关键点，从那个异端书籍的所有亵渎中挑选出来，因此缩小了论题的范围。\par

然而，如果有人期望我们的答复与他的论点阵列完全对应，就请他告诉我们这一做法的效用。如果我像悲剧舞台上的斯芬克斯那样，去解开他在开篇就抛给我们的那个复杂的标题谜语——即这个\WorkTitle{为辩护辞的新辩护辞}——以及他围绕它写下的所有废话，再讲述他的梦呓长篇，读者又能得到什么好处呢？我认为，读者早已厌倦了其标题中关于\Quote{新}字的小虚荣——这已保留在欧诺弥自己的文本中——以及他在对自己事迹的叙述中所显露的粗俗乏味，他遍历大地和海洋，在全世界被\Quote{传扬}。如果这一切必须重新写下——而且还要添枝加叶，因为对这些谎言的驳斥自然会使叙述膨胀——谁能有铁石心肠而不厌恶这种浪费功夫呢？假如我逐字逐句写下对他那个疯故事的说明；例如，假如我解释那个尤克辛海岸的亚美尼亚人是谁——他起初因与自己同名而惹恼了他；他们的生活和追求如何；他如何因性格极其相似而与那亚美尼亚人争吵；然后那两人如何和好，以至于共同同情他那迷人而最荣耀的老师阿埃提乌（因为他的赞美是如此浮夸）；之后，又有什么阴谋对付他，使他们以他极受爱戴的罪名将他提交审判：假如我要解释这一切，我岂不是像那些因频繁接触眼疾患者而自己也染上此病的人一样，也染上了这种琐碎啰嗦的毛病吗？我将一步步追述他废话连篇的故事的每个细节；找出那些\Quote{被释放获得自由的奴隶}是谁，\Quote{受启蒙者的阴谋}和\Quote{雇来奴隶的呼唤}是什么，以及\Quote{蒙提乌斯、加卢斯和多米提安}、\Quote{假证人}、\Quote{愤怒的皇帝}和\Quote{某些被放逐者}与论点有何关系。对于一个希望不仅仅是叙述故事，而是要反驳那些反对其异端者的论点的人来说，还有什么比这些故事更无用的呢？故事的后续更加无益；我认为作者自己再读一遍也会打哈欠，尽管每个父亲对自己孩子\Emph{确实}怀有强烈的天然喜爱。他假装在那里展示自己的功绩和苦难；文风自命不凡，传说膨胀到悲剧的音调。\par

% structure: p0026 source=h2 target=subsection reason=relative-heading-rank; base=h2

\subsection{五、他对亚美尼亚主教欧斯塔修斯和加拉提亚主教巴西略的独特漫画并不成功。}

但是，为了不在拒绝提及这些荒谬言论时仍沉溺其中，也不因迫使我的主题遍历他书写的所有回忆录而玷污本书——就像一个人策马穿过泥沼而沾满污泥一样——我认为最好以我的思想所能达到的最高最快的跳跃，跳过他那堆垃圾，因为迅速远离令人厌恶之物乃是一大益处；让我们赶快讲到他的故事的结局，以免他自己话语的苦味渗入我的书中。让\Person{欧诺弥}独占这些恶劣品味的话语，他论及天主的司铎时如此说：\Quote{悭吝的地主、差役和爪牙，到处搜查，不许逃亡者继续躲藏，}以及其他一切他居然不羞于用来描述白发司铎的东西。正如在世俗学问的学校里，为了训练少年伶俐口才，他们会提出命题作文，其中被评论的人物是无名的；同样，\Person{欧诺弥}一上来就攻击所提示的事实，放出谩骂之舌，丝毫不提任何实际恶行，只是堆砌一切陈腐的轻蔑短语和所有可想象到的辱骂言辞：其中还夹杂着不相协调的观念，例如\Quote{附庸风雅的士兵}、\Quote{被咒骂的圣人}、\Quote{因禁食而苍白，因仇恨而凶残}，以及许多这类的粗俗言词；就像一个世俗游行中的狂欢者，当他想将放肆推向极致时，便摘下面具大喊大叫；\Person{欧诺弥}同样毫不掩饰恶意，以铜锣般的喉咙喊出车夫的粗话。然后他揭示自己如此愤怒的原因：\Quote{这些司铎采取一切预防措施，以免许多人}被这些异端者的错误所迷惑；因此他恼怒这些人不能随意留在他们喜欢的地方，而是根据当时\Place{弗里吉亚}总督的命令给他们指定了住处，以便多数人能免受如此邪恶邻居的危害；他对这件事的愤慨爆发成这些话语：\Quote{我们试炼的过度严酷}、\Quote{我们的悲惨苦难}、\Quote{我们对此的高贵忍受}、\Quote{从祖国被流放到\Place{弗里吉亚}}。确实如此：这位\Place{奥利图塞里亚人}很可能为所发生的事感到骄傲，它终结了他所有的家族骄傲，并给他的家族带来了如此耻辱，以至于他那远负盛名的祖父\Person{普里斯库斯}——他从他那里得到那些光辉灿烂、最引人注目的传家宝，\Quote{磨坊、皮革和奴隶的仓库，}以及在\Place{客纳罕}其余遗产——绝不会选择这份如今使他如此愤怒的命运。他辱骂那些执行这次流放的人，这也是意料之中的。我完全理解他的感受。确实，这些不幸的始作俑者——如果现在或过去存在这样的人——理当受到这些人的责难，因为这样一来，他们先前生活的名声被掩盖了，他们丧失了提及并大肆炫耀自己更令人印象深刻的前尘往事的机会：每个人在人生起步时的伟大殊荣；他们从父亲那里继承的职业；每个人意识到的或大或小的贵族标记，甚至在他们变得如此广为人知、如此受重视，以至于皇帝们也以结识他们为荣之前——如他现在在其著作中所夸耀的——并且所有高级政府都为他们而骚动，世界充满了他们的所作所为。\par

% structure: p0028 source=h2 target=subsection reason=relative-heading-rank; base=h2

\subsection{六、关于其异端师尊\Person{艾提乌斯}及其本人的情况，并描述各人的出身与职业。}

确实，这极大地损害了我们这位演说辞作者，或者更确切地说，损害了他的赞助人和人生导师\Person{艾蒂乌斯}；我认为，他的热忱与其说是为了传播谬误，不如说是为了谋生。我这样说并非仅凭臆测，而是从那些与他相熟的人口中听闻的。我曾聆听过先前担任加拉达人的主教\Person{阿塔纳修斯}谈论\Person{艾蒂乌斯}的生活；\Person{阿塔纳修斯}是一个视真理高于一切的人；他还出示了\Place{劳迪塞雅}的\Person{乔治}的书信，以便让许多人能证实他话语的真实性。他告诉我们，起初\Person{艾蒂乌斯}并未试图教导他那怪异的教义，只是在经过一段时间之后，才将这些新奇的东西作为一种谋生的伎俩提出来；他先是逃离了自己所属的葡萄园中的奴役——我不想说是如何逃离的，以免被认为我在恶意记述他的历史——最初成了一个补锅匠，对于这种肮脏的机械手艺非常熟练，坐在一个用山羊毛织成的帐篷下，拿着小锤子和一个小铁砧，以此勉强糊口，生活艰辛。靠修补铜器上的松动之处、焊接漏洞、将锡片敲成碎片、用铅夹住锅腿，能赚到什么像样的收入呢？我们被告知，他在这一行中遇到的某件事促使他改变了生活。他从一个女兵那里收到一件金饰，一条项链或手镯，因受到撞击而断裂，需要修理；但他欺骗了这位可怜的女子，私吞了她的金饰，换了一个同样大小、外表相同、只是表面镀了金的铜饰给她；她一时被蒙骗了，因为他在铜匠手艺以及其他技艺上都足够聪明，能用行业中的伎俩误导顾客；但最终她发现了这卑劣行径，因为铜上的镀金被擦掉了；她家族和民族中的一些士兵被激怒了，于是她起诉了这个窃取她饰物的人。经过这次尝试后，他当然受到了小偷应受的惩罚；然后他离开了这一行，发誓说这不是他的本意，而是生意诱惑他犯下了这偷窃之罪。此后，他成了一个江湖医生的助手，以免完全失去生计；在这种身份下，他攻击了那些较不显赫的家庭和最卑微的人。财富逐渐从他针对某个\Person{亚美尼乌斯}的阴谋中得来，这个\Person{亚美尼乌斯}是外乡人，容易被骗，被说服让他做自己的医生，并经常预付给他钱；于是他认为在别人手下做事有失身份，想自称医生。从此以后，他参加医学会议，与那里好争论的辩手为伍，成了夸夸其谈者之一，总是在天平倾斜时加上自己的砝码，从而成为那些为了党争而购买铜制嗓音的人的重要人物。\par

尽管他因此日子过得很好，但他认为不应继续从事这样的职业；于是他逐渐放弃了医学，就像之前放弃补锅匠一样。天主的敌人\Person{亚流}早已播下了那些邪恶的莠子，结出了\Emph{欧诺米派}作为果实；当时医学院中回响着关于这个问题的争论。因此\Person{艾蒂乌斯}研究了这一争论，并根据他对亚里士多德的记忆设置了一系列三段论，他因甚至在异端之父\Person{亚流}之上、以其猜测的新颖性而臭名昭著；更确切地说，他意识到了\Person{亚流}所提出的一切的后果，并因此获得了善于发现不明显真理的名声；他揭示出，由无中生有的受造物，与从无中引出祂的造物主是\Emph{不相似}的。\par

他用这些命题搔痒那些渴望这些新奇之事的耳朵；于是埃塞俄比亚人\Person{德奥斐洛}了解到了它们。\Person{艾蒂乌斯}在\Person{加卢斯}处理某些事务时已经与这个人有过联系；现在借助他的帮助潜入了宫殿。在\Person{加卢斯}对总督\Person{多米提安}和\Person{蒙提乌斯}犯下罪行之后，所有其他参与者自然都与他一同遭殃；然而这个人却逃脱了，没有被与他们一同受罚。此后，当伟大的\Person{阿塔纳修斯}被皇帝命令逐出亚历山卓教会，而塔尔巴斯忒尼特人\Person{乔治}撕碎他的羊群时，另一次转变发生了，\Person{艾蒂乌斯}成了亚历山卓人，在那些在卡帕多细亚人的餐桌上养肥自己的人中得到了他的一份；因为他没有忘记对\Person{乔治}阿谀奉承。\Person{乔治}实际上自己也来自\Place{客纳罕}，因此对一个同胞感到亲切：事实上，他长期被他的错谬观点所占据，以至于对他十分迷恋，并且容易成为\Person{艾蒂乌斯}的恩赐，只要他愿意。\par

这一切都没有逃过他真诚的崇拜者，我们的\Person{欧诺米乌斯}的注意。后者察觉到他生父——一个优秀的人，只是有了这样一个儿子——过着非常诚实且受人尊敬的生活，但却是辛劳贫困、充满无数艰辛的生活。（他是那种总是弯腰扶犁、为他的小块田地费尽心思的农夫之一；冬天，当他不用干农活时，他为男孩子们精心雕刻字母表，让他们组成音节，用卖这些的钱来谋生。）看到父亲生活中的这一切，他与犁、锄头以及所有父亲的工具告别，打算再也不像那样辛苦劳作；然后他开始学习\Person{普鲁尼库斯}的速记技巧，精通之后，最初我相信，他进入了自己家族中的一个家，以抄写劳役换取食宿；然后，在指导房东的男孩子们时，他产生了成为演说家的雄心。我略过接下来的间隔期，无论是在家乡的生活，还是在\Place{君士坦丁堡}被发现的那些事和那些人。\par

此后，他忙于\Quote{衣裳和钱袋}的事，但看出这一切收效甚微，靠这种工作所能积累的一切都不足以满足他的野心。于是他放弃所有其他行径，专心一意崇拜\Person{阿埃提乌斯}；这或许并非没有计算：他选择的这项投入全部精力的追求，可能有助于他谋生的图谋。事实上，从他要求分享如此深奥的智慧那一刻起，他就不再劳作，也不再纺织；因为他对自己所从事的事确实很精明，并且知道如何赢得人类中较易动情的部分。看到人的本性通常很容易被快乐俘获，而且它在这方面的自然倾向非常强烈，从行为的严峻高处降落到舒适的平坦地带，为了使尽可能多的人改信他那有害的观点，他对自己正在启蒙的人变得非常和蔼可亲；他完全避开了德行的陡峭艰辛之路，因为那并非让人接受他的秘密的说服方式。但若有人有余暇查问他们这秘密教义是什么，那些被欺骗接受这败坏咒诅的人毫不保留地说出什么，以及在神秘入门仪式中，可敬的圣师教给他们什么——洗礼的方式、\Quote{本性的帮助}等等，就让他去问那些不因口出秽语而内疚的人；我们将保持沉默。因为即使我们是控告者，提及这些事也并非无罪，而且我们被教导要在言语和行为上尊重纯洁，不让我们的书页沾染模棱两可的故事，即使我们所说的是事实。\par

但我们提起我们当时所听到的（即，正如\Person{亚里士多德}的恶劣技艺为\Person{阿埃提乌斯}提供了他的不敬，同样，他那些受骗者的单纯也为训练有素的\Person{门徒}及其师傅保证了丰裕的生活），是为了提出一些问题。毕竟，\Person{巴西略}在\Place{黑海}一带，或\Person{尤斯塔修斯}在\Place{亚美尼亚}，给他造成了什么大损害？他的故事中那冗长的离题都回到了这两人身上。他们如何破坏了他的人生目标？他们不是反而滋养了他和他的同伴新获得的名声吗？他们的广泛知名度从何而来，如果不是通过这两人的作用——假设控告者说的是真话？因为，像我们作者所承认的那样，那些本身杰出的人屈尊与那些尚未找到出名途径的人争斗，这自然给了那些要与这些公认的英雄较量的人一个雄心勃勃的起步；从而掩盖了他们卑微的过往。事实上，他们后来的名声恰恰归功于此——这对一个思考的心灵来说确实可憎，它决不会选择把名声建立在恶行之上，但对于这类人物，却是福乐之巅。他们讲到\Place{亚细亚}省有一个人，最卑微、最卑劣，渴望在\Place{以弗所}成名；某种伟大而辉煌的成就完全超出他的能力，他从未想过；但通过想出最严重损害\Place{以弗所}人的办法，他留下的痕迹比那些最伟大行为中的英雄还要深刻；因为在他们的公共建筑中，有一座以其特有的壮丽和昂贵而引人注目；他把这座巨大的建筑烧为平地，当人们事后追查这一恶行的心理原因时，他表现出珍视知名度，并设计了让灾难的规模确保其作者的名字与之共存。这两个人的隐秘动机是同样的渴求名声；唯一的区别是他们造成的损害更大。他们毁坏的并非无生命的建筑，而是教会活生生的建筑，用他们教导的缓慢腐蚀代替了火。但教义问题我将留到适当的时候再谈。\par

% structure: p0035 source=h2 target=subsection reason=relative-heading-rank; base=h2

\subsection{七、\Person{欧诺米乌斯}本人证明祂所做的信德宣认并未受到诘难。}

现在让我们暂且看一看这个人处理的是什么样的真理，他在\WorkTitle{引言}中抱怨说，正是因为他讲真话，才被不信者憎恨；我们很可以用他处理真理的方式（在教义之外）来作为检验他教义本身的试金石。\Quote{\BibleQuote{在最小的事上忠信的，在许多事上也必忠信；在最小的事上不义的，在许多事上也必不义。}}\BibleRef{路 16:10}\newline{} 现在，当他开始写这篇\WorkTitle{为辩护词的辩护}（这是他这部著作新颖惊人的标题和主题）时，他说，我们必须到别处去寻找这一惊人宣告的原因，只能在他那第一部著作的回应者中寻找。那部书题为\WorkTitle{辩护词}；但当我们那位神学大师告诉他，辩护只能来自那些受到某种指控的人，而且如果一个人只是出于自己的倾向写作，他的作品就不是辩护时，他并不否认——否认会过于明显荒谬——辩护需要先有指控；但他宣称，他的\Quote{辩护词}已在对他的审判中使他摆脱了非常严重的指控。这多么虚假，从他自己的话中可见一斑。他抱怨说，\Quote{那些定他罪的人使他遭受了许多沉重的痛苦}；这一点我们可以在他的书中读到。\par

但是，如果他以\Quote{辩词}清除了这些指控，他又怎会如此受苦呢？如果他成功地采用辩词逃脱了这些指控，那么他那可怜的抱怨就是虚伪的伪装；反之，如果他确实如他所说的那样受苦，那么显然，他受苦是因为他并\Emph{没有}通过辩词为自己辩解；因为每一个辩词，作为辩词，都必须达成这一目的，即防止投票权被任何虚假陈述所误导。当然，他现在不会试图说，在审判时他提出了他的辩词，但因无法赢得陪审团而败诉给原告。因为他在审判时完全没有提到\Quote{关于提出他的辩词}；他也不太可能这样做，考虑到他在书中明确声明他拒绝与那些心怀恶意和敌意的法官有任何瓜葛。\Quote{我们承认，}他说，\Quote{我们被缺席判罪：那里有一个邪恶之人的有偏见的陪审团，本应是陪审团坐席的地方。}他在这里非常费力，并且我认为他的注意力被他的论证转移了，否则他会注意到他在他的句子后面添加了一个巧妙的语法错误。他试图用他的短语\Quote{有偏见的陪审团}显得仿佛是阿提喀风格；但语言正确的人使用这些词，正如那些熟悉法庭词汇的人所知，与我们的新阿提喀主义者截然不同。\par

稍后他又补充道：\Quote{如果他认为，因为我拒绝与实际上是我控告者的陪审团打交道，他就可以辩倒我的辩词，那么他必定对自己头脑简单视而不见。}那么，我们这位尖刻的朋友是在何时、在谁面前提出辩词的呢？他因陪审团是\Quote{敌人}而拒绝他们，并且正如他自己坚持的，他没有对任何审判说一个字。请看这位真理的勇士，如何一点一点地转向虚假一边，在言辞上尊崇真理，在行动上却与之斗争。但可笑的是，即使在自己谎言的支持上，他也是多么软弱。同一个人怎么可能既\Quote{在他被提起的审判中以辩词洗清了自己}，又\Quote{因为法庭掌握在敌人手中而谨慎地保持沉默}？不，他在其辩词前言中的措辞本身清楚地表明，根本没有对他开启任何法庭。因为他没有将前言指向任何具体的陪审团，而是指向某些未指明的人，无论是当时活着的，还是后来将要来到世上的；我承认，对于这样的听众，需要一份非常有力的辩词，但不是像他实际写的那样——那份辩词还需要另一份来支撑——而是一份广泛易懂的辩词，能够证明这一点，即：他写这辩词时并不处于通常的理智状态，他敲响集合钟召集那些从未到来、也许从未存在过的人，在一个想象中的法庭面前发表辩词，并向一个看不见的陪审团恳求，不要让人数决定真理与谬误，也不要将胜利归于单纯的数量。实在的，他应当向那些还在他父亲腰中的陪审员发表那种辩词，并向他们解释，他何以认为采纳与普世信仰相矛盾的意见是正当的，并且宁愿相信自己错误的幻想，也不相信那些在全世界光荣基督之名的人。\par

请让他再写一份辩词，除了这第二份之外；因为这一份不是对他所受错误指责的纠正，而是对那些指责真理的证明。每个人都知道，一份适当的辩词旨在反驳指控；因此，一个被指控偷窃、杀人或其他罪行的人，要么完全否认事实，要么将责任转嫁给另一方，或者，如果两者都不可能，则向将要判决的人恳求慈悲或怜悯。但在他的书中，他既不否认指控，也不将其转嫁给别人，也不求饶，也不承诺将来改正；反而通过一种异常费力的论证证实了对他的指控。他自己承认，这一指控实际上是一份亵渎的控告，而且并没有将这种亵渎的性质模糊化，而是宣告了具体种类；而他的辩词证明这种亵渎是一种积极的义务，不但没有消除指控，反而加强了它。如果我们的信仰教义尚有任何模糊之处，尝试新事物可能不那么危险；但我们的神学导师的教导已牢固地扎根在信徒的灵魂中；因此，一个大声喊出与大家同样已经下定决心的事情相矛盾的人，是在为自己所受的指控辩护，还是在激怒听众，使他的控告者更加愤怒？我倾向于后者。所以，如果正如我们的作者所说，既有他辩词的听众，也有对他攻击信仰的控告者，那么让他告诉我们，那些控告者现在如何可能妥协此事，或者陪审团必须作出何种裁决，既然他的罪行已由他自己的\Quote{辩词}证明了。\par

% structure: p0040 source=h2 target=subsection reason=relative-heading-rank; base=h2

\subsection{八、事实表明他用来攻击巴西略的辱骂之词更适合他自己。}

但这些话是顺便提及的，源自我们没有紧扣自己的论证。我们当初需要查问的不是他应当如何为自己辩解，而是他是否曾经作过任何辩解。但现在让我们回到先前的立场，即：他被自己的陈述所定罪。这位憎恨谎言的人首先告诉我们，他被定罪是因为指派给他的陪审团藐视法律，他被驱赶上天下海，饱受烈日和尘土之苦。然后，为掩盖他的谎言，他用一个钉子钉出另一个钉子，如俗语所说的，用一个谎言纠正另一个谎言。既然每个人都像他一样清楚，他在法庭上一言未发，他便宣称他曾请求免于进入敌对的法庭，并被缺席审判定罪。还有比这更明显的自相矛盾、既违背真理又违背自己的事例吗？当别人追问他的书名时，他便将他的审判作为这一\Quote{辩解}的逼迫原因；但当别人追问他未曾向陪审团说过一句话的事实时，他否认有过任何审判，并说他是拒绝了那样的陪审团。看，这位勇敢的真理捍卫者如何英勇地与谎言作战！然后他竟敢称我们伟大的\Person{巴西略}是\Quote{恶毒的流氓和说谎者}，此外还有\Quote{大胆无知的新贵}、\Quote{不深的神学家}，并在他的谩骂清单上加上\Quote{彻底疯狂}，在他的书页中散播无数这样的词语，仿佛他以为他那恶毒的咒骂能压倒人类的共同见证——人们尊崇那个伟大的名字如同尊崇古代的圣人之一。事实上，他认为，即使没有别人，他也能用诽谤触及那个从未被诽谤触及的人；但太阳在天空中并非低到任何人都可以用石头或其他投掷物击中它；它们只会回击投掷者，而预期的目标则远非其所能及。此外，如果有人指责太阳缺乏光明，他的嘲笑并不能减弱阳光的亮度；太阳依然是太阳，挑剔者只会证明自己视觉器官的软弱；而且，如果他试图仿效这种\Quote{辩解}的方式，说服所有他遇见并愿意听他讲话的人不要屈从于关于太阳的普遍看法，也不要因\Quote{将胜利归于单纯的数量}而更重视所有人的经验而非一个人的猜测，那么他的胡言乱语对能用眼看见的人将是徒劳。\par

那么，让某人劝告\Person{欧诺弥}勒住他的舌头，不要放纵如此狂野的言论，不要在对一个受尊敬的名字的傲慢辱骂中踢刺棒；而是让对\Person{巴西略}的单纯回忆充满他的灵魂，使他敬畏和恐惧。他通过这种无度的粗俗话能得到什么？被辱骂的对象将保留所有其生命、言语以及文明世界的普遍评价所宣告他拥有的品格。着手辱骂的人显露出自己的性情：因为邪恶而不能说出好话，只能\BibleQuote{从心里所充满的说话}，并从那邪恶的宝库里取出恶来。现在，他的那些表达不过是与事实完全脱节的辱骂，这可以从他自己的著作中得到证明。\par

% structure: p0043 source=h2 target=subsection reason=relative-heading-rank; base=h2

\subsection{九、在指控\Person{巴西略}于\Quote{审判}时期没有为自己的信仰辩护时，他使自己面临同样的指控。}

他暗示了某个地点，那里发生了这场关于异端的审判；但他没有给出确切指示在哪里，读者不得不在黑暗中猜测。他告诉我们，一个从各地挑选出来的代表会议被召集到那里；他在这里处于最佳状态，用一些有力的笔触将他认为发生的事件的准备情况呈现在我们眼前。然后，他说，一场他必须为保命而奔波的审判被交给了某些仲裁者，我们的导师和师傅当时在场，向他们作了指示；由于所有的投票权就这样被争取到敌人一边，他放弃了阵地，逃离了那里，四处寻找某个炉灶和家园；在这生动的描绘中，他在指责我们英雄的怯懦方面很出色；任何喜欢的人都可以通过阅读他所写的内容看到。但我不能停下来在这里给出他话语的苦胆的例子；我必须继续前进到那个我提及这一切的目的所在。\par

那么，那个未命名的地点，即对他的教导进行审查的地方，在哪里？那次为审判而聚集最优秀之人的场合是什么？那些匆忙跨越陆地和海洋来分担这些劳苦的人是谁？这个\Quote{悬念投票结果的期待世界}是什么？谁是\Quote{审判的安排者}？然而，我们可以认为他编造了这一切以夸大其故事的重要性，就像学校里的男孩们在\Emph{他们的}这种虚构的对话中容易做的那样；并请他只告诉我们，那位我们的师傅躲避与之相遇的\Quote{可怕战士}是谁。如果这也是虚构的，那就让他再次成为赢家，并从他的空话中获利。我们将保持沉默：在与阴影的无用战斗中，真正的胜利是拒绝在\Emph{那一点上}取胜。但如果他指的是\Place{君士坦丁堡}的事件，指的是那里的集会，并出于文学上的愤怒而对那里上演的悲剧感到激动，并将自己意味为那位伟大而令人敬畏的运动员，那么我们将展示理由，说明为何尽管我们当时在场，我们并未投入战斗。\par

如今，这位以懦弱斥责那位英雄的人，请告诉我们：\Emph{他}是否深入了战场的中心？\Emph{他}是否为自己的正统信仰说出了哪怕一个字？\Emph{他}是否做出了任何有力的结辩？\Emph{他}是否胜利地与敌人交锋？他不能告诉我们这些，否则他显然自相矛盾，因为他承认由于自己的缺席，他受到了不利的裁决。如果在审判当时发言是应当的（因为他在书中为我们制定这条规定），那么他为何当时因缺席而被判罪？另一方面，如果他在这样的法官面前保持沉默是做得好，那么他如此随意地称赞自己，却责备我们在这样的时刻沉默，这是何等荒谬的不公！自从审判以来发表了两篇论文，他声称他的辩护虽然写作于很久之后，却及时，却指责那回应他辩护的文章太过迟延！他本当在\Person{巴西略}的预期反驳尚未实际提出之前就加以指责，但在他其他的抱怨中并未发现这一点。既然他知道\Person{巴西略}在审判过后将会写什么，他为何不当时当场就指责它？事实上，从他自己的供认可以清楚看出，他在审判本身中从未做过那篇辩护。我再重复他的话：——\Quote{我们承认我们因缺席而被判罪；}他补充了原因：\Quote{心怀恶意的人被选为陪审员}，或者更确切地说，用他自己的话：\Quote{在应当有陪审团的地方，有一个串通好的小组。}而另一方面，从书中另一处可以清楚看出，他证明他的辩护是\Quote{在适当的时候}做出的。经文如下：——\Quote{我被敦促在适当的时候以适当的方式做出这篇辩护，不是出于虚假的理由，而是为了那些为我担保的人被迫如此，这从事实以及这个人的话中都很清楚。}他巧妙地将他的话转来转去以应对各种可能的反对；但对这一点他会怎么说？\Quote{在审判期间保持沉默是不对的。}那么，\Person{欧诺弥}在同一个审判中为何无言？而他的辩护，在审判之后才出现，为何就及时了？如果及时，那么\Person{巴西略}与他的论战为何不及时？\par

然而，那位教父的话尤其真实：\Person{欧诺弥}假装道歉，实际上却给了他希望给予其教导的支持；那位效法\Person{丕乃哈斯}热忱的真正模仿者，用圣言之剑消灭每一个精神上的淫乱者，在\WorkTitle{回答他的亵渎}中刺出一剑，既能治愈灵魂，又能摧毁异端。如果他抵抗那一击，因背道而麻木的灵魂不愿接受治疗，那么责任在于选择恶的一方，如同外邦人的谚语所说。以上是关于\Person{欧诺弥}对待真理和我们自己的态度：如今，古时的律法允许对首先伤害者以同样的回报，这可能促使我们用一场反斥责的暴雨回击他，而且由于他很容易成为这样的目标，我们应当慷慨地回击，以胜过他所施加的痛苦：因为如果他对一个没有提供诽谤机会的人如此无理谩骂，那么对那些讽刺过那位圣洁生活的人，我们还能期待多少这样的谩骂呢？但我们从那位真理的学者那里一开始就被教导要成为福音的学者，因此我们不会\BibleQuote{以眼还眼，以牙还牙}；我们清楚知道，所有发生的恶都可以通过其对立面被消灭，如果没有一句好话或一件好事来打断邪恶之流的连续性，任何坏话或坏事都不会发展成如此绝望的邪恶。因此，无礼和辱骂的循环因忍耐而停止重复；而如果无礼遇到无礼，辱骂遇到辱骂，你就只会滋养这个怪物般的恶习，并大大增加它。\par

% structure: p0048 source=h2 target=subsection reason=relative-heading-rank; base=h2

\subsection{十、事实表明，他所有的侮辱性称呼都是虚假的。}

因此，我略过其余一切，只视之为傲慢的嘲弄与辱骂，急忙转向其教义问题。若有人以为我拒绝辱骂只是因为我无力以其人之道还治其人之身，那么让这样的人在自己身上思考：一般而言，人何等易于倾向邪恶，何等机械地滑入罪恶，无需任何练习。成为恶的能力在于意志；一个意愿的行为往往足以成就一种完整的邪恶；而这种操作的轻易性在口舌之罪中尤为致命。其他类别的罪恶需要时间、机会和合作才能犯下；但言语的倾向随时可以犯罪。现在在我们手中的\Person{尤诺米乌斯}的论著足以证明这一点；仔细思考的人会察觉在措辞方面滑入罪恶的迅速性：模仿这些是最容易的事，即使一个人完全不习惯于惯常的诋毁。何必费力把我们意图的侮辱编造成名目，当我们可以用他自己词句来攻击这个毁谤者？事实上，他在其著作的这一部分，串起了各种谎言和恶言，全都由他自身中的模型塑造而成；每一种过度都能在这些文字中找到。他写道：\Quote{狡猾}、\Quote{争论}、\Quote{真理之敌}、\Quote{高傲}、\Quote{江湖骗子}、\Quote{对抗普遍意见与传统}、\Quote{无视事实的反驳}、\Quote{不顾法律的恐怖与人的指责}、\Quote{不能区分对真理的热忱与单纯推理技巧}；他补充道：\Quote{缺乏敬畏}、\Quote{急于辱骂}，然后又是\Quote{喧嚷}、\Quote{充满矛盾的怀疑}、\Quote{结合不可调和的论点}、\Quote{对抗自己的言论}、\Quote{肯定矛盾}；然后，尽管他渴望说尽对方一切坏话，却无法找到其他新奇辱骂来发泄他的恶毒，常常在穷尽一切后重复相同的措辞，第三、第四次甚至更多次回到他曾经说过的话；在这个言辞的跑马场里，他反复驱车上下，一遍又一遍地跑同一圈傲慢辱骂的跑道；以至于最后，对这种无耻表演的愤怒也因极度厌倦而消退。这些低俗不雅的街头男孩的嘲笑确实引起厌恶而非愤怒；它们并不比某个醉酒老妇含混不清的咕哝好多少。\par

那么，我们必须对此逐条深究，费力驳斥他所有的辱骂，证明\Person{巴西略}并非他那想象中的怪物吗？如果我们这样做，满足于证明他身上没有任何卑劣或犯罪之处，那我们就似乎是在侮辱一位对他那世代而言\Quote{明亮的特别之星}。但我记得，他用那神圣的声音，引用先知关于他的话\BibleRef{耶 3:3}，将他比作一个无耻的妇人，将自己的责难投向贞洁者。他的这些推理宣告谁是真理之敌、与普遍意见为敌？是谁恳求他书籍的读者\Quote{不要着眼于宣称信仰的人数，或单纯的传统，也不要让判断被偏见所左右，以至于认为仅被怀疑为更强的一方是可信的？}同一个人能写出这样的文字，然后做出那些指控，策划让读者追随他自己的新奇之说，而就在此时他却指责他人反对普遍信念？至于\Quote{厚颜无耻地无视事实的反驳和人们的指责}，我让读者判断这适用于谁：是适用于那个通过最谨慎的自我克制，使自己以及周围人的生活以节制、安静和完全纯洁为准则的人，还是适用于那个建议：当自然乐意通过身体的欲望前进时，不应打扰她，不应阻挠放纵，也不应在修练生活方面如此讲究；而一个自我选择的信德足以使人达至成全？如果他否认这是他的教导，我和任何正直的人都会为他在这种否认中说了实话而欣喜。但他的真正追随者不会允许他做出这样的否认，否则他们的主导原则就会丧失，那些由此原因接受他教义者的平台也会崩塌。至于对他人的指责无耻地漠不关心，你可以看看他的青年或此后生活，你会发现他在两者中都易受这种责备。这两个人的生平，无论是青年时期还是成年时期，都讲述了一个大相径庭的故事。\par

让我们的言论撰写者，在回忆他在故土及后来在\Place{君士坦丁堡}的青年行为时，从那些能告诉他他们所知的、他所毁谤之人的事情的人那里听取。但如果有人想询问他们随后从事的职业，让这样的人告诉我们，他认为两人中哪一个更值得拥有如此高的声誉：是那个在成为司铎之前就慷慨地将祖产用于穷人的人，尤其是在饥荒时期——那时他已是教会的治理者，虽然仍是长老职位的司铎——之后甚至没有积攒剩余之物，以便他也能作出宗徒的夸耀\Quote{我们也没有白吃人的饭}\BibleRef{得后 3:8}；还是另一个人，将捍卫一个教义主张视为收入来源，潜入人家，不因呆在家中而掩盖他那令人厌恶的疾患，也不考虑健康人对此自然产生的厌恶，尽管按照古律，他属于那些因明显疾病的传染而被逐出有人居住的营地的人。\par

\Person{巴西略}被此人称为\Quote{鲁莽的}和\Quote{傲慢的}，并在这两种角色中称为\Quote{说谎者}；此人声称要\Quote{以忍耐和温柔教导那些持相反意见的人}；因为他谈到巴西略时就摆出这种姿态，只要有机会，他就不放过任何夸张的刻薄言辞。那么，他凭什么指控巴西略鲁莽和傲慢呢？因为\Quote{他称我为加拉太人，而我本是卡帕多细雅人}；然后，因为他把一个住在\Place{科尔尼亚斯皮内}这样偏僻角落边界上的人称为加拉太人，而实际上他是奥尔蒂塞里人；当然，这是假设已经证明他说过这话。我在我的抄本中没有找到；但就算如此。就为此，他被称为\Quote{鲁莽的}、\Quote{傲慢的}，以及一切恶名。但有智慧的人深知，挑剔者的小指控反而为被指控者的正义提供了强有力的论据；否则，如果他急于指控，就不会放过大错，而把恶意用在微小的事上。在后者上，他无疑很在行，夸大罪过的严重性，对谎言做出庄严的反思，并认为无论大事还是微不足道的小事，谎言同样可憎。就像他那异端的祖师们，经师和法利塞人，他知道怎样仔细滤出蠓虫，却一口吞下满载邪恶的驼峰骆驼。但对他这样说也并非不妥：\Quote{不要在我们的体系中制定这样的规则；不要再命令我们，认为衡量谎言之罪可以依据情况的轻微或重大而不予考虑。}\Person{保禄}说了谎，并按照犹太人的方式洁净自己，以满足他巧妙地欺骗的人的需要，他所犯的罪不同于\Person{犹达斯}，后者为了背叛的需要而装出和蔼可亲的表情。\Person{若瑟}出于爱弟兄而用谎言欺骗了他们；而且那还是指着法老的性命起誓\BibleRef{创 42:15}时说的；但他的弟兄们确实向他撒了谎，因嫉妒而图谋害死他，后又把他卖为奴隶。还有许多类似的例子：\Person{撒辣}撒谎，因为她因发笑而羞愧；蛇撒谎，引诱人类悖逆并变为神性的存在。谎言因其动机而大不相同。因此，我们接受圣神藉先知所说关于人的那句普遍陈述：\Quote{每个人都是说谎者。}而这位天主的人也未能完全避免谎言，因为他碰巧因疏忽或不知其真名而将一个地方说成了邻近地区的名称。但\Person{欧诺弥}也说了谎，是什么呢？正是对真理本身的错误陈述。他断言永远存在的那一位曾经不存在；他证明真正是子的一位却被称为子是假的；他将造物主定义为受造物和作品；他将世界的主称为仆人，并将本质上统治的存在与臣服的存在并列。难道谎言之间的差别如此微不足道，以至于可以认为无论谎言是以这种方式还是那种方式显现都无关紧要吗？\par

% structure: p0053 source=h2 target=subsection reason=relative-heading-rank; base=h2

\subsection{十一、他用来证明我们承认他曾受审、且他的信仰宣认并非无可指责的诡辩是软弱的。}

他指责别人诡辩；请看他自己多么小心，要让他的证明成为真正的证明。我们的导师在他写给此人的书中说，当我们的案件败诉时，\Person{欧诺弥}赢得了\Place{基齐库斯}作为他亵渎的奖赏。那么，这位诡辩的揭露者做了什么呢？他立刻抓住\Emph{奖赏}这个词，声称我们这一方承认他做了辩护，并因此得胜，通过这些努力获得了胜利的奖品；然后他把他的论证构成三段论，他认为是由无可辩驳的命题组成的。但我们将逐字引用他所写的话：\Quote{如果奖品是对胜利的承认和冠冕，而审判意味着胜利，并且，作为与之不可分离的，还有控告，那么，那在论证中承认奖品的人必然允许曾有过辩护。}那么我们的回答是什么呢？我们不否认他以最充沛的精力进行了这场不敬神的下流战斗，并且在这些对抗真理的汗流浃背的努力中远远超过了同伙；但我们不允许说他战胜了对手；而只是说，与那些和他一样通过异端奔向错误的人相比，他在谎言数量上名列前茅，因此赢得了\Place{基齐库斯}作为邪恶高成就的回报，击败了所有为同一奖赏而与真理作战的人；并且为这场亵渎的胜利，他的名声被响亮地传扬，当时他同派的裁判们为他选择了\Place{基齐库斯}作为他放纵的赏报。这就是我们的意见陈述，也是我们所承认的；我们现在主张\Place{基齐库斯}是异端的奖赏，而非辩护的成功结果，这正表明了这一点。这难道不像他自己那堆幼稚的诡辩吗，以至于他可以借此希望有理由证明他受审和辩护的事实？他的方法就像一个在酒宴中的人，他比其他人干掉了更多的烈酒，然后向他的酒友索要赌注，并试图以此胜利来证明他在法庭上打赢了官司。那个人可能会切出同样的逻辑：\Quote{如果奖品是对胜利的承认和冠冕，而法庭审判意味着胜利，并且，作为与之不可分离的，还有控告，那么我赢得了我的诉讼，因为我在这场酒宴中以我的酒量被加冕。}\par

肯定有人会对这样的自夸者说，法庭审判与酒宴竞赛截然不同，靠酒杯获胜的人并不会因此比他的诉讼对手更有优势，尽管他得到一顶漂亮的花冠。因此，那个在亵渎的辩护中胜过同辈的人，也无法证明他赢得了裁决。我们这方作证说他是亵渎之首，这并不能为他虚构的\Quote{辩护}提供理由。如果他确实在法庭上说了那番话，并且以此战胜了对手，因此得到\Place{基齐库斯}作为荣誉，那么他或许有理由用我们自己的话来对付我们；但是，既然他在书中不断抗议说他屈从于投票者的意见，默默接受他们施加的惩罚，甚至没有等待这个不利的判决，他又何必自欺欺人，把这个\Emph{奖}字当作成功辩护的证据呢？我们这位优秀的朋友不明白这个\Emph{奖}字的力量；\Place{基齐库斯}被交给他，作为他极端不虔敬的美德的奖赏；既然他愿意接受这样的赏，并视之为胜利者的报酬，那就让他也接受那胜利所隐含的东西，即亵渎罪责的最大份额。如果他坚持用我们自己的话来对付我们，他就必须同时接受这两者，否则就都别接受。\par

% structure: p0056 source=h2 target=subsection reason=relative-heading-rank; base=h2

\subsection{十二、他对怯懦的指控毫无根据：因为\Person{巴西略}在皇帝及其总督面前展现了最大的勇气。}

他这样对待我们的话；在他其他狂妄的陈述中，能看出有一丝真理吗？在这些话中，他称他是\Quote{怯懦的}、\Quote{无精神的}、\Quote{逃避更艰苦劳动的}，用尽了这类词汇，并煞费周章地详细列出这怯懦的每一个症状：那\Quote{退隐的小屋，紧闭的门，对闯入者的焦虑恐惧，声音，表情，泄露秘密的脸色变化}，以及诸如此类的一切，由此可见恐惧之情。如果他除了这个谎言之外未被发现任何其他谎言，仅此一点就足以揭示他的倾向。因为谁不知道，当\Person{瓦林斯}皇帝被煽动起来反对主的教会时，我们那位伟大的斗士如何以崇高的精神超越了那些压倒性的环境和敌人的恐怖，展现出一种凌驾于一切恫吓手段之上的心灵？东方居民以及我们文明世界最遥远地区的人，谁没有听说过他为真理而与皇座本身的战斗呢？谁看见他的对手而不感到惊惶？因为他的对手并非凡庸之辈，仅拥有在诡辩把戏中获胜的能力，在那里胜利无荣耀，失败也无害；但他拥有使整个罗马政府服从他意志的权力；而且，除了这种帝国的傲慢，他对我们的信仰怀有偏见，这偏见是由\Place{日尔曼尼基亚}的\Person{欧多克修}狡猾地灌输给他的，\Person{欧多克修}已将他争取到自己一边；他在当时所有掌权者中找到了实现其计划的盟友，有些人已经因精神上的共鸣而倾向他们，而其他人（他们是大多数）则出于恐惧而准备纵容皇帝的喜好，并且看到对坚持信仰的人所施行的严厉手段，便大肆炫耀他们对皇帝的效忠。那是流放、没收、放逐、罚款威胁、生命危险、逮捕、监禁、鞭打的时期；没有什么太可怕而不能对那些不愿屈服于皇帝这突发奇想的人施行的；对于信徒来说，在天主殿中被抓住，比被发现在最可憎的罪行中更糟糕。\par

但是，详细叙述那段历史会过于冗长，需要单独处理；此外，既然那个悲伤时期的苦难众所周知，仔细用文字陈述它们对我们目前的目的毫无裨益。这种尝试的第二个弊端是，在叙述那段忧郁历史的细节时，我们不得不提及自己；如果我们在这场为信仰而进行的斗争中做过什么值得一提的光荣事迹，智慧命令我们留给他人去讲述。\BibleQuote{让另一个人赞美你，而不是你自己的口；} \BibleRef{箴 27:2}而我们那位无所不知的朋友恰恰没有意识到这一点，竟将书中大半篇幅用于自我夸耀。\par

那么，略过这一切细节，我只谨慎陈述我们导师的成就。他必须对抗的敌人不亚于皇帝本人；那个敌人的副手是政府中仅次于他的人；执行其意志的助手是朝廷。让我们也考虑时间点，以检验和说明我们高贵斗士的刚毅。那是什么时候？皇帝正从\Place{君士坦丁堡}向东行进，因近来对蛮族的胜利而得意洋洋，不愿忍受任何对其意志的阻碍；而他的总督调整路线，推迟一切必要的国政管理，只要信仰的追随者还有家可归，直到每一个人，无论在哪里，都被驱逐，而他亲自挑选的人被引入以亵渎我们神圣的圣统。然后\Place{普罗蓬提斯}的势力像乌云般愤怒地冲向教会；\Place{比提尼亚}完全被蹂躏；\Place{加拉太}很快被他们的洪流冲走；所有中间地区都随他们成功；现在我们的羊群下一个就要受攻击。这时我们伟大的\Person{巴西略}表现如何，这个\Person{欧诺弥}所说的\Quote{畏缩的懦夫}，\Quote{逃避危险，依赖隐蔽的小屋保全自己}？他是否在这邪恶的进攻面前畏缩？他是否让先前牺牲者的遭遇建议他确保自身安全？他是否听从建议对这场邪恶的洪流稍作退让，以免公然阻挡那些在屠杀中已老练的人的道路？相反，我们发现所有语言、所有思想和言语的高度，都不足以表达他的真相。无人能描述他对危险的蔑视，从而在读者眼前展现这场新战斗，可以公正地说，这不是人与人之间的战斗，而是一方面基督徒的坚定和勇气，另一方面是沾满鲜血的权力之间的战斗。\par

总督不断援引皇帝的命令，使一个本身就因巨大力量而令人恐惧的权力，因其报复的无情残忍而更加可怕。在\Place{比提尼亚}上演悲剧后，\Place{加拉太}以其特有的反复无常不战而降，他认为我们的地区会轻易落入他的计划。残忍的行为之前先有用言语提出，混合威胁和承诺：服从则给予皇家恩惠和教会权力，抵抗则是残忍灵魂在获得权力后所能想出的一切。这就是敌人。\par

我们的斗士非但没有被他所见所闻吓倒，反而像医生或谨慎的顾问一样，被请来纠正错误，吩咐他们为鲁莽悔改，并停止在主的仆人中进行杀戮；\Quote{他们的计划，}他说，\Quote{不可能成功，因为这些人只关心基督的帝国和不朽的权力；尽管他们想虐待他，却找不到任何言语或行为能使基督徒痛苦；充公不会触及他，因为他唯一的财产是信德；流放对在各地行走时怀着同样情感的人没有恐惧，并因寄居短暂而视每座城市为异乡，却又因所有人类都与自己同样受束缚而视为家乡；忍受殴打、折磨或死亡，如果是为了真理，连妇女也不畏惧，而对于每个基督徒来说，为此希望而受最苦的苦是至高的福乐，他们只遗憾本性只允许一次死亡，无法想出在这场为真理的战斗中多次死去的方法。}\par

当他如此面对他们的威胁，并看穿那令人敬畏的权力，仿佛一切都不存在时，他们的恼怒，就像舞台上快速变化的场景，一个面具接着另一个面具戴上，转而以威胁变成谄媚；那个之前精神如此坚定可怕的人，使用了最温和顺从的言辞；\Quote{请不要，我求你，认为我们伟大的皇帝与你们的人共融是一件小事，相反，请愿意也被称为他的导师；也不要阻挠他的愿望；他渴望这和平，只要书写的信经中删去一个词，即\OriginalTerm{同质}{Homoousios}这个词。}我们的导师回答说，皇帝成为教会的一员极为重要；也就是说，他拯救自己的灵魂，不是作为皇帝，而是作为单纯的人；但减少或增加信德，\Person{巴西略}连想都没想过，他甚至不会改变书写文字的顺序。这就是这个\Quote{畏缩的懦夫，听到门吱声就发抖}的人对这位伟大统治者说的话，他用行动证实了这些话；因为他亲自挡住了冲向教会的皇权毁灭洪流，并使之转向；他独自一人抵挡了这次攻击，如同海中巨大不动的岩石，击碎那可怕冲击的巨浪。\par

他的角力并未就此止息；皇帝本人接着发动攻击，因第一次尝试未能如愿达成一切而恼怒。于是，正如亚述人藉着厨子\Person{拿步匝辣当}摧毁了耶路撒冷以色列人的圣殿，我们这位君主也将他的事务托付给一位\Person{德默斯德内斯}——他的厨房总管、厨师长——认为他比别人更为进取，指望借此完全实现其计划。此人搅动汤锅，一位来自\Place{伊里利孔}的亵渎者手持信件，以此为目的召集权贵，而\Person{摩德斯督斯}则比先前更炽热地煽动激情，于是众人都随从皇帝的怒气，将其狂怒据为己有，屈从于权威的脾性；另一方面，众人对未来可能发生之事都感到希望沉落。那位总督再次登场；比先前更可怕的恐吓开始了；他们的威胁抛出；怒气升到更高点；审判的悲剧排场再现：传唤员、随从、扈从、帷幕后的法庭——这些事物自然会使即使完全准备好的心灵也感到畏惧；我们再次看到天主的斗士在这场战斗中超越了他先前的荣光。若你要求证据，请看事实。何处有教会，那灾难不曾到达？哪个民族未被这些异端命令触及？哪间教会中的显赫人物未被逐出他劳作的场所？哪一民族逃脱了他们的欺凌？它到达了整个叙利亚、直到边界的美索不达米亚、腓尼基、巴勒斯坦、阿拉伯、埃及、直到文明世界边界的利比亚部落；以及所有近处：本都、基利家、吕基亚、吕底亚、彼西底、旁非利亚、卡里亚、赫勒斯滂、直到普罗庞提斯本身的岛屿；色雷斯沿海，直到色雷斯延伸之处，以及直到多瑙河的邻国。这些国家中有哪一个保持了先前的面貌，除非它们已为邪恶所占据？唯有\Place{卡帕多细亚}的人民没有感受到教会这些苦难，因为我们强大的斗士拯救了他们于试炼之中。\par

这就是我们这位\Quote{懦弱}导师的成就；这就是那位\Quote{逃避一切更艰苦劳作}者的成功。这确实不是那\Quote{在贫穷老妇中赢得名声，并操练欺骗那自然落入一切陷阱的性别}者，以及\Quote{认为被罪犯和堕落者钦佩是大事}者的成就；而是那以其行为证明了他灵魂的刚毅和坚定不移、高贵精神之人。他的成功带来了整个国家的救恩、我们教会的和平、为有德者树立的每种卓越的典范、敌人的覆灭、信德的维护、软弱弟兄的坚固、热心者的鼓励——一切被认为属于得胜一方的事物；在纪念中，除了这些事件，听觉和视觉在实现的事实中不再结合；因为在这里，用言语叙述他的高贵行为，和以事实显示我们言语的确认，以及以彼此相互印证——通过眼前所见记录，通过所说之语确认事实——是同一回事。\par

% structure: p0065 source=h2 target=subsection reason=relative-heading-rank; base=h2

\subsection{十三、教义教导摘要。逐条反对意见。}

但我们的论述不知怎的已大大偏离了目标；它不得不转向并逐一面对这个诽谤者的侮辱。对\Person{欧诺弥乌斯}来说，讨论停留在这些点上确实不无小利，而对他冒犯人的控告拖延我们接近他更重罪行的时机。但辱骂一个因谋杀而受审的人说话轻率是无益的（因为后者的证明足以使死亡判决通过，即使说话轻率没有被同时证明）；同样，看来最好只将他亵渎的言论付诸检验，而把他的侮辱搁置一边。当他在最重要问题上的可憎已被揭露时，他的其他过失就无需细究而被潜在地证明了。那么，在他的所有论证中，首当其冲的就是对信德定义的亵渎——无论是他以前的著作还是我们现在正在批评的这部——以及他竭力摧毁、取消并完全颠覆关于天主的独生子及圣神的一切虔敬概念的顽强努力。那么，为了表明他针对这些真理教义的论证是多么虚假和矛盾，我将首先逐字引用他的整个陈述，然后重新开始分别检视每一部分。\Quote{我们教义的全部概况总括如下：存在一个至高无上且绝对的存在，以及另一个因第一个而存在、但在它之后且先于所有其他存在的存在，以及第三个不与这两个中任何一个同等的存在，但次于一个（作为其原因），次于另一个（作为产生它的能力）；当然，在这个概况中必须包括每个存在之后的能力以及与这些能力相关的名称。再者，由于每个存在是绝对单一的，并且事实上和思想上是单一的，其能力以其作品为界，其作品与其能力相称，必然地，这些存在之后的能力相对更大或更小，有些是较高层次，有些是较低层次；一言以蔽之，它们的差异等于其作品之间的差异：事实上，说同一种能力产生了天使或星辰、以及诸天或人，是不合法的；但一个虔敬的心灵会得出结论，正如一些作品优于和更尊贵于另一些一样，一种能力超越另一种，因为相同的能力产生相同的作品，而作品的差异表明能力的差异。既然这些情况如此，并且它们之间的关系保持着不间断的联系，看来那些按照与主题相称的顺序进行研究、并且不坚持将所有事物混杂混淆在一起的人，当讨论涉及存在时，应通过首要的能力和附属于存在的能力来证明正在论证的内容，并解决争议点，而当能力成为问题时，又通过存在来解释，但仍然认为从第一个到第二个的过渡是两者中更合适和各方面更有效的方式。}\par

这就是他有系统的亵渎！愿真天主——真天主的圣子——藉圣神的引导，将我们的讨论引向真理！我们将逐一重复他的陈述。他断言\Quote{他的教义的全部概况总括在至高无上且绝对的存在，以及另一个因第一个而存在、但在它之后且先于所有其他存在的存在，以及第三个不与这两个中任何一个同等的存在，但次于一个（作为其原因），次于另一个（作为能力）。}那么，这个陈述中不公处理的第一点值得注意：在声称阐述信德的奥迹时，他仿佛纠正福音中的说法，拒绝使用我们的主在完善我们的信德时传达这奥迹所用的词语：他压制了\Quote{圣父、圣子和圣神}的名字，而代之以\Quote{至高无上且绝对的存在}代替父，\Quote{另一个通过它而存在、但在它之后}代替子，\Quote{第三个与这两个中任何一个都不等同}代替圣神。然而，如果这些是更合适的名称，真理本身不会无法发现它们，那些相继受托宣讲奥迹的人——无论是从起初就是圣言的目击见证者和仆役，还是作为他们的继承者以福音教义充满普世、并在此后各个时期在公会议中定义了关于教义所引起的歧义的人——也不会无法发现它们；他们的传统一直在教会中以文字保存。如果那些是合适的术语，他们就不会像他们所做的那样提及圣父、圣子和圣神——当然，为了这种革新而改造信德的术语是虔敬或安全的；否则他们全都是无知的人，未受教于奥迹，且不认识他所谓的合适名称——那些人确实既不知道也不希望将他们自己的观念优先于天主口中传授给我们的东西。\par

% structure: p0068 source=h2 target=subsection reason=relative-heading-rank; base=h2

\subsection{十四、他在提及救恩教义时犯了错，因为采用自己选择的术语，而不用传统的术语圣父、圣子和圣神。}

我认为这种新词发明的原因对每个人都是明显的——即：每当\Quote{父}和\Quote{子}这些词被说出时，每个人立刻便认出它们所暗示的彼此间恰当而自然的关系。这种关系通过称呼本身立刻传达。为了防止它被理解关于圣父和独生子，他剥夺了我们这个随词语进入耳中的关系概念，并放弃受启示的术语，而用其他一些为了伤害真理而设计的术语来阐述信德。\par

然而，他所说的一句话是真实的：即总结起来的是他自己的教训，而非天主教的教训。的确，任何深思的人都能轻易看出他言论的不敬。现在详尽讨论他的意图何在，并非不当：他把那至高与本然的最高程度只归于圣父的存有，而不承认圣子与圣神的存有是至高与本然的。依我看来，这是他对独生子和圣神的\Quote{存有}完全否认的前奏，他的这一体系暗中旨在取消对其位格的一切真实信仰，而在表面上和言语中承认它。对他的话稍加反思，任何人都能察觉出确是如此。这不像是认为独生子和圣神确实以明确位格存在的人，会对表达全能天主伟大的名称如此挑剔。承认事实，然后对适当措辞作细微区分，这确实是极端的愚昧。因此，在将最高程度的至高与本然的存有只归于圣父时，他通过对其他两位的沉默让我们猜测（在他看来）他们并非本然存在。被否认了本然存有的事物，怎么能说真实存在呢？当我们否认它的本然存有时，我们必然要把它的一切相反术语加给它。不能被适当言说的，就是被不适当言说的，因此证明它不能适当言说，就是证明它并不实际存有。欧诺弥似乎就意在如此，在他的教训中引入这些新名词。因为没有人能说他出于无知而陷入幼稚的幻想，将至高者与低者作空间上的分离，把好似山巅的位置留给圣父，而把圣子置于低处的洼谷。当讨论理智与属灵的事物时，没有人会幼稚到设想空间上的差异。地域位置是物质的性质；而理智与非物质的显然脱离了地域的概念。那么，他说只有圣父享有至高存有，理由何在？因为他多次展示中声称自己聪明，甚至\Quote{使自己过于聪明}，正如圣经所禁止的\BibleRef{训 7:16}。很难认为他是出于无知而飘入这些概念。\par

% structure: p0071 source=h2 target=subsection reason=relative-heading-rank; base=h2

\subsection{十五、他错了，因他使圣父的存有独为本然与至高，藉他对圣子和圣神的省略，暗示他们的存有是被不当言说的，而且较低。}

但无论如何，他会承认这存有的至高并不表示能力、善性或任何这类事物的超越。每个人都知晓——且不说那些被认为学识渊博的人——即独生子与圣神的位格在完美的善性、完美的能力和一切类似品质上毫无欠缺。善，只要不能成为其对立面，它的善性就没有界限：只有它的对立面才能限制它，如同我们从具体例子所见。力量仅在虚弱侵袭时才被阻止；生命仅被死亡所限；黑暗是光的终结：总之，每个善都受其对立面限制，且仅受此限制。那么，如果他假设独生子和圣神的性能会变坏，他就明显减损了对他们善性的概念，使他们能与对立面关联。但如果神圣而不变的性能不可能堕落，正如连我们的仇敌也承认的，我们就必须视其在善性上绝对无限：无限与无穷相同。但在无穷与无限中设想过度与不足，是极其不合理的：因为如果我们在其中设定了增加和减少，无穷的概念又如何能保留？我们只有通过对界限的比较才能得到过度的概念：没有界限，我们就无法设想任何过度。然而或许这并不是他的用意，而他只是凭时间先后的特权来赋予这种超越，并仅以此观念宣称唯有圣父的存有是至高的。那么他必须告诉我们，他凭什么标准给圣父量出了更多的生命长度，而在圣子的位格中从未预先构想任何时间上的区别。\par

而且，暂且假设为了论证起见，情况如此，那么单就纯粹存在而言，时间上在先的存在对后来的存在有何优越性，以致他能说一个是至高和真正的，另一个却不是？因为年长者与年幼者相比，寿命虽然更长，但其存在却并不因此有所增减。用一个例子可以清楚说明这一点。就存在而言，与\Person{亚巴郎}相比，生活在十四代之后的\Person{达味}有什么不利？就人性而言，后者发生了什么变化？他因时间上较晚而成为较少的人吗？谁会愚蠢到如此断言？他们存在的定义对两者相同：时间的流逝并不改变它。没有人会断言，时间上在先的人更是一个人，而后来才寄居在世上的则较少；仿佛人性已在前者身上耗尽，或者时间已将其主要力量花费在死者身上。因为时间无权为每个人界定本性的尺度；但本性保持自足，在随后的世代中保存自身；而时间有它自己的进程，无论围绕还是流经这个在本性界限内稳固不动地存留的本性。因此，即使如我们的论证暂且假设，在时间上被允许有优势，他们也不能恰当地将存在的至高无上仅仅归于圣父；但既然在时间的优先权上实际上没有任何差别，那么对于这些先于时间的存在，人怎么可能会产生这样的想法呢？我们所能发现的一切距离尺度都低于天主性体：所以，对于那些试图通过优越与低劣的区别来分割这个先于时间且不可理解的存在的人，没有任何根据。\par

我们也毫不犹豫地断言，他们教条式教导的乃是犹太教义的辩护，因为他们主张只有圣父的存在有自立体，并坚持只有这个有真正的存在，而将圣子和圣神的存在算作非存在，因为那并不真正存在的只能名义上并因用词不当而说存在。例如，人的名字不是给代表一个人的肖像，而是给绝对如此的那一位，即画的原型，而不是画本身；而肖像只是在言语上是人，并不绝对拥有归属于它的性质，因为其本性并不是所称呼的。在我们面前的情况也是如此，如果存在被真正地归于圣父，但到了圣子和圣神就停止了，这简直就是对救恩讯息的公然否认。让他们离开教会，退回到犹太人的会堂，因为他们通过否认圣子真正的存在来证明圣子的非存在。不真正存在的东西与非存在是同一回事。\par

再者，他本意是想在此显得非常聪明，并且对那些试图写作却缺乏逻辑力量的人评价很低。那么，让他告诉我们，尽管我们如此可鄙，他是用何种技巧在纯粹存在中发现了更大和更小。他建立\Quote{一个存在比另一个存在更是存在}的方法是什么——取存在最朴素的意义——因为他不能提出那些在存在的概念中被理解并围绕它的各种性质与属性，但它们不是主体本身？阴影、颜色、重量、力量或声誉、独特方式、性情、任何与身体或心灵相关的性质，在此都不予考虑：我们只需探究所有这一切的实际主体，即绝对称为存在的那个，是否在存在的程度上与另一个不同。我们还不知道，在两个已知的仍然存在的存在物中，一个更多存在，另一个更少存在。只要它们都在存在的范畴内，两者都是同等的存在，并且当所有关于更多或更少价值、更多或更少力量的概念被排除后。\par

因此，如果他否认我们可以将独生子视为完全存在——因为他的陈述似乎导向这一深度——通过拒绝赋予他真正的存在，那么让他连较少的程度也否认吧。然而，如果他的确承认圣子以某种实体的方式自立存在——我们现在不争论具体方式——他为何又夺回他已承认圣子所是的，并证明他并不真正存在，正如我们所说，这等于不存在？因为正如人性不属于那不具备\Quote{人}一词完全内涵的东西，并且在缺少任何属性的情况下，其整个概念都被取消；同样，在任何其完全和真正的存在被否认的东西中，对其存在的部分肯定根本不是证明其自立存在；事实上，对其不完全存在的证明就是对其完全抹除的证明。因此，如果他明智，他就会转向正统信仰，并从他的教导中移除圣子和圣神的本性中较少和不完全的观念：但如果他决心亵渎，并因某种不可测知的原因如此回报他的造主、天主和恩主，那么无论如何让他放弃拥有一些浮夸学识的自负，不哲学地，如他所做，把存在堆积在存在之上，一个在另一个之上，一个真正，另一个不真正，却找不出任何理由。我们从未听说过任何外教哲学家犯过这种愚行，也从未在受启示的著作中或人类的普遍理解中遇到过它。\par

我以為從以上所述已可明瞭這些新造名稱的目的是什麼。他把它們當作所有危害信德之事的據點或基石：一旦他能使\Quote{唯獨那一位存有才是至高而真正意義上的存有}這一觀念流傳開來，他就可以攻擊其他二位，說牠們屬於較低層級，不被視為真正的存有。他在後文尤其顯示了這一點，當他討論對聖子和聖神的信仰時，他不再繼續使用這些名稱，以免藉由那些稱呼向我們呈現牠們本性的特有標記：這個總是告訴我們要藉由合適的名稱和語句引導聽眾心靈的人，卻讓牠們被忽略過去。然而，還有什麼名稱比\Quote{真理本身}所賜予的更加合適呢？他以自己的觀點反對福音，不稱聖子，卻稱\Quote{一個藉由第一位而存在，但在第一位之後、卻在萬有之前的存在}。為了摧毀對獨生子的正確信仰而說這話，將由他後續的論證更加明顯。然而，這些話中的危害尚屬有限：一個對基督毫無不敬之意的人有時也可能使用它們；因此我們暫且省略所有關於我主的討論，將答覆保留給他對基督更明顯的褻瀆。但關於聖神的褻瀆則是明白而不加掩飾的：他說聖神不應與聖父或聖子並列，而是隸屬於兩者。因此，我將盡可能仔細地審視這一說法。\par

% structure: p0078 source=h2 target=subsection reason=relative-heading-rank; base=h2

\subsection{十六、探討\Quote{服從}的意義：他說聖神的本性隸屬於聖父與聖子的本性。證明聖神與聖父、聖子同等，而非較低。}

那麼，我們首先確定\Quote{服從}一詞在聖經中的意義。它適用於誰？造物主為尊崇按自己肖像受造的人，\Quote{使}動物受造物\Quote{服在他腳下}；正如偉大的達味在聖詠中敘述這份恩寵時呼喊說：\BibleQuote{祂使萬有，} 他說，\BibleQuote{服在他的腳下，} 並一一列舉那些被服從的受造物。聖經中還有另一種\Quote{服從}的意義。聖詠作者將戰爭勝利的緣由歸於天主自己，說：\BibleQuote{祂使萬民萬國服在我們腳下，} 和\BibleQuote{祂使萬民服在我以下。} 聖經中常見此詞，表示勝利。至於將來萬有對獨生子的服從，並藉著祂對聖父的服從——在宗徒以深奧智慧論述天主與人之間的中保服從聖父的那段經文中，藉著那共享人性的聖子的服從，暗示了人類的真實服從——我們現在不討論，因為它需要充分而詳盡的審查。但僅就服從一詞明白無歧義的意義而言，他怎能宣稱聖神的存有服從於聖子和聖父的存有？如同聖子服從聖父，按宗徒的思想？那麼在這種觀點下，聖神應與聖子並列，而非在祂之下，因為兩位位格都處於這較低層級。這不是他的意思？那麼是怎樣？如同動物受造物服從理性受造物，如在聖詠中？那麼其間差異有如動物受造物服從於人時所隱含的那樣。或許他也會拒絕這種解釋。那麼他就不得不來到僅剩的一種解釋，即聖神起初處在反叛行列，後來被更高的力量強迫向征服者屈服。\par

任凭他选取这两者中他喜欢的那一个：无论他选取哪一个，我都看不出他如何能避免那必然的亵渎之罪：无论他说圣神如同受造之物\OriginalTerm{受造之物}{brute creation}（如鱼、鸟、羊）那样\OriginalTerm{隶属}{subject}于人，还是说祂像叛乱者那样被掳\OriginalTerm{俘虏}{captive}于一个更高的权能。或者他两者都不是，而是完全用与圣经不同的意义来使用这个词？那么，那意义是什么？他是否规定我们必须将祂视为低等而非平等，因为祂是由我们的主按顺序第三位赐给门徒的？按同一推理，他应使圣父低于圣子，因为圣经常把我们的主之名放在首位，把全能圣父放在第二位。我们的主说\BibleQuote{我与父}。\BibleQuote{我们的主耶稣基督的恩宠\OriginalTerm{恩宠}{grace}，和天主的爱}，以及勤奋研读圣经证言者可能收集的无数其他段落：例如，\BibleQuote{恩赐虽有区别，却是同一的圣神；职分虽有区别，却是同一的主；功效虽有区别，却是同一的天主}。这样，那么，让那被排在第三位的全能圣父成为\OriginalTerm{隶属于}{subject}圣子和圣神者。然而，我们从未听说过这样一种哲学，它仅仅因为某物出于某种序列理由而被排在第二位或第三位，就将其归入低等和依附的范畴——而这正是我们的作者想要做的，他争辩说在传授位格中的秩序意味着尊严和本性上的多少差异。事实上，他断定次序上的顺序表明本性的不同：他从何处得到这种奇想，什么必要性迫使他如此，并不清楚。单纯的数字阶位并不产生不同的本性：我们计数的东西，无论计数与否，在本性上保持不变。数字只是事物纯粹数量的标记：它只将那些本性价值较低的东西放在第二位，而是按照计数者的意图造成所指示的数值对象的顺序。\InnerQuote{保禄和息耳瓦诺和弟茂德}是三个按特定意图被提及的人。息耳瓦诺位于保禄之后第二位，是否表明他不是人？或者弟茂德因为是第三位，就被如此排列的作者认为是一种不同的存在？并非如此。每个人在这安排之前和之后都是人。言语不能同时说出所有三个名字，因此按照合适的顺序分别提及每一个，但用连词连接它们，以便名字的连接显示三者朝向同一目的的和谐行动。\par

然而，这并不令我们的新教义学家满意。他反对圣经的安排。他将祂从那我们的主亲自宣告的、与圣父和圣子同等的固有自然地位和关系分离出来，将祂列在\InnerQuote{隶属者}中：他宣告祂是两位位格的工程，圣父提供祂构成的根源，独生子作为祂存在的创造者：并将此定义为祂\OriginalTerm{隶属}{subjection}的理由，却尚未阐明\InnerQuote{隶属}的含义。\par

% structure: p0082 source=h2 target=subsection reason=relative-heading-rank; base=h2

\subsection{十七、讨论关于这位人士所宣称的\Quote{随伴}圣父与圣子之存有的\Quote{能量}的确切本质。}

然后他说：\Quote{当然必须在此叙述中包括随伴每个\OriginalTerm{存有}{Being}的能量，以及相应于这些能量的名称。}笼罩在如此模糊的迷雾中，其意义远非清楚：但可以推测如下。所谓存有的能量，他意指那些产生圣子和圣神的能力，以及第一存有用以造出第二存有、第二存有用以造出第三存有的能力：他意指所产生的结果的名称已以一种相应于那些结果的方式被赋予。我们已经揭露了这些名称的危害，并且当我们回到那部分问题、若需要进一步讨论时，我们还会再次揭露。\par

然而，现在值得花片刻思考\OriginalTerm{能量}{energies}如何\OriginalTerm{随伴}{follow}存有：这些能量本质上是什么：是与它们所随伴的存有不同，还是存有的一部分，属于它们的内在本质：然后，如果不同，它们如何以及从哪里产生：如果相同，它们如何被从存有中切断，而不是共同存在，仅仅从外部随伴它们。这是必要的，因为我们无法一下子从他的话语中得知，某种自然必然性是否迫使\Emph{能量}（无论它是什么）\Emph{随伴}存有，就像热和蒸气随伴火，以及各种蒸汽随伴产生它们的物体一样。不过，我认为他不会肯定我们应把天主的存有视为某种异质和复合的东西，把能量作为像某个主体中的\OriginalTerm{依附体}{accident}那样不可分离地包含在自己的概念中：他必定是指，那些存有，故意和自愿地运作，由自身产生所期望的结果。但是，如果如此，谁会把这自由意向的结果称为其外在后果之一？我们从未听说过在日常用语中有这样表达的；任何工人的能量都不会被说成\Emph{随伴}那个工人。我们不能将一个与另一个分离，让一个单独留下：相反，当人提到能量时，他就在观念中包含了那随能量运作的；当人提到工人时，他立刻暗示了那未被提到的能量。\par

试以一例阐明其意。吾人云某人操铁工、木工或它艺，此一语即兼含工作与工匠之意，若去其一，则另一亦不存在。若能力与施能者如此共思，则在此情形中，岂能谓产生第二存有之能力\InnerQuote{随}第一存有之后，如同一中间物介于两者，既不与第一存有之本性融合，亦不与第二存有结合？其与第一存有分离，因其非其本性，仅为本性之运作；与后生者分离，因其后生者非单纯之能力，乃一活动之存有。\par

% structure: p0086 source=h2 target=subsection reason=relative-heading-rank; base=h2

\subsection{十八、彼无理由区分天主圣三中多数存有。彼未提供证明其如此。}

亦当察其下文。彼称一存有为另一存有之工，第二为第一之工，第三为第二之工。此论基于何种先前证明？用何证据、何方法使人信服后存有乃前存有之果？纵使可能从受造物得此比喻，以低级存在推想超性者亦不尽稳妥，然从自然现象论至不可理解者，其误或可宽恕。但实则无人敢断言：天为天主之工，太阳为天之工，月亮为日之工，星辰为月之工，其余受造物为星辰之工，盖万物皆为一天主之工：因为天主子、万物之父只有一位，万物皆出于祂。若亦由相互传递而生，如动物之族类，在此亦非一造另一，因本性通过各代延续。既不能在受造世界作此断言，彼如何能对超性存有宣告第二为第一之工，以此类推？若彼思及动物生殖，并以为纯粹存有间亦行此过程，使长者生幼者，则仍未一贯：因此类产物与亲代类型相同；而彼赋予其继位者以陌生而非继承的属性：如此徒显其多言虚妄，却欲以巧拳师之姿双拳齐中真理。为显示圣子与圣神之低等及内在价值之衰减，彼宣称\Quote{一由\Emph{另一}而生}；为使识得互生者不生亲属之情，彼自相矛盾地宣称\Quote{一由\Emph{另一}而生}，同时显明圣子与圣父之本性相较如私生子。\par

然我认为，在论及此一切之前，便可指摘其过。若有一人，素未习讨论、不谙逻辑表达，以此漫然方式陈述己见，或可因其未用公认方法而立论而稍加宽容。但既然欧诺弥如此富于此种能力，能以\InnerQuote{不可抗拒}之论证方法进至超自然之境，彼岂能不知此\InnerQuote{不可抗拒}之对隐藏真理的知觉，在一切逻辑推演中从何起点出发？众所周知，凡此类论证必须从明显且公认之真理出发，以自身使人信服仍存疑之真理：若无显然者引导我们趋近未知，则无一未知者可被把握。反之，若用以阐明此未知之起点与普遍信念相悖，则未知者将久不得阐明。\par

是故，教会与异质派之间整个争论在于：是否应将圣子与圣神视为受造或非受造之存有？吾人对手主张众人否定之事：彼大胆断定——不寻任何证明——每一存有为前存有之工。其教育方法或思想学派何以允许如此，实难明了。每一证明过程必始于无可否认或攻驳之公理，以便从所假设者经中间三段论合法推导出未知者。是故，将本应探究之对象本身作为证明的前提，乃是借晦暗证晦暗，借幻象证幻象。此乃\InnerQuote{瞎子领瞎子}，因断言万物之创造者、造物主为受造物，实为盲目无据之论；由此他们又引出盲目结论：即圣子在本质上与圣父相异，存有上\Emph{不相似}，完全欠缺其固有特性。然此已足。其思想之显明亵渎处，我等亦可稍缓驳斥。今须返回审视其下文。\par

% structure: p0090 source=h2 target=subsection reason=relative-heading-rank; base=h2

\subsection{十九、彼承认天主性为\InnerQuote{单一}仅止于言辞。}

\Quote{\OriginalTerm{存有}{Being} 实际上和在概念上，都具有一个纯一、单一且绝对同一的本性，以其尊荣来衡量；并且由于工作被每个操作者的运作所限定，运作又被工作所限定，那么每个存有所随之而来的运作不可避免地会有大小之别，有些属第一等，有些属第二等。} 贯穿这一切的意图，尽管表达冗长，却是同一而一致的；即，要确立\Person{圣父}与\Person{圣子}之间，或\Person{圣子}与\Person{圣神}之间并无关联，而是这些\OriginalTerm{存有}{Being}彼此分离，拥有彼此陌生而不熟悉的本性，不仅在此方面有差异，而且在尊荣的大小和隶属上也有差异，因此我们必须认为一个比另一个更大，并呈现其他各种差别。\par

在许多人看来，沉溺于如此明显之事，并对他们来说看似虚假、可恶且毫无根据的东西加以讨论，似乎是无用的；然而，为了避免甚至让人以为我们因缺乏反驳而不得不放过这些陈述，我们将尽全力迎击它们。他说：\Quote{它们中间的每个\OriginalTerm{存有}{being}都是纯一、单一且绝对同一的，以它的尊荣来衡量，无论是在事实上还是在概念上。} 然后，他将这个非常可疑的陈述作为公理预设，并将自己的\Quote{ipse dixit}视为足以代替任何证明，便自以为言之成理。\Quote{有三个\OriginalTerm{存有}{Beings}：}因为他说的\Quote{它们中间的每个\OriginalTerm{存有}{being}}就暗示了这一点；如果他只指一个，就不会用这些词。现在，如果他这样谈论\OriginalTerm{存有}{Beings}之间的相互差别，是为了避免与\Person{撒伯里乌}的异端有牵连，即把三个称号应用于同一个主体，那么我们可以默认他的说法；也没有任何信徒会反对他的观点，只是他似乎在用词上，以及单纯在表达方式上——用\OriginalTerm{存有}{beings}而非\OriginalTerm{位格}{persons}——有所不妥：因为就存有而言相同的事物，在位格上的定义却不完全一致。例如，\Person{伯多禄}、\Person{雅各伯}和\Person{若望}，就\OriginalTerm{存有}{beings}而言是相同的，每个人都是人；但在各自位格的特质上，他们并不相同。那么，如果他只是为了证明混淆\OriginalTerm{位格}{Persons}、把三个名字都套在同一个主体上是不对的，那么他的\Quote{这句话}，用宗徒的话来说，就是\BibleQuote{确实的，值得完全接纳}\BibleRef{弟前 1:15}。但这并非他的目的：他这样说，不是因为仅凭公认的特质将\OriginalTerm{位格}{Persons}彼此分开，而是因为他使每个\OriginalTerm{存有}{Being}实际的实体性存在与其他的不同，或者更确切地说，与自身不同：因此他谈到多个\OriginalTerm{存有}{beings}，它们之间具有使彼此疏离的独特差异。因此我宣布，他的观点毫无根据，且缺乏原则：它从未经认可的数据出发，然后仅仅通过逻辑便以此构建出亵渎之辞。它并未试图进行任何能够吸引人接受这种教义概念的论证：它仅仅包含一种未经证明的不敬之辞，仿佛在向我们讲述一个梦。虽然教会教导我们不可将\OriginalTerm{信德}{faith}分散在多个\OriginalTerm{存有}{beings}之间，而必须承认三个主体或\OriginalTerm{位格}{Persons}中没有存有的差异，然而我们的对手却主张它们作为\OriginalTerm{存有}{Beings}之间存在多样性和不相似性，这位作者自信地假定那些从未、也永不可能通过论证得到证明的事情已经得到证明：也许他甚至尚未为自己的言论找到听众；或者他可能从某位智能聆听者那里得知，任何随意、无根据的陈述都是\Quote{老妇之谈}，本身就无力证明问题，无力借助任何来自\WorkTitle{圣经}或人类推理的诉求。关于这一点就说到这里。\par

但让我们继续审视他的话。他宣称他在论述中隐约描绘的这些\OriginalTerm{存有}{Being}中的每一个，都是\Quote{单纯且绝对单一}。我们相信，即使最为粗鄙愚昧的人也不会否认，那祝福且超越的\OriginalTerm{天主性体}{Divine Nature}是\Quote{单一的}。那无形、无相、无量的，不能被设想为多形且\OriginalTerm{复合的}{composite}。然而，只需极轻微的反思便可明了，将至高存有视为\Quote{\OriginalTerm{单纯}{simple}}的观点，无论他们说得多么巧妙，都与他们精心构建的系统全然不一致。因为谁不知道，确切地说，在\OriginalTerm{天主圣三}{Holy Trinity}的情形中，\OriginalTerm{单纯性}{simplicity}没有程度之分？在此情形中，无需考虑性质的混合或汇合；我们理解一种无部分无组合的能力；那么，如何以及凭什么，有人能在其中感知更少和更多的差异呢？因为凡在那裡标出差异者，必然想到某些性质在主体中的偶有性。他实际上必须已感知到其中大小差异，才将量的概念引入这个问题；或者他必须设定美善、能力、智慧或任何可虔敬地与天主相联系的事物的丰盈或减损：无论如何，他都将无法逃脱组合的观念。任何拥有智慧、能力或其他美善者，若非作为外在恩赐，而是根植于其本性，便不能在其中遭受减损；因此，若有人说他在天主性体中察觉更大和更小的存有，他便是在无意识中建立一个复合且异质的天主，并将主体视为一物，而将分享那性质能使先前不善的成为善的性质视为另一物。假如他想到的是一个真正单纯且绝对单一的存有，与美善同一而非拥有美善，他根本不能在其中数出更大和更小。此外，上文已说，美善唯有因邪恶的临在而减损，而在本性不会退化的地方，对美善无从设想界限：事实上，无限并非凭借任何关系而成为无限，而是就其本身而言超越限制。确实，很难看出一个反思的心智如何能设想一个无限大于或小于另一个无限。因此，如果他承认至高存有是\Quote{单一的}且同质的，就让他承认它与这单纯和无限的普遍属性相连。另一方面，如果他将这些\Quote{存有}彼此分离和疏远，认为\OriginalTerm{独生子}{Only-begotten}的存有不同于圣父的存有，圣神的存有不同于独生子的存有，在每种情况下都有\Quote{更多}和\Quote{更少}，那么就让他现在被揭露为只在表面上承认单纯性归于天主，但实际上证明祂是复合的。\par

但让我们按顺序重新审视他的话。\Quote{每一个\OriginalTerm{存有}{Being}在实际和概念上都有一个纯粹、单纯且绝对单一的本性，按其\OriginalTerm{尊严}{dignity}评定。}为什么\Quote{按其尊严评定？}如果他是在这些存有的共同尊严中思考，这个附加就是不必要的和多余的，并且停留在明显的事情上：尽管一个如此不当的词可能被原谅，如果有一丝敬虔的感觉促使他不拒绝它。但这里的祸害实际上并非由于词句的误解（那很容易纠正），而是与他的邪恶意图有关。他说这三个存有中的每一个都是\Quote{单一的，按其尊严评定}，为的是，凭借他对第一、第二和第三存有先前的定义，它们单纯性的观念也会受损。在断言唯独圣父的存有是\Quote{\OriginalTerm{至上}{Supreme}}和\Quote{\OriginalTerm{本有}{Proper}}，并拒绝将这两个称号归于圣子和圣神之后，据此，当他谈到它们全体都是\Quote{\OriginalTerm{单纯}{simple}}时，他认为自己有责任只按它们本质价值的比例赋予它们单纯性的观念，以至于唯独\OriginalTerm{至上}{Supreme}被设想为处于单纯性的顶峰和完美，而第二位，按其从至上地位下降的比例，也接受减损量的\OriginalTerm{单纯性}{simplicity}，而第三存有的情形，按两极端的价值减少的程度，与完美单纯性的偏差同样大：由此导致，圣父的存有被设想为纯粹单纯性，圣子的存有在单纯性上并非如此无瑕疵，而是混合了\OriginalTerm{复合}{composite}，圣神的存有则在复合中不断增加，而单纯性的量逐渐减少。正如不完美的美善必须被承认在某种程度上分享相反倾向，不完美的单纯性也不能不被视为复合的。\par

% structure: p0095 source=h2 target=subsection reason=relative-heading-rank; base=h2

\subsection{二十、他错误地假定，为了解释独生子的存在，有一种\Quote{能力}产生了\Person{基督}的位格。}

他将如此使用这些短语的意图，从下文将清楚看出，在那里他更加露骨地将我们对\OriginalTerm{圣子}{Son}和\OriginalTerm{圣神}{Holy Spirit}的观念物质化和贬低。\Quote{由于\OriginalTerm{能力}{energies}被\OriginalTerm{作品}{works}所限定，作品与能力相称，因此必然得出，这些伴随这些\OriginalTerm{存有}{Beings}的能力相对地有大小之分，有些属于较高层次，有些属于较低层次。}尽管他刻意用术语的迷雾包裹其含义，使得多数人难以发现，但按照我们已经检验过的内容，这一点将很容易澄清。\Quote{能力，}他说，\Quote{被作品所限定。}他以\Quote{作品}意指圣子和圣神，以\Quote{能力}意指产生它们的能力，这些能力，他稍早前说过，\Quote{伴随}\OriginalTerm{存有}{Beings}。\Quote{被……限定}这一措辞表达了被产生的存有与产生能力之间的平衡，或者更确切地说，与那权能的\Quote{能力}（用他自己的话）之间的关系，暗示被产生的东西并非操作者全部权能的效果，而只是其某种能力的局部效果，仅施展了足以与结果相称的那部分权能。随后他颠倒其陈述：\Quote{作品与操作者的能力相称。}其含义通过一个例子将变得更清楚。让我们思考一个皮匠的工具：即一把割皮刀。当它在需要切出某种形状的皮革上旋转时，被切除的部分受工具大小的限定，切出的圆的半径等于工具的长度，反过来说，工具的跨度将测出并切出相应的圆形。这就是我们的神学家对\OriginalTerm{独生子}{Only-begotten}的神圣位格所持的想法。他宣告，某种\Quote{伴随}\OriginalTerm{第一存有}{the first Being}的\Quote{能力}，以这样一种工具的方式，产生了相称的作品，即\OriginalTerm{我们的主}{our Lord}：这就是他荣耀天主子（祂现今已在父的光荣中被光荣，并将在审判日被揭露）的方式。祂是一件\Quote{与产生能力相称的作品}。但什么是这个\Quote{伴随}\OriginalTerm{全能者}{the Almighty}、并被设想为存在于独生子之前、且限定了祂存有的能力呢？某种自存的本质权能，它以自发的冲动行使自己的意愿。那么，这才是我们主的真正\OriginalTerm{父}{Father}。如果并非祂，而是属于外在伴随祂之物的某种能力产生了圣子，那我们为什么还要继续谈论全能者是父呢？如果按照\Person{欧诺弥}的说法，是别的东西给了圣子存在，使得祂像一个私生子（愿我主宽恕这种说法！）般潜入与父的关系，仅在名义上被尊为圣子，那么圣子又如何能继续是子呢？\Person{欧诺弥}怎能将我们的主列在紧随全能者之后呢？当他只将主排列在第三位，而将那中介的\Quote{能力}放在第二位？圣神照此顺序也将被发现不是在第三位，而是在第五位，因为那伴随独生子、且圣神借以产生的\Quote{能力}必然插在它们之间。\par

同样，万物借\OriginalTerm{圣子}{Son}受造也将发现毫无根据：我们的新派神学家发明了另一个先于祂的位格，世界的创造权必须归于该位格，因为按照他们，圣子自身的存在来自那\Quote{\OriginalTerm{能力}{energy}}。然而，如果他为了避免如此亵渎，而将那产生圣子的\Quote{能力}说成某种非实体的东西，那他就必须向我们解释，\OriginalTerm{非存有}{non-being}如何能\Quote{伴随}\OriginalTerm{存有}{being}，\OriginalTerm{非本体}{non-substance}如何能产生\OriginalTerm{本体}{substance}：因为，如果他那样做了，我们将会发现一个\OriginalTerm{非实在}{unreality}伴随天主，一个不存在者是一切存在的作者，根本非实体的东西限定了一个本体的本性，创造的操作力量最终被包含在非实在之中。这就是这位神学家教导的结果，他肯定天地及一切受造的工匠主（那在起初就存在的天主的\OriginalTerm{圣言}{Word}，万物借祂而受造）的存在归功于这样一个无根据的实体或概念，即他刚刚编造的那个不可名状的\Quote{能力}，并且祂被它限定，如同被一个非现实的监狱所包围。那位\Quote{窥视不可见之物}的人看不到他的教导所导致的结论。那就是：如果天主的这个\Quote{能力}并没有真实存在，并且这个非实在所产生的\OriginalTerm{工作}{work}也被它限定，那么很清楚，我们只能认为这工作中具有与这幻想中的工作生产者相同的性质：事实上，从非实在产生并被非实在所包含的东西，本身只能被设想为虚无。按事物的本性，对立物不能被对立物所包含：如水不能被火包含，生不能被死包含，光不能被暗包含，存有不能被非存有包含。但是，尽管他极其聪明，却没有看到这一点：否则他就是故意对真理视而不见。\par

某种必然性迫使他看到圣子有减损，并进而主张在圣神方面更进一层。他说：\Quote{这些伴随着这些存有的能力必然有相对的大小之分。}欧诺米并未向我们解释这种在天主性中定出大小强弱的必然性究竟是什么，至今我们自己也无法理解。迄今，在那些以纯朴之心接受福音的人中，普遍相信在天主性之上没有必然性迫使独生子如同奴隶般屈居低微。但他完全无视这一信念——尽管它值得认真考虑——反而教条式地断言我们必须设想这种低微。然而，他的这种必然性并未止步于此，更将他引向亵渎：正如我们逐一的检视已经表明的。也就是说，如果圣子并非生于圣父，而是生于某种无实体的\OriginalTerm{能力}{energeia}，那么祂就不仅被视为次于圣父，而且这种学说最终必然沦为纯粹的犹太教。这种必然性若推演下去，会显示出从虚无产生之物不仅微不足道，而且是一种连控告者也不敢妄称的亵渎。因为正如由实有而生者必然实有，由虚无而生者必然恰恰相反。当事物本身并非自存时，它如何能产生另一物呢？\par

既然如此，如果那\Quote{伴随}天主性并产生圣子的能力自身并无实存，那么没有人会盲目到看不见结论——即他的目的是否认我们救主的天主性。并且，如果圣子的位格被他们的学说从信仰中偷走，只留下一个名号，那么圣神也将因来自一连串虚无之物的血统而难以再被相信。那\Quote{伴随}天主性的能力自身并无实存：那么常识要求其产物也必然是虚无的；接着，一种无实体的第二能力随之产生；然后它宣称圣神由这个能力形成：因此他们的亵渎是显而易见的：无非是否认在无始无终的天主之后有任何真实的存在；他们的学说陷入了虚幻无实的虚构，其中没有任何真实基体\OriginalTerm{实有}{hypostasis}的基础。他们的教导就这样将论证搁浅在如此荒谬的结论上。\par

% structure: p0100 source=h2 target=subsection reason=relative-heading-rank; base=h2

\subsection{二十一、这些异端者的亵渎比犹太人的不信更恶劣。}

但我们且假设并非如此：因为他们姑且在理论上对人类施以仁慈，承认独生子和圣神有一些位格存在。而如果他们在承认这一点时，也承认了关于他们的合宜概念，他们就不会与教会教义交战，也不会断绝基督徒的希望。但是，如果他们给予圣子和圣神某种存在，只是为了给他们的亵渎提供材料，那么或许对他们来说——尽管说出这话有些大胆——更好的是背弃信仰，改信犹太教，而不是用这种假意的认同来侮辱基督徒的名号。无论如何，犹太人虽然至今仍拒绝圣言，他们的不虔敬只到否认基督已经来临，但仍盼望祂来临的程度；我们并未听到他们对所期盼的那位的荣耀有任何恶毒或毁灭性的想法。但这新\OriginalTerm{割损}{concision}派（或更确切地说\Quote{割损}派）却一面承认祂已来临，一面却像那些以恣意不信侮辱主肉身临在的人。他们想要用石头砸我们的主；这些人用亵渎的称呼砸祂。他们强调祂卑微卑贱的出身，拒绝祂在万世之前天主性的诞生：这些人同样否认祂从圣父而来的伟大、崇高、不可言传的出生，企图证明祂的存在源于受造，如同人类和所有受造之物那样。在犹太人眼中，我们的主被视为至高者之子是罪行；这些人也对真诚宣认祂的人感到愤慨。犹太人以为排除圣子的同尊可以荣耀全能者；这些人通过消灭圣子的荣耀，以为能更敬重圣父。然而，要完全说清他们加诸独生子身上的侮辱之数量与性质是十分困难的：他们捏造了一种先于圣子位格的\OriginalTerm{能力}{energeia}，并说圣子是这能力的产物和成果——这是犹太人迄今不敢说的。然后他们限制祂的本性，将其封闭于造祂之能力的某种界限内：这种生产能力的多寡是他们用来围困祂的尺度；他们设计它如同一件斗篷将祂包裹起来。我们不能指责犹太人做这些事。\par

% structure: p0102 source=h2 target=subsection reason=relative-heading-rank; base=h2

\subsection{二十二、他无权在天主性中主张大小之别。教会教义的系统陈述。}

然后他们在祂的存在中发现某种贫乏——以欠缺的方式，却未告诉我们他们是用什么方法度量那既无量也无大小的存在；他们能精确断定独生子的大小比完美缺少多少，因而必须被归类为低劣和不完美；他们还有许多其他主张，有的公开断言，有的暗中推论：始终把他们对于圣子和圣神的宣认仅仅当作其不信精神的练习场。那么，我们怎能不更怜悯他们，甚至超过怜悯被判罪的犹太人呢？因为连犹太人都不敢提出的看法，竟被他们推论出来。那使圣子和圣神之存有相对较小的人，在言辞上或许只是轻微不敬；但若严格检验其观点，便会发现亵渎之极。那么，让我们审视这一点，并请容我（为着教义并为了清晰展示他们所论证的谎言）进一步陈述我们自己对真理的理解。\par

如今，一切存有的终极划分在于可理解者与可感觉者。宗徒大体上称可感觉的世界为\Quote{那可见的}。因为一切物体都有颜色，而视觉捕捉到它，他便用\Quote{那可见的}这一粗略而现成的名称来称呼这世界，略去了所有其他本质上固存于其结构中的属性。同样，宗徒又用\Quote{那不可见的}作为整个理智世界的通称\BibleRef{哥 1:16}：藉着撤去一切感官理解的概念，他引导心灵趋向非物质和可理解者。理性再次将\Quote{那不可见的}划分为非受造与受造，以推理方式把握它：非受造者是那创造万有的，受造者则是由非受造者获得其存在与力量。因此，在可感觉的世界中，包含着一切我们藉身体感官所理解的，以及其中品质的差异涉及多与少的概念，诸如数量、质量及其他属性。\par

但在理智世界——我指的是其受造部分——中，无法容纳在可感觉世界中所感知到的诸如此类差异的概念；于是设计了另一种方法来发现更大与更小的程度。一切美善的泉源、根源、供应，被视为存在于非受造世界中，整个受造界倾向它，并且只因其分有首要美善而触及并分享至高的存在；因此，根据各者意志自由的程度上，对至高福乐的分有程度不同，随之在这受造界中，更大与更小便依据各者倾向的比例而显明。受造的理智性体处于善恶之间的边界，因而能趋向两者，并随意倾向于其所选择的事物，正如圣经所教导的；所以我们能说，它处于卓越的高峰多少，仅与其远离恶而接近善的程度成正比。而非受造的理智性体则远离此类区分：它不是藉获得而拥有美善，也不是分有某一高于自身的美善之善；在其本身的本性中它就是善，并被如此理解；它是善的泉源，是单纯、单一、不加复合，即使我们的敌人也承认这一点。但天主圣三自身内在具有与其本性的威严相称的区分，却并非像欧诺弥所设想的那样以数量来理解：（实际上，凡将较少的善之概念引入天主圣三中所信的任何一位，必然由此承认在那缺乏善者中有相反性质的混杂；而无论是对独生子还是对圣神作此想象，都是亵渎的）：我们认为天主圣三是至为完美且不可测度的卓越，但又在其自身内包含清晰的区分，这些区分存在于每一位的独特属性中：藉着非受造的共同属性而拥有不变性，但又因每一位的独特性质而区分。在每一位中所观察到的这种独特性质，将一位与另一位鲜明而清楚地区分开：例如，圣父既是非受造也是非受生；祂既从未受生也从未受造。这非受造性为祂与圣子、圣神所共有，而祂作为圣父还是非受生的。这是独特而不可传递的，在其他各位中无法见到。圣子在祂的非受造性上与圣父和圣神相连，但作为圣子和独生子，祂拥有并非全能者或圣神的属性。圣神藉其本性的非受造性与圣子和圣父相连，但凭其自身的标记与它们区分。祂最独特的特征是，祂既不是我们在圣父和圣子中所分别思考的任何一个。祂\Emph{是}简单地存在，既非非受生，也非独生子：这构成了祂的主要独特性质。因非受造性与圣父相连，又因不是\Quote{父}而与祂分离。因非受造性的纽带以及与从至高者领受存在而与圣子相连，又因不是圣父的独生子、而藉圣子自身得以显示的这一特性而与祂分离。再者，既然受造界是由独生子造成，为了确保圣神不因其藉圣子显示而被认为与这受造界有某种共通之处，祂便因自身的不变性以及对一切外在美善的独立而与受造界区分。受造界在其本性中并不拥有这种不变性，正如圣经在描述晨星堕落时所言，主向门徒启示了对此事的奥妙：\BibleQuote{我看见撒殚如同闪电从天坠落\BibleRef{路 10:18}。}但正是那些使祂与受造界分离的属性，构成了祂与圣父和圣子的关系。一切不可能堕落的事物，都拥有同一个\Quote{不变}的定义。\par

已经说了这些作为前言，我们现在可以讨论敌人教导的其余部分。他在论圣子和圣神的体系中说道：\Quote{必然的结果是，存有是相对地较大和较小。}那么让我们探究这种差异必然性的意义何在。它是来自以某种物质方式彼此测量比较，还是来自从道德卓越的多寡这一精神基础，或者来自纯粹存有？但就最后一点而言，有能力的思考者已经表明，如果从属性与特性中抽象出存有，只按其赤裸的定义来看它，则不可能设想任何差异。此外，将这种差异设想为存在于独生子和圣神的情况中，在于道德卓越的强度或减弱，并因此暗示他们的本性容许向任一方向变化，以致同样能容纳对立面，并被置于道德美与其对立面之间的边界地带——这是严重的亵渎。这样想的人将证明他们的本性本身是一回事，而凭借参与这种美或其对立面而变成另一回事：例如铁的情况：如果某时靠近火，它获得热的性质而仍是铁；如果放在雪或冰中，它的性质改变为占主导的影响，让雪的冷气渗入其孔隙。\par

正如我们不能根据铁上现时观察到的性质来称呼其材料（因为我们不把经火或冰淬炼的东西称为火或冰），同样，一旦我们承认这些异端者的观点，即对于赋予生命的大能而言，善并非本质上内在于它，而只是被赋予它，那么就不可能恰当地称它为善：这样的概念将迫使我们视它为某种不同的东西，并非永恒地展现善，并非本身被归入真正的善，而是善有时不在它内，有时也不太可能在它内。如果这些存有只是通过分享某种超越自身的东西才成为善，那么显然在分享之前它们并不善，而如果它们本非善，后来被善的影响所染色，那么如果再次与善隔离，就必定被视为非善。因此，如果这个异端得势，就不能视神圣本性为传递善的，反而它自身需要善：因为一个人如何能将自己不拥有的东西给予他人呢？如果它处于完美状态，则不能想象任何减退，谈论较少的完美是荒谬的。另一方面，如果它对善的参与是不完美的，而这就是他们所谓\Quote{较小}的含义，那么注意后果：任何处于那种状态的东西永远不能帮助较低者，而将忙于满足自己的匮乏。因此，按照他们的说法，眷顾是虚构的，审判和独生子的救恩计划也是虚构的，所有其他被认为是祂所做、仍然由祂所做的工作也是虚构的：因为祂必然要忙于照顾自己的善，而必须放弃对整个宇宙的监督。\par

那么，如果这种猜测得势，即我们的主并未在每一种善中得以成全，那么亵渎的结论就很容易看到。既然如此，我们的信仰便是虚空的，我们的宣讲也是虚空的；我们的希望，从信仰取得实质，也是不实在的。为什么他们受圣洗归于基督，如果祂自己没有善的能力？愿天主宽恕我这样说！为什么他们相信圣神，如果对祂也给出同样的解释？如果重生的大能在其本性中不拥有无误性和独立性，他们如何藉着圣洗从必死的出生重生？他们如何能改变自己\Quote{卑贱的身体}，而他们认为那位将要改变祂自身的也需要改变，即需要另一位来改变祂？因为只要一个本性在善方面有缺陷，那较高的存有就对这较低的存有施加一种不断的吸引，使之趋向自己：这种对更多的渴求永不会停止：它将伸向尚未把握的东西；这有缺陷的主体将总是需要供给，总是转变为更伟大的本性，却永不会触及完美，因为它找不到一个目标来把握，并停止向上的冲动。第一善在其本性上是无限的，所以必然的结果是，对它的享受的分享也将是无限的，因为总是会把握更多，而已经把握的东西之外总有东西被发现，这探索永不能赶上它的对象，因为它的资源如同参与者的增长一样无穷无尽。\par

那么，这些就是他们区分位格在善方面的差异所生出的亵渎。另一方面，如果多与少的程度是在某种物质意义上被理解，那么这种猜测的荒谬立刻显而易见，无需详细考察。性质、距离、重量、形状等所有构成物体概念的观念，将必然连同这种猜测被引入对神圣本性的看法中：而一旦假设了复合物，就必须承认该复合物的解体。这种可怕的教导，竟敢在无形无质的东西中发现较小和较大，使我们陷入这些及诸如此类的结论，这里只指出少数例子：要想揭示隐藏其下的所有危害确实不易。然而，从他们亵渎的前提所导致的惊人荒谬，从这简短的说明中就会清楚。现在，在简短地界定和确认我们自己的教义之后，我们继续讨论他们的下一个立场。因为受默感的见证是任何教义真理的可靠检验：所以在我看来，我们的教义可以通过引用神圣的话语得到很好的保证。\par

在一切存在物的划分中，我们便找到这些区分。有一个诉诸我们感知的感官世界；在此之外，有一个由理性藉着感官对象所观看到的世界：我指的是理知世界；在其中我们又探得进一步区分，分为受造者和非受造者；我们已界定天主圣三属于后者，而凡在祂之后可能存在或可思及的一切属于前者。但为使这陈述不致无凭，而由圣经加以确认，我们将补充：我们的主并非受造，而是出自圣父，正如圣言在福音中以自己的亲口所证，以一种不可言宣、奥秘的诞生或发出的方式；有什么比我们的主在整部福音中不断宣示：圣父本身是祂的父亲，而非创造者，祂不是天主的化工，而是天主子，这更真实的见证吗？正如祂愿按其血肉关系称呼祂与人类的联系时，称祂存在的那个阶段为人子，由此指明祂与生祂者按血肉本性的亲属关系；同样，祂以儿子这个称号表达了祂与全能者之间真实的关系，以儿子之名显示这种自然的联系。尽管有些人为了抵触真理，确实按字面不加解释地接受了以比喻晦涩形式含有隐意的言语，并援引箴言作者放在智慧口中的\Quote{受造}一语，以支持他们扭曲的观点。他们实际上说，\BibleQuote{主创造了我}是证明我们的主是受造物的证据，仿佛独生子本人藉这句话承认了这一点。但我们无需理会这种论证。他们没有给出任何理由说明为什么我们必须将那经文指归我们的主；也无法证明希伯来文中的词义只能导致这一意义而非其他，因为其他译者将其译作\Quote{拥有}或\Quote{设立}。最后，即使这是原文的含义，其真实意思也不会如此明显和表面。因为这些箴言式的言论并不立即显明其目标，反而隐藏它，只通过间接的含意揭示它。从上下文他所说的\BibleQuote{当祂在风中设立祂的宝座}以及所有类似表达，我们可以判断这段特殊经文的晦涩。天主的宝座是什么？是物质的还是意识的？风是什么？是我们所熟悉的风，自然哲学家告诉我们是由蒸气与湿气形成的？还是以另一种不为人类熟悉的方式理解，当它们被称为祂宝座的根基时？这非物质的、不可理解的、无形状的神性的宝座是什么？谁能按字面理解这一切？\par

% structure: p0111 source=h2 target=subsection reason=relative-heading-rank; base=h2

\subsection{二十三、我们信德的这些道理由圣经章节所见证和确认。}

因此很明显，这些都是隐喻，包含着比表面更深的意义：因此，虔敬的探究者不应从中产生任何怀疑我们的主是受造物的理由，这些人已按照福音作者伟大的话语受教：\BibleQuote{凡受造的，都是藉着祂造的}以及\BibleQuote{在祂内得以存立}。\BibleQuote{没有祂，什么受造的也不曾受造}。如果福音作者相信我们的主是受造之物中的一个，他就不会这样定义。除非它们的制造者拥有与它们不同的本性，因而产生它们而不是祂自己，否则万有怎能藉祂受造并在祂内存立？如果受造界是藉着祂，而祂并非由自己所出，显然祂是受造界之外的某物。在福音作者藉这些话如此清楚地宣告受造之物都是由圣子所造，而非藉任何其他渠道进入存在之后，\Person{保禄}跟随而来，为杜绝这种甚至将圣神也列于受造物之中的亵渎言论，他逐一列举福音作者的话所暗示的一切存在物：正如\Person{达味}在说了\Quote{万有}都置于人之下之后，逐一加上了那\Quote{万有}所包含的每一物种，即地上、水中、空中的生物；同样，\Person{保禄}宗徒，神圣教义的解释者，在说了万有都由祂所造之后，通过列举它们来定义\Quote{万有}的含义。他说\BibleQuote{可见的事物}\BibleRef{哥 1:16}和\BibleQuote{不可见的事物}：藉前者，如我们所见，他给一切感官可认知的事物一个总称；藉后者，他暗示了理知世界。\par

关于第一点，无需详述。没有人会如此属血肉、如此如禽兽，以致想象圣神居于可感世界。但是，在保禄提及了\Quote{那看不见的事物}之后，他继续（为了使没有人猜测圣神，因为祂是属于理智且非物质的界域，因而在此关联中存在）转向另一个最为清晰的分割：按照受造方式被造的事物，以及超越受造的存在。他列举了这些受造理智体的几类：\Quote{\BibleRef{哥 1:16} 上座天使，} \Quote{宰制天使，} \Quote{率领天使，} \Quote{掌权天使，}，用极其广泛的术语传达关于这些不可见影响的道理；但是，正是通过保持静默，他将超越它们的事物从他所列受造物中分开。这好比有人被要求列出某军队中的分区低级军官，在列举了所有之后（十夫长、百夫长、队长、上校，以及所有其他给分部权威的名称），最后却不提统率所有其他军官的最高统帅：并非故意忽略或遗忘，而是因为当需要或打算只列举其下属的各级别时，若将最高统帅列入低级军官名单中将是一种侮辱。保禄正是如此，他曾有一次在乐园中被引入奥秘，当他被提升到那里，成为天上奇观的目睹者，看见并听见了\Quote{人不可说出的言语\BibleRef{格后 12:4}。}这位宗徒打算告诉我们由我们主所创造的一切，并以某些广泛的术语列出它们；但是，在遍历所有天使与超世界之后，他止步于那里，拒绝将超越受造之物拖到受造的层面。因此，圣经中有明确的见证：圣神高于受造界。若有人试图反驳这一点，辩称革鲁宾也未在保禄的列举中出现，他们与圣神同样被遗漏，因此这一省略必须证明要么他们也在受造界之上，要么圣神并不比他们更被认为在受造界之上，就让他衡量列表中每个名字的全部含义：他将在其中发现那些因未被实际提及而看似（也只是看似）被遗漏的东西。在\Quote{上座天使}名下，他包含了革鲁宾，给他们取了这个希腊名称，因为比希伯来名更易理解。他知道\Quote{天主坐在革鲁宾之上：}因此他将这些掌权者称为坐于其上之祂的宝座。同样，依撒意亚的色辣芬也包含在列表中\BibleRef{依 6:6}，当他们发出那奇妙的呼喊\Quote{圣，}时，他们光辉地宣报了天主圣三的奥秘，因他们被圣三每位位格中的美所震慑。他们被大能者保禄和先知达味用\Quote{掌权者}的名号称呼。后者说：\Quote{凡祂的掌权者，你们这些执行祂旨意的臣仆，请赞美主：}，而依撒意亚没有说\Quote{赞美}，却写下了他们赞颂的言语：\Quote{圣、圣、圣，万军的上主，祂的光荣充满大地}。他通过以色辣芬之一对他所做的事揭示了这些掌权者是执行天主旨意的臣仆，按照派遣者之意志施行\Quote{涤除罪过}，因为这就是这些属灵存有的职务，即被派遣为那些得救者的得救。\par

那位神圣的宗徒察觉了这一点。他明白两位先知用不同名称指示同一事物，于是采用了两个词中更常见的一个，称那些色辣芬为\Quote{掌权者}：这样就没有留给我们批评者任何口实，说这些存有中任何一个与圣神一样被排除在受造物目录之外。我们从保禄详述的存有中得知，虽然一些存有被提及，另一些却被略过；并且当他如实地对受造物进行大致计数时，他（在其他地方）也作为单位提及那些被单独设想的事物。因为天主圣三的特性在于它要作为个体被宣报：一位圣父，一位圣子，一位圣神；而上述那些存有则是按群体计数的：\Quote{宰制天使、} \Quote{率领天使、} \Quote{统权天使、} \Quote{掌权天使、}，这样便排除了圣神是其中之一任何猜疑。保禄明智地对我们奥秘保持沉默；他明白，在乐园中听到那些不可言传的话语后，如何在提及较低存有时克制自己不去宣告那些秘密。\par

但这些真理的敌人却冲入不可言说之事；他们将圣神的尊威严降到受造物的层次；他们仿佛从未听说过，天主的圣言在向门徒们传授认识天主的奥秘时，亲自说过，重生者的生命要在他们内完成，并因圣父、圣子、圣神之名而赐予，从而将圣神与圣父及自己并列，使圣神不致与受造物混淆。因此，我们从两方面可以获得对圣神虔敬而正确的认识：从\Person{保禄}在提及受造物时略去圣神的存在，以及从主在提及赋予生命之力时将圣神与圣父及自己联合。这样，我们的理性在圣经的引导下，不仅将独生者，也将圣神置于受造物之上，并促使我们遵照救主的命令，在赐予生命和不受造的存在的福乐世界中，藉信德默观圣神；因此，我们所信的这超越受造物、分享至高和绝对完美本性的统一体，绝不能被视作有任何\Quote{较小}之处，即使这位异端导师试图通过引入程度的概念来限制其无限性，从而通过界定在诸位格中存在较大和较小来收缩天主的完美。\par

% structure: p0116 source=h2 target=subsection reason=relative-heading-rank; base=h2

\subsection{二十四、他在天主圣三内关于\Quote{工程}和\Quote{运行}的程度与差异的详尽叙述是荒谬的。}

现在让我们看看他由此得出的结论是怎样的。在说了我们必须迫使自己将存有视为较大和较小，且一些因卓越的规模和价值而占据领先地位，另一些则因其本性和价值是次要的而必须被贬至较低地位之后，他又补充说：\Quote{它们之间的差异相当于它们工程之间的差异：事实上，说同一运行产生了天使或星辰、天穹或人类，那是不敬的；但人们肯定会对此坚信，即工程越是比别的工程更古老、更尊贵，一个运行就越超越另一个运行，因为运行相同产生工程相同，工程不同表明运行不同。}\par

我怀疑这些话的作者自己也难以说出他写这些话时的意图。它们的思想被修辞的泥浆所遮蔽，厚到几乎无法透过任何线索来解读它们。\Quote{它们之间的差异相当于它们工程之间的差异}这句话可能让人怀疑是来自异教故事中的某个\Person{罗克西亚斯}，在迷惑听者。但如果我们通过我们已经审查过的内容来猜测他此处言论的意图，那么他的意思可能是：如果我们知道一项工程与另一项工程之间的差异大小，我们就会知道相应运行之间的差异大小。但他这里说的是哪些\Quote{工程}，从他的话中无法发现。如果他说的是在受造物中观察到的工程，我看不出这与前文有何关联。因为问题涉及圣父、圣子和圣神：那么，一个理性的人有什么必要一一探究地、水、气、火以及不同动物的本性，并区分一些工程比另一些工程更古老、更尊贵，并说一种运行超越另一种运行呢？但如果他将独生者和圣神称为\Quote{工程}，那么他所说的产生这些工程的运行之间的\Quote{差异}又是什么意思？这位作者的那些奇妙的运行超越其他运行的是什么？他既没有解释他打算让它们如何彼此\Quote{超越}的具体方式；也没有讨论这些运行的本性：而是没有朝任何一个方向推进，既没有证明它们的实际存在，也没有证明它们是一些意志的非实质的运用。自始至终，他的意思悬在这两个概念之间，摇摆不定。他补充说：\Quote{说同一运行产生了天使或星辰、天穹或人类，那是不敬的。}我们再次问道，有什么必要从他的先前言论中得出这个结论？我看不出这一点得到更多的证明，因为运行之间差异如同工程之间一样大，并且并非所有工程都来自同一源头，而是被他陈述为来自不同源头。至于天穹和每一位天使、星辰和人，或者任何被\Quote{受造物}一词所理解的东西，我们从圣经中知道它们都是一位的工程；而在他们的神学体系中，圣子和圣神并非同一位的工程，圣子是那\Quote{跟随}首位存有的运行的工程，而圣神又是那工程的进一步的工程。那么，这个陈述与他随意引入的天穹、人、天使、星辰之间有什么关联，必须由他自己或某个被他引入其深邃哲学的人来揭示。他的话所意图的亵渎是显而易见的，但陈述亵渎的方式前后矛盾。假设在天主圣三内存在着我们可以观察到的、如同包围整个受造物的天穹与其中闪耀的一个人或一颗星辰之间的巨大差异，这是公开的亵渎：但这种思想的联系和这种比较的恰当性对我来说是个谜，而且我怀疑对其作者本人也是如此。事实上，如果他对受造物的描述是这样的：天穹是某种超越性运行的工程，其中的每颗星辰是伴随天穹的运行的成果，然后天使是那星辰的成果，人是那天使的成果，那么他的论证就会由相似过程的比较构成，并可能在一定程度上证实他的教义。但既然他承认万物都由一位所造（除非他想彻底反驳圣经），而同时他以不同的方式描述诸位格的产生，那么这个新引入的观点与前文有什么关联呢？\par

但让我们姑且同意他，这个比喻确实与证明存有之间的变化有关（这正是他想要确立的）；然而，让我们看看他下面的话如何与刚才所说的衔接：\Quote{根据一项工作比另一项更在先、更珍贵，一个虔诚的心灵便会肯定一项行动超越另一项。}如果这是他指涉的可感世界，那么这句话与手头的问题相去甚远。没有任何必要要求一个以神学为主题的人去哲学化地论述世界造成中不同结果出现的秩序，并规定造物主的行动依当时所造之物的体量而高低有别。但若他所指的是那些位格本身，并且以\InnerQuote{更古老、更可敬的}工作意指他刚刚在自己的信经中捏造的那些\InnerQuote{工作}，即圣子和圣神，那么或许最好保持沉默，略过这种可憎的观点，而不是通过与之纠缠而制造出它貌似一种论证的表象。因为，哪里没有较不尊贵的，哪里还能发现\InnerQuote{更尊贵的}？如果他竟能如此轻率地亵渎，暗示\InnerQuote{较不珍贵的}这一表述和概念可以用于我们关于天主圣三所相信的任何事物，那么我们最好掩耳，尽快远离这种邪恶及其论证的感染——这种论证如同从装满不洁的器皿中渗入心灵。\par

有谁敢这样谈论神圣至高的存有，以至于通过论证证明在相形之下有较少的尊荣？福音者说：\BibleQuote{为使所有人尊敬子，如同他们尊敬父。}\BibleRef{若 5:23} 这句话（而如此一句话就是我们的法律）立定了这尊荣上的平等；然而此人却废除了法律及其颁布者，并按他自己发现的某种测量超额丰盈的奥秘方法，将更多的尊荣分给一位，较少的尊荣分给另一位。按照人类的习惯，价值的差异是每位掌权者所受尊崇程度的尺度；因此属下不会以完全相同的方式接近较低的官长和君主，他们表现出或多或少的畏惧或敬重，便表明其对象更值得或不那么值得崇拜：事实上，我们可从属下的这种态度中发现谁是特别可敬的；例如，当我们看到某人比邻人更受畏惧，或比其余人获得更多敬重。但在神圣本性方面，因为每一种善的完美都隐含在天主这名号中，我们——至少就我们所看到的——无法发现任何尊荣等级的依据。哪里在能力、荣耀、智慧、爱情或任何其他可想象的善上没有大小之别，而子所拥有的善也是父的，凡父所有的都在子内看见，什么心态能促使我们在父的情形中显示更多的敬畏？如果我们想到王权和价值，子是君王；如果想到审判者，\InnerQuote{一切审判都交给了子}；如果想到创造的宏伟工程，\InnerQuote{万物都是藉着祂造的}；如果想到我们生命的作者，我们知道真生命曾降到我们的本性中；如果想到我们被带出黑暗，我们知道祂是真光，使我们脱离黑暗；如果智慧对任何人珍贵，基督是天主的德能和智慧。\par

因此，我们的灵魂既然如此自然地、且按各自能力、又如此奇妙地被导向承认基督内如此多且大的奇迹，还有什么额外的尊荣是留给父的，作为不适合子的？人对天主的敬畏，就其最浅显的意义而言，无非是对祂的爱慕态度，以及对祂内完美之事的宣认；我认为\Quote{子应如父一样被尊敬}这条诫命是由圣言代替爱情而颁布的。因为法律命令我们藉着以全心、全力爱天主而向祂献上这合宜的尊荣；而爱情之等价物就在于，圣言作为立法者如此说：子应如父一样被尊敬。\par

伟大的达味所充分献上的正是这种尊荣，他在圣咏的序曲中向主承认他爱慕主，并讲述了他爱慕的一切理由，称祂为\Quote{\InnerQuote{磐石}、\InnerQuote{堡垒}、\InnerQuote{避难所}、\InnerQuote{拯救者}、\InnerQuote{天主的助佑者}、\InnerQuote{希望}、\InnerQuote{盾牌}、\InnerQuote{救恩的角}、\InnerQuote{保护者}}。如果独生子对人类不是这一切，那么就让尊荣的过度按这异端所规定的程度减少吧；但如果我们始终相信祂是这一切，且配得这一切，甚至更多，并在每一行动和对善的概念上都与父之善的尊威平等，那么怎能说是合宜的：要么不爱这样的存在，要么爱祂却轻视祂？没有人会说我们应当以全心全力爱祂，却只用一半来尊敬祂。因此，如果子应以全心被尊敬——即以我们全部的爱献给祂——那么当这样的尊荣以我们全心所能容纳的爱情之币付给祂时，还有什么方法能发现胜过祂所享尊荣的事物呢？因此，在本质上可尊荣的存有身上，任何关于更优越之尊荣的教条化，并由此暗示较低尊荣，都是徒劳的。\par

再者，只有在受造物的情况下，谈论\Quote{先}才是正确的。工程的顺序是按着日子显明的；可以说，诸天先于人的创造，其间隔可以用日子来衡量。但在完全超越时间概念、超出一切思维所及的天主性中，谈论时间上的\Quote{先}与\Quote{后}，只有这种标新立异的哲学才敢如此。简言之，谁宣称圣父比圣子的存有更\Quote{先}，便无异于说：圣子比圣子所制造的万物还晚（如果说所有世代和一切时间都是在圣子之后、并由圣子所造，这话是确实的）。\par

% structure: p0124 source=h2 target=subsection reason=relative-heading-rank; base=h2

\subsection{二十五、凡断言圣父与圣子之间有间隔的\Quote{先}者，必然也承认圣父并非无始。}

但不仅如此：进一步暴露此观点之站不住脚的是，它不仅在圣子的存在上设定了时间上的开端，而且推至极致，连圣父也不放过，反倒证明圣父也有其在时间上的开端。因为任何为圣子的诞生所预设的识别标记，也必然界定圣父的开端。\par

为明此事，更仔细地讨论是好的。当他宣称圣父的生命比圣子的生命更\Quote{先}时，他在二者之间放置了一段间隔；那么，他必定是指此间隔或是无限的，或是包含在固定界限之内。但中项的原理不容许他称之为无限；否则，只要这无限在两端都不受限制，既无圣父的概念在上面截断它，也无圣子的概念在下面阻止它，他就会取消圣父和圣子本身的概念以及任何连接二者的思想。无限的本性正是在两个方向上延伸，并且毫无边界。\par

因此，如果要使圣父和圣子的概念保持稳固不动，他将找不到理由认为这间隔是无限的：他的学派必须将一段确定的时间间隔置于独生子和圣父之间。我所说的便是：他们的观点将使我们得出圣父并非永远存在，而是从时间上的某一点开始的结论。我将用常见的比喻来表达我的意思；以已知之事阐明未知。当我们依据梅瑟的经文说，人是诸天之后的第五天被造的，我们便含蓄地暗示了在那些日子之前，诸天也不存在；随后的一个事件，通过其前面的间隔，也界定了先前事件的发生。如果这个例子没有把我们的论点说清楚，我们可以再举一例。我们说\Quote{梅瑟所赐的法律比亚巴郎的许诺晚了四百三十年}。如果我们一步一步向上追溯先前的时间，到达那个年数终点，我们也就牢牢把握了一个事实：在那日期之前，天主的许诺也不存在。可以举出许多这样的例子，但我不愿琐碎而令人厌烦。\par

因此，以这些例子为引导，让我们审视眼前的问题。我们的对手设想圣父、圣子和圣神的存在分别涉及长幼之分。那么，若按这异端的吩咐，我们越过圣子的受生，接近那纯粹由这些教条主义者的想象所臆造出的、介于圣父与圣子之间的那段延续，然后到达他们用以结束那段延续的另一个至高时间点，我们便在那里发现圣父的生命如同固定在一个顶点上；由此我们必然得出结论：在此之前，圣父并非常常存在。\par

如果你对此仍有困难，让我们再举一个例子。假设有两根标杆，一根比另一根短。如果我们把两根标杆的底端对齐，从顶端我们便知道长出的部分；从与之并排的较短标杆的末端，我们量出这超出部分，通过计算弥补较短标杆的不足，从而使其达到较长标杆的末端；例如一肘，或无论一端到另一端的距离是多少。同样，如果正如我们的对手所言，圣父的生命比圣子的生命有某种超出，那么此超出必然存在于某种确定的延续间隔中；他们也会承认这超出间隔不在将来，因为敌视真理者也会同意两者都是不朽的。不，他们设想这种差异在过去，并且在那里不是使圣父与圣子的生命相等，而是用一个活的间隔延伸了圣父的概念。但每一间隔必为两端所限：因此，对于他们所发明的这个间隔，我们必须把握标示两端的那两点。在他们看来，一部分从圣子的受生开始；另一部分则必须在某一点结束，间隔从此点开始，并以此限定自身。这究竟是什么，应该由他们来告诉我们；除非他们对自己假设的后果感到羞愧。\par

因此无可置疑，他们将完全无法找到与那臆想间隔的初始部分相对应的另一部分，除非他们假设某种他们所谓\Quote{\OriginalTerm{无始者}{Ungenerate}}的开端，以便构想那与圣子生成相连接的中段起点。我们断言，当他藉某种生命间隔的延伸将圣子置于圣父之后时，他必须承认圣父的存在也有一个固定开端，这开端受制于他所设计的同一间隔；如此，他们大肆吹嘘的父之\Quote{\OriginalTerm{无始性}{Ungeneracy}}将被发现遭到其自身拥护者的论证所动摇；他们将不得不承认，他们那无始的天主曾经不存在，而是从一起点开始：事实上，凡存在有开端者，皆非无始。但若我们无论如何都必须承认圣父的无始性，那么就不要在确定圣子生命的某个端点时如此精确，这端点作为祂存在的终结，必将其与另一侧的生命隔开；只要在原因的层面上将圣父视为先于圣子便已足够；而且，不要在圣子生成之前将圣父的生命视为一种分离且独特的存在，以免我们不得不接受由此必然产生的观念，即圣子显现之前有一间隔，它度量那生养祂者的生命，并随之带来必然结果，即也必须假设圣父生命的开端，从而使他们臆想的间隔能在向上进展中停住，为圣父先前的生命也设定一个界限和开端：当我们宣认\Quote{出自祂}时，也请允许——尽管看似大胆——\Quote{与祂同存}；因为圣经的启示引导我们如此相信。我们已从智慧中学会，在那原始之光的永恒性之内并与之一起，默观永恒之光的辉耀，将光辉与其原因合为一念，不承认任何先后。如此，我们将保全我们信仰的教义：圣子的生命在向上观照中不缺失，圣父的永恒性也不因假设圣子有任何确定开端而受侵削。\par

% structure: p0131 source=h2 target=subsection reason=relative-heading-rank; base=h2

\subsection{二十六、不可将上述关于圣父与圣子的概念应用于受造界，如同他们坚称的那样；相反，我们必须将圣子与圣父分开默观，并相信受造界是从一个确定的时间点开始的。}

然而，反对者中或许有人说：\Quote{受造界也有一个公认的开端；然而其中的事物并不在思想上与圣父的永恒性相联，它也不因其自身的开端而阻断神圣生命的无限性——而这正是这场讨论在圣父与圣子问题上所揭示的荒谬结论。因此，两者必居其一：要么受造界是永恒的，要么必须大胆承认，圣子在时间上（较圣父）为晚。时间间隔的概念会导致荒谬的结论，即使是从受造界量到造物主。}\par

如此异议者，或许由于未仔细关注我们信仰的意义，用毫不相干的外来比较来攻击它。若能指出受造界之上有任何事物其起源以某种时间间隔为标记，且所有人都认为在受造界之前存在某种时间间隔，那么他或有机会通过这类攻击来摧毁我们上述证实的圣子之永恒性。但既然所有信友一致同意：万有之中，一部分是受造的，一部分先于受造界，且神圣本性被认为是非受造的（尽管在祂之内，正如我们信仰所教导的，有一个原因，和一个出自原因而无分离的存有），而受造界则被视作处于距离的延伸中——一切事件中时间的秩序与连续性仅能在此世的诸世代中被感知，但先于那些世代的本性却逃避所有先后的区分，因为理性在那神圣而蒙福的生命中看不到它仅在受造界中所观察的事物。受造界，如我们所说，按照秩序的顺序进入存在，并与诸世代的期间相称，以致若沿着受造物之链上溯至其开端，检讨便止于那些世代的基础。但受造界之上的世界，远离一切距离的概念，逃避所有时间顺序：它没有这种开端，也没有终点可据任何可发现的秩序方法停止其进展。穿越诸世代及其中所产生的一切之后，我们的思想瞥见神圣本性，如同无际的汪洋；但当想象向前伸展去把握它时，它自身并不显明任何开端；因此，一个好奇探究诸世代之\Quote{在先}并试图上溯至万有之源的人，将永远无法找到一个可立足的计算点；他所寻求的总是向前移动，不为思想的探索提供任何基础。\par

即使对事物的本性有中等的洞察，也清楚知道，没有任何东西可以衡量那神圣而蒙福的生命。它不在时间之内，而是时间从它流出；反观受造界，始于一个明显的开端，在时间中前行至其适当终点；因此，正如撒罗满在某处所说，我们能在其中发现开端、终结和中间，并以时间的划分来标记其历史的次序。但至高而蒙福的生命并没有时间延伸伴随其进程，因此没有长度或度量。受造之物被限制在适当的尺度内，如同在界限内一样，顾及智慧造物主按其喜悦对整体的良好安排；因此，尽管人的理性因软弱不能完全达到受造界的内容，我们仍不怀疑创造之能已为所有受造之物指定了界限，且它们不超越受造界。然而，这创造之能本身，虽自行限制万物的生长，却没有任何限制的界限；它埋没一切想要上升到天主生命源头的思想努力，并躲避那些忙碌而野心勃勃、试图达到无限之终点的尝试。每一种试图追溯到年代之前的思想的推论努力，只能上升到看见所寻求的永远无法穿越：时间及其内容似乎成了人类思想运动和工作的度量和界限，但那超越之物仍在其触及范围之外；那是一个不可涉足的世界，未受任何人类可理解之物的玷污。那里没有形式、没有位置、没有大小、没有时间的计算，也没有任何可知之物：因此，我们那总是寻求某个可把握之对象的理解官能，必然从这不可理解之存在的任何一侧退回，并在年代及其所包含的受造界中寻找其同类的和相宜的领域。\par

凡对事物的本性有某种洞察——即使是中等程度的洞察——的人都知道，世界的创造主设立了时间和空间作为背景，以接纳将要存在之物；祂在这基础上建造宇宙。任何借着受造方式已经存在或正在形成之物，都不可能独立于空间或时间。但是那全备、永恒、包罗世界的存在，不在空间内，也不在时间中：它以不可言喻的方式在它们之前、在它们之上；自我含容，唯独借着信德可知；不以年代衡量；无时间相伴；安坐并安息于自身之内，没有过去或将来的关联，因为除祂自身外无物旁侧或超越，其经过不会造成过去和将来。这样的属性只限于受造界，其生命随着时间的划分被分为记忆和希望。但在那超越而蒙福的大能之内，万有如同在一瞬间同等临在：过去和将来都在其全包围的掌握和综合视野之内。\par

这存在，用宗徒的话来说，万有在其中形成；我们各自分享存在的部分，生活、行动、存在。祂超越开端，不呈现其最深本性的任何标记：唯独在无法察觉祂的可能性中才可被认知。这确实是祂最特殊的特征：祂的本性太高，无法有任何区分性属性。对受造物必须给出与不受造者非常不同的描述：正是这一点使受造物与造物主之间没有任何比较和关联；我的意思是，本质的这种差异，以及容许一种特殊的解释来说明其本性的属性，这解释与造物主的本性毫无共同之处。天主性对这些受造物中的特殊标志是陌生的：它把时间的分段——'之前'与'之后'——以及空间的观念留在自己之下：事实上，'更高'根本不能恰当地说它。关于那不受造大能的每一个概念都是一个崇高的原则，并且包含最高程度上适当之物的观念。\par

因此，我们通过所说的话表明，不可在受造界中寻求独生者和圣神，而应相信它们超越受造界；并且，虽然受造界或许可以通过野心搜寻者的坚持不懈的努力，在其自身的开端（无论是什么）中被把握，但超自然之物并不会因此进入知识的领域，因为找不到在年代之前指示其本性的任何标志。那么，如果在这不受造的存在中，那些奇妙的实在连同它们奇妙的名称——圣父、圣子和圣神——要在我们的思想中，我们怎能想象，在那种先于时间的世界中，我们的忙碌、不安的心灵在此下事物中所察觉的东西，即通过比较它们彼此并用时间间隔赋予优先地位？因为在那里，与圣父——无始、无生、永远是圣父——一起，圣子作为出自祂并与祂并列的观念被不可分离地结合；并通过圣子并与圣子一同，在任何模糊而空泛的观念介入之前，圣神立即在以最紧密的结合中被发现：在存在上不后于圣子，好像圣子曾被设想为没有圣神一样；但祂自己也拥有同一存在的原因，即统御万有的天主，如同独生之光，并在那同一光中照耀出来，既不被时间延长，也不被外来本性从圣父或独生者分离。在那先于时间的世界中没有间隔：在存在方面也没有差异。甚至不可能将不受造者与不受造者比较而看到差异；而圣神是不受造的，正如我们先前所表明的。\par

既然所有以纯朴之心接受未掺杂的福音的人都持有这种观点，那么有什么必要通过受造物来试图分解圣子与圣父之间这种紧密的结合呢？仿佛必须假设圣子要么与受造物一同从永远存在，要么他也同样比受造物晚出现？因为圣子的受生并不落于时间之内，正如受造物也不在时间之前：因此，以任何方式分割那不可分割的，并通过宣称曾有一段时期万有的创造者不存在，而将这个虚假的时间观念插入宇宙的创造之源，是绝对不合理的。\par

因此，我们先前的主张是正确的，即圣子的永恒性连同其受生的概念一同被包含在圣父的非受生之内；而且，倘若人们设想某种间隔将两者分开，那么这同一间隔将为全能者的生命设定一个开端——这是一种荒谬的假设。但受造物按其本性，原本就与造物主不同，并且绝不触及那纯洁的前时间领域，因此，正如我们所说的，相信它有自身的开端并无不妥。说天、地及受造物中的其他内容是由不存在的材料，或如宗徒所说，由\Quote{看不见的东西}而来，这无损于宇宙的创造者；因为圣经告诉我们，这一切都不是永恒的，也不会永远存续。另一方面，若能相信天主圣三之中有某物不与圣父共存，若依从这种异端而设想剥去全能者与圣父的荣耀，以及圣神和圣子的荣耀，则结果只能是一位明显远离一切良善和神圣行为与思想的天主。但是，倘若在万世之前就已存在的圣父常享荣耀，且前时间的圣子就是祂的荣耀，同样，基督的圣神也是圣子的荣耀，常与圣父和圣子一同受默想，那么这位具学识的人经过什么培养竟会宣称，在无时间的事物中有\Quote{之前}，在本质上全部可尊敬的事物中有\Quote{更可尊敬的}，并通过比较来偏爱一方，进而以这种偏袒来羞辱另一方呢？\Quote{更可尊敬}的对立面更清楚地显示了他的倾向。\par

% structure: p0140 source=h2 target=subsection reason=relative-heading-rank; base=h2

\subsection{二十七、他错误地认为，同样的能力产生同样的作品，而作品的变化表明能力的变化。}

与此同类的是他在下一段中所补充的：\Quote{相同的活动产生相同的工程，而不同的工程也表明活动之间的差异。}这位高尚的思想家为其学说辩护得巧妙而无可辩驳。\Quote{相同的活动产生相同的工程。}让我们用事实来检验这一点。火的活动总是一样的，就是加热；但它的结果显示出何种一致性？铜在其中熔化，泥在其中变硬，蜡消失；所有其他动物都被它毁灭，而火蜥蜴却得以存活；亚麻燃烧，石棉被火焰洗涤如同用水；这就是他所谓的\Quote{来自同一种活动的工程相同性。}太阳又如何？它的热力不总是一样的吗？然而它使一株植物生长，却使另一株枯萎，根据每种植物的内在力量改变其运作的结果。\Quote{那落在岩石上的}枯萎了；\Quote{那落在深土里的}结出百倍的果实。探究自然的工作，你就会明白，在他所称的\OriginalTerm{相同的活动产生相同的结果}{sameness of energy effects sameness of result}这一说法中，自然艺术性地产生的那些身体有多少精确性。一个单一的操作是受孕的原因，但其中内在所成就的组合是如此多样，以至于任何人都很难甚至数算身体的所有不同品质。同样，婴儿吸吮乳汁是一个单一的操作，但由此滋养的结果却复杂得无法一一详述。当这食物从口腔的通道进入分泌管时，自然的转化力将其按需输送到各个部分；因为通过消化，她将总量分解为无数差异的细微变化，并供应与她所处理的材料相宜的养分；因此，同样的乳汁滋养动脉、静脉、脑及其膜、骨髓、骨、神经、筋腱、肌腱、肉、表面、软骨、脂肪、毛发、指甲、汗液、蒸气、痰、胆汁，此外，还有源自同一来源的一切无用废物。你无法说出任何一个器官，无论是运动还是感觉的，还是构成身体体积的任何其他东西，不是（尽管有惊人的差异）由这一相同的滋养活动所形成的。如果也将机械技艺进行比较，就会明白他的陈述的科学价值；因为在它们中我们看到相同的活动，即手的运动；但结果有什么共同之处？建造圣殿与制作外套有什么关系，尽管两者都使用手工劳动？破门而入者和挖井者都移动他们的手；挖掘泥土和杀人都是手部运动的结果。士兵用他的手杀死敌人，农夫用他的手挥动叉子打碎土块。那么，这位教条主义者如何能断定\Quote{相同的活动产生相同的工程}？即使我们承认他的这一观点有几分真理，圣子与圣父，以及圣神与圣子的本质合一，仍然更加充分地得到证明。因为如果他们的活动存在任何变化，以至于圣子以与圣父不同的方式实行祂的旨意，那么（在上述假设下）就有理由根据这种变化来猜测作为这些不同活动结果的存在物也有变化。但是，如果圣父的运作方式同样也是圣子运作的方式，这是根据我们主自己的话以及\Emph{从先验的角度}我们可以预期的——（因为一个不是无形体的而另一个是有形体的，一个不是来自这种质料而另一个来自那种质料，一个不是在此时此地实行祂的旨意而另一个在彼时彼地实行祂的旨意，它们之间也没有器官的差异导致结果的差异，而是它们意愿和意志的单一运动就足够了，在创立宇宙时由能够创造万物的能力辅助）——如果，我说，在一切方面，万物所从出的圣父和万物藉以造成的圣子在运作的实际方式上相同，那么，这个人如何能指望通过圣子和圣父运作的任何差异和区分来证明圣子和圣神之间的本质差异？正如我们刚才所看到的，恰恰相反的情况得到证明；既然在圣父和圣子的运作之间看不出任何差异；那么，圣子的存在与圣神的存在之间没有任何鸿沟，这由赋予它们本体的能力的同一性所表明；而这位小册子作者自己也证实了这一点；因为这是他的\Emph{原话}：\Quote{相同的活动产生相同的工程。}如果工程的相同性确实是由活动的相似性产生的，并且如果（如他们所说）圣子是圣父的工程，圣神是圣子的工程，那么圣父和圣子活动方式的相似性将证明这些各自结果的存在物是相同的。\par

但他补充说：\Quote{工程的变化表明活动的变化。}他的这一论断又如何得到事实的证实？请看，如果您愿意，举几个简单的例子。那位凭祂的独一意志创造世界及其中万物的天主的命令的\Emph{活动}，不是一个单一的活动吗？\BibleQuote{祂发言，万物形成；祂下令，万物受造。}在每种情况下，所命令的不是同样地被赋予存在吗？祂单一的意志不足以使无中生有吗？那么，当从这单一的命令活动中看到如此巨大的差异时，这个人怎能对现实视而不见，并宣称工程的差异表明活动的差异？如果我们的教条主义者坚持这一点，即工程的差异意味着活动的差异，那么我们就应当期望与实际情况完全相反；也就是说，世上的一切都应属于同一类型。难道他在这里看到了普遍的相似性，而只在圣父与圣子之间发现了不相似吗？\par

让他观察，若他从前未曾注意到，世界元素之间的不相似性，以及组成宇宙框架的每一事物如何依附于其天然对立面。有些物体轻而浮，有些重而下沉；有些永远静止，有些永远运动；在后者中，有些按一种不变的计划运动，如天空与行星，其运行都与宇宙反向旋转；其他则向所有方向渗透并随意冲撞，例如空气与海洋，以及每一种天生具有穿透性的物质。何须提及热与冷、湿与干、高与低位置之间的对照？至于动物与植物在形状、大小以及其产物与品质的一切变异上的无数不相似性，人的心智将无法穷尽。\par

% structure: p0144 source=h2 target=subsection reason=relative-heading-rank; base=h2

\subsection{二十八、他错误地想象我们可以拥有一系列不变的和睦性本性并列存在。}

但这位博学之士仍然宣称，多样的作品需要多样的能力来产生。他要么还不明白圣经所教导的天主能力的本性——\BibleQuote{万有都是因祂的命令之言而造成的}——要么就是对现有事物的差异视而不见。他为了我们的益处说出这些轻率的言论，并在神圣教义上制定法律，仿佛他从未听说过：任何仅仅被断言的东西——若关于讨论的事物没有完全不容置疑的明确陈述，且断言者自行说出谨慎的听众无法赞同的话——无异于说梦或酒后故事。因此，他的这个论断虽不符合事实，然而——就像一个人被梦境迷惑，以为看到了他清醒时努力追求的对象，急切地抓住这个幻影，被这虚幻的欲望欺骗的眼睛以为抓住了它——他面前有这教义的梦影，便以为自己的话有力，坚持其真理，并试图借此证明其余一切。值得引用这段话。\Quote{情况既然如此，并且它们之间保持着不间断的联系，那么，对于那些按照与主题相宜的顺序进行研究、不坚持将所有事物混合混淆的人，若讨论涉及存有，似乎宜通过首要的能力和附着于存有的能力来证明正在论证的东西，并解决争论点；而当涉及能力时，则再次通过存有来解释。}我认为，他这些不敬的话本身足以证明这种教导是多么荒谬。如果有人要描述某种疾病如何毁坏人的面容，他最好实际解开病人的绷带，这样眼睛看到他的模样后便无需言语。同样，某种心灵的眼睛可以觉察到这异端所造成的可怕残缺：只需阅读即可揭开面纱。但由于有必要使这教导的潜在危害对众人显明，我必须用手指指出它，我将再次重复每个字。\Quote{情况既然如此。}这个做梦者是什么意思？\Quote{如此}是什么？它是如何陈述的？\Quote{圣父的存有是唯一固有的、至高的；因此，其次的存有是依赖的，第三个更依赖。}他用这样的话制定法律。但为什么呢？是因为首要的存有伴随着一种能力，其效果和工作——独生子——被限定在这产生原因的范围之内？还是因为这些存有被认为有大小之分，较小的被较大的包含和包围，像桶一个套一个，因为他在无限存有中发现了大小的程度？还是因为产品的差异暗示了生产者的差异，仿佛不可能用相似的能力产生不同的效果？好吧，没有一个神智如此沉睡的人会在听到这样的陈述后立即同意接下来的断言：\Quote{情况既然如此，并且它们之间保持着不间断的联系。}说这样的话以及不加质疑地听这样的话，同样是疯狂。它们被置于\Quote{系列}和\Quote{彼此不变的关系}中，却又因本质的不相似而彼此分离！要么，如我们自己的教义所坚持的，它们在存有上合一，那么它们确实保持着彼此不变的关系；要么，如他所幻想的，它们在本质的不相似中彼此分离。但外来的实体之间能有什么系列和不变的关系呢？它们又如何呈现那\Quote{与主题相宜的秩序}——按照他的说法，这种秩序应主导研究？如果他只着眼于真理的教义，他计算差异的次序只是信德在天主圣三中所见的属性——一种如此\Quote{自然}和\Quote{相宜}的秩序，使位格不能混淆，位格虽分开，但存有合一——那么他根本不会被视为我们的敌人，因为他所意指的正是我们教导的同一教义。但实际上，他看的正相反，他使他在\Emph{那里}所幻想的次序变得完全不可理解。意志行为的完成与机械自然法则的完成之间有天壤之别。热是火固有的，光芒是阳光固有的，流动性是水固有的，下落趋势是石头固有的，等等。但如果一个人建造房屋、谋求官职、载货出海或尝试任何需要思考和准备才能成功的事，我们不能说在这种情况下，他的运作中固有某种等级或秩序：它们的次序在每个情况下都是指导他选择的动机或他所成就的事物的效用之后的结果。那么，既然这异端使圣子与圣父的任何本质关系分离，并对圣神也采取同样的看法，即圣神与圣父或圣子无关，并且它始终断言圣子是圣父的工程，圣神是圣子的工程，且这些工程是意志的结果，而非本性的结果，他有什么根据宣称这意志的工程是\Quote{事物固有的秩序}？这种教导的意图是什么？它使全能者制造了圣子和圣神中如此的本性，使超然的存在彼此低劣？如果真是他的意思，他为什么不清楚说明他凭什么假定在天主身上，结果的小将是更大能力的证据？但有谁会真正允许从结果的小中寻找伟大而有力的原因？好像天主无法在自己所出的事物中建立完美！他又如何能赋予天主至高的至高特权，同时显示祂的能力如此不及祂的意志？欧诺米肯定似乎认为，完美甚至不是天主工作的目标，因为怕受敬拜者因其优越而享有的荣耀和尊崇因此减少。然而，有哪个心胸狭窄的人会认为有福的天主本身也免不了嫉妒之情？那么，还有什么合理的理由使至高天主在圣子和圣神中构成了这样的\Quote{秩序}？\Quote{但我并不是说那\InnerQuote{秩序}来自祂，}他回答说。但如果不具备这种\Quote{秩序}的存有彼此没有本质关系，它又从何而来？也许他称圣子和圣神存有的低劣本身为这种\Quote{固有秩序}。但我恳请他告诉我这事的理由，即为什么圣子在存有方面低劣，而在这存有和能力上却具有相同的特征和属性？另一方面，如果存有和能力的定义不同，各自表示不同的东西，他为什么要用完全不同的事物来证明所讨论的事物？这就像在讨论人的存有本身——他是否是能笑的动物，或是能接受识字教育——有人引述木匠建造房屋或造船匠建造船只作为解决方式，坚持巧妙的推理：我们通过运作知道存有，而房屋和船只是人的运作。那么，最单纯的先生，我们通过这些前提是否得知人是能笑的，同时也是宽指甲的？有人也许会反驳：\Quote{人是否拥有运动和运作不是问题所在：问题是运作原理本身是什么；而我从你的决定方式中无法得知。}的确，如果我们想知道风的本性，你不会通过指向被风扬起的沙堆或糠秕，或吹散的尘土来给出满意的回答：因为对风的说明完全不同；你的这些例证与问题无关。那么，他有什么根据试图通过能力来解释存有，并根据实体的结果来定义实体？\par

我们也要看看，按他的说法，圣父的存有借以被认知的，究竟是圣父的什么工程。他若照常说话，想必会说：最确实的，就是圣子。但是，最博学的先生，按您的体系，这位您的圣子只与产生祂的那个能力大小相当，且只表明那个能力；而我们探求的对象仍隐于暗处，因为您自己也承认，这能力只是\Quote{随之而来}于第一存有的事物之一。这种能力，如您所说，延展到它所产生的工作中，但并未在其中揭示其自身的本性，只揭示我们能在该工作中窥见的那一部分。一个铁匠不会动用全部手艺来造一枚锥子；这位工匠的技术只运作到足以造出此工具的程度，尽管他能制造多种其他工具。因此，能力的界限可在它所产生的工作中找到。但目前的问题不是关于能力的多少，而是关于发出能力者的存有。同理，若他声言能在圣神（他称之为\Quote{随之}于圣子的能力的工作）中感知独生者的本性，他的断言是无根据的；因为这里又是能力延展到其工作中，却没有在其中揭示自身或运用此能力者的本性。\par

但让我们让步；姑且允许他：存有是在其能力中被认知的。第一存有通过祂的工作而被认识；第二存有则在从\Emph{祂}发出的工作中显明。但是，我博学的朋友，什么可显示这第三存有？找不到这种第三存有的工作。如果你坚持存有是通过其能力被感知的，你必须承认圣神的本质是不可感知的；你不能从祂为延续连续性而发出的任何能力推断祂的本性。请指出圣神任何实质性的工作，通过它您认为已察觉到圣神的存有，否则您所有的蛛网都将被理性一触即溃。如果存有是由随后的能力所认识，而圣神又没有实质性的能力（如您说圣父在圣子中显示，圣子在圣神中显示的那种），那么圣神的本质必须被承认是不可知的，且不能通过这些能力被理解；没有与本体相关联的能力被设想出来，甚至无法瞥见它一丝。但如果圣神逃避理解，那么凭借本身不可知的东西，如何能感知更崇高的存有？如果圣子的工作——即按他们的说法，圣神——是不可知的，那么圣子自身永远不能被认识；祂将陷入那证明祂者的模糊之中；而如果圣子的存有就这样被隐藏，那么那位最真正是存有、最崇高的存有，如何能通过本身被隐藏的存有而被显明？圣神的这种模糊通过回溯经过圣子传递给圣父；因此，按照这个观点，即使我们的对手也承认，圣父的存有的不可知性被清楚显明。那么，这个人无论他的眼睛多么\Quote{敏锐地看见无形实体}，怎能凭借自身辨认那不可见、不可理解者的本性？他怎能命令我们通过它们的工作把握存有，再通过它们把握工作呢？\par

% structure: p0148 source=h2 target=subsection reason=relative-heading-rank; base=h2

\subsection{二十九、他虚妄地以为，关于能力的疑问可由存有解决，反之亦然。}

现在让我们看看接下来是什么。\Quote{关于能力的疑问要由存有来解决。}有什么办法能把这个人从虚妄幻想中带回常识？如果他认为通过理解存有本身就能解决关于能力的疑问，那么若这些存有未被理解，他如何能将这疑问转变为任何确定性？如果存有已被理解，何必把能力弄到如此重要，好像\Emph{它}能引导我们理解存有？但若正是这一点使对能力的考察成为必要，即我们借此被引导去理解发出能力者的存有，那么这尚不为人知的本性如何能解决关于能力的疑问？任何存疑之事的证明必须借由已知的真理进行；但当我们所探求的两个对象都同样不确定时，欧诺弥怎能说它们彼此借由对方而被理解，两者本身都超出我们的认识？当讨论圣父的存有时，他告诉我们可借着随其而来的能力和这能力完成的工作来解决疑问；但当探讨独生者的存有时（无论欧诺弥称祂为能力还是能力的产物，因为他两者都做），他则告诉我们，只要看着祂的生产者的存有就容易解决！\par

% structure: p0150 source=h2 target=subsection reason=relative-heading-rank; base=h2

\subsection{三十、没有任何圣言命令这类探究：由此证明从事这种哲学的无用。}

我还想问问他：他是否只认为在天主性方面，行动由产生它们的存在者所解释？还是他承认，对于任何具有有效力量的事物，都可以通过产生者的存在来认识被产生者的本性？如果只有在天主性方面他持这种观点，那么请他指出，他如何通过匠人的本性来解决关于天主工程的问题。以天主的无容置疑的工程为例：天、地、海，整个宇宙。假设我们正在探究其中之一的存在者，并选定\Quote{天}作为我们推理的主题。问题是：天的本体是什么？对此，各位自然哲学家根据各自的理解提出了广泛不同的见解。那么，对天的制造者的默观如何能解决这个问题呢？那位制造者是非物质的、无形的、无形式的、无生的、永恆的、不朽坏的、不变易的、不改动的，以及其他一切如此的性质。任何怀有这种对匠人的观念的人，如何能引领自己认识天的本性呢？他如何从无形者中认识有形者？从不朽坏者中认识可朽坏者？从无始存在者中认识有始存在者？从永恆者中认识暂存者？事实上，他从与所寻对象完全相反的一切中认识所寻的对象？让这位深入洞察事物奥秘的人告诉我们，两个不同的事物如何能彼此认识。\par

% structure: p0152 source=h2 target=subsection reason=relative-heading-rank; base=h2

\subsection{三十一、通过观察天主眷顾所做的观察足以给予我们认识存在同一性的知识。}

然而，如果他能够看见自己陈述的后果，他就会因此被引导去同意教会的道理。因为如果制造者的本性是指示被造物的标志，如他所肯定的，并且按他的学派，圣子是天父所制造的被造物，那么任何观察到圣父本性的人必定由此认识圣子的本性——如果确实如他所说，工匠的本性是其所制造之物的标志。但按他们的说法，独生子因与圣父的\Emph{不相似}，将被排除在通过天主眷顾施行作为之外。欧诺米毋需再为祂的被生而烦恼，也毋需从中强行推出圣子不相似的另一个证据。意愿的差异本身就足以揭示祂异类的本性。因为即便是我们的对手也承认，第一存在是单一而唯一的，祂的旨意必然被想为遵循其本性的倾向；但天主眷顾显明旨意是良善的，因此产生那旨意的本性也显明是良善的。所以圣父单独行善；而圣子不与祂同怀同一旨意，如果我们采纳对手的假设；那么，他们本性的差异就会由他们旨意的这种变化得到清楚的证实。但是，如果圣父眷顾宇宙，圣子也同样眷顾宇宙（因为\Quote{凡祂看见圣父所做的，圣子同样也做}），他们旨意的同一性表明那些怀有同一旨意者的本性共融。那么，为什么他完全忽略对天主眷顾的提及，好像它无助于我们寻找的东西？然而，许多熟悉的例证支持我们的看法。凡曾注视过火的明亮并体验过其温暖能力的人，当他接近另一同样的明亮和温暖时，必定会被引导去想火；因为他的感官通过这些类似的现象会引领他认识到产生两者的同类元素；不是火的任何东西不能像火那样总是在所有情况下工作。同样，当我们察觉圣父和圣子中相似且同等的眷顾能力时，我们便通过由此进入我们认识范围的事物，对我们无法理解的事物做出猜测；我们感到，在同等和相似的效果中，不能发现异类的原因。如同观察到的现象彼此相关，那些现象背后的主体也是如此：如果前者彼此相反，我们必须认为所揭示的实体也是如此；如果前者相似，那么后者也一定相似。我们的主曾用比喻说，它们所结的果子是树木品格的标志，意思是这标志不欺骗那品格：坏树不结好果子，好树也不结坏果子——\BibleQuote{你们可凭他们的果实辨别他们}——所以，当果实即天主眷顾呈现无差别时，我们便发觉产生这果实的单一本性，即使结出果实的树木可能不同。因此，通过我们认知所能把握的，即在圣子身上与圣父身上同样可见的天主眷顾计划，独生者与圣父的共同相似性便无可置疑；而且我们正是通过天主眷顾果实的同一性认识这一点的。\par

% structure: p0154 source=h2 target=subsection reason=relative-heading-rank; base=h2

\subsection{三十二、他所说的\Quote{相似的方式必须随生成的方式}是不可理解的。}

但是，为了防止有人产生这种想法，并假装被某种方式迫使他离开它，他说他退出所有这些天主眷顾的结果，回到圣子生成的方式，因为\Quote{祂相似的方式必然跟随祂生成的方式。}多么无可辩驳的证明！这种措辞如何强有力地迫使人们赞同！这一论断的措辞何等巧妙和精确！那么，如果我们知道生成的方式，我们就由此知道相似的方式。好吧；既然所有——或者至少大多数——通过分娩出生的动物都有相同的生成方式，而根据他们的逻辑，相似的方式跟随这种生成的方式，那么这些动物，既然它们在生产中遵循同样的模式，就会完全相似那些同样生成的动物；因为相似于同一事物的东西彼此相似。那么，按照这个异端观点，生成的方式使每一生成物恰好与自己相似，而事实上，这种方式在多样化的动物种类中并无变化，在绝大部分动物中保持不变，我们就会发现，他这种笼统无条件的论断，凭借这种出生的相似性，建立了人、狗、骆驼、老鼠、大象、豹子以及自然以同样方式产生的每一种动物之间的相互相似。或者他的意思不是指以相似方式来到世界上的东西都彼此相似，而是指它们每一个都只相似于那给予它生命的源头？但如果这样，他本应宣告孩子相似父母，而不是\Quote{相似的方式}相似于\Quote{生成的方式。}然而，这个在本身如此可能、并且在自然中被观察到的事实——即所生者相似生育者——他却不承认为真理；这会使他全部论证变成与他本意相反的证据。如果他允许后裔相似父母，那么他辛苦积累的用以证明实体\Emph{不相似}的论据就会被驳斥为转瞬即逝和毫无根据的。\par

所以他说\Quote{相似的方式跟随生成的方式。}若由一个精确批评任何观念意义的人来检验，这一点将完全不可理解。显然不可能说出\Quote{生成的方式}可能意味着什么。它是指父母的形像？还是他的冲动？还是他的性情？还是时间、地点？还是受孕完成胚胎？还是生育的器皿？或者不是这类东西，而是\Quote{生成}中所观察到的其他东西？要找出他的意思是不可能的。\Quote{方式}一词的不当和模糊导致对其在此处的含义感到困惑；每一种可能的猜测都同样开放，同时也与当前主题缺乏同样多的联系。同样，他所说的\Quote{相似的方式；}没有任何意义的痕迹，如果我们把注意力集中在我们熟悉的例子上。因为所生成的东西在那里并不被比作它出生的种类或方式。就动物出生而言，出生在于身体与身体的分离，其中在子宫中完全形成的动物被生出来；但所生的东西是一个人、马、牛，或任何它因出生而存在的东西。因此，\Quote{后裔相似的方式跟随其生成的方式}必须由他来解释，或者由他某个在助产术方面的学生来解释。出生是一回事；所生的东西是另一回事：它们完全是不同的概念。任何有理智的人都不会否认，他在动物出生的情况下所说的完全不真实。但是，如果他称实际的制造和实际的成形为一种\Quote{生成的方式，}所造之物的\Quote{相似的方式}要\Quote{跟随，}那么即使如此，他的陈述也远离一切可能性，正如我们将从一些例子中看到的那样。铁被工匠的锤击打成形为某种有用的工具。那么，它的边缘的轮廓——如果碰巧有这样的话——如何能被说成是相似于工人的手，或者相似于其制造的方式——例如，锤子、煤炭、风箱和铁砧，工匠借助它们来塑造它——没有人能解释。一种情况成立，所有情况都成立；任何产生结果的操作，都不能说所造之物相似于\Quote{其生成的方式。}一件衣服的形状与梭子、杆子、梳子或织布工具的形式有什么关系？实际的座位与木块的操作有什么关系？任何成品与其制造者的体格有什么关系？——但我认为甚至我们的对手也会承认，他的这条规则在可感觉的和物质的事例中并不适用。\par

尚需查看：这对其亵渎之证明是否还有进一步的助益。他意在何为？即：必须按照圣子与圣父的相似或不相似来相信；而既然我们无法从上智安排的角度认识这种\Quote{是}，就必须诉诸\Quote{生成的方式}，由此我们得以知道——并非绝对知道受生者与生者相似，而仅知道二者之间某种\Quote{相似的方式}；又因这方式对多数人而言是隐秘的，就必须较为详尽地探讨生者的\Quote{是}。那么，他岂非忘记了自己关于从工程认识\Quote{是}的定义？但这受生的\Quote{是}——他称之为至高\Quote{是}的工程——（依他之见）尚未被光照；那么又如何处理其本性？他又怎能\Quote{超越这较低因而较直接可理解者}，从而持守于绝对和至高的\Quote{是}？此外，他在整篇论述中自诩对神圣言词有精确认识；然而在此却对其缺乏敬意，忽视了一个事实：除非通过圣子，否则无法接近对圣父的认识。\BibleQuote{除了子，以及子愿意启示的人，没有人认识父。}\BibleRef{玛 11:27}然而，欧诺米乌斯在每一次可以侮辱我们虔诚而崇拜天主的圣子观念时，都明确断言圣子的卑下，却在这认识神性的策略中不知不觉地肯定了圣子的优越；因为他假定圣父的\Quote{是}更易于为我们理解，然后试图由此追溯并论证圣子的本性。\par

% structure: p0158 source=h2 target=subsection reason=relative-heading-rank; base=h2

\subsection{三十三、他虚假地宣称\Quote{生成的方式可从生者内在的价值得知}。}

例如，他回到受生的\Quote{是}，并由此审视受生者；\Quote{因为,}他说，\Quote{生成的方式可从生者内在的价值得知。}再次，我们发现他这大胆的无限制概括使得探究者的思想向各种可能方向消散；这类一般性陈述的本性是：其含义扩展到每个实例，不容任何事物逃脱其广泛断言。如果\Quote{生成的方式可从生者内在的价值得知}，而根据可发现的众多分类，生者的价值有许多差异（因为有人生为犹太人、希腊人、野蛮人、西徐亚人、为奴的、自由的），结果会怎样？那么，我们必须期望有多少生者价值上的差异，就有多少\Quote{生成的方式}；他们的出生不会都以同样方式实现，而是其本性会随父母的价值而变化，并且根据这些不同的评价，每种出生方式都会被设计出来。因为每个个体父母身上可观察到某种不可剥夺的价值；即根据种族、评价、宗教、国籍、权力、奴役、财富、贫困、独立、依赖等等构成终生价值差异的因素，而有较好或较差的区别。如果\Quote{生成的方式}由父母的内在价值显示，而价值有许多差异，那么跟随这位意见贩卖者，我们将不可避免地发现生成的方式也各不相同；事实上，这种价值的差异将指令本性采用出生的方式。\par

但若他不承认这种价值是本性的，因为它们可以在思想中被置于其主体本性之外，我们并不反对他。但无论如何他会同意：人的存在以其内在特征与禽兽的存在相区别。然而这两种情况下的出生方式在内在本性上并无差异；本性以完全相同的方式（即通过生育）将人和禽兽带到世上。但如果他只在最恰当和至高的存在中领会这种本然尊贵，那么让我们看看他究竟意指什么。依我们之见，天主的\OriginalTerm{本然尊贵}{native dignity}在于天主性本身、智慧、能力、美善、审判、正义、力量、慈悲、真理、创造力、宰制、不可见性、永恒性，以及圣经中为荣耀祂而提及的一切其他属性；我们宣称，这些属性每一项都恰当且不可剥夺地存在于圣子中，只承认在无始性方面有区别；甚至那区别，我们也不根据其\Emph{所有}含义将圣子排除在外。但不要让吹毛求疵的批评者攻击这一陈述，仿佛我们试图表明圣子本身是非受生的；因为我们认为，坚持这点的人与异类派同样不虔敬。然而，既然\OriginalTerm{起源}{origin}的含义有多种，并暗示许多观念，其中有些含义下，\OriginalTerm{无始的}{unoriginate}这一称号并非不适用于圣子。例如，当这个词具有\OriginalTerm{从任何原因都不获得存在}{deriving existence from no cause whatever}的意思时，我们承认这是圣父特有的；但当涉及\Quote{起源}的其他含义时（因为任何受造物、时间或秩序都有起源），我们将超越起源归于圣子，并相信那造万物的祂超越受造的起源、时间的观念和次序的顺序。因此，祂就其自立而并非没有起源，但在其他一切方面具有不容置疑的\Emph{无始性}；圣父是无始的、非受生的，而圣子在如上所述的方式下是无始的，尽管不是非受生的。\par

那么，圣父那本有的尊荣是什么，他将要察看它，从而推断出\Quote{受生的方式}？\Quote{当然就是祂的非受生，}他会回答说。那么，\Person{欧诺米}啊，如果我们所学会用来颂扬天主荣耀的所有那些名号，在你看来都是无用和无意义的，那么只是列举这些表达就是一项无益且多余的任务；你说，这些别的词没有一个表达了那超越万有的天主的固有价值。但是，如果这些词每一个都有一种特别的效力，适合我们对天主的观念，那么天主的固有尊荣必然要与这个列举联系起来看待，而两个存有之间的相似性也将由此得到证明；也就是说，如果那不可与存有分离的特性是那些特性之主体的标记。每个存有的特性被发现是相同的；那么这两个具有相同尊荣的主体在存有上的同一性就被最清楚地显示出来。因为如果一个名称的变异就要被视为一个异质存有的标记，那么这无数名称的同一性更应该有助于证明本性的共融！\par

那么，为什么其他的名号全都被忽略，而受生却唯独通过一个名号来表示？为什么他们宣称这个\OriginalTerm{非受生性}{Ungeneracy}是圣父唯一的内在特性，而把其余一切都搁置一旁？这是为了藉着与受生者的对比，来确立他们那有害的圣父与圣子相异的方式。但是我们会发现，当我们到适当的地方检验时，他们的这个企图与之前的企图一样软弱、无根据且无价值。\par

尽管如此，他所有的推理都指向这一方向，这由下文表明。在文中他称赞自己恰当地采用了这种方法来证明他的亵渎，却又没有一下子泄露他的意图，也没有在他那欺骗性论证的连锁完成之前，用他的不敬来震惊毫无防备的听众，也没有在他的论述早期就把这个\OriginalTerm{非受生性}{Ungeneracy}展示为天主的存有，也没有用关于存有差异的谈论来使我们厌倦。以下是他的原话：\Quote{或者，正如\Person{巴西略}所命令的，从要证明的事情开始，而不连贯地断言非受生性就是存有，并谈论本性的差异或相同，这是正确的吗？}然后他插入了很长的一段话，由嘲笑和侮辱性的辱骂组成（这就是这位思想家用来捍卫他自己学说的武器），然后他恢复论证，转而攻击他的对手，把他所说的话的责任归咎于对方，用这些话：\Quote{因为你们一方，比其他任何人都更犯有这种过错；你们把这同一个存有分配给了\OriginalTerm{生者}{Begetter}和\OriginalTerm{受生者}{Begotten}；所以你们所给予的责骂不过是一个无法逃避的套索，是你们为自己的脖子编织的；正如可以预料的，正义用你们自己的话记录了对你们自己的判决。要么你们首先把这些存有设想为分离的、彼此独立的；然后藉着受生，把其中一个降低到圣子的等级，并主张一个独立存在的存有却是由另一个存有所造成的；这样你们就暴露在你们自己的谴责之下：因为你们所想象为无始的那一位，你们却归因于由另一个所受生。——要么你们首先承认一个单一的无原因存有，然后藉着一个因果行为把它标记为圣父和圣子，你们就宣称这个非受生的存有藉着它自身而存在。}\par

% structure: p0164 source=h2 target=subsection reason=relative-heading-rank; base=h2

\subsection{三十四、他攻击\OriginalTerm{同质}{\GreekText{ὁμοούσιον}}的段落，以及回应它的论辩。}

我将略过不谈论所引段落之前出现的那些话。它们不过是对我们在天主内的导师和父的无耻辱骂，与手头的事情毫无关系。但关于所引的段落本身，当他藉着这个可怕的两难困境的诡计提出一个双刃的驳斥时，我们不能沉默；我们必须接受这一智力挑战，用我们所有的力量为信德而战斗，并显示他磨利的这把可怕的双刃剑比布景画中的假象还要软弱。\par

他用两个假设攻击性体的共融；他说，我们要么把两个彼此平行伸展开来的独立本原称为圣父和圣子，然后说这些存在体中的一个是由另一个存在体所产生的；要么我们说同一个本质被设想为依次参与这两个名称，既是圣父，又成为圣子，并且它自己在受生中由自身所产生。我用我自己的话来表达这一点，从而不曲解他的思想，只是纠正其表达上浮夸的夸张，以便用更清楚的词语揭示他的含义，并提供一个全面的理解。他指责我们缺乏润饰，并带到争论中的学问不足，他用如此闪光的风格装饰他自己的作品，用他自己的话说，如此频繁地打磨他的句子，使他的辞藻如此精致漂亮，以至于立刻用语言的吸引力迷住了读者；我们刚才作为序言引用的段落就是众多例子之一。请允许我们再次引用它。\Quote{所以你们所给予的责骂不过是一个无法逃避的套索，是你们为自己的脖子编织的；正如可以预料的，正义用你们自己的话记录了对你们自己的判决。}\par

观察这些古阿提卡的词句；他的作品中闪耀着何等精炼的文采；一种何等细致微妙的风格之魅力在那里绽放！然而，让人们这样想吧。我们的道路需要我们再次转向那些话语中的思想；让我们再次投入这位小册子作者的句子中。\Quote{要么你将存有理解为彼此分离和独立，然后通过受生，将其中一个降格为\Person{圣子}的等级，并主张一个独立存在的存有却是由另一个存有所造成。}目前这就够了。那么，他说我们宣讲两个无因的存有。这位总是指责我们拉平混淆的人，怎么能因为我们相信两者的单一实体就断言这一点呢？如果我们的信仰宣讲两种不同的本性，正如异类派所宣讲的那样，那么有理由认为这种本性的区别导致了对两个无因存有的假定。但事实上，我们承认一个本性，只是位格有别，如果\Person{圣父}被相信，\Person{圣子}也被光荣，那么对手怎么能歪曲这样的信仰，说它宣讲两个元始因呢？接着他说，这两个原因中有一个被我们\Quote{降格为子的等级}。让他指出这种教义的一个维护者吧；无论他能否证明有任何人这样说话，或者他只知道教会中任何地方有这种教导，我们都不作声。因为谁在推理中如此疯狂，如此缺乏思考，以至于在谈论了\Person{圣父}和\Person{圣子}之后，却想象两个无始的存有，然后又假设其中一个是由另一个产生的？此外，他显示了什么逻辑必然性，将我们的教导推向这样的假设？他通过什么论据显示这种荒谬必然来自我们的教导？如果他真的引用了我们信仰的一条，然后无论是作为诡辩还是真正的论证力量对此提出批评，那么他这样做或许有理由以否定那条信仰。但既然教会中没有、也永远不可能有这样的教义，既找不到教师也找不到听众，而且这种荒谬也不能被显示为任何东西的严格逻辑结果，我不明白他与影子作战的意义。这就像一个发狂的人以为自己与一个想象中的对手搏斗，然后费了很大力气把自己摔倒，却以为敌人就躺在那里的；我们这位聪明的小册子作者处于同样的状态：他虚构我们不知道的假设，与他自己的大脑活动所描绘的影子作战。\par

因为我挑战他说，为什么相信\Person{圣子}是由\Person{圣父}所生的人必须走向认为有两个元始因的观点；让他告诉我们，谁在建立两个元始因上最有罪？是那个断言\Person{圣子}虚假地如此称呼的人，还是那个坚持当我们这样称呼祂时，这名称代表一个实在的人？第一个，拒绝\Person{圣子}真正的受生，而简单地肯定祂存在，更容易被怀疑把祂当作一个元始因，如果祂确实存在，但并非由受生而来；而第二个，使\Person{独生者}的位格的代表标记在于从\Person{圣父}通过受生而存在，绝不可能被拉入认为\Person{圣子}是无始者的错误。然而，按照你们这些思想家的看法，只要\Person{圣子}不被\Person{圣父}所生，\Person{圣子}本身就可以恰当地被称为无始的，在无始的众多含义之一中；因为既然有些事物通过出生而存在，另一些通过被塑造而存在，没有什么阻止我们称后者中并非通过\Emph{受生}而存在的一个为无始的，只看受生的概念；而你们的说法，将\Person{主}定义为受造物，确实确立了这一点。所以，非常博学的先生们，按照你们的观点，而不是我们的，当这样推导时，\Person{独生者}可以被叫做无始的：而你们会发现，\Quote{正义,}——无论你们指什么——用\Emph{你们自己的}话记录了对我们的不利判决。\par

此后，从他的话语中也很容易找到泥土来抛向这可憎的教导。因为他两难推理的另一端也参与了同样的精神错乱：他说，\Quote{或者你首先承认一个单一的无因存有，然后通过一个生行为将其划分为\Person{圣父}与\Person{圣子}，你宣称这个非受生的存有是由自身而产生的。}这是什么新奇奇妙的故事？一个人如何由自己而生，以自己为父，成为自己的儿子？这是何等的眩晕和迷惑！就像设想屋顶在自己脚下旋转，地板在头顶上；就像醉酒而感觉迟钝之人的精神状态，他持续大喊地面不稳，墙壁在消失，他看到的一切都在旋转而不静止。也许我们的小册子作者在写作时灵魂中有这样的骚动；如果是这样，我们必须怜悯他而不是憎恶他。因为谁如此远离我们神圣的教义，如此远离教会的奥迹，以至于接受这种损害信仰的观点？不，甚至很难说，没有人曾梦想过这种荒谬来损害信仰。为什么，就人的本性或任何其他落入感官把握的实体而言，谁在听到实体相通时，会梦想要么所有基于实体进行比较的事物都没有原因或开始，要么某物从自身产生，同时产生自己并被自己产生？\par

第一人，以及从他而生的人，以不同的方式获得存在：后者通过交合，前者由基督亲手塑造；然而，尽管他们被视为两个，他们在存有的定义中是不可分离的，并不被看作两个无始无因、平行并行的存在；也不能说已有的存在由已有的存在所生，或者这两个在一种怪诞的意义上被视为一个，即各自是自己的父亲和自己的儿子；而是因为两者都是人，所以两者有相同的存有定义：各自都是有死的、理性的、能直观和能科学的。如果亚当和亚伯身上的人性观念不因其起源的差异而变化，他们进入存在的次序和方式丝毫不改变他们的本性——这本性在两者中是相同的，每个有理智的人都可以证明，没有一个精神错乱得需要治疗的人会否认——那么，我们有什么必要非得在神圣本性上接受这种奇怪的思想呢？我们从真理那里听到了父和子，在这两个主体中，我们被教导他们本性的一体；这些名字所表达的自然关系指示了那本性；我们的主自己的话也是如此。因为当祂说：\Quote{我与父原是一体\BibleRef{若 10:30},}，祂通过承认父恰恰传达了自己不是第一因的真理，同时通过祂与父的合一表明了他们的共同本性；这样，祂的话从两方面保障了我们的信仰免受异端错误玷污：因为撒伯里乌没有理由混淆每个位格的独特性，当独生子用\Quote{我与父;}这样的话语如此明确地将自己与父区分开来时；而亚略也找不到支持他关于两者本性相异之说的依据，因为两者的合一不允许本性上的区分。因为在这话语中，父与子的合一所表示的，无非是他们实际存有本身所蕴含的；所有其他在他们身上观察到的超乎本性的道德卓越，都可以无误地视为所有受造物所共享。例如，先知称主为慈悲的和怜悯的，祂也愿意我们成为并被称为同样的：\Quote{你们应当慈悲\BibleRef{路 6:36},}，以及\Quote{怜悯人的人是有福的\BibleRef{玛 5:7},}，还有许多类似的章节。因此，如果有人通过勤奋和注意，按照天主的旨意塑造自己，变得仁慈、怜悯、同情，或温良和心谦，如同许多圣人在追求这类卓越时被证成的那样，难道他们就因此与天主合一，或凭其中任何一种德行与祂联合吗？并非如此。那在每一方面都不相同的东西，不能与本性如此不同的祂成为\Quote{一}。因此，一个人与另一个人成为\Quote{一}，是在意志上，如主所说，他们\Quote{被成全成为一,}，这种意志的合一添加到本性的联系上。同样，父与子是一，本性的共融和意志的共融在祂们身上合而为一。但如果圣子只是在意愿上与父相合，而在本性上与祂分离，我们怎能发现祂见证自己与父合一，而事实上祂在最属于祂自己的方面与父分离了呢？\par

% structure: p0171 source=h2 target=subsection reason=relative-heading-rank; base=h2

\subsection{三十五、证明欧诺米派的教导倾向于摩尼教。}

我们听到主说：\Quote{我与父原是一体,}，我们在这句话中得知我们的主依附于一个原因，同时又得知圣子与父的本性绝对同一；我们不使关于祂们的概念融化成一个位格，而是保持位格的属性有别，另一方面又不分割位格中本体的合一；这样就避免了在原因范畴中设想两个不同原则的理论，也没有给摩尼派异端留下进入的口子。因为受造与不受造正如它们的名称一样截然对立；因此，如果两者都被列为第一因，摩尼教的危害就会暗地里被带进教会。我说这话，是因为我反对对手的热忱使我非常仔细地审视他们的教义。现在我认为，没有人会否认，当我们说：如果受造与不受造具有同等能力，那么这些本性不同的事物之间就会有某种对立，而且只要两者在能力上都不欠缺，它们就会陷入一种相互不和谐的状态——因为我们必须承认，意志与本性相应并密切相连；如果两个事物在本性上不相似，它们在意志上也会不相似。但当两者能力相等时，任何一个都不会在实现其愿望上懈怠；如果每个的能力都与其愿望相等，那么这两个之间的首要地位就变得可疑；结果将因能力无穷无尽而成为平局。这样，摩尼派异端就潜入进来，两个对立的原则在原因范畴中竞争出现，因本性和意志的差异而分离和对立。因此，他们将发现，断言（神圣存在中的）缩减是摩尼教的开始；因为他们的教义在存有内部组织了一种不和谐，这导致了两个首要原则，正如我们的叙述所表明的：即受造和不受造。\par

但或多数人会指责这过于强烈的\Emph{\OriginalTerm{归谬法}{reductio ad absurdum}}，并希望我们连同其他反对意见一起根本不要写下它。就算如此；我们也不反驳他们。迫使我们将论证推到如此细微结果的，不是我们的冲动，而是我们的对手自己。但如果这样论证是不对的，那么我们的对手的教导，即给了这种反驳以机会的，更不应该被听到。压制对坏教导的回答只有一种方式，那就是消除需要答复的主题。但我最乐意做的，是建议那些如此倾向的人稍微摆脱竞争精神，不要成为他们所拥有的私见的极其热心的斗士，并确信这赛跑是为了他们的（属灵）生命，只关注其利益，并将胜利让给真理。那么，如果一个人停止这种野心勃勃的争斗，直视我们面前的实际问题，他很快就会发见这一教导的明显荒谬。\par

因为让我们假定我们对手的系统所要求的，即无生状态是\OriginalTerm{存有}{Being}，同样地，有生状态也被承认为存有。那么，如果一个人仔细地按照全部含义来追踪这些陈述，就连摩尼教异端也将被重建，因为摩尼教徒惯于把善与恶、光明与黑暗以及所有这类自然对立的东西作为公理。我想，任何不满足于肤浅看待此事的人都会相信我说的是真的。让我们这样看。每个主体都有某些内在特性，通过它们可以认识那底层性质的特殊性。无论我们是在研究动物界还是其他任何东西，都是如此。树和动物不是通过相同的标志被认识的；人的特征在动物界中也不延伸到畜类；同样，生命和死亡也不以相同的征象表示；在每种情况下，毫无例外，正如我们所说的，主体的区别抵制任何混淆它们、使它们相互融合的努力；我们在每个事物上观察到的标志不能被传达以致破坏那种区别。让我们在检查我们对手的立场时追踪这一点。他们说无生状态是存有；他们也使有生状态成为存有。但正如人和石头没有相同的标志（在定义有生命物和无生命物的本质时，你不会对每个给出相同的描述），所以他们也必然承认，无生者要通过与有生者不同的记号来认识。那么，让我们审视一下无生的天主的那些独特品质，即圣经教导我们可以提及和思想、而不致对祂不敬的品质。\par

它们是什么？我想没有哪个基督徒不知道祂是美善、仁慈、圣洁、公正而神圣的，不可见而不朽的，不会朽坏、变化和改变，全能、智慧、恩惠、主宰、审判者，以及所有此类事物。何必耽误我们的讨论，在已知事实前徘徊？那么，如果我们在无生的本性中发现这些品质，而有生状态在其概念本身中与无生状态相对立，那些将这两种状态都定义为存有的人，必须被迫承认，有生之物的标志，遵循有生与无生之间存在的这种对立，必须与无生之物的标志相反；因为如果他们宣布标志相同，这种相同性就会破坏作为这些观察对象的两者之间的差异。不同的事物必须被视作拥有不同的标志；相似的事物通过相似的记号被认识。那么，如果这些人在独生子身上见证相同的标志，他们就无法设想标志的主体有任何差异。但如果他们坚持他们的亵渎立场，并在断言有生者与无生者的差异时维持本性的变化，那么很容易看到必然的结果：即，正如在追踪名称的对立时，这些名称所指事物之本性必须被认为处于与自身的对立状态，那么每个被观察的品质就必须被引出彼此对立；以至于必须将那些与圣父的称谓——即天主性、圣洁、美善、不朽、永恒以及每一种对虔敬心灵代表天主的品质——相对立的品质归给圣子；事实上，这些品质的每一种否定，每一种与美善相对的意念，都必须被认为属于有生的本性。\par

为了确保清晰，我们必须就此点稍作停留。正如热与冷的特有现象——它们本身本性上彼此对立（让我们用火和冰作为各自事例），每一个都不是另一个所是——彼此冲突，冷却为冰的特有性，加热为火的特有性；因此，如果按照名称所表达的对立，这些名称所揭示的本性被分开，就不能承认这些伴随自然\Quote{\OriginalTerm{次对立}{subcontraries}}的能力彼此相似，如同冷却不能属于火，或者燃烧不能属于冰。那么，如果美善与无生本性的概念不可分离，而那一本性如他们所宣称因存有而与有生本性分开，那么前者的特性也将与后者的特性分开：因此，如果美善存在于前者中，则与美善相对立的品质就应在后者中被察觉。这样，多亏了我们聪明的系统化者，\Person{摩尼}又复活了，他带着与美善对立的那一排邪恶，以及他的对立能力居于对立本性的理论。\par

诚然，若要毫无保留地大胆直言，据说\Person{Manes}是第一个敢于接纳摩尼教观点的人，并以其名为该异端命名，那么相比之下，或许可以公允地认为他还不那么令人反感。我这样说，好比要在毒蛇和角蝰之间选择哪一种对人更友善；然而，仔细考量，某些野兽之间确实存在\Emph{某种}差异。比较二者的教义，不是更能说明那些更古老的异端者比这些更为可容忍吗？\Person{Manes}以为是在为善之源辩护，因为他主张恶绝不能从善之源获得任何原因；因此他将恶之名单上的事物之链条与一个分离的原理相连，以全能者的护卫者自居，且出于虔敬的厌恶，不愿将任何不当偏差的责任归咎于善之源；却没有以其狭隘的理解察觉到，既不可能设想天主是恶的制造者，也不可能在祂之外还有任何其他第一原理。关于这一点可以长篇讨论，但偏离我们当前的目的。我们提及\Person{Manes}的陈述，仅为了表明他至少认为有义务将恶与任何关乎天主的事物分开。然而，这些人所系统化的关于圣子的亵渎性错误则更为可怕。如同他人一样，他们也用存在上的对立来解释恶的存在；但除此之外，他们还宣称宇宙的天主实际上是这种异己产物的制造者，并说祂所形成之\Quote{受生}成了一种实体，拥有与其制造者本性相异的本性，因此他们在这方面表现出比前述教派更多的不虔敬；因为他们不仅赋予本性上对立于善者以实际存在，而且还说一位善的神祇是另一位在本性上偏离于祂的神祇的原因；他们的教导几乎公开宣称，存在某种相反于善的本性的事物，且其位格源自于善本身。因为当我们知道圣父的本体是善的，从而发现圣子的本体因其在本性上不像圣父（这是本异端的信条）而属于相反的谓词时，这证明了什么？不仅证明了相反于善的事物存在，而且这种相反来自于善本身。我宣称这比摩尼教徒的非理性更为可怕。\par

但若他们从自己的体系中否决了这种亵渎（尽管这是他们教导的逻辑结果），又说独生子继承了圣父的卓越，并非作为真正的儿子，而是——正如这些误信者所喜好的——通过受造行为获得其位格，那么我们也应审视这一点，看看这种想法是否合理。那么，若我们承认他们所想的属实，即万有之主并非作为真子继承，而是祂自己也被造和被创造，却统治着一群受造的同类，那么其余受造物如何会接受这种统治而不反抗？它们就这样被从同类关系推入隶属状态，并被判定（尽管在本性特权上丝毫不亚于祂，两者皆为受造）要服事并屈从一个同类。那就像篡权——不是将统治权赋予存在上的优越性，而是将本因自然权利拥有同等特权的受造界分为奴隶和统治力量，一部分发号施令，另一部分服从；仿佛通过任意分配的结果，这些同样的特权被随意堆砌在分配后优于同侪的那一位身上。甚至人在接受对野兽的统治之前，也没有与它们分享尊荣；理性的特权给了他发号施令的资格；他凌驾于它们之上，是因为他的本性向优越方向有了差异。人类政府之所以经历如此频繁的革命，正是因为这个原因：那些自然受赐平等权利的人不可能被排除在权力之外，因为当所有人同源时，自然的本能冲动便是与统治方平起平坐。\par

再者，若圣子本身也是受造之物之一，那么\BibleQuote{万物是藉着祂造的}这句话又如何为真呢？要么祂必须造了祂自己，才能使那经文为真，这样他们为损害我们的信德所设计的这个不合理之处将全部反噬他们自身；要么，若这是违情悖理的，那么\Quote{整个受造界都是藉着祂造的}这一断言将被证明毫无根据。抽去一个，就使\Quote{万有}成为虚假陈述。因此，从圣子是受造之物的这一定义，必然得出两个邪恶且荒谬的后果之一：要么祂不是一切受造之物的创造者，因为他们坚持认为祂是受造之工之一，必须从\Quote{万有}中抽去；要么祂被显明是祂自己的创造者，因为\BibleQuote{凡受造的，没有一样不是藉着祂造的}这一宣讲并非谎言。关于他们的教导，就说这些。\par

% structure: p0180 source=h2 target=subsection reason=relative-heading-rank; base=h2

\subsection{三十六、教会训导的简要复述。}

但如果一个人持守纯正的教义，相信圣子是那毫无掺杂的神性本性，他将发现一切与其他信仰真理和谐一致，即：我们的主是万物的造主，他是宇宙的君王，并非凭任意无常的权力凌驾其上，而是凭藉更高的本性来统治；此外，他还将发现，我们所教导的那个第一因并未因实体的不相似而被分割成多个分离的第一因，而是相信一个天主性、一个原因、一个统御万有的权能，这神性藉着这些相似存在之间的和谐而被发现，并引领思想从一个相似者到另一个相似者，以致万有的原因，即我们的主，藉着圣神照耀在我们心中（因为宗徒说，\BibleQuote{除非藉着圣神}\BibleRef{格前 12:3}，主耶稣不可能真正被认识）；然后，那超越的原因，即万有之上的天主，通过万有的原因——我们的主——被发现；事实上，除非藉着那不可见者的（可见）肖像，不可能获得对原型之善的准确认识。但随后，越过那神学的顶峰（我指的是万有之上的天主），我们仿佛在思想的跑道上折返，从圣父通过圣子到圣神，经由结合而相关的观念飞驰。因为一旦我们立足于对不生之光的领悟，我们立刻从那有利位置感知到他发出的光，就像与太阳共存的射线，其根源确实在太阳内，但其存在与太阳同时，并非后来添加，而是在看见太阳本身的一刹那显现；或者更确切地说（因为不必拘泥于这个比拟，从而因我们形象的不足而给批评者提供攻击我们教导的把柄），我们将感知的并非太阳的一道光线，而是另一个太阳，作为后裔从不生的太阳闪耀而出，与我们对第一位的观念同时，且在各方面与他相似：在美感、能力、光泽、大小、光辉上，以及我们在太阳中观察到的一切特质上。然后，我们又以同样的方式看见另一个这样的光，与那后裔之光没有时间间隔，虽然藉着它照耀，却将其存在的来源追溯到最初之光；然而它自身也是以同样方式照亮的光，它自身也是光的来源，并且做一切光所做的事。确实，光与光之间并无差别，就其作为光而言（\Emph{就其作为光而言}），当一方在照耀的恩宠上没有缺乏或减少，而是以其完满的完美成为一切至高之光的一部分，并与圣父和圣子一同被看见，尽管在他们之后被计数，凭其自身的能力使所有能分享它的人得以接近那在圣父和圣子内被感知的光。至此为止。\par

% structure: p0182 source=h2 target=subsection reason=relative-heading-rank; base=h2

\subsection{三十七、为圣巴西略的陈述辩护，优诺米曾攻击该陈述，即\Quote{圣父}与\Quote{不生者}可以有相同的含义。}

他的辱骂之流非常汹涌；傲慢是他所提出每项原则的基础；他把诽谤当作对可疑论点的任何证明。因此，让我们简要讨论一下关于\Quote{非受生}一词的许多歪曲陈述，他用这些陈述侮辱我们的导师本人及其论著。他引用了我们导师的以下话语：\Quote{依我之见，我倾向于说\InnerQuote{非受生}这个称号，尽管它似乎很适合表达我们的观念，但由于圣经中无处可寻，并且构成了欧诺米亵渎的基本要素，当我们有\Quote{父}这个表示相同含义的词时，完全可以省略它；因为那位本质上且独自为父者，并非来自其他任何一位；而不来自其他任何一位者，就等同于\InnerQuote{非受生}。}现在，让我们听听他如何证明这些话的\Quote{愚昧}：\Quote{仓促与无耻的不诚实促使他将这剂异常使用的词汇放入他的尝试中；他完全转向，因为他的判断摇摆不定，他的推理能力薄弱。}请注意那一击多么准确；他凭借其全部逻辑技巧，何等巧妙地剖析了巴西略的话语，并用一种更符合虔诚的观念取而代之！\Quote{措辞异常，}\Quote{判断仓促而不诚实，}\Quote{因推理薄弱而摇摆不定、转向。}这是为何？是什么激怒了这位判断如此坚定、推理如此正确的人？他最谴责巴西略话语中的什么？是他接受了\Quote{非受生}的\Emph{观念}，但说实际这个词，由于被那些曲解它的人滥用，应当被抑制？那么，信仰仅仅在词语和外在表达上处于危险之中，我们就不需要考虑内在思想的正确性吗？圣言难道不是首先劝勉我们拥有一颗远离邪恶思想的纯洁之心，然后，为了表达灵魂的情感，使用任何能够表达这些心灵秘密的话语，而不必细致地关注这个或那个特定的声音吗？因为这样或那样的说话方式并不是我们内在思想的原因；而是心中隐藏的概念为这样或那样的词语提供了动机；\Quote{因为心里所充满的，口里就说出来。}我们让词语解释思想；我们不会反过来从词语中收集思想。如果两者兼备，人当然可以在两方面都准备好，既有聪明的思想，也有巧妙的表达；但如果其中一方欠缺，对没有学问的人来说损失很小，只要他灵魂中的知识在道德美善方面是完善的。\Quote{这百姓用嘴唇尊敬我，心却远离我\BibleRef{依 29:13}。}那是什么意思？意思是灵魂对真理的正确态度，在天主眼中比措辞的恰当更宝贵，天主垂听\Quote{那说不出来的叹息}。词语可以被用于相反的意义；舌头乐意按说话者的意图服务；但灵魂的性情，无论它如何，都被那察看一切隐秘的主所看见。那么，他凭什么该被称为\Quote{异常}、\Quote{仓促}和\Quote{不诚实}，只因他吩咐我们在\Quote{非受生}一词中抑制所有那些能帮助那些背弃信仰的人在亵渎中得势的因素，同时保留并欢迎所有可以虔诚持有的含义？如果他确实说过我们不应将天主思想为\Quote{非受生}，那或许有些理由用这些甚至更恶劣的辱骂来攻击他。但如果他认同信徒们的普遍信仰并接受这一点，然后发表一个完全符合导师心意的意见，即\Quote{避免使用这个词，因为颠覆性的异端正是由此进入并被引出，}并吩咐我们通过其他名称来珍视\Quote{非受生}的天主的观念——那么他就不应受到他们的辱骂。真理本身难道没有教导我们如此行事，甚至不要紧抓极其宝贵的事物，若其中任何一样倾向于造成伤害吗？当祂吩咐我们砍掉右眼或右脚或右手，若其中一样成为绊脚石时，祂用这个比喻暗示了什么？岂不是说，任何事物，无论看起来多么美好，如果因轻率使用而引人作恶，就应使其无效且不被使用，并且向我们保证，通过切除那些导致犯罪的肢体而得救，比保留它们而丧亡更好？\par

同样，基督的追随者保禄说了什么？他也以深奥的智慧教导了同样的道理。他宣称\Quote{凡天主所造的物都是好的，若感谢着领受，就没有一样可弃绝的，}但有时，因着\Quote{软弱弟兄的良心，}他从他所接受的事物中退回一些，并命令我们回避它们。\Quote{若食物叫我弟兄跌倒，我就永远不吃肉\BibleRef{格前 8:13}。}这正是我们这位保禄的追随者所做的。他看到那些试图教导位格不平等之人的欺骗力量，因着他们有害的异端观点下的\Quote{非受生}一词而增强，因此他建议，虽然我们在灵魂中珍视对这位\Quote{非受生}天主的虔诚意识，但我们不应特别喜爱实际这个词，它是被弃绝之人犯罪的缘由；因为\Quote{父}这个称号，如果我们追随其全部含义，将会向我们提示\Quote{非受生}的意义。因为当我们听到\Quote{父}一词时，我们立刻想到万有的根源；因为如果祂自身之上还有某种超越的原因，祂就不会被恰如其分地称为\Quote{父}；因为那个称号本应转移到更高的、这个预设的原因上。但如果祂自身就是那个万有由之而来的原因，正如宗徒所说，那么显然，没有东西可以被认为存在于祂的存在之外。而这正是相信那存在并非受生。但这个人，他声称连真理本身也不应被认为比他本身更有说服力，不会默认这一点；他大声地教条化地反对；他嘲笑这个论证。\par

% structure: p0185 source=h2 target=subsection reason=relative-heading-rank; base=h2

\subsection{三十八、驳斥他诡辩式三段论的几种方法。}

如果您愿意，让我们来审视他那无可反驳的三段论，以及他在自己的虚假前提中巧妙的术语换位，他希望通过这些来动摇那个论点；不过，我确实担心他话中可怜的诡辩在某种程度上也会对他所要纠正的言论产生偏见。当年轻人挑战打斗时，人们与这样的对手厮打所受到的指责多于他们因展示胜利而获得的荣誉。尽管如此，我们想说的是这个。我们确实认为，他用我们现已熟知的、那众所周知的演说方式所说的话，仅仅是为了蛮横无理，最好沉默并遗忘；这些话可能适合他，但对我们来说，它们只是提供了锻炼极大耐心的练习。而且，我认为，将他的可笑说法插入我们严肃的争论中，并因此使这份对真理的热忱消散于粗俗、低级的笑声中，是不恰当的；因为确实，当他以如此崇高而华丽的冗词说\Quote{当增加言辞等于增加亵渎时，沉默比说话更加令人安宁一半}时，听到而不动容是不可能的。让那些知道哪些说法可信、哪些只配嘲笑的人对这些说法发笑吧；而我们则审视那些他试图用之以撕碎我们体系的三段论的敏锐性。\par

他说：\Quote{如果\InnerQuote{圣父}与\InnerQuote{非生者}含义相同，而含义相同的词语自然在各方面具有同样的力量，并且按他们的承认，\InnerQuote{非生者}意指天主出自虚无，那么必然得出：\InnerQuote{圣父}意指天主不是从任何事物而来，而不是意指生了圣子。}那么，这有什么逻辑必然性，阻止了生了圣子这一点被\Quote{圣父}这个称号所表示，如果那个称号本身也向我们表达了圣父的无始性？如果确实一个观念完全破坏了另一个，那么从矛盾物的本性来看，肯定一个就必然否定另一个。但是，如果世界上没有任何事物阻止同一个存在既是圣父又是\OriginalTerm{非生者}{Ungenerate}，当我们试图在\Quote{圣父}这个称号下思考非生成性作为其中隐含的一个观念时，有什么必然性能阻止圣父与圣子的关系再被\Quote{圣父}这个词所标记？其他表达相互关系的名称并不总是局限于那些关系观念；例如，我们称皇帝为\OriginalTerm{专制君主}{autocrat}和\OriginalTerm{无主者}{masterless}，我们也称同一位为其臣民的统治者；而且，虽然\Quote{皇帝}这个词确实也意味着无主性，但并不能因此而认为，这个词因为表示\OriginalTerm{专制的}{autocratic}和\OriginalTerm{不受统治的}{unruled}，就必须不再意味着对下属拥有权力；事实上，\Quote{皇帝}这个词介于这两个概念之间，时而表示无主性，时而表示对较低阶层的统治。那么，在我们目前的情况下，如果除了我们主的圣父之外，还有别的圣父可以想象，那么就让这些夸耀自己深刻智慧的人把他指给我们看，然后我们就会同意他，非生者的观念不能由\Quote{圣父}这个称号代表。但是，如果最初的圣父没有超越自身状态的原因，而圣子的\OriginalTerm{存有}{subsistence}始终隐含在圣父的称号中，为什么他们试图像吓唬小孩一样，用这些专业的命题扭曲来吓唬我们，试图说服或者说是诱骗我们相信：如果在圣父的称号中承认了非生成的特性，我们就必须切断圣父与圣子的任何关系？\par

那么，我们当轻视他们这种愚蠢肤浅的尝试，勇敢地承认他们所引述为荒谬绝伦的事，即我们相信：不但\Quote{圣父}与无生同义，而且后者确立了圣父非由任何所出，同时\Quote{圣父}一词也随同自身引进了独生者的概念，作为与之相联的关系者。如今，在我们的导师的论著中能找到的以下段落，被这位聪明且无往不胜的辩论家从上下文中删除了；因为他删去了\Person{巴西略}为防患而附加的部分，以为这样会使他自己的回答变得容易得多。这段文字逐字如下：\Quote{就我而言，我倾向于说，这个无生的称号，尽管似乎很容易符合我们自己的观念，但因圣经中无所记载，且构成了\Person{尤诺米乌斯}亵渎的字母表，便大可被压制，当我们有了圣父这个词，其含义相同，并随同自身引进了圣子的概念，作为与之相联的关系者。}这位真理的慷慨捍卫者，以天生的善意，删除了这句为防患而附加的句子，即：\Quote{并随同自身引进了圣子的概念，作为与之相联的关系者；}经过如此篡改后，他便与剩余部分短兵相接，切断了活生生的整体联系，从而使之，如他所想，更易于屈服于他的逻辑，并以那冰冷无力的谬误推理误导他自己的一方，即：\Quote{在一件单一事上与别物有共同意义者，在每一可能点上均保留这种意义的共同性；}并以此攻陷他们浅薄的智力。因为我们只肯定了\Quote{圣父}一词在\Emph{某种}含义上等同于无生，而此人却使意义的吻合在每一点上完全一致，大大偏离了任一词的通常理解；如此他便将事情引向荒谬，妄称如果\Quote{圣父}一词传达了未受生的观念，便不能再指与圣子的任何关系。这恰如一人对一块面包有两种认识——一为它是由面粉制成，二为它是食用者的食物——却对告诉他这件事的人争辩，用与\Person{尤诺米乌斯}相同的谬误：即\Quote{由面粉制成是一回事，是食物则是另一回事；那么，如果承认面包是由面粉制成，它的这一性质便不能再严格地称为食物。}\Person{尤诺米乌斯}的三段论的思想正是如此：\Quote{如果圣父一词隐含了未受生，那么这个词便不能再传达生了圣子的观念。}但我认为，现在我们反过来，应当将他那（已引用的）夸张圆滑的句子应用于他的论证。对于这样的话，宜说：如果他已将此等论证性防卫的界限设定为彻底沉默，则他更有理由被视为清醒。因为\Quote{当增加的话语意味着增加的亵渎时，}或更准确地说，表明他已完全丧失理智，则不仅\Quote{多一半，}而是整个多出部分，\Quote{沉默比说话更令人安宁。}\par

但也许人通过个人的示例更易被引入正确的观点；所以让我们撇开这种前后顾盼和虚假前提的扭曲，以不那么深奥、更通俗的方式讨论此事。\Person{尤诺米乌斯}，你的父亲当然是一个人；但同一个人也是你生命的作者。那么，你是否也曾在他身上使用过你所用过的这种巧妙诡辩，以致你的\Quote{父亲}，一旦接受了关于他存在的真正定义，便因他是\Quote{人}而不能再指与你自身的任何关系；\Quote{因为他必须是两者之一，要么是人，要么是\Person{尤诺米乌斯}的父亲？}那么，你不得以违背亲密含义的方式使用亲密关系的名称。然而，尽管你会控告任何以这种意义的篡改来轻蔑嘲讽你的人诽谤，你岂不畏惧嘲讽天主吗？当你嘲笑我们信仰的这些奥迹时，你安全吗？正如\Quote{你的父亲}指示与你的关系，同时人性并不被该词排除，并且没有一个清醒的人会用\Quote{人}来称呼那生你的人而非\Quote{你的父亲}，或者反过来，若被问及其种类而回答\Quote{人}，会断言该回答妨碍了他作为你的父亲；同样，在对全能者的默观中，敬畏的心灵不会否认：藉着圣父的称号，意指祂是无生的，同时在另一种意义上，它也代表祂与圣子的关系。然而\Person{尤诺米乌斯}公然蔑视真理，竟断言：一旦这个词向我们传达了\Quote{从未受生}的观念，它便不能再意指\Quote{生了圣子}了。\par

让我们再举一例说明他论断的荒谬。这是一个所有人都熟知的事例，即便是刚开始在文法教师指导下学习语词的孩子。请问，谁不知道有些名词是绝对且不表示任何关系，另一些则表达某种关系？在这后一类名词中，有些又可根据说话者的意愿倾向于任何一方；它们自身有简单的意向，但可被转化为关系名词。我不愿在无关我们主题的例子中停留。我将从我们信仰的话语本身来解释。\par

天主在圣经中被称作圣父、君王以及无数其他名号。这些名号中，一部分可以绝对地，即单纯地按其自身而言，不再多言：例如\Quote{不朽的}、\Quote{永存的}、\Quote{不死的}等等。每一个这样的名号，无需我们引入另一个思想，本身就包含着关于天主性体的完整思想。另一些则仅表达相对的有用性：如帮助者、护卫者、拯救者，以及类似意义的词语；若从中除去需要帮助之人的概念，则此词所表达的全部力量便消失。另一方面，有些名号，正如我们所说，既是绝对的，也属于关系词之列；例如\Quote{天主}、\Quote{善}以及许多其他此类。在这些词中，思想并非始终停留在绝对之内。那普世的天主常常成为呼求祂之人的私有财产，正如圣人们所教导的，他们将那独立自存者占为己有。\Quote{上主天主是圣的}；至此没有关系；但当人加上\Emph{我们的}天主上主，从而在相对于自身的关系中收归含义时，便使此词不再被视为绝对。再者，\Quote{阿爸，父啊}是圣神的呼声；这是没有任何片面指涉的宣告。但我们被命令称天上的父为\Quote{我们的天父}；这是此词的关系用法。一个人将那普世的天主据为己有，并不减损祂至高无上的尊位；同样，当我们指出圣父和那出自祂者、在一切受造物之前首生者时，我们同时藉圣父这一称号既表示生了那圣子，又表示祂并非出自任何更高超越的原因，这并无妨碍。因为凡谈及第一位圣父者，是指那先于一切存在而被预设者，那超越者。正是祂，没有在祂之前任何可观望之物，也没有祂将止息之终点。无论我们向哪边观看，祂都同样永恒存在；祂超越任何终点的界限、任何开始的观念，因祂生命的无限性；无论祂有何称号，永恒都必与之俱含。\par

然而，欧诺米，虽精于思辨那逃脱思想者，却拒绝这来自不科学头脑的看法；他不承认\Quote{圣父}一词有双重含义：其一，万物出自祂，万物之中首出者是独生子；其二，祂自己并无更高原因。他或许轻蔑此说，但我们仍将勇敢面对他的嘲笑，并重复我们已说过的话：\Quote{圣父}与那无始者是同一的，既表示生了圣子，也代表出自虚无。\par

但欧诺米，与我们这一说法争辩，说道（与他之前所说的完全相反）：\Quote{如果天主因为是生了圣子而被称为圣父，并且\InnerQuote{圣父}与无始者意义相同，那么天主是因为生了圣子而是无始者，但在生祂之前，祂并非无始者。}看他的转弯方法；他如何拆解自己最初的诡辩，并将其转向完全相反，以为这样就能用无可逃脱的结论困住我们。他的第一个三段论提出了这样的荒谬：\Quote{如果\InnerQuote{父}意味着出自虚无，那么它必然不再表示生了圣子。}但这后一个三段论，通过将前提转向其反面，用另一种荒谬来威胁我们的信仰。那么，他如何拆解自己先前的结论呢？\Quote{如果祂是\InnerQuote{父}是因为生了子。}他的第一个三段论完全没给我们这样的说法；相反，其逻辑推理旨在显示，如果\InnerQuote{父}这个词意指父未被生，那么它\Emph{不能}同时也表示生了圣子。因此他的第一个三段论丝毫没有暗示天主因为生了子而是父。我不明白这富于争论、工于职业技巧的逆转有何用意。\par

但让我们留意其言语中的思想。\Quote{如果天主因为生了圣子而是无始者，那么在生祂之前，祂并不是无始者。}对此的回答是清楚的：只需简单陈述真理，即圣父一词既表示生了圣子，也表示生者本身不被认为出自任何原因。如果你看效果，圣子的位格在圣父一词中显明；如果你寻找先前原因，则生者无始由该词表明。在说\Quote{在生圣子之前，全能者并非无始者}时，这位小册子作者使自己面临双重指控：即歪曲我们并侮辱信仰。他仿佛确凿无误地攻击我们老师从未说过、我们现在也未断言的事情，即全能者在时间过程中变成了父，此前曾是别的样子。此外，在嘲笑这虚构教义的荒谬时，他显出了自己的教义狂乱。他假定全能者曾一度是别的样子，然后通过进步才有资格被称为父，于是认为在此之前祂也不是无始者，因为无始性隐含在父的观念中。这愚蠢几乎无需指出；对任何反思者都将极为清楚。如果全能者在成为父之前曾是别的样子，那么这理论的拥护者若被问及他们设想祂处于什么状态，该说什么呢？在此存在阶段，他们将给祂什么名号？儿童、婴儿、幼孩、少年？他们会不会因如此明显的荒谬而脸红，说绝非如此，并承认祂起初就是完美的？那么，当祂尚不能成为父时，怎能是完美的呢？或者他们不会剥夺祂这能力，而只说同时存在父职\Emph{不适合}？但如果从一开始就作为如此圣子的父并非善、并非适合，那么祂如何获得那并非善的呢？\par

但事实是，天主成为如此圣子的圣父，这符合天主的尊威，是美好而恰当的。因此，他们将断言，起初祂在这美好之事上毫无份，并且只要祂还没有这位圣子，他们就不得不断言（愿天主宽恕我这样说！）祂没有智慧、没有能力、没有真理，也没有其他任何荣耀——从各种角度讲，这些荣耀正是独生子所是并被称呼的。\par

但让这一切都落在始作俑者的头上吧。我们将回到原处。他说：\Quote{如果天主因为生了圣子而成为圣父，并且如果圣父的意义就是\OriginalTerm{无生性}{Ungenerate}，那么天主在生育之前就不是无生性的。}现在，如果他在这里能像谈论人类生活时惯常的那样说话——在人类生活中，无法想象一个人能一下子获得许多技能，而是要按照一定的时间和顺序逐一赢得所追求的对象——如果我们能在全能者的情况下也如此推理，以至于可以说祂在某时拥有无生性，之后获得能力，然后是不朽，然后是智慧，然后进步而成为圣父，再之后是公义，然后是永恒，从而按照某种顺序进入所有关于祂的哲学概念——那么，认为祂的一个名号优先于另一个名号，谈论祂先是无生性，然后才成为圣父，就不那么明显荒谬了。\par

然而，实际上，没有人如此囿于尘世，如此不了解我们信仰的崇高，以至于在一旦领悟了万有的原因之后，不能将所有虔诚归于天主的属性汇集到一个统一而紧凑的整体中，反而构想出先后的属性，以及介于其中的某种顺序。事实上，不可能在思想中遍历一个属性然后到达另一个属性，无论是实在还是概念，这第一个在时间上超越另一个。天主的每一个名号，每一个关于祂的崇高概念，每一个与我们对祂的一般观念相协调的言说或思想，都与其同伴紧密相连；所有这些概念在我们的理解中都被聚合为一个统一而紧凑的整体——即，祂的父性、无生性、能力、不朽、良善、权威以及其他一切。你不能取其中一个并借着任何时间间隔在思想中将其与其余分离，仿佛它先于或后于其他东西；在祂身上无法发现任何一个崇高或可敬的属性不是在其永恒性中同时表达的。因此，正如我们不能说天主曾不是良善的、或不是有能力的、或不是不朽的、或不是不死的，同样地，不将父性永远归于祂，并说那后来才出现，是一种亵渎。真正是圣父的天主永远是圣父；如果永恒不包括在这一宣认中，并且如果一种愚蠢的先入之见追溯性地限制和阻碍了我们对圣父的概念，那么真正的父性就不能再恰当地归于祂，因为那种关于圣子的先入之见会取消祂父性的连续性和永恒性。现在被称为圣父的祂，怎么能被视为在其他属性之后才存在的东西呢？如果祂先是无生性，然后才成为圣父，并接受这名号，那么祂并不总是完全像现在这样被称呼的。但现在存在的天主，祂永远是这样；祂不会因任何增加而变得更坏或更好，祂不会因从另一来源取得什么而改变。祂始终与自身同一。因此，如果祂起初不是圣父，那么祂以后也不是圣父。但如果祂被承认为圣父（现在），我将回到同一个论证：如果祂现在是圣父，那么祂永远是圣父；如果祂永远是圣父，那么祂将永远如此。因此，圣父永远是圣父；并且由于圣子必须始终与圣父一起被思考（因为没有圣子使之为真，圣父的名号无法成立），我们在圣父内所默观的一切也必须同样在圣子内被观察到。\BibleQuote{凡圣父所有的，都是圣子的；凡圣子所有的，都是圣父的。}经文说：\Quote{圣父\Emph{拥有}那属于圣子的}，所以挑剔的批评者将无权从\Quote{一切}这个词的内容中推断出圣子的无生性，当说圣子拥有圣父所有的一切时；另一方面，也不能从中推断出圣父的被生性，当说凡圣子所有的都在圣父内被观察到时。因为圣子确实\Emph{拥有}圣父所有的一切，但祂\Emph{不是}圣父；同样，凡圣子所有的都在圣父内被观察到，但祂\Emph{不是}圣子。\par

那么，如果凡属于圣父的一切都在独生者内，而祂又在圣父内，并且圣父的身份不与\Quote{未生性}分离，就我个人而言，我看不出在圣父本身那里有什么东西可以思议，以至于被任何间隔分开，从而在我们对圣子的理解之先。因此，我们可以勇敢地面对那个吹毛求疵的三段论所引发的困难；我们可以轻视它，如同一个吓唬孩子的稻草人，仍然断言天主是圣的、不朽的、圣父、\OriginalTerm{未生者}{Ungenerate}、永在的，以及一切同时；并且，如果可能设想你能够撤去虔诚归给祂的属性之一，那么一切都会随着那一个而毁灭。因此，在祂内没有更老或更年轻；否则祂就会被发现比祂自己更老或更年轻。如果天主不是永远地是祂所有属性，而是在祂内有些是，有些则成为，遵循某种顺序（我们必须记住天主不是复合的；祂是什么就是祂的全部），并且按照这个异端，祂先是未生者，而后成为圣父，那么，鉴于我们不能用一种品质的堆积来思考祂，必然的结果是祂的全部必须比祂的全部既更老又更年轻，前者因祂的\OriginalTerm{未生性}{Ungeneracy}，后者因祂的圣父的身份。但是，如果正如先知论及天主所说的，祂\Quote{是同一的，}那么说在祂生养之前祂自己不是未生者，便是徒然的；我们无法发现这两个名字——圣父和未生者——彼此分离；这两个概念一起升起，彼此提示，在虔诚推理者的思想中。天主是从永远为圣父，也是永远的圣父，以及虔诚归给祂的一切其他称谓都按同样意义给予，时间和时间的流逝，正如我们所说，在永恒中没有任何位置。\par

现在我们来看他处理词语的熟练技巧所产生的其余结果；这些结果，他本人确实说，既可笑又可悲。诚然，对于他所说的话，人必须放声大笑，如果为那浸透他灵魂的错误而深深哀叹不是更合适的话。虽然按照我们的教导，圣父这个称谓根据其一个含义包含了未生性的概念，但他却将圣父这个词的全部意义转给未生者，并宣称：\Quote{如果圣父与未生者是相同的，我们可以放弃圣父，改用未生者；这样，子的未生者就是未生者；因为正如未生者是子的圣父，反过来说，圣父是子的未生者。}在此之后，一种对我们朋友的灵巧的钦佩之情油然而生，深信他的神学训练的多方面精妙之处远非大多数人所能及。我们的导师所说的包含在一句简短的话里，大意是藉着\Quote{圣父}这个称号可以表示未生性；但\Person{欧诺米}的言辞在数量上不取决于思想的多样性，而取决于任何类似名称范围内的东西如何可以被翻转。就像蒙眼绕圈推磨的牛，虽然走了一程，却仍留在原地，他也一样围着同一个主题打转，从未离开。有一次他讥讽我们说，\Quote{圣父}不表示生养，而表示出自虚无。又一次他编织了一个类似的难题：\Quote{如果圣父表示未生性，那么在他生养之前，他不是未生者。}然后第三次他又使用同样的伎俩：\Quote{我们可以放弃圣父，改用未生者}；接着他直接重复那多次呕吐出来的逻辑：\Quote{因为正如未生者是子的圣父，反过来说，圣父是子的未生者。}他多少次回到他的呕吐物，多少次又吐出来！难道我们不会惹恼大多数人吗，如果我们伴随着这种愚蠢的文字卖弄而拖拽我们的论证？在这种情况下，保持沉默或许更体面；然而，免得有人以为我们因为论证无力而拒绝讨论，我们就对他所说的话这样回答。\Quote{\Person{欧诺米}，你无权称圣父为子的未生者，即使圣父这个称号\Emph{确实}表示生养者本身不是出于任何原因。因为，以已经举过的例子来说，当我们听到\InnerQuote{皇帝}一词时，我们理解两件事：既指那在权威上卓越的人不受任何人的支配，也指他统治属下。同样，圣父这个称号给我们提供了关于天主性的两个概念：一个与祂的儿子有关，另一个与祂不依赖于任何可预先设想的原因有关。因此，正如在\InnerQuote{皇帝}的情况下，我们不能因为那个词表示两件事，即统治臣民和不拥有任何先于他的人，就合理地说\InnerQuote{臣民的未被统治}，而不说\InnerQuote{民族的统治者}，或者允许如此之多，以致我们可以使用这样一组词语的并列，仿效\InnerQuote{民族的王}，如\InnerQuote{民族的无王}；同样，当\InnerQuote{圣父}指示一个儿子，同时也表示未生性的概念时，我们不能不正当地转移这后一个含义，以便将未生性的概念紧紧附着于父子关系，并荒谬地说\Quote{未生者是子的未生者}。}\par

他自认为，在说出这些话之后，便踏上了真理之地；他已揭露了对手立场的荒谬。他何等自夸地喊道：\Quote{请问，哪个理智清醒的人曾经希望自然的思想被压制，而欢迎谬论？} 最卓越的先生，没有一个理智清醒的人会这样。因此，我们的论述也不会如此。我们教导说，\OriginalTerm{非受生者}{Ungenerate}一词确实适合我们的思想，我们应当完好地保存在心中，然而\OriginalTerm{圣父}{Father}一词足以替代您所曲解的那一个，并将心灵引向正确的方向。请记住您自己引用的话；\Person{巴西略}（Basil）并没有\Quote{希望自然的思想被压制，而欢迎谬论}，如您所说；而是建议我们通过压制\Quote{非受生者}这个单纯的词（即那个以多个音节表达的词语）来避免一切危险，因为它曾被恶意解释，而且也不见于圣经；至于其含义，他宣称它确实完全适合我们的思想。\par

以上便是我们的陈述。但是，这位谴责一切诡辩者、以真理完全武装自己的论点、并依逻辑指控我们罪过的人，在他任何关于教义的论证中，却毫不脸红地沉迷于非常漂亮的诡辩；这与那些在餐后为了引人发笑而讲出的精妙笑话如出一辙。请想想那个复杂三段论所展现的推理分量；我现在重新引用它：\Quote{如果\InnerQuote{圣父}等同于非受生者，那么我们就可以放弃\InnerQuote{圣父}，而改用非受生者；如此，非受生者便是圣子的非受生者。因为正如非受生者是圣子的圣父，反过来，圣父便是圣子的非受生者。} 好吧，这很像下面这样一个例子。假设有人对\Person{亚当}提出了正确而健全的看法：即我们称他为\Quote{人类之父}或是\Quote{天主所造的第一个人}（因为两者意思相同）都无关紧要。然后，一个属于\Person{欧诺米}（Eunomius）学派推理者的人抓住了这个陈述，并从中得出了同样的复杂混乱：如果\Quote{天主所造的第一个人}和\Quote{人类之父}是相同的事，那么我们就可以放弃\Quote{父亲}这个词，而改用\Quote{第一个被造的}，称\Person{亚当}为\Quote{第一个被造的}而不是\Quote{父亲}，关于\Person{亚伯尔}（Abel）；因为正如第一个被造的是儿子的父亲，反过来，那个父亲便是那个儿子的第一个被造的。如果这话是在酒馆里说的，那么从饮酒的人群中会爆发出何等笑声和掌声，为了如此精妙绝伦的笑话！这位博学的神学家所依赖的就是这些论证；当他攻击我们的教义时，他实际上需要一个导师和一根棍棒来教导他：所有属于某人的谓词，并不必然在意义上指向同一个对象；这从上述\Person{亚伯尔}和\Person{亚当}的例子可以清楚看出。同一个\Person{亚当}既是\Person{亚伯尔}的父亲，也是天主的造物，这是真理；然而，并不因为他两者都是，就因此对\Person{亚伯尔}而言两者都是。同样，称全能者为\Quote{圣父}既包含这个词的特殊含义，即生了儿子，也包含对圣父本身没有前因的思想；然而，当我们提到圣子时，我们并不必须称祂为\Quote{非受生者}而非\Quote{圣父}；另一方面，如果\Quote{无始}这个属性在涉及圣子时没有被表达，我们也不必将\Quote{非受生性}这个属性从我们对天主的思考中驱逐出去。但他抛弃了通常的接受方式，像喜剧演员一样，把整个主题当成了笑话，并凭借他诡辩的古怪性，使我们信仰的问题变得可笑。我必须再次重复他的话：\Quote{如果圣父等同于非受生者，那么我们就可以放弃圣父，而改用非受生者；如此，非受生者便是圣子的非受生者。因为正如非受生者是圣子的圣父，反过来，圣父便是圣子的非受生者。} 但让我们反过来用他的诡辩来嘲笑他吧：如果\Quote{圣父}与\Quote{非受生者}不同，那么圣父的圣子就不会是非受生者的圣子；因为祂只与圣父有关系，所以祂将与不同于圣父且不适合那个概念的本性完全疏远。因此，如果圣父是不同于非受生者的某物，并且\Quote{圣父}这个称号不包含那个含义，那么圣子作为一个，就不能被分配在这两种关系之间，同时既是圣父的圣子，又是非受生者的圣子；正如从前说天主是圣子的非受生者被认为是荒谬的，同样，在这个逆命题中，称独生子为非受生者的圣子也会被发现同样荒谬。因此，他必须在两者中选择其一：要么圣父等同于非受生者（这是必要的，以便圣父的圣子也能是非受生者的圣子），那么我们的教义就被他无端嘲笑了；要么圣父是不同于非受生者的某物，那么圣父的圣子就与一切与非受生者的关系疏远。但若这样认为独生子不是非受生者的圣子，那么逻辑必然指向一个\Quote{受生的圣父}；因为存在但不从受生而来的存在，必须有一个受生的实体。因此，如果按照这些人的说法，圣父是不同于非受生者的，因此是受生的，那么他们常说的非受生性在哪里？他们异端构筑的那个基础和根基在哪里？他们刚才以为抓住了的\Quote{非受生者}，却逃脱了他们，在一些空洞的三段论作用下完全消失了；他们想要证明\Emph{不相似}的企图，像一场梦一样，在批判的触碰下滑走，并随着这个\Quote{非受生者}一起飞走了。\par

因此，每当谬误被接纳而胜过真理时，它确能因它所造成的幻象而繁盛一时，但很快就将崩蹋；它自己的论证方式将消解它。但我们提出这一点，只是为了对他们那\Emph{\OriginalTerm{不相似}{Unlikeness}}所可能施行的美妙报复付之一笑。我们现在必须回到我们论述的主线上来。\par

% structure: p0203 source=h2 target=subsection reason=relative-heading-rank; base=h2

\subsection{三十九、答复他始终提问的问题：\Quote{那\InnerQuote{是者}能为受生者吗？}}

\Person{欧诺弥}并不喜欢\Quote{圣父}这一术语传达\OriginalTerm{无生者}{Ungenerate}的含义，因为他要确立曾有一段时间圣子不存在。事实上，在他的门徒中不断有人提问：\Quote{那常存的\InnerQuote{是者}如何能被受生？}我认为，这来自于当我们思考天主时，没有使自己摆脱从人类用语的角度考虑。但让我们毫不苦毒地立即揭露他这\Quote{隐秘动机}的实际虚假，首先陈述我们在这件事上的结论。\par

\Person{欧诺弥}，这些名称在我们这里有不同的意义；当我们来到超越性能力时，它们产生另一种意义。的确，在一切其他事物中，分隔人类与天主的鸿沟是广阔的；经验无法在此下界指出任何在量上类似于我们可能在那里猜测和想象的东西。同样，就我们术语的意义而言，尽管在语词层面上，人与永恒者之间可能有些相似，但这两个世界之间的鸿沟才是意义分离的真正尺度。例如，我们的主在比喻中称天主为一个作\Quote{家主}的\Quote{人}；但是尽管这个称号对我们十分熟悉，我们所想的人物和那里所指的人物会是同样的描述吗？我们的\Quote{家}会与那大屋——宗徒说，其中有金器、银器，以及所列举的其他材料的器皿\BibleRef{弟后 2:20}——相同吗？难道\Emph{那些}器皿不是超越我们当下的理解，并要在有福的不朽中默观，而我们的器皿却是土造的，并将消解为尘土吗？所以，几乎在一切其他术语中，人类事物与神圣事物之间名称相似，然而在这种相同之下揭示出意义的巨大差异。我们在两个世界中都同样提到身体肢体和感官；如同我们一样，天主的生活——众人都承认是超越感官的——也被依次列出手指、手臂、手、眼、眼睑、听、心、脚、凉鞋、马、骑兵和战车；以及无数其他取自人类生活的隐喻，用来象征性地说明神圣事物。因此，这些名称的每一个都有人类的声音，但没有人类的意义；同样，\Quote{圣父}这一名称，尽管同样适用于神圣和人类生活，却隐藏着所表达意义之间的区别，这一区别恰好与这一称号的拥有者之间的差异成比例。我们对人类的产生以一种方式思考；对神圣的产生以另一种方式推测。一个人在一个特定的时间出生；一个特定的地方必须是他生命的容器；没有它，他在本性上不可能有任何具体的实体：因此，时间片段不可避免地包围着他的生命；有他之前、与他同时和他之后。对于任何生于这世界的人，说\Quote{曾有一个时间他并不存在}\Emph{是}真的，说他现在存在，并且又会有一个时间他将不再存在；但这些时间观念并不进入永恒世界；对清醒的思想家来说，它们与那个世界毫无相似之处。思想神圣生命真正为何的人，将超越\Quote{某时}、\Quote{之前}和\Quote{之后}，以及时间中这种延伸的一切标记；他将在如此崇高的事上持有崇高见解；他也不会认为那绝对者受制于他观察到在人类诞生中生效的那些法则。\par

受动先于人的具体存在；为形成生物奠定了某些物质基础；在这底下是本性，它凭天主的旨意，施行奇工，使万物为即将诞生者恰当地提供营养，从每个土元素中提取特定情况所需的量，接受适当时间的协作，以及父母食物中为形成子女所必需的量：总之，本性通过这些构建人的生命的过程前进，将不存在者带到诞生；因此我们说，曾经不存在，它现在诞生了；因为，一度不存在，在另一时它开始存在。但当涉及神圣产生时，心智拒绝本性的这一服务，拒绝时间在发展中贡献的圆满，以及我们论证中所设想的在人类产生中发生的一切；而那个不带肉性概念进入神圣主题的人，不会再次堕落到任何那些卑贱思想的水平，而是寻求一种与所表达事物的庄严相称的思想；他不会认为与无受动者相关的受动，或认为万物的创造者需要本性的帮助，或承认时间延伸到永恒生命之中；他将看到神圣产生应摆脱所有这些观念，并且只允许\Quote{圣父}这一称号具有这样的意义：即\OriginalTerm{独生子}{Only-begotten}本身并非无源头，而是从那一位领受祂存在的原因；至于祂实际开始存在，他不会计算，因为他无法看到所谈之事的任何迹象。\par

既然那\Emph{是}与圣父同在者，在某种不可思议的方式下，在万世之前不承认有\Quote{曾有一时}，祂确实由生育而存在，但祂从未开始存在。祂的生命既不在时间中，也不在地方。但当我们摒除这些以及所有类似的概念来默观圣子的自立体时，我们所能想到的只有一样东西在祂之先——即圣父。然而独生子，正如祂亲口告诉我们的，是在圣父内，因此，按祂的本性，不能假设祂曾不存在。如果圣父曾不存在，那么圣子的永恒性必定因圣父的这种虚无而被追溯性地取消；但如果圣父常是，圣子怎么可能会不存在呢？因为祂根本无法脱离圣父而独自被思考，而总是在\Quote{圣父}这个名称中被默默地暗示。事实上，这个名称同等地传达了两位格；\Quote{圣子}的概念必然由那个词暗示出来。那么，圣子什么时候不存在呢？我们在哪个范畴中发现祂的不存在呢？在地方？没有。在时间？我们的主在万有之先；如果是这样，祂什么时候不存在呢？如果祂在圣父内，祂在什么地方不存在呢？请你告诉我们，你这如此善于看见看不见之事的人。你的思考给什么样的间隔赋予了形状？你在圣子身上，无论是按实体还是按概念，你能够想到什么样的空虚，以至于显示出圣父的生命在被平行延伸时超过了独生子的生命？哎呀，甚至对于人，我们也不能绝对地说某人先不存在，然后出生。\Person{肋未}，在他肉身出生的许多代之前，就向\Person{默基瑟德}纳了什一之税；所以宗徒说：\Quote{连那接受什一之税的\Person{肋未}也（在\Person{亚巴郎}内）纳了什一之税。}并加上了证明：\Quote{因为当\Person{亚巴郎}会见至高者的司铎时，他还在他祖先的腰中。}如果这样，一个人按某种意义先不存在，然后出生，然而因着在他的先祖中的实体亲缘关系而预先存在，按照宗徒的见证，那么关于神圣的生命，他们怎么敢说出祂先不存在，然后被生这种想法？因为祂\Quote{在父内}，正如主告诉我们的：\BibleQuote{我在父内，父在我内}\BibleRef{若10:38}，当然两者以两种不同的方式彼此内住：圣子在圣父内，如同肖像的美存在于它所被描绘的形式中；圣父在圣子内，如同那原初的美存在于它的肖像中。现在，在所有手工制作的肖像中，时间的间隔是模型与接受其形式之物之间的分离点；但在那里，二者不能相互分离，既不能将\Quote{本体的真像}与\Quote{位格}分离，用宗徒的话\EditorialNote{希伯来书第一章}，也不能将\Quote{光辉}与天主的\Quote{光荣}分离，也不能将代表与美善分离；一旦思想把握了其中之一，它就也接纳了相关的真理。他说祂\Quote{\Emph{是}，}（而不是成为）\Quote{祂光荣的光辉}；这样，我们显然可以永远摆脱潜藏在这两种概念中的亵渎；即，认为独生子可以被想作\OriginalTerm{非生者}{Ungenerate}（因为他说\Quote{祂光荣的光辉}，光辉来自光荣，而不是相反，光荣来自光辉）；或者认为祂曾开始存在。因为\Quote{是}这个字是一个见证，向我们解释了圣子的连续性、永恒性以及对一切时间标记的超越。\par

那么，我们的敌人有什么理由提出那个无关紧要的问题来损害我们的信德呢？他们认为这个问题无法回答，因而证明自己的教义，并且他们不断问这个问题，即\Quote{那存在者能否被生？}我们可以立刻勇敢地回答他们：那在\OriginalTerm{非生者}{Ungenerate}内的那一位\Emph{被生}自祂，并\Emph{确实}从祂那里获得自己的根源。\BibleQuote{我因圣父而生活}\BibleRef{若4:57}：但不可能指出祂开始的\Quote{何时}。当没有中介的质料、概念或时间间隔来分隔圣子的存在与圣父时，也不可能想到任何记号，藉此独生子可以从圣父的生命中脱离，并被显示为来自祂自己的某个特殊根源。如果这样，没有别的原则指导圣子的生命，除了圣父之外，虔诚的心在默观圣子的自立体之前（但不可分开）不能想到任何东西；如果圣父是无始无生的，甚至我们的对手也承认，那么那位只能在无始的圣父内被默观者，怎么能够祂自己有一个开始呢？\par

再者，如果我们承认我们的对手提出作为荒谬的那些说法，即当他们问\Quote{那存在者能否被生？}时，我们的信德会受到什么损害呢？我们并不主张这在\Person{尼苛德摩}提出冒犯性问题\BibleRef{若3:4}的意义上是可能的，在那里他\Emph{认为}已经存在的人不可能有第二次出生；但我们主张，祂的存在依附于一个永恒且无始的存在，并伴随着每一位探索时间古旧者，且抢先于思想进入超越世界的好奇心，且与我们对圣父的所有概念密切结合，祂的存在没有开始，正如祂也不是\OriginalTerm{非生者}{Ungenerate}；但祂既是被生的，又是存在的，在因果关系上显示出从圣父的出生，但因祂永远的生命而排除了任何不存在的时刻。\par

但是这位思想家以其极端的精细，反对这一陈述；他以一个是被生成的、另一个是非生成的为由，将独生子的存有与圣父的本性分离；尽管有如此之多的名号可以虔敬地用于天主，且所有这些名号都同样适用于两位，他却对它们视而不见，因为这些名号表明的是两者共有的东西；他紧抓住\Quote{非生成者}这个名号，并且只这一个；甚至对这个名号，他也不采用通常和公认的含义；他彻底改变了它的概念，取消了它通常的关联。这究竟是什么原因？因为如果没有非常有力的理由，他不会把语言从其公认的意义中强行扭出，并通过改变词语的意义来创新。他完全知道，如果它们的意义局限于惯常的意义，他就无力颠覆正确的道理；但如果这些词语从其普遍和通常的含义中被歪曲，他就能随同词语一起败坏道理。例如（具体谈谈他所滥用的词语），如果按照我们信德的一般思想，他承认天主之所以被称为非生成者，只是因为祂从未被生成，那么他的整个异端体系就会崩溃，因为他关于\Quote{非生成者}的诡辩就被撤除了。也就是说，如果他能通过效法教会所用的大多数天主名号的类比，而被说服将万有的天主视为非生成者，正如祂是不可见、无情感、非物质的；并且他同意在这些术语的每一个中，都只表示那些无论如何不属于天主的属性——例如，身体、情感、颜色以及源于一个原因——那么，如果他对事情的看法如此，他那一派的\Emph{不相似性}信条就会失去意义；因为在所有其他方面（除了非生成性）关于万有之天主所思想的，甚至连这些反对者也承认独生子与圣父之间的相似性。但为了防止这一点，他将\Quote{非生成者}这个术语置于所有这些指示天主超越本性的名号之前；并使之成为攻击我们信德的据点；他将\Quote{生成的}和\Quote{非生成的}这两个实际表达之间的对立转移到这些词语所指的位格本身；因此，通过词语之间的这种差异，他诡辩地论证存有之间的差异；不同意我们所说的\Quote{生成的}只用于因为圣子是被生成的，而\Quote{非生成的}是因为圣父存在而没有经过生成；却肯定他认为前者是通过被生成而获得了存在；尽管我不知道是什么样的哲学导致他得出这种看法。如果人们只注意这些词语本身的含义，在心中抽象出这些名称所代表的位格，便会发现他们这些说法的无根据。那么，考虑一下：并非由于圣父在概念上先于圣子（正如信德所真实教导的），这些名字本身的顺序就必须按照它们所代表者的价值和顺序来安排；而是独立地看待它们本身，看看它们之中哪一个（我再说一遍，是词语，而非它所代表的事实）在心思的概念中应置于另一个之前；哪一个传达了某种观念的确立，哪一个又是同一观念之否定；例如（为清楚起见，我认为类似的词语对会表达我的意思）：知识和无知、情感与无感——以及诸如此类的对比——其中哪一个在概念上优先于其他？是那些确定否定者，还是那些确定对该性质之肯定者？我认为是后者。知识、愤怒、情感首先被构想；然后才是这些观念的否定。但愿没有人出于过度的热心而责备这一论证，好像它要把圣子置于圣父之前。我们并非主张圣子在概念上应置于圣父之前，因为论证只是在区分\Quote{生成的}和\Quote{非生成的}的含义。所以，\Quote{生成}表示对某个实在或某个观念的确立；而\Quote{非生成性}表示其否定；因此完全有理由将\Quote{生成}首先加以考虑。那么，他们为什么坚持在这两个名字中把在概念顺序上第二个的名字归给圣父？为什么他们不断地认为否定能够界定并包含所讨论术语的整个实体，并对那些指出他们论证无根据的人激怒不已？\par

% structure: p0211 source=h2 target=subsection reason=relative-heading-rank; base=h2

\subsection{四十、他在\Person{巴西略}驳斥他之后试图与自己陈述一致的尝试未成功。}

请看他何等尖刻地对待那个识破他那恶毒作品的腐朽和软弱的人；他尽其所能地报复，而这一切不过是谩骂和污蔑：在这方面，他倒是能力非凡。那些想把文体装扮优雅的人，有一种方法，就是使用某些连接词来填补缺乏节奏的地方，从而使他们的短语组合更加和谐连贯；同样，\Person{欧诺弥}在他大多数段落中都用辱骂的修饰语来装点他的作品，仿佛要炫耀他那滔滔不绝的谩骂能力。我们又成了\Quote{愚人}，我们又\Quote{推理不当}，又\Quote{没有以重要性所要求的准备来介入争论}，又\Quote{错过了发言人的意思}。这位文雅的演说家对我们的导师使用了这样以及更多的话。不过，或许他的愤怒并非没有理由；这位小册子作者义愤填膺也是情有可原。因为\Person{巴西略}凭什么要揭穿他教义的弱点？凭什么要在淳朴弟兄们眼前揭露隐藏在他那貌似合理的诡辩之下的亵渎？为什么他不该让沉默掩盖这观点的错误？为什么他要挂着这个可怜的人，本应怜悯他，并用面纱遮盖他论点的无耻呢？他居然找出并展示了一个在私人学生中被誉为聪明机智的人！\Person{欧诺弥}在著作中的某处说过，未生的属性\Quote{跟随}神性。我们的导师评论他这句话说，一个\Quote{跟随}的东西一定是外在物之一，而实际存有却不是外在物之一，而是任何事物就其存在而言的实在。然后这位温和却不可战胜的对手大怒，滔滔不绝地谩骂，因为我们的导师一听到那句话，就同样领会了它的意思。但他做错了什么，如果他只是坚持你自己著作的意思？如果他确实不合逻辑地抓住了所说的话，那你说的就都是真的，我们就得忽略他所做的；但既然你对他的指责感到脸红，为什么不从你的小册子中删掉那个词，反而辱骂指责者？\Quote{是的，但他没有理解论证的要点。}那么，我们做错了什么，如果作为人，我们从你的实际话语中猜测意思，却无法理解你内心深处隐藏的东西？洞察不可察之事、检阅我们无法理解之物的特性、并认知不可见世界中之\Emph{不同}，乃是天主的事。我们只能凭听到的来判断。\par

% structure: p0213 source=h2 target=subsection reason=relative-heading-rank; base=h2

\subsection{四十一、跟随者与被跟随者不同。}

他先说：\Quote{未生的属性跟随神性。}由此我们理解他认为未生性是天主外在事物之一。然后他说：\Quote{或者不如说，这未生性就是祂的实际存有。}我们不明白这种\Quote{推论}的合理性；事实上我们注意到一些非常奇怪和不协调之处。如果未生性跟随天主，却又构成祂的存有，那么按照这种观点，同一主体将被赋予两种存有；因此，天主仍是以前和历来所信的那样，但除此之外还有另一个存有伴随，他们称之为未生性，与祂（其所跟随的那位）截然不同，正如我们的导师所言。那么，如果他命令我们这样想，就必须原谅我们思想贫乏，无法追随如此微妙的思辨。\par

但如果他否认这观点，并不承认在天主内有两种存有，一种由天主性代表，另一种由\OriginalTerm{无始性}{Ungeneracy}代表，那就请这位既非\Quote{轻率}亦非\Quote{恶意}的朋友，在对真理的战斗进行之际，勉励自己不要过于偏好谩骂，而是向我们这些如此缺乏教养的人解释：何以那\Quote{跟随}的并非一物，那\Quote{引领}的亦非另一物，而是两者如何融合为一？因为，尽管他为自己的陈述辩护，其中的荒谬仍旧存在；且那寥寥数词的添加，如他所说，并未纠正其中的矛盾。我至今未能看出其中能找到任何解释。但我们将原文照录。\Quote{我们说，\InnerQuote{或更确切地说，\OriginalTerm{无始性}{Ungeneracy}是祂的实际\OriginalTerm{存有}{being}，}并非有意要将我们已证明为跟随存有的东西\Emph{收缩}入存有，而是将\InnerQuote{跟随}应用于称号，而将\Emph{是}应用于存有。}据此，当这些放在一起时，整个得出的论点就会是：\OriginalTerm{无始}{Ungenerate}这个称号跟随，因为是无始乃是祂的实际存有。然而，我们将得到这位解释者的什么解释者呢？他说：\Quote{并非有意要将我们已证明为跟随存有的东西收缩入存有。}也许有些猜谜者会告诉我们，他说的\Quote{收缩入}意思是\Quote{连接在一起}。但谁能看到其余部分有任何可理解或连贯的东西？他告诉我们，\Quote{跟随}的结果不属于存有，而属于称号。但最博学的先生，什么是称号？它与存有不和，还是在思想中不与它一致？如果称号不适合于存有，那么存有怎能由称号代表？但如果如其本人所述，存有由无始的称号合适地限定，那么此后它们怎能分离？你让存有的名字跟随一物，而存有本身跟随另一物。那么\Quote{整个观点的构造}是什么？\Quote{\OriginalTerm{无始}{Ungenerate}的称号跟随天主，因为祂本身是\OriginalTerm{无始的}{Ungenerate}。}他说，那\Quote{跟随}天主的——天主是与无始者不同的某物——正是这个称号。那么，他怎能将天主性的定义置于无始性之内？再者，他说这称号\Quote{跟随}天主，如同存在于没有先前受生之物中。谁来为我们解开这些谜语的奥秘？\Quote{无始}先出现，然后跟随；先是存有的合适附着的称号，然后像陌生者一样跟随！还有，是什么原因导致对这个名字的过分躁动？他将天主性的全部内容赋予它；好像如果天主如此被命名，我们的朝拜将毫无欠缺；而如果祂不如此命名，我们信仰的整个体系将受到危害？现在，如果对此的简要陈述不应被视为多余或无关，我们将如此解释此事。\par

% structure: p0216 source=h2 target=subsection reason=relative-heading-rank; base=h2

\subsection{四十二、\Quote{无始}之解释与永恒之\Quote{研究}。}

天主的生命之永恒，仅作轮廓描述，是这样的：祂常被领悟为在存有中；祂不允许一个祂尚未存在和将不再存在的时间。那些在平面几何中从中心到圆周距离画圆的人告诉我们，他们的图形没有确定的起点；且那线条既无确定终点亦无可见起点：他们说，由于它自身形成一个单一整体，四周半径相等，它避免了任何起点或终点的指示。那么，当我们把无限存有与这样一个图形比较时——尽管它是有界限的——愿无人指责这一描述；因为我们关注的并非圆周，而是图形与我们断言为我们对永恒者的直观的那种在每个方向上都难以把握的生命之间的相似性。从当下瞬间，如同从一个中心和\Quote{一点}，我们延伸思想至各个方向，达至那生命的无限。我们发现我们被不间断、均匀地吸引着转圈，且我们总是跟随一个无一物可把握的圆周；我们发现神性生命在其自身内以一种无间断的连续性返回，其中无法识别任何终点或部分。关于天主的永恒，我们述说从先知所听到的：即天主是\Quote{自古}的君王，统治直到万代，永远且超越。因此，我们定义祂早于任何开始，且超越任何终点。那么，持有对全能者的这种观念，作为一种充分的观念，我们用两个称号来表达它；即\OriginalTerm{无始}{Ungenerate}和\OriginalTerm{无尽}{Endless}代表了天主性的这种无限性、连续性和永恒性。如果我们只采用其中之一用于我们的观念，而丢弃另一个，我们的意思将因这一省略而受损；因为单独使用两者中的任何一个都无法表达存在于两者之中的概念；但当一个人说到\Quote{无尽}时，只指示了关于终点的缺席，并不意味着关于开端有任何暗示；而当一个人说到\OriginalTerm{无源}{Unoriginate}时，超越了开端的事实已被表达，但关于终点的情况则存疑。\par

既然这两个名号同样有助于表达天主生命的永恒，那么现在是时候探究，为什么我们的朋友们将这永恒的完整意义一分为二，并宣告，那个否定开始的含义构成天主的存有（而非仅仅构成永恒定义的一部分），而他们却认为另一个否定终结的含义属于该存有之外的事物。很难看出将否定开始归于存有领域而将否定终结排斥在该领域之外的理由。这两者是我们对同一事物的概念；因此，要么两者都应被纳入存有的定义，要么，如果其中一个被认为不可接受，另一个也应被拒绝。然而，如果他们执意如此分割永恒的思想，并使其中一个属于该存有的领域，而将另一个归入神性的非实在物中（因为他们在此问题上所采纳的思想是卑下的，如同脱去羽毛的鸟儿，无法翱翔于神学的崇高境界），我建议他们颠倒他们的教导，将无终视为存有，反而忽略无始，并将胜利归于未来且激发希望的事物，而非归于过去与陈腐的事物。我说，鉴于（我如此说是因为他们心胸狭隘，并将讨论降低到孩童理解的水平），过去的人生时期对经历过它的人而言已无意义，他所有的兴趣都集中在未来和可期待的事物上；那么，无终之物将比无始之物更有价值。所以，让我们对天主本性的思想配得上且崇高；否则，如果他们要按照人类的试炼来判断它，就让他们更看重未来而非过去，并让他们将神性的存有限制在未来，因为时间的流逝将过去的一切存在一并冲刷，而期盼的存在则从我们的希望中获得实体。\par

现在我提出这些荒谬幼稚的建议，如同孩子们坐在市场上玩耍\BibleRef{路 7:32}；因为当人审视他们异端教导的卑下世俗时，不可能不陷入一种戏谑的幼稚。然而，将这一点补充到我们所说的内容中是恰当的，即：既然永恒的观念只有凭借两者（正如我们已经论证的），即通过对开始的否定和对终结的否定才能完成，那么，如果他们将天主的存有局限于其中之一，他们对这一存有的定义将明显不完善，被削减一半；它只被视为没有开始，而不将没有终结作为本质要素包含在自身之内。但如果他们确实将两种否定结合起来，从而完善他们对天主存有的定义，那么再次观察这个观点中立即显现的荒谬；在他们所有的努力之后，将会发现它不仅与独生子冲突，也与自身冲突。情况是清楚的，不需要多费唇舌。开始的观念与终结的观念彼此对立；各自的含义如同其他截然对立的概念一样相距甚远，其间没有中间命题。如果有人被要求定义\Quote{开始}，他不会给出一个与终结相同的定义；而是会将它的定义带入相反的极端。因此，这两个否定也将以同样的对立距离相互分离；无始之物，由于与通过开始而见到的事物相反，将与无终之物（或对终结的否定）非常不同。那么，如果他们将这些属性都引入天主的存有中，我指的是对终结和对开始的否定，他们会证明他们的这一神性是两种矛盾且不和谐之物的结合，因为与开始和终结相反的\Emph{相反者}也在它们那一边复制了开始与终结之间存在的\Emph{矛盾}。矛盾之物的相反者彼此也是矛盾的。事实上，这是一个永远正确的公理：自然与两个彼此对立的事物相对立的两样事物，它们本身也彼此对立；我们可以通过例子看到。水与火相对立；因此，破坏它们的力量也彼此对立；如果湿润易于灭火，而干燥易于破坏水，那么火与水的对立便延续到与它们相反的品质本身；所以干燥明显与湿润相对立。因此，当开始与终结必须被截然对立地放置时，与它们相反的术语在其含义上也相互矛盾，我指的是对终结和对开始的否定。那么，如果他们确定这些否定中只有一个指示存有（重复我之前的断言），他们将只见证天主存在的一半，将其局限于没有开始，而拒绝将其延伸到没有终结；然而，如果他们两者都引入对其定义中，他们实际上会像上述方式那样将其展示为矛盾的组合；因为对开始和对终结的这两个否定，凭借存在于开始与终结之间的矛盾，将把它分裂开来。因此，他们的神性将被发现是一种拼凑的复合物，一种矛盾的聚集体。\par

但是在天主的教会中，现在没有，将来也不会有一种教导，能使那唯一且单纯的不仅变成多样和拼凑，而且变成对立面的结合。真信仰的单纯性假定天主就是祂之所是，即：不能被任何词语、任何观念或我们领悟的任何其他方法所把握，不仅超越人类理智，也超越天使和所有超世界智慧的领域，不可思议，不可言说，超越一切言语表达，只有一个名字能代表祂本性的恰当，就是那\Quote{超越一切名号}的单一名字；这也赐给了\OriginalTerm{独生子}{Only-begotten}，因为\Quote{凡父所有的，都是子的}。正统理论允许这些词，即\Quote{\OriginalTerm{无始}{Ungenerate}}、\Quote{\OriginalTerm{无终}{Endless}}，用来指示天主的永恒，而非祂的本质；所以\Quote{无始}意味着在祂之外没有任何源头或原因，\Quote{无终}意味着祂的国度不会在任何终点止步。先知说：\Quote{你是一样的，你的年岁不会穷尽}，藉着\Quote{你}表明祂存在于没有任何原因之中，而藉着随后的话表明祂生命的福乐是永不停息、无穷无尽的。\par

但是，也许即使是很虔诚的人中，也会有人对我们关于天主\OriginalTerm{永恒}{Eternity}的这些探讨踌躇不前，并说从我们所说的话来看，保全对\OriginalTerm{独生子}{Only-begotten}的信仰将是困难的。他会说，我们的论述必然不可避免地会触及两种不可接受的教义之一。要么我们得出子是\OriginalTerm{无始}{Ungenerate}的，这是荒谬的；要么我们完全否认祂的永恒，而那种否认正是那\Emph{一伙}亵渎者的专长。因为如果永恒的特征是无始无终，那么必然要么我们不敬地否认子的永恒，要么我们在关于祂的内心思想中被引向\OriginalTerm{无始}{Ungeneracy}的观念。那么，我们将如何回答呢？那就是：如果在因果关系上设想父先于子时，我们在独生子存有之前插入了任何时间的标记，那么我们关于子永恒的信仰就有理由说是受到危害。但实际上，永恒的本性，无论是在父的生命中，还是在子的生命中，以及在我们关于圣神的信仰中，都不允许它会有终止的想法；因为时间不在的地方，\Quote{何时}也随之消灭。如果子总是与对父的思想一同显现，并且总是在存在的范畴中被发现，那么承认独生子的永恒有什么危险呢？祂\BibleQuote{没有生日的起始，也没有生命的终结} \BibleRef{希 7:3}。因为正如祂是光出自光，生命出自生命，善出自善，并且同样地，智慧、正义、力量等所有一切，所以祂也确实是永恒自永恒。\par

但是爱好争论争吵的人抓住这个论点，认为这样的推论会使祂成为无始自无始。然而，让他冷静他那好斗的心，坚持正确的表达，因为当他承认祂是\Quote{出自父？}时，他已经消除了关于\OriginalTerm{独生子}{Only-begotten}的一切\OriginalTerm{无始}{Ungeneracy}观念；那么称祂为永恒而非无始就没有危险了。一方面，因为子的存在不以任何时间间隔为标记，祂生命的无限性在万世之前和万世之后以渗透一切的潮流奔流，所以祂被恰当地冠以永恒的称号；另一方面，因为事实上和称号上关于祂作为子的思想，使我们不可分离地联想到父，因此祂避免了被归为\OriginalTerm{无始}{ungenerate}的存在，而祂总是与永远存在的父同在，正如我们导师受默感的话语所表达的：\Quote{藉着生成的方式与祂父的\OriginalTerm{无始}{Ungeneracy}相连}。我们对圣神的论述也将相同；区别仅在于秩序上安排的位置不同。因为正如子与父相连，从父获得存在，但在本质上并不在父之后，同样，圣神也与独生子相连，独生子在理论上的原因意义上，被设想在圣神存有之前。时间的延伸在永生中找不到入口；因此，当我们消除了原因的概念，天主圣三在任何方面都不显示与自身不和谐；荣耀应归于祂。\par

\SourceRecord{Translated by William Moore and Henry Austin Wilson.；Nicene and Post-Nicene Fathers, Second Series；Vol. 5.；Edited by Philip Schaff and Henry Wace.；Buffalo, NY: Christian Literature Publishing Co.,；1893.；<http://www.newadvent.org/fathers/290101.htm>.}

% structure: p0001 source=h1 target=section reason=page-title

\section{驳斥欧诺米（卷二）}

% structure: p0002 source=h2 target=subsection reason=relative-heading-rank; base=h2

\subsection{一、卷二宣明天主圣言的降生成人，以及主传授给祂门徒的信德；并指出那些试图推翻此信德并捏造其他附加名称的异端者，是出于他们的父亲魔鬼。}

基督徒的信德，按照我们主的命令，由祂的门徒传于万民，既非出于人，亦非借着人，而是出于我主耶稣基督自己。祂原是圣言、生命、光明、真理、天主和智慧，以及一切祂本性所是；祂为此尤其成为人的样式，并分担了我们的本性，在一切事上相似我们，只是没有罪过。祂在一切事上相似我们，因为祂取了完整的人性，包括灵魂和身体，这样我们的救恩便借着二者得以完成。\newline{}——我说，祂出现在世上，\BibleQuote{与人同居共处 \BibleRef{巴 3:37}}，为使人不再按照自己的观念对自有着持有意见，将模糊臆测的提示形成为教义，而为要使我们确信天主确已在肉身上显现，并相信那唯一真正的\BibleQuote{虔诚的奥迹 \BibleRef{弟前 3:16}}，乃是借圣言和天主亲自传给我们，祂亲自向宗徒们说话，并使我们领受关于天主超越本性的教导，这教导正如从较古老的圣经——法律、先知和智慧书——\BibleQuote{好像借着镜子观看，模糊不清 \BibleRef{格前 13:12}}，作为向我们完全启示的真理的证据，虔敬地接受所言之事的含义，以符合整部圣经之主所制定的信德，这信德我们按所领受的、字字纯粹地持守，不加篡改，甚至认为与所传给我们的话语稍有偏离便是极端的亵渎和不敬。\newline{}我们便这样相信，正如主向祂门徒制定信德时所说：\BibleQuote{你们要去使万民成为门徒，因父、及子、及圣神之名给他们授洗 \BibleRef{玛 28:19}}。这是奥迹之言，借着从上而来的新生，我们的本性从可朽坏转变为不朽坏，从\BibleQuote{旧人}更新为\BibleQuote{按那创造者的肖像}，即起初按天主的相似受造。\newline{}因此，在这由天主传给宗徒的信德中，我们既不删减，也不更改，也不增加，确知那妄图以诡诈的巧辩曲解天主之言的人，\BibleQuote{是出于他的父亲魔鬼}，他离开了真理的话语，\BibleQuote{凭自己的私欲说话}，成为谎言之父。\BibleRef{若 8:44} 因为凡所说的若不与真理完全相符，那必是虚假而非真实。\par

% structure: p0004 source=h2 target=subsection reason=relative-heading-rank; base=h2

\subsection{二、额我略随后详加说明关于永恒圣父、圣子和圣神。}

既然这道理是由真理本身所提出，那么，凡是制造有害异端的人为了推翻这神圣言论而发明的任何说法——例如，称圣父为子的\Quote{\OriginalTerm{造物主}{Maker}}和\Quote{\OriginalTerm{创造者}{Creator}}而非\Quote{圣父}，称圣子为\Quote{\OriginalTerm{结果}{result}}、\Quote{\OriginalTerm{受造物}{creature}}、\Quote{\OriginalTerm{产品}{product}}而非\Quote{圣子}，称圣神为\Quote{\OriginalTerm{受造物的受造物}{creature of a creature}}和\Quote{\OriginalTerm{产品的产品}{product of a product}}而非其本有称号\Quote{圣神}，以及那些对抗天主的人所乐意对祂说的任何话——所有这一切妄想，我们都称之为对这道理所启示给我们的天主性的否认和违背。因为我们曾一劳永逸地从主那里学到（借着祂，我们的本性从必死转变为不死）——我说，从祂那里我们学到应该用理智的眼目注视什么——就是圣父、圣子和圣神。我们说，曲解这些神圣言论并编造断言来取代以推翻它们——这些断言妄图纠正天主圣言，是祂规定我们应将这陈述持守为信德的一部分——这是一件可怕且毁灭灵魂的事。因为每一个称号按其自然意义理解，对基督徒而言，都成为真理的准则和虔诚的法则。因为虽然在历史书、先知书和法律中有许多其他名字用来指示天主性，但我们的主基督却越过所有这些，将这三个称号交给我们，认为它们更能引领我们达致关于\OriginalTerm{自有者}{Self-Existent}的信仰，宣告说，我们只需紧守\Quote{圣父、圣子、圣神}这个称号，就足以领悟那绝对存有者，祂是一位，却又不是一位。就本质而言，祂是一，所以主规定我们应仰望一个名字；但就标示位格的属性而言，我们对祂的信仰分为对圣父、圣子和圣神的信仰；祂被分开而不分离，结合而不混淆。因为当我们听到\Quote{圣父}这个称号时，我们领悟到其含义是：这个名字不只是指自身，而且也以其特殊含义指示与圣子的关系。因为\Quote{圣父}这个术语本身没有意义，除非在说出\Quote{圣父}一词时也暗含着\Quote{圣子}。因此，当我们学到\Quote{圣父}这个名字时，借着同一个称号，我们也同时被教导了对圣子的信仰。既然天主性就其本质而言，在一切与本质相关的事上，永恒不变地保持同一，从未在任何时候缺乏它现在所是的任何东西，也不会在未来的任何时候成为它现在所不是的任何东西；既然那正是圣父的一位被圣言称为圣父，既然在圣父中蕴含着圣子——既然这些事如此，我们必然相信，祂的本性不容任何变化或改变，祂始终完全是祂现在所是，或者说，如果有什么祂过去不是的，那东西现在肯定也不是。既然祂被圣言本身称为圣父，祂肯定永远是圣父，现在是，将来也是，就如祂曾是。因为确实，在谈论神圣且无瑕疵的本质时，否认美善的东西永远属于它，是不合法的。因为如果祂不总是祂现在所是，祂肯定是从较好变为较差，或从较差变为较好；这两种说法，无论哪一个关于天主本性的陈述，其不虔敬程度都是一样的。但事实上，天主性是不能变化和改变的。因此，一切美善的事物总是被默观于美善的泉源中。但\BibleQuote{\OriginalTerm{独生子天主}{Only-begotten God}，在圣父怀里的那位}是美善的，且超乎一切美善：——请注意，祂说：\BibleQuote{那位\Emph{是}在圣父怀里的}，而不是\Quote{那位在那里成了}。\par

那么，这些证明已经表明：圣子从永远就在圣父内被默观，祂在圣父内，是生命、光明、真理，以及一切尊贵的名号和概念——若说圣父曾独自存在，脱离这些属性，那便是最大的不敬和愚妄。因为如果圣子就如圣经所说，是天主的德能、智慧、真理、光明、圣化、和平、生命等，那么按异端者的观点，在圣子存在之前，这些也根本不存在。若这些不存在，他们必然认为圣父的胸怀缺乏这些美善。因此，为使圣父不被视为缺乏自己的美善，为使教义不陷入这种极端，对圣子的正确信仰必然包含在主的话语中，与对圣父永恒的默观相联。为此，祂略过了所有用于表示天主性超越尊荣的名号，而将\Quote{父}这一称号作为我们信仰宣认的一部分传授给我们，这称号更适于表明真理；正如所说，这称号通过其相对意义本身就蕴含着圣子，而圣子在圣父内，始终是祂本质所是，正如已经说过的，因为天主性按其本性不容增衍。它不察觉自身之外有任何美善，能借参与而获得任何增加，而是永远不变，既不丢弃所有，也不获取所无：因为它的属性没有一样是可丢弃的。如果任何有福、无玷、真实、美善的事物与祂相连并在祂内，我们必然看到，善而圣的圣神必属于祂，并非通过增添。那圣神无疑是君王的圣神、赋予生命的圣神、一切受造的控制和圣化力量，随祂意愿而\BibleQuote{在一切内运行一切}。\BibleRef{格前 12:6} 这样，我们在受傅的基督与祂的傅油之间、在君王与祂的王权之间、在智慧与智慧之神之间、在真理与真理之神之间、在德能与德能之神之间，不设想任何空隙；而是如同从永远在圣父内默观到圣子——祂是智慧、真理、谋略、德能、知识、理解——同样在祂内也默观到圣神，祂是智慧之神、真理之神、谋略之神、理解之神，以及圣子是并被称为的一切。为此，我们说，神圣的宗徒们接受了虔敬的奥迹，以一种同时表达合一与区别的形式——要我们相信父、子及圣神之名。因为\OriginalTerm{自立}{subsistence}的区分使位格的区别清晰而不混淆，而信仰宣认开头的那独一之名清楚地向我们阐明了信仰所宣认的位格——即父、子及圣神——的本质合一。因为这些称谓教导我们的不是本性的差异，而只是标志\OriginalTerm{自立}{subsistence}的特殊属性，使我们知道：父不是子，子不是父，圣神既不是父也不是子，而是凭借各自的\OriginalTerm{位格自立}{Personal Subsistence}的独特标志认识每一位，在无限圆满中，既各自被默观，又不与所关联者分离。\par

% structure: p0007 source=h2 target=subsection reason=relative-heading-rank; base=h2

\subsection{三、\Person{国瑞}继而讨论天主圣三那\OriginalTerm{不可言说的名}{unnameable name}的相对力量与各位格的相互关系，此外还有本质的不可知性，以及祂对我们方面的俯就、祂由童贞女所生、祂的第二次来临、从死者中复活及将来的报应。}

那么，主所说的那不可名之名是什么意思呢？主说：\BibleQuote{给他们授洗，因那名}，却没有加上\Quote{那名}所指的实际意义之词。关于此事，我们有这样的看法：受造界中一切事物都由各自的名称来界定。因此，每当一个人提到\Quote{天}，他就将听者的概念引向这名所指的受造物；提到\Quote{人}或某种动物，一说出名称就立刻在听者心中印上那受造物的形象；同样，其他一切事物都借着加于其上的名称，描绘在听者的心中，他通过听闻接受了加于该事物的称谓。唯独不受造的自然，我们在父、子及圣神内所承认的，超越了一切名称的意义。因此，圣言在传授信仰时说到\Quote{那名}，并没有加上它是什么——因为怎能给那超乎万名之上的找到名称呢？——而是授权：无论我们的理智借着虔诚的努力能够发现什么名称来指示那超越的自然，那名称都应同等地应用于父、子及圣神，无论是\Quote{善}还是\Quote{不朽}，无论各人认为什么名称合适用以指示神性无玷的自然。藉着这一传授，圣言在我看来为我们立下了这条法则：我们该确信天主本质是不可言说、不可理解的；因为显然，\Quote{父}的称号并不向我们呈现本质，而只指示与子的关系。那么，如果可以教给人类本性天主本质，那\BibleQuote{愿意所有的人都得救，并得以认识真理\BibleRef{弟前 2:4}}的一位就不会压下关于此事的认识。但事实上，祂对天主本质只字未提，表明对此的认识超越我们的能力；当我们学到了我们所能及的，我们就不需要那超越我们能力的认识，因为我们在所交付给我们的教义的信德宣认中，拥有足以拯救我们的东西。因为认识到祂是绝对存在者，同时借着这名称的关联意义也宣示了子的尊威，就是最完全的虔敬教导；子，如前所述，在自己内紧密地结合着生命与真理的圣神，因为祂本身就是生命与真理。\par

这些区分既已确立，我们在神圣教义领域内弃绝一切异端的幻想，我们相信，正如我们曾受主的教诲，因父及子及圣神之名，在承认这一信仰的同时，也承认创造主为人类所实行的救恩安排。因为祂\BibleQuote{虽具有天主的形体，却没有将自己与天主同等视为应当强夺的，反而使自己空虚，取了奴仆的形体\BibleRef{斐 2:6}}，并在童贞玛利亚内降生成人，将我们从\Quote{我们被囚于其中}的死亡中赎回，\Quote{卖给了罪恶}，为赎回我们的灵魂献出了祂宝贵的血，祂藉十字架倾流了这血，并亲自为我们显明了从死者中复活的道路，祂要在自己的时候，以父的荣耀来临，按正义审判每一个灵魂，那时\BibleQuote{凡在坟墓里的，都要听见祂的声音，并要出来：行过善的，复活进入生命；作过恶的，复活进入审判}。但为了使欧诺弥目前广泛散布的有害异端不致落入某些较单纯之人的心中，且未经查考而损害无伪的信德，我们不得不陈述他们所传播的信条，并努力揭露他们教导的邪恶。\par

% structure: p0010 source=h2 target=subsection reason=relative-heading-rank; base=h2

\subsection{四、他接着巧妙驳斥了欧诺弥关于绝对存在者的片面、空洞且亵渎的陈述。}

他们教义的措辞如下：\Quote{我们相信独一真实的天主，按照主自己的教导，不以虚谎的尊号尊崇祂（因为祂不能撒谎），而是真实存在的，在本性和荣耀上是独一的天主，无始、无终、永恒、独自存在。} 那自称按照主的教导而相信的人，切莫按自己的幻想曲解那关乎万有之主的信仰解释，而应亲自追随真理之言。既然信仰的表述包含圣父、圣子和圣神的名，那么他们这种构造与主的言论有何相符之处，竟将这样一种教义归诸那些言论的教导？他们无法证明主在\WorkTitle{福音}中何处说过我们应当相信\Quote{独一真实的天主}，除非他们有什么新的福音。因为从古至今在教会中一直诵读的\WorkTitle{福音}，并不包含这个告诉我们应当相信或受洗归于\Quote{独一真实的天主}的说法，正如这些人所说，而是\Quote{因父、及子、及圣神之名。}但我们按照主的亲口教导说：\Quote{一}这个词所指的不仅是圣父，而是在其含义中包含了圣子与圣父，因为主说：\BibleQuote{我与父原是一体。}同样地，\Quote{天主}这个名也同等属于\Quote{圣言所在之元始}和\Quote{在元始的圣言。}因为福音作者告诉我们：\BibleQuote{圣言与天主同在，圣言就是天主。}因此，当表达天主性时，圣子与圣父同被包括在内。此外，真实者不能被视为与真理相异且无关的东西。但是，主就是真理，这是没有人会争辩的，除非他是与真理疏远的人。那么，如果圣言在\Quote{一}之内，且是天主和真理，正如\WorkTitle{福音}所宣告的，那么这些使用区别性用语的人以主的什么教导为基础来建立他的教义？因为对立在于\Quote{独}与\Quote{不独}之间，在于\Quote{天主}与\Quote{非天主}之间，在于\Quote{真实}与\Quote{不真实}之间。如果他们对措辞的区分是针对偶像的，我们也同意。因为\Quote{神性}一名是在歧义的意义上给予外邦人的偶像，因为\Quote{外邦人的众神都是魔鬼}，而另一个意义则标示出独一与众多、真实与虚假、那些不是神的与那位是神的之间的对比。但如果这对比是与独生天主相对，愿我们的贤士明白：真理仅以虚假为对立，天主仅以非天主为对立。但既然为主且为真理的那一位是天主，祂在父内，并相对于父是\Quote{一}，那么在真教义中就没有这些措辞区分的余地。因为真正相信\Quote{一}的人，在\Quote{一}内看见那在真理、天主性、本质、生命、智慧以及一切属性上与祂完全合一的那一位；否则，如果他在\Quote{一}内没有看见那拥有一切的那一位，那么他什么也不信。因为没有圣子，圣父便既不存在也无名，正如有能力者没有能力，或有智慧者没有智慧一样。因为基督是\BibleQuote{天主的德能和天主的智慧} \BibleRef{格前 1:24}；因此，那自以为看见独一天主却脱离能力、真理、智慧、生命或真光的人，要么什么也看不见，要么所见的一定是恶。因为善的属性之撤除，成为恶的设定与起源。\par

他说：\Quote{不尊崇祂}，\Quote{以虚谎的尊号，因为祂不能撒谎。} 我祈求\Person{欧诺弥}能恪守那句话，从而为不能撒谎的真理作证。因为如果他持这种心态，认为主所说的一切都远离虚谎，那么他当然会被说服：那位说\BibleQuote{我在父内，父在我内}——显然，一者完全在另一者内，父并不溢出于子，子也不缺失于父——又说子应受尊敬如父所受的尊敬，以及\BibleQuote{看见我的，就是看见了父}，以及\BibleQuote{除了子，没有人认识父}，在所有这些段落中，对于那些将这些宣告视为真实的人来说，丝毫没有提示父与子之间在荣耀、本质或任何其他方面有任何差异。\par

\Quote{真实存在，}他说，\Quote{本性与光荣中的独一天主。}真实存在是与不真实存在相对立的。每个存在的事物就其存在而言是真实存在的；但那些仅凭表象和暗示看似存在而实际不存在的，则不是真实存在的，例如梦中的景象或画中的人。因为这些以及类似的事物，尽管就表象而言存在，却没有真实存在。那么，如果他们按照犹太人的意见，主张独生子天主根本不存在，那么他们将真实存在仅归于圣父是对的。但如果他们不否认万物的创造者的存在，就让他们满足于不剥夺那自有者的真实存在，祂在天主向\Person{梅瑟}显现时自称为自有者，说：\Quote{我是自有者；}正如\Person{欧诺弥}在后面的论证中也同意这一点，说那向\Person{梅瑟}显现的就是祂。然后他说天主是\Quote{本性上与光荣中的独一。}天主是否存在而不本性上是天主，说这些话的人或许知道；但如果本性上不是天主的就不是天主，那么让他们从伟大的\Person{保禄}那里学到，那些侍奉不是天主的假神的人并不是侍奉天主。但我们\Quote{侍奉生活而真实的天主，}正如宗徒所说\BibleRef{得前 1:10}；我们所侍奉的就是耶稣基督。因为宗徒\Person{保禄}甚至以侍奉祂为荣，说：\Quote{\Person{保禄}，耶稣基督的仆人\BibleRef{罗 1:1}。}那么，我们不再侍奉那些本性上不是天主的假神，已经认识了那本性是天主的一位，万膝都向祂跪拜，\Quote{天上、地上和地下的一切。\BibleRef{斐 2:10}}但我们不会成为祂的仆人，除非我们相信这就是生活而真实的天主，\Quote{万口都承认耶稣基督是主，以光荣天主圣父。\BibleRef{斐 2:10}}\par

\Quote{天主，}他说，\Quote{无始，永恒，无终，独一。}再一次，\Quote{愚蒙人，你们要明白，}正如\Person{撒罗满}所说，\Quote{他的诡诈，}免得你们被欺骗，坠入否认独生子天主性的深渊。无终是指不容死亡与腐朽；同样，永恒不仅指暂时。因此，那既非永恒亦非无终的，必然在可朽坏、可死的本性中显现。因此，那将\Quote{无终}归于独一天主，却不将圣子包含在\Quote{无终}和\Quote{永恒}的宣告中的人，凭借这样的命题，主张祂所如此与永恒无终相对照的那位是可朽坏而暂时的。但我们，即使被告知天主\Quote{独有不死，\BibleRef{弟前 6:16}}，也把\Quote{不死}理解为圣子。因为生命就是不死，而主就是那生命，祂说：\Quote{我是生命。}如果祂被称为住在\Quote{不可接近的光中，\BibleRef{弟前 6:16}}，我们同样毫不困难地理解为那不可为虚假所接近的真光就是独生子，在祂内我们从真理本身知道圣父就在其中。在这些意见中，让读者选择更虔诚的：我们应以配得上天主性的方式思考独生子，还是像异端所规定的那样称祂为可朽坏而暂时的。\par

% structure: p0015 source=h2 target=subsection reason=relative-heading-rank; base=h2

\subsection{五、接着他惊人地推翻了\Person{欧诺弥}那些不可理解的陈述，这些陈述断言圣父的本质不分隔、不分开，也不变成别的东西。}

\Quote{我们信天主，}他告诉我们，\Quote{在祂作为独一的本质中，不分隔为多个，或有时成为这个，有时成为那个，或从祂之所是改变，或从一个本质转变为另一本质而采取三重位格的样式：因为祂永远是绝对独一，一致而不变地存留为独一天主。}从这些引文中，谨慎的读者首先可以将插入陈述中毫无意义的废话与那些似乎有些意义的区分开，然后检验他陈述中剩余部分可发现的含义，看它是否与对基督应有的尊敬相符。\par

因此，所引述的言论中第一个完全脱离了任何有意义的意思，无论好坏。因为欧诺米本人无法告诉我们这些话语中有什么意义：\Quote{就祂是一体的本质而言，不分化为多个，或有时成为一个，有时成为另一个，或从祂所是改变}，我认为他的任何盟友也无法在这些话语中找到任何含义的影子。当他说祂\Quote{就祂是一体的本质而言，不分隔}，他要么说祂不与自己的本质分隔，要么说祂自己的本质不与祂分开。这个无意义的陈述不过是噪音和空声的随意组合。为什么要花时间研究这些无意义的表达呢？因为当一个人与他自己的本质分隔时，他如何还能存在？或者，任何东西的本质如何能被分开并分别显示？或者，一个人如何能离开他所在之处，变成另一个，走出自身？但他又加上：\Quote{不从一种本质转变，假装三个位格：因为祂永远是绝对的一，恒常不变地是唯一的天主。}我认为他陈述的无意义对每个人而言不言自明，无需我一言：让任何认为他所说的话有意义、有道理的人来反驳吧；凡是能辨别语句力量的人都会避免与无形的影子争斗。因为说\Quote{就祂是一体的本质而言，不分隔或分化为多个，或有时成为一个，有时成为另一个，或从一种本质转变，假装三个位格？}，这对我们的教义有何作用？——这些东西基督徒既不说也不信，也无法从我们所宣认的真理中推断出来。因为在天主的教会中，谁曾说过或听别人说过，圣父就其本质而言被分隔或分裂，或有时成为一个、有时成为另一个，走出自身，或假装三个位格？这些是欧诺米对自己说的，不是与我们辩论，而是把他自己的垃圾串在一起，在他言谈的不敬中混杂了大量荒谬。因为我们说，称受造界的主宰为受造物，以及认为圣父在祂之所是中分隔、分裂、走出自身，或像粘土或蜡般被塑造成各种形状而假装三个位格，同样是邪恶和亵渎天主的。\par

但让我们审视接下来的话：\Quote{祂永远是绝对的一，恒常不变地是唯一的天主。}如果他在谈论圣父，我们同意他，因为圣父最真实地是一，独自且永远是绝对统一而不可改变，无论现在还是将来，从不停止是祂所是。因此，如果这样的断言是关于圣父的，就让他不要与虔诚的教义争辩，因为在这点上他与教会一致。因为凡宣认圣父常是且不变地相同，是一且唯一的天主的人，就持守了虔诚的话，只要他在圣父中看见圣子，没有圣子，圣父既不存在也不被命名。但如果他在圣父之外发明另一位天主，就让他与犹太人，或那些被称为\OriginalTerm{希普西斯特派}{Hypsistiani}的人争辩，在这种人与基督徒之间的区别在于：他们承认有一位他们称为至高者或全能者的天主，但不承认祂是圣父；而基督徒，若不信圣父，就根本不是基督徒。\par

% structure: p0019 source=h2 target=subsection reason=relative-heading-rank; base=h2

\subsection{六、他随即显示圣子与圣父的合一，以及欧诺米对圣经缺乏理解和认识。}

他接着在此之后添加的话如下：——\Quote{没有分享者，}他说，\Quote{在他的天主性中，没有他荣耀的分割者，没有人有份于他的权能，或在他的王座上有一份：因为他是独一的天主，全能者，万神之神，万王之王，万主之主。}我不知道\Person{欧诺米}在抗议圣父不容许任何人分享他自身的天主性时，指的是谁。因为如果他使用这样的说法是指虚妄的偶像和那些敬拜偶像者的错误观念（正如\Person{保禄}向我们保证基督与\Person{贝里雅尔}之间没有和谐，天主的殿与偶像之间没有相通\BibleRef{格后 6:15}），我们就同意他。但如果他通过这些断言意在将独生天主与圣父的天主性分开，那么就让他知道，他正给我们提供一个可以用来反对他自己以驳斥他自己不虔敬的两难处境。因为他要么完全否认独生天主是神，以便为圣父保留那些（按照他的说法）不能与圣子分享的神性特权，因而因否认基督徒所敬拜的天主而被定罪为悖逆者；要么如果他承认圣子也是神，却不同意在本质上与真天主一致，那么他就必然不得不承认他主张诸神因本性不同而彼此分离。让他选择他愿意选择哪一个——要么否认圣子的天主性，要么在他的信条中引入众多神祇。因为无论他选择哪一个，在渎神方面都是一样的：因为我们这些通过圣经中天主默感的话语被引入虔敬奥秘的人，在圣父与圣子之间看到的不是天主性的合伙关系，而是合一，因为主用他自己的话教导了我们这一点，当他说：\BibleQuote{我与父原是一体}，以及\BibleQuote{看见了我就是看见了父}。因为如果他不与圣父同性，他怎么能在自己内拥有那不同的事物\BibleRef{若 17:10}？或者，如果那外来的、异己的本性没有接受那与自己不同类者的印记，他又怎么能在自己内显出那不像者呢？但是他说：\Quote{他也没有他荣耀的分割者。}在此他按事实说话，尽管他不知道自己在说什么：因为圣子不与圣父分割荣耀，而是拥有圣父的全部荣耀，正如圣父拥有圣子的一切荣耀。因为他对圣父这样说：\BibleQuote{我的一切都是你的，你的一切都是我的}\BibleRef{若 17:10}。因此他也说，他将在审判日\BibleQuote{带着父的荣耀}\BibleRef{谷 8:38}显现，那时他将按各人的行为予以回报。通过这句话，他显示存在于他们之间的本性合一。因为正如\BibleQuote{日有日的荣耀，月有月的荣耀}\BibleRef{格前 15:41}，因为那些光体本性的差异（因为如果两者有同样的荣耀，就不会被认为在本性上有任何差异），同样，那预言自己将带着父的荣耀显现的那一位，藉着荣耀的同一性，指示了他们的本性相通。\par

但说圣子在祂圣父的王座上无分，这显示欧诺弥对天主的神谕进行了极大量的研究；他虽然极其热爱默感的圣经，却尚未听到\BibleQuote{你们要追求上面的事，那里有基督坐在天主右边}\BibleRef{哥 3:1}，以及许多类似段落，其数量不易计算，但欧诺弥从未学过，因此否认圣子与圣父同坐王座。再次，对于\Quote{在祂的权能中无分}这样的说法，我们宁愿视其为无意义而不加理会，也不愿将其驳斥为不虔敬。因为从该词的通常用法中，不难发现\Quote{有分}这一措辞附着何种意义。如圣经告诉我们，那些人为主的衣服拈阄，他们不愿撕裂祂的衣服，而愿意将衣服交给他们中阄所决定的那一人。\BibleRef{若 19:23}因此，那些这样为\Quote{长衣}彼此拈阄的人，或许可说是在它里面\Quote{有分}了。但这里，关于圣父、圣子、圣神，既然他们的权能居于他们的本性中（因为圣神\Quote{随意}吹，且\Quote{按祂的意愿在一切内行一切}；圣子，万物藉祂而造，看得见的、看不见的，在天上的、在地上的，\Quote{按祂的喜悦行了万事}，且\Quote{祂愿意复活谁就复活谁}；圣父将\BibleQuote{时期置于祂自己的权能下}\BibleRef{宗 1:7}，而从\Quote{时期}的提及，我们推断所有在时间内所行的事都隶属于圣父的权能），如果我这样说，已经表明圣父、圣子、圣神皆处于有能力行其所愿的地位，那么\Quote{在祂的权能中有分}这句话能有什么意义，就不可理解了。因为万有的嗣子、万代的创造者，祂闪耀着圣父的荣耀，在祂自己内表达圣父的位格，拥有圣父自己所拥有的万有，并拥有一切属于圣父的权能；并非权能从圣父转移给圣子，而是它既留在圣父内，也居于圣子内。因为那在圣父内的，显然带着祂自己的全部大能而在圣父内；那在自己内有圣父的，包含圣父的全部权能与大能。因为祂在自己内有整个圣父，而不仅是祂的一部分：那完全拥有祂的，也必然拥有祂的权能。那么，欧诺弥以何种意义断言圣父\Quote{没有人有分于祂的权能}？那些是他愚昧的门徒的人或许能告知；一个懂得鉴赏语言的人承认，他无法理解脱离意义的短语。他说：\Quote{圣父没有人有分于祂的权能}。为何？有谁说圣父与圣子为权能彼此争竞，并借拈阄决定此事？但神圣的欧诺弥前来作他们之间的调停者，通过一个友好的协定，不经抽签就将权能的优越性派给了圣父。\par

我请你们留意，他对自己信条的这种卑下阐释是何等荒谬幼稚。什么！那\Quote{以自己大能的话支撑万有，}\BibleRef{希 1:3}，祂说什么就成就什么，以命令本身的力量成就祂所愿的，祂的能力不迟于祂的意愿，祂的意愿就是祂能力的尺度（因为\Quote{祂一发言，万有造成；祂一下令，诸天形成}），祂独自创造万有，并使万有在祂内得以存立，没有祂，任何存在之物既不能产生也不能存留——祂竟要靠某种分配的过程来获得祂的能力！请听者判断，说这话的人是否神志清醒。他说：\Quote{因为祂是唯一的天主，全能者。}如果他用\Quote{全能者}的称号是指圣父，那么他用的语言是我们的，并非陌生的语言；但如果他是指圣父以外的另一位天主，那么让这位犹太教义的赞助者随他高兴也宣讲割礼吧。因为基督徒的信德是朝向圣父的。圣父就是这一切——至高者、全能者、万王之王、万主之主，总之，一切最高意义的词语都属于圣父。但圣父的一切也是圣子的；所以，基于这种理解，我们也接纳这个短语。但如果他离开圣父，谈论另一位全能者，那么他是在说犹太人的语言，或是追随柏拉图的思想——因为据说那位哲学家也肯定有一位创造者，在至高之处造了一些次等的神祇。正如在犹太和柏拉图的观点中，不相信天主圣父的人不是基督徒，即使他在信经中声称一位全能的天主；同样，欧诺米也虚假地自称基督徒之名，心向犹太人，或在披着基督徒称号的外衣下宣扬希腊人的教义。至于接下来的几点，他说同样的解释也适用。他说祂是\Quote{万神之神}。我们通过加上圣父的名字使这宣告成为我们自己的，因为我们知道圣父是万神之神。但属于圣父的一切当然也属于圣子。\Quote{万主之主}。同样的解释也适用于此。\Quote{全地之上的至高者}。是的，因为你无论想到三位中的哪一位，祂都是全地之上的至高者，因为对地上事物的眷顾由圣父、圣子和圣神同样从上行使。同样，上面的话之后是\Quote{诸天之上的至高者，至高的至高者，属天的，在祂所是上真实并如此持续，言语真实，作为真实}。为何，基督徒的目光在圣父、圣子和圣神身上同样看到这一切。如果欧诺米确实把这些只归于信经中所承认的一位，让他胆敢称那说过\Quote{我是真理}者为\Quote{言语不真实}，或称真理之神为\Quote{言语不真实}，或者拒绝将\Quote{作为真实}的称号给予那施行正义和审判的祂，或给予那随己意在众人内行万事的圣神。因为如果他不承认这些属性属于信经中传授给我们的各位，那么他就是在完全取消基督徒的信经。因为怎能有人会认为那言语虚假、作为不真实者堪当信德的对象呢？\par

但我们继续看下文。他说：\Quote{超越一切统治、服从和权柄。}这话是我们的，也恰当地属于天主教会——相信天主性超越一切统治，并相信在万有中一切可设想的事物都隶属于祂。但圣父、圣子和圣神构成天主性。如果他把这一属性只归于圣父，如果他断言只有祂是免于可变和改变，如果他说只有祂是无玷的，那么其暗示之意是明显的，即那些不具备这些特性的东西是可变的、可朽的、会改变和腐朽。那么，这就是欧诺米对圣子和圣神的断言：因为如果他不是对圣子和圣神持这种观点，他就不会使用这种对立，将圣父与他们对比。至于其余，弟兄们，请判断，持这些观点的人是否不是基督徒信德的迫害者。因为谁会允许认为那变化、改变、会腐朽的东西堪当崇敬呢？所以，怀有这样观念的人——这些观念使他断言真理本身和真理之神不是无玷的、不变的、不可改变的——其全部目的就是要把对圣子和圣神的信仰从教会中驱逐出去。\par

% structure: p0024 source=h2 target=subsection reason=relative-heading-rank; base=h2

\subsection{七、格列高利进一步表明，独生子不仅由圣父所生，而且由圣神从童贞玛利亚无痛苦地受生，并不分割本体；因为人的本性也不是因生育而从父母分割或分离，这从亚当和亚伯拉罕的例子被巧妙地证明。}

现在让我们看看他在之前的陈述上又添加了什么。他说：\Quote{不是藉着\OriginalTerm{生育}{begetting}分割自己的\OriginalTerm{本质}{essence}，并且同时是生育者与被生育者，同时是圣父与圣子；因为他是\OriginalTerm{不朽}{incorruptible}的。}或许，先知论及不敬虔之人时所说的\BibleQuote{他们编织蜘蛛网}\BibleRef{依 59:5}正是此意。因为在蜘蛛网中，有编织的形貌，但形貌中没有实质——触摸它的人触摸不到实质，因为蜘蛛丝触指即破——空泛言语的毫无实质的纹理也正是如此。\Quote{不是藉着生育分割自己的本质，并且同时是生育者与被生育者。}我们应当称他的话语为论证，还是称它们为某种水肿膨胀所分泌的体液？因为\Quote{藉着生育来分割自己的本质，并且同时是生育者与被生育者}是什么意思？谁如此心神散乱，谁如此癫狂，以至于说出欧诺弥认为他正在与之作战的陈述？因为教会相信真正的圣父确实是其独生子的真正圣父，如宗徒所说，而不是一个与他疏离的儿子的父。因为他在一封书信中如此宣告：\BibleQuote{他连自己的儿子都没有顾惜}\BibleRef{罗 8:32}，藉着\Quote{自己的}这一附加词，将他与那些因恩宠而非因本性而堪当义子名分的人区分开来。但是，贬低我们这信仰的人说了什么呢？\Quote{不是藉着生育分割自己的本质，也不是同时是生育者与被生育者，同时是圣父与圣子；因为他是不朽的。}一个在福音中听到\BibleQuote{圣言在起初就存在，且是天主}，并且圣言出自父的人，竟然用这些卑劣恶臭的观念玷污无瑕的教义，说\Quote{他不是藉着生育分割自己的本质}？可耻啊，这些卑劣肮脏念头的憎恶！这样说的人怎么不明白，天主在肉身显现时，为了他自身身体的形成，并没有接纳人性的条件，而是藉着圣神和至高者的能力，为我们生为婴孩；童贞玛利亚也不受那些条件的影响，圣神也没有减少，至高者的能力也没有分割？因为圣神是完整的，至高者的能力保持不减；婴孩以我们本性的丰满而生，并没有玷污他母亲的不朽。那时，肉体从肉体生出，没有肉欲；然而欧诺弥却不承认荣耀的光辉是从荣耀本身发出的，因为荣耀在生出光时既不减少也不分割。再者，人的言语从心灵生出，没有分割，但天主圣言却不能从圣父生出而不导致圣父的本质被分割！谁如此愚蠢，以至于看不出他的主张不合情理？\Quote{不是藉着生育分割自己的本质。}为什么？谁的本质是由生育分割的呢？因为在人身上，本质意味着人的本性；在畜生身上，一般而言意味着畜生的本性；但在牛、羊和所有畜生身上，具体而言，是按照它们种类的区别来看待的。其中哪一个生物藉着生殖过程分割了自己的本质呢？难道本性不总是通过后代延续保持完整吗？再者，人从自己生出一个人，并没有分割自己的本性，反而保持完整，既在生育者身上，也在被生育者身上，不是从一方分割转移给另一方，也不是当一方完全形成时在另一方被截断，而是同时完整地在先者中存在，完整地在后者中可被发现。因为既在生育孩子之前，人是理性的、可朽的、有智力和知识的动物，在生育具有这些特质的人之后也是如此；因此，他身上显示了本性的一切特殊属性；他并没有因为生育出自自己的人而失去自己作为人的存在，反而在那事件之后仍保持先前所是，而由他所出的人的出现并没有给源自他的本性造成任何减少。\par

然而，人由人生出，生育者的本性并未被分割。但欧诺弥却不承认那位在圣父怀里的\OriginalTerm{独生子天主}{Only-begotten God}真正出自圣父，恐怕是由于\OriginalTerm{位格}{subsistence}而损伤了圣父不可侵犯的本性：但他在说了\Quote{不是藉着生育分割自己的本质}之后，又添加了\Quote{或自身是生育者与被生育者，或自身成为圣父与圣子}，并以为藉着这些松散不连贯的短语可以削弱对虔诚的真实宣认，或为自己的不虔诚提供某些支持，却没有意识到，正是他用来构造\OriginalTerm{归谬法}{reductio ad absurdum}的手段，暴露了他其实是真理的辩护者。因为我们也说：凡拥有属于自身圣父的一切的那一位，他就是他所是的一切，除了是圣父；凡拥有属于圣子的一切的那一位，在他自身内完全显示出圣子，除了是子：所以，欧诺弥在此发明的归谬法，当我们遵照福音的指引，扩展这个概念以更清楚地展示它时，反而成了真理的支持。因为如果\BibleQuote{看见了子的人就看见了父}，那么圣父就生出了另一个自我，没有走出自身之外，同时在他内完整地显现：因此，从这些考虑来看，那些似乎针对虔诚而说的话，被证明是健全教义的支持。\par

但他却说：\Quote{圣父藉着生子并不分散自己的本质，而是同时是生者与受生者，同时是圣父与圣子；因为祂是不朽坏的。} 多么有力的结论！高明的先生，您这是什么意思？因为祂是不朽坏的，所以祂在生圣子时并不分散自己的本质：祂也不生于自己，也不生自己，也不因为是自己的父亲和儿子而成为自己的父亲和儿子，因为祂是不朽坏的。那么，如果任何人具有可朽坏的本性，他就在生子时分散自己的本质，并且由自己而生，也生自己，是他自己的父亲和儿子，因为他是可朽坏的。如果这样，那么亚巴郎因为是可朽坏的，就并没有生依市玛耳和依撒格，而是由婢女和合法妻子生了自己；或者，按照这论点的其他诡辩，他把自己分散在从他而生的儿子们中间，首先，当哈加尔给他生了一个儿子时，他被分成了两部分，一半成了依市玛耳，而另一半仍然是一半的亚巴郎；随后，亚巴郎本质的残余部分再次被分，在依撒格中取得了实在。因此，亚巴郎本质的四分之一被分成依撒格的双胞胎儿子，于是每个孙子里有八分之一！那么，如何将八分之一再细分，切成微小碎片分给十二位圣祖，或分给与雅各伯一同下埃及的七十五个人呢？而我为什么这样说呢？其实我应该从第一个人开始，反驳这种观念的愚妄。因为如果只有不朽坏者才在生子时不分散本质，而亚当是可朽坏的（经上对他说：\Quote{你原是灰土，将来还要归于灰土。\BibleRef{创 3:19}}），那么，按照欧诺弥的推理，他当然分散了自己的本质，被切割在那些从他而生的人中，并且由于他后裔的庞大数目（必须在每个人都发现的他的本质碎片根据后裔数目再次细分），亚巴郎还没有开始存在时，亚当的本质就已经用尽了，分散在这些微小到无穷的粒子中，在他的无数后代中；而通过分割过程到达亚巴郎及其后裔的亚当的微小碎片，在他们中已不能再发现为本质的残余，因为他的本性已经被先前无数万众中分割成无穷小碎片用尽了。请看那\Quote{不明白自己所讲的，也不知道自己所确定的是什么}的人是何等愚妄！因为他说：\Quote{既然祂是不朽坏的}，祂就不分散本质，也不生自己，也不成为自己的父亲，他暗中认定：凡是他声称只有不朽坏者才免于的那些事，我们必须设想它们发生在每一个受朽坏支配者身上。虽然还有许多其他考虑足以证明他论点的空虚，但我认为以上所说的已足够揭示其荒谬。然而，所有留意逻辑一致性的人必定已经承认：当他独把不朽坏归于圣父时，他就把在圣父之后被思考的一切事物都归入可朽坏的一类，通过与不朽坏的对立，以至于连圣子也不能免于朽坏。那么，如果他把圣子置于与不朽坏对立的位置，他不仅定义圣子为可朽坏的，而且断言祂具有所有那些他声称为不朽坏者所免于的属性。因为必然的结果是：如果圣父既不生自己，也不由自己而生，那么一切不是不朽坏的事物就生自己，也由自己而生，成为自己的父亲和儿子，从自己的本质转变到这些关系的每一个。因为如果不朽坏只属于圣父，而不具有那些属性是不朽坏的特殊性质，那么，根据这个异端论点，圣子当然不是不朽坏的，所以所有这些情况当然都发生在祂身上——本质被分散，生自己和由自己而生，成为自己的父亲和儿子。\par

然而，沉溺于这种愚行已久或许是浪费时间。让我们转向他论述的下一点。他在已说的话中补充道：\Quote{祂在创造行为中不需求物质、部分或自然工具，因为祂一无所缺。}这一命题，虽然\Person{欧诺弥}表述时措辞有些松散，但我们并不因其与虔敬教义不符而予以拒绝。因为我们既已得知\Quote{祂一发言，万物即成；祂一命令，万物受造}，便知道圣言是物质的造物主，且在创造物质的同时也产生了物质的属性，因此对祂而言，全能意志的冲动就是一切，并取代一切——物质、工具、地方、时间、本质、属性，以及受造之物中所想象的一切。因为祂同时愿意那应当存在者存在，且祂那产生万有的大能与祂的意志齐步并进，将祂的意志化为行动。因为伟大的\Person{梅瑟}在创造记录中就是这样教导我们关于天主的权能，将创造中显现的每个对象的产生归于命令它们存在的言语。因为他告诉我们：\BibleQuote{天主说：\InnerQuote{要有光，就有了光}}\BibleRef{创 1:3}：其余的事也是如此，丝毫没有提及物质或任何工具媒介。因此，\Person{欧诺弥}在这点的说法不应被拒绝。因为天主在创造一切源于受造之物时，既不需求任何供其运作的物质，也不需求任何帮助其建造的工具：因为天主的德能和智慧不需要任何外在协助。然而基督是\BibleQuote{天主的德能和天主的智慧}\BibleRef{格前 1:24}，万物藉祂而造，离了祂一物不存在，正如\Person{若望}所证。如果万物都是由祂造的，无论是可见的还是不可见的，且祂的意志独自足以使存在物得以存立（因为祂的意志就是德能），则\Person{欧诺弥}虽表述松散，却直说了我们的教义。因为祂用祂权能的话托住万有，在藉祂全能的话维持万物结构时，还需要什么工具和物质呢？但如果他坚持认为我们所相信的关于独生子在创造中的事，同样适用于圣子——即圣父如同圣子创造了受造物一样创造了祂——那么我们就收回先前的说法，因为这种假设是对独生子神性的否认。因为我们从\Person{保禄}的宏论中得知，崇拜和事奉受造物超过造物主乃是拜偶像的突出特征；我们也从\Person{达味}那里得知，祂说：\Quote{在你中间不可有新的天主，也不可崇拜任何异邦的天主。}我们以此准绳和规则来辨识朝拜的对象，从而确信：唯有那既不\Quote{新}也不\Quote{异邦}的才是天主。既然我们受教相信独生子就是天主，我们就因相信祂是天主而承认祂既不\Quote{新}也不\Quote{异邦}。那么，如果祂是天主，祂就不是\Quote{新}的；如果祂不是新的，祂就必定是永恒的。因此，永恒者既不是\Quote{新}的，那出自圣父并在圣父怀里且圣父在祂内的，也并非与真神性\Quote{异邦}。因此，凡将圣子与圣父的本性分离的人，要么全然禁止对圣子的朝拜，以免朝拜异邦的天主；要么跪拜在一偶像前，使受造物而非天主成为他朝拜的对象，并将基督之名赋予他的偶像。\par

如今，要明白他在关于独生子的观念中所倾向的这层含义，通过考察他论及独生子本人时所使用的语言会变得更加清楚。他的语言如下：\Quote{我们也信天主子、独生子天主、一切受造物的\OriginalTerm{首生者}{first-born}、真实之子、非\OriginalTerm{不受生}{ungenerate}者、实是在万世之前\OriginalTerm{受生}{begotten}、被称为子并非在存在之前不受生、在一切受造物之前生成、非不受造。}我认为，仅仅阅读他对自己信仰的阐述，就足以使其不敬昭然若揭，无需我们再进行任何调查。因为尽管他称祂为\Quote{首生者，}但为了不让读者心中产生祂不是受造物的任何疑虑，他立即补充说\Quote{非不受造}，以免读者领会了\Quote{子}这一术语的自然含义后，脑中产生任何关于祂的虔诚概念。正是出于这个原因，他在起初宣认祂是天主子、\OriginalTerm{独生子}{Only-begotten}天主之后，立即通过自己的补充，将读者的思想从他们的虔诚信仰扭曲到他的异端观念中。因为凡是听到\Quote{天主子}和\Quote{独生子天主}称号的人，必然会提升到关于子更高层次的断言，被这些术语的含义所引导，因为使用\Quote{天主}这个称号和\Quote{子。}这个术语的含义并没有引入本质上的差异。因为那真正是圣父的儿子且本身就是天主的那一位，怎么可能被设想为与圣父的本性不同的某种东西呢？但是，为了使这些神圣概念不会通过这些名号预先印在读者心中，他立即称祂为\Quote{一切受造物的首生者、被称为子，并非在存在之前不受生，在一切受造物之前生成、非不受造。}那么，让我们在他的论证上稍作停留，以便表明这个恶徒只是用他最初的说法作为诱饵，引诱人们接受他用虔诚倾向的短语，就像用蜜糖包裹的毒药。谁不知道\Quote{独生子}和\Quote{首生者？}这两个术语在含义上的巨大差别？因为\Quote{首生者}意味着有兄弟，而\Quote{独生子}意味着没有其他兄弟。因此，\Quote{首生者}不是\Quote{独生子，}因为\Quote{首生者}当然是兄弟中的首生者，而\Quote{独生子}没有兄弟：因为他若与兄弟并列，他就不是独生子了。此外，无论首生者的兄弟们的本性是什么，首生者本人的本性也是如此。这个称号所表示的还不止这些，还包括首生者和在他之后出生的人都是从同一个源头获得存在，而首生者对后来者的出生毫无贡献。因此，这就支持了\Person{若望}所说\BibleQuote{万物是藉着祂造的}这一陈述的虚假性。因为如果祂是首生者，祂与后来者的区别仅在于时间上的优先，而必须另有某一位，将存在的权能同样赋予祂和其余的一切。但是，为了避免我们的反对给任何不公正的对手以借口，暗示我们不接受圣经的默感之言，我们将首先向读者表明我们自己对这两个称号的看法，然后由他们判断哪一个更好。\par

% structure: p0030 source=h2 target=subsection reason=relative-heading-rank; base=h2

\subsection{八、他还非常恰当地解释了\Quote{独生子}和宗徒四次使用的\Quote{首生者}一词的含义。}

伟大的保禄知道那在一切事上居首位的独生子天主，是一总美善的创始者和原因，便为祂作证说：不仅万有的创造是藉祂而成，而且当人的原初创造已经腐朽消逝（用他自己的话来说），并在基督内成就了另一个新的创造时，在此新创造中也是祂率先引领；祂本身就是那藉福音而成就的全新人类创造的首生者。\newline{}为使我们对这点的看法更加清晰，让我们如此划分我们的论述。\newline{}这位受默感的宗徒有四次使用了这一称呼：一次如这里，称祂为\BibleQuote{一切受造物的首生者},\BibleRef{哥 1:15}；另一次称祂为\BibleQuote{在众多弟兄中为首生者},\BibleRef{罗 8:29}；再次称祂为\Quote{从死者中首生者}；还有一次他绝对地使用这一称呼，没有与其他词连用，说道：\BibleQuote{再者，当祂引领首生者进入世界时，祂说：天主的众天使都要朝拜祂}\BibleRef{希 1:6}。\newline{}因此，无论我们对其他组合中的这一称号持何种观点，我们都会一致地将其应用于\Quote{一切受造物的首生者}这一短语。\newline{}因为既然称号是同一个，它所传达的意义也必然是同一个。\newline{}那么，祂在何种意义上成为\Quote{在众多弟兄中为首生者？}祂在何种意义上成为\Quote{从死者中首生者？}\newline{}诚然，这是显而易见的：因为我们按出生是血肉之躯，正如圣经所说：\Quote{祂为我们而诞生在我们中间，有分于血肉之体},祂意在藉着由上而生的重生，即藉着水和圣神而生的重生，将我们从朽坏变为不朽；祂自己在这重生中率先引领，藉着祂自己的圣洗，将圣神引到水上；如此，祂在一切事上成了那些在灵性上重生者的首生者，并将弟兄的称号赐给那些藉着水和圣神而分享与祂相似之重生的人。\newline{}但由于祂也应当在我们本性中植入从死者中复活的能力，祂便成了\BibleQuote{长眠者初果}\BibleRef{格前 15:20}和\BibleQuote{从死者中首生者}\BibleRef{哥 1:18}，因为祂首先以自己的行动解除了死亡的痛苦，以致祂从死者中复生也为我们开辟了一条道路，因为我们被死亡痛苦所束缚，因主的复活而获得解除。\newline{}因此，正如祂因分享了重生的圣洗而成为\Quote{在众多弟兄中为首生者}，又因使自己成为复活的初果而获得\Quote{从死者中首生者}的称号，同样，祂在一切事上居首位，在\Quote{旧事都已过去}（如宗徒所说）之后，祂藉着双重重生——即藉着圣洗圣事和从死者复活所带来的重生——成了在基督内新人类的首生者，在这两重意义上，祂都为我们成了生命之君、初果和首生者。\newline{}这位首生者，因此也有弟兄；关于他们，祂对玛利亚说：\BibleQuote{去告诉我的弟兄们，我升到我的父和你们的父那里，到我的天主和你们的天主那里。}这些话总结了祂作为人所行救恩计划的全部目的。因为人类背叛了天主，\Quote{侍奉了那些本来不是神的},虽然原是神的儿女，却依附于一个虚假称为恶父。因此，天主与人之间的中保，既已取了整个人性的初果，便向祂的弟兄传达关于祂自己的宣告，不是在祂的神性身份中，而是在祂与我们共享的人性中，说道：\Quote{我离去，是要藉我自己使你们与之分离的那位真父成为你们的父，并藉我自己使你们背叛的那位真天主成为你们的天主，因为我所取的那初果，我是在我自己内将全人类呈献给它的天主和父。}\par

既然如此，\OriginalTerm{初果}{first-fruits}使真天主成为其天主，良善的圣父成为其父，这祝福便确保了整个人性，藉着初果，真天主和圣父成为所有人的父和天主。现在，\BibleQuote{如果初果是圣的，面团也是圣的。}但哪里有了初果——基督（而且初果不外乎是基督），那里也就有属于基督的人，正如宗徒所说。因此，在他提及\Quote{\OriginalTerm{首生者}{first-born}}并与其他词语连用之处，他提示我们应按我所指出的方式理解这短语；但在没有这类补充的地方，他说：\Quote{再者，当祂引领首生者进入世界时\BibleRef{希 1:6}，}其中\Quote{再者}一词宣告了万有之主在末日将要显示的那次显现。因为正如\Quote{因耶稣之名，一切在天上的、地上的和地下的，无不屈膝\BibleRef{斐 2:10}，}虽然那人性的名称并不属于那超乎万名之上的子，同样，祂说那为我们而被称为首生者的，在祂再度进入世界时，将受到所有超世受造物的朝拜，那时祂\Quote{要按公义审判世界，按正直审判万民。}因此，\Quote{首生者}和\Quote{\OriginalTerm{独生子}{Only begotten}}这两个称号的不同含义，藉着虔敬之言得以区分，各自的意义得以保持。但是，若有人将\Quote{首生者}的名称关联于子在时间之前的存在，他又怎能保留\Quote{独生子}这一术语的适当含义？让有辨识力的读者思考这些是否彼此一致，因为\Quote{首生者}必然暗示弟兄，而\Quote{独生子}必然排除弟兄的观念。因为当圣经说：\BibleQuote{在起初已有\OriginalTerm{圣言}{Word}}，我们理解为是指独生子；当它补充说\BibleQuote{圣言成了血肉，}我们心中便接受了首生者的概念；如此，虔敬之言保持清晰，为每个名称保留其自然意义，以致在\Quote{独生子}中我们看到时间之前的存在，而通过\Quote{\OriginalTerm{万物的首生者}{the first-born of creation}}我们认识到时间之前者在血肉中的显现。\par

% structure: p0033 source=h2 target=subsection reason=relative-heading-rank; base=h2

\subsection{九、额我略再次讨论独生子的产生，以及其他不同的产生方式，物质的与非物质的方式，并高贵地证明圣子是神圣光荣的闪耀，而非受造物。}

现在，让我们再次回到\Person{Eunomius}的精确陈述。\Quote{我们也信天主子，\OriginalTerm{独生子天主}{only begotten God}，\OriginalTerm{万物的首生者}{first-born of all creation}，\OriginalTerm{真子}{very Son}，不是\OriginalTerm{非受生者}{Ungenerate}，\OriginalTerm{在万世之前确实受生}{verily begotten before the worlds}。}那么，他将\Emph{产生}的意义转化为\Emph{受造}，这一点从他明确地称祂为受造的可以清楚看出，当他称祂为\Quote{成为存在}和\Quote{不是\OriginalTerm{非受造}{uncreate}}时。但是，为了让教义中表现出轻率的鲁莽和缺乏训练得以显明，让我们省略对他明显亵渎的愤慨之词，并在讨论此事时采用科学的划分。因为我认为，进行较为仔细的考察以了解\Quote{\OriginalTerm{产生}{generation}}这一术语的确切含义是好的。这个表达传达了因某种原因而存在的含义，这对所有人都是显而易见的，我想无需为此争论：但由于作为原因结果的存在方式有多种，这种差异我认为应该通过科学划分在我们的讨论中得到彻底的解释。在作为某种原因结果而存在的事物中，我们认识到以下区别。有些是物质和技艺的结果，如房屋的结构和所有通过各自材料产生的其他作品，其中某种技艺给予指导并将其目的导向其适当目标。另一些是物质和自然的结果；因为自然安排动物彼此产生，通过父母身体中的物质存有实现其自身工作；还有一些是通过物质流出。在这些中，原初物保持不变，而流出的东西被单独观察，如太阳及其光芒，或灯及其光辉，或香料和香膏，以及它们散发出的气味。因为这些，自身保持不变，每种都有它们自然产生的特殊而独特的效果，如太阳的光线，灯的明亮，以及香气在空气中产生的芬芳。除了这些，还有另一种产生方式，其原因是非物质和非形体的，但产生是可感知的，并通过身体作为工具发生；我指的是由心灵产生的言语。因为心灵本身是非形体的，通过可感知的工具产生言语。这些就是我们从哲学角度审视的产生一词的众多差异，而这本身就是，可以这样说，产生的结果。\par

现在，我们既已如此区分了各种生成模式，便当注意，圣神仁慈的经纶在将神圣奥迹传授给我们时，如何以我们所能接受的方式传授那超越理性的教导。因为默感的教导为了陈明天主不可言喻的德能，采用了人类理智所认识的一切生成形式，却又不包含附着于这些词语上的物质性的感官理解。当它论及创造之能时，将此能力称为生成，因为其表达必须俯就我们卑下的能力；然而，它并不借此传达我们在创造性的生成中所包含的一切——如时间、空间、质料的预备、工具的适宜、受造物中的设计——而是撇开这些，以崇高而雄辩的语言断言天主创造了万有，它说：\Quote{他一发言，万有造成；他一命令，万物生成。}再者，当它向我们解释独生子从圣父而来的不可言喻和超越的存在时，因人类理智的贫乏无法接受那超越一切言语和思想的教义，于是它也借用我们的语言，称他为\Quote{子}——这一名称，在我们的用法中，是赋予那些由物质和本性所生者的。但是，正如圣经在论及由创造而来的生成时，在天主的情形中并不暗示这种生成是以任何质料为媒介，而是肯定天主旨意的能力代替了物质实体、空间、时间及一切此类情况，同样地，这里在使用\Quote{子}一词时，它也摒弃了人类本性在现世生成中所注意到的一切其他东西——我是指情感、倾向、时间的协作和空间的必要性——尤其摒弃了物质，没有这些，现世自然生成便不会发生。但是，当一切这样的物质、时间和空间的存在都被排除在\Quote{子}一词的意义之外时，唯独\OriginalTerm{本性相通}{\GreekText{κοινωνία} \GreekText{φύσεως}}被保留下来，因此，藉着\Quote{子}这一称号，关于独生子，宣告了那显示他出自圣父的密切亲缘和关系的确凿性。既然这类生成不足以在我们心中植入对独生子不可言喻的存在方式的充分概念，圣经便也采用第三种生成方式来表明关于圣子天主性的教义——即由\OriginalTerm{物质流溢}{\GreekText{ῥεῦσις} \GreekText{ἐξ} \GreekText{ὕλης}}所产生的那一类——并称他为\BibleQuote{光荣的辉映}\BibleRef{希 1:3}、\Quote{香膏之气}、\BibleQuote{天主的嘘气}\BibleRef{智 7:25}；这些比喻，在我们所采用的科学用语中，通常称作物质流溢。\par

但是，正如在这些被引用的例子中，无论受造的诞生还是\Quote{子}一词的力量，都不容许时间、物质、空间或情感，同样地，在此，圣经也只采用\OriginalTerm{光照}{\GreekText{ἔλλαμψις}}及其他我已提过的比喻，脱离一切物质观念，针对这类生成方式的神圣适宜性，表明我们必须藉着这一表达的意义，理解一种既出自圣父又与圣父同有的存在。因为嘘气的形象并非旨在向我们传达由构成它的物质消散于空气中的概念，香气的形象也并非旨在表达香膏的特质消散于空气中，光照的形象也并非旨在表达光线从太阳的物体发出的流溢；而是如已说的，在一切情形中，藉着这类生成方式，只指示这一点：子出自圣父，并且与他一同被认识，在圣父与出自圣父的那一位之间没有间隔。因为圣神恩宠的极深仁慈如此安排，使关于独生子的神圣概念从多方面传到我们这里，从而被植入我们心中，于是它又增加了剩下的生成方式——即\OriginalTerm{出自理智的言语}{\GreekText{λόγος} \GreekText{ἐκ} \GreekText{νοῦ}}的那种。在此，崇高的若望用了非凡的先见之明。为使读者不致因疏忽和不配的概念，沉沦到对\Quote{言语}的普通理解，以致认为子不过是圣父的一个声音，他便断言圣言（\Person{圣言}）在本体上自立存在于那第一而蒙福的本性自身之中，高声宣告：\BibleQuote{在起初已有圣言，圣言与天主同在，圣言就是天主，且是光，是生命}，以及凡起初所是的，圣言也是。\par

由于这些受生的类型——我指的是那些因某种原因而产生并在我们日常经验中被认识的类型——也被圣经用来教导关于超越性的奥秘，以合理的方式转移来表达神圣的观念，我们现在也可以继续检查欧诺弥的陈述，看看他如何理解\OriginalTerm{受生}{generation}的含义。他说：\Quote{真子，不是非受生，在万世之前确实受生。}我认为，可以迅速掠过他在区分中违背逻辑顺序的暴力，因为所有人都容易识别。因为谁不知道，虽然正确的对立是圣父与圣子、受生与非受生，但他却忽略了\Quote{圣父}一词，将\Quote{非受生}与\Quote{圣子}对立；而如果他关心真理，就应该避免将他的措辞偏离关系的适当顺序，并说\Quote{真子，而非圣父}？这样，对虔敬和逻辑一致性都会给予应有的关注，因为本性不会在区分位格时被撕裂。但他在信仰陈述中交换了真实的、符合圣经的\Quote{圣父}一词的使用——这由圣言亲自交付给我们——并用\Quote{非受生}代替\Quote{圣父}，目的是通过将祂与圣子那种在\Quote{圣父}称号中自然被想到的亲密关系分开，将祂放在与所有受造物同等的水平上，这些受造物同样与\Quote{非受生}相对立。他说：\Quote{在万世之前确实受生。}让他说祂是由谁受生的。他当然会回答：\Quote{由圣父}，除非他准备无耻地反驳真理。但由于不可能将圣子的永恒与永恒的圣父分开，因为\Quote{圣父}一词本身的意义就隐含了圣子，因此他拒绝了圣父的称号，将措辞转向\Quote{非受生}，因为后者的意义与圣子没有任何关系或联系；通过用一词代替另一词误导读者，使他们不再将圣子与圣父一同默观，他为自己的诡辩开辟了道路，通过偷偷引入\Quote{非受生}一词为不虔敬铺路。因为那些按照主的命令相信圣父的人，当他们听到圣父的名字时，同时也在思想中接纳圣子，因为心灵从圣子转向圣父，不会踏在介于两者之间的空虚实无上。但那些被转向\Quote{非受生}称号而不是圣父的人，只获得这个名称的赤裸裸的概念，只知道祂从未被生成，而不是祂是圣父。尽管如此，即使在这种概念模式下，那些有辨别力阅读的人的信仰仍然保持清晰。因为\Quote{从未被生成}这个表述对所有非受造的本性都有相同的意义：圣父、圣子和圣神同样是非受造的。因为跟随圣言的人一直相信，所有受造物，可感的和超世界的，都从圣父、圣子和圣神获得存在。听到\BibleQuote{诸天藉主之圣言造成，万象藉祂口中的气息造成}的人，既不会把\Quote{圣言}理解为单纯的话语，也不会把\Quote{气息}理解为单纯的呼气，而是通过所说的话形成对天主圣言和天主圣神的概念。现在，创造和被创造并不等同，一切存在物分为创造者和被造者，各自在本性上互不相同，因此被造的不是非受造的，产生被造之物的也不是受造的。因此，那些按照主亲自赐予我们的信仰阐释、相信圣父、圣子和圣神之名的人，承认这三个位格都是同样无始的，而\Quote{非受生}所传达的意义不会损害他们的纯正信仰；但对于那些迟钝和模糊的人来说，这个词成为偏离纯正教义的起点。因为不理解这个词的真正含义，即\Quote{非受生}仅仅表示\Quote{从未被生成}，而\Quote{从未被生成}是所有超越受造本性的共同属性，他们放弃了对于圣父的信仰，用\Quote{非受生}的短语代替\Quote{圣父}；并且，正如所说，独生子的位格存在并不隐含在这个名称中，他们断定圣子的存在始于时间中的某个确定起点，确认（欧诺弥在此处附加到之前陈述中的）祂之所以被称为圣子，并非没有先于祂存在的受生。\par

这是什么虚妄的文字游戏？他不知道他所谈论的是天主吗？祂在起初就存在，且存在于父内，从来没有一个时候祂不存在。他不知道自己在说什么，也不知道自己肯定的是什么，却如同在编造一个普通人的家谱一样，试图将本属于我们下界本性的语言应用于万有的主。例如，\Person{依市玛耳}在产生他的出生之前并不存在，在他的出生之前当然有一段时间间隔。但论到那\BibleQuote{光荣的反映}，以前和以后没有位置：因为在反映之前，当然也没有光荣，因为与光荣的存在同时，它的反映必然发出；按事物的本性，二者不能被分开，也不可能在反映之前单独看见光荣。因为这样说的人会认为光荣本身是暗淡模糊的，如果它的反映没有同时照耀出来。但异端的不公平方法在于，试图通过运用于独生子天主的观念和术语，将祂从与父的一体性中移开。为此他们说：\Quote{在产生祂的出生以前，祂不是子}。但是先知所说的\BibleQuote{公绵羊的儿子}——他们不也是在出生之后才被称为儿子的吗？那么，理智在\BibleQuote{公绵羊的儿子}身上注意到的特性，即他们在产生他们的出生之前不是\BibleQuote{公绵羊的儿子}——现在，我们可敬的神学家将其归于万有和所有受造物的创造者，祂拥有永恒之父在自己内，并在父的永恒中被默观，正如祂自己所说：\BibleQuote{我在父内，父在我内}。然而，那些不能察觉隐藏在他陈述中的诡辩，也没有受过任何逻辑训练的人，跟随这些前后不连贯的陈述，并把接下来的一切当作前面陈述的逻辑结果而接受。因为他说：\Quote{在一切受造物之前产生}，好像这还不足以证明他的不虔敬，他在接下来的短语中还保留了一句亵渎的话，称子为受造的。那么，他称那受造的为真子是什么意思？如果称那受造的为真子是合适的，那么天当然也是真子，因为它也是受造的。同样，太阳也是真子，受造物中所有的一切，无论大小，当然都有权称为真子。他称那已产生者为独生子又是什么意思？因为一切已产生的事物彼此之间无疑是弟兄关系，我的意思是，就它们的产生而言。祂是由谁产生的？因为毫无疑问，一切曾产生的事物都是由子产生的。因为\Person{若望}这样作证说：\BibleQuote{万物是借着他造的}。那么，如果按照\Person{欧诺米}的信条，子也是产生的，祂当然被列入已产生的事物类别中。那么，如果一切已产生的事物都是由祂造的，而圣言是已产生的事物之一，谁那么愚钝，不能从这些前提推出那个荒谬的结论：我们这位新的信条贩卖者竟然明说万有的主和创造者是受造的？让他告诉我们他的这种大胆断言来自何处。来自哪句受默感的话语？哪位福音作者，哪位宗徒曾说过这样的话？哪位先知，哪位立法者，哪位族长，或是其他任何蒙圣神感动、其话语被记载下来的人，曾提出过这样的说法？由真理所传授的信仰传统教导我们相信父、子和圣神。如果相信子是受造的是正确的，那么真理在向我们传授这奥迹时，为何命令我们相信子，而不是相信受造物？受默感的宗徒自己朝拜基督，为什么规定崇拜受造物而不崇拜造物主的人犯有偶像崇拜之罪？因为，如果子是受造的，他要么就不会朝拜祂，要么就不会把那些崇拜受造物的人与偶像崇拜者归为一类，以免自己朝拜受造物而显得像个偶像崇拜者。但他知道他所朝拜的是\BibleQuote{在万有之上的天主}\BibleRef{罗 9:5}，因为他在致罗马人的书信中这样称呼子。那么，那些把子与父的本质分开、并称祂为受造物的人，为何要嘲弄地给祂一个虚假的天主称号，白白地把天主之名加给一个与真天主无关者，就像把此名加给贝尔、达贡或巨龙一样？因此，那些断言祂是受造的人，要承认祂根本不是天主，这样他们就被看出不过是伪装的犹太人；或者，如果他们承认一个受造的是天主，就让他们不要否认自己是偶像崇拜者。\par

% structure: p0039 source=h2 target=subsection reason=relative-heading-rank; base=h2

\subsection{十、祂解释\BibleQuote{上主创造了我}这话，以及关于子起源的论证、\Person{欧诺米}推理的欺骗性，以及那段经文\BibleQuote{我必不将我的光荣归给另一个}，从不同角度加以考查。}

当然，他们引用了 \WorkTitle{箴言篇} 中的那句话：\BibleQuote{主在祂道路的开端创造了我，为祂的工程。}如今，若要完全解释这段经文的真正含义，需要冗长的讨论；然而，即使寥寥数语，也足以向善心的读者传达其用意。有些精于神学的人确实指出，希伯来文本中并不是\Quote{创造}，而我们自己在更古老的抄本中读到的是\Quote{占有}，而非\Quote{创造}。毫无疑问，\WorkTitle{箴言篇} 中寓意性语言的\Quote{占有}指的是那为我们的缘故而取了奴隶的形状\BibleRef{斐 2:7}的奴隶。但如果有人主张这段经文中使用的是教会中通行的读法\Quote{创造}，我们也不拒绝。因为这同样是以寓意性语言暗示\Quote{奴隶}，因为正如宗徒告诉我们：\BibleQuote{所有受造之物都处于束缚之中}\BibleRef{罗 8:20}。因此我们说，这个说法与另一个说法一样，都可以作正统的解释。因为那为我们的缘故而成为像我们一样的祂，在末世真正地成了\Emph{受造的}——祂在起初就是圣言和天主，后来成了肉躯和人。因为肉躯的本性是受造的；祂因在各方面与我们一样地分享肉躯，却没有罪，所以当祂成为人时，祂被创造了；而且祂是\BibleQuote{按照天主}\BibleRef{弗 4:24}被创造的，而非按照人，正如宗徒所说，以一种新的方式，而非依照人的习惯。因为我们被告知，这个\Quote{新人}是\Emph{受造的}——尽管是出于圣神和至高者的能力——\Person{保罗}，奥秘的揭示者，命令我们\Quote{穿上}，用两个短语来表达这要穿上的衣服，在一处说：\BibleQuote{穿上那按照天主所创造的新人}\BibleRef{弗 4:24}，在另一处说：\BibleQuote{穿上主耶稣基督}\BibleRef{罗 13:14}。因为正是如此，那曾说过\BibleQuote{我就是道路}的祂，对于我们这些穿上祂的人来说，成为了救恩之路的开端，以便祂使我们成为祂亲手做成的工程，将我们从罪的那恶模具中重新塑造，再一次成为祂自己的肖像。祂同时是我们在未来世界之前的根基，按\Person{保罗}的话说：\BibleQuote{没有人能奠立别的根基，除非已奠立的那一位}\BibleRef{格前 3:11}，而且确实：\BibleQuote{在诸水之泉涌出之前，在山岳奠定之前，在祂造深渊之前，在一切丘陵之前，祂就已生了我。}因为根据 \WorkTitle{箴言篇} 的用法，这些短语每一个都可以从比喻的意义上应用于圣言。伟大的\Person{达味}称正义为\BibleQuote{天主的山岳}，祂的审判为\Quote{深渊}，教会中的教师为\Quote{泉源}，说：\BibleQuote{你们要从以色列的泉源中祝福天主主}；他称纯朴为\Quote{丘陵}，如同他提到它们像羊羔跳跃时所显示的那样。因此，在这些之前，那为我们之故作为人受造的祂在我们内诞生，为的是这些事物也能在我们内被创造。但我想，可以不再讨论这些点，因为真理已经用几句话向善心的读者充分指明了；现在让我们继续看\Person{欧诺米}接下来所说的。\par

\Quote{在起初就已存在}，他说，\Quote{并非无始}。那自诩有卓越洞察力的人是如何理解天主的圣言的？他宣称那在起初的祂自己有一个开始：却没有意识到，如果那在起初的祂有一个开始，那么起初本身必然有另一个开始。凡他对起初所说的，必然也必须承认是对那在起初的祂为真的：因为那在起初的如何能与起初分离？任何人都不能想象一个\Quote{非有}先于\Quote{有}。因为无论人的思想追溯多远以领悟起初，他必然同时领会到那在起初的圣言（因为它不能与它所处的起初分离）在其中不存在任何时间点上开始或停止存在。然而，不要因我这些话，就将我们所承认的唯一起初分裂为二。因为起初确是一，在其中不可分割地辨识出那与圣父完全联合的圣言。这样思想的人绝不会给异端留下漏洞，以\OriginalTerm{未生}{ungenerate}这个新奇术语损害他的虔敬。\newline{}\newline{}但在\Person{欧诺米}接下来的主张中，他的陈述如同掺杂了大量沙子的面包。因为，他将异端见解与正统教义混合，使得原本富有营养的食物也因他掺入的砂砾而难以下咽。因他称主为\Quote{活的智慧}、\Quote{运作的真理}、自立的能力和\Quote{生命}：——这部分是有营养的。但他将这些论断注入了异端的毒药。因为当他称\Quote{生命}为\OriginalTerm{受生}{generate}时，他通过隐含与\OriginalTerm{未生}{ungenerate}生命相对照而有所保留，并没有肯定圣子就是那生命本身。接着他说：\Quote{作为天主子，使死者复生，真光，照亮那进入世界的每个人的光，良善，美好之物的赐予者。}所有这些他都当作蜂蜜献给头脑简单的人，将致命的毒药隐藏在这些甜言蜜语之下。紧接着这些陈述，他立刻引入了他的有害原则，说：\Quote{不与那生他的祂分尊荣，不与另一者分父的本质，而是透过受生成为了荣耀的，是的，光荣的主，并从圣父那里接受荣耀，不与圣父分享祂的荣耀，因为全能者的荣耀是不可传递的，正如祂曾说：\BibleQuote{我的荣耀我不给与另一个人}\BibleRef{依 42:8}}\newline{}\newline{}这些是他的致命毒药，只有那些灵魂感官受过训练的人才能发现：但这些话的致命危害却由其结论暴露出来：——从圣父接受荣耀，不与圣父分享荣耀，因为全能者的荣耀是不可传递的，正如祂曾说：\Quote{我的荣耀我不给与另一个人。}\newline{}\newline{}那\Quote{另一个人}是谁？天主曾说过不把自己的荣耀给他。先知说的是天主的仇敌，而\Person{欧诺米}却将预言指向独生天主自己！因为当先知以天主的位格说话时，说过：\Quote{我的荣耀我不给与另一个人}，接着又说：\Quote{也不将我的赞美给予雕刻的偶像。}因为当人们被诱骗，将唯独归于天主的敬拜和钦崇献给天主的仇敌，在雕刻的偶像的形像中向天主的仇敌致敬——这仇敌在人间以偶像提供的形状显现于多种形态中——那医治病人的祂，怜悯人类的沉沦，藉着先知预告了末世祂要废除偶像的慈爱，说：\Quote{当我的真理显明时，我的荣耀不再给与另一个人，我的赞美也不再归于雕刻的偶像；因为人们，当认识我的荣耀时，将不再受那些本来不是神的束缚。}\newline{}\newline{}因此，先知以主的位格所说关于那仇敌权势的一切话，这位与天主动武者（指\Person{欧诺米}）却指向主自己，那位藉先知说话的主！有记载说暴君中有谁如此迫害信仰？有谁坚持如此亵渎的话：我们相信那为拯救我们灵魂而显现于肉身的那位，不是真天主，而是天主的仇敌，藉着偶像和雕刻的像对人类施行诡计？因为\Person{欧诺米}将先知论及那仇敌的话转用于独生天主，甚至没有考虑到这些正是独生自己藉先知所说的，正如\Person{欧诺米}自己后来也承认说：\Quote{这就是那藉先知说话者。}\par

我何需更详细地追究这主题的这一部分？因为前面的言语也沾染了同样的亵渎——\Quote{从圣父接受光荣，不与圣父共享光荣，因为全能天主的光荣是不可分享的。}就我而言，即使他的话是指\Person{梅瑟}——他在法律的服役中受了光荣——即便如此，我也不容忍这样的说法，即便承认\Person{梅瑟}没有内在的光荣，因着天主赐予他的恩宠，在以色列人面前显得完全光荣。因为赋予这立法者的光荣，不外乎是天主自己的光荣；主在福音中命令众人寻求这光荣，祂责备那些高举人的光荣、而不寻求惟独来自天主的光荣的人。因为祂命令他们寻求来自惟一天主的光荣，祂就声明了他们获得所寻求之物的可能性。那么，全能者的光荣怎会不可分享呢？如果连我们的本分也是向惟一的天主祈求光荣，而且按照我们主的话：\Quote{凡求的，就必得到？}但谁说关于圣父光荣的光辉，祂是通过接受而拥有光荣，就实际上说光荣的光辉本身缺乏光荣，并且需要从另一个那里接受光荣，好使自己终于成为某种光荣的主。那么我们如何处置真理的言语呢？——一个告诉我们，祂要在圣父的光荣中显现\BibleRef{谷 8:38}，；另一个则说：\Quote{凡圣父所有的，都是我的}？听者应当听从谁呢？是听从那说\Quote{按宗徒所说，那在圣父内的\InnerQuote{万有的承继者\BibleRef{希 1:2}}在祂圣父的光荣中无分无关}的人呢？还是听从那宣告凡圣父所有的，祂自己也都有的人呢？在这些\Quote{所有事物}中，光荣当然包括在内。然而\Person{欧诺弥}却说全能者的光荣是不可分享的。这种观点，\Person{岳厄尔}并不证实，强而有力的\Person{伯多禄}也不证实，他在对犹太人的演讲中采用了先知的话。因为先知和宗徒都以天主的口吻说：\Quote{我要将我的圣神倾注在一切血肉上\BibleRef{岳 2:28}。}那一位不吝惜将自己的圣神分享给一切血肉的，又怎能不将自己的光荣赋予那在圣父怀中的独生子、那拥有圣父所有一切的呢？也许有人会说，\Person{欧诺弥}在此说了实话，尽管他并非有意如此。因为\Quote{赋予}一词严格用于那些没有内在光荣、拥有光荣是来自外在的增益而非其本性的一部分的人。但是，当在两位位格中观察到同一本性时，那位在本性上就是圣父所相信的一切的，就不需要有人将每一个属性赋予祂。这一点最好更清楚、更精确地加以解释。那位有圣父整个居住在他内的，当圣父内所考量的属性没有一样离开祂时，祂对圣父的光荣有何需要呢？\par

% structure: p0043 source=h2 target=subsection reason=relative-heading-rank; base=h2

\subsection{十一、在阐述了全能者的崇高地位、圣子的永恒、以及\Quote{成了服从的}这一短语之后，他指出\Person{欧诺弥}的愚蠢，因为\Person{欧诺弥}声称圣子并非通过服从而获得祂的子身份。}

再者，欧诺弥肯定圣子没有份的全能者的高位是什么？那么，让那些在自己眼中是智慧，在自己眼中是精明\BibleRef{依 5:21}的人发表他们尘埃般的见解吧——他们正如先知所说的，\BibleQuote{从地里说话。\BibleRef{依 29:4}} 但是，我们这些敬畏圣言、作真理的门徒，或更确切说，自认如此的人，就连这个断言也不要放过而不加审查。我们知道，在一切用来指示神性的名号中，有些是表达神圣尊威的，是绝对地使用和理解的；有些则是根据对我们及一切受造界的作为而赋予的。因为当宗徒说：\Quote{现归与那不朽、不可见、独一智慧的天主，}等等时，他用这些称号提示我们某些概念，向我们代表那超越的能力；但是当圣经中称天主为仁慈、怜悯、充满同情、真实、良善、主、医生、牧者、道路、食粮、泉源、君王、造物主、工匠、保护者，是超越万有、贯乎万有、在万有之内的一位时，这些以及类似的称号都包含了神圣慈爱在创造中的作为的宣告。因此，那些精确查考\Quote{全能者}一词含义的人，会发现它关于神圣能力所宣告的，无非是那控制受造之物的作为，并且\Quote{全能者，}这个词所表明的，是与某物有某种关系。因为正如祂不会被称为医生，除非是因病人之故，也不会被称为仁慈、怜悯等，除非是因需要恩宠和怜悯者之故，同样，若一切受造物不需要一位管理者来管理和保存它们，祂也不会被称为全能者。因此，正如祂向那些需要医治者显示自己为医生，同样祂在全需要受治理者身上是全能者；并且正如\Quote{健康的人不需要医生，}所以我们可以说，那本性中包含无误和不动摇的正义原则者，不像其他人一样，不需要一位管理者在其上。因此，当我们听到\Quote{全能者，}这个名号时，我们的概念是这样的：天主保持一切可理解之物以及一切物质性事物存在。因此祂坐在地球之圈上，因此祂握着地极在祂手中，因此祂\Quote{用掌量苍天，用手心量水}；因此祂在自身内包含一切可理解的受造物，以致万物可以在祂怀抱大能控制下继续存在。那么，让我们探询，谁是那位\Quote{在万有内工作万有。}者？谁创造了万有，且没有祂，没有现存事物存在？谁是那万有在其中被创造，并且万有在其中持续存在的那一位？我们在谁内生活、行动、存在？谁拥有圣父所拥有的一切？所说过的话，是否使我们仍旧不认识那位被称为\BibleQuote{超越万有的天主\BibleRef{罗 9:5}，}由\Person{圣保禄}如此称呼的——我们的主耶稣基督，祂亲自说，手中握着\Quote{圣父所有的一切，}确实在包含万有的掌心中掌握一切，并对所掌握的有主权，没有人能从那手中拿走什么，因为祂手中握着万有？那么，如果祂拥有万有，并对所拥有的有主权，为什么这位对万有如此主权的祂，却是别的什么，而不是全能者？如果异端回答说，圣父对圣子和圣神都有主权，让他们先证明圣子和圣神具有可变的本性，然后在这可变性之上设立管理者，以便藉由从上植入的帮助，那被如此管理的能继续不能转为恶。另一方面，如果神圣本性不能为恶，不改变，不变化，永恒不变，那么它为什么需要一位管理者？它自己控制一切受造物，而它本身因不变性不需要控制它的管理者。正是为此，在基督的名前，\Quote{一切在天上的、地上的和地底下的膝盖都屈膝。}因为确实，不是所有的膝盖都会这样屈膝，除非他们认出在基督内是那为他们的救恩管理他们的那一位。但是说圣子因圣父的良善而成为存在，这无异于将祂与最卑微的受造物并列。因为有什么事物不是因造它者的良善而诞生？人的形成归因于什么？归因于造物主的恶，还是祂的良善？我们将动物的繁衍、植物和香草的出产归因于什么？没有一样事物不是因造它者的良善而开始的。那么，理性辨别出为万有共有的属性，欧诺弥竟仁慈地允许给永生的圣子！但是祂没有与圣父分享祂的本质或祂的地位——这些断言和其余的空话，我在处理他关于圣父的陈述时已经预先反驳了，并且表明他是随机冒险而说的，没有任何可理解的意义。因为即使在我们这些彼此出生的人中间，也没有本质的分割。表达本质的定义完整地保留在每一个，在生者和受生者中，对生者没有减少，对受生者没有增加。但是在拥有圣父所有一切的那一位的情况下，谈论地位或主权的分割，是没有意义的，除非这表明提议者的不敬。因此，纠缠于这样的讨论，并将我们的论文延长到不合理的长度，是多此一举。让我们继续下面的内容。\par

\Quote{受光荣的，}他说，\Quote{被圣父在万世之前。}真理之言已被证明，因仇敌的见证得以证实。因为这就是我们信仰的总纲：圣子从永远就存在，被圣父所光荣；因为\Quote{在万世之前}在意义上等同于\Quote{从永远}，因为先知用这句话向我们说明天主的永恒，当它论及他说\Quote{他从万世之前即存在。}如果存于万世之前是超乎一切元始，那在万世之前将光荣加于圣子的人，便由此断言了他在那光荣之前从永远就已存在：因为受光荣的当然不是非存在者，而是存在者。于是他继续为自己播下亵渎圣神的种子；并非为了光荣圣子，而是为了妄加凌辱圣神。因为他意图造出圣神是天使队伍的一部分，便插入了\Quote{被圣神永远光荣，并被每一个有理性和被造的存有永远光荣}这句话，以致在圣神与一切受造物之间不再有区别：也就是说，如果圣神光荣主的方式与先知所列举的其他一切存有相同，即\Quote{天使与掌权者，天上的诸天，天上的水，地上的一切，龙与深渊，火与冰雹，雪与雾，暴风，山岳与一切丘陵，果树与一切香柏，走兽与一切牲畜，爬虫与有羽的飞鸟。}那么，他说，既然圣神也与这些一同光荣主，他那敌对天主的舌头便使圣神自身也成了它们中的一个。\par

接下来这些支离破碎、毫无连贯的话，我想最好略过不谈，并非因为它们完全无可指摘，而是因为若将它们与其恶意的上下文分离，其语言可能被虔诚者使用。如果他在这里或那里使用了一些有利于虔诚的表述，那不过是作为诱饵抛给纯朴的灵魂，以便不敬神的钩子随之被吞下。因为在使用了一个教会成员可能使用的语言之后，他又接着说：\Quote{关于所有存在之物的创造和产生，服从；关于每一种服务，服从；祂并非因服从而获得子性或神性，而是作为圣子且作为独生子天主受生的结果，在言语上表现出服从，在行为上表现出服从。}然而，熟悉天主圣言的人中，有谁不知道大能的\Person{保禄}曾论及祂在哪个时间点（且仅一次）说祂\BibleQuote{成了服从的}\BibleRef{斐 2:8}？因为那是当祂取了奴仆的形象，为藉十字架完成救赎的奥迹时——祂曾空虚自己，因取人的样式和形态而谦卑自己，以人的卑微本性显现为人——那时，我说，祂成了服从的，甚至祂\BibleQuote{承担了我们的软弱，背负了我们的疾病}\BibleRef{玛 8:17}，以祂的服从医治了人的不服，以便因祂的伤痕我们得医治，因祂的死亡消除全人类共同的死亡——那时，祂为我们成了服从的，正如祂为了我们而成了\BibleQuote{罪}\BibleRef{格后 5:21}和\BibleQuote{诅咒}\BibleRef{迦 3:13}，因着祂为我们所做的安排，并非本性如此，而是因着祂对人的爱而成为这样。但有什么神圣的言语曾教导祂如此多的服从的清单呢？相反，每一篇受启示的圣经都证明祂独立自主的权能，说：\BibleQuote{祂一发言，万物造成；祂一命令，万物形成。}——因为显然圣咏作者是这样论及那位\BibleQuote{以自己大能的话支撑万有}\BibleRef{希 1:3}的，祂的权威凭其意志的单一冲动就构成了每一个存在和本性，以及受造界中由理性或感官所理解的一切。那么，\Person{欧诺弥}是从哪里被感动，以如此多样的方式将服从的属性归于宇宙的君王，说祂\Quote{在一切创造工作中服从，在每一种服务中服从，言语上服从，行为上服从}呢？然而，每个人都清楚，只有在行为和言语上服从另一个的人，是还没有在自己身上完全实现精确行动或无懈可击的言谈状态，而是始终注视着他的老师和向导，通过他的建议被训练成在言行上恰如其分。但认为智慧需要一个主人和老师来正确引导祂模仿的尝试，只是\Person{欧诺弥}的幻想，而且仅是他自己的幻想。关于圣父，他说圣父在言语上忠信，在行为上忠信；而对于圣子，他却不肯定祂在言语和行为上的忠信，只肯定服从而非忠信，因此他的亵渎无差别地贯穿他所有的陈述。但也许最好对穿插在上述言论之间的那种欠考虑愚行的主张保持沉默，以免一些不假思索的人讥笑它的荒谬——他们应当为灵魂的丧亡而哭泣，而非讥笑言语的愚蠢。因为这位聪明而谨慎的神学家说，祂并非因服从的结果而成为圣子！注意他的洞察力！他以何等令人信服的力量为我们定下：祂不是先服从而后成为圣子，我们不应认为祂的服从先于祂的受生！现在，如果他未曾加上这一限定条款，那么没有它，谁会有足够愚蠢和白痴到幻想祂的受生是圣父赐给祂的，作为祂在受生之前就已经表现出的应有的顺从和服从的奖赏呢？但为了不让任何人轻易地从这些评论中提取笑料，让每个人都想想，即使是言语的愚蠢也有令人哭泣之处。因为他想通过这些观察建立的是这样一种东西：祂的服从是祂本性的一部分，因此即使祂愿意，祂也不可能不服从。\par

他说，祂被如此构成，以致祂的本性只适于服从，就如器械中，按照某种形状制造的器械，必然在承受其作用的物体上产生工匠在器械构造中所植入的形式；若器械本身呈曲线，则不可能在接受印记的物体上画出直线；同样，若器械被制成画直线，也不可能以其印痕产生圆。何须我们多言来揭示这种思想是多么亵渎？异端的话本身已在大声宣告其荒谬。因为如果祂服从仅是因为祂被这样造成，那么祂甚至不与人平等，因为按照这个理论，我们的灵魂是自主和独立的，以自主权随意选择其所喜悦的，而祂却相反地在本性强制律的约束下行使，或更确切地说，经历服从，而祂的本性不许祂不顺服，即使祂愿意。因为\Quote{作为子和受生者，祂在言语和行动上都显示了顺服。}唉，这教义的兽性愚蠢！你使圣言服从言语，并假设在真正是圣言的那位之前有其他言语，而另一个元始之言在元始与在元始之内的圣言之间作中保，将决定传达给祂。这还不止一个：有好几个言语，欧诺弥把它们当作元始与圣言之间的许多链条，任意利用祂的服从。但何必在这闲谈上停留？任何人都能看到，即使在圣保禄说祂成了顺服的那个时刻（他告诉我们祂是这样成为顺服的：为我们成了肉身、仆人、诅咒和罪），即使在那时，我说，荣耀之主，祂鄙视羞辱，甘愿在肉身中受苦，却并未放弃祂的自由意志，祂说：\Quote{你们拆毁这座圣殿，三天之内我要把它重建起来}；又说：\Quote{没有人夺去我的性命，我有权舍掉，也有权再取回来}；并且当那些带着刀剑棍棒的人在祂受难前夕接近祂时，祂说\BibleQuote{我就是} \BibleRef{若 18:5}，他们就都退后；还有，当垂死的盗贼求祂记念祂时，祂说\BibleQuote{今天你就要同我在乐园里了} \BibleRef{路 23:43}，以此显明祂的普世主权。那么，如果即使在祂受难时期，祂也不曾与祂的权能分离，异端还能从哪里看出荣耀之王的权威从属呢？\par

% structure: p0048 source=h2 target=subsection reason=relative-heading-rank; base=h2

\subsection{十二、他继而以雄辩的言论论述\Quote{中保}、\Quote{相似}、\Quote{无始}和\Quote{受生}的解释，以及\Quote{全能者的作为及其工作的肖像与印记}。}

再者，他以令人厌倦的重复将多种中保归于天主，称祂为\Quote{教义中的中保，律法中的中保}。但我们从宗徒崇高的言论所学并非如此，宗徒说，祂以自己的教义废除了诫命的律法，成为天主与人之间的中保，并以此语宣告：\Quote{只有一个天主，在天主与人之间也只有一个中保，就是那人基督耶稣}；其中通过\Quote{中保}一词的区分，他向我们启示了虔诚奥秘的全部目的。目的如下：人性曾因敌人的恶意而背叛，沦为罪恶的奴役，也与真生命隔绝。此后，受造物的主将祂自己的受造物召回，成为人，同时仍然是天主，在两种不同本性的全部完整中既是天主又是人，如此人性便与天主不可分解地结合，那在基督内的人执行了中保的工作，借着为我们所取的初果，全团已潜在地与他联合。既然中保不只属于一方\BibleRef{迦 3:20}，而天主是唯一的，不在我们信仰所宣认的各位之间分裂（因为圣父、圣子、圣神中的天主性是一），主便一次而永远地成为天主与人之间的中保，亲自将人与天主性结合。但即使从中保的概念，我们也学到信经中所含的虔诚教义。因为天主与人之间的中保，可以说是进入了人性的共融，不是仅仅被视为人，而是真实地成了人：同样地，祂是真实的天主，并未如欧诺弥想要我们相信的那样，仅被授予天主性的虚名。\par

他对前面陈述所加的内容，同样缺乏意义，更确切地说，同样具有恶毒的意味。因为他称呼那不久之前他明确宣称是受造者的那一位为\Quote{子}，又称呼那被他与其余受造之物同列的那一位为\Quote{独生天主}，却断言祂与生祂的那一位相似，只是\Quote{以一种特殊的方式，在独特的意义上相似}。因此，我们必须首先区分\Quote{相似}这个术语的各种含义，即它在日常用法中有多少种意义，然后再讨论\Person{Eunomius}的主张。\newline{}\newline{}首先，一切欺骗我们感官的事物，虽然在本性上并不真正相同，却借着各自主体的一些偶性（如形状、颜色、声音，以及由味觉、嗅觉或触觉所产生的印象）制造幻觉，它们在本性上确实不同，却被误以为与它们真实的样子不同；习俗将这类关系称为\Quote{相似}。例如，当无生命的材料通过艺术（无论是雕刻、绘画还是塑造）被制成活物的仿制品时，该仿制品就被说成与原物\Quote{相似}。因为在这种情况下，动物的本质是一回事，而材料的本质是另一回事，后者仅仅凭借颜色和形状欺骗视觉。同样的相似还包括镜中原始形象的形象，它呈现运动的表象，但在本质上与其原物并不相同。同样，我们的听觉也可能经历同样的欺骗，例如，当某人用自己的声音模仿夜莺的歌声时，它说服我们的听觉，使我们似乎正在听鸟鸣。味觉也会遭遇同样的错觉，当无花果的汁液模仿蜂蜜的怡人甜味时：因为果实的汁液中确实有与蜂蜜的甜味某种相似。同样，嗅觉有时也可能因相似而受骗，当甘菊草的气味模仿芳香的苹果本身时，它欺骗了我们的感知。触觉也是如此，相似以各种方式违背真相，因为银币或铜币，若与金币大小相同、重量相近，若我们的视觉不能辨别真相，就可能被当作金币。\par

我们已用几句话一般描述了若干情况，在这些情况中，对象因为被认为与其实在面目不同，而在我们的感官中产生幻觉。当然，通过更费力的研究，我们有可能将考察扩展到所有那些在种类上彼此真正不同，却因某种偶然的相似而被认为彼此相似的事物。难道他（指欧诺弥）不断归于子的，就是这种形式的\Quote{相似}吗？不，他绝不会如此愚妄，竟在身为真理的那一位身上发现欺骗性的相似。再者，在受默感的圣经中，那说\BibleQuote{让我们按我们的肖像，照我们的模样造人\BibleRef{创 1:26}}的一位，告诉我们另一种相似；但我不认为\Person{Eunomius}会在父与子之间看出这种相似，以致把独生天主说成与人相同。我们还知道另一种相似，就是圣经在创世纪论及舍特时所说的：\BibleQuote{亚当生了一个儿子，按他自己的模样，照他的肖像\BibleRef{创 5:3}}；如果这就是\Person{Eunomius}所说的那种相似，我们并不认为他的说法应该被拒绝。因为在这种情况下，两个相似对象的本质并无不同，印记和原型暗示着本质的相通。这些（或与此类似的）就是我们对\Quote{相似}各种含义的看法。\newline{}\newline{}那么，让我们来看，\Person{Eunomius}断言子对父有\Quote{特殊相似}是出于什么意图，他说祂\Quote{以一种特殊的方式，在独特的意义上与父相似，不是如父对父，因为他们不是两个父}。他应许向我们展示子对父的\Quote{特殊相似}，却通过他的定义确立了不应将祂设想为相似的主张。因为他说\Quote{祂不像父对父那样相似}，从而认定祂不相似；当他加上\Quote{也不像\OriginalTerm{无生者}{Ungenerate}对\OriginalTerm{无生者}{Ungenerate}那样相似}时，这句话也禁止我们设想子与父的相似；最后，他又补充\Quote{也不像子对子那样相似}，引入了第三种概念，从而完全推翻了\Quote{相似}的含义。他就是如此跟从他的陈述，通过建立不相似来进行相似的论证。\newline{}\newline{}现在，让我们考察他在这些区分中展现的洞察力和坦率。在说子与父相似之后，他又加上一句作为防备：我们不应认为子与父相似，\Quote{如父对父那样}。唉，世间有哪个愚人会如此，一旦得知子与父相似，就被任何推理过程引导去思想父对父的相似呢？\Quote{也不像子对子那样相似}：在这里，区分的敏锐同样显著。当他告诉我们子与父相似时，他又附加定义说，不应把祂理解为像祂与另一个子相似那样。这些就是\Person{Eunomius}可怖教义的奥秘，他的门徒借此变得比世间其他人更有智慧，他们学到：子因与父相似，并不像子，因为子不是父；祂也不像\Quote{无生者对无生者那样相似}，因为子不是无生的。但我们所接受的奥秘，当它谈到父时，当然命令我们理解那为子之父，当它命名子时，教导我们领悟那为父之子。直到现在，我们从未感到需要这些哲学上的精巧区分，即由父和子这两个词暗示出两个父或两个子，可以说是两个无生者。\par

现\Person{欧诺弥}对\OriginalTerm{无生者}{\GreekText{αγέννητος}}过度关切之立意已多次解释，此处将再次略为揭示。因为既然\Person{圣父}这一称号并不指与\Person{圣子}在本质上有何不同，他的不虔敬若在此止步，便无支持，因\Person{圣父}与\Person{圣子}的自然含义排除了本质相异之念。但事实是，他使用\Quote{受生者}和\Quote{无生者}这些术语，由于二者之矛盾对立不容中间状态，正如\Quote{可朽的}与\Quote{不朽的}、\Quote{理性的}与\Quote{非理性的}以及所有那些意义相互排拒而对立之术语——我重复说，他借此放纵其亵渎，以便在\Quote{受生者}与\Quote{无生者}之间，想象出如同\Quote{可朽的}与\Quote{不朽的}之间那样的差异；而且正如可朽者的本性是一，不朽者的本性是另一，理性与非理性的特殊属性本质上不相容，他同样想证明无生者的本性为一，受生者的本性为另一，以表明正如非理性之本性被造于理性之下的从属地位，受生者也因其存在的必然性而处于对无生者的属从状态。为此，他将\Quote{全能者}之名归于无生者，而他应用此名并非如论证引导他所暗示的那样表示眷顾运作，而是将该词转用于专横主权，以使圣子成为附属与服从宇宙的一部分，与其余一切同作奴仆，归于那位以专横绝对的主权同样统驭万有的主。而这一点，他运用这些论证区分的用意，将从下文清楚显明。因为在他那些看似智慧而深思熟虑的表述之后——他说圣子既非如父相似于父，也非如子相似于子——然而并没有必要父总相似于父，或子总相似于子：例如，假定在\Place{厄提约丕雅}人中有一位父亲，在\Place{西徐亚}人中另有一位父亲，各有一子，厄提约丕雅人的儿子皮肤黝黑，而西徐亚人的儿子皮肤白皙、头发金黄，然而并不因各自都是父亲，西徐亚人就因厄提约丕雅人而变黑，厄提约丕雅人的身体也不因西徐亚人而变白——说这些话之后，\Person{欧诺弥}却不尊重他聪明的读者，凭自己臆想补充说：\Quote{圣子相似于圣父就如同子相似于父。}但是，虽然这种说法表明本性的亲属关系，正如受默感的圣经在\Person{舍特}与\Person{亚当}之事上所证实，我们的博士却不尊重他聪明的读者，引入了对\Quote{圣子}称号的空洞解释，定义圣子是全能者的\OriginalTerm{像}{\GreekText{εἰκών}}和\OriginalTerm{印}{\GreekText{σφραγίς}}，即其\OriginalTerm{能力}{\GreekText{ἐνέργεια}}的像和印。\Quote{因为圣子，}他说，\Quote{是全能者能力的像和印。}我请凡有耳可听的首先思考这一点：什么是\Quote{能力的印}？每一种能力都在表现它的主体中被视为一种活动，当活动完成时，它便没有独立的存在。例如，跑者的能力就是双脚的运动，当运动停止时，能力也不复存在。同样，凡一切活动都可以这样说——当从事某事者的活动停止时，能力也停止，并不独立存在，不论那人是在进行他所从事的活动，抑或停止该活动。那么，他告诉我们这能力本身是什么？它既非本质，也非像，也非\OriginalTerm{位格}{\GreekText{ὑπόστασις}}。所以他称圣子是非位格的相似物，而相似于非存在者的事物自身也毫无存在。这就是他玩弄空谈所达到的——相信虚无！因为相似于非存在者的事物本身也非存在。哦，\Person{保禄}、\Person{若望}，以及你们宗徒与福音作者队伍中的诸位，是谁用他们的毒舌武装起来攻击你们的话语？是谁以青蛙般的聒噪对抗你们天上的雷声？那么，这位雷霆之子说什么？\BibleQuote{在起初已有圣言，圣言与天主同在，圣言就是天主。}那么，那位——曾进入天上的圣殿，在乐园中得闻不可言传奥秘的——他说什么？他说：\BibleQuote{他是天主光荣的反映，是他位格的真像}\BibleRef{希 1:3}。在这些人如此讲说之后，我们的腹语者说了什么？\Quote{印，}他说，\Quote{是全能者的能力的印。}他使圣子位列第三，次于圣父，其间有那非存在的能力作为中介，或者更确切地说，被非存在任意塑造。起初原有的天主圣言，就是那\Quote{能力的印}——独生子天主，祂在存在之物的起初之永恒中被注视，祂在圣父怀中，以自己大能的话支撑万有，是诸世的创造者，万物从祂、藉祂、在祂内存在，那\BibleQuote{坐于大地之穹窿，以手掌量天，以手心测水}\BibleRef{依 40:12}，手中握有万有，居高位而垂顾卑微，或不如说曾垂顾以使祂脚下踏遍全地，被圣言的足迹所印——天主的形状就是某一\Quote{能力}的\Quote{印}。那么天主是能力，而不是位格吗？诚然，\Person{保禄}在解释这真理时说，祂是\Quote{真像}，不是祂能力的真像，而是\Quote{祂位格的}真像。祂光荣的反映是天主能力的印吗？哀叹他这种不虔敬的无知！在天主与祂自己的形像之间有什么居间的？位格究竟以谁为中介而与祂自己的精确形像相通？在光荣与其反映之间又能想象出什么中介？然而，既然有如此众多且有力的见证——那些受托传扬福音的人宣告了受造之主的伟大——这末次叛教的先驱竟用怎样的语言论及祂？他说什么？\Quote{作为像，}他说，\Quote{和全能者一切能力与德能的印。}他怎敢擅自修改伟大的\Person{保禄}的话语？\Person{保禄}说圣子是\BibleQuote{天主的德能}\BibleRef{格前 1:24}；\Person{欧诺弥}却称祂为\Quote{一种德能的印}，而非德能本身。然后，重复他的措辞，他在之前陈述之外还加上了什么？他称祂为\Quote{圣父的作为、言语和计谋的印。}圣子相似于圣父的哪些作为？他当然会说：世界及其中的一切。但福音已见证这一切都是独生子的作为。那么，圣子相似于圣父的哪些作为？祂成了哪些作为的印？哪部圣经曾称祂为\Quote{圣父作为的印}？如果有人允许\Person{欧诺弥}随己意造词，即使圣经不同意，也请他告诉我们：除了圣子所行的作为之外，圣父还有哪些作为是他所说圣子成为其印的？一切可见与不可见之物都是圣子的作为：可见的包括整个世界及其中的一切，不可见的包括超越世界的受造物。那么，在可见与不可见之物以外，还有哪些圣父的作为可以单独被思考，而他称圣子成了其\Quote{印}？也许，当他被逼入困境时，他会再次回到异端的那股臭气，说圣子是圣父的作为？那么，圣子自己既是圣父的作为（如\Person{欧诺弥}所说），又如何成为这些作为的印？或者，他说同一位既是作为，又是作为的相似？就算如此：让我们假设他说的是其他作为，他说圣父是那些作为的创造者，如果他果真要用\Quote{印}一词来理解相似的话。但是，除了那常在圣父内的圣言——就是那绝对、真实、首要意义上的圣言——\Person{欧诺弥}还知道圣父的哪些其他\Quote{言语}，竟称祂为\Quote{印}？他又可能提到哪些计谋，除了天主的智慧之外，天主的智慧相似于那些计谋而成为其\Quote{印}？请看他的陈述多么缺乏辨别和谨慎，多么混乱不清，他将这奥秘引入嘲笑，却不明白自己说什么或争辩什么。因为那在自身内完全拥有圣父，也完全自在于圣父内的祂，作为圣言、智慧、德能、真理，作为祂的精确形像和荣光，祂自己在圣父内就是一切，而非成为某些在祂之前于圣父内被辨认之物的像、印和相似。\par

然后，欧诺弥将诺厄时代洪水灭世、索多玛降下火雨、以及向埃及人所施的正义报复归于祂，仿佛他向那位手中握着世界尽头、万物——如宗徒所说——\BibleQuote{万有都靠祂存在 \BibleRef{哥 1:17},}的主做出了巨大让步；他岂不知道，对于那统摄万有、按其美善旨意引导和治理一切已往和将来之事的主，提及两三件奇事并不会增加祂的荣耀，正如省略其余的事意味着剥夺或损失。然而，即便对这一切只字不提，保禄的一句话本身就足以全部涵括——那句话是说祂\BibleQuote{在万有之上，贯通万有，且在万有之内。}\par

% structure: p0054 source=h2 target=subsection reason=relative-heading-rank; base=h2

\subsection{十三、他诠释福音中\BibleQuote{父不审判任何人}的段落，并进一步论及主所完成的带有身体和灵魂的人性接纳、亚当的过犯、死亡与死者复活。}

接着他说：\Quote{祂按照永恒天主的命令立法。}永恒的天主是谁？那在颁赐法律时服务于祂的又是谁？这一点众所皆知：天主是通过梅瑟将法律赐给接受它的人。既然欧诺弥本人承认与梅瑟交谈的是独生天主，那么，面前的论断将万有的主置于梅瑟的位置，并将永恒天主的属性仅归于圣父，以致藉着将祂与永恒者对比，得出万有的创造者——独生天主——不是永恒的吗？我们那位用功的朋友，记忆力甚佳，似乎忘记了保禄曾将所有这些术语用于自己，宣布在人间传播福音的宣告是奉天主的命令。因此，宗徒所自认的，欧诺弥竟不羞耻地归于先知和宗徒的主，以便将师父与祂自己的仆人保禄等同。然而，我何必详细反驳每一断言，来延长我的论述呢？因为欧诺弥著作中那些过于不设防的读者可能认为作者所说的符合圣经的许可，而每一个能批判性地分析各陈述的人都会发现它们全然染有异端的诡诈。因为教会人士和异端者都肯定\BibleQuote{父不审判任何人，但将一切审判交给了子，}但他们对这个断言分别赋予不同的含义。教会人士藉此理解至高的权威，而另一者则坚持从属和臣服。\par

但除了已经说过的话之外，还应当补充一点，即关于他们在论及降生成人时作为不敬之基础的那个立场，即：并非完整的人被祂拯救，而只是人的一半，我指的是肉身。他们如此恶意曲解真道的目的，是要证明主在祂的人性中所说的那些较低层次的话，应被视为出自天主性本身，这样他们就能使自己的亵渎显得更有力，如果它得到主亲自承认的支持的话。正因为如此，\Person{欧诺弥}说：\Quote{祂在最后的日子成为人，并没有取那由灵魂和肉身组成的完整的人。}然而，在查阅所有受默感且神圣的圣经之后，我并未发现任何这样的说法，即万有的创造者在祂为世人在地上服务之时，只取了没有灵魂的肉身。因此，在必要的情况下，考虑到救恩计划所针对的目标、教父们的教导以及受默感的圣经，我将努力驳斥在这方面捏造的亵渎谬论。主来\Quote{寻找并拯救迷失了的。}迷失了的不仅仅是肉身，而是完整的人，即由灵魂和肉身所组成的；事实上，更确切地说，灵魂比肉身更先迷失。因为悖逆不是肉身的罪，而是意志的罪；而意志本就属于灵魂，我们本性的整个灾难正是从灵魂开始的，正如天主那不容虚假的警告所证明的，祂宣布在那天他们吃了禁果，死亡会立即随之而来。既然对人的惩罚是双重的，死亡相应地在我们本性的每一部分造成双重生命的丧失，这双重生命运行在遭受致命打击的人身上。因为肉身的死亡在于感官知觉的灭绝，以及肉身分解为同类的元素；但\Quote{犯罪的人的灵魂，}他说，\Quote{必要死亡}\BibleRef{厄 18:20}。罪不是别的，就是远离天主，祂是真实而唯一的生命。因此，第一个人在他悖逆之后还活了数百年，但天主并没有说谎，祂说：\Quote{你吃的那一天，必定要死。}因为由于他远离了真实生命，死亡的判决在同一天就对他生效；此后，在很久以后，才发生了\Person{亚当}肉身的死亡。因此，祂来是为了寻找并拯救迷失了的（比喻中的牧人称之为羊），祂找到那迷失的，并将整只羊扛在肩上带回家，不只是羊皮，以便使天主的人完整，在灵魂和肉身中与天主性结合。因此，祂在各方面都像我们一样受到试探，却没有犯罪，我们没有一点本性是祂没有取来的。灵魂本身不是罪，尽管它可能因受不良建议而接纳罪；祂通过与自己结合而圣化灵魂，目的是使面团的初果一同成为圣的。因此，当天使告知\Person{若瑟}关于主的仇敌的毁灭时，也说：\Quote{那些谋害孩子性命的人已经死了}（或\Quote{灵魂}）；主对犹太人说：\Quote{你们想要杀我，一个将真理告诉你们的人。}这里\Quote{人}不仅指人的肉身，而是指由灵魂和肉身两者组成的整体。祂又对他们说：\Quote{因为我使一个人在安息日完全痊愈，你们就对我生气吗？}祂在别的福音中说明了\Quote{完全痊愈}的意思，当时祂对从中被放在床上的人说：\Quote{你的罪赦了}，这是灵魂的医治；以及\Quote{起来行走吧}，这关乎肉身；在\BibleBook{圣若望}{福音}中，祂在赐予肉身健康之后，也释放了灵魂脱离其自身的疾病，祂说：\Quote{你已经痊愈了，不要再犯罪}，你是在两者中都被治愈了，我指的是灵魂和肉身。\Person{圣保禄}也这样说：\Quote{为把双方在自己身上造成一个新人}\BibleRef{弗 2:15}。同样，祂预言在祂受难时，祂会自愿将自己的灵魂与肉身分开，说：\Quote{没有人夺去我的灵魂，是我自己舍弃的；我有权舍弃，也有权再取回。}是的，\Person{达味}先知根据伟大的\Person{伯多禄}的解释，也预见了祂，说：\Quote{你不会将我的灵魂遗弃在阴间，也不让你的圣者见朽坏}，而宗徒\Person{伯多禄}这样解释这话：\Quote{祂的灵魂没有被遗弃在阴间，祂的肉身也没有见朽坏。}因为祂的天主性，无论是在取肉身之前、在肉身中，还是在祂受难之后，都是永恒不变的，始终是祂本性所是，并永远如此。但在祂人性的苦难中，天主性通过暂时使灵魂与肉身分离而完成了为我们益处的救恩计划，但并未与它曾一次结合的任何一部分分离，并将已分开的部分重新结合，以便给整个人性一个从死者中复活的起始和效法的榜样，使所有朽坏的穿上不朽，所有必死的穿上不死，我们的初果已通过与天主的结合而转化为天主性，正如\Person{伯多禄}所说：\Quote{你们钉在十字架上的这位耶稣，天主已经立祂为主，为基督了。}我们可以引用许多圣经章节来支持这一立场，说明主如何通过基督的人性使世界与自己和好，将祂的仁爱工作分派给祂的灵魂和肉身，通过祂的灵魂愿意，通过祂的肉身接触他们。但详细列举每一个细节可能会使我们的论证过于冗长。\par

然而，在继续讨论下文之前，我还要提及一处经文：\Quote{你们拆毁这座圣殿，三天之内，我要把它重建起来。}正如我们藉着灵魂和身体，成为那\Quote{住在我们内，在我们中间往来 \BibleRef{格后 6:16}}的天主的圣殿，同样，主也将二者的结合称为\Quote{圣殿}，其\Quote{拆毁}表示灵魂与身体的分离。如果他们引用福音中的话：\Quote{圣言成了血肉}，试图证明肉身被纳入天主性时没有灵魂，理由是灵魂没有与肉身一同被明确提及，让他们知道，圣经通常以部分来暗示全体。因为那说过\Quote{凡有血肉的，都要来到你面前}的，并非意味着肉身可以不带着灵魂呈现在审判者面前：并且当我们读到圣史记载\Person{雅各伯}带着七十五人下到\Place{埃及}时，我们也理解这事也同时指肉身与灵魂。所以，圣言成了血肉时，带着肉身取了整个人性；因此，饥饿、干渴、恐惧、战栗、欲望、睡眠、眼泪、精神困扰，以及一切这类事情，都可能存在于祂内。因为天主性按其本性不容纳这样的情感，而如果灵魂不与身体一同受影响，肉身本身也不会陷入其中。\par

% structure: p0058 source=h2 target=subsection reason=relative-heading-rank; base=h2

\subsection{十四、他进而讨论\Person{厄诺弥乌斯}和教会关于圣神的观点；并证明圣父、圣子和圣神不是三位神，而是一位天主。他还讨论了\Quote{屈服}的不同含义，并指出万物对子的屈服与子对父的屈服是相同的。}

关于他对圣子的亵渎，就说到这里。现在让我们看看他关于圣神说了什么。他说：\Quote{在他之后，我们相信，} \Quote{护慰者、真理之神。}我认为，凡读过这段文字的人，都会明白他在关于圣子和圣父的言论中，如此歪曲主所传授给我们的信仰宣言，其用意何在。虽然这种谬误已经暴露，但我仍将尽力用几句话说清他的奸诈意图。如先前情形，他避免使用\Quote{圣父}的名号，以免因此将圣子包含在圣父的永恒之内；同样，他避免使用圣子的称号，以免因此暗示圣子与圣父的本性关系；同样在这里，他也避免说\Quote{圣神}，以免因此名承认祂荣耀的尊威，以及祂与圣父和圣子的完全合一。因为圣经同样将\Quote{神}和\Quote{圣}的名称应用于圣父和圣子（因为\BibleQuote{天主是神}，\BibleQuote{在我们面前的受傅者主是神}，\BibleQuote{我们的主天主是圣的}，并且\BibleQuote{唯一的圣者，唯一的主耶稣基督}），因此，倘若用这些名号，就会在读者心中产生对圣神的正确观念，这种观念本应由于圣神与圣父和圣子共享荣耀的称号而自然产生。为此，他欺骗愚昧者的耳朵，改变了天主在传授此奥迹时所设立的信仰之言，通过这一系列安排，从而为他对圣神的不敬打开了入口。因为如果他曾说\Quote{我们相信圣神}，以及\BibleQuote{天主是神}，那么任何受过神圣事物教导的人都会插话指出：如果我们相信圣神，而天主被称为神，那么祂在本质上与那个在真正意义上接受同样名称的实体肯定没有区别。因为对于所有那些并非虚假地、也不是比喻地，而是真实地、绝对地由同一名称所指的事物，我们必须承认，由这些名称的等同所表示的本性也是同一的。正是为此，他压制了主在信仰公式中所指定的名称，说：\Quote{我们相信护慰者。}然而我得知，受默感的圣经也将同样的名称应用于圣父、圣子和圣神。因为圣子将\Quote{护慰者}的名号同样给予自己和圣神；而圣父，在祂被说成施行安慰的地方，当然也拥有\Quote{护慰者}的名号。因为那位实行护慰之事的，当然不会轻视属于这事的名称：因为\Person{达味}对圣父说：\BibleQuote{主，祢帮助了我，安慰了我}；伟大的宗徒也将同样的话应用于圣父，说：\BibleQuote{愿我们的主耶稣基督的天主和父受赞美，祂在一切困苦中安慰了我们}\BibleRef{格后 1:3}；\Person{若望}也在他的公函中明确将护慰者的名称给予圣子。不仅如此，主自己在谈到圣神时说，将赐给我们\Emph{另一位}护慰者，这明显首先将这称号归于自己。但因着\OriginalTerm{安慰}{\GreekText{παρακαλεῖν}}一词有两种含义——一种是通过尊敬的话语和姿态\Emph{恳求}，使我们所求的对象，对于我们恳求之事产生同情；另一种是\Emph{安慰}，对身体和灵魂的疾病采取治疗性的关注——圣经以两种含义都断言护慰者的概念属于天主性。因为有一次，\Person{保禄}用\OriginalTerm{安慰}{\GreekText{παρακαλεῖν}}一词向我们呈现天主的治愈能力，如他所说：\BibleQuote{天主，安慰那些沮丧的人，藉着弟铎的到来安慰了我们}\BibleRef{格后 7:6}；另一次，他使用该词的另一含义，在写信给格林多人时说：\BibleQuote{我们现在是基督的使者，好像天主藉着我们恳求你们；我们替基督请求你们：与天主和好吧}\BibleRef{格后 5:20}。既然这些事如此，无论你如何理解用于圣神的\Quote{护慰者}称号，你在任何一种含义上都不能将祂从与圣父和圣子共享此称号的共同性中分离出来。因此，即使他愿意，他也无法通过将圣经也归给圣父和圣子的属性归于圣神，来贬低圣神的荣耀。但是，\Person{欧诺弥}称祂为\Quote{真理之神}，我以为他本人的愿望是通过这个短语暗示从属关系，因为基督是真理，而他将祂称为真理之神，好像有人说祂是真理的财产和物品，却不知道天主被称为公义之神；因此我们当然不会理解为天主是公义的财产。因此，当我们听到\BibleQuote{真理之神}时，我们通过这个短语获得了合乎天主性的概念，被随后的话语引导到更高的解释。因为当主说\BibleQuote{真理之神}时，祂立刻加了\BibleQuote{祂从父而来}，这一事实主的圣言从未对任何受造物声称过，无论可见或不可见，无论宝座、宰制者、掌权者、统治权威，还是今生和来世任何其他被称呼的名号。因此，显然那一切受造物都不能分享的，是非受有者的特殊和专有之物。但这个人命令我们相信\Quote{虔诚的引导者}。那么，让人相信\Person{保禄}、\Person{巴尔纳伯}、\Person{弟铎}、\Person{息尔瓦诺}、\Person{弟茂德}，以及所有那些引领我们进入信仰之路的人吧。因为如果我们要相信\Quote{引导我们走向虔诚的那一位}，与圣父和圣子一起，那么所有先知、立法者、圣祖、传报者、福音传道者、宗徒、牧者和教师，都与圣神有同等的尊荣，因为他们对后辈曾是\Quote{虔诚的引导者}。他接着说：\Quote{祂由唯一的天主通过独生子而受造。}这句话凝聚了他所有的亵渎。他再次称圣父为\Quote{唯一的天主}，祂使用独生子作为产生圣神的工具。他在圣经中找到了什么这种概念的影子，竟敢如此断言？他从什么前提推导出这样的结论，将他的亵渎带到如此地步？哪位福音传道者说过？哪位宗徒？哪位先知？不，恰恰相反，每部由天主默感、藉圣神的嘘气写成的圣经，都证明了圣神的天主性。例如（因为最好用实际的证据来证明我的立场），那些接受权能成为天主的子女的人，为圣神的天主性作证。谁不知道主的这句话：由圣神而生的人是天主的子女？因为祂明确将天主子女的诞生归于圣神，说：如同由肉而生的就是肉，由圣神而生的就是神。但凡是生于圣神的，就称为天主的子女。同样，当主向门徒嘘气，将圣神赐予他们时，\Person{若望}说：\BibleQuote{我们众人都从祂的满盈中领受了。}而\BibleQuote{在祂内住有整个天主性的圆满}\BibleRef{哥 2:9}，伟大的\Person{保禄}证实；不仅如此，通过先知\Person{依撒意亚}，关于恩赐给他的天主显现也被证实，当他看见那位坐在\BibleQuote{崇高而举扬的宝座}\BibleRef{依 6:1}上时；诚然，较古老的传说说是圣父显现给他，但福音传道者\Person{若望}将预言归于我们的主，提到那些不相信先知关于主所说话语的犹太人，说：\Quote{依撒意亚这样说，当他看见祂的光荣时，就讲论了祂。}但伟大的\Person{保禄}在罗马对犹太人的讲论中将同一段经文归于圣神，他说：\BibleQuote{圣神藉依撒意亚先知对你们说得很好：你们听是听见，却不明白}，依我之见，他藉着圣经本身表明：每一个特别神圣的异象，每一次神视，每一个以天主位格说出的话语，都应理解为指向圣父、圣子和圣神。因此，当\Person{达味}说：\BibleQuote{他们在旷野中激怒了天主，在沙漠中使祂忧伤}，宗徒将以色列人对天主所做的侮辱归于圣神，说：\BibleQuote{因此，圣神说：你们今天不要再心硬，像在激怒之时，像在旷野中试探之日；你们的祖先在那里试探了我}\BibleRef{希 3:7}，并继续将预言中所有归于天主的话，都归于圣神。那些因为我们持这些观点而不断对我们重复\Quote{三个天主}的人，也许还没有学会如何数数。因为如果圣父和圣子不被分为两数（因为按照主的话，他们是一，不是二），并且圣神也是一，那么一加一怎么能分成三位天主的数目呢？难道不是很清楚，没有人能指控我们相信三位天主的数目，除非他自己先在他的学说中主张一对话神的数目？因为正是通过加到二上，那一个才完成了三神的组合。但是，对于那些崇拜一位天主——即由圣父、圣子和圣神之名所表达的天主——的人，哪里还有空间指控三神论呢？\par

然而，让我们完整地再次引用优诺米的陈述。\Quote{这位圣神也是由唯一的天主通过独生子而产生——}有什么证据可以证明\Quote{这位圣神也是}是独生子所造之物之一呢？他们当然会说\BibleQuote{万物是藉着祂造的}，而在\Quote{万物}一词中包含了\Quote{这位圣神也是}。我们对此的回答是：万物是藉着祂造的，就是那些受造之物。按照保禄告诉我们的，受造之物包括有形和无形的、宝座、权威、宰制者、掌权者、异能者；而在宝座和权能的名目下，保禄也算入了革鲁宾和色辣芬。如此，\Quote{万物}一词的范畴便达到这里。但对于圣神，由于祂超越受造之物的本性，保禄在列举万有时对祂只字未提，既没有用祂的话表明圣神的从属地位，也没有表明祂的产生；然而，正如先知称圣神为\Quote{良善的}、\Quote{正直的}和\Quote{引导的}（用\Quote{引导}一词表示管治的能力），同样，宗徒也赋予圣神尊威独立的权柄，当他肯定圣神随己意运行一切在所有之中时。\BibleRef{格前 12:11}再者，主在与尼苛德摩的谈话中彰显了圣神的独立能力和行动，祂说：\BibleQuote{圣神随意向哪里吹}。那么，优诺米怎么敢妄下定论，说圣神也是通过圣子而产生的万物之一，被判定永远服从呢？因为他描述圣神是\Quote{一次而永远地服从}，我不知道他把这引导和治理的圣神束缚在何种形式的服从之下。因为\Quote{服从}这一术语在圣经中有多种含义，并被以多种不同的方式理解和运用。因为圣咏作者说，连无理性的自然也被置于服从之下，并且将同一术语用于战争中被打败的人；而宗徒吩咐仆役要服从自己的主人\BibleRef{铎 2:9}，并且那些被置于司祭职位之上的人要叫他们的子女服从\BibleRef{弟前 3:4}，因为他们的无纪律行为会给他们的父亲带来耻辱，如同司祭厄里儿子们的例子。再者，他也论及万人对天主的服从，那时我们众人因信德彼此结合，成为那在万有之中的主的同一身体，正如圣子对圣父的服从，那时万有一致对圣子的朝拜，包括天上、地上和地底下的一切，都归光荣于圣父；正如保禄在其他地方说：\BibleQuote{上天、地上和地下的一切，无不向耶稣屈膝，无不口认耶稣基督是主，以光荣天主圣父}。因为当这事发生时，保禄的深刻智慧肯定说，那在万有之中的圣子，由于那些祂所在之人的服从，而服从于圣父。优诺米对圣神所主张的\Quote{一次而永远的服从}究竟是哪一种，从他抛出的短语中无法得知——他是指无理性受造物的服从，还是俘虏的服从，还是仆役的服从，还是受管束的孩子的服从，还是那些藉服从而得救之人的服从？因为人对天主的服从，对于那些如此服从的人而言就是救恩，按先知的声音，他说自己的灵魂服从于天主，因为救恩藉着服从来自天主，因此服从是避免丧亡的途径。因此，正如病人迫切寻求医术的帮助，同样，需要救恩的人也迫切寻求服从。然而，那赋予万有生命的圣神，还需要什么生命，以致祂要藉服从为自己获得救恩呢？既然他并非凭任何天主的话语来主张圣神的这种属性，也不是因任何可能的论证结果而对圣神发出这种亵渎，那么无论如何，有理智的人必能看清，他完全没有任何根据，既无圣经权威也无逻辑推论，就对圣神发泄他的不敬。\par

% structure: p0061 source=h2 target=subsection reason=relative-heading-rank; base=h2

\subsection{十五、最后，他详细展示了优诺米的愚昧，优诺米时而说圣神是受造的，是圣子最美的工程，时而又因归给圣神的行动而承认祂是天主，如此结束了本书。}

他继续加上说：\Quote{既不与圣父同等，也不与圣父同列（因为万有的天主是独一圣父），也不与圣子平等，因为圣子是独生子，没有与祂同生的弟兄。}就我而言，如果他在之前的陈述上仅仅加上一句话，说圣神不是圣子的父，那么我甚至认为他浪费时间去思索无人怀疑的事，并禁止人们对祂形成连最愚笨的人都不会有的概念，实属无聊。但是，因为他试图用不相关、不连贯的陈述来建立他的不敬，幻想通过否认圣神是独生子的父来证明祂是隶属的、从属的，所以我提这些词，作为证明这个人的愚蠢的证据，他幻想自己正在证明圣神是从属于圣父的，理由在于圣神不是独生子的父。因为有什么强制性的结论，如果祂不是父，祂就一定是隶属的？如果已经证明\Quote{父}和\Quote{主宰}在意义上相同，那么无疑会得出这样的结论：既然绝对主权是父这个概念的一部分，我们就应当肯定圣神是隶属于那位在权威上超越祂的。但是，如果\Quote{父}仅仅暗示祂与圣子的关系，而使用这个词并不包含绝对主权或权威的概念，那么，从圣神不是圣子的父这一事实，如何能得出圣神是隶属于圣父的结论？他说：\Quote{也不与圣子平等。}他怎么这样说？因为存在、不变、不容任何邪恶、不可改变地持守于善之中——所有这些都在圣子和圣神身上显示出毫无差别。因为圣神不朽的本性同样远离朽坏，正如圣子一样；在圣神内，正如在圣子内，祂本质的善完全与它的对立面隔绝；在两者中，他们在一切善上的完美都不需要任何增添。\par

现在，受默感的圣经教导我们对圣神肯定所有这些属性，当它把\Quote{善}、\Quote{智慧}、\Quote{不朽坏}、\Quote{不朽的}以及所有这样崇高且恰当地应用于天主性的概念和称谓归诸圣神时。那么，如果祂在这些方面毫不逊色，\Person{欧诺弥}以什么方式确定圣子和圣神的不平等呢？他告诉我们：\Quote{因为圣子是独生子，没有与祂同生的弟兄。}关于我们不应把\Quote{独生子}理解为有弟兄这一点，我们在评论\Quote{一切受造物的首生者}这一短语时已经讨论过了。但是，我们不应让\Person{欧诺弥}现在不合理地附加在这一术语上的意义未经检验。因为教会的教义宣告：在圣父、圣子和圣神内，有一个能力、一个善、一个本质、一个荣耀等等，但保持位格的区别，而这个人，当他想要使独生子的本质与受造物共有时，祂在时间先存方面称祂为\Quote{一切受造物的首生者}，通过这种表达方式宣称受造物中一切可想之物都与主有弟兄关系；因为首生者肯定不是那些以其他方式出生者的首生者，而是那些与祂相似出生者的首生者。但是，当他致力于切断圣神与圣子的结合时，他称祂为\Quote{独生子，没有与祂同生的弟兄}，目的并非为了设想祂没有弟兄，而是为了通过这一断言建立关于圣神本质与圣子疏离的论点。诚然，我们从圣经中得知，不可称圣神为圣子的弟兄；但是，不可说圣神与圣子是同类同质的，这在神圣的圣经中没有任何地方显示。因为如果圣父和圣子内驻有赋予生命的能力，那么根据福音的话，这能力也归于圣神。如果人们可以同样在圣父、圣子和圣神内辨别出这些属性：不朽坏、不变、不容邪恶、善、正直、引导、随祂意愿在一切内行作一切，以及所有类似的属性，那么如何能通过在这些方面的同一性而推论出种类上的差异呢？因此，虔诚的话语一致肯定，我们不应认为任何种类的弟兄关系归于独生子；但是，说圣神与圣子不是同类同质的——正直者与正直者、善者与善者、赋予生命者与赋予生命者——这已通过逻辑推理被清楚证明为异端的欺诈行为。\par

为何圣神的威严竟被这样的论据所削减呢？因为没有甚么会导致他在适合天主性的概念上产生过度或不足的偏差，而且既然这一切都由圣经同样地归于圣子和圣神，他也不能告诉我们他在哪里看出了不平等。但他却以赤裸裸的形式对圣神发出亵渎，毫无准备，也没有任何连贯的论据支持。他说：\Quote{既不与任何其他者并列，}他又说：\Quote{因为祂在存在方式、本性、光荣和知识上，已超越了所有藉圣子而存在的受造物，作为独生子的首要且最高贵的工程，最伟大、最光荣的。}然而，我将把嘲笑其风格的粗俗和冗余的任务留给别人，因我认为，当处理我们面前的论题时，以表达粗俗来反对一个不虔敬的人，对白发苍苍的老年人来说是不合适的。我只在我的探讨中补充这一点：如果圣神已\Quote{超越}了圣子的一切受造物（因为我要用他自己不合语法且无意义的措辞，或者更确切地说，为了更清楚，我用我自己的话表达他的观念），如果祂超越圣子所做的一切事物，那么圣神就不能与其余的受造物同等；并且，如果如欧诺弥所说，祂凭着出生的先后性而超越它们，那么他必须承认，在其余受造物中，按产生顺序最先的事物比其后的事物更受重视。现在，无理性动物的受造先于人。因此，他当然会宣称无理性的本性比有理性的存在更尊贵。同样，按照欧诺弥的论点，加音将被证明比亚伯尔更优越，因为他在出生时间上先于他；星辰将被证明比所有从地生长出来的东西更低劣、更不卓越；因为后者在第三天从地生出，而所有星辰据梅瑟记载是在第四天创造的。好吧，肯定没有人会傻到因草地的先时性而认为它比天上的奇迹更值得尊敬，或将胜利奖赏给加音而不是亚伯尔，或将后来才存在的人置于无理性动物之下。因此，我们作者的主张毫无意义，即圣神的本性优于后来存在的受造物，理由是祂比它们先存在。现在让我们看看那使祂与圣子脱离共融的人愿意授予圣神的光荣：他说：\Quote{因为祂也是独一的、首要的、独一无二的，在本质和本性的尊贵上超越圣子的一切受造物，按照圣子的喜悦完成一切作为和一切教导，被祂派遣，从祂领受，向受教者宣告，并引导人进入真理。}他称圣神\Quote{完成一切作为和一切教导。}甚么作为？是指圣父和圣子所执行的，按照主自己所说\Quote{迄今工作}人的救恩，还是另有所指？因为如果祂的作为是所说的那一个，祂必定与行那作为的圣父和圣子有相同的德能和本性，而在这样一位内，与天主性种类的差别不可能存在。正如如果任何事物行使火的功能，以完全相同的方式发光和发热，它本身当然是火；同样，如果圣神行使圣父的作为，祂必定被承认与圣父有相同的本性。另一方面，如果祂运作别于我们救恩的某事，并在相反方向显示祂的运作，祂将被证明有不同的本性和本质。但欧诺弥的陈述本身作证：圣神与圣父和圣子一样使人生活。因此，从作为的同一性，确实得出圣神与圣父和圣子的本性并非疏离。至于圣神按照圣子所喜悦的完成圣父的作为和教导，我们同意。因为本性的共融确保圣父、圣子和圣神的旨意是同一个；因此，如果圣神愿意圣子所喜悦的，旨意的共融明确指向本质的统一。但他继续说：\Quote{被祂派遣，从祂领受，向受教者宣告，并引导人进入真理。}如果他没有先前关于圣神说过甚么，读者肯定会以为这些话是指某个人类教师。因为接受使命与被派遣是同一回事，一无所有，而是从派遣者的慷慨中领受，并将他的话传达给受教的人，并为迷途者引导真理。所有这些，欧诺弥慷慨地允许圣神拥有，都属于教会目前的牧者和导师——被派遣、领受、宣布、教导、提示真理。现在，如他上面所说：\Quote{祂是独一的、首要的、独一无二的，超越一切，}如果他只停留在这里，他会显得是真理教义的捍卫者。因为那不可分割地在独一者中被默观的那一位，最真实地是独一者；那在首位中的，是首位；那在独一者中的，是独一者。因为就如人内的人灵与人本身只是一个人，同样，在祂内的天主之神和天主本身，适当地被称为独一天主、首位和独一者，因为不能与祂所在的那一位分离。但实际情况是，由于他添加了他的亵渎性短语\Quote{超越圣子的一切受造物，}他产生了混乱，将那股\Quote{任意而吹}并\Quote{在一切内行一切 \BibleRef{格前 12:6}}的那一位只赋予与其余受造物相比的优越性。\par

现在让我们进一步看他添加到这\Quote{圣化圣人}的内容。若有人也这样论圣父和圣子，他必说真话。因为那圣者居住在其中的人，祂使之成圣，正如善者使人成为善。而圣父、圣子及圣神都是圣而善的，如已证明的。\Quote{为那些接近奥迹者作向导。}这很可用来论及\Person{阿颇罗}，他浇灌了\Person{保禄}所栽种的。因为宗徒借其引导而栽种，而\Person{阿颇罗}在施洗时，借圣事性的重生而浇灌，将那些受过\Person{保禄}教导的人引向奥迹。如此，他便将那借洗礼成全人的圣神置于与\Person{阿颇罗}同等的地位。\Quote{分施一切恩赐。}对此我们也同意；因为一切美善都是圣神恩赐的一部分。\Quote{与信友合作，为理解并默观所指定之事。}既然他没有说明是由谁指定，他便使其含义含混，不知是否正确还是相反。但我们将稍加补充，使他的陈述符合虔诚。因为无论是智慧之言、知识之言、信德、援助、治理，或是救恩恩赐清单中所列举的其他任何事物，\BibleQuote{这一切都是同一而自有的圣神所运行，随祂心愿分给各人}\BibleRef{格前 12:11}，因此我们不拒绝\Person{欧诺米乌斯}的陈述，他说圣神\Quote{与信友合作，为理解并默观祂所指定之事}，因为一切美好教训都由祂指定给我们。\Quote{为祈祷者发出伴奏。}认真考察这表达的含义是愚蠢的，其荒谬无意义之处对所有人都一目了然。因为谁这样癫狂失心，等待我们告诉他圣神不是一只铃铛或一个空桶，被祈祷者的声音像打击一样敲响而发出伴奏？\Quote{引导我们走向对我们有益的事。}圣父和圣子也同样这样行：因为\BibleQuote{祂引领\Person{若瑟}如羊群}，和\BibleQuote{引领祂的百姓如羊群}，以及\BibleQuote{善神在正义之地引领我们}。\Quote{坚固我们以致虔诚。}\Person{达味}说，坚固人以致虔诚是天主的作为；圣咏作者说：\BibleQuote{因为祢是我的力量和避难所}，和\BibleQuote{上主是自己百姓的力量}，以及\BibleQuote{祂必赐力量与权能给祂的百姓}。如果\Person{欧诺米乌斯}的表达符合圣咏作者的意思，它们便是圣神天主性的见证；但如果它们与先知的话对立，那么正因此，\Person{欧诺米乌斯}就犯了亵渎之罪，因他以自己的意见对立于圣先知。接着他说：\Quote{以知识之光减轻灵魂的负担。}虔诚的教义也将这恩宠同样归于圣父、圣子和圣神。因为祂被\Person{达味}称为光，从此知识的光照耀在那些被光照的人身上。同样，这段陈述所谈到的我们思想的洁净，也属于主的能力。因为祂是\BibleQuote{圣父荣耀的光辉，和祂本体的真像}，祂\BibleQuote{洗净了我们的罪}\BibleRef{希 1:3}。再者，驱逐魔鬼（\Person{欧诺米乌斯}说这是圣神的属性），这同样也是独生天主（祂曾对魔鬼说：\Quote{我命令你}）所归于圣神的能力，因为祂说：\BibleQuote{我若靠天主的圣神驱魔}\BibleRef{玛 12:28}，所以驱逐魔鬼并非损害圣神的荣耀，反而是祂神圣超越之能力的明证。\Quote{医治病人}，他说，\Quote{治愈软弱者，安慰受苦者，扶起跌倒者，挽回痛苦者。}这是那些对圣神心怀尊敬之人的话，因为没有人会将这些效果的任何一件归于除天主之外的任何一位。如果异端断言那些唯独天主才能成就的事，是由圣神的能力所行，那么连我们的敌人也为我们所辩护的真理作了见证。圣咏作者如何向天主求医治？他说：\BibleQuote{上主，求祢怜悯我，因为我软弱；上主，求祢医治我，因为我的骨头战栗！}\Person{依撒意亚}向天主说：\BibleQuote{从祢而来的露水是他们的医治。}再者，先知的话证实，使迷途者回转是天主的作为。因为圣咏作者说：\BibleQuote{他们在旷野在干燥之地迷失了路}，他又加上：\BibleQuote{祂便带领他们走直路，使他们可往可居之城}，以及\BibleQuote{上主带领\Place{熙雍}被掳的人回归。}同样，安慰受苦者也被归于天主，\Person{保禄}这样说：\BibleQuote{愿我们的主耶稣基督的天主和父受赞美，祂在我们一切的磨难中安慰我们}\BibleRef{格后 1:3}。再者，圣咏作者以天主的口吻说：\BibleQuote{你在困苦中呼求我，我便拯救你。}而扶起跌倒者，圣经无数次归功于主的能力：\BibleQuote{你重重地推我，叫我跌倒，但上主帮助了我}，\BibleQuote{他虽跌倒，也不致被弃，因为上主用手扶持他}，以及\BibleQuote{上主扶助跌倒的人。}至于挽回痛苦者，明认属于天主的慈爱，如果\Person{欧诺米乌斯}指的是我们在预言中所学的同一件事，正如圣经所说：\BibleQuote{祢将苦难放在我们的腰间，祢使人骑在我们的头上；我们经过火与水，祢却领我们到丰盛之地。}\par

至此，圣神的尊威已藉我们对手的证词得到证明；但在接下来的部分，虔诚的清流再次被异端的淤泥玷污。因为他说圣神\Quote{给那些正在竞技的人加油}：这句话使他陷入极度愚蠢和不虔敬的指责。因为在运动场上，有些人负责安排那些打算展示运动活力的人之间的比赛；另一些人则在力量和技巧上超越其余的人，为胜利而奋斗，并赤身露体彼此较量；而其余的人，则根据自己的好恶，分别支持某一位竞争者，在较量时为他加油，嘱咐他防备某种伤害，或记住某种摔跤技巧，或靠技巧使自己不被摔倒。从上述内容可以看出，\Person{欧诺弥乌斯}将圣神贬低到何等低下的等级。因为在赛场上，有些人安排比赛，另一些人判断比赛是否按规则进行，还有一些人实际参与比赛，另有一些人为参赛者加油——这些人公认远逊于运动员本身——\Person{欧诺弥乌斯}却将圣神视为旁观的人群之一，或视为随从运动员的人之一，因为圣神既不决定胜负，也不授予胜利，也不与对手较量，而只是加油，对胜利毫无贡献。因为祂既不参与争斗，也不注入斗争的力量，而只是希望祂所关心的运动员不会在较量中失利。因此，\Person{保禄}与\BibleQuote{对抗率领者，对抗掌权者，对抗这黑暗世界的统治者，对抗在诸天界里的邪恶属灵势力\BibleRef{弗 6:12}}摔跤，而大能的圣神并不坚固作战者，也不随己意将恩赐分给各人，\BibleQuote{随祂的意愿分给各人\BibleRef{格前 12:11}}，但祂的影响仅限于为那些正在竞技的人加油。\par

他又说：\Quote{使胆小者壮胆。}在这里，尽管他按照自己的方式延续先前对圣神的亵渎，真理却照样藉着敌人的口显明自己。因为除天主外，没有人能向恐惧者注入勇气，祂对胆小者说：\BibleQuote{不要害怕，因为我与你同在；不要惊慌\BibleRef{依 41:10}}，正如圣咏作者所说：\Quote{纵使我走过死荫的幽谷，我也不怕凶险，因为祢与我同在。}不仅如此，主亲自对恐惧者说：\Quote{你们心里不要烦乱，也不要胆怯，}又说：\BibleQuote{为什么你们胆怯，小信德的人哪？\BibleRef{玛 8:26}}，又说：\Quote{放心吧！是我，不要怕！}，又说：\Quote{你们放心！我战胜了世界。}因此，即使\Person{欧诺弥乌斯}无意于此，正统信仰仍藉着敌人的声音显示自身。接下来的句子也与前面一致：\Quote{关心一切，并显示所有关怀和远见。}因为事实上，唯独天主才拥有一切关怀和挂虑，正如伟大的\Person{达味}所说：\Quote{我是贫苦和匮乏的，但主照顾我。}如果剩下的部分似乎化为空洞的言语，有声无义，那么谁都不必责怪，因为在他所说的大部分话中，就任何健全的意义而言，他都虚弱且未经训练。因为他说\Quote{为了引导更善意的人和守护更忠信的人}究竟是什么意思，连他自己，以及那些愚蠢地赞慕他的愚行的人，也无法告诉我们。\par

\SourceRecord{Translated by H.C. Ogle and H.A. Wilson.；Nicene and Post-Nicene Fathers, Second Series；Vol. 5.；Edited by Philip Schaff and Henry Wace.；Buffalo, NY: Christian Literature Publishing Co.,；1893.；<http://www.newadvent.org/fathers/290102.htm>.}

% structure: p0001 source=h1 target=section reason=page-title

\section{\WorkTitle{驳斥欧诺米（第三卷）}}

% structure: p0002 source=h2 target=subsection reason=relative-heading-rank; base=h2

\subsection{一、这第三卷展示了欧诺米的第三次跌倒：他自相矛盾，有时说圣子因自然生育而当被称为\OriginalTerm{独生子}{Only-begotten}，且圣经从一开始就证明了这一点；有时又说，由于他是受造的，他不应被称为圣子，而是一个\Quote{\OriginalTerm{产物}{product},}或\Quote{\OriginalTerm{受造物}{creature}.}}

如果有人\Quote{合法地竞赛\BibleRef{弟后 2:5}}，他在比赛中的挣扎会因对手要么拒绝挣扎、自动退出以支持征服者，要么按照摔跤规则被三次摔倒（这样，在裁判的判决中证明得胜的人，就藉着宣告的荣耀获得冠冕的光荣）而找到其限度；那么，既然\Person{欧诺弥}在我们之前的论证中已被两次摔倒，却不承认真理应持有战胜谎言的标记，反而第三次用他的作品在他惯常的谎言的竞技场中扬起尘土反对敬神的教义，加强自己在欺骗一方的挣扎，那么，我们的真理陈述也必须被唤起击溃他的谎言，将希望寄托于那赐予胜利并审判胜利的主，同时从对手摔跤技巧的不公正中获得力量。因为我们不以承认我们没有准备任何藉修辞磨利的论证武器来帮助我们与那些列阵反对我们的人作战为耻，也没有辩证法的巧妙或锋利，这些在无经验的法官面前甚至会使真理蒙上谎言的嫌疑。我们的推理对抗谎言只有一个力量——首先是圣言本身，祂是我们言语的力量，其次是那反对我们的论证的腐朽，它自行倒塌。\newline{}\newline{} 现在，为了使所有人都尽可能清楚地看到，\Person{欧诺弥}的努力本身为那些与他争辩的人提供了挫败自己的手段，我将向读者陈述他的幻影学说（因为我认为那完全在真理之外的学说可以这样称呼），并且我希望你们所有参与这场战斗并观看我的教义和与之匹敌的教义之间正在进行的交锋的人，成为我们论证的合法竞赛的公正裁判，以便通过你们的公正裁决，敬神的推理得以向教会的整个剧场宣告为胜利者，在击败不敬神方面赢得无争议的胜利，并因其敌人的三次摔倒而被装饰以得救者的不朽冠冕。现在，这个陈述作为他第三篇演讲的前言被提出来反对真理，其形式如下：——\Quote{保持自然秩序，}他说，\Quote{并坚守那些从上而来为我们所知的事物，我们不拒绝用\InnerQuote{\OriginalTerm{受生的产物}{product of generation}}这一名称来谈论圣子，因为祂是受生的，既然受生的本质和圣子的称谓使得这样的词语关系是合适的。}我请读者仔细注意这一点，即当他称天主为\Quote{受生}和\Quote{圣子}时，他将这些名称的理由归于\Quote{自然秩序}，并援引自上而来的知识作证这一概念：因此，如果在接下来的内容中发现任何与他所立立场相悖的东西，那么所有人都清楚，他是在被自己的论证驳倒之前就已被自己推翻。因此，让我们用他自己的话来审视他的陈述。他承认，圣子的名称若非如他所说\Quote{自然秩序}证实了这一称谓，就绝不会恰当地应用在独生子天主身上。那么，如果从对\Quote{圣子}这一称谓的考虑中抽去自然秩序，他使用这一名称就会失去其恰当和自然的含义，变得无意义。此外，他说这些陈述之所以得到确认，是因为它们坚守从上而来的知识，这一事实为有关\Quote{圣子}这一称谓的正统观点提供了强有力的额外支持，因为从上面而来的圣经的默感教导证实了我们在这些问题上的论证。如果情况如此，并且这是不容欺骗的真理标准，即这两者一致——如他所说的\Quote{自然秩序}和从上而来的知识的见证证实了自然的解释——那么显然，断言任何与这些相悖的东西，无非是公开与真理本身作对。让我们再听听这位以自然为师的作者，在关于这个名称的事情上，并说他坚守圣人教导从上而来赐予我们的知识，在我刚才引用的段落之后，他在稍后进一步详细阐述的内容。因为我现在暂时省略他论文中接下来连续的叙述，以免他著作中的矛盾因中间内容的阅读而被掩盖。\Quote{同样的论证，}他说，\Quote{也将适用于受造和创造的事物，因为自然的解释和事物之间的相互关系，以及圣人的用法，都给予我们自由使用这个公式的权威：因此，如果有人将受造物与造物主相对应，将受造物与创造者相对应，他不会错。}他所说的受造或创造的产品是什么？它与其造物主或创造者在其名称中自然具有关系？如果是指我们在受造界中默观的可见和不可见之物（正如保禄所说，藉着祂万物被创造，无论是可见的还是不可见的），以至于这种名称的相对连接具有恰当和特殊的应用——受造物与造物主相对，受造物与创造者相对——如果这是他的意思，我们同意他。因为事实上，既然主是天使的造物主，天使当然是由造他的主所造的受造物；既然主是世界的创造者，显然世界本身及其中的一切都是造祂的创造的受造物。然而，如果他是有意用这个意图来解释\Quote{自然秩序}，将相对术语的专有化系统化以针对它们在言语意义上的相互关系，那么这就很奇怪了，因为每个人都意识到这一点，他竟会离开他的教义陈述，为我们勾画出一套文法琐事的系统。但如果他将这样的措辞应用于独生子天主，以致说祂是由造祂的那位所造的受造物，由创造祂的那位所造的受造物，并将这种术语归于\Quote{圣人的用法}，那么让他首先在他的陈述中向我们指明，他说有哪一位圣人宣称万物的造物主是产物和受造物，并且他追随谁而如此大胆措辞。教会知道那些心被圣神神圣引导的人是圣人——族长、立法者、先知、福音作者、宗徒。如果在这些人中，有人在其默感的言语中宣称统御万有的天主，祂\Quote{以祂权能的话支撑万有}，用手掌握万有，并藉着祂的意愿本身创造了宇宙，是受造物和产物，那么他将被原谅，因为他正如他所说，在着手系统化这样的教义时追随了\Quote{圣人的用法}。但如果圣经的知识自由地提供给所有人，并且任何选择分享神圣教导的人都不被禁止或隐藏，那么他如何能通过曲解圣经来误导听众，将应用于独生子的\Quote{受造物}一词归于\Quote{圣人的用法}呢？因为万物都是藉着祂受造的，你几乎可以从所有圣人的神圣言论中听到，从梅瑟和后来的先知及宗徒，他们特定的表述在这里详细列出会很冗长。对我们来说，与其他人和超越其他人，崇高的若望就足够了，他在关于独生子天主性的讲话的序言中大声宣告，万物中没有一样不是藉着祂受造的，这是一个无可争辩且确凿的证明，证明祂是受造界的主，不被算在受造物之列。因为如果所有受造物都只藉着祂而存在（若望见证，在受造界中，没有一样东西不是藉着祂受造的），那么谁会在理智上如此盲目，以至于看不到福音作者宣告中的真理：那位创造整个受造界的天主，确实是与受造界不同的某种东西？因为如果所有被算在受造物之中的事物都藉着祂而存在，而祂自己是\Quote{在起初}，是\Quote{与天主同在}，是天主、圣言、生命、光明、肖像、光辉，并且如果受造界中没有一样东西用同样的名称来称呼——圣言、天主、生命、光明、真理、肖像、光辉，或任何其他属于天主性的专有名称在受造界中都不被使用——那么显然，祂是这些事物，在本性上是与受造界不同的某种东西，受造界既不是也不是被称为这些事物中的任何一个。如果确实在这些措辞中存在受造界与其造物主之间的名称同一性，那么他或许可以因为其他名称的共同性而原谅自己也将\Quote{受造物}这一名称共同应用于受造物和造它的那位；但是，如果藉着名称在受造的本性和非受造的本性中所默观的特征，在任何情况下都不相容或共同，那么那个人敢于将奴役的名称应用于那位，正如圣咏作者所宣告的，\Quote{以祂的权威永远统治}，并藉着\Quote{受造物}的名称和概念，将那位正如宗徒所说\Quote{在一切事物中居首位\BibleRef{哥 1:18}}者，降低到奴役的本性水平，他的曲解怎能不昭然若揭呢？因为伟大的保禄宣告\Quote{整个受造界在奴役之中\BibleRef{罗 8:21}}——他在诸天之上的学校中被教导那不可言传的知识，在那里，任何藉着言语表达意义的嗓音都静默，沉默的默想成为教导的话语，藉着思想的无声光照，向纯洁的心传授那些超越言语的真理。那么，如果一方面保禄大声宣告\Quote{受造界在奴役之中}，另一方面独生子天主是真正的主和万有的天主，若望见证所有受造物都是藉着祂受造的，那么任何一个在某种意义下被算作基督徒的人，怎能保持沉默，当他看到\Person{欧诺弥}通过不一致和不连贯的系统化，藉着一个倾向于奴役的名称的同一性，将超越一切管辖和一切权威的主宰权能贬低到受造物的卑微状态？如果他说他有某些圣人宣告祂是奴隶、受造物、被造物或任何这些卑微和奴役的名称，那么这里就是圣经。让他，或别人代他，向我们提出一个这样的措辞，我们就保持沉默。但如果没有这样的措辞（并且在那些我们相信的默感圣经中，永远找不到任何支持这种不敬虔的思想），那么对于一个不仅曲解圣人的话，而且甚至反对自己定义的人，还有必要在已承认的问题上进一步争辩吗？因为如果\Quote{自然秩序}，正如他自己承认的，因着祂是受生的而额外证实了圣子的名称，并且名称的对应是根据受生者与生育者的关系，那么他如何将\Quote{圣子}一词的意义从其自然的应用扭转，并将关系改变为\Quote{受造物与其造物主}——这种关系不仅适用于宇宙元素，而且也可以断言于蚊子或蚂蚁——就这些每一个都是受造物而言，其名称与其造物主的关系同样等价？他的教义的亵渎性质，不仅从许多其他段落，而且从所引用的段落中已显而易见；至于他声称在这些表达中所追随的\Quote{圣人的用法}，显然根本不存在这样的用法。\par

% structure: p0004 source=h2 target=subsection reason=relative-heading-rank; base=h2

\subsection{二、他随后再次卓越、恰当且清晰地探讨并阐述了这句经文：\Quote{主创造了我。}}

也许箴言篇中的那段经文可能被拿来反对我们，而异端的拥护者常引用它作为主受造的证据——就是那句经文：\BibleQuote{主在祂道路的起始，在祂造化的起头，就造了我。}因为这番话是由智慧说出的，而伟大的保禄称主为智慧\BibleRef{格前 1:24}，他们便引用这段经文，仿佛独生天主自己，以智慧的名义，承认祂是由万物的创造者所造的。然而，我想，对于相当注意和勤勉的人而言，这句陈述的虔敬含义是清楚的，因此，对于那些受过箴言篇隐晦言辞教导的人，信德的教义并未受到损害。但我认为最好简要讨论一下关于这个话题应说些什么，以便当这段经文的意图被更清楚地解释时，异端教义便没有空间依据受默感的作者著作来大胆发言。普遍承认，\Quote{箴言}这个名称在圣经使用中，并非根据明显的意义，而是用于某种隐藏的含义，正如福音因此而将隐晦难懂的言论称为\Quote{箴言}；因此，如果一个人通过定义来阐述这个名称的解释，那么\Quote{箴言}是一种言谈形式，它通过直接呈现的一组观念，指向某种隐藏的东西，或者是一种不直接指出思想目标，而是通过间接的意指来进行教导的言谈形式。现今，这部书特别以这样一个名称作为标题，而智慧的撒罗满在前言中立刻解释了这称谓的含义。因为他并未称这部书中的言论为\Quote{格言}、\Quote{劝告}或\Quote{清晰教导}，而是称为\Quote{箴言}，并接着加以解释。这个词语的含义是什么？\Quote{认识}，他告诉我们，\Quote{智慧和教诲 \BibleRef{箴 1:2}}；并非按照其他学问中常见的方法，将智慧教导的课程摆在我们面前；相反，他要求一个人先通过预备训练成为有智慧的人，然后才能接受箴言所传达的教导。因为他告诉我们，有一些\Quote{智慧之言}通过\Quote{转折}来揭示其目标。因为那不能直接理解的东西需要某种转折才能把握隐藏的事；正如保禄在即将把历史的字面意义转换为寓意默观时，说他将要\Quote{改变声音 \BibleRef{迦 4:20}}；同样，在这里，撒罗满将隐藏意义的彰显称为\Quote{言语的转折}，仿佛思想的美无法被感知，除非一个人通过将表面意义翻转过来，获得思想启示的光辉，就像孔雀身后装饰的羽毛那样。因为在他身上，看见他羽毛背面的人会因其缺乏美和色彩而轻视它，视其为卑陋的景象；但如果有人将其翻转过来，向他展示另一面，他便能看见自然的多彩画卷，半圆形中央闪耀着紫色染料，金色的雾霭环绕圆圈，边缘闪烁着许多彩虹般的色调。既然言语中明显部分没有美（因为\Quote{君王女儿的一切荣耀都在内在}，以隐藏的装饰在金黄的思想中闪耀），撒罗满必然建议这部书的读者\Quote{言语的转折}，以便他们能\Quote{理解比喻和隐晦的话语，智慧者的言辞和谜语。}既然这种箴言式的教导包含了这些要素，一个理性的人不会不经检验和审视就接受从这部书中引用的任何段落，即便它初看起来非常清晰明白；因为肯定有一些奥秘的默观隐藏在那看似明显的段落之下。如果这部书中明显的段落都必然需要相当细致的审视，那么那些即使直接理解也向我们呈现许多晦涩难懂之处的段落，岂不更需要如此？\par

那么，让我们从所讨论经文的上下文开始考察，看看邻近语句的读法是否提供了任何清晰的意义。这段论述描述智慧以她自己的位格说出某些话语。每一位学生都知道，在智慧以谋略为居所、呼召知识和明达的经文中，她说她拥有力量和审慎（而她本身被称为明达），她行走在正义的道路上，她的交往在于公正审判的途径，并且她宣称：藉着她，君王执政，王侯书写公平法令，君主赢得他们土地的拥有权。现在，每个人都会看到，体贴的读者不会按照字面意义不加审视地接受所引用的任何短语。因为如果藉着她，君王被提升到统治地位，如果从她那里君主国获得力量，那么就必须得出结论：智慧向我们显示为立王者，并将那些在其王国中施行恶政者的责任转移到她自己身上。但是我们知道，有些君王确实在智慧的引导下走向那没有终点的统治——神贫的人，他们的产业是天国，正如主所应许的，祂就是福音的智慧；我们也认出这样的王侯，他们统治自己的情欲，不被罪的权势所奴役，他们将公平的法令铭刻在自己的生命上，如同写在版上。同样，那种可赞美的专制，它藉着智慧的联盟，将情欲的民主制改变为理性的君主制，使那些不受约束地奔向有害自由的事物——我指一切肉体和属世的思想——受到束缚：因为\BibleQuote{肉体相反圣神\BibleRef{迦 5:17}，}并反抗灵魂的治理。因此，这样的一位君主赢得了这片土地的拥有权，根据起初的创造，他是由圣言指定为统治者的。\par

有鉴于此，一切理性的人都承认，这些表述应当按照这样的意义来理解，而非按字面初看所呈现的意义，因此，我们所讨论的这句话，既然与它们紧密相连，审慎之人就不会毫不经审查地全盘接受。\Quote{我若向你述说，}她说，\Quote{每日发生的事，我就要记念从亘古以来述说的事：主创造了我。}请问，那字句的奴隶，像犹太人那样贴近音节的声音，对这句话能说什么呢？难道连词\BibleQuote{我若每日向你述说发生的事，主创造了我}在留心听的人耳中不是显得奇怪吗？仿佛她若不每日述说发生的事，就必然完全否认自己受造。因为说\Quote{我若述说，我就受造}的人，借着沉默让你理解为\Quote{我若不述说，我就不受造}。\Quote{主创造了我，}她说，\Quote{在祂道路的开端，为祂的工程。祂从亘古立了我，在起初，在未有大地以前，在未有深渊以前，在未有水泉涌出以前，在未有山岳安放以前，在未有丘陵以前，祂生了我。}这是受造物的何种新次序呢？首先受造，然后被立，最后受生。\Quote{主造了地，}她说，\Quote{甚至无人居住之地，以及天底下有人居住的地极。}她所说的主是谁，是无人居住和有人居住之地的创造者？当然是那创造了智慧的主。因为这两句话都出自同一位：一句说\BibleQuote{主创造了我}，另一句接着说\BibleQuote{主造了地，甚至无人居住之地。}因此，主同样是智慧本身和有人居住及无人居住之地的创造者。那么，我们该如何看待\BibleQuote{万物是藉着祂造成的；凡受造的，没有一样不是藉着祂造成的}这句话呢？因为如果同一位主既创造了智慧（他们建议我们理解为圣子），又创造了包含在受造界中的个别事物，那么崇高的\Person{若望}说万物是藉着祂造成的，如何为真呢？因为这段圣经听起来与福音相悖，它将无人居住和有人居住之地的创造归于智慧的创造者。同样，其后的一切也是如此：她论及天主在风中所设立的宝座，并说天上的云彩被造得坚固，天底下的泉源稳固；上下文包含许多类似的表述，强烈要求以精细明亮的智慧来解释，正如在已经引用的段落中应当注意的那样。什么是设立在风中的宝座？什么是天底下泉源的稳固？天上的云彩如何被造得坚固？若有人将此段按可见之物来解释，他会发现事实与话语颇不一致。因为谁不知道，天底下地极由于过冷或过热——或因过于接近太阳的热力，或因过于远离——而无法居住；有些地方过于干燥焦枯，另一些地方水分过多、为寒霜所冻，只有与两种对立极端距离相等之处才有人居住？但如果人类居住的是大地中央，那么箴言说天底下的地极有人居住又有何意义？再者，人能在云彩中感知到什么坚固，使这段经文按其表面意图而言为真，即说天上的云彩被造得坚固？因为云的本质是一种散布在空气中的非常稀薄的水汽，由于极其精微而轻盈，被空气的流动托起，当因压缩而被迫聚集时，从托住它的空气中以沉重的雨滴形式落下。那么，在这些一触即散的事物中，坚固何在？因为在云中，你能辨别出空气的稀薄和易分解的特性。再者，天主的宝座如何设立在本性不稳定的风上？至于她起初说自己是\Quote{受造}，最后说自己是\Quote{受生}，在这两句话之间说自己是\Quote{被立}，谁能给出一个与常识和显而易见的意义相一致的说明呢？我们先前在论证中质疑过的一点，即每日述说发生的事，并记念从亘古以来述说的事，这仿佛也是智慧宣称自己由天主所造的一个条件。\par

因此，既然以上所述已清楚表明，这段经文中没有任何部分应不经考察与反思便接受其字面意思，那么，或许正如对待其他经文一样，我们也不宜按字面所呈现的意义来解释\Quote{主创造了我}这句经文，而应以全部的注意与关怀，寻求从这话语中能虔诚理解的含意。现在，要完全把握眼前这段经文的意义，似乎只属于那些藉圣神之助探寻深奥之事、并知晓在圣神内讲述神圣奥秘的人；然而，我们的论述将仅限于讨论这段经文，以免完全忽略其要旨。那么，我们的论述是什么呢？我认为，由神圣光照而生于任何人心中的智慧，不可能离开圣神的其他恩赐而单独出现，反而必须同时有预言的恩宠进入其中。因为如果对现存事物真理的领悟是智慧的独特能力，而预言包含对将来事物的清晰认识，那么，一个人若未能藉预言之助将其知识进一步涵盖将来事物，他便不会拥有圆满的智慧恩赐。现在，既然\Person{撒罗满}所自称的并非单纯属人的智慧——他说\Quote{天主教导了我智慧}，又在说\Quote{我所有的话都出自天主}时，将自己所说的一切归于天主——那么，在这段\BibleBook{}{箴言}中，我们或许应探寻那与他的智慧相融合的预言。但我们认为，在本书前一部分，当他说\Quote{智慧为自己建造了一座房屋}时，他藉这些话隐约指向主肉身的准备；因为常言的智慧并不居住在他人所建的居所内，而是从童贞女的身体中为自己建造了那住所。然而，在此处，他于其论述中加入了由两者——即那房屋与建造房屋的智慧——所合成的整体，也就是说，人性与那与人相结合的神性；并且他对每一者都使用了适当且相称的措辞，正如你们在\BibleBook{}{福音书}中所见那样，那里的论述依照其主题，使用更高超且神圣的用语来指明天主性，而使用谦卑低下的用语来指明人性。同样，我们在此段中也可看到\Person{撒罗满}受预言感动，将降生成人的奥迹完满地传授给我们。因为我们首先谈论的是智慧的永恒能力与运行；在此，圣史在一定程度上与他用词一致。正如圣史以其概括性的话宣称祂是万物的原因与创造者，同样，\Person{撒罗满}也说，包含在万有中的那些个别事物都是由祂造成的。因为他告诉我们，天主以智慧奠定了大地，以睿智准备了诸天，以及按次序随后的一切，都保持同样的含义；并且，为了不显得忽略对人类美善恩赐的提及，他又以智慧的身份接着说我们稍早前提到的话，即\Quote{我以谋略、知识与明智为居所}以及所有与理智和知识的教导有关的话。\par

在叙述这些及类似的事情之后，他继续介绍他关于人的救恩计划的教导，即为何圣言成了血肉。因为显然，那超越万有的天主在祂内并无任何受造或从外进入之物，无论是德能、智慧、光明、圣言、生命、真理，或任何在神圣怀中圆满境界所默观的事物（这一切，那在父怀里的独生子就是），因此\Quote{受造}之名不能恰当地用于任何默观于天主内的事物，以致那在父内的圣子、在起初的圣言、在光明中的光明、在生命中的生命、在智慧中的智慧，不应说：\Quote{主创造了我。}因为如果天主的智慧是受造的（而基督是天主的德能和天主的智慧\BibleRef{格前1:24}），那么天主，按推论，便拥有祂的智慧作为从外进入之物，后来借着制作而获得祂起初没有的某种东西。但那在父怀里的，绝不容许我们认为父怀曾一度缺乏祂自己。那在起初的，肯定不属于从外在进入那怀中之物，而是作为一切圆满的圆满，被认为常是在父内，不等待借受造而在祂内升起，所以父不应被认为曾一度缺乏善，但那被认为在父的天主性永恒中的，常是在祂内，是德能、生命、真理、智慧等。因此，\Quote{创造了我}这些话并不出自神圣不朽的本性，而是出自那在降生成人中与祂混和的本性，即出自我们受造的本性。那么，为何同一位被称为智慧、明达、聪明的，建立大地、布置穹苍、剖开深渊，而在这里却是\Quote{在祂工作的开始受造}呢？他告诉我们，这样的救恩计划并非无重大原因而提出。但因为人接受了当遵守的诫命后，却因背命抛弃了记忆的恩宠，变得健忘，因此，\Quote{好叫我可以每日宣布为你们得救发生的事，并借着回忆你们已忘记的从亘古以来的事来提醒你们（因为我所宣告的不是新福音，而是致力于使你们恢复原初的状态）——为此我，永在者，受造了，并不需要受造才能存在；如此我就是天主工作的道路的开始，即对人的道路。因为第一条道路被毁坏了，必须再次为亡羊祝圣一条新的活路，即我自身，我就是道路。}这个观点，即\Quote{创造了我}的意思指向人性，神圣的宗徒更清楚地向我们展示，当他嘱咐我们\Quote{穿上主耶稣基督\BibleRef{罗13:14}，}以及同一位（使用同样的词）说：\Quote{穿上那按天主创造的新人。\BibleRef{弗4:24}}因为如果救恩的衣裳只有一件，那就是基督，就不能说\Quote{那按天主创造的新人}是基督以外的另一位，而是显然，那\Quote{穿上了基督}的，也\Quote{穿上了那按天主创造的新人}。因为实际上，只有祂被适当地称为\Quote{新人}，祂不是借着已知和惯常的自然方式出现在人类生活中，而是唯独在祂的情况中，受造以一种奇特而特殊的形式被重新建立。因此，当论及祂奇妙的诞生方式时，他称同一人为\Quote{那按天主创造的新人}；当注视那混和在创造这\Quote{新人}中的神圣本性时，他称祂为\Quote{基督}；如此，这两个名称（我是指\Quote{基督}之名和\Quote{那按天主创造的新人}之名）应用于同一位上。\par

既然基督就是智慧，让有理智的读者思考我们对手的说法和我们自己的说法，判断哪一个更虔诚，哪一个更好地保存经文中适合天主性体的概念；是那宣示万有的造物主和主被造成、并将祂置于与受束缚的受造物同等地位的说法，还是那顾念降生成人、并保持我们对天主性和人性概念之适当比例的说法，同时铭记伟大的保禄赞同我们的观点，他在\Quote{新人}中看到受造，在真正的智慧中看到创造的能力。并且，经文顺序也与此一致，与它所传达的教理一致。因为如果\Quote{途径的起始}未曾在我们中间受造，我们所期待的诸世代的基础就不会被奠立；主也不会为我们成为\Quote{未来世代之父}，除非依照依撒意亚，有一个婴孩为我们诞生，祂的名字被称为先知所赐的一切其他名号，并且也称为\Quote{未来世代之父}。如此，首先生效的是在童贞中所成就的奥迹和苦难的救恩计划，然后信德明智的建筑师奠定了信德的基础：这基础就是基督，未来世代之父，在祂身上建立了无尽世代的生命。当这一切成就，为使在每个信徒心中能成就福音法律的圣令和圣神的多样恩赐——（圣经形象地、以适当的含义将这一切命名为\Quote{山}和\Quote{丘陵}，称义德为天主的\Quote{山}，称祂的审判为\Quote{深渊}，并将那由圣言播种并结出丰硕果实者称为\Quote{地}；或者按达味教导我们的意义，通过\Quote{山}理解和平，通过\Quote{丘陵}理解义德）——智慧在信徒心中受生，这句话便应验了。因为那位在已接受祂的人中的，尚未在不相信的人中受生。因此，为使这些事在我们身上成就，它们的创造者必须在我们的内中受生。因为如果智慧在我们内中受生，天主就在我们每人内预备了土地和无人居住之地——土地，即接受圣言的播种和耕耘的；无人居住之地，即清理了邪恶居民的心——这样，我们的居所就将在极地之上。因为在地上，有深处，有表面，当一个人不埋在地里，或像住在地洞里，由于思念地下的事（如同生活在罪中的人，他们\Quote{陷入深淤泥中，没有立足之地}，他们的生命确实是坑，如圣咏说：\Quote{求坑口不要在我身上合上}）——我若说，当智慧在他内受生时，一个人思念上面的事，只按必要接触地，这样的人居住于\Quote{穹苍下地的极处}，不深陷于世俗思想；智慧与他同在，因为他在自己内预备了天而非地：并且，当人将诫命付诸实行，为自己加强高处云彩的教导，并以其精确的言谈，像用海滩围住广阔无边的罪恶大海，阻止混乱的水从他的口中流出；如果因教导的恩宠，他得以居于泉源之间，极其谨慎地倾泻言论的河流，以免给人喝的是毁灭的浊水而非纯净之水；如果他超越所有尘世途径，生活如气，向那被称为\Quote{风}的属灵生命迈进，以致他被分别出来，成为那位在他内安坐者的宝座（正如保禄为了福音被分别出来，成为承担天主之名的选定器皿，如另一处所说，他成为宝座，承载那坐在他上面的）——我说，当人通过这类方式如此建立，以致那已完全在自己内组成天主所居住之地的人，现在欢喜快乐，因为他被造成父亲，不是野性和无理性的野兽，而是人（这些将是神性的思想，按天主的肖像塑造，通过信德在祂身上——那位受造、受生、并在我们内被建立者——而信德，按照保禄的话，被理解为根基，智慧借此在信徒中受生，我所言的一切都成就）——我说，如此建立的人的生活真是有福的，因为智慧时时与他一致，并与他一同喜乐，他每日唯独在她内找到喜乐。因为主在祂的圣者中喜乐，天上为那些得救者也有喜乐，基督如父亲为祂获救的儿子设宴。虽然我们匆忙讨论了这些事，但让谨慎的人阅读圣经的原文，并将其晦涩说法与我们的反思相配，检验是否远更好认为这些晦涩说法之意涵有此指涉，而非一目了然所归于的。因为若望的神学不可能被认为真实，它陈述\Quote{所有受造物都是圣言的工程}，如果在此段落中，那创造智慧者被认为也连同她一起创造了一切其他事物。因为那样，万物将不是藉她而成，而她自己将被算在受造物中。\par

而谜语之言的所指，由后续经文清楚揭示，经文说：\Quote{现在，我儿，你要听从我；遵守我道路的，是有福的。}此处以\Quote{道路}显然指通往德行的途径，其开端即拥有智慧。那么，有谁瞻仰了神圣的圣经，会不同意真理的敌人既是不敬的又是诽谤的呢？——不敬，是因为他们尽力贬低独生天主不可言说的荣耀，并将其与受造物结合，试图表明那具有独生权能的主是由祂自己所造之物中的一员；诽谤，是因为尽管圣经本身并未给他们提供此类意见的根据，他们却武装自己反对虔敬，仿佛从圣经中引证。既然他们绝无法指出任何一段圣经经文引导我们将独生天主在时间之前的荣耀与受造界相提并论，那么在此得到证明之后，宜将战胜谎言的标记作为虔敬道理的见证提出，并扫除他们那些使受造物与造物主、受造之物与制造者相对应的言辞体系，我们应如来自天上的福音所教导的那样，宣认那蒙爱之子——不是私生子，不是赝品；而是接受\Quote{子}这个名号所自然包含的一切，说那出自真天主的乃是真天主，并相信凡我们在圣父身上所见的一切都在祂内，因为祂们是一体，在一体之中，另一体被理解，既不超越祂，也不低于祂，在任何神圣卓越的属性上都不改变或受变动。\par

% structure: p0012 source=h2 target=subsection reason=relative-heading-rank; base=h2

\subsection{三、然后，他从亚当和亚伯尔的例子以及其他例证，指明在\Quote{\OriginalTerm{受生}{generate}}与\Quote{\OriginalTerm{非受生}{ungenerate}}的事上，本质并无疏离。}

如今，既然欧诺米与自身的矛盾已显明，他曾说：\Quote{祂照本性应被称为\InnerQuote{子}，因为祂是受生的}，而后又说，因为祂是受造的，所以不再称为\Quote{子}，而是\Quote{产物}；我认为，细心留意的读者——因为两个彼此矛盾的陈述不可能同时为真——应从两者中弃绝那不敬亵渎的，即关于\Quote{受造物}与\Quote{产物}的说法，而只赞同那正统倾向的，即宣认为\Quote{子}的称号自然归于独生天主；这样，真理之言甚至由其敌人的声音所推荐。\par

然而，我重新开始论述，拿起我们原先搁置的那个论点。他说：\Quote{我们不拒绝}，\Quote{称圣子（因祂是受生的）为\InnerQuote{受生的产物}，因为受生的本质本身，以及\InnerQuote{圣子}的称号，使得这样的词语关系是恰当的。}同时，请批判地跟随论证的读者记住这一点：在论及独生子的情况下说到\Quote{受生的本质}，他由此允许我们在父的情况下说到\Quote{非受生的本质}，这样，无论是非受生还是受生，都不能再被认为构成本质，而必须分别看待本质，并且它的受生或非受生必须通过在其中思考的独特属性来分别理解。然而，让我们更仔细地考虑他关于这一点论证。他说有一个本质被生了，而这受生的本质的名称是\Quote{圣子}。好了，在这一点上，我们的论证将在两个方面定我们对手的罪：首先，是试图欺骗；其次，是他们对反对我们的尝试的懈怠。因为他是在耍花样，当他谈论\Quote{本质的生成}时，为的是要在本质之间建立对立，一旦它们按\Quote{受生}和\Quote{非受生}之间的本性差异而分开；而他们尝试的懈怠正由他们欺骗试图建立的立场本身所显示。因为那说本质是受生的人，清楚地定义生成是不同于本质的某种东西，这样生成的含义就不能归于\Quote{本质}这个词。因为他在这段话中并没有像他经常做的那样，说生成本身就是本质，而是承认本质是受生的，这样在他读者中每个词都产生不同的概念：因为听者听到它被生时产生一个概念，而\Quote{本质}这个名称唤起另一个概念。我们的论证通过例子可以更清楚。主在福音中说，妇人临近产期时忧愁，但后来因欢喜快乐，因为一个人生到了世上。因此，正如在这段话中我们从福音获得两个不同的概念——一个是出生，我们设想其为通过生成；另一个是出生所产生的东西（因为出生不是人，而是人因出生而有）——同样在这里，当\Person{尤诺米乌斯}承认本质是受生的，我们通过后一个词得知本质从某物而来，通过前一个词我们设想那从某物获得真实存在的实体本身。那么，如果本质的含义是一回事，而表达生成的词向我们提示另一个概念，那么他们聪明的计谋就完全化为乌有，如同瓦器互相投掷，彼此撞得粉碎。因为如果他们将对\Quote{受生}和\Quote{非受生}的对立应用于父和圣子的本质，他们将不可能同时将这些名称之间的互相冲突应用于事物本身。因为正如\Person{尤诺米乌斯}所承认的，本质是受生的（鉴于福音的例子解释了这样一个短语的含义，在那里，当我们听说一个人被生了，我们并不设想人和他的生成是同一回事，而是在这两个词中各得一个不同的概念），异端必定不再被允许用这样的词语来表达其关于本质差异的教义。然而，为使我们关于这些事的叙述尽可能清楚，让我们再次以下列方式讨论这一点。那创造宇宙的那一位在起初造了人的本性与万物，在\Person{亚当}被造之后，祂为人定下了彼此生成的律法，说：\BibleQuote{你们要生育繁殖} \BibleRef{创 1:28}。现在，当\Person{亚伯尔}通过生成而存在时，有理性的人谁会否认，在人类生成的实际意义上，\Person{亚当}是非受生地存在呢？然而，第一个人在自己身上拥有人的本质本性的完整定义，而从他所生的人被登记在同一个本质名称下。但如果那受生的本质被造成不同于那非受生的本质，那么同一个本质名称就不会适用于两者：因为那些本质不同的事物，其本质名称也不相同。因此，既然\Person{亚当}和\Person{亚伯尔}的本质本性以相同的特征为标记，我们必定同意一个本质在两者之中，并且两者在同一个本性中显示。因为就它们的本性定义而言，\Person{亚当}和\Person{亚伯尔}都是一个，但通过在每个个体中观察到的个体属性而彼此无混杂地区分。因此，我们不能恰当地说\Person{亚当}产生了另一个不同于自己的本质，而宁可说他自己产生了另一个自我，与这另一个自我一起产生了那生他者的全部本质定义。那么，我们在人类本性的事例中通过定义所给予我们的推论引导所学到的，我认为我们也应当用来引导对神圣教义的纯粹领会。因为当我们从神圣而崇高的教义中抖落一切属血肉和物质的观念时，当它被清除这些观念后，剩下的概念必最确实地引导我们到达崇高而不可接近的高处。甚至我们的对手也承认，那超越万有的天主是独生子的父，而且也被称为独生子的父，并且他们还称那出自父的独生子天主为\Quote{受生者}，因为祂是受生的。因此，既然在人中，\Quote{父亲}这个词带有某些含义，而纯全的本性与这些含义无关，人应当放下一切通过联想\Quote{父亲}一词的血肉含义而进入的物质观念，并在天主和父的情形下形成一个适合神性本性的观念，只表达关系的实在。因此，既然在人的父亲的概念中，不仅包括一切血肉提示给我们的思想，而且毫无疑问地，一种间隔的概念也与人的父亲身份的观念一同被设想，那么，在神圣生成的情形下，最好在拒绝身体污染的同时，也拒绝间隔的概念，以便那恰属于物质的东西可以被完全清除，而超越的生成不仅清除了一切情感的观念，也清除了间隔的观念。现在，那说天主是父的人，将把天主是的思想与另一个思想祂是某种东西结合起来：因为那从某个开端获得存在的事物，当然也从此事物获得其存在的开端，无论它是什么；但祂——在祂的情形中，存在没有开端——不从任何事物获得祂的开端，即使我们在祂内思考某些异于单纯存在的属性。好了，天主是父。由此可知，祂从永恒就是祂所是的：因为祂并不是成为父，而是\Emph{是}父：因为在天主内，那曾经是的，现在和将来都是。另一方面，如果祂曾一度不是任何事物，那么祂现在和将来也不是那事物：因为祂不被相信是那样一个存在的父，以至于可以虔诚地断言天主曾一度单独存在而没有那存在。因为父是生命、真理、智慧、光明、圣化、德能以及独生子所是或被称为的一切同类事物的父。因此，当对手断言那光明\Quote{曾一度不存在}时，我不知道对哪一方造成了更大的伤害：是光明，因为光明不存在，还是那拥有光明的，因为祂不拥有光明？对生命、真理、德能以及独生子以自己圆满充满父怀的其他一切特性也是如此。因为无论哪种方式，荒谬都是相等的，而对父的不敬将等于对圣子的亵渎：因为如果说主\Quote{曾一度不存在}，你将不仅断言德能的不存在，而是说天主的德能，即德能之父的德能，\Quote{曾不存在}。因此，你们的教义所主张的圣子\Quote{曾一度不存在}，无非确立了父方面一切善的缺乏。请看这些智者的敏锐导致何等结局，主的话如何因他们而应验，祂说：\BibleQuote{那轻视我的，轻视那派遣我者}。因为正通过他们轻视独生子任何时候存在的论证，他们也羞辱父，用他们的教义从父的光荣上剥夺每一个好名字和概念。\par

% structure: p0015 source=h2 target=subsection reason=relative-heading-rank; base=h2

\subsection{四、他如此显明永恒圣子与圣父的一体性、本质的同一性和本性的共通（其中包含对产酒的自然探讨），并指出在独生子的命名中，\Quote{圣子}和\Quote{产物}这两个词包含相似的关系概念。}

因此，以上所述已清楚暴露我们作者诡辩中的松懈之处；他试图通过称呼一方为\Quote{非受生}，另一方为\Quote{受生}的方法来确立独生子与圣父本质的对立，却因其自相矛盾的论证而被证明是在愚弄自己。因为从他的话语中首先显明，\Quote{本质}之名意谓一件事，\Quote{受生}之名意谓另一件；其次，与圣子一同出现的并非任何异于圣父本质的新异本质，而是父就其本性的定义所是的，那出自父的也同样是，因为按照我们对\Person{亚当}和\Person{亚伯}的思考所展示的论证真理，本性在圣子的位格中没有变成多样性。因为正如在那事例中，非由相似方式受生的，就本质的定义而言，与那受生的是同一的，并且亚伯的受生没有在本质上产生任何变化；同样，在这些纯正教义中，独生天主并没有藉自己的受生在自己内对非受生者的本质产生任何变化（他如福音所说出自父，并在父内），而是按照我们所宣认的简单而平实的信经语句，是\Quote{光明来自光明，真天主来自真天主}；那一位是另一位所是一切，只是不是那另一位。然而，关于他进行这种体系化的目的，我认为目前无需表达任何意见，无论是大胆而危险，还是容许而无危险，将用以指称天主本性的短语从一种转变为另一种，并用\Quote{受生之产物}的名称来称呼那受生的。\par

我放过这些事情，以免我的论述过于纠缠于次要之点的争斗而忽略了更重大的；但我认为我们应当仔细考虑这个问题：自然关系是否引入了这些术语的使用？因为\Person{欧诺弥}确实主张，名称的相近也意味着本质的关系。因为我想他不会说，名称本身，脱离所意指的事物的意义，就有任何相互的关系或相近；但所有人都通过词语所表达的意义来辨别名称的关系或差异。因此，如果他承认\Quote{圣子}与\Quote{圣父}有自然关系，让我们放下名称，考虑其含义中的力量，看在其相近中我们是否辨别出本质的差异，或是同质而特有的东西。说我们发现差异，简直是疯狂。因为没有亲属关系或共通的东西，如何在名称中\Quote{保持秩序}，相连而一致，其中\Quote{那受生的本质本身，如他所说，与\InnerQuote{圣子}的称号使这样的词语关系成为恰当}？反之，如果他说这些名称意指关系，他必然显得是本质共通性的拥护者，并维护这个事实：名称的相近也意指主体的联系；而这一点他在其作品中经常不自觉地在做。因为借着那些他试图摧毁真理的论证，他常常在不知不觉中被吸引去拥护他正反对的教义。历史关于\Person{撒乌耳}告诉了我们类似的事：有一次，他满怀愤怒反对先知，却被恩宠所胜，发现自己在被灵感之中（我想是预言的圣神愿意藉着他教训那叛徒），因此这事件的惊人性质成了他后来生活中的谚语，正如历史记载那样惊奇地表达：\Quote{撒乌耳也列在先知中吗？}\BibleRef{撒上 19:24}\par

那么，\Person{欧诺弥}在何时同意真理呢？当他说主自己\Quote{作为永生天主之子，不以自己从童贞玛利亚所生为耻，在言谈中常自称为\InnerQuote{人子}}？因为我们也引用这话作为本质共通性的证据，因为\Quote{圣子}的名称显示本性在两方面的共通性是同等的。因为正如他因肉身与所出生的那位的亲属关系而被称为人子，同样，他确实也因他的本质与那来自者的联系而被视为天主之子，而这一论据是真理最强大的武器。因为再没有比\Quote{圣子}这个名称（同样适用于神性与人性两种本性）更能清楚地指出那位\BibleQuote{天主与人之间的中保}\BibleRef{弟前 2:5}（正如大宗徒所称他的）。因为同一位既是天主之子，又在降生成人中成为人子，为藉着他与每一方的共融，亲自将本性所分开的联结起来。如果他在成为人子时没有参与人性，那么逻辑上可以说，他虽然为天主之子，也不分享天主本质。但如果整个人性复合体在他内（因为他\BibleQuote{在各方面曾与我们相似而受过试探，只是没有罪}\BibleRef{希 4:15}），那么必然相信超绝本质的一切属性也在他内，因为\Quote{圣子}这一词语为他要求两者的相似——在人内的人性，在天主内的神性。\par

那么，倘若如欧诺弥所说，名称表明关系，而关系的存在是在事物本身中观察到，而非仅在词语的声音中（我用事物指其自身所构想的事物，若如此谈论圣子和圣父不算过分鲁莽的话），谁会否认，这位亵渎的斗士已被自己的行动拖入为正统辩护的行列，用自己的手段推翻自己的论证，并在神圣教理中宣扬本质的相通呢？因为他不情愿地投向真理那边的论证，在这一点上并未说错——即，若名称的自然概念不能证实\Quote{圣子}这一称呼，他就不会被称为圣子。因为就如一张长凳不被称作工匠的儿子，也没有理智的人会说建筑者生出了房屋，我们也不说葡萄园是葡萄园丁的\Quote{产品}，而是把人制作的东西称为他的作品，把他所生的人称为人的儿子（我想，这是为了通过名称将适当的意义赋予各自的对象），同样，当我们被告知独生子是天主之子时，我们并不是通过这个称呼理解天主的一个受造物，而是理解\Quote{圣子}一词在其含义中真正显示的东西。即使圣经称葡萄酒为葡萄树的\Quote{产品}，即便如此，我们关于正统教理的论证也不会因名称的相同而受损。因为我们不把葡萄酒称为橡树的\Quote{产品}，也不把橡子称为葡萄树的\Quote{产品}，只有当\Quote{产品}与其来源之间有某种自然的共通性时，我们才使用这个词。因为葡萄树中的水分，通过髓从根部经茎干抽出，在其自然能力中是水；但它沿着自然的途径有序地经过，从最低处流向最高处，就转变为葡萄酒的品质，太阳光线在某种程度上促成了这一变化，它们通过温暖将水分从深处抽到嫩芽，并通过适当成熟的进程使水分成为葡萄酒：因此，就其本性而言，存在于葡萄树中的水分与从中产生的葡萄酒之间没有区别。因为一种水分来自另一种水分，不能有人说葡萄酒的原因除了自然存在于嫩芽中的水分之外还有别的。但是，就水分而言，品质的差异并不产生变化，只有当某种特性将存在于葡萄酒形式中的水分与嫩芽中的水分区分开来时，才被发现，这两种形式之一伴随着涩味、甜味或酸味，因此，在实体上两者是相同的，但由性质上的差异区分开来。因此，正如当我们从圣经中听到独生天主是人子时，我们通过名称所表达的同族关系了解他与真人的亲属关系，同样，即使圣子被对手的措辞称为一个\Quote{产品}，我们仍然通过这个名称了解到他在本质上与\Quote{产生}他的圣父的亲属关系，因为被称为葡萄树\Quote{产品}的葡萄酒已被发现，就水分的概念而言，并不异于存在于葡萄树中的自然能力。事实上，如果有人明智地审查我们对手所说的话，它们倾向于我们的教理，其意义反对他们自己的捏造，因为他们竭力在每一点上确立他们的\Quote{本质差异}。然而，要猜测他们是从何处被引导到这样的概念，绝不是一件容易的事。因为如果\Quote{圣子}的称呼不仅仅表示\Quote{出自某物}，而是通过其含义特别向我们呈现，如欧诺弥本人所说，本性上的关系，并且葡萄酒不被称作橡树的\Quote{产品}，而那些\Quote{产品}或\Quote{毒蛇的种类}，福音某处所说的，是蛇而不是羊，那么，很明显，在独生子身上也是如此，\Quote{圣子}或\Quote{产品}的称呼不会传达与另一种类事物之关系的含义：但即使按照我们对手的措辞，他被称作\Quote{生育的产品}，而\Quote{圣子}的名称，如他们所承认的，涉及本性，那么圣子确实出自那生育或\Quote{产生}他的圣父的本质，而不是出自我们视为外在于那本性的其他事物中的某一个的本质。而且，如果他真是从他而来，那么他就不会与他的来源所属于的一切相异，正如在其他情况中也已表明，凡通过生育而存在的一切，显然与其来源同属一类。\newline{}\newline{}（译者注：本段中\Quote{产品}对应原文\OriginalTerm{产品}{product}，\Quote{毒蛇的种类}对应圣经经文\BibleRef{玛 3:7}或\BibleRef{路 3:7}，此处未直接引用，故不加圣经标记。）\par

% structure: p0020 source=h2 target=subsection reason=relative-heading-rank; base=h2

\subsection{五、他论述天主本质的不可理解性，以及对撒玛黎雅妇人的话：\Quote{你们朝拜你们所不知道的。}}

现在如果有人要求对天主本质作某种解释、描述和说明，我们并不否认在此种智慧上我们是无知的，只承认这一点：那本性无限者不可能被任何言辞所表达的概念所把握。天主伟大无边这一事实由预言所宣告，它明确宣称祂的光辉、祂的荣耀、祂的圣洁\Quote{没有穷尽}；如果祂的周围没有界限，那么祂本身——无论其本质为何——更在任何方面不受任何限制所把握。因此，如果通过言辞和名称的解释按其含义意味着对主题的某种把握，而另一方面，无限者不能被把握，那么没有人可以合理地责怪我们的无知，如果我们对于无人应冒昧之事不敢妄为的话。因为我能用什么名称描述那不可理解者？用什么言语宣告那不可言说者？所以，既然天主性过于卓越崇高，无法用言语表达，我们学会了以沉默尊敬那超越言语和思想的事。如果那\Quote{看得过高，超过他所当看的}\BibleRef{罗 12:3}的人践踏我们这句谨慎的话，嘲笑我们对不可理解之事的无知，并在那无形、无限、无量、无穷者（我指圣父、圣子、圣神）中认出了不相似的差异，并且提出那句常被欺骗之徒所引用的短语来责备我们的无知：\Quote{\InnerQuote{你们崇拜你们所不认识的,} 如果你们不认识你们所崇拜之物的本质}，我们将听从先知的劝告，不怕愚人的辱骂，也不因他们的诽谤而大胆谈论不可言说之事，而以那位不擅言辞的\Person{保禄}为我们超越知识之奥秘的老师。他远不认为天主本性在人的感知范围内，以至于称天主的审判为\Quote{不可探究的}，祂的道路为\Quote{不可追踪的}\BibleRef{罗 11:33}，并断言为爱祂的人因今生善行所应许的事物超越理解，以致不能用眼睛看见，不能藉耳朵听见，也不能用心容纳。因此，从\Person{保禄}学习了这些，我们大胆宣告：不仅天主的审判对试图探究它们的人太高，而且那导向认识祂的道路甚至至今仍无人踏足、不可通行。因为我们理解宗徒想表达的就是这个意思：当他称那导向不可理解者的道路为\Quote{不可追踪的}，藉此短语表明那知识是人间的推理无法达到的，并且从未有人将他的理解置于这样的推理路径上，或藉感知的方式为那不可理解之事显示出任何接近的痕迹或迹象。\par

我们从宗徒的高超言语中学习了这些，然后用所引的经文这样推理：如果祂的审判不可探究，祂的道路不可追踪，祂美善的应许超越我们猜测所能形成的任何表象，那么祂实际的天主性，就其不可言说和不可接近而言，比那些被视为伴随它的属性更高更崇高，而那些属性，受天主启示的\Person{保禄}宣称是无法认识的。藉此我们证实了我们所持的、他们嘲笑的道理，承认自己在认识那超乎认知范围的事物方面不如他们，并宣告我们真正崇拜我们所认识的。现在我们认识我们所崇拜之主的荣耀崇高，正是因为我们不能通过推理在我们的思想中领悟祂伟大无比的特质；而我们敌人用来攻击我们的主对撒玛黎雅妇人说的话，倒可能更适合对他们说。因为主对撒玛黎雅妇人说：\Quote{你们崇拜你们所不认识的，}她当时在关于天主的观念上受到物质思想的偏见。这句话对她很适用，因为撒玛黎雅人以为他们崇拜天主，但同时认为天主性物质地固定在某一地方，所以他们只是名义上朝拜，实际崇拜的是别的东西，不是天主。因为任何被视为有限的东西都不是神圣的；属于天主性的是无所不在、渗透万物、不受任何事物限制。所以那些反对基督的人发现他们用来攻击我们的这句话反过来成了对自己的控告。因为，正如撒玛黎雅人以为天主性被某种地点的范围所包围，因而受到他们听到的话语的责备：\Quote{\InnerQuote{你们崇拜你们所不认识的,} 你们的敬礼对你们毫无益处，因为被认为固定在某个地方的天主不是天主。}——同样，有人可以对新撒玛黎雅人说：\Quote{你们以为天主性因缺乏受生而受到限制，如同被某种地方界限限制，\InnerQuote{你们崇拜你们所不认识的,} 你们确实向祂敬礼如同天主，却不知道天主的无限性超越一切名称所能提供的意义和理解。}\par

% structure: p0023 source=h2 target=subsection reason=relative-heading-rank; base=h2

\subsection{六、随后他阐释了\Quote{子}的名称、\Quote{受生者}的名称，以及多种多样的\Quote{子}——天主的、人的、公羊的、丧亡的、光明的、白昼的。}

但我们的论述已偏离当前主题太远，在不断追索由推论而产生的各种问题时。因此，让我们再次回到论述的连贯性上，因为我认为，通过我们所说的，正在审查的语句已被充分证明不仅与真理相矛盾，而且也与自身相矛盾。因为，按照他们的观点，与圣父的自然关系是由\Quote{子，}以及由\Quote{受生者，}对他所生的那位（正如这些人的智慧根据某种语法上的轻浮，将表示天主性本质的术语错误地塑造成一种文字排列）而确立的，那么没有人再能怀疑：由本质确立的名词相互关系是其同类，更确切地说是本质同一的证据。但不要让我们的论述仅仅围绕反对者的话打转，以免正统教义似乎只是因反对者的软弱而获胜，而应显得自身拥有充足的力量。因此，让我们自己以更有力的辩护尽可能加强相反的论点，以便人们能充分自信地认识到我们力量的优越性，因为我们将反对者所遗漏的那些论点也带到真理无误的检验中。为反对者辩护的人或许会说，\Quote{子，}或\Quote{受生者，}的名称丝毫不能确立本质上的同类关系。因为在圣经中，使用了\Quote{义怒之子}一词，以及\Quote{丧亡之子 \BibleRef{若 17:12}，}和\Quote{毒蛇的种类；}在这些名称中，显然看不出任何本质的共同性。因为被称为\Quote{丧亡之子}的\Person{犹达斯}，其实体与丧亡并不相同，按我们对该词的理解。因为在\Person{犹达斯}身上，\Quote{人}的含义是一回事，\Quote{丧亡}的含义是另一回事。同样地，从相反的例子也可以确立这一论点。因为那些在某种意义下被称为\Quote{光明之子，}和\Quote{白日之子 \BibleRef{得前 5:5}，}的人，就其本性的定义而言，与光明和白日并不相同；石头也成为\Person{亚巴郎}的儿女，当他们因信德和行为声称与他有亲属关系时；而那些\Quote{被天主的圣神引导}的人，如宗徒所说，被称为\Quote{天主的儿女 \BibleRef{罗 8:14}，}却在本性上与天主并不相同；人可以从受默感的圣经中收集许多这样的例子，藉着这些例子，欺骗如同某种以圣经的见证装饰的肖像，伪装成真理的样式。\par

那么，我们对此如何回应？神圣的圣经知道如何从两种意义上使用\Quote{子}一词，以致在有些情况下这种称谓源于本性，在其他情况下则是附加的、人为的。因为当它说到\Quote{人的儿子们，}或\Quote{公绵羊的儿子们，}时，它标记了受生者与其所源出自者的本质关系；但当它说到\Quote{有能者的儿子们，}或\Quote{天主的儿女，}时，它向我们呈现那种因选择而产生的亲属关系。而且，在相反的意义上，同一些人也被称为\Quote{厄里的儿子们，}和\Quote{贝里雅耳的儿子们，}\Quote{儿子们}的称谓容易适用于这两种观念。因为当他们被称为\Quote{厄里的儿子们，}时，他们被宣告与他有自然关系；但当被称为\Quote{贝里雅耳的儿子们，}时，他们因选择的邪恶而受责备，因为他们在生活中不再效法他们的父亲，而是将自己的意志依附于罪恶。那么，在我们这低下的本性以及与我们相关的事物的情况下，由于人的本性同样倾向于两边（我指的是恶与善），我们有能力成为黑夜或白日之子，同时我们的本性，就其主导部分而言，仍在其适当界限内。因为无论是因罪恶而成为\Quote{义怒之子}的人，并未脱离他的人性出身；也无论是因选择而投身于善的人，并未因其习惯的精炼而拒绝他的人性起源；而是在每种情况下，他们的本性保持不变，而他们意向的差异承担了那亲属关系的名称，按照他们因德行而成为天主的儿女，或因恶习而成为相反的儿女。\par

但是，欧诺弥——至少在天主道理上，他是那位\Quote{保存自然秩序}（因为我将使用我们作者的原文），\Quote{并且坚持那些从起初就为我们所知的事物，不拒绝称那受生者为\InnerQuote{受生之物}的名称，因为受生本质本身}（如他所说），\Quote{和\InnerQuote{圣子}的称谓使这样的词语关系成为合适的}——他又如何使受生者与生他者在本质上疏远亲属关系呢？因为对于因责备而被称为\InnerQuote{儿子}或\InnerQuote{产物}的人，或者在某些赞美伴随这些名称的情况下，我们不能说任何人被称为\InnerQuote{义怒之子}，而同时实际上是由义怒所生；也没有人曾有日子作他的母亲，在肉身意义上，以至于他被称为它的儿子；而是他们意志的差异造成了这种亲属关系的名称。然而，欧诺弥说，\Quote{我们不拒绝称圣子——既然他是受生的——为\InnerQuote{受生之物}的名称，因为受生本质}，他告诉我们，\Quote{和\InnerQuote{圣子}的称谓使这样的词语关系成为合适的}。那么，如果他承认这样的词语关系是由圣子确实是\Quote{受生之物}这一事实而成为合适的，那么将这种名称的原理同样应用于那些因比喻而不精确使用的词，以及那些如欧诺弥所说，自然关系使名称使用成为合适的词，又怎会恰当呢？确实，这样的说法只在那些本性处于德行与恶行之间的人身上成立，他们常常轮流分享相反类别的名称，由于与善或与其对立面的亲近，成为光明之子，或黑暗之子。但在没有对立的地方，就不能再说\Quote{圣子}一词是比喻性地应用的，就像那些通过选择将称号归于自己的人那样。因为无法得出这种观点：正如一个人丢弃黑暗的行为，通过端正的生活成为光明之子，同样，独生子天主也是从低劣状态改变后获得更尊贵的名称。因为一个人通过属灵重生与基督结合而成为天主的儿子；但那亲自使人成为天主子嗣的那位，并不需要另一位圣子来将儿子的义子关系赐予他，而是拥有他本性所是的名称。人自己改变自己，将旧人换成新人；但天主会改变成什么，以便接受他所没有的呢？人脱去自己，穿上神性；但那始终如一者，他脱去什么，或穿上什么呢？人成为天主的儿子，接受他所没有的，抛弃他所拥有的；但那从未处于恶习状态者，既没有什么要接受，也没有什么要放弃。再者，当按照本性说话时，人可以真实地被称为某人的儿子；另一方面，当生活的选择赋予名称时，也可以不精确地这样称呼。但天主，作为唯一的善，具有单一而纯一的本性，永远看向同一方向，从不被选择的冲动所改变，而是永远愿望他所是的，并且必然是他所愿望的：所以他在两方面都适当而真实地被称为天主子，因为他的本性包含善，他的选择也永不脱离那更优越的，因此这个词被用作他的名称而没有不精确性。因此，这些论证（我们假借对手之口对自己提出的）没有余地被我们的对手作为对本性上亲缘关系的抗辩而提出。\par

% structure: p0027 source=h2 target=subsection reason=relative-heading-rank; base=h2

\subsection{七、然后他以阐释独生子之神性与人性名称，并讨论\Quote{受生}与\Quote{非受生}这两个词来结束本书。}

但他们不知为何憎恶真理，确实称祂为\Quote{子}，但为了避免这个称呼所见证的\OriginalTerm{本质的相通}{community of essence}，他们将这个称呼与名称所含的意义分开，只把\Quote{子}这个空名赐予独生子，仅将这个词的声音施予祂。我所说的真实不虚，且我并非错误地瞄准对手的标志，这可以从他们对真理的实际攻击中清楚得知。他们提出这样的论证来建立自己的亵渎主张，即圣经教导我们独生子有许多名字——石头、斧头、磐石、根基、饼、葡萄树、门、道路、牧者、泉源、树、复活、导师、光，以及许多类似的名字。但我们不可虔诚地按照这些名称的直接意义来理解并称呼主。因为确实，认为那无形无体、单纯无像者，应按照这些名称的明显意义——无论它们是什么——来塑造，这极为荒谬。以至于我们听到斧头时，就想到某种铁的形状；听到光时，就想到天空的光；听到葡萄树时，就想到由插条生长的植物；或听到其他任何一个名称时，就像其通常用法提示我们所想的那样。但是我们将这些名称的意义转移到更适合天主本性的事物上，形成另一种概念。如果我们这样称呼祂，并非因为祂的本性定义就是这些事物中的任何一个，而是因为祂被这样称呼，同时通过所采用的名称而被理解为与这些事物本身不同的某物。但是，如果这些名字确实真地用来描述独生子天主，而不包含祂本性的宣告，那么他们说，同样地，我们也不应承认\Quote{子}这个名称按照普遍用法作为表达本性的意义，而应也为这个词找到一种不同于平常和明显的意义。这些以及诸如此类的，就是他们用来建立\Quote{子不是祂所是和所被称为的}这一观点的哲学论证。我们的论证本在奔向一个不同的目标，即表明\Person{欧诺弥}的新言论是虚假且不一致的，既不与真理相符，也不与自身相符。然而，由于我们用来攻击他们学说的论点被引入讨论，反而成了他们亵渎的某种支持，那么最好先简要讨论他的观点，然后再有序地审查他的著作。\par

那么，对这些无关的事我们能说什么呢？正如他们所说，圣经应用于独生者的名号虽多，但我们坚持，其他名号中没有一个与指向生祂者的指涉紧密相连。因为我们并不像使用\Quote{圣父之子}之名那样，以指向万有之天主的方式来使用\Quote{磐石}、\Quote{复活}、\Quote{牧人}、\Quote{光}或其余任何名号。我们可以如同按科学法则，将神圣名号的意义作双重划分：一类属于祂崇高而不可言喻的光荣的指示；另一类则表明眷顾安排的多样性。因此，我们以为，若那领受祂恩惠者不存在，那么那些表明祂恩惠的词语也不会被应用于它们。另一方面，所有表达天主属性的名号，都合适而恰当地应用于独生者天主，无关乎眷顾的对象。但为清楚地阐述这一道理，我们将审视这些名号本身。若非为栽植那些扎根于祂的人，主就不会被称为葡萄树；若非以色列家的羊迷失，祂就不会被称为牧人；若非为病人之故，祂就不会被称为医生；除非祂以某种眷顾的行动，使这些称号有利于那些蒙受祂恩惠的人，否则祂也不会为自身接受其他这些名号。有何必要一一列举，并在公认之点上延长我们的论证呢？另一方面，祂确实被称为\Quote{子}、\Quote{右手}、\Quote{独生者}、\Quote{圣言}、\Quote{智慧}、\Quote{德能}以及所有其他这类关系名号，因为它们与圣父在某种关系连接中一同被命名。因为祂被称为\Quote{\Emph{天主的}德能}、\Quote{\Emph{天主的}右手}、\Quote{\Emph{天主的}智慧}、\Quote{\Emph{圣父的}子与独生者}、\Quote{\Emph{与天主同在的}圣言}，其他名号也是如此。因此，从我们所述的可以推出，在每个名号中，我们都应思考某种适合该主题的适当意义，以免误解它们，偏离虔诚的教训。因此，正如我们将其他每个术语转移到它们可能应用于天主的意义上，并抛弃它们的直接意义，从而不理解为物质的\Quote{光}、被践踏的\Quote{路}、农耕生产的\Quote{饼}或言语表达的\Quote{言}，而是代之以向我们呈现天主圣言德能之伟大的一切思想——同样，如果有人抛弃\Quote{子}一词的普通和自然意义（藉此我们得知祂与生祂者同体），他当然会将这名号转移到某种更神圣的解释。因为既然在其他每个术语中，向更荣耀意义的转变使它们适于表达天主的德能，那么这名号的意义当然也应转移到更崇高的事物上。但是，若我们按反对者的观点，抛弃与生祂者之间的自然关系，那么我们在\Quote{子}的称谓中能找到什么更神圣的意义呢？我想无人如此胆大不敬，以至于认为在论及天主性体时，卑微低下者比崇高伟大者更为适宜。因此，如果他们能发现任何比这更崇高的意义，以至于认为独生者拥有圣父的性体是一件不值得思考的事，那就让他们告诉我们，他们是否在隐秘的智慧中知道任何比圣父的性体更崇高的事物，以便将独生者天主提升到这一层次，也将祂提升到与圣父的关系之上。但若天主性体的崇高超越一切高度，超出一切令人惊叹的能力，那么还有什么观念能将\Quote{子}之名的意义带到更伟大的事物上呢？因此，既然承认每一应用于独生者的重要短语，即使这名号源于我们低级生活的日常用法，都恰当地应用于祂，其意义朝着更崇高的方向有所改变，并且既然显示我们在\Quote{子}的称号上找不到比向我们呈现祂与生祂者之间关系的实在性更高贵的观念，我认为我们无需在此主题上花费更多时间，因为我们的论证已充分表明，不应像解释其他名号那样来解释\Quote{子}的称号。\par

但我们须将我们的探讨再次带回到那本书。它不适合同一些人既\Quote{不拒绝}（因为我要用他们自己的话）\Quote{称那受生者为\InnerQuote{\OriginalTerm{受生产物}{product of generation}}，因为受生的本质本身和圣子的名号都使这样的措辞关系相称}，而又将自然属于他的名号改换为隐喻的解释：以致他们遭遇了两种情况之一——要么他们的初次攻击已经失败，他们徒然逃向\Quote{自然秩序}，以确立称那受生者为\Quote{受生产物}的必要性；要么，如果这个论据成立，他们将会发现他们的第二个论据被他们已经确立的东西化为乌有。因为那由于受生而被称为\Quote{受生产物}的那一位，绝不能因同样理由而被称为\Quote{\OriginalTerm{受造产物}{product of making},}或\Quote{\OriginalTerm{受造之物}{product of creation}.}因为几个术语的含义相去甚远，一个审慎使用措辞的人应当根据主题适当地使用语词，以免我们不恰当地互换词语的含义而陷入观念的混乱。因此，我们把由技艺制成的东西称为工匠的作品，把由人所生的人称为那人的儿子；没有一个理智清醒的人会把作品称为儿子，或把儿子称为作品；因为那是通过错误使用名称而混淆和模糊真实意义的人的语言。因此，我们必须对独生子天主真实地确认以下两者之一：如果他是圣子，他就不能被称为\Quote{受造之物}；如果他是受造的，他就与\Quote{圣子}的名号无关，正如天、海、地以及一切个别事物，因为是受造的，就不采用\Quote{圣子}的名称。但是，既然欧诺米证明独生子天主是受生的（而敌人的证言对于确立真理更有价值），他藉着说他是受生的，也必然证实了他不是受造的。关于这些点，已经足够了；因为虽然许多论证涌到我们面前，但我们将满足于我们已经就眼前的主题所提出的那些，免得数量过多导致比例失调。\par

\SourceRecord{Translated by H.A. Wilson.；Nicene and Post-Nicene Fathers, Second Series；Vol. 5.；Edited by Philip Schaff and Henry Wace.；Buffalo, NY: Christian Literature Publishing Co.,；1893.；<http://www.newadvent.org/fathers/290103.htm>.}

% structure: p0001 source=h1 target=section reason=page-title

\section{驳斥欧诺弥（第四卷）}

% structure: p0002 source=h2 target=subsection reason=relative-heading-rank; base=h2

\subsection{一、第四卷讨论关于\OriginalTerm{受生之物}{product of generation}本性的论述，以及独生子的无受情受生，以及经文：\BibleQuote{在起初已有圣言}，以及童贞女的诞生。}

也许，是时候在我们的论述中审视那关于\Quote{受生之物}本性的论述了，这乃是他那可笑哲学思辨的主题。他说道（我将逐字重复他反对真理的优雅论证）：— \Quote{谁如此漠不关心、不注意事物的本性，以至于不知道，在地上所有身体中，在它们生成与被生成、主动与被动的过程中，那些生成者经考察被发现传递自己的本质，而那些被生成者自然地接受同一本质，因为物质因和从外部流入的供应对两者是共同的；而被生之物是通过受情而生，而那些生者自然地有并非纯正的行动，因为它们的本性与各类受情相联？}\newline{}请看，他是以何等适宜的文体，在他的思辨中讨论那在起初就已存在的天主圣言在时间之前的受生！他仔细审查事物、地上物体、物质因、生成之物与被生成之物的受情，以及所有其余诸如此类——任何有理解力的人都会对此脸红，即使这些话是用来说我们自己的，如果我们那为受情所束缚的本性就这样被他的话语暴露于嘲笑之中。然而，这就是我们的作者对关于独生天主进行的辉煌的探究。然而，让我们放下抱怨（因为叹息又能帮助我们推翻敌人的敌意什么呢？），并尽我们所能，使我们所引用的话的意义广为人知——这思辨乃是针对什么样子的\Quote{产物}而提出的——是按肉体而存在的呢，还是应当沉思于独生天主之内的呢？\par

由于这思辨具有双重性，一关于神圣、单纯且非物质的生命，一关于物质且受情的存在，且\Quote{受生}一词用于两者，我们必须使我们的区分清晰明确，以免\Quote{受生}一词的含糊性以任何方式歪曲真理。既然，通过肉体的进入存在是物质的，并受情的推动，而非形体、不可触摸、无形体、且没有任何物质混杂的存在，则与任何容受情的状态都无关，那么，考虑我们是由什么样子的受生进行询问就是恰当的——是纯洁神圣的受生，还是受情和玷污的受生。\newline{}现在，我想，无人会否认，关于独生天主，供\Person{欧诺弥}的论述考虑的乃是时间之前的存在。那么，他为何徘徊于这有形本性的论述，借其论证的可憎呈现玷污我们的本性，并公开陈列围绕人类受生的种种受情，而遗弃摆在他面前的主题呢？因为我们需要学习的，并不是这种借助肉体完成的动物性受生。当一个人审视自己并思考自身内的人性时，他如此愚蠢，以至于寻求另一位自己本性的解释者，并需要被告知所有包含在肉体受生观念中的不可避免的受情——即生者以一种方式被影响，被生者以另一种方式——从而人应从这一教导中得知他自己是通过受情而生的，且受情乃其自身受生之开端？因为这些事不论是略过还是说出，不论是人详述这些秘事，还是将不当说的事静默隐藏，我们都不无知于这一事实，即我们的本性是通过受情而进展的。但我们所寻求的是，关于独生子的崇高而不可言宣的存在，给出一个清晰的论述，祂因此被视为出自圣父。\par

现在，这就是摆在他面前的探究，我们这位新神学家却用\Quote{涌流}、\Quote{受情}、\Quote{物质因}，以及某种\Quote{行动}它\Quote{并非纯正}而脱离污染，以及所有其他此类短语来丰富他的论述。我不知道他是受什么影响，他因智慧高超，声称没有任何不可理解之物超出他的知识，并承诺解释圣子不可言宣的受生，却放弃了他面前的问题，像一条鳗鱼一样一头扎入他论证的污浊淤泥中，效法那位夜间来访的\Person{尼苛德摩}。当我们的主教导他关于从上而来的受生时，尼苛德摩在思想中冲向子宫的凹陷，并提出了怀疑，说一个人怎能再次进入母腹，用这些话：\BibleQuote{怎能成就这事？}认为他会以老人不能再在母亲腹内重生的事实证明属神的受生的不可能。但主纠正了他的错误观念，说肉体和圣神的属性是不同的。让欧诺弥也，如果他愿意，通过类似的反思纠正自己。因为我想，沉思真理的人应当根据其本身的属性来沉思其主题，而不是通过对物质事物的指控来诽谤非物质事物。因为如果一个人，或一头公牛，或任何其他由肉体所生之物，在生成或被生成时并非没有受情，那么这与那无受情、无朽坏的本性有何相干？我们是可朽的这一事实，并不构成对独生子不死性的反对；人们倾向于邪恶，也不使神圣本性中所有的不变性变得可疑；我们本身的任何其他属性也不转移给天主；反而，人类生命与神圣生命的本性是分离的，没有共同基础，它们的区别属性完全分开，以致后者的属性不从前者的属性中被理解，反之亦然，前者的属性也不从后者的属性中被理解。\par

那么，当谈论天主的受生时，\Person{欧诺米乌斯}为何放下主题，长篇大论地讨论地上的事物，而在这一点上我们与他并无争论？显然，这位工匠的目的很清楚——他通过诽谤性地暗示情感，来对主的受生提出异议。在此，我略过他那观点亵渎的本质，而钦佩他的敏锐——他是多么谨记自己的热心努力，他先前通过陈述确立了圣子必须是、也必须被称为一个\Quote{\OriginalTerm{受生的产物}{product of generation}}的理论，现在却又主张我们不应怀有关于祂受生的概念。因为，如果如这位作者所想象的，一切受生都与情感的状态相连，那么我们就绝对被迫承认：与情感无关的，也与受生无缘；因为如果情感与受生被视为结合在一起，那么凡不与其中之一有分者，也不会与其他有分。那么，他如何因祂的受生而称祂为一个\Quote{产物}？而他用现在的论据试图证明祂并没有受生。他又为何要攻击我们的导师，这位导师在天主道理上劝诫我们不可擅自造名，而应承认祂是受生的，而不将此概念转化为一个名称的公式，以致称那受生者为\Quote{一个受生的产物}，正如这个术语在圣经中正好应用于无生命之物，或应用于那些被称为\Quote{作为邪恶的预表}的事物？当我们提到避免使用\Quote{产物}一词的妥当性时，他准备好他那战无不胜的修辞术，并搬出他那冷冰冰的语法术语来支持自己，通过他巧妙地误用名称、或一词多义（或者随便人们怎么称呼他的手法），我说，他就通过这些手段推出他的三段论结论：\Quote{不拒绝称那受生者为\InnerQuote{受生的产物}的名字。}然后，我们刚承认了这个术语，并着手考察这个名称所包含的概念（根据这种理论，\OriginalTerm{本质的共通性}{community of essence}由此得到辩护），他就又收回自己的话，主张\Quote{受生的产物}并非受生，并通过他对肉体受生的污秽描述来提出反对，以及对圣子纯洁、神圣、无情感的受生提出异议，理由是：这两件事——与圣父的真正关系，以及祂本性的无情感——不可能在天主身上同时存在；相反，如果没有情感，就没有受生；如果有人承认真正的父子关系，那么他在承认受生的同时，当然也承认了情感。\par

崇高的若望并非如此说话，那宣告神学奥秘的雷声也并非如此。他称祂为天主子，并从一切激情观念中净化自己的宣告。请看，福音的开端，老师何等谨慎地预备我们的耳朵，何等周详，唯恐听众中有人因无知而陷入对主题的低下观念，滑入任何不协调的想法。为了将未经训练的听觉尽可能远离激情，他不在开场言词中提及\Quote{圣子}或\Quote{圣父}或\Quote{\OriginalTerm{受生}{generation}}，免得有人先听到\Quote{圣父}便急于联想到该词明显的含义，或得知\Quote{圣子}的宣告便按通常意义理解此名，或在\Quote{受生}一词上绊倒，如\BibleQuote{绊脚石}\BibleRef{伯前 2:8}；他以\Quote{元始}代替\Quote{圣父}，以\Quote{在}代替\Quote{受生}，以\Quote{圣言}代替\Quote{圣子}，并宣告\BibleQuote{在\OriginalTerm{元始}{Beginning}已有圣言}。请问，这些词\Quote{元始}、\Quote{在}、\Quote{圣言}中有何激情？\Quote{元始}是激情吗？\Quote{在}暗示激情吗？\Quote{圣言}藉激情而存在吗？抑或我们必须说，既然所用术语中找不到激情，宣告中也不表达亲密关系？然而，还有什么能比这些词更有力地显示圣言与元始在本质上的共融、真实的关系及同永恒？因为他没有说\Quote{从元始圣言被生}，以免将圣言与元始因任何时间延展的观念分开，而是同时宣告元始与在元始中的祂，使\Quote{在}一词共通于元始与圣言，好叫圣言不落后于元始，而是在对元始的信德一同进入时，其宣告先于我们的听闻，在听闻单独接纳元始本身之前。然后他宣告：\BibleQuote{圣言与天主同在。}福音作者再次为我们未经训练的状态担忧，再次畏惧我们幼稚未受教导的状况：他尚未将\Quote{圣父}的称谓托付给我们的耳朵，免得任何更属血肉的人，得知\Quote{圣父}后，由理解引导，进而想象一位母亲。他也尚未在宣告中提及圣子；因为他仍怀疑我们倾向于低下本性的习惯，担心有人听到圣子，会以激情的观念将天主性人性化。为此，他再次重复宣告，又称祂为\Quote{圣言}，将祂的本性向不信的你们如此陈明。正如你的言语出自你的心智，无需激情介入，同样，在此听到圣言时，你应领悟那出自某物的，而不领悟激情。因此，再次重复宣告，他说：\BibleQuote{圣言与天主同在。}哦，他如何使圣言与天主相称！不如说，他如何将无限与无限并论！\BibleQuote{圣言与天主同在}——圣言的整个存有，无疑，与天主的整个存有同在。因此，天主有多大，显然，与祂同在的圣言也有多大；所以，若天主是有限的，那么圣言也必受限制。但若天主的无限超越界限，那么与祂一同被默观的圣言也不被界限和尺度所包含。因为无人能否认圣言与圣父的整个天主性一同被默观，以致认为天主性的一部分在圣言内，另一部分没有圣言。若望属神的声音再次发言，福音作者在宣告中再次温柔地顾念孩童之人的听闻：我们尚未因听闻他的最初话语而成长到足以听到\Quote{圣子}仍能站稳，不被习惯的意义所动摇。因此，我们的传报者再次高声呼喊，在第三次宣告中仍宣告\Quote{圣言}，而非\Quote{圣子}，说：\BibleQuote{圣言就是天主。}他先宣告祂在哪里，然后宣告祂与谁同在，如今宣告祂是什么，藉第三次重复完成了他宣告的对象。因为他说：\Quote{我要向你们宣告的，不是容易理解的圣言，而是以圣言之名表示的天主。}因为这位在元始、并与天主同在的圣言，并非天主之外的任何东西，而是祂自己也是天主。随即，这位传报者达到其崇高言论的顶峰，宣告他所宣告的天主，是万物藉祂而造的那位，是生命，是人的光，是照耀黑暗的真光，却不被黑暗遮蔽，住在自己的民中，却不被自己的民接纳，并成为血肉，藉着血肉在人的本性中支搭帐幕。在经历了这些数目和多样的陈述之后，他才提及圣父和独生子，此时已无危险，因经过如此多防范的洁净，不会因\Quote{圣父}一词的含义而沉沦于任何沾染污秽的意义。因为他说：\BibleQuote{我们见了祂的光荣，正是圣父独生子的光荣。}\par

那么，\Person{欧诺弥}，请重复，重复你向福音作者提出的这个巧妙异议：\Quote{你如何在你论述中给予\InnerQuote{父}这名号，又如何给予独生子，既然所有肉体的受生都是藉情欲运作的？}真理确实替你回应了他：神学的奥秘是一回事，变动不居的身体的生理学是另一回事。它们彼此之间有广阔的间隔隔开。你为什么在你论证中将不能混合的东西结合起来？你为什么用你污秽的论述玷污天主受生的纯洁？你如何用影响身体的情欲为无形体者制定体系？停止从下面的事物中引出对上面事物本性的叙述。我宣告主是天主子，因为来自天上的福音，通过光明的云彩，如此宣告了他；因为他说：\BibleQuote{这是我的爱子\BibleRef{玛 17:5}.}然而，虽然我被教导他是子，我并没有被这名号拖到\Quote{子}的世俗意义上，而是既知道他是从父而来，又不知道他是从情欲而来。此外，我还要对所说的话补充这一点：我知道甚至有一种肉体的受生是没有情欲的，因此即使在这点上，\Person{欧诺弥}关于肉体受生的生理学也被证明是错误的，也就是说，如果能找到一种不允许情欲的肉体诞生的话。告诉我，圣言成了血肉，还是没有？我想你不会说没有。那么，它成了血肉，没有人否认。那么，\BibleQuote{天主在肉身内显现}是怎么发生的？你当然会说：\Quote{藉着诞生。}但你说的是什么样的诞生？显然，你说的是来自童贞的诞生，并且\BibleQuote{那在她内受孕的是出于圣神}，以及\BibleQuote{她分娩的日期满了，就生了\BibleRef{路 2:6},}而她的纯洁在她生育时丝毫没有受损。那么，你若相信，你就相信那从妇女而来的诞生是没有情欲的，但你却拒绝接受从父而来的神圣而不朽的受生，为要避免受生中情欲的观念。但我深知，他在他的教义中并非要避免情欲，因为他根本不在神圣而不朽的本性中分辨情欲；而是为了使万有的创造者被视为受造物的一部分，他编造这些论据以否认独生子天主，并利用他假装对情欲的谨慎来帮助他完成任务。\par

% structure: p0009 source=h2 target=subsection reason=relative-heading-rank; base=h2

\subsection{二、他证明\Person{欧诺弥}将适用于大地存在的术语用于独生子，从而显示他的意图是要证明子是可变、受造的存在物。}

他通过反对我们论点的争辩非常清楚地表明了这一点，他说：\Quote{子之本质是从父而来受生，不是以延伸的方式发出，不是通过流溢或分裂与那生他的结合分离，不是以成长的方式得以成全，不是以变化的方式转变，而是仅凭生者的意志获得存在。}什么人的心智感官没有被封闭，以至于不明白这些话表明\Person{欧诺弥}正通过这些陈述将子表现为受造物的一部分？有什么阻止我们逐字逐句地说这一切，关于我们在受造界中观察的每一件事物？如果你愿意，让我们将这定义应用于出现在受造界中的任何事物，如果它不容许同样的推论，我们就该谴责自己，因为我们轻率地检查了这定义，没有以适合真理的细心。那么，让我们更换子的名称，并逐字阅读这定义。我们说\Emph{大地}的本质是从父而来受生的，不是以延伸或分裂的方式与那生它的结合分离，不是以成长的方式得以成全，不是以变化的方式发出，而是仅凭那生它的意志获得存在。我们所说的有什么不适用于大地的存在吗？我想没有人会说有。因为天主并没有通过延伸发出大地，也没有通过流溢或从与自身的结合中分离而使它的本质存在，也没有通过从微小到规模的完整性的逐渐成长来带出它，也没有通过经历突变或变化而被塑造为大地的形状，而是他的意志足以使一切受造之物存在：\BibleQuote{他说了，它们就受生,}因此连\Quote{受生}这个名称也与大地的存在相符。如今，如果这些话可以真实地用于宇宙的各部分，那么对我们的反对者的教义还有什么怀疑呢？即尽管在言辞上他们称他为\Quote{子,}但他们却表现他像是通过创造而存在的事物之一，只是在次序上先于其余的一切？正如你可能会说关于铁匠的工艺：从中产出所有由铁打造的东西；但夹钳和锤子这件工具，铁藉以被塑造成有用之物，其存在先于其他器物的制造；然而，尽管它优先于其他，但在材料方面，塑造的工具和被工具塑造的铁之间并无差别（因为两者都是铁），只是形状一个比另一个早而已。这就是异端关于子的神学——以为主自身和他所造之物之间没有差别，除了在次序方面的差别。\par

在任何意义上被列为基督徒的人中，有谁承认世界的各部分和创造世界的那位的本质定义是相同的呢？就我而言，我因这亵渎而战栗，因为知道若事物的定义相同，则其本性也没有差别。正如伯多禄、若望及其他人的本质定义是共同的，其本性是同一的；同样，如果主在本性上与世界的各部分一样，那么他们必须承认主也受制于那些它们所察觉的属性，无论这些属性是什么。世界并不永存：因此，按照他们的说法，主也将与天地一同消逝，如果正如他们所说，祂与世界是同一类的。另一方面，如果祂被承认为永恒的，那么我们也必须认为世界也并非完全不具有天主性，如果如他们所说，世界在受造方面与独生子相对应的话。你可以看到这精妙的推论过程如何使论点倾向于，如同一块从山脊脱落、凭自身重量冲下山坡的石头。因为要么世界的元素必须是神圣的（按照希腊人的愚昧信仰），要么圣子不应当受钦崇。让我们这样思考。我们说受造物，无论是心智所察觉的还是感官所察觉的本性，都是从无中生有的：他们也宣称主也是这样。我们说一切受造之物因天主的旨意而存立：他们也告诉我们独生子也是如此。我们相信天使的受造物和世界的受造物都不是造物主的本质：他们也使主远离圣父的本质。我们承认万物都服事创造它们的那一位：他们也持有关于独生子的同样观点。因此，必然地，无论他们对受造物有何种观念，他们也会将这些属性归于独生子：无论他们对祂相信什么，他们也会对受造物作同样的观念：以致如果他们承认主为天主，他们也会将受造界的其余部分神化。另一方面，如果他们界定这些受造物不分享天主性，那么他们也不会拒绝关于独生子的同样概念。再者，没有一个理性的人宣称受造物具有天主性。那么……我也不说出其余的话，以免我的舌头助长敌人的亵渎。让那些口舌善于亵渎的人说出随之而来的结论吧。但即使他们保持沉默，他们的教义也是显而易见的。因为两件事中必有一件发生：——要么他们将废黜独生的天主，以致对他们而言，祂既不再存在，也不再被称为天主；要么，如果他们宣称祂具有天主性，他们也将同样宣称所有受造物具有天主性；——或者（这仍为他们保留着），他们将回避两端都显而易见的邪恶，而投靠正统教义，并必然同意我们：祂不是受造的，以便他们能承认祂是真实的天主。\par

需要花时间从这开端起一一列举他的教义所隐含的其他亵渎言论吗？因为凭我们所引用的，考虑其推论的人将会明白：谎言之父、死亡之制造者、邪恶之发明者，虽然被造于一种理性无形体的本性中，但并未被那本性所阻碍，而通过变化成为他现在所是的。因为本质的可变性，随意朝任一方向运动，涉及一种本性能力，它跟随意志的冲动，从而成为其意志所导向者。因此，他们将界定主甚至能具有相反的情操，藉着受造的概念将祂拉低到与天使同等的地位。但让他们听从保禄伟大的声音吧。为什么他说唯独祂被称为子？因为祂\Emph{不是}天使的本性，而是更卓越的。\BibleQuote{因为祂曾对哪一个天使说过：\InnerQuote{你是我的儿子，我今日生了你}？当祂再次引领首生子到世上来时，祂说：\InnerQuote{天主的众天使都要钦崇祂。}关于天使，祂说：\InnerQuote{祂使祂的天使成为风，使祂的仆役成为火焰}；但关于圣子，祂说：\InnerQuote{天主啊，你的宝座是永永远远的；你国度的权杖是正直的权杖}，以及预言在宣告祂的天主性时与这些言辞一同列举的其他所有内容。}他还从另一篇圣咏加上相宜的话：\BibleQuote{主啊，你在起初奠定了大地的根基，诸天是你手的化工}，以及其余，直到\BibleQuote{但是你却相同，你的岁月没有穷尽}，借此描述祂本性的不变和永恒。那么，如果独生子的天主性超越天使的本性，如同主人高于他的奴仆，他们又如何通过详述关于祂存在方式的、本应适用于我们在受造物中所见的个别事物的论述，而将祂这位受造界的主宰与可感的受造物视为同类，或者将祂这位被天使所钦崇者与天使的本性视为同类呢？正如我们已经表明，异端关于主的描述密切且恰当地适用于大地的制造。\par

% structure: p0013 source=h2 target=subsection reason=relative-heading-rank; base=h2

\subsection{三、随后他再次卓越地讨论了宗徒四次使用的术语 \OriginalTerm{首生子}{\GreekText{πρωτότοκος}}。}

但为使本作品的读者不再有任何模糊之处，以致为异端学说提供支持，我们或许值得在已审查的段落之外，再补充圣经中的这一点。他们可能会从我们引用的宗徒著作本身提出一个问题：\Quote{如果祂不是受造之物，怎能被称为\InnerQuote{受造物的首生者}呢？因为每个首生者不是与他种事物，而是与同类事物的首生者：如\Person{勒乌本}，在出生上优先于那些在他之后被数算的人，他是首生者，一个人是人的首生者；还有许多其他人被称为与他们同列的弟兄的首生者。}他们说：\Quote{我们断言，那\InnerQuote{受造物的首生者}具有我们视为一切受造物本质的同一本质。如果整个受造物与万有的圣父是同一本质，我们就不否认受造物的首生者也如此；但如果万有的天主在本质上不同于受造物，我们就必须说，受造物的首生者在本质上与天主也没有共通之处。}我认为，这个异议的结构，在我们这里引述的形式中，并不比对手可能用来反对我们的形式逊色。但我们在这点上应当知道的，现在将尽我们所能，在我们的论述中清楚地阐明。\par

宗徒在所有著作中四次使用了\Quote{首生者}或\Quote{首生者}这个名号；但他以不同的意义，并非以相同的方式提及这名号。因为有时他说\Quote{一切受造物的首生者}，又说\Quote{在许多弟兄中作首生者，\BibleRef{罗 8:29}}，然后说\Quote{从死者中首生者；\BibleRef{哥 1:18}}；在致希伯来人书中，\Quote{首生者}这名号是绝对的，单独提及：因为他这样说：\Quote{再者，当天主引领首生者进入世界的时候，说：\InnerQuote{天主的众天使都要崇拜祂。}}由于这些段落如此不同，最好分别解释每一个：祂如何是\Quote{受造物的首生者}，如何是\Quote{在许多弟兄中}，如何是\Quote{从死者中}，以及如何，单独提及而与这些无关，当祂再次被带入世界时，被所有天使崇拜。那么，如果你愿意，我们就从最后提到的开始审视我们面前的段落。\par

\Quote{再者，当祂引领，}他说，\Quote{首生者进入世界。}添加\Quote{再者}表明，借此词的力度，这件事并非第一次发生：因为我们用这个词表示已经发生过的事情的重复。因此，他借此短语指末世时审判者可畏的显现，那时祂不再以仆人的形象显现，而是坐在祂国度的宝座上的荣耀中，并被环绕祂的所有天使崇拜。因此，那一次进入世界，成为\Quote{从死者中}、以及\Quote{祂的弟兄中}、以及\Quote{一切受造物中}的首生者，当祂再次进入世界，作为以正义审判世界者，如预言所说，并未抛弃祂曾为我们而接受的首生者名号；但正如因耶稣之名，天上一切无不屈膝，同样，众天使的团体都崇拜那以首生者之名而来的祂，因他们为人类的复原而欢喜，祂藉着在我们中间成为首生者，使我们回复到起初的恩宠。因为既然天使为那些从罪恶中被拯救的人而欢喜（因为直到如今，受造物因影响我们的虚幻而一同叹息、同受产痛\BibleRef{罗 8:19}，认为我们的丧亡是他们的损失），当天主子女的显现——他们所期待和盼望的——发生时，当羊被安全地带到那上百只羊之上（而我们——即人类——正是那只被善牧藉成为首生者所拯救的羊），那时他们尤其会以对我们深厚的感恩，向天主献上崇拜，因祂藉成为首生者，使那从父家迷失的人复原。\par

既然我们已经明白了这些话语，再没有人会对其他经文犹豫不决，即祂因何缘故是首生者，是\Quote{死者中的}，或是\Quote{受造物中的}，或是\Quote{众多弟兄中的}。因为所有这些经文都指向同一个要点，尽管每段经文都提出某种特殊的概念。祂是死者中的首生者，祂自己首先解除了死亡的痛苦\BibleRef{宗 2:24}，为要使复活的诞生也成为众人的道路。再者，祂成为\Quote{众多弟兄中的首生者}，祂通过水中重生的新生在我们之前出生，这新生由鸽子盘旋作助产，祂使那些与祂有相同出生的人成为祂的弟兄，并成为那些在祂之后由水和圣神而生者的首生者。简而言之，正如我们身上有三次出生，人性借此得以活着：一次是肉身的，一次在重生圣事中，另一次是我们所盼望的死者复活；祂在所有这些中都是首生者——对于由两者（圣洗与复活）所完成的二重重生，祂是每一样的领导者；而在肉身中祂是首生者，因为祂首先且独自为自己谋划了那个本性所不知的诞生，那是人类许多世代中从未有人创始的。如果这些经文被正确理解，那么祂是首生者的\Quote{受造物}的含义，也不会为我们所不知。因为我们认识到我们本性的双重受造：第一次是使我们被造，第二次是使我们被再造。但我们若没有因不顺从而使第一次归于无效，就不需要第二次受造了。因此，当第一次受造已经衰老消逝时，就有必要在基督里有一个新受造（如宗徒所说，他断言我们在第二次受造中不应再看到衰老之物的任何痕迹，他说：\Quote{脱去旧人及其行为和情欲，穿上新人，这新人是按照天主而创造的}，又说：\Quote{若有人在基督内，他就是一个新受造：旧事已过，看，一切都成了新的}）——因为起初和后来的人性创造者是同一位。\Emph{那时}祂取尘土从大地造人；后来，祂从童贞女取尘土，不只造人，更是将自己包裹在人当中；\Emph{那时}，祂创造；后来，祂被创造；\Emph{那时}，圣言成了肉躯；后来，圣言成了血肉，为要改变我们的血肉成为神，因为祂与我们一样有分于血肉。因此，对于这基督内的新创造，祂自己是开始的，祂被称为首生者，作为万物的初果，无论是被生入生命的，还是因死者复活而活过来的，\Quote{祂可以为死人和活人的主，\BibleRef{罗 14:9}}，并借着在祂自己内的初果圣化整个面团。至于\Quote{首生}这个称谓不适用于圣子在其先于时间的存有中，这由\Quote{独生子}的称号证明。因为真正是独生子的没有弟兄，因为若与弟兄同列，又怎能是独生子呢？但是，正如祂被称为天主和人、天主子和人子——因为祂有天主的形状和奴仆的形状，有些是根据祂至高的本性，有些是在祂爱人的施恩中成为的——同样，祂是独生子天主，却成为一切受造物的首生者——独生子，那在父怀里的，却在新受造得救的人中，成为并被称为受造物的首生者。然而，如果如异端所主张，祂被称为首生者是因为先于其他受造物而被造，那么这个称号并不符合他们对独生子天主的坚持。因为他们不是这样说——圣子和宇宙同样来自父——而是说，独生子天主由父所造，其余万物由独生子所造。因此，基于同样的理由，当他们主张圣子是被造的，便称天主为被造者的父；同样，当他们说万物由独生子天主所造，他们给祂的名称不是祂所造之物中的\Quote{首生者}，而更恰当地是它们的\Quote{父}，因为在这两种情况下，对受造物的相同关系逻辑上导致相同的称谓。因为如果那超越万有的天主不被恰当地称为\Quote{首生者}，而是祂自己所造之物的父，那么独生子天主也必然被同样的推理称为\Quote{父}，而不是恰当地称为祂自己受造物的\Quote{首生者}，因此\Quote{首生者}这个称谓将是全然不当且多余的，在异端的观念中毫无地位。\par

% structure: p0018 source=h2 target=subsection reason=relative-heading-rank; base=h2

\subsection{四、他继续讨论主的受生之无受难性；以及\Person{尤诺米乌斯}的愚昧，他声称受生的本质包含圣子的称谓，却又忘记这一点，否认圣子与圣父的关系：此处他谈到\Person{喀耳刻}和曼德拉草的毒。}

然而，我们必须回到那些把情欲与天主的受生联系起来的人，并因此而否认主真是受生的，以避免情欲的观念。说受生必然与情欲相连，并因此，为使天主性保持在纯洁中、不受情欲影响，我们应当认为圣子与受生的观念无关，这或许在那些容易被欺骗的人看来是合理的，但受过天主奥秘教导的人，基于公认的事实，有一个现成的答案。因为谁不知道，正是受生引领我们回到真正而蒙福的生命，它不同于那种藉着\BibleQuote{不是由血慾，也不是由肉慾}而发生的受生，其中充满了流变与改变，逐渐成长至圆满，以及我们在尘世受生中所观察到的一切；而另一种则被认为是来自天主，是属天的，正如福音所说：\BibleQuote{由上}，这排除了血肉的情欲？我推测他们都承认这种受生的存在，并且在其中找不到情欲。因此，并非所有受生在本性上都与情欲相连，而是物质的受生受制于情欲，非物质的则纯洁无情欲。那么，是什么迫使他将本属于肉身的东西归于圣子不朽的受生，并通过用那不雅的生理学嘲笑低级的受生形式，而将圣子排除在圣父的亲密关系之外呢？因为即使在我们自身的情形中，也是受生作为两种生命的开端——通过肉身的受生，带来有情欲的生命；属神的受生，带来纯洁的生命（任何多少算是名列基督徒中的人都不会反驳这一说法）——当思考涉及不朽本性的受生时，怎么能允许怀着情欲的观念呢？此外，除了我们已经提到的那些，让我们再审视这一点。如果他们因为肉身所受的情欲而不相信天主受生的无情欲特性，那么让他们也从同样的记号（我指的是我们自身可以找到的记号）拒绝相信天主作为造物主是既无情欲也无劳累地行动的。因为如果他们通过我们自身条件的比较来判断天主性，他们就不能承认天主受生或创造；因为这两项活动我们自己都不能无情欲地进行。所以，让他们要么从天主性中分离创造和受生，以便他们能在两方面维护天主的无情欲性，并让他们为了确保圣父安全地远离情欲的范围，既不受创造的劳累，也不被受生玷污，而完全从他们的教义中拒绝相信独生子；或者，如果他们同意一种活动是由天主的能力无情欲地进行的，就让他们不要争论另一种活动：因为如果祂无需劳力和物质就创造，祂也肯定无需劳累和流变就受生。\par

在此，我又一次在论证中得到了欧诺弥的支持。我将简洁扼要地陈述他的谬论，概括他的全部意思。人并不能为我们制造材料，而只是通过其技艺赋予物质以形式——这是他在大量无稽之谈中所说的要旨。那么，如果他将概念和形成理解为属于较低级的受生，并因此禁止对受生的纯粹概念，那么，按照同一推理，既然尘世的创造忙于形式，却不能连同形式一起提供物质，他也就同样应禁止我们设想圣父是创造者。反之，如果他拒绝按照人的能力尺度来设想天主内的创造，那么他也就应停止以人的不完美来诽谤神圣的受生。然而，为了更清楚地证明他在论证中的精确与周密，我将再次回到他陈述中的一个小点。他断言\Quote{主动者与被动者彼此共享本性}，并在提及肉身的受生之后，提到\Quote{工匠在材料上展示的作品}。现在，让敏锐的听众注意他如何在此偏离了他的正确目标，如何在他碰巧编造的任何陈述中游荡。他在按肉体方式生成的事物中看到了\Quote{主动与被动被设想为具有同一本质，一方赋予本质，另一方接受本质}。因此，他知道如何准确地辨别现存事物本性的真理，从而将赋予者与接受者从本质中分离出来，并说每一个都独立于本质而自身有别。因为接受或赋予者当然是与被给予或被接受者不同的另一者，所以我们首先必须设想一个自身单独存在、在其独立存在中被观看的某一位，然后才能说他给予他所拥有的，或接受他所没有的。当他以如此可笑的方式喷出这个论证时，我们这位博学的朋友并没有察觉，下一步他再次推翻了自己。因为那个凭其技艺随意塑造面前材料的人，当然在此活动中是\Emph{主动}的；而材料在接受成形时，受了技艺者的塑造，是被动地受影响：因为材料并非通过保持不受影响、无动于衷而获得其形式的。那么，即使在技艺作品中，若没有被动与主动共同作用，也无物能够生成，我们的作者如何能认为他在此坚守了自己的话？既然他宣称行动与受动的关系包含本质共通，他似乎不仅以某种方式见证了受生者与生养者之间的本质共通，而且也使整个受造界与其造物主成为同一本质——如果如他所说，主动与被动应按本性界定为相互关联。因此，正是通过他用来确立其愿望的那些论证，他推翻了自己努力的主要目标，并通过自己的论点使同体圣子的荣耀更加稳固。因为如果从某物中起源表明生养者的本质存在于被生者中，并且如果人工制造（通过主动与被动完成）按照他的说法，使制造者与产物都归于本质共通，那么我们的作者在他自己的著作中多次主张主是受生的。因此，正是通过他试图证明主与圣父本质不同的那些论证，他为主建立了紧密的联系。因为如果按照他的说法，无论是在受生中还是在制造中，本质分离都不被观察到，那么，无论他允许主是什么，无论是\Quote{受造之物}还是\Quote{受生之物}，他都以这两个名称同样主张了本质的亲缘关系，因为他把主动与被动、生养者与受生者中的本性共通作为其体系的一部分。\par

现在让我们转向论证的下一步。我恳请读者不要对检查的细致感到不耐烦，因为这使我们的论证超出了我们期望的长度。因为我们所面临危险的并非寻常小事，以致如果我们匆忙掠过某个需要更仔细注意的点，损失会很小；我们赌上的恰恰是我们望德的总和。我们面临的选择是：是做基督徒，不被异端的毁灭性诡计引入歧途，还是被完全卷入犹太人或外邦人的观念中。因此，为了不使我们遭受这两样禁止之事，既不至于因否认真实受生之子而同意犹太人的教义，也不至于因崇拜受造物而陷入偶像崇拜者的堕落，让我们勉强花些时间讨论这些问题，并陈述欧诺弥的原话如下：—\par

\Quote{既然这些事物这样被划分，人们可以合理地说，最合适、最原初的本质，以及那唯独藉圣父运作而存在的本质，为自己接受了\InnerQuote{受生之物}、\InnerQuote{制作之物}和\InnerQuote{受造之物}的名称：}紧接着他又说：\Quote{但圣子独自藉圣父的运作而存在，拥有祂的本性以及与生养祂者的关系，无共通。}这就是他的话。但是，让我们如同观看敌人内部发生派别斗争的人一样，首先考虑我们对手之间的自相矛盾，然后从另一边陈述敬虔的真实教义。\Quote{圣子独自，}他说，\Quote{藉圣父的运作而存在，拥有祂的本性以及与生养祂者的关系，无共通。}但在先前的陈述中，他说他\Quote{不拒绝称那受生者为\InnerQuote{受生之物}，因为受生的本质本身以及圣子的称号使这种词语的关系成为适当。}\par

这些段落之间明显的矛盾如此昭然，我不禁钦佩那些赞美这一学说的人的敏锐。因为很难说他们能转向他哪一个陈述而不发现与其他陈述相矛盾。他早先的陈述表明，受生的本质和\Quote{圣子}的称谓使得这种关系措辞恰当。他现在的体系则相反：\Quote{圣子与生祂者的关系没有共通性。}如果他们相信第一个陈述，他们肯定不接受第二个；如果他们倾向于后者，他们会发现自己与早先的观念对立。谁能制止这场战斗？谁能在这内战中调解？谁能将这不和化为一致，当灵魂本身因对立的陈述而分裂，被引向相反学说的不同方向？或许我们在此看到了\Person{达味}论及犹太人的那句晦涩预言——\Quote{他们分裂了，却没有刺心。}因为看，即使他们在学说的对立中分裂，也感觉不到自己的分歧，而是像酒罐一样被耳朵牵着走，随着操纵他们者的意愿被带到各处。他乐于说受生的本质与\Quote{圣子}的称谓紧密相连：他们立刻像睡着的人一样点头同意他的话。他又改变说法，转向相反的主张，否认圣子与生祂者的关系：他亲爱的朋友们再次对此也表示赞同，随他的意愿转向任何方向，如同身体的影子随着前进之人的运动以自发的模仿改变形态，即使他自相矛盾，他们也接受这一点。这是\Person{荷马}告诉我们的另一种形式的魔饮，不是将饮下毒药者的身体变成兽形，而是作用于他们的灵魂，使之变成无理性的状态。因为关于那些人，传说他们的心智健全，但形态变成了野兽；而在这里，他们的身体保持自然状态，灵魂却变成了野兽的状态。正如诗中奇事所说，那些饮下药水的人被诱惑他们本性的人随意变成了各种野兽的形态，现在从这位\Person{刻耳刻}的酒杯中也发生同样的事。因为那些从同一著作中饮下巫术欺骗的人，变成了不同学说的形态，一时变成一种，一时变成另一种。与此同时，这些非常可笑的人，按照这寓言的修订版，仍然对引领他们走向如此荒谬的人感到满意，俯身捡拾他散播的话语，如同捡拾山茱萸果实或橡子，像猪一样贪婪地奔向洒在地上的学说，天性上无法注视那些崇高而属天的事物。因此，他们看不到他的论证倾向于相反的立场，而是不加检验地抓住遇到的东西：正如他们所说，被曼德拉草麻痹的人的身体处于某种昏睡和无法移动的状态，这些人的灵魂的感官也是如此，因对欺骗的理解而变得迟钝。确实，因隐藏的诡诈而陷入无意识，是可怕的事；然而，如果是无意的，不幸尚可原谅；但作为某种预谋和热切渴望的结果，而非不知后果，去尝试邪恶，则超过了所有极度的悲惨。我们当然可以抱怨，当我们听说甚至贪婪的鱼也会在无饵的钩子靠近时避开，只有在获得食物的希望诱骗它们上钩时才吞下钩子；但是，当邪恶明显时，自愿走向这种毁灭，比鱼的愚昧更可悲：因为那些鱼是被贪婪引向隐藏的毁灭，而这些人在不敬神的赤裸钩子前张开嘴吞下，因某种非理性的激情而满足于毁灭。还有什么比这矛盾更清楚——说同一位既是受生又是受造物，并且某物与\Quote{圣子}的名字紧密相连，又远离\Quote{圣子}的含义？但这些事已经够了。\par

% structure: p0024 source=h2 target=subsection reason=relative-heading-rank; base=h2

\subsection{五、他再次显示，\Person{欧诺弥}被真理所迫，成为正统学说的代言人，宣认不仅圣父的本质是最恰当和首要的，而且独生子的本质也是如此。}

然而，或许有益的是，查看\Person{欧诺弥}的陈述中按顺序呈现在我们面前的这段话的含义，回到他陈述的开头。因为我们刚才考察的点显然是促使我们从最后一段开始回答，因为他的言辞中明显的矛盾。\par

这就是\Person{欧诺弥}在开头所说的：——\par

\Quote{既然事物如此区分，有人或许合理地说，那\OriginalTerm{至为本体和元始的本质}{most proper essence}，以及唯独因圣父的运作而存在的本质，可为自己接受 \InnerQuote{\OriginalTerm{受生之物}{product of generation}}、\InnerQuote{\OriginalTerm{受造之物}{product of making}} 和 \InnerQuote{\OriginalTerm{受创之物}{product of creation}} 的称号。}首先，我请求这些正留意此论述的人记住，在他的第一部著作中，他说圣父的本质也是\InnerQuote{至为本体的}，并用以下这些话引出他的陈述：\InnerQuote{我们教导的整体论述，是藉着那至高和至为本体的本质而完成的。}此处，他称\OriginalTerm{独生子}{Only-begotten}的本质为\InnerQuote{至为本体和元始的}。因此，若将欧诺弥从各卷书中的措辞放在一起，我们可把他自己召来作为\OriginalTerm{本质的共通}{community of essence}的见证，他在另一处作出这样的宣告：\InnerQuote{凡具有相同称号的事物，其本性也毫无区别。}因为这位自相矛盾的朋友不会用相同的称号来指示本性不同的事物；相反，正因圣父和圣子中本质的定义是同一的，所以他说前者是\InnerQuote{至为本体的}，后者也是\InnerQuote{至为本体的}。而人类的普遍用法也支持我们的论点：当名称与本性不真正相符时，就不使用\InnerQuote{至为本体的}这一称号。例如，我们不太准确地称一幅肖像为\InnerQuote{人}，但我们用此名称合宜地指称的是自然呈现于我们面前的动物。同样，圣经的语言也承认\InnerQuote{神}的称号可用于偶像、魔鬼和肚腹，但在此名称也没有其\OriginalTerm{合宜的含义}{proper sense}。所有其他情况也如此。一个人据说在梦中幻想象中吃了食物，但我们不能称这种幻象食物为合宜意义下的食物。因此，正如在两个自然存在的人中，我们合宜地同样称他们两个为\InnerQuote{人}；但如果有人将一尊无生命的肖像与一个真实的人并列计数，他或许会称那真实存在的人和那肖像为\InnerQuote{两个人}，但再也不会将这个词的合宜含义归于二者。同样，若设想独生子的本性是圣父本质之外的某种东西，欧诺弥就不会称两个本质均为\InnerQuote{至为本体的}。因为一个人怎能用相同的名称来指示本性不同的事物呢？诚然，真理甚至藉着那些反对它的人而显明，因为谬误即使以敌人的言辞表达，也无法完全胜过真理。因此，正统信仰的教义藉着对立者的口得以宣告，他们自己不知道在说什么，正如\Person{盖法}预言主为我们所作的救赎性受难，却不知道自己在说什么。所以，如果真正的本质本体性是二者（我指圣父和圣子）所共有的，那么还有什么空间去说它们的本质彼此分歧？或者，在它们中又如何能看到因优越能力、伟大或荣耀而来的差异呢？既然\InnerQuote{至为本体的}本质不容任何缩减？因为那不完全地是其所是的东西，无论它是本性、能力、位阶或任何其他个别可思考的对象，都不能被\InnerQuote{最合宜地}称为那物，因此异端所主张的圣父本质的超越性证明了圣子本质的不完全。如果它不完全，它就不是合宜的；但如果它是\InnerQuote{至为本体的}，它也一定是完全的。因为不可能把有欠缺的称为完全的。而在比较时，既不可能将完全的事物与完全的事物并列而觉察到任何过与不足的差异：因为完全在二者中是同一的，正如一把尺规，既不是因不足而凹陷，也不是因过剩而凸起。因此，从这些段落中可以充分看到欧诺弥对我们教义的辩护——我宁愿说，不是他为我们热心，而是他与自己冲突。因为他用那些他用来建立我们教义的论证来攻击自己。然而，让我们再次逐字逐句地跟随他的作品，使所有人都能明白，他们的论证除了作恶的欲望之外，没有能力造成邪恶。\par

% structure: p0028 source=h2 target=subsection reason=relative-heading-rank; base=h2

\subsection{六、他接着揭露关于 \Quote{受生之物}、\Quote{受造之物} 和 \Quote{受创之物} 的论证，并揭示欧诺弥和特奥格诺斯图斯关于本质的 \Quote{直接} 和 \Quote{不可分} 特性及其 \Quote{与其创造者和造物主的关系} 的言辞之不敬本质。}

那么，让我们听听他说了些什么。\Quote{有人可能会合理地说，最恰当和最首要的本质，即唯独藉圣父的运作而存在的本质，允许自身接受\InnerQuote{受生之物}、\InnerQuote{制造之物}和\InnerQuote{受造之物}的称谓。}谁不知道，使教会与异端分开的正是这个应用于圣子身上的术语\Quote{受造之物}呢？因此，既然教义上的差异是普遍承认的，那么一个试图证明他的观点比我们的更真实的人，应当采取什么合理的做法呢？显然，他应当通过尽其所能的论证，表明我们应当认为主是受造的，从而确立自己的陈述。或者，如果他忽略这一点，他是否应该为他的读者制定一条律法，让他们把有争议的事当作公认的事实来谈论呢？就我个人而言，我认为他应该采取前一种做法；或许所有稍有理智的人都会要求对手首先在某种无可辩驳的基础上确立他们论证的第一原则，然后再通过推理来推进他们的理论。而我们的作者却放弃了确立\Quote{我们应当认为祂是受造的}这一观点的任务，直接推进到下一步，将他论证的推理过程与这个未经证明的假设衔接起来，正好比那些心思深陷于愚妄欲望中的人，思想飘忽于一个王国或某种其他追求目标之上。他们不去想他们所渴望的事情如何可能实现，而是随心所欲地为自己安排和整理他们的好运，仿佛它已经是他们的了，带着一种愉悦在不存在的事物中游荡。同样，我们这位聪明的作者不知何故将自己那著名的辩证法催眠了，在尚未对争议点给出论证之前，就像对孩子们讲故事一样，讲述他自己教义中这个欺骗性且前后不一的蠢话，就像在酒席上讲述一个故事一样。因为他说，那\Quote{藉圣父的运作而存在}的本质允许接受\Quote{受生之物}、\Quote{制造之物}和\Quote{受造之物}的称谓。有什么推理向我们表明，圣子是藉某种构造性的运作而存在，并且圣父的本性在圣子的位格存在方面保持不动呢？这正是争论中的焦点：圣父的本质是生了圣子，还是制造了祂，如同伴随祂本性的外在事物之一？如今，既然教会依照神圣教导，相信独生子确实是天主，憎恶多神论的迷信，并为此不承认本质的差异，以免诸神性因本质的不同而落入数的概念之下（因为这无异于将多神论引入我们的生活）——我说，既然教会用明白的语言教导：独生子在本质上是天主，从真实天主的本质而来的真实天主，那么，一个反对教会决议的人应当如何推翻先入之见呢？难道他不应该通过确立相反的陈述，从某个公认的原则来论证争议点吗？我想没有一个明智的人会期待别的做法。然而我们的作者却从争议点出发，把这些有争议的事情当作已被承认的，作为后续论证的原则。如果首先已经证明圣子是通过某种运作而获得存在的，那么他随后说那通过运作而存在的本质允许接受\Quote{制造之物}的名称，我们又能有什么异议呢？但是，让错误的倡导者告诉我们，只要前提尚未确立，结论又如何能有力量呢？因为假如有人以假设的方式承认人是有翅膀的，那么对于随后的事情就不会有任何让步的问题：因为那变得有翅膀的人会以某种方式飞翔，并高高升起于大地之上，用他的翅膀在空中翱翔。但我们必须看到，那本性不属于空气的人如何能变得有翅膀，而如果这个条件不存在，讨论下一步就是徒劳的。那么，让我们的作者首先证明这一点：教会相信独生子真实存在，并非被一个虚假称为父的收养，而是按照本性，藉来自那自有者的生成而存在，与生祂者的本质不相异，这是徒劳的。但是，只要他的首要命题仍未得到证明，停留在次要命题上就是无意义的。也不许任何人打断我，说我们所宣认的也应当通过构造性的推理来确认：因为证明我们的陈述足够了，即传统从我们的父辈传给我们，如同某种遗产，经由宗徒及后来的圣人们依次传递。另一方面，那些将自己的教义改为这种新奇的学说的人，如果想要说服的不是轻如尘土、摇摆不定的人，而是有分量、稳重的人，就需要大量论证的支持。但是，只要他们的陈述是在未经确立、未经证明的情况下提出的，谁还会如此愚蠢和粗野，以至于认为福音作者和宗徒以及那些在教会中相继如明光照耀之人的教导，还比不上这未经证明的胡言乱语呢？\par

让我们进一步看看我们作者聪明才智最突出的例子；他是如何凭借其丰富的辩证技巧，巧妙地将较为单纯的人引向相反的观点。他在\Quote{受造之物}和\Quote{被造之物}的称号之外，又增加了\Quote{受生者}这一说法，声称圣子的本质\Quote{为自身接纳这些名号}；并且认为，只要他像在酒徒聚会中那样滔滔不绝，他在处理教义方面的狡诈就不会被任何人察觉。因为将\Quote{受生者}与\Quote{受造之物}和\Quote{被造之物}放在一起，他以为通过把毫无共同之处的东西拼凑起来，就偷偷消除了这些名号在含义上的差别。这些就是他辩证法的巧妙把戏；但我们在论证中不过是些平信徒，并不否认就声音和舌头而言，我们正如他的演讲所描述的那样；但我们承认，正如先知所说，我们的耳朵也预备好了进行明智的聆听。因此，我们不会被那些毫无共同之处的名号的并列所动摇，以致混淆它们所表示的事物：即使伟大的宗徒把木材、干草、麦秸、金、银和宝石列在一起\BibleRef{格前 3:12}，我们也能大致数出他提到的事物的数量，却仍能不混淆地认出所列每种物质的性质。同样在这里，当\Quote{受生者}和\Quote{受造之物}被并列提起时，我们从声音转向意义，并不会在每个名号中看到相同的含义；因为\Quote{被造之物}意味着一回事，而\Quote{受生者}意味着另一回事：因此即使他试图混合不能混合的东西，明智的听众会加以分辨，并指出任何单一本质都不可能\Quote{为自身接纳}\Quote{受生者}和\Quote{被造之物}的称号。因为，如果其中一个是真实的，另一个就必然是虚假的；因此，如果某物是被造的，它就不是受生的；反之，如果它被称为受生者，它就与\Quote{被造之物}的称号无关。然而，欧诺弥却告诉我们，圣子的本质\Quote{为自身接纳\InnerQuote{受生者}、\InnerQuote{受造之物}和\InnerQuote{被造之物}的称号}！\par

他是否通过剩余的内容，使他这个既无头又无根的陈述更加稳固？这个陈述在其最初阶段，关于他要建立的观点，没有提出任何有说服力的内容。还是剩余部分也紧抓着同样的愚蠢，不是从论证中获得力量，而是像讲述梦境一样用模糊的细节来展开亵渎的论述？他说——（这是他附加在我已经引用的话之后的）——\Quote{其受生没有中介，并保持与它的生者、制造者和创造者的不可分割的关系。}如果我们暂时不考虑\Quote{没有中介}和\Quote{不可分割}，只看这些词语本身的意义，我们会发现，他那荒谬的教导在每一处都灌输给被他欺骗的人，却没有一个论证的支持。他说：\Quote{它的生者、制造者和创造者。}这些名称虽然看起来是三个，但只包含两个概念的意义，因为其中两个词含义相同。因为制造与创造是相同的，但受生与所说的那些是不同的。既然这些词语的意义结果是将人的普通理解划分为不同的概念，那么有什么论证向我们证明制造与受生是相同的事，以便我们能够将这唯一的本质适应于这些术语的差异呢？因为只要这些词语的普通含义仍然成立，并且没有找到任何论证将术语的意义转移到相反的含义，那么任何单一本质就不可能被分割为\Quote{受造之物}和\Quote{受生者}的概念。既然这些术语每个单独使用都有其自身的含义，我们也必须假定它们所处的相对联系对于这些术语是恰当且贴切的。因为所有其他相对术语的联系，不是与陌生的和异质的东西相联，而是即使对应的词被省略，我们也会在听到主要词的同时，自发地听到与之相连的词，例如\Quote{制造者}、\Quote{奴隶}、\Quote{朋友}、\Quote{儿子}等等。因为所有被视为与另一个词相关的名称，通过提及它们，向我们呈现了每个名称与其所宣告之物的恰当且紧密相连的关系，同时避免了一切异质之物的混合。因为\Quote{制造者}这个名称既不与\Quote{儿子}这个词相连，\Quote{奴隶}这个术语也不指向\Quote{制造者}，\Quote{朋友}也不向我们呈现\Quote{奴隶}，\Quote{儿子}也不呈现\Quote{主人}，而是我们清楚分明地认识到每一个与其相关者的联系，通过\Quote{朋友}这个词想到另一个朋友，通过\Quote{奴隶}想到主人，通过\Quote{制造者}想到工作，通过\Quote{儿子}想到父亲。同样地，\Quote{受生者}有其恰当的相关意义；与\Quote{受生者}相连的当然是\Emph{生者}，与\Quote{被造之物}相连的则是\Emph{创造者}；如果我们不想通过替换名称来引入事物的混乱，就必须为每个相对术语保留其恰当的内涵。\par

既然这些词语的意义倾向是明显的，那么一个借助逻辑体系来推进其理论的人，怎么未能察觉到这些名称中它们各自恰当的相关含义呢？但他认为，通过说圣子的本质\Quote{为自身接纳\InnerQuote{受生之物}、\InnerQuote{受造之物}和\InnerQuote{被造之物}这些称谓}，并且\Quote{保持与其生者、造者、创造者不可分割的关系}，他就把\Quote{受生之物}与\Quote{造者}联系起来，把\Quote{受造之物}与\Quote{生者}联系起来。因为单一的事物分裂成不同的关系，这是违反本性的。但圣子正是与圣父相关，那受生的与生他的相关，而\Quote{受造之物}则与它的\Quote{造者}相关；除非有人认为在日常用语中某种不精确的、不加区分的用法可以凌驾于严格的含义之上。\par

那么，根据他那无敌的逻辑，他是凭借什么推理和什么论点，来赢回大众的意见，并随心所欲地顺着这条思路推导出：既然统御万有的天主既被构想和称为\Quote{创造者}又被构想和称为\Quote{圣父}，那么圣子就与这两个称号紧密相连，同样既被称为\Quote{被造之物}又被称为\Quote{受生之物}呢？因为正如惯常精确的言谈区分这类名称，对于从本质本身生出的事物使用\Quote{受生}之名，而对于那些在其造物主本性之外的事物理解\Quote{被造}之名；又因为神的教义在传授天主的认识时，交付给我们\Quote{圣父}和\Quote{圣子}这些名称，而不是\Quote{创造者}和\Quote{作品}，以免产生倾向于亵渎的错误（如果后一种称谓将圣子归到外来者与陌生人的位置上，就可能发生），并使那些将独生子与圣父的本质联系切断的邪恶学说无从进入——看到这一切，我说，他宣称\Quote{受造之物}这样的称谓也适合于圣子，那将安全地推出结论说，\Quote{圣子}之名也恰当地适用于那受造之物；因此，如果圣子是\Quote{受造之物}，那么天也被称为\Quote{圣子}，而个别被造的事物，按照我们作者的观点，也恰当地用\Quote{圣子}这个称谓来命名。因为如果祂拥有此名，不是因为祂与生祂者共享本性，而是因为祂是被造的而被称为圣子，那么同样的论证将允许一只羊、一只狗、一只青蛙，以及一切因造物主意愿而存在的东西，都被冠以\Quote{圣子}的名号。反之，如果其中一个不是圣子，也不被称为天主，因为它在圣子的本性之外，那么，那位真正是圣子的，的确是圣子，并且因为祂出自生祂者的本性而被承认为天主。但欧诺米厌恶受生的概念，并将其排除在神学教义之外，用他肉性的思辨来诽谤这个术语。那么，我们在前面的论述中已经充分表明，正如圣咏作者所说：\Quote{他们在没有恐惧的地方感到恐惧。}因为如果在人身上已经表明，并非所有受生都存在于情欲的方式中，而是质料性的出于情欲，属神的则是纯洁不朽的（因为由圣神生的就是神，不是肉体，而在神中我们看不到任何受情欲支配的状况），既然我们的作者认为必须通过我们自身的例子来估计天主的能力，那么让他说服自己从另一种受生的模式来设想神性受生的无情感性。此外，他把这三个名称混在一起，其中两个是同义的，他认为他的读者会因这两个词组含义的共通性而推断第三个也是同义的。因为既然\Quote{受造之物}和\Quote{被造之物}这些称谓表明被造之物在造物主的本性之外，他就把\Quote{受生之物}这一词组与之并列，以便这一点也能与上述词组一起被理解。但这种论证被称作欺诈、谎言和欺骗，而不是深思熟虑、技巧熟练的证明。因为只有那种从公认的事物中揭示未知事物的才被称为证明；但欺诈地、谬误地推理，掩盖你自己的羞辱，并用肤浅的骗术混淆人的理解，如宗徒所说\Quote{心地败坏的人\BibleRef{弟后 3:8}}，没有一个理智的人会称此为技巧熟练的证明。\par

然而，让我们继续按顺序看接下来的内容。他说本质的受生是\Quote{无中介}，并且它\Quote{保持了它与祂的生成者、制造者、创造者的不可分割的关系}。如果他只是谈及本质的直接性和不可分性，并在此止步，那么他的言论并未偏离正统观点，因为我们也承认圣子与圣父的紧密联系和关系，以至于在他们之间没有任何插入物被视为中介，没有任何间隔的概念，即便是那细微而不可分的间隙——当时间被划分为过去、现在和未来时，现在本身被不可分地构想，因为它既不能被看作过去的一部分，也不能被看作未来的一部分，因为它完全无维度、不可分割、不可察觉，无论加到哪一边。因此，那完全直接的，我们说不允许任何这样的中介；因为任何被间隔分开的就会不再是直接的。所以，如果我们的作者也同样说圣子的受生是\Quote{无中介}，排除了所有这些观念，那么他就确立了关于那一位与圣父结合的正统教义。然而，仿佛后悔似的，他立刻又补充说，本质\Quote{保持了它与祂的生成者、制造者、创造者的关系}，从而用第二个陈述玷污了第一个，将亵渎的言论呕吐在纯洁的教义上。因为很明显，他那里的\Quote{无中介}并没有正统的意图，而是如同有人会说锤子是铁匠与钉子之间的中介，但锤子本身的制造是\Quote{无中介}的，因为当工具尚未被技艺发现时，锤子首先通过某种创造性的过程从工匠手中产生，并非借助任何其他工具，然后通过它制造了其他工具；同样，\Quote{无中介}这一短语表明这就是我们的作者关于独生子的概念。而且在这里，欧诺米并非唯一在教义的严重错误上犯错的人，你可以在德奥格诺斯图斯的著作中找到相似的观点，他说天主想要创造这个宇宙，首先创造了圣子作为创造的某种标准；他没有意识到他的陈述中包含了这种荒谬：即那为某物而存在、而非因自身而存在的东西，其价值当然低于它为之存在的那物：就像我们为了生计提供农具，但犁的价值当然不能与生命等同。所以，如果主也是因为世界而存在，而非万物因祂而存在，那么他们所说祂为之存在的万物整体，其价值将超过主。而这正是他们通过论证在此确立的，他们坚持说圣子与祂的创造者和制造者的关系是\Quote{无中介}的。\par

% structure: p0035 source=h2 target=subsection reason=relative-heading-rank; base=h2

\subsection{七、接着，他清晰而巧妙地批评了不能与圣子之后所造之物比较的教义，并揭露了欧诺米所捏造、并以\Quote{圣子}和\Quote{独生子}的术语掩盖以欺骗读者的偶像崇拜。}

在段落的剩余部分，他却变得安抚，说那本质 \Quote{不与任何由它并后于它造成的事物相比较。} 仇视真理者献给主这样的礼物，由此他们的亵渎更加明显。请告诉我，受造界中还有什么事物容许与另一不同事物相比较，既然在每一事物中显现的特有本性绝对排斥与不同类的事物共融？天不容与地相比较，地不容与星辰，星辰不容与海洋，水不容与石头，动物不容与树木，陆生动物不容与飞禽，四足兽不容与游泳生物，无理性者不容与有理性者。其实，何必逐一举例，表明我们可以对在受造界中看到的每一事物恰当地说出那仿佛作为特殊之物扔给独生子的话——即祂与任何由祂并后于祂产生的事物都不容比较？显然，你独自设想的每一事物都与宇宙及其组成的个体事物不可比较；而仇视真理者将这本可真正适用于任何受造物的话，当作足够并充分的尊荣和荣耀，分配给了独生子天主！接着，他在段落剩余部分组合了同类短语，以空洞的尊荣尊崇祂，称祂为 \Quote{主} 和 \Quote{独生子}：但这些名号并未向读者传达正统含义，他立即将亵渎与那些更显著的名号混合。他的表述如下：他说，\Quote{既然受生的本质不让其他任何事物有共融的余地（因为它是独生的），造物主的行动也不被视为共同的。} 啊，惊人的傲慢！仿佛在对野兽或无知之物讲话，如同那些 \Quote{没有理解} 之物，他任意以相反方式扭曲他的论证；或者，他像那些失明者一样受苦；他们常在看得见的人面前做出不得体的事，以为因为自己看不见，也就不会被看见。因为什么样的人看不出他话语中的矛盾？他说，因为它是 \Quote{受生的}，所以本质不让其他事物有共融的余地，因为它是独生的；然后他说出这些话后，真的仿佛他看不见或以为不被看见，他将与所说之事毫无共同之处的东西附加上去，将 \Quote{造物主的行动} 与独生子的本质连结起来。受生者与生养者是相互关联的，而独生子当然与圣父相关联；仰望真理的人看见的是与圣子相协调的，不是 \Quote{造物主的行动}，而是生养祂者的本性。但他，仿佛在谈论植物或种子，或受造秩序中的其他事物，将 \Quote{造物主的行动} 置于独生子存在之旁。为何？如果考虑的是石头、棍棒或类似之物，那么预先假定 \Quote{造物主的行动} 是合乎逻辑的；但如果连仇敌也承认独生子是天主，是圣子，并且是通过受生而存在，那么适用于受造界最低微部分的同一类话语如何能适用于祂？他们如何看待把确实可对蚂蚁或蚊虫说的话用来论及主，认为是虔敬的？因为如果有人了解蚂蚁的本性及其与其他生物的特殊关系，他说 \Quote{它的造物主的行动不被视为共同的} 相对于其他事物，这并不超出真理。因此，对这类事物所肯定的，他们也用来肯定独生子；正如猎人们据说用陷阱拦截猎物，并用一些不牢固且无实质的材料掩盖陷阱口，使坑看起来与周围地面齐平，异端同样设计对付人，用这些好听且虔敬的名号，如同用平坦的茅草遮盖他们不敬的陷阱口，让那些较不聪明的人，以为这些人的宣讲因用词一致而与真信仰相同，就朝着圣子、独生子的单纯名号奔去，踏入了坑中的空虚，因为这些称号的意义无法承受他们脚步的重量，将他们坠入否认基督的陷阱。这就是为什么他谈论受生的本质不让任何事物有共融的余地，并称之为 \Quote{独生子}。这些是陷阱的遮盖。但当任何人在落入深渊前停下，并像伸手一样将论辩的考验放在他的话语上时，他就会看到偶像崇拜的危险堕落隐藏在教义之下。因为当他走近，如同走向天主和天主子，却发现一个天主的受造物被设定为他的崇拜对象。这就是为什么他们高喊低呼叫独生子的名号，以使毁灭容易被欺骗的受害者接受，就像有人将毒药混入面包，向求食者致致命的问候，他们本来不会单独服毒药，若非被所见之物引诱。因此，他敏锐地注视着他努力的目标，至少就他自己的意见而言。因为如果他完全从教导中抛弃圣子的名号，当他的否认在明确的宣告中公开陈述时，他的谎言就不会被人接受；但现在只留下名号，并将其含义改为表示受造，他立刻树立了他的偶像崇拜，并欺诈地隐藏了其耻辱。但由于我们被命令不要用嘴唇尊敬天主，虔诚不是由词语的声音来检验，而必须先在心里相信圣子以至于成义，然后用口宣认以至于救恩，而那些心里说祂不是天主的人，即使口中宣认祂为主，也成了败坏和可憎的，如先知所说——因此，我们必须看那些假装提出信仰言辞之人的心意，不要被他们的声音所引诱。因此，如果谈论圣子的人并非用这个词指受造物，他就是站在我们这边，而非敌人那边；但如果有人将圣子的名号应用于受造物，他就应被归为偶像崇拜者。因为他们也曾将天主的名号给予达贡、贝尔和龙，但他们并非因此崇拜天主。因为木头、铜和怪物不是天主。\par

% structure: p0037 source=h2 target=subsection reason=relative-heading-rank; base=h2

\subsection{八、他继而指出，圣父与圣子的本质并没有\Quote{差异}；其中他阐释了差异与和谐的多种形式，并解释了\Quote{形式}、\Quote{印记}与\Quote{确切形象}。}

但在我们的论述中，有何必要仅凭猜测他的意图来揭露他隐藏的诡诈，并可能给听众以反对的理由，认为我们是不实地指控我们的敌人？因为看哪，他赤裸裸地向我们展示了他的亵渎，没有用任何面纱遮掩他的诡计，而是无所顾忌地以放肆的声音谈论他的荒谬。他所写的如下：——他说：\Quote{我们这一方，既然找不到圣子本质之外任何可容受生的东西，认为我们必须将这些称谓归于本质本身，不然，如果我们真的要将它们与本质分离，那么谈论\InnerQuote{圣子}和\InnerQuote{受生的}便是徒劳，不过是空话；从这些名称出发，我们也确信本质彼此有差异。}\par

我想，无需我们提出论据来驳斥这里所陈述的荒谬。仅阅读他所写的就足以暴露他的亵渎。但让我们这样检验它。他说圣父与圣子的本质是\Quote{有差异的}。什么是\Quote{有差异的}？首先让我们考察这个术语单独使用时的含义，以便通过这个词的解释，更清楚地揭示其亵渎性质。\Quote{差异}一词，在习俗许可的不精确意义上，用于物体，当因瘫痪或任何其他疾病，某个肢体偏离其自然的协调时。因为我们将痛苦状态与健康状态比较，称一个经历恶化变化者的状况为从其通常健康的\Quote{差异}；在那些在德行与恶习方面不同的人中，将放荡的生活与纯洁和节制的生活比较，或不义的生活与正义的生活比较，或激情、好战、易怒的生活与温和、平安的生活比较——一般来说，所有被指责为恶习的，与更卓越的比较，都被说成表现出与它的\Quote{差异}，因为在两者——我指善的与较差的——中观察到的标记互不一致。再者，我们称在元素中观察到的那些性质为\Quote{有差异的}，它们是作为相反者彼此对立，具有相互破坏的力量，如热与冷，或干与湿，或一般说来，任何作为相反者与另一者对立的；这些缺乏统一我们用\Quote{差异}一词来表达；并且一般来说，任何在可观察特征上与另一者不相谐和的，都被说成与它\Quote{有差异}，如健康与疾病，生命与死亡，战争与和平，德行与恶习，以及所有类似的情况。\par

既然我们已经这样分析了这些表述，也让我们考虑一下，就我们的作者而言，他在什么意义上说圣父和圣子的本质是\Quote{彼此差异的}。他是什么意思？是指圣父是依照本性，而圣子\Quote{偏离}了那本性？还是他用这个词表达德行的扭曲，用\Quote{变异}之名将邪恶与更卓越者分开，从而视一个本质为善，另一个为相反的方面？还是他断言一个神圣本质与另一个也是差异的，如同元素对立的方式？或者，如同战争之于和平、生命之于死亡，他是否也在本质中看到这种存在于所有此类事物中的冲突，以致于它们无法彼此合一，因为相反者的混合对混合之物施加一种消耗的力量，正如\WorkTitle{箴言}的智慧论及这样的道理，说水和火从不言\Quote{够了}，隐秘地表达了势均力敌、平衡对立的相反者及其相互毁灭的本性？或者，他是否以这些方式中的任何一种来看待本质中的\Quote{差异}？那么，让他告诉我们，除了这些他还想到了什么。我认为，即使他重复他惯常的说法\Quote{圣子与生育祂的那位是差异的}，他也无法说出别的；因为他陈述的荒谬因此而更加清楚地显明。因为什么相互的关系能像\Quote{圣子}这个词所表达的与圣父的关系那样紧密而和谐地契合和连结在一起？一个证明是：即使这两个名字都没有被说出，那被省略的也被说出的那一个所暗示，如此紧密地一个隐含在另一个之中，并与之和谐；两者都在一个中被识别，以致一个不能脱离另一个而被设想。现在，那\Quote{差异的}肯定是在与那\Quote{和谐的}相对的意义上被如此设想和称呼的，正如铅垂线与直线和谐，而弯曲的东西放在直的东西旁边就不与之和谐。音乐家们也习惯称音符的一致为\Quote{和谐}，而把走调和刺耳的称为\Quote{不和谐}。因此，说事物是\Quote{差异的}等于说它们是\Quote{不和谐的}。所以，如果独生天主的本性与圣父的本质是\Quote{差异的}（用异端的话说），它当然不与之和谐：而不和谐之处不可能存在。因为情况正如：当蜡和印模的图案是一个时，被印模压过的蜡，当它再次与印模贴合时，使它自己上面的印痕与周围的部分一致，填充凹陷并使雕刻的凸起与它自己的纹路相适应；但如果某种奇怪而不同的图案贴合到印模的雕刻上，它会使它自己的形状粗糙而混乱，因为在不相应的雕刻面上蹭掉了自己的图形。然而，祂\BibleQuote{具有天主的形像}\BibleRef{斐 2:6}，不是由与圣父不同的任何印模所形成，因为祂是圣父位格的\Quote{确切印像}\BibleRef{希 1:3}，而\Quote{天主的形像}当然与祂的本质是同一回事。因为，正如\BibleQuote{成为仆人的形像}\BibleRef{斐 2:7}，祂是在仆人的本质中形成的，不仅仅取了一个形像而与本质分离，而是本质包含在\Quote{形像}的意义中；同样，那说祂\Quote{具有天主的形像}的人，由\Quote{形像}指明本质。所以，如果祂\Quote{具有天主的形像}，并且子在圣父内，被圣父的光荣所印证（正如福音的话所宣告的，它说\Quote{天主圣父印证了祂}——由此也\Quote{看见我的就是看见了父}），那么\Quote{美善的肖像}和\Quote{光荣的光辉}以及所有其他类似的称呼，都证明圣子的本质不与圣父不和谐。因此，所引用的经文显示了对手亵渎的空洞性。因为如果\Quote{差异的}事物是不和谐的，并且那被圣父印证、在自己的内显示圣父、既在圣父内又有圣父在自己内的祂，在各方面都显示祂的亲密关系和和谐，那么反对观点的荒谬性就由此被压倒性地显明了。因为正如\Quote{差异的}被显示为不和谐的，相反地，那和谐的当然没有争议地被承认不是\Quote{差异的}。因为正如\Quote{差异的}是不和谐的，那和谐的也不是\Quote{差异的}。此外，那说独生者的本性与圣父的善本质是\Quote{差异的}的人，显然着眼于善本身中的变异。但至于那与善差异的是什么——\Quote{愚蒙人啊，}正如\WorkTitle{箴言}所说，\Quote{要明白他的诡诈！}\par

% structure: p0041 source=h2 target=subsection reason=relative-heading-rank; base=h2

\subsection{九、接着，在区分本质与受生之后，他宣告尤诺米乌斯空洞轻浮的语言如同拨浪鼓。他进而指出，伟大的巴西略关于独生者受生问题所用的语言遭到了尤诺米乌斯的严重诽谤，并以此结束本书。}

然而，我将略过这些事，因为其中涉及的荒谬是明显的；让我们检查之前的内容。他说，什么也找不到，\Quote{除了圣子的本质，即容许受生的本质。}他说这话是什么意思？他区分了两个名称，并通过他的论述将它们所意指的事物分开，将每一个单独地分别开来。\Quote{受生}是一个名称，而\Quote{本质}是另一个。他告诉我们，本质\Quote{容许受生，}因此当然是与受生不同的事物。因为如果受生就是本质（这正是他不断声称的），以致两个称谓在意义上等同，他就不会说本质\Quote{容许受生}了：因为那等于说本质容许本质，或受生容许受生——如果受生与本质是同一回事的话。那么，他理解受生是一回事，而本质是另一回事，本质\Quote{容许受生：}因为被容纳的东西不能与容纳它的东西相同。好了，这就是我们这位贤哲有条理的论述所说的；但他的话语是否有意义，让精通判断的人去考虑吧。我将复述他的原话。\par

他说，他发现\Quote{除了圣子的本质，即容许受生的本质之外，什么也没有}；然而他的话没有意义，这对每一个听到他陈述的人来说都是显而易见的：剩下的任务似乎是揭示他试图借助这些无意义的词语构建的亵渎。因为他希望，即使他无法达到目的，也要通过这种模糊的表达方式，在听众心中产生一种想法，即圣子的本质是构造的结果：但他把这种构造称为\Quote{受生，}用最美好的词语装点他那可怕的亵渎，以便如果\Quote{构造}是\Quote{受生}一词所传达的含义，那么主被造的观念就可以轻易得到赞同。于是，他说，本质\Quote{容许受生，}这样每一种构造都可以被视为在某种基质中。因为没有人会说那不存在的东西是被构造的，从而将\Quote{制作}在他的论述中延伸为某种构造的作品，应用于独生者天主的本性。\Quote{那么，}他说，\Quote{如果它容许这种受生，}——他想表达的意思大约是这样：它本来不会存在，如果没有被构造的话。但在受造物中我们所领会的事物中，还有什么东西是\Emph{存在}而不被造的呢？天、地、空气、海，以及所有存在的事物，当然都是由于被造而存在。那么，他何以认为独生者的本性有一个特性，即\Quote{容许受生进入其实际本质}（因为他用这个名称来称呼制作），就好像大黄蜂或蚊子不容许受生进入自身，而是进入自身之外的某种东西？因此，他的著作承认，通过它们，独生者的本质被置于与受造物中最微小的部分同等的地位：而他试图证明圣子与圣父分离的所有证据，对于个体事物也同样有效。那么，他何必需要这种多样的敏锐来确立本性的差异呢？他本应采取否定的捷径，公开宣称不应宣认圣子的名字，也不应在教会中宣讲独生者天主，而应认为犹太人的崇拜优于基督徒的信德，并且在我们承认圣父是宇宙唯一的创造者和造物主的同时，将所有其他事物都归入受造物的名称和概念，并把这些事物中先于其余的那一个作品称为\Quote{被造之物，}它是通过某种构造活动而存在的，并称祂为\Quote{首造者，}而不是\Quote{独生者}和\Quote{真子。}因为当这些意见得胜时，就很容易使教义与它们所追求的目标达成一致，当所有人都被引导——正如你可以从这样一条原则中预料的——得出这样的结论：祂既非受生，也非儿子，而是通过某种能力而存在，不可能与天主共享本质。然而，只要福音的宣告仍然有效——福音宣告祂为\Quote{子}、\Quote{独生者}、\Quote{出自圣父}、\Quote{出自天主}等——欧诺弥将徒劳地胡说八道，用这种无聊的喋喋不休误导他自己和他的追随者。因为当\Quote{子}的头衔大声宣告与圣父的真实关系时，谁敢如此愚蠢，以至于当若望、保禄以及圣人合唱团中的其他人都宣告这些话——真理的话，指出密切亲属关系的话——时，他却不看他们，反而被欧诺弥诡辩的空话所引导，认为欧诺弥比那些藉圣神讲述奥秘\BibleRef{格前 14:2}，并在自己身上承载基督的人的教导更真实？唉，这个欧诺弥是谁？他从哪里被树立为基督徒的导师？\par

但让这一切都过去吧，让我们对眼前之事的热忱平息我们因维护信仰而对抗亵渎者所膨胀的心。因为当我们的天主、主、生命施予者和救主被这些可怜的人侮辱时，我们怎能不义愤填膺、心怀憎恨？若有人辱骂我按肉身之父，或与我的恩主为敌，我岂能无动于衷地忍受他对所爱之人的怒气？至于我灵魂的主，祂赐予不存在者存在，救赎被奴役者，让我尝到现世生命，又为我预备来世生命，召叫我们进入天国，赐下诫命以使我们逃脱地狱的永罚——我说的这些都是小事，不足以表达我们共同之主的伟大——祂受一切受造物朝拜，包括天上、地上和地下的，无数天朝供职者侍立在祂面前，一切在此统治下的和一切向往美善的都转向祂——若祂遭人辱骂，而那些人不仅自己加入叛教者的行列，还认为若不能借其笔墨将他人拖入同一深渊、使后来者不缺少引向毁灭的手便是损失，那么，谁还会责怪我们对这些人的愤怒？但让我们回到他论述的脉络。\par

他接着再次诽谤我们以人的类比侮辱圣子的受生，并提到我们教父就此所写的文字，其中说：虽然\Quote{子}一词可指两件事，即由情欲所形成的和与生养者的真实关系，但在讨论神性事务的言论中，他不接受前一种不雅而肉性的意义；而就后者能证明独生者的光荣而言，唯有这一点才适用于崇高的教义。那么，究竟是谁以人的观念侮辱圣子的受生？是将情欲和人性之物远远排除于神圣受生之外，并将圣子无痛苦地与生养祂者结合的人？还是那将创造万有者置于与下界受造物同等地位的人？然而，这种观念——将圣子置于圣父的尊威之中——似乎被这种新智慧视为侮辱；相反，它却认为将祂降至与我们一同被奴役的受造物同等是伟大而崇高的。空洞的抱怨！\Person{巴西略}被诽谤为侮辱圣子，而他却尊敬圣子如同尊敬圣父；\Person{欧诺弥}成了独生者的捍卫者，却将祂与圣父的善性分离！这种指责\Person{保禄}也曾一度在雅典人那里承受过，他们指控他为\BibleQuote{传讲外邦鬼神}\BibleRef{宗 17:18}，而\Person{保禄}正在斥责那些在偶像崇拜中迷乱于诸神的人，并引导他们走向真理，宣讲藉圣子而来的复活。如今这些指控被新的斯多葛派和伊壁鸠鲁派加在\Person{保禄}的追随者身上，他们正如史书论雅典人所说：\BibleQuote{不愿别事，只将新闻说说听听}\BibleRef{宗 17:21}。因为还有什么比这更新奇呢——能力之子、受造物之父、从无中生出的新神、善与善相争？这些人正是那些声称以恰当尊荣尊敬祂，却否认祂与生养祂者的本性相同的人。\Person{欧诺弥}难道不羞愧于这种尊荣的形式？若有人说他自己与本性的父亲并非同族，却与其他异类有相通，他作何感想？若将创造之主与受造物等同的人宣称如此是尊敬祂，那么让他也接受与无理性的畜类等同的尊荣吧。但若他认为与低劣本性等同是难堪和冒犯，那么将那位如先知所说\Quote{以大能永远统治}者置于臣服和奴役的本性之下，又怎能算是尊荣？但就此止住。\par

\SourceRecord{Translated by H.A. Wilson.；Nicene and Post-Nicene Fathers, Second Series；Vol. 5.；Edited by Philip Schaff and Henry Wace.；Buffalo, NY: Christian Literature Publishing Co.,；1893.；<http://www.newadvent.org/fathers/290104.htm>.}

% structure: p0001 source=h1 target=section reason=page-title

\section{\WorkTitle{驳斥欧诺弥（第五卷）}}

% structure: p0002 source=h2 target=subsection reason=relative-heading-rank; base=h2

\subsection{一、第五卷承诺讨论宗徒伯多禄所说的话语，但将其解释推迟。他首先论述受造界，说明其中没有任何值得崇拜之物，然而人因着无知和软弱的理智而误入歧途，惊叹于其美丽，将宇宙的各部分神化。在此他精彩地解释了依撒意亚的段落：\Quote{我是天主，最初的一位。}}

现在或许是时候探究欧诺弥本人关于宗徒伯多禄的话语所说的内容，以及我们的教父关于后者所说的内容。如果详细的考察使我们的讨论扩展到相当长的篇幅，公正的读者无疑会原谅这一点，并不会责怪我们浪费时间在言语上，而是将责任归于那个引起这些言语的人。请允许我在预定的探讨之前也做一些简短的评论：也许这些评论也会被发现与我们讨论的目标并不相违。\par

任何受造物都不值得人的崇拜，天主圣言如此明确地颁布了法律，以至于几乎整个受默感的圣经都可以学到这样的真理。梅瑟、约版、法律、随后的先知、福音、宗徒的谕令，都同样禁止崇拜受造物。详细列举与此相关的具体段落是一件冗长的工作；但尽管我们只从许多受默感的见证中举出少数几个，我们的论证无疑同样具有说服力，因为每一句天主的话，尽管是最小的，都有同样的力量来宣告真理。有鉴于此，我们对于存在的观念分为两种，即受造物和非受造的本性，如果我们的对手目前的争论占了上风，以至于我们应该说天主子是受造的，那么我们将绝对被迫要么藐视福音的宣告，拒绝崇拜那在起初就与天主同在的天主圣言，理由是我们不可崇拜受造物，要么，如果福音中记载的这些奇迹对我们来说过于紧迫，通过这些奇迹我们被引导去崇敬和崇拜在其中显示的那一位，在这种情况下，我们将使受造的和非受造的处于同样的尊荣等级；因为，按照我们对手的意见，即使受造的天主也被崇拜，尽管他在本性上并不比其他受造物优越，如果这种观点占了上风，宗教的教义将完全转变为一种无政府状态和民主式的独立。因为当人们相信他们所崇拜的本性并非一个，而是他们的思想被转向各种不同的神性时，将没有人阻止神性的概念在受造物中前进，而一旦在受造物中识别出神性，就会成为下一个被默观的事物中类似概念的垫脚石，然后又为下一个秩序，作为这种推断过程的结果，错误将扩展到所有事物，因为第一个欺骗通过相邻的案例一直传播到最后。\par

为表明我并非随意提出超越或然性之论断，我将引证外邦人中仍然流行的谬误，作为有利于我主张的可信证言。因为他们未经训练而狭隘的理智，倾向于惊奇地注视自然之美，却不以所见之物为引导和向导，去趋向那超越它们的本性之美；反而在抵达其领悟的对象时，使自己的理智停滞不前，逐一惊叹创造的各个部分——为此，他们未能将天主的观念停留在所见任何单一之物上，而认为他们在创造中所看到的一切皆为神圣。于是，在埃及人中间，由于这谬误在理智对象方面更强有力地发展，无数精神体的形态被视为同样多的神性本质；而在巴比伦人中间，穹苍那无误差的运转被视为一位天主，他们甚至给祂起名为贝耳。同样，外邦人的愚昧将七个连续的天体分别神化，一人向这一个膜拜，另一人向那一个膜拜，各依某种独特的谬误形式。因为他们觉察所有这些天体彼此相联地运行，又见他们已在高天之事上走迷，便依逻辑次序保持同样的谬误，直至最后者。此外，以太本身，及其下散布的大气，大地、海洋与地下区域，以及在大地之中一切有益或必需于人类生命之物——所有这些，他们没有一样认为与神圣本性无分无关，反而向每一样膜拜，借着创造中可见的某些对象，使自己依次受制于创造的各个部分；假若从一开始，连他们也知道敬畏创造是显然不合法的事，他们就不会被引入这多神论的欺骗之中。那么，让我们省察，免得我们也遭遇同样的命运——我们这些受圣经教导，敬畏真天主的人，曾被训练视所有受造存在为外在于神圣本性，并仅崇拜与敬拜那未受造的本性，其特性与标志是它既无开始，也永不终止；因为伟大的依撒意亚关于这些教义，以其崇高的言论述及神圣本性时，这样以天主位格说话：\BibleQuote{我是元始，以后仍然是我；在我以前没有天主，在我以后也没有天主。}因为比众人更完美地知晓福音信仰的奥秘，这位伟大的先知，他甚至预言了关于童贞女的那奇妙征兆，给我们传报了婴孩诞生的喜讯，并清晰地为我们指明了圣子的那一个名——总之，他藉圣神将一切真理包含于自身——为使神圣本性的特性，即我们藉以辨识真实\Emph{是}者与那成为之者的区别，尽可能清晰地为人所知，他遂以天主的位格发出此言：\BibleQuote{我是元始，以后仍然是我；在我以前没有天主，在我以后也没有天主。}既然那在天主之前的不是天主，那在天主之后的也不是天主（因为在天主之后的是受造物，在天主之前的是虚无，而虚无不是天主——或更确切地说，那在天主之前的是天主在其永恒福乐中，与虚无相对而定义），——既然，我说，这句受默感的话由先知的口说出，我们便藉他学到了这一教义：神圣本性是唯一的，与自身连续且不可分割，在自身中不接纳先后，尽管它被宣告于天主圣三中，且我们所默观的每一事物彼此之间没有更古或更新。既然这话是天主的话，无论你承认这话是圣父的还是圣子的，正统教义同样由两者所维护。因为如果是圣父如此说话，祂便为圣子作证，证明祂不\BibleQuote{在圣父之后}：因为如果圣子是天主，而凡是\BibleQuote{在圣父之后}的不是天主，显然这话就为真理作证，即圣子在圣父之内，而非在圣父之后。另一方面，如果有人承认这言说是圣子的，那么\BibleQuote{在我以前没有一位}这句话，便是清晰的暗示：我们默观\BibleQuote{在元始}之中的那位，是与元始的永恒一同被理解的。因此，如果任何事物\BibleQuote{在天主之后}，由所引经文可发现那是受造物，而非天主；因为祂说：\BibleQuote{在我之后的不是天主。}\par

% structure: p0006 source=h2 target=subsection reason=relative-heading-rank; base=h2

\subsection{二、然后他解释圣伯多禄的话：\BibleQuote{天主已立祂为主，为基督。}在此他陈述欧诺弥对圣巴西略因该话所作出的反对陈述，及其暗藏的谩骂与侮辱。}

既然我们已经获得了对存有的这一初步看法，现在正宜审查面前的这段经文。于是，伯多禄对犹太人说：\BibleQuote{祂，天主立为主及基督，就是你们钉死的这位耶稣，\BibleRef{宗 2:36}}，而我们方面则说，将\Quote{受造}一词归于独生子的天主性是不虔诚的，它应归于那\BibleQuote{奴隶的形状，\BibleRef{斐 2:7}}，这形状是因降生成人、在祂显现于肉身的适当时候而形成的；另一方面，那些将这句话反其道而行之的人则说，宗徒在\Quote{受造}这个词里指的是圣子先在的受生。因此，我们将把这段经文摆在中位，在对两种假设进行详细审查之后，将真理的判断留给读者。至于我们对手的观点，欧诺米本人足可作为代言人，因为他在此事上勇敢地辩论，因此我们逐字考察他的论点，就能完全跟随那些与我们敌对者的推理：而我们自己则将尽力充当己方教义的捍卫者，尽可能追随大巴西略先前阐述的论证路线。但你，即因你的阅读而在此案中担任法官者，要\BibleQuote{执行真实的判断，}如一位先知\BibleRef{匝 7:9}所说，不将胜利授予好辩的成见，而是授予通过检验而显明的真理。现在，让我们的教义的控告者上前，像在法庭上一样宣读他的诉状。\par

此外，除了我们已提到的，由于他拒绝将\Quote{受造}一词理解为指圣子的本质，同时又因他对十字架感到羞耻，他便将连那些最努力以愚蠢之名诽谤宗徒的人也无人归咎于他们的事，归咎于宗徒；同时，他借着其教义与论证，清楚地引入了两个基督和两个主；因为他说，天主没有使那在起初就存在的圣言成为主和基督，而是使那\BibleQuote{空虚自己，取了奴隶的形状}，并\BibleQuote{因软弱而受钉\BibleRef{格后 13:4}}的祂成为了主和基督。无论如何，大巴西略明写如下：\Quote{而且，宗徒的意图也不是要向我们展示那在万世之前的独生子的存在（这正是我们现在议论的主题），因为他清楚谈论的，不是那在起初就与天主同在的天主圣言的本质，而是那位\BibleQuote{空虚自己，取了奴隶的形状}、并\BibleQuote{成为与我们卑微的身体相似的\BibleRef{斐 3:21}}，且因软弱而受钉的祂。}又说：\Quote{即使稍微留意宗徒话语含义的人都知道，他不是在向我们陈述天主存在的方式，而是在引入属于降生成人的术语；因为他说：\InnerQuote{\Emph{祂}，天主立为主及基督，就是你们钉死的这位耶稣}，显然借着指示词强调祂内那属于人性且为众人所见的方面。}\par

\Quote{那么，这就是那人的话，他曾换——因我们不可称之为\InnerQuote{应用}，以免有人因这样的疯狂而责备那些为宣讲虔敬而蒙拣选的圣人，以致指责他们的道理陷入如此极端——他曾以自己心智取代宗徒们的意向！他们将那些无聊之言归于圣人们的记忆，是何等混乱啊！他们幻想那人\InnerQuote{贬抑自己}成为人，并宣称那位因服从而\InnerQuote{卑抑自己}取了奴仆形像的，甚至在取那形像之前就已与人相似，这是何等荒谬啊！请问，你这最冒失的人，当他已有奴仆形像时，如何还能取奴仆形像？一个人又怎能\InnerQuote{贬抑自己}而成为他本来就是的东西？你找不出办法应对这一点，尽管你大胆说出或想出那些无法设想的事。你岂不是所有人中最可怜的吗？你竟以为一个人为所有人死了，并将你的救赎归于他？因为如果真福伯多禄所说的不是那起初就与天主同在的圣言，而是那\InnerQuote{被看见的}，并且如巴西略所说，那\InnerQuote{贬抑了自己}的，如果那被看见的人\InnerQuote{贬抑自己}以取\InnerQuote{奴仆的形像}，而那位\InnerQuote{贬抑自己}以取\InnerQuote{奴仆的形像}的，为了成为人而贬抑了自己，那么那被看见的人为了成为人而贬抑了自己。事物的本性本身与此相悖；而且那位在论及天主性时颂扬这救恩计划的人明确驳斥了这一点，他说不是被看见的人，而是那起初就与天主同在的圣言取了肉身，这相当于取了\InnerQuote{奴仆的形像}。那么，如果你认为这些事应当相信，就离开你的错误，不再相信那人\InnerQuote{贬抑自己}成为人。如果你不能说服那些不愿被说服的人，就用另一句话、另一个反对他们的裁决摧毁他们的不信。请记住那说\BibleQuote{祂虽具有天主的形像，却没有以自己与天主同等为应当把持不舍的，反而贬抑自己，取了奴仆的形像}的人。在人中没有一个会把这话归于自己。从没有一个活着的圣人是独生天主而成为人——因为取\InnerQuote{奴仆的形像}、\InnerQuote{具有天主的形像}的意义正在于此。那么，如果真福伯多禄说到那位\InnerQuote{贬抑自己}以取\InnerQuote{奴仆的形像}的，如果那位\InnerQuote{具有天主的形像}的确实\InnerQuote{贬抑了自己}以取\InnerQuote{奴仆的形像}，如果那起初就是天主、是圣言和独生天主的，就是那位\InnerQuote{具有天主的形像}的，那么真福伯多禄对我们所说的就是那起初就与天主同在的圣言，并向我们阐明祂成了主和基督。那么，这就是巴西略与自己进行的冲突，他显然既没有\InnerQuote{将自己的心智用于宗徒们的意向}，也无法保持自己论点的连贯性；因为，根据这些论点，如果他意识到它们不可调和，他就必须承认那起初就是天主的圣言成了主；或者，如果他试图将相互矛盾的陈述拼凑在一起，并固执地坚持它们，他将增添更多敌对的陈述，并主张有两位基督和两位主。因为如果起初就是天主的圣言是一位，那位\InnerQuote{贬抑自己}并取了\InnerQuote{奴仆的形像}的是另一位，如果天主圣言——万物藉祂而成——是主，而这位耶稣——万有形成后被钉十字架的——也是主，那么按照他的观点，就有两位主和两位基督。因此，我们的作者无法用任何论点为自己辩护，摆脱这明显的亵渎。但如果有人为他辩解说，起初的圣言确实就是那成了主的一位，但祂是在肉身的临在中成了主和基督，那么他必然会被迫说圣子在肉身临在之前不是主。无论如何，即使巴西略和他的不信随从们虚假地宣称有两位主和两位基督，为我们来说却只有一位主和基督，万物藉祂而成，祂不是通过提升而成为主，而是在一切受造物和万世之前就已存在的主耶稣，万物藉祂而成，而所有圣人们以和谐的声音教导我们这一真理，并宣布它是最卓越的道理。在这里，真福若望教导我们，万物藉祂而成的天主圣言成了肉身，说：\BibleQuote{圣言成了血肉}；在这里，最卓越的保禄劝勉听从他的人谦卑，论及基督耶稣，祂虽具有天主的形像，却贬抑自己，取了奴仆的形像，且谦卑至死，且死在十字架上\BibleRef{斐 2:7}；在另一处，他又称被钉十字架者为\BibleQuote{光荣的主}：他说，\BibleQuote{因为他们若知道，就不致将光荣的主钉在十字架上了}\BibleRef{格前 2:8}。的确，关于本质本身，他用\InnerQuote{主}这个名称说得更为明了，他说：\BibleQuote{主就是圣神}\BibleRef{格后 3:17}。因此，如果那起初的圣言，就祂是圣神而言，是主，是光荣的主，如果天主要祂成为主和基督，那么天主所造的正是圣神本身和天主圣言，而不是巴西略所梦想的某位别的主。}\par

% structure: p0010 source=h2 target=subsection reason=relative-heading-rank; base=h2

\subsection{三、对这些言论的卓越而独创的回应，以及对被钉者权能的展示，并证明这服从是属于人性，而非独生者由圣父所享有。另加上十字架图形的解释，以及\Quote{基督}这一称号，并论及与人性结合的天主性所赐予人性的美好恩惠。}

好了，这就是他的指控。但我认为有必要首先以提要方式简要地过一遍他所提出的论点，然后以我的论证来纠正他所说的，使那些判断真理的人易于记住我们所必须回应的控诉，并且我们能够依次处理每一项指控。他说我们以基督的十字架为耻，诽谤圣人，说一个人\Quote{\OriginalTerm{空虚自己}{\GreekText{κένωσις}}}而成为人，又以为主在降生成人之前已具有\Quote{\OriginalTerm{奴仆的形象}{\GreekText{μορφὴ} \GreekText{δούλου}}}，将我们的救赎归功于一个人，在我们的教义中谈论两个基督和两个主；或者，如果我们不这样做，我们就否认\OriginalTerm{独生子}{\GreekText{μονογενής}}在受难之前是主和基督。为了避免这种亵渎，他要我们承认圣子的本质是受造的，因为宗徒\Person{伯多禄}亲口证实了这种教义。这就是指控的实质；至于他不厌其烦地对我们所说的谩骂话，我将默然不予理会，因为它们完全与要点无关。或许这种按照某种矫揉造作的理论形成的辞藻性笔法，是那些擅耍修辞术之人的惯常伎俩，一种为了夸大控诉篇幅的杜撰。让我们的诡辩家施展他的技艺，显示他的傲慢，在责备我们时炫耀他的力量，在竞赛间隙卖弄他的笔法；让他称我们为愚笨的，称我们为众人中最鲁莽的，众人中最可怜的，充满混乱和荒谬，并随意以任何他喜欢的方式轻慢我们，我们将承受；因为对明理人而言，耻辱不在于听到辱骂他的人，而在于回击他所说的话。他耗费气力攻击我们或许也有些益处；因为当他用辱骂的舌头谴责我们时，至少他在对抗天主的战斗中会暂时休战。所以，让他尽情傲慢吧：无人会回应他。因为如果一个人因身体失调或某种恶疾而口臭难闻，他不会激起任何健康的人仿效他的不幸，以至于选择自己染上疾病，用同样的恶行来回报此人臭气的不适。我们共通的本性要求我们怜悯这样的人，而非仿效他们。因此，让我们放过所有这类他在论证中不断掺入的嘲笑、愤慨、挑衅和辱骂，只审查那些涉及争议教义点的论证。那么，我们将从头开始，依次回应他的每一项指控。\par

他指控的开端，是说我们为我们而受苦的祂的十字架感到羞耻。他难道不也打算指控我们宣讲本质差异的教义吗？其实，反倒正是那些偏向这个意见的人，才负有让十字架成为可耻之事的责难。因为如果双方都同样将受难的天主经世视为信仰的一部分，而我们则认为必须尊敬那位藉十字架显现的天主，如同尊敬圣父一样，而\Emph{他们}却认为受难妨碍了将独生子天主与生祂的圣父同等光荣，那么，我们这位诡辩家的指控就反噬自身，他用那些自以为是控告我们的话，宣扬了他自己的教义不虔。因为显然，他将圣父置于圣子之上，以至高尊荣高举圣父的原因在于——在圣父身上看不到十字架的耻辱；而他断言圣子的本性在卑微意义上有所差异的原因在于——十字架的耻辱只归于圣子，不触及圣父。\newline{}但愿无人以为我说这话只是追随他作品的一般思路，因为我通读了他言论中所有亵渎之辞（这些言辞在那里被费力地汇集），在较前这段稍后的一处，我发现了这种亵渎以毫不掩饰的语言清楚表达出来；我打算在我自己论证的有序进程中，摆出他们所写的，内容如下：——他说：\Quote{如果，}他说，\Quote{他能证明那统治万有的天主，即那不可接近的光，降生成人，或能降生成人，服从权威，听从命令，受人法律约束，背了十字架，那么让他说光是同等的光。}那么，是谁对十字架感到羞耻？是那位即使在受难后仍与圣父同等敬拜圣子的人，还是那位甚至在受难前就侮辱祂的人，不仅将祂与受造物并列，而且坚持祂具有可受苦的本性，理由是：若非祂具有能承受这些受苦的本性，祂就不可能经历祂的苦难？我们方面断言：甚至祂用来受苦的身体，通过与天主性的混合，也因这混合而变成了那接纳本性所是的。我们远非对独生子天主有任何卑下想法，反而认为，若在祂爱人经世中接纳了任何我们卑微本性之物，连这也被转化成了神圣不朽的；但欧诺米却将十字架的苦难视为本质差异（在卑微意义上）的标记，我不知何故，竟将祂能完成此事的超越大能之举视为软弱的证据；他没有察觉到，凡是按其本性运动之物，都不被视为惊人的奇妙，凡超越其本性界限之物，则尤其成为赞叹的对象，每一只耳朵都转向它们，每一个心灵都专注它们，惊叹于这奇事。因此，所有宣讲圣言的人都指出奥秘在这方面奇妙之处——就是\Quote{天主在肉身显现}，就是\Quote{圣言成了血肉}，就是\Quote{光在黑暗中照耀}，\Quote{生命尝了死亡}，以及所有此类宣告，即信仰的宣扬者惯于作出的，由此增加了那位藉自身本性以外的事物彰显其大能丰盈者的奇妙性。\newline{}但是，尽管他们以此为傲慢的题目，尽管他们以十字架的经世为由将圣子与圣父的光荣平等分开，我们却相信，正如那些\Quote{从起初就亲眼看见并做了圣言仆役的人}藉圣经传给我们的，那位在起初就有的天主，后来，如巴路克所说，\Quote{在地上显现，与人往来}\BibleRef{巴 3:37}，并成为我们死亡的赎价，藉自己的复活解开了死亡的锁链，亲自使复活成为一切血肉的道路，并与自己的圣父同享同一宝座、同一光荣，在审判之日将按各人所行的赏罚来判决受审判者。这些就是我们关于被钉十字架者所相信的，为此我们不停地极力颂扬祂，按我们的能力，因为那位因不可言说、不可接近的伟大而除祂自己、圣父和圣神外无人能领悟的祂，我说，竟能下降到与我们软弱相通的地步。但他们却引证主藉肉身和十字架显现，作为圣子本性异于圣父的证据，理由是：圣父的本性在不受苦中保持纯洁，丝毫不能接纳倾向于受苦的相通，而圣子，因其本性在卑微方面的差异，并非不能受肉身和死亡的经验，因为状况的改变并不大，而是从一种类似状态到另一种相近而同类的状态，因为人的本性是受造的，独生子的本性也是受造的。那么，谁该公正地被指控为对十字架感到羞耻？是那卑下谈论它的人，还是那为它更高超方面争辩的人？我不知道我们这位控告者，如此贬低那在十字架上被认识的天主，是否听过保禄崇高的讲话，他用何等措辞、何等篇幅以他崇高的嘴唇谈论那十字架。因为他，能藉如此众多而伟大的奇迹使自己为人所知，却说：\Quote{但愿我别无所夸，只夸我们主耶稣基督的十字架。}\BibleRef{迦 6:14}他又对格林多人说：十字架的话\BibleQuote{为那正在得救的人，是天主的大能。}\BibleRef{格前 1:18}而且，对厄弗所人，他以十字架的图形描述了那掌管和维系宇宙的大能，当他渴望他们被高举以认识这大能的超越光荣时，称之为\BibleQuote{高、深、广、长}\BibleRef{弗 3:18}，用它们特有的名称谈论我们在十字架图形上看到的几个突出部分，以至于他称上部为\Quote{高}，下方接合处对侧为\Quote{深}，而以\Quote{长和广}之名指示向两侧突出的横木，借此以显明这伟大的奥秘：天上、地下以及一切存在的最远边界，都由那位在十字架图形上给予这不可言说又强大大能榜样者所统治并维系。\newline{}但我认为无需再与这类异议争辩，因为只要寥寥数语便能显明真理，我却焦虑于驳斥诬蔑的论证，实属多余。因此，让我们转向另一项指控。\par

他说我们诽谤了圣人。好吧，如果他亲耳听见，就请他把我们诽谤的话说出来；如果他以为我们是对别人说的，就请他用见证人证明他指控的真实性；如果他根据我们所写的来证明，就请他把那些话读出来，我们愿意承担责备。但他提不出任何这类证据：我们的著作是公开的，任何人都可以查阅。如果这话不是对他说的，他也没有从别人那里听到，更无法从我们的著作中提出证据，我认为在这个问题上需要答辩的人完全可以保持沉默：沉默无疑是对于无根据的指控最合适的回答。\par

宗徒\Person{伯多禄}说：\BibleQuote{天主使你们钉死的这位耶稣，成为主和基督。}\BibleRef{宗2:36}我们从他那里学到这一点，便说这段经文的整个上下文都指向同一个方向——十字架本身、人的名字、陈述的指示性语气。因为圣经的话是说，就同一个人而言，两件事同时发生——犹太人施加了苦难，天主则赐予了尊荣；并不是说一个人受苦，另一个人因被高举而得尊荣：他在接下来的话中更清楚地解释了这一点：\Quote{他被天主的右手高举。}那么，谁被\Quote{高举}了？是那卑微的，还是那至高者？卑微的除了人性，还能是什么？至高者除了天主性，还能是什么？诚然，天主不需要被高举，因为他本是至高者。由此可见，宗徒的意思是人性被高举了：它的高举是通过成为主和基督而实现的。这事发生在苦难之后。因此，宗徒用\Quote{成为}一词所指的，并非主在时间之前的存在，而是那卑微者被\Quote{天主的右手}提升为崇高者的变化。即使这句话也揭示了虔敬的奥迹；因为说\Quote{被天主的右手高举}的人，显然揭示了这奥迹那不可言说的安排：那创造万有天主的右手（即主，万物藉着他而被造，离了他，存在的无一存留），亲自将与之结合的人提升到自己的高度，也使他成为他自己本性所是的那一位。如今他是主和君王：基督是君王的称号；这些也使他成了。因为他既然因在至高者中而被高举，同样也成了其他一切——在不朽者中成为不朽，在光明中成为光明，在不朽坏者中成为不朽坏，在不可见者中成为不可见，在基督中成为基督，在主中成为主。即使在物质的结合中，当一方大大超越另一方时，低劣者往往完全改变成更强大者。这一点，宗徒\Person{伯多禄}在他的奥秘言论中清楚地教导了我们：那因软弱而被钉死者的卑微本性（软弱，正如我们从主听说的，是肉身的标记），我说，那卑微的本性，由于与无限无边的善结合，不再停留在自身的限度和属性中，而是被天主的右手同自己一起高举，由仆人变成了主，由臣民变成了君王基督，由卑微变成了至高，由人变成了天主。那么，那位假装为了捍卫宗徒而警告我们提防别人的人，从我们著作的材料中找到了什么攻击圣人的把柄呢？让我们也默默地略过这项指控；因为我认为，为虚假无据的指控辩护是一件卑鄙可耻的事。让我们继续讨论他指控中更为紧迫的部分。\par

% structure: p0015 source=h2 target=subsection reason=relative-heading-rank; base=h2

\subsection{四、他显示\Person{欧诺米}诬蔑\Person{大巴西略}说\Quote{人被倒空而成为人}的谎言，并证明独生子的\OriginalTerm{贬抑}{emptying}是为了使那曾受苦的人得以恢复生命。}

祂断言我们说，人虚空了自己成为人，并且那位藉服从谦卑自下取了奴仆形象的，甚至在取得那形象之前就已拥有了人的形象。措辞上没有任何改动；我们只是照搬了他话语中的原词。如果我们的著作中有任何此类说法——因为我称我老师的著作为\Emph{我们的}——那么没有人会责怪我们的演说者诽谤。我要求一切顾及真理：我们自己也会作证。但如果我们的著作中根本没有这一切，而他的语言不仅指责我们，而且义愤填膺，仿佛事情已确凿证明，称我们充满荒谬、胡说、混乱、矛盾等等，那我真不知该走哪条路了。正如那些对疯子的无端狂怒感到困惑的人无法决定采取什么计划一样，我自己也找不出应对这一困惑的计策。我们的老师说（让我再次逐字复述他的论点）：\Quote{祂并非向我们阐述天主存在的模式，而是属于降生成人的言词。}我们的控告者由此出发，声称我们认为人虚空了自己成为人！这两句话之间有什么共同之处呢？如果我们说宗徒并未向我们阐述天主存在的模式，而是藉其措辞指向苦难的救恩安排，那么我们便因此被指控谈论人的\Quote{虚空自己}以成为人，并说\Quote{奴仆形象}是先在的，且生于玛利亚的人在降生之前就已存在！我认为，花时间讨论已承认之事是多余的，因为真理本身已免除我们的指控。确实，在一个人可能给诽谤者提供把柄的情况下，应当反驳控告者；但若没有遭受某种荒谬指控的嫌疑，则指控所证明的，并非被诽谤者的虚假罪行，而是控告者的疯狂。然而，正如我们在处理以十字架为耻的指控时，通过审查表明该指控回击了控告者一样，我们也将表明这项指控如何也回击了那些提出者，因为正是他们，而非我们，在苦难的救恩安排中确立了圣子从相似到相似的改变理论。我们将简要审查，将双方的说法并排比较。我们主张，独生子天主亲自创造万有，祂自己享有对万有的全权，而人的本性也是祂所创造的万物之一；当这本性堕落于恶，陷入死亡的毁灭时，祂亲自藉由祂所居住的人，完全取了人性，将祂的生命之能与我们这必朽坏的本性混合，并藉与祂自身的结合，将我们的死亡改变为生命的恩宠和德能。我们认为这是主按肉体的奥迹：那不变者进入可变者之中，旨在使其改善，从较劣转为较优，从而消灭与我们可变状态相混合的邪恶，在祂自身内摧毁邪恶。因为\BibleQuote{我们的天主是吞噬之火}\BibleRef{希 12:29}，一切邪恶之质都被祂消除。这就是我们的说法。我们的控告者怎么说呢？不是说那不变、不受造者与受造而因此向恶变化的本性结合，而是说祂自身是受造的，来到与自己同族同类的存在，并非出于爱人之心从超越的本性降卑穿戴低微的本性，而是成为祂所是的那位本身。\par

论到\Quote{\OriginalTerm{受造物}{creature}}这个称谓的一般特性，它是用来指称一切从无中生成之物的名称；而我们视为包含在\Quote{受造物}一词中的事物，其不同部分因属性的差异而彼此区分。因此，如果祂是受造的，而人也是受造的，那么祂便\Quote{\OriginalTerm{空虚}{emptied}}了自己（借用欧诺弥的用语）以成为祂自己，并且改变了祂的位置——不是从超越降为卑微，而是从同类之物变为类似尊严之物（除了身体与无形体之特殊性质外）。那么，现在谁会在审判我们案件的人那里得到公正的投票？谁看来是背负着这些指控？是说受造物被非受造的天主拯救的人，还是将我们得救的原因归于受造物的人？虔敬之人的判断无疑是不容置疑的。因为任何清楚知道受造者与非受造者之间区别的人（这两个术语的分歧由统治与奴役所标志，因为非受造的天主，如先知所说，\BibleQuote{以祂的权能永远统治}，而受造的一切都是祂的仆人，正如同一先知的声音所言：\BibleQuote{万物都事奉祢}），我说，仔细考虑这些事的人，必定不能认出那使独生子从奴役变为奴役的人。因为如果按照保禄，整个受造界\Quote{在奴役之中}\BibleRef{罗 8:21}，而按照欧诺弥，独生子的本质是受造的，那么我们的对手无疑是在他们的教义中主张，不是主人与仆人混合，而是一个仆人来到仆人中间。至于我们说主在祂肉身显现之前就有\OriginalTerm{奴仆形象}{form of a servant}，这就像指控我们说星辰是黑的、太阳是雾蒙蒙的、天空是低的、水是干的一样：一个不基于所听来指控，而是随心所欲编造自己认为合适之事的，不必吝于对我们提出这类指控。我们被要求对这组指控负责，与对另一组指控负责，就它们在我们所说的话中毫无根据这一点来说，是完全相同的。一个清楚说真子原在圣父荣耀中的人，怎能因认为祂曾有\Quote{奴仆形象}而侮辱独生子的永恒荣耀呢？当我们的作者认为合适诽谤我们，同时又小心使其说法看起来有些真实时，反驳他的无端指控也许并非多余或无益。\par

% structure: p0018 source=h2 target=subsection reason=relative-heading-rank; base=h2

\subsection{五、此后，他表明并非有两个基督或两个主，而是一个基督和一个主；神性与人性结合后，无混淆地保持了每一本性的属性；并宣告由于结合，两性的行动可共同言说，即主承担了仆人的苦难，人性在主的荣耀中同受光荣，且藉神性本性的能力被更新，与神性本性自身相符合。}

他的下一个指控也具有同样类型的荒谬。因为他指责我们说有两个基督和两个主，却无法从我们的话中证实他的指控，而是随意使用谎言迎合自己的幻想。既然他认为自己有权说想说的话，为何如此细致地发出谎言，并坚持我们只说两个基督的？他若愿意，不妨说我们宣讲十个基督，或一百个，或把数目扩展到一千，以便更激烈地处理他的诽谤。因为亵渎同样包含在两个基督的教义中，也包含在更多的教义中，而这两个指控的特点也同样缺乏证据。那么，当他证明我们确实说两个基督时，就让他的判决对我不利，就好像他证明了有一万个一样。但他声称他通过我们自己的陈述来定罪我们。好吧，让我们再看一次我们老师的那些话，他正是通过这些话来指控我们的。他说：\Quote{他（即说出\InnerQuote{天主已立祂为主、为基督，就是你们所钉死的这位耶稣}的人）并非向我们陈述天主性存在的方式，而是属于\OriginalTerm{降生成人}{Incarnation}的术语……藉指示词强调祂身上那属人性并为众人所见的一面。}这就是他所写的。但欧诺弥如何通过这些话将他的\Quote{两个基督}搬上舞台？说指示词强调那可见的一面，就证明了主张\Quote{两个基督}吗？难道为了避免被指控说\Quote{两个\OriginalTerm{至高者}{Highest}}，我们就应当否认事实，即祂在苦难后被天主高举？因为那在起初就与天主同在的圣言本是至高者，但照宗徒所说，祂在苦难后复活时也被高举。我们必须在两者中择一：要么说祂在苦难后被高举（这与说祂被立为主和基督相同），因而被欧诺弥指控；要么，如果我们避免这指控，就否认对受难者被高举的承认。\par

此时，似乎应当再次提出我们控告者的陈述，以支持我们自己的辩护。因此，我们将逐字重复他所提出的陈述，该陈述如下支持我们的论点：——\Quote{真福若望}，他说，\Quote{教导我们：万有藉着祂而受造的天主圣言降生成人，说\InnerQuote{圣言成了血肉}}。当他将这句话加入自己的论点时，他明白自己在写什么吗？我自己几乎难以相信，同一个人能同时知晓这些话的意义，却又反对我们的陈述。因为若有人仔细审视这些话，就会发现在我们所说与他所说之间并无冲突。因为我们二人都将肉身中的安排视为分开的，并思考天主性本身；而他同我们一样，说起初就有的圣言已在肉身中显现：然而从未有人控告他宣讲\Quote{两个圣言}——起初就有的那位和成了血肉的那位——他也不控告自己；因为他一定知道，圣言与圣言是同一的，在肉身中显现的那位与在天主面前的那位原是同一的。但是，血肉与天主性本不相同，直到这血肉也被转化到天主性中，以至于必然有一些属性属于天主圣言，而另一些属性属于\Quote{奴仆的形体}。那么，既然他在承认这些之后并不因圣言的二元性而责备自己，为何我们被诬告将我们信仰的对象分为\Quote{两个基督}呢？——我们主张，祂在受难后被高举，因与那位真实的主和基督结合而成为主和基督；我们知道从所学的一切中，天主性始终是同一且不变的，具有同一的存在方式，而血肉本身是理性与感官所认知的，但与天主性混合后，便不再停留在自己的界限与属性内，而被提升到那超越而卓越的境地。然而，我们对血肉与天主性各自属性的默观，只要各自独自被思考，便保持清晰不混：例如，\Quote{圣言在万世以前，但血肉在末期才存在}；但人不能颠倒这话，说后者是先于时间的，或者说圣言在末期才存在。血肉属于可受苦的本性，圣言属于运作的本性；血肉不能造出万有，天主性所具有的能力也不能受苦。圣言起初就与天主同在，这人却要经历死亡的考验；人性并非从永远存在，天主性也不是可朽的：其他所有属性也以同样方式默观。使拉匝禄复活并非人性，那不能受苦的能力也不会在拉匝禄躺在坟墓时为他哭泣：眼泪出于这人，生命出于真生命。喂养数千人并非人性，全能的力量也不会急着去无花果树那里。是谁因旅途疲倦？又是谁以祂的圣言使全世界存在？什么是光荣的辉耀，什么又被钉子穿透？什么形体在受难中被掌击，什么形体从永远就被光荣？这一点很清楚（即使人不详细思考论证）：打击属于主所在的奴仆，荣耀属于奴仆所环绕的主，以致于因接触与本性的结合，各方的固有属性同时属于两者，就如主承受奴仆的鞭痕，而奴仆因主的荣耀而受光荣；这正是为何十字架被称为光荣之主的十字架\BibleRef{格前 2:8}，也是为何众口都宣认耶稣基督是主，归于天主父的光荣\BibleRef{斐 2:11}。\par

但是，如果我们要以同样的方式讨论其他点，就让我们思考什么是死的，什么是毁灭死亡的；什么是更新的，什么是空虚自己的。天主性\Quote{空虚}自己，为能进入人的本性的能力范围；而人的本性通过与神圣者混合而成为神圣的，因而被更新。就如空气被重物带入水中，留于水深之处，但它不会停留在水中，而是迅速升到同类的元素那里；水在空气上升时也常随之一同升起，被空气的环圈塑造成带有轻微薄膜表面的凸形状。同样，作为肉体基础的真正生命，在受难后迅速升到自己那里，肉体也与祂一同被举起，因神圣的不朽性而被迫从必朽上升到不朽。就如隐藏在木头深处的火，常不为人所见，甚至触摸也感觉不到，但燃烧时就显露出来；同样，在祂的死亡中（这死亡是祂按自己的旨意带来的，祂将灵魂与身体分开，祂向自己的圣父说：\BibleQuote{我把我的灵魂交在祢手中}\BibleRef{路 23:46}，祂又说：\Quote{我有权放下它，也有权再取回它}），祂，因为祂是光荣的主，藐视在人看来可耻的事，将祂生命的火焰，藉死亡的照顾，隐藏在祂的肉体本性中，然后用祂自己天主性的德能再次点燃并使之燃烧，使那已归死亡的生命得以复苏；祂以神圣德能的无限性注入我们本性卑微的初果，使之也成为祂自己所是的——使奴仆形象成为主，使\Person{玛利亚}所生的人成为基督，使因软弱而被钉死的那位成为生命和德能，并使一切虔敬地默想存在于天主圣言中的属性，也存在于圣言所取的人性中；这样，这些属性不再被视为分属两个本性，而是那必朽的本性通过与神圣者的混合，被重新塑造，符合那淹没它的本性，参与天主性的德能，就如有人说，混合物使一滴醋在大海中变为海水，因为这液体的天然性质在淹没它的无限性中不再持续。这就是我们的教义，并不像\Person{欧诺弥}指控的那样，宣讲多个基督，而是宣讲人与天主性的合一；它用\Quote{造成}的名称来形容从必朽到不朽、从奴仆到主、从罪恶到正义、从诅咒到祝福、从人到基督的转化。我们的诽谤者还有什么可说，以显示我们在教义中宣讲\Quote{两个基督}呢？如果我们拒绝说那从起初从圣父而来、不因受造而为主、为基督、为圣言、为天主的那位是\Quote{被造成的}，并宣称蒙福的\Person{伯多禄}仅是一时顺便指出降生成人的奥迹，按照现在解释的意思，即那因软弱而被钉死的本性，也因内住者的压倒性德能，已成为内住者本身实际和名义上所是的，即基督和主。\par

\SourceRecord{Translated by H.A. Wilson.；Nicene and Post-Nicene Fathers, Second Series；Vol. 5.；Edited by Philip Schaff and Henry Wace.；Buffalo, NY: Christian Literature Publishing Co.,；1893.；<http://www.newadvent.org/fathers/290105.htm>.}

% structure: p0001 source=h1 target=section reason=page-title

\section{\WorkTitle{驳斥欧诺弥（第六卷）}}

% structure: p0002 source=h2 target=subsection reason=relative-heading-rank; base=h2

\subsection{一、第六卷表明：为人的救恩而来的那一位并非纯然的人——如同\Person{欧诺弥}在诽谤中所妄称的，伟大的\Person{巴西略}曾如此说——而是天主的独生子，披上人的肉身，成为天主与人之间的中保；我们相信祂，在肉身内受苦难，但在祂的天主性内不受苦难；并证明了\Person{欧诺弥}的诽谤。}

但我觉察，由于主题的必然性迫使我依循这一思路，我在此处已停留过久。现在必须继续他抱怨的脉络，以免我们对他加于我们的任何指控不予答复。首先，我提议我们应研究这一点：他指控我们声称一个普通人完成了世界的救恩。尽管通过我们已进行的探讨，这一点在某种程度上已得澄清，但我们仍将再次简要处理，以便那些在此诽谤指控中担任我们审判者的人能完全摆脱误解。我们绝不将如此伟大而不可言喻的恩宠的原因归于一个普通人，即使有人将如此巨大的恩赐归于\Person{伯多禄}和\Person{保禄}，或归于从天而来的天使，我们也要与\Person{保禄}一起说：\BibleQuote{让他受诅咒。} \BibleRef{格前 1:13} 当然，当我们承认基督的救赎之能比人性更强大时，我们对手反对真理的学说并未因此得到加强——然而，这似乎可能如此，因为他们的目标是在一切方面维护圣子与圣父本质的差异，他们竭力不仅通过受生者与不受生者的对比，而且通过可受难者与不可受难者的对立来显示本质的不相似。这一点在他们的论证最后部分中更公开地得到维护，而在他们当前的论述中也清楚显示。因为如果他指责那些将苦难归于人性的人，他的意图无疑是使天主性本身屈从于苦难。因为我们的观念是双重的，并允许两种发展，即根据神性或人性被视为处于受难状态，攻击其中一种观点显然就是维护另一种。因此，如果他们指责那些视苦难为关乎那\Quote{人}的人，他们显然会赞同那些说圣子的天主性受苦难的人，而后一立场便成为他们荒谬学说的论据。因为根据他们的说法，圣子的天主性受苦难，而圣父的天主性保持绝对不受苦难，那么不受苦难的本性在本质上与接纳苦难的本性不同。因此，鉴于我们面前的这段论述，尽管就字面数量而言是简短的，却为各种教义上的邪僻提供了原则和假设，我们的读者似乎应该要求我们的回应与其说简短，不如说健全。因此，我们既不将我们的救恩归于一个人，也不承认不朽和神圣的本性能够受难和死亡；但由于我们必须确信那些向我们宣告的神圣宣言：在起初就有的圣言就是天主，后来圣言成了血肉，在地上显现并与人同住，我们在信经中接纳那些与神圣宣言相符的观念。因为我们听到祂是光、能力、正义、生命、真理，万物由祂而造，我们就把所有这些以及类似的说法视为可信的事，归于天主圣言；但是当我们听到痛苦、睡眠、匮乏、忧愁、束缚、钉子、长枪、血、伤口、埋葬、坟墓，以及其他一切这类的事，即使它们与先前所说有些矛盾，我们仍同样承认它们是可信的、真实的，就肉体而言，我们因信接纳肉体与圣言结合。因为正如无法在起初就有的圣言中默观肉体的特有属性，同样，我们也不能认为那些属于天主性的属性存在于肉体的本性中。因此，正如福音关于我们主的教导是混合的，一部分是崇高和神圣的观念，一部分是卑微和人的观念，我们相应地将每一个特定的说法归于我们在奥迹中所默观的这两个本性之一：属于人的归给人性，属于崇高的归于天主性，并说，作为天主，圣子当然是不可受难、不可朽坏的；而福音关于祂所声称的任何苦难，祂无疑是通过祂那接纳苦难的人性来实行的。因为天主性确实藉着环绕祂的身体来施行世界的救恩，如此，苦难属于身体，而运行属于天主；即使有人将宗徒的话曲解为支持相反的学说：\BibleQuote{天主没有怜惜自己的儿子，}\BibleRef{罗 8:32}，以及\BibleQuote{天主派遣了自己的儿子}，以及其他类似的说法，这些说法似乎在苦难问题上指涉天主性而非人性，我们仍不会放弃正确的教导，因为\Person{保禄}自己更清楚地向我们宣告了这主题的奥迹。因为他处处将苦难的经纶归于基督内的人性因素，当他说：\BibleQuote{因为死亡是藉着人来的，死者的复活也是藉着人来的，}\BibleRef{格前 15:21}，以及\BibleQuote{天主派遣了自己的儿子，带着罪恶肉身的形状，就在肉身上判定罪过}（因为他说\BibleQuote{在\Emph{肉身上}，}而非\Quote{在天主性上}）；以及\BibleQuote{祂因软弱被钉在十字架上}（此处借\Quote{软弱}意指\Quote{肉体}），\BibleQuote{却因天主的德能活着}\BibleRef{格后 13:4}（而藉\Quote{德能}他指明天主性）；以及\BibleQuote{祂因罪而死}（即就身体而言），\BibleQuote{却为天主而活}\BibleRef{罗 6:10}（即就天主性而言，如此由这些话确立：那\Quote{人}尝了死亡，而不朽的本性并未接纳死亡的苦难）；以及\BibleQuote{祂曾使那不认识罪的，替我们成了罪}\BibleRef{格后 5:21}，再次赐予\Quote{罪}这一名称为肉体。\par

% structure: p0004 source=h2 target=subsection reason=relative-heading-rank; base=h2

\subsection{二、随后他再次提及圣伯多禄所说的\Quote{\OriginalTerm{受造}{made}}，以及\WorkTitle{希伯来书}中称耶稣被天主\Quote{造为宗徒和大司祭}的段落；在对欧诺弥向他提出的指控给出充分答复后，他表明欧诺弥本人反倒支持了巴西略的论点，并指出独生子在穿上血肉后成为了主。}

虽然我们顺便做了这些评论，但这个插入的补充似乎不亚于眼前的主要问题。因为，当圣伯多禄说\BibleQuote{祂立祂为主和基督\BibleRef{宗 2:36}}，此外，当宗徒保禄对希伯来人说祂造祂为司铎时，欧诺弥抓住\Quote{\OriginalTerm{受造}{made}}这个词，认为它适用于祂的先存存在，并试图借此建立他的教义，即主是\OriginalTerm{受造之物}{a thing made}。现在让他听从保禄所说的话：\BibleQuote{祂使那不认识罪的，替我们成了罪\BibleRef{格后 5:21}}。如果他把在希伯来书和伯多禄的话中用于主的\Quote{\OriginalTerm{受造}{made}}这个词，参照到先存的概念，那么他同样可以把那段说天主使祂成为罪的经文中\Quote{受造}这个词，参照到祂本体的最初存在，并试图借此，如同他其他证据一样，表明祂是\Quote{\OriginalTerm{受造}{made}}，从而将\Quote{\OriginalTerm{受造}{made}}这个词归结到本体，自相一致，并在那本体中看出罪来。但如果他因这明显的荒谬而退缩，并辩称，宗徒说\Quote{祂使祂成为罪}是指末世的安排，那么让他以同样的推理说服自己，即\Quote{\OriginalTerm{受造}{made}}这个词在其它段落中也指那安排。\par

然而，让我们回到我们离题之处；因为我们能从同一圣经中收集无数其他经文，除已引用的之外，它们都与这事有关。且不要以为神圣的宗徒自相矛盾，并藉他自己的言论为那些在教义上争论的人提供双方争论的材料。因为仔细检查会发现，祂的论证准确地指向一个目标；祂在自己的意见上并不踌躇：因为祂虽处处宣告人性与神性的结合，却仍辨别各自的本性，即人的软弱因与不朽的共融而变得更好，而神能另一方面却不因与卑微的本性形态接触而降低。因此，当他说：\Quote{祂没有怜惜自己的圣子，} 他将真圣子与其他众子——生养的、高举的或义子（我是指那些按照祂的命令受造的）——作对比，藉着加上\Quote{自己的}来标示本性的特殊性。并且，为避免任何人将十字架的苦难与不朽的本性联系起来，他用另外的话相当明确地纠正了这种错误，当他称祂为\BibleQuote{在天主与人之间的中保\BibleRef{弟前 2:5}}和\BibleQuote{人\BibleRef{弟前 2:5},}和\BibleQuote{天主,}，这样，由于两者都归于同一个存有，就能对每个本性持有合适的观念——关于神性，是不能受苦；关于人性，是苦难的安排。因此，鉴于祂的思想将那因爱世人而成为一体、但在概念上有区别者分开了，当祂宣告那超越并超过一切理智的本性时，祂使用更崇高的名号次序，称祂为\BibleQuote{在万有之上的天主\BibleRef{罗 9:5}}、\BibleQuote{伟大的天主\BibleRef{铎 2:13}}、\Quote{德能}和\Quote{智慧}\BibleRef{格前 1:24}等等；但当祂提及那因我们的软弱而必然与我们的本性一起被摄取的一切苦难经验时，祂给两性结合赋予一个源于我们的名号，称祂为人，并不是用这个词将祂向我们展示的那个人置于与其他本性同等的层面，而是为使正统在每个本性的意义上得到保护，即人性因祂的摄取而得荣耀，神性不因祂的屈尊而受玷污，但使人性因素服从于苦难，同时藉其神性德能，使那受苦者复活。因此，死亡的经验不归于那因人的结合而与我们可受苦的本性相通的那一位，同时崇高和神圣的名号下降到人身上，以致那在十字架上显现的那一位甚至被称为\BibleQuote{光荣的主\BibleRef{格前 2:8}}，因为藉其本性（神性）与那卑微的本性的混合，这些名号所含的威严从神性传递到人性。为此，祂用多变不同的语言描述祂，时而作为从天降下者，时而作为由女人所生者，作为永恒的天主，在末日作为人；同样，独生者天主被认为不能受苦，而基督能受苦；祂的言论在这些对立的陈述中不说谎言，因为它在概念上使属于每个本性的术语适应于该本性。如果这些是我们从受默感的教导中学到的教义，我们如何将我们得救的原因归于一个普通人呢？如果我们宣称蒙福的\Person{伯多禄}所用的\Quote{被立}一词不是指先于时间的存在，而是指降生成人的新安排，这又与我们被指控的有什么相干呢？因为这位伟大的宗徒说，那以奴仆的形状显现的，藉着被摄取，成了那摄取者在其自己的本性中的所是。此外，在希伯来书中，我们也能从\Person{保禄}学到同一真理，他说耶稣被天主立为宗徒和大司祭，\BibleQuote{忠于那立祂为这样的人\BibleRef{希 3:1}}。因为在那段经文中，在将大司祭的名号赋予那用自己的血作了我们罪过的司祭赎罪的那一位时，他并不是用\Quote{被立}一词来宣告独生者的最初存在，而是说\Quote{被立}意在表示那通常与司祭任命相关的恩宠。因为耶稣，这位大司祭（如\Person{匝加利亚}所说），祂献上自己的羔羊，即自己的身体，为世界的罪；祂因那些同有血肉的儿女，自己也照样亲身与他们有分于血（不是按照祂在起初的所是，即圣言和天主，具有天主的形象，与天主同等，而是按照祂以奴仆的形象贬抑自己，并为我们献上祭品和牺牲），我说，祂后来在许多世代之后，按照\Person{默基瑟德}的品位，成了大司祭。当然，一位对希伯来书有超过泛泛了解的读者知道此事的奥迹。因此，正如在那段经文中祂被称为被立为司祭和宗徒，同样，这里祂被称为被立为主和基督——后者是为了我们而作的安排，前者是藉着人性向神性的改变和转化（因为宗徒用\Quote{立}的意思是\Quote{再造}）。这样，我们敌人们的奸诈就显明了，他们无礼地曲解那些指安排的话，将其应用于先于时间的存在。因为我们从宗徒学到，不要像从前那样认识基督，正如\Person{保禄}这样说：\Quote{纵使我们曾按肉体认识基督，现今却不再这样认识祂了，}其意义是，一种认识给我们显示祂暂时的安排，另一种认识显示祂永恒的存在。因此，我们的言论对祂的指控作了不小的回答：——我们既不持守两位基督，也不持守两位主，我们不因十字架为耻，我们不荣耀一个纯粹的人为世界受苦，我们确实不相信\Quote{被立}一词指本质的形成。但是，既然这是我们的观点，我们的论证从控告者本人得到了不小的支持，他在言论中长篇大论攻击我们，并引用以下句子：\Quote{那么，这就是\Person{巴西略}与自己进行的争战，他明显既没有\InnerQuote{将自己的心思用于宗徒们的意向}，也不能保持自己论证的顺序；因为根据这些论证，他若意识到它们不可调和的性质，就必须承认那在起初就存在、就是天主的圣言成为了主，或者他拼合\InnerQuote{互相矛盾的陈述}。}哎呀，这实际上是\Person{欧诺米}重复的我们的陈述，他说：\Quote{那在起初就存在、就是天主的圣言成为了主。}因为，祂所是的（天主、圣言、生命、光明、恩宠、真理、主、基督，以及每一个崇高和神圣的名号），祂确实在祂所摄取的人身上成了这些——那人原不是这任何一样——而圣言曾是的其他一切，并且其中成了主和基督，按照\Person{伯多禄}的教导，也按照\Person{欧诺米}的自认；——不是神性藉推进获得了什么，而是（一切崇高的尊严都在神性中被默观）祂就这样成为主和基督，不是通过就其神性获得任何恩宠的增加（因为神性的本性被承认一无所缺），而是藉着引导人性参与天主性，这参与由\Quote{基督}和\Quote{主}这两个词所表示。\par

% structure: p0007 source=h2 target=subsection reason=relative-heading-rank; base=h2

\subsection{三、他随后对主向\Person{斐理伯}所说的话作出了卓越的解释：\Quote{看见了我，就是看见了圣父}；在此他极好地论述了主因对人的爱所受的苦难，以及圣父的不受苦性、创造能力和眷顾，又论述了人的复合本性，以及人复归为其构成元素的分解。}

关于这些点已经提供了充分的辩护。至于\Person{欧诺弥}以诽谤方式攻击我们教义所说的话，即\Quote{基督自我空虚，成为祂自己}，在前文已有充分讨论，其中已指明他将自己的亵渎归于我们的教义。因为承认不变的本性披上了受造的可朽坏者的那位，并不是在谈论从相同向相同的转变，而是认为本性的尊威没有降为更为卑微者的那位。因为如果如他们教义所宣称，祂是受造的，人也是受造的，那么教义的奇妙就消失了，所声称的就没有任何可惊叹之处，因为受造的本性自存。但我们从预言中学到\Quote{至高者右手的改变}——而我们以圣父的\Quote{右手}理解那创造万物的天主的能力，那就是主（不是作为部分依赖于整体那样依赖于祂，而是确实出于祂，却在其个体存在中被默观）——我们这样说：就本性的概念而言，右手与祂所是的右手并无不同，除了肉身的经世之外，也不能说其中有其他改变。因为天主的右手确实是天主自己；在肉身中显现，被那些目光清澈的人透过同一天肉身而看见；祂既实行圣父的作为，在事实和思想上都是天主的右手，然而就祂所被包裹的肉身之帷幔而言，就所见到的而论，祂从祂本性所是的那位改变为默观的客体。因此祂对只注视着那个改变之物的\Person{斐理伯}说：\Quote{透过那改变之物看向那不可改变者，如果你看见了这个，你就看见了那位你寻求看见的圣父本人；因为看见了我——不是那个在改变状态中出现的那位，而是我自身，即在圣父内的那位——将看见我在其中存在的那位圣父本人，因为完全相同的神性特征在两者之中都被看见\BibleRef{若 14:9}。}那么，如果我们相信那不死、不受苦、非受造的本性进入了受造物的可受苦的本性之中，并认为\Quote{改变}在于此，那么那些不停念叨自己关于我们教义的说法的人，凭什么指控我们说祂\Quote{自我空虚，成为祂自己}？因为受造者与受造者的共融并不是\Quote{右手的改变}。说非受造本性的右手是受造的，这唯独属于\Person{欧诺弥}和那些采纳他观点的人。因为用真理之眼观看的人，将辨明至高者的右手正如他所见的至高者本身——非受造者出自非受造者，善者出自善者，永恒者出自永恒者，祂经由受生而在圣父内，而不损及其永恒性。因此，我们的控告者不知不觉地使用了本应落在自己身上的指责来攻击我们。\par

但是，对于那些因\Quote{受难}的概念而绊倒，并以此主张本体差异的人——他们认为，圣父因本性崇高而不接纳受难，而圣子则因缺陷和歧异而屈尊参与祂的苦难——我想在已有论述之外补充以下几点：凡不倾向于罪的，并非真正的\Quote{受难}，严格来说，也不会将本性的必然进程称为\Quote{受难}，即认为复合本性以某种次序和序列进行。因为在我们身体形成中，异质元素的相互结合是一种由若干不同元素和谐联结的组合；但当适当时刻，将这元素结合的纽带松开时，结合的本性便再次分解为其组成的元素。这与其说是本性的受难，不如说是本性的运作。因为我们只将\Quote{受难}之名给予与有德的无情状态相对立的东西，并且我们相信，赐予我们救恩的那位始终没有这种受难，祂\Quote{在各方面都和我们一样受过试探，只是没有罪}\BibleRef{希 4:15}。至少，对于真正的受难，即意志的病态，祂并无分受；因为经上说：\Quote{祂没有犯罪，口中也没有诡诈}\BibleRef{伯前 2:22}；但本性中特有的属性——由于习惯上的术语滥用，也被称为\Quote{受难}——我们承认主确实分受了这些：出生、养育、成长、睡眠、劳苦，以及灵魂因身体不适而常有的所有自然倾向——对匮乏的渴望，当渴望从身体传到灵魂时，疼痛感、对死亡的恐惧，以及诸如此类，惟独那些若被追随就会导向罪恶的除外。因此，当我们察觉祂的能力贯穿天上、空气、地上、海中的一切时，我们相信祂无所不在，却不说祂是祂所在其中的任何事物；因为祂不是天，祂以怀抱之量标定天；祂不是地，祂支撑地之圆环；祂也不是水，祂环绕液态的本性，同样，我们也不说祂在经历我们所说的那些肉身的苦难时是\Quote{被动的}，而是如我们说祂是万有的原因，祂掌握宇宙，祂引导一切运动之物，并依祂自己尊威的不可言喻的能力，使一切静止之物立于固定的基础之上，同样，我们说祂为我们而诞生，为医治罪的疾病，以与苦难相应的方式调整祂医治能力的运用，以祂知道对祂所知道患病的受造部分有益的方式施用医治。既然用触摸医治苦难是适宜的，我们便说祂如此医治了；然而，祂并不因为是我们软弱的医治者，就因此被认为本身是可受难的。因为在人的情形中，日常用法也不允许我们肯定这样的事。我们不说触摸病人以医治他的人本人也分有病症，而是说他赐给病人恢复健康的恩惠，而不分有病症：因为苦难不接触他，是他接触疾病。如果凭技艺在人的身体上产生益处的人不被称作迟钝或软弱，而被称作爱人者、恩人等等，那么他们为何诽谤向我们施行的救恩安排，视其为低贱无荣，并用以主张圣子的本体\Quote{因低下而歧异}，理由是圣父的本性超越苦难，而圣子的本性并非不受难呢？如果降生成人的救恩安排的目的不是让圣子受难，而是彰显祂为爱人者，而圣父也无疑是爱人者，那么如果我们关注目的，圣子与圣父便处于同样的情形。但如果不是圣父亲自成就了死亡的消灭，不必惊奇——因为祂也将一切审判交给了圣子，祂自己不审判任何人；祂并非因为不能拯救迷失者或审判罪人而藉圣子行一切事，而是因为祂也藉自己的能力行这些事，祂藉此工作一切。那么，凡藉圣子得救的人，都是因圣父的能力得救；凡受祂审判的人，都是因天主的正义受审判。因为\Quote{基督}，如宗徒所说，\Quote{是天主的正义}\BibleRef{罗 1:17}，这正义藉福音启示出来；无论你观看整个世界，还是构成整体的世界各部分，这一切都是圣父的工程，因为都是祂能力的工程；因此，那说圣父造了一切，又说没有一样存在之物不是藉圣子而成的言语，在两方面都是真实的；因为能力的运作与能力所属的那位有关。因此，既然圣子是圣父的能力，圣子的一切工程都是圣父的工程。祂进入受难的救恩安排，不是由于本性的软弱，而是由于意志的能力，这一点可用无数福音经文证明；但既然事情清楚，我就不再赘述，以免我的论述因在承认之点上停留而冗长。因此，如果所发生的是恶，那么我们不仅要将圣父，也要将圣子从这恶中分离；但如果拯救迷失者是善，并且所发生的不是\Quote{受难}，而是爱人之心，那么你为什么将藉自己的能力即基督为人类成就脱离死亡自由的圣父，从我们对我们救恩的感谢中疏远呢？\par

% structure: p0010 source=h2 target=subsection reason=relative-heading-rank; base=h2

\subsection{四、然后他回到伯多禄的话：\Quote{天主立祂为主和基督}，他用许多论据巧妙地解释这一点，并在此显示欧诺弥是正统教义的辩护者，并通过指出神性和人性的名称因混合而应用于两种本性，来结束本书。}

但我们须再次回到我们那位激烈的演说撰稿人，并重新拿起他对我们的严厉指责。他埋怨我们否认圣子的本质是被造的，认为这与伯多禄的话相矛盾，即\Quote{天主立祂为主和基督，就是你们所钉死的这位耶稣}\BibleRef{宗 2:36}；他对此事大动肝火，出言辱骂，并且坚持某些观点，以为这些观点能驳斥我们的教义。现在，就让我们看看他论点的力度。\Quote{请问，你们这些最不顾一切的人，}他说，\Quote{当一个人具有仆人的形象时，会再取仆人的形象吗？}\Quote{没有一个有理性的人，}这是我们给他的答复，\Quote{会使用这种语言，除非他完全背离了基督徒的希望。但你们就属于这一类，你们指责我们不顾一切，因为我们不承认造物主是被造的。因为如果圣神不会说谎，当祂藉先知说：\InnerQuote{万物都事奉祢}，而整个受造界都在奴役之下，而圣子，如你们所说，是被造的，那么祂显然与万物同为奴仆，因祂分享受造界而下降到也分享奴役。而你们一定将那位处于奴役中的祂赋予仆人的形象：因为当你们承认祂本性是仆人时，你们当然不会以奴役的外貌为耻。请问，我最敏锐的修辞家，是谁把圣子从奴役的形象转到另一种仆人的形象？是那位为祂主张非受造的存在，从而证明祂不是仆人的人，还是你们，不停呼喊圣子是父的仆人，并在祂取仆人的形象之前实际上就在祂的统治之下？我不需要其他审判者；我将这些问题交由你们自己投票决定。因为我想，没有人会如此无耻地对待真理，以至于纯粹出于厚颜无耻而反对公认的事实。我们所说的对任何人都很清楚：奴役特有的属性标记着本性上为奴的，而被造则是奴役的适当属性。因此，宣称祂虽然是仆人却取了我们的形体的人，肯定就是那个将独生子从一个奴役转到另一个奴役的人。}\par

然而，他试图攻击我们的话，并在稍后说（我现在略过他的中间评论，因为它们在我先前的论证中已或多或少充分讨论过），当他指责我们\Quote{大胆地说或想不可设计之事，}并称我们为\Quote{最可悲的人，}——我再说，他补充道：——\Quote{因为如果圣伯多禄所说的不是那在起初就已存在且是天主的圣言，而是那被\InnerQuote{看见}的、\InnerQuote{空虚了自己}的那一位，如巴西略所说，并且如果那被\InnerQuote{看见}的人\InnerQuote{空虚了自己}以取\InnerQuote{仆人的形象}，而那\InnerQuote{空虚了自己}以取\InnerQuote{仆人的形象}的那一位\InnerQuote{空虚了自己}而成为人，那么那被\InnerQuote{看见}的人便\InnerQuote{空虚了自己}而成为人。}或许我的读者早已从上述引文中立即识破了他所坚持的论点的奸诈和愚蠢：然而，我们仍将附上对他所说之话的简短驳斥，与其说是为了推翻他拙劣的诡辩——这诡辩对有耳可听的人而言本身已不攻自破——不如说是为了避免留下我们不讨论他的指控，假装因鄙视其论点的无价值而不予理睬。因此，让我们这样看问题。宗徒的话是什么？\Quote{你们应当知道，}他说，\Quote{天主立祂为主和基督。}\BibleRef{宗 2:36}然后，好像有人问他，这样的恩宠赐给了谁，他仿佛用手指指着对象说：\Quote{这位耶稣，就是你们所钉死的。}巴西略对此怎么说？他指出，指示词宣告\Emph{那个}人被立为基督，就是被听众钉死的那位——因为他说：\Quote{你们钉死了，}而且很可能那些要求杀害祂的人就是这次演讲的听众；因为从钉死到伯多禄讲话的时间并不长。那么，欧诺弥对此又拿出什么回答？\Quote{如果圣伯多禄所说的不是那在起初就已存在且是天主的圣言，而是那被\InnerQuote{看见}的、\InnerQuote{空虚了自己}的那一位，如巴西略所说，并且如果那被\InnerQuote{看见}的人\InnerQuote{空虚了自己}以取\InnerQuote{仆人的形象}——}且慢！谁说那被看见的人再次空虚自己以取仆人的形象？或者谁主张十字架的苦难发生在肉身显现之前？十字架并未先于身体，身体也未先于\Quote{仆人的形象。}然而，天主在肉身显现；而那个在其自身内彰显了天主的肉身，在亲自完成了死亡的伟大奥秘之后，藉混合转化成为那崇高而神圣的，成为基督和主，被转移和改变为祂原先所是的，即那在肉身中显现自己的那一位。但如果我们这样说，我们真理的拥护者又坚持说，我们认为那位被显示在十字架上的祂\Quote{空虚了自己}成为另一个人，并将他的诡辩措辞如下：\Quote{如果，}他说，\Quote{那被\InnerQuote{看见}的人\InnerQuote{空虚了自己}以取\InnerQuote{仆人的形象}，而那\InnerQuote{空虚了自己}以取\InnerQuote{仆人的形象}的那一位\InnerQuote{空虚了自己}而成为人，那么那被\InnerQuote{看见}的人\InnerQuote{空虚了自己}而成为人。}\par

他多么牢记眼前的任务！他的论证结论多么切中要点！巴西略宣称，宗徒说那\Quote{被看见的}人被立为基督和主，而这位置清词敏、推翻他人陈述的人却说：\Quote{如果伯多禄并没有说那\InnerQuote{在起初就有的那位}本质被造，那么\InnerQuote{被看见的}人\InnerQuote{空虚了自己}以取得\InnerQuote{奴仆的形象}，而那\InnerQuote{空虚了自己}以取得\InnerQuote{奴仆的形象}的祂，空虚自己成了人。}欧诺弥，我们被这无敌的智慧征服了！宗徒的言论是指那\Quote{因软弱而被钉在十字架上的}\BibleRef{格后 13:4}，当我们得知若我们如此相信，那\Quote{被看见的}人就成了另一位，再次为人的新生成而\Quote{空虚自己}时，这一事实岂非被有力地驳倒了？你难道永不停止嘲弄那些本应免受此类攻击之事吗？你难道不因以如此荒谬的诡辩摧毁围护神圣奥迹的敬畏而脸红吗？你难道现在还不——如果从未——转向认识：那在圣父怀中的独生子天主，祂是圣言、是君王、是主，以及一切言语和思想所高举的，并不需要\Emph{成为}任何善的事物，因为祂自身就是万善的满全？那么，祂通过变为原来不是的东西而成其所不是的，那是什么呢？看，正如那不知罪的成了罪，为要除去世人的罪，同样，另一边，那接受了主的肉躯成了基督和主，通过混合而转变为它本性所不是的；由此我们学到：若非圣言成了血肉，天主也不会在肉身中显现；若非那感官可见的成了基督和主，那环绕祂的人性肉身也不会转变为神圣。但那些用三段论践踏天主的存有、并试图证明那藉创造使万有存在者本身是受造的一部分的人，鄙视我们所宣讲的纯朴，并曲解伯多禄对犹太人的话以助其建立亵渎：\Quote{以色列全家应确知，天主已立祂为主、为基督，就是你们钉在十字架上的这位耶稣}\BibleRef{宗 2:36}。这就是他们提出来支持独生子天主的本质是受造的论断的证据！什么？告诉我，被说这些话的犹太人，在万世之前就已存在吗？十字架在世界之前吗？比拉多先于一切受造吗？耶稣先存在，而后才有圣言吗？肉躯比天主性更古老吗？加俾额尔在世界之前向玛利亚传报喜讯吗？那在基督内的人不是在西泽奥古斯都的日子才开始通过诞生而有的吗？而作为圣言的那位，在起初就是天主，正如先知所证，是万世之前我们的君王？你难道看不出你给这事带来何等混乱，如同俗语所说的，颠倒是非吗？那是受难后第五十天，伯多禄向犹太人讲道说：\Quote{你们所钉死的这位，天主立祂为基督和主。}你难道没有留意他说话的顺序？他的话中哪个在先，哪个在后？他没有说：\Quote{天主所立为主的那位，你们钉死了}，而是说：\Quote{你们所钉死的这位，天主立祂为基督和主}；由此可见，伯多禄所说的不是万世之前的事，而是降生奥迹之后的事。\par

那么，你怎么会看不见你关于这个问题的整个论点正被推翻，却继续用你幼稚的诡辩之网使自己可笑，说如果我们相信那感官所显明的那位已被天主立为基督和主，那么主必然再次\Quote{空虚自己}重新成为人，并经历第二次诞生？你的教义从中得到什么好处？你所说的话如何表明造物之王是被造的？就我而言，我断言另一方面，我们的观点得到那些反对我们的人的支持，而这位修辞学家，由于对这件事的过分关注，未能看到当他如他所想将论证推向荒谬时，他正用他用来推翻对手的武器，站在他所攻击的人一边。因为如果我们相信耶稣的情况中的状态变化是从崇高到卑微，而只有神圣和非受造的本性超越受造物，那么他也许在彻底审视自己的论证时，会转向真理的阵营，并同意非受造者因其对人的爱而成为在受造者中。但如果他想象他通过表明祂，身为天主，参与了人性，来证明主是受造的，他会找到许多类似的经文来确立相同的意见，这些经文以类似的方式支持他的论点。因为祂既然是圣言，就是天主，并且\Quote{后来}，如先知所说，\BibleQuote{在地上显现，与人交谈}\BibleRef{巴 3:37}，祂将由此被证明是受造物之一！而如果这被认为离题，那么类似的经文也与主题不太相合。因为在意义上，说\Quote{在起初的圣言通过肉身显现给人}与说\Quote{祂本有\OriginalTerm{天主的形象}{form of God}，却取了\OriginalTerm{奴仆的形象}{form of a servant}}是完全相同的：并且如果这两种说法之一对他的亵渎之成立毫无帮助，那么他也必须放弃另一种。然而，他好心劝我们放弃我们的错误，并指出他自己所坚持的真理。他告诉我们，宗徒伯多禄宣告那在起初就是圣言和天主的祂被造成了。好吧，如果他是在编造梦境供我们消遣，并给我们讲关于睡眠异象的先知解释，那么允许他随意讲述他想象的谜语或许没有风险。但当他告诉我们他是在解释神圣言词时，我们不再安全地任由他随意解释这些话语。圣经说了什么？\BibleQuote{天主立为主和基督的，就是你们钉死的这个耶稣}\BibleRef{宗 2:36}。因此，当一切都一致时——指示性的词指明那以祂的人性之名被谈论的祂，对染有血罪之人的指控，十字架的苦难——我们的思想必然转向那感官所显明的。但他断言，当伯多禄使用这些词时，\Quote{立了}这个词所指示的是先于时间的存在。我们可以放心地让保姆和老妇人跟孩子开玩笑，并按自己的意思解释梦的含义；但当默感的圣经摆在我们面前要我们解释时，伟大的宗徒禁止我们求助于老妇的无稽之谈。当我听到谈论\Quote{十字架}时，我理解十字架；当我听到提到一个人名时，我理解那个人名所暗示的本性。所以当我从伯多禄那里听到\Quote{这个}人被立为主和基督时，我不怀疑他是在谈论那曾经在世人眼前的那位，因为圣人们在这一点上也彼此一致，如同在其他事上一样。因为正如他说那被钉十字架者已被立为主，保禄也说祂在受难和复活后\Quote{被高举}，并非祂身为天主而被高举。有什么高度比神性的高度更崇高，以致他说天主被高举到那里？而是他意指人性的卑微被高举了，我想这个词表示被所取的人与神性本性的崇高状态的同化和结合。即使有人允许他肆意曲解神性的言词，即使如此，他的论证也不会符合他的异端目的。因为就算承认伯多禄\Emph{确实}关于那在起初的祂说了：\BibleQuote{天主立为主和基督的，就是你们钉死的这个耶稣}，我们也会发现即使如此，他的亵渎也不会获得任何反对真理的力量。\Quote{天主立了祂}，他说，\Quote{为主和基督。}我们要将\Quote{立了}这个词归属于哪个词？与句子中所用的哪个词连接？我们有三个词：\Quote{这个}、\Quote{主}和\Quote{基督}。他将用这三个中的哪一个来构造\Quote{立了}这个词？没有人会如此大胆反对真理，以至于否认\Quote{立了}与\Quote{基督}和\Quote{主}有关；因为伯多禄说，祂（既然已经是祂所是的无论什么）被圣父\Quote{立为基督和主}。\par

这些话不是我的；它们出自那对抗圣言者之口。因为他就在我们面前正在审查的这段文字中，确切地这样说：——\Quote{真福伯多禄论及那起初就存在且是天主的那一位，向我们解释正是祂成了主和基督。}于是，欧诺米说，那原本就是其所是者成了主和基督，正如达味的历史告诉我们，他作为叶瑟的儿子，又是牧羊人，被傅油为王：并非那傅油使他成为人，而是他按其本性所是，从普通人转变为王。接下来呢？难道就因此更能确定圣子的本体是受造的，如果如欧诺米所说，天主使那起初存在且是天主的那一位成了主和基督？因为主位不是祂\Emph{存在}的名称，而是祂\Emph{拥有权威}\Emph{的存在}之名，而基督的称谓指明祂的王国，而祂的王国概念是一回事，祂的本性概念是另一回事。假定圣经确实说这些事发生在天主圣子身上。那么让我们考虑哪个观点更为虔诚、更为合理。我们容许说哪一个经由进步而成为优越的分享者——天主还是人？谁会有如此幼稚的心思，以至于认为神性经由添加而达至成全？但就人性而言，这样的假设并非不合理，因为福音的话清楚地讲我们的主在人性方面的增长归于祂：因为经上说，\BibleQuote{耶稣在智慧和身量上，并在天主和人前的恩宠上，渐渐地增长。}那么，从宗徒的话中引申出哪个建议更为合理呢？——那在起初就是天主的那一位经由进步而成了主，还是人性的谦卑因其与神性的共融而被提升至尊荣的高峰？因为先知达味也以主的口吻说，\BibleQuote{我被祂立为王，}其含义非常接近\Quote{我被造成基督：}再者，他以圣父对主的口吻说，\BibleQuote{在祢的仇敌中为主，}其含义与伯多禄所说的\Quote{成为祢仇敌的主}相同。因此，正如祂王国的建立并不表示祂本体的形成，而是祂尊荣的提升，而那命令祂\Quote{为主}的那一位并非在那特定时刻吩咐不存在者开始存在，而是将统治权赐予那已存在者，使其统治那些不顺从者——同样，真福伯多禄，当他说某人被造成基督（即万有的王）时，加了\Quote{祂}字，以将这个概念与本体以及与之相关的属性区分开来。因为祂使之成为所宣告的那一位时，祂已经是其所是的了。如今，如果我们可以断言那超越的本性经由进步而成为任何事物，如从普通人成为王，或从卑微成为崇高，或从仆役成为主，那么将伯多禄的话应用于独生子或许是恰当的。但既然神性——无论它被相信是什么——常是同一，超越一切增补且不能被减少，我们绝对必须将他的话归于人性。因为天主圣言现在且永远保持为起初之所是，常是王，常是主，常是天主和至高者，并非经由进步而成为这些的任何一种，而是凭其本性就是祂所被宣告的一切；另一方面，那通过被摄取而从人提升至神性者，原是此物而成为彼物，严格而真正地被称为成了基督和主。因为祂使祂从仆役成为主，从臣民成为王，从隶属成为基督。祂极度高举了那卑微的，并将那超乎万名之上的名赐予那拥有人的名字者。于是成就了那不可言喻的混合与结合，人的渺小与神性的伟大相混，以致那些伟大而神性的名字被适当地应用于人性，另一方面，神性也被以人的名字称呼。因为同一位既拥有那超乎万名之上的名，又在耶稣的人性之名中被所有受造物朝拜。因为他说，\BibleQuote{致使上天、地上和地下的一切，一听到耶稣的名字，无不屈膝叩拜，一切唇舌无不明认耶稣基督是主，以光荣天主圣父。}然而关于这些事已经说得够多了。\par

\SourceRecord{Translated by H.A. Wilson.；Nicene and Post-Nicene Fathers, Second Series；Vol. 5.；Edited by Philip Schaff and Henry Wace.；Buffalo, NY: Christian Literature Publishing Co.,；1893.；<http://www.newadvent.org/fathers/290106.htm>.}

% structure: p0001 source=h1 target=section reason=page-title

\section{\WorkTitle{驳斥欧诺弥（卷七）}}

% structure: p0002 source=h2 target=subsection reason=relative-heading-rank; base=h2

\subsection{一、第七卷通过致格林多人书和希伯来人书中的不同论述，以及主的话，表明\Quote{主}这个词按欧诺弥的解释，并非表达本质，而是表达尊位。并在许多关于\Quote{圣神}和主的著名评论之后，他表明欧诺弥从自己的话中，虽然无意，却被发现支持正统信仰，并被自己的标枪击中。}

然而，既然欧诺弥断言\Quote{主}这个词是用来指称独生子的\Emph{本质}，而非其\Emph{尊位}，并引用宗徒向格林多人所说的话作为这一观点的见证：\BibleQuote{主就是那圣神}\BibleRef{格后 3:17}，那么现在或许是合适的时候，我们不应放过他这方面的错误而不加纠正。他断言\Quote{主}这个词表示本质，并且为证明这一假设，他引用了上述经文。这经文说：\Quote{主就是那圣神\BibleRef{格后 3:17}}。但是我们的这位朋友，即凭己意解释圣经者，将\Quote{主权}称为\Quote{\Emph{本质}}，并以为透过所引用的经文可以证明他的说法。好吧，如果保禄说过\Quote{主就是本质}，我们也会同意他的论证。但是鉴于受默感的著作一方面说\Quote{主就是那圣神}，而欧诺弥另一方面说\Quote{主权就是本质}，我不知道他从哪里找到支持他声明的依据，除非他准备再说\Quote{圣神}一词在圣经中代表\Quote{本质}。那么，让我们思考，宗徒在任何地方使用\Quote{圣神}一词时，是否用它来指示\Quote{本质}。他说：\BibleQuote{圣神亲自与我们的心神一同作证}\BibleRef{罗 8:16}，和\BibleQuote{除了人里面的心神，谁也不知道人的事}\BibleRef{格前 2:11}，和\BibleQuote{文字叫人死，圣神却叫人活}\BibleRef{格后 3:6}，和\BibleQuote{如果你们靠着圣神治死身体的恶行，你们必要生活}\BibleRef{罗 8:13}，和\BibleQuote{如果我们生活在圣神内，也当随从圣神而行}\BibleRef{迦 5:25}。谁能数得清宗徒在这方面的论述呢？在这些话中，我们找不到这个字表示\Quote{本质}。因为说\Quote{圣神亲自与我们的心神一同作证}的人，所指的无非是圣神，祂降临在信徒的心里；因为在他的著作的许多其他段落中，他将心神称为精神，当心神领受了圣神的共融时，领受者便获得义子的尊位。再者，在\Quote{除了人里面的心神，谁也不知道人的事}这段经文中，如果\Quote{人}是用来指本质，而\Quote{心神}也一样，那么从这话就会推出，人被认为有两个本质。再者，我不知道说\Quote{文字叫人死，圣神却叫人活}的人，如何将\Quote{本质}与\Quote{文字}对立；也不知道这位作者如何想象，当保禄说我们应当\Quote{靠着圣神}治死\Quote{身体的恶行}时，他正引导\Quote{圣神}一词的表达指向\Quote{本质}；至于\Quote{生活在圣神内}和\Quote{随从圣神而行}，如果\Quote{圣神}一词的意义是指\Quote{本质}，那么这些话就完全不可理解了。因为我们所有活着的人，若不是在本质中，又在哪里分享生命呢？——因此，当宗徒就此事给我们提出建议，说我们应当\Quote{生活在本质中}时，就好像他说了\Quote{靠你们自己、而非靠别人来分享生命}。如果这样的理解在任何经文中都不可能成立，那么欧诺弥怎能再次效法解梦的人，吩咐我们将\Quote{圣神}当作\Quote{本质}，以便他能以正当的三段论形式得出\Quote{主}这个词适用于本质的结论呢？——因为如果\Quote{圣神}是\Quote{本质}（他论证说），而\Quote{主就是那圣神}，那么显然\Quote{主}就是\Quote{本质}。这个尝试的力量多么无可辩驳！我们怎能躲避或化解这无可辩驳的论证必然性呢？他说，\Quote{主}这个词是谈论本质的。他如何维护这一点？因为宗徒说：\Quote{主就是那圣神}。那么，这与本质有什么关系？他进一步教导我们，\Quote{圣神}被用来代替\Quote{本质}。这些就是他论证方法的技巧！这些就是他的亚里士多德科学的成果！这就是为什么，按你们的看法，我们这些未入门这种智慧的人是如此可怜！而你们当然应被视为有福的，你们用这样一种方法追踪真理——也就是说，宗徒的意思是这样，以至于我们必须假设\Quote{圣神}被他用来代替独生子的本质！\par

那么你如何使之与下文相符呢？因为当保禄说\BibleQuote{现在主就是圣神}，他接着说\BibleQuote{哪里有主的圣神，那里就有自由}。如果\BibleQuote{主就是圣神}，并且\Quote{圣神}意指\Quote{本质}，那么\Quote{本质的本质}我们该作何理解？他又提到另一个主的圣神，即那圣神——也就是说，根据你的解释，是另一个\Emph{本质}。因此，依你之见，宗徒在明写\BibleQuote{主圣神}和\BibleQuote{主的圣神}时，所指的无非是本质的本质。好吧，让\Person{欧诺弥奥}随意理解所写的内容；我们对事情的看法如下。圣经，正如宗徒所称，是\BibleQuote{由天主所默感}的，是圣神的圣经，其用意是为人的益处。因为\BibleQuote{凡圣经，}他说，\BibleQuote{都是由天主所默感，并且是有益的}；而益处是多种多样的，正如宗徒所说——\BibleQuote{为教训、为督责、为改正、为教导人学正义}\BibleRef{弟后 3:16}。然而，这样的恩惠并非任何人所能把握，但天主的旨意隐藏在圣经的身体之下，如同藏在帕子之下，有些立法规定或历史叙述覆盖在心灵所默观的真理之上。因此，宗徒告诉我们，那些看圣经身体的人，有\BibleQuote{帕子蒙在心上}\BibleRef{格后 3:15}，不能看见属神法律的荣光，因被盖在立法者面上的帕子所阻碍。因此他说：\BibleQuote{文字叫人死，圣神却叫人活}，表明通常明显的解释，若不按适当的意义来理解，会产生与圣神所指的生命相反的效果，因为圣神为众人定下在无激情中的完美德行，而书中所含的历史有时甚至包含不一致的事实的阐述，并且可以说，被理解为与我们本性的激情一致，如果有人按明显的意思应用自己，他将使圣经成为死亡的道理。因此，他说，对于以过于肉体方式处理所写之事的人，帕子蒙在他们灵魂的感知能力上；但对于那些将默观转向理智对象的人，就有如揭下面具，启示出隐藏在文字底下的荣光。而这由更高超的知觉所发现的东西，他说就是主，即圣神。因为他说的：\BibleQuote{几时它转向主，帕子就除去了：现在主就是圣神}\BibleRef{格后 3:16}。这样说，他就在圣神的主权和文字的束缚之间作了对比；因为正如赐生命的与杀人的相对，他也将\Quote{主}与束缚相对。并且为了使我们受教关于圣神时不致混乱（因被\Quote{主}这个词引向对独生者的思想），他以此缘故通过重复来保守护这个词，既说\BibleQuote{主是圣神}，又进一步提到\BibleQuote{主的圣神}，以便祂本性的至高性可由统治所隐含的尊荣显明，同时避免在论证中混淆祂位格的个别性。因为那既称祂为\Quote{主}又为\BibleQuote{主的圣神}的，教导我们将祂视为独生者之外的一个独立个体；正如别处他提到\BibleQuote{基督的圣神}\BibleRef{罗 8:9}，恰当地并在其奥秘意义上使用这个按照福音传统虔诚用于教义体系的术语。因此我们，\Quote{众人中最可怜的人}，由宗徒在奥秘中引领，从杀人的文字转向赐生命的圣神，从那位在乐园中领受了不可言传的奥秘的那一位学习，得知一切圣经所说的都是圣神的话语。因为\BibleQuote{圣神藉先知以赛亚向你们的祖宗所说的话是不错的}\BibleRef{宗 28:25}——这是他对在罗马的犹太人说的，引入了依撒意亚的话；又对希伯来人，援引圣神的权威，说：\BibleQuote{所以，正如圣神所说的}\BibleRef{希 3:7}，他引用了圣咏中以上主的口吻所说的话。并且我们从主自己学习同样的事——达味宣布天上的奥秘不是\Quote{凭}自己（即不是按人性说话）。因为一个人，既只是人，怎能知道天父与圣子之间的超天对话呢？但他是\Quote{在圣神内}，他说主对主说了那些祂所说的话。因为，祂说，\BibleQuote{达味\Emph{在圣神内}称他为主，他怎么又是他的儿子呢？}因此，藉着圣神的德能，受天主感动的圣人们被默感，而每部圣经因此被称为\BibleQuote{由天主所默感}，因为它是天主所嘘的教导。如果言语的肉体帕子被除去，所剩下的就是主、生命和圣神，按照伟大的保禄的教导，也按照福音的话。因为保禄宣告，那从文字转向圣神的人，不再领会杀人的束缚，而是领会那赐生命的圣神；崇高的福音说：\BibleQuote{我对你们所说的话，就是圣神，就是生命}，因脱去了肉体的帕子。然而，认为\Quote{圣神}是独生者本质的想法，我们留给我们的做梦者：或者不如说，我们要利用他们所说的，\Emph{ex abundanti}，用敌人的武器装备真理。因为埃及人被以色列人掠夺是允许的，我们应当使他们的财富成为我们的装饰。如果子的本质被称为\Quote{圣神}，而天主也是圣神（因为福音这样告诉我们），显然圣父的本质也被称为\Quote{圣神}。但如果他们的特别论点是：由不同名称引入的事物在本质上也是不同的，那么结论必然是：称呼相似的事物在本质上也不是彼此相异的。因此，按照他们的说法，圣父的本质和圣子的本质都称为\Quote{圣神}，这就清楚地证明了在本质上没有任何差异。因为稍后欧诺弥奥说：\Quote{在那些发散的本质上，表示本质的名称也必定是发散的，但倘若有一个且同一个名称，那由同一称号所宣告的也必定是同一个}：——所以处处\BibleQuote{祂以智慧俘获智者于其诡计中}，已将我们作者的长久辛劳和用于其推敲的无限努力，转化为对我们所持道理的建立。因为如果天主在福音中被称为\Quote{圣神}，而独生者的本质被欧诺弥奥主张为\Quote{圣神}，既然一个名称与另一个没有明显区别，那么毫无疑问，名称所意指的事物在本质上也不会相互不同。\par

如今，我已揭穿这虚妄而无谓的伪论证，似乎可以略过不议他随后拼凑出来攻击我们师长言论的那些东西。因为他的言论之愚昧，已由其论证本身充分证明，这论证自行高呼其软弱无力。纠缠于这等事物之争，犹如践踏亡者。当他满怀信心地提出我们师长的某段言论，先加以诽谤和轻蔑，然后许诺要证明它毫无价值时，他遭遇的正是年幼孩童的命运——他们不完善、不成熟的智力和未经训练的感知能力，不能准确理解所见之物。因此，他们常以为自己头顶的星辰离得很近，在幼稚的愚昧中，看见星辰便用土块投掷；土块落下时，他们拍手大笑，向同伴夸耀，仿佛他们的投掷真的击中了星辰。此人正是如此：用他那孩子气的投射物掷向真理，陈列我们师长那些如星辰般璀璨的言论，然后从他那低贱而匍匐于地的理解力中，抛出他那属世的、不稳固的论证。这些论证上升到无处可坠的高度，又因自身重量自行折返。伟大的巴西略的那段话是这样写的：\newline{}\par

\Quote{然而，有哪个心智正常的人会同意这个说法：名称不同的事物，其本质也必定相异？因为伯多禄和保禄的称谓是不同的，一般而言，人的称谓也是不同的，但所有人的本质却是同一个；因此，我们在大多数方面彼此相同，仅在个体中观察到的那些特殊属性上彼此相异；所以，称谓并不指示本质，而是指示标示特定个体的属性。因此，当我们听到伯多禄时，我们并非通过名称理解本质（我在此所说的\InnerQuote{本质}是指物质基体），而是被我们在他身上所观察到的属性的概念所打动。}这是那位伟人的话。而针对这个说法进行争论的人对我们显示了何等技巧，我们——任何有闲暇在无益之事上浪费时间的人——从欧诺米的实际作品中便可得知。\par

从他的著作，我说，因为我不喜欢在自己作品中插入这位修辞家所吐的令人作呕的东西，或是在我的论证中展示他的无知和愚蠢以供鄙视。他继续对表达主题的那类有意义词语大加赞美，并以他惯常的风格，拼凑粘贴一些破旧词句的碎屑：可怜的\Person{伊索克拉底}再次被蚕食，被剪去词语和辞藻以支持所提出的论点——甚至希伯来人\Person{斐洛}也在各处遭受同样待遇，并提供他本人作品中的短语——然而即便如此，这篇缝缀多色之网仍未完成，反而是他的每一个攻击、每一个防御、所有艺术准备都自行崩溃，就像通常气泡那样，当水珠从上穿过水体撞击障碍物时，产生那些泡沫般的鼓胀，刚一形成便立即消散，在水面上不留下任何自身形成的痕迹——这位作者的思想气泡也是如此，刚一发出便不触而灭。因为在这所有无可辩驳的论述和他那梦幻般的哲学推理之后——其中他断言本质的独特特征是通过名称的差异来把握的——就像一团顺流而下的泡沫撞上任何更坚固的物体时就会破裂，他的论证也顺其自然，意外地与真理相撞，并将其虚幻的、气泡般的谬误结构化为虚无。因为他这样说道： \Quote{谁如此愚蠢，如此远离人类的构造，以致在谈论人时，竟将一人称为人，而将另一人称为马，如此比较？} 我将回答他： \Quote{你把使用名称时犯这种错误的人称为愚蠢，这很正确。我将用你亲自给出的证言来支持真理。因为，如果假设二者都是人，却将一人称为马、另一人称为人，这极其愚蠢；那么，当圣父被宣认为天主，圣子被宣认为天主时，将一方称为\InnerQuote{受造}，另一方称为\InnerQuote{非受造}，同样也是愚蠢的，因为正如在另一情形中人性不接纳表示另一种类的名称变化，在此情形中天主性也不接纳。因为非理性相对于人如何，受造物相对于天主性也如何，同样不能接受与超越它的本性相同的名称。正如不可能对理性动物和四足动物应用相同的定义（因为每一方都因其特殊属性而与另一方自然不同），同样你也无法用相同的术语来表达受造的本质与非受造的本质，因为后面那种本质的属性在前面那种中找不到。因为正如理性在马中找不到，蹄的坚硬在人中找不到，同样天主性在受造物中也找不到，受造的性质在天主性中也找不到：但如果他是天主，他当然不是受造的，如果他是受造的，他就不是天主；除非有人通过某种误用或习惯表达方式，仅仅应用天主性的名称，就像一些马被主人赋予人的名字；然而，马并不是人，尽管它被冠以人的名字，受造物也不是天主，即使有人为他声称天主性的名字，并给他一个双音节空响的便宜。} 既然\Person{欧诺米}的异端言论被发现自发地符合真理，让他听从自己的劝告，坚持自己的话，决不收回自己的言论，而应认识到，那不以事物本相来命名主体，却将\Quote{马}说成\Quote{人}，将\Quote{海}说成\Quote{天}，将\Quote{受造物}说成\Quote{天主}的人，确实是愚妄愚蠢的。并且不要有人认为将受造物与天主对立是不合理的，而要留意先知和宗徒。因为先知以圣父的身份说： \Quote{我的手造了这一切}，在他的隐晦说法中，以\Quote{手}意指独生子的能力。现在宗徒说万有来自圣父，万有藉圣子而成\BibleRef{格前 8:6}，先知的精神在某种程度上与宗徒的教导一致，宗徒的教导也是藉圣神赐予的。因为在一处，先知说万有都是那超越万有者的手的工作，他在与造物者的关系中展示了那些已存在之物的本质，而造物者是超越万有的天主，他拥有手，并藉手造了一切。在另一处，宗徒对实体做了同样的划分，使万有依赖于它们的产生原因，却没有将产生者算在\Quote{万有}之中：这样我们便由此得知受造与非受造之间的本质区别，并表明在其本性中，造物者是一回事，被产生者是另一回事。既然万有来自天主，而圣子是天主，受造界便恰当地与天主性相对立；既然独生子与宇宙的本性不同（甚至那些与真理争战的人也不否认这一点），那么圣子也必然同样与受造界对立，除非圣人们的证言——他们证明万有是藉他造成的——是虚假的。\par

% structure: p0008 source=h2 target=subsection reason=relative-heading-rank; base=h2

\subsection{二、他随即宣称名称与事物之间的密切关系是不可改变的，然后以最卓越的方式，继续讨论关于\Quote{受生}与\Quote{非受生}的论述。}

既然在圣经中独生者被宣称为天主，就让欧诺米考虑他自己的论证，并判定那将天主性分为受造与不受造的人为极其愚蠢，正如他判定那将\Quote{人}分为\Quote{马}与\Quote{人}的人一样。因为他在稍后、在他的中间废话之后，自己说道：\Quote{名称与事物的密切关系是不可改变的}，他本人借此声明同意称谓与其主体之真正联系的固定特性。那么，如果天主性的名称被适当地用于与独生者天主密切相连（而欧诺米，尽管他可能想要与我们不合，但必定会承认圣经并不说谎，且天主性的名称并非不协调地归于独生者），就让他以他自己的推理说服自己：如果\Quote{名称与事物的密切关系是不可改变的}，并且主被称为\Quote{天主}之名，他就不能领悟在父与子之间在天主性概念上的任何差异，因为这名称为两者所共有——或者更确切地说，不仅这个名称，还有一长串名称，子在其中毫无意义分歧地分享了父的称谓——\Quote{善}、\Quote{不朽}、\Quote{公义}、\Quote{审判者}、\Quote{宽忍}、\Quote{仁慈}、\Quote{永恒}、\Quote{永久}，所有那些表明本性之尊威与能力之表达的称谓——在子身上，就尊高概念的本性而言，没有任何称谓是有所保留的。但欧诺米却像是闭着眼睛一样，忽视了大量神圣的称谓，只盯住一点，他的\Quote{受生与不受生}——将其教义托付给一根细弱易断的绳索，它被错谬的狂风吹得飘摇不定。\par

他断言：\Quote{没有一个尊重真理的人会把任何受生之物称为\InnerQuote{不受生}，也不会称那统御万有的天主为\InnerQuote{子}或\InnerQuote{受生}。}这句话本身无需我们再做更多的论证来反驳。因为他不像往常那样用任何面纱遮掩其狡诈，而是将其荒谬陈述的颠倒视为等同，他说：既没有任何受生之物被称为\InnerQuote{不受生}，也没有统御万有的天主被称为\InnerQuote{子}或\InnerQuote{受生}，而并没有为子的独生者天主性与其他\InnerQuote{受生之物}作出任何特别区分，而是不加区别地将\Quote{所有生成之物}与\Quote{天主}对立起来，且不将子从\Quote{所有生成之物}中排除。在他荒谬的颠倒中，他显然将子与天主性分离开来，他说：既没有任何受生之物被称为\InnerQuote{不受生}，也没有天主被称为\InnerQuote{子}或\InnerQuote{受生}，由此对比明显揭示出他亵渎的可怖性质。因为当他将\Quote{所有生成之物}与\Quote{不受生者}区分开来之后，在那反向归纳中，他接着说：不能称（不是\InnerQuote{不受生者}，而是）\Quote{天主}为\InnerQuote{子}或\InnerQuote{受生}，试图以此表明，不是不受生的就不是天主，并且独生者天主因受生的事实，距离是天主的程度正如不受生者距离实际上或名义上是受生的程度一样远。因为他并非不知道其论证的后果而做出这样不协调、不适当的词语颠倒：而是为了攻击正统教义，他将\Quote{天主性}与\Quote{受生者}对立起来——这正是他试图通过其言语确立的观点，即不是不受生的就不是天主。他的论证本应如何？即说过\Quote{没有受生之物是不受生的}之后，应该接着推断：\Quote{当然，如果某物在本性上是不受生的，它也不能是受生的}。这样的说法既包含真理，又避免亵渎。但现在，根据他的前提——没有受生之物是不受生的，以及他的推断——天主不是受生的，他明确地将独生者天主排除在天主之外，断定：因为祂不是不受生的，所以祂也不是天主。那么，我们还需要更多的证据来揭露这骇人的亵渎吗？这本身难道不是足够作为反对基督仇敌的见证吗？他通过所引用的论证坚持那起初与天主同在的圣言不是天主。与这样的人进一步争论又有何必要呢？因为我们不会与那些沉溺于偶像和在祭台上流血的人纠缠争论，并非我们默许那些执迷于偶像的人的毁灭，而是因为他们的病对我们的治疗来说太重了。因此，正如事实本身宣告偶像崇拜，而人胆大妄为的恶行先于控诉者的谴责，同样，我认为正统教义的拥护者也应该对那位公开宣扬不敬虔以致自取其辱的人保持沉默，就像医学在面对癌症时也无能为力，因为疾病已经超出了医学技艺所能处理的范围。\par

% structure: p0011 source=h2 target=subsection reason=relative-heading-rank; base=h2

\subsection{三、此后，他讨论名称与事物的差异，论及那不受生者是无原因的，以及那不存在者，如\OriginalTerm{辛达普苏斯}{Scindapsus}、\OriginalTerm{弥诺陶洛斯}{Minotaur}、\OriginalTerm{布里提里}{Blityri}、\OriginalTerm{独眼巨人}{Cyclops}、\OriginalTerm{斯库拉}{Scylla}，这些从未受生过，并表明本质上不同的事物相互毁灭，如火与水，以及其他各自的关系。但在父与子那里，由于本质是共同的，属性相互可交换，本性不受损害。}

然而，在引用了上述段落之后，他声称将提出更强有力的证据，让我们也来审视这一点以及所引段落，以免我们似乎在面对压倒性的力量时退却。\Quote{然而，如果}他说，\Quote{我要放弃所有这些立场，退回到我更有力的论点，我会说：即使他提出的所有反驳术语都成立，我们的陈述仍将明显显示为真。如果，如所承认的，标示特性的名字的差异标示事物本身的差异，那么当标示本质的名字有差异时，必然也必须承认本质也有差异。而且这在所有情况下都成立，我的意思是本质、能力、颜色、形状和其他性质。因为不同的本质我们用不同的称呼来标示，如火与水、气与土、冷与热、白与黑、三角形与圆形。何须提及可理解的本性呢？在列举它们时，宗徒正是通过名字的差异标示了本质的差异。}\par

谁不会因这不可抗拒的攻击力量而感到沮丧？论证超越了应许，经验比威胁更可怕。\Quote{我要来，}他说，\Quote{转向我更强有力的论证。}这是什么？既然属性的差异是通过标示特殊属性的名字被认识的，他说，当然我们必须承认，本质的差异也是通过名字的差异来表达的。那么，这些标示本质的称谓是什么？通过它们我们了解圣父和圣子之间本性的差异？他谈到火和水、气和土、冷和热、白和黑、三角形和圆形。他的例子使他赢了：他的论证无往不利：我不能反驳那个说法，即那些完全不可共通的名称指示本性的差异。但我们这位敏锐而目光锐利的智者却忽略了这些要点：即在这个例子中，圣父是天主，圣子是天主；\Quote{正义的}和\Quote{不朽的}，以及所有属于天主性体的名字，相同地用于圣父和圣子；因此，如果称谓的差异特性指示本性的差异，那么名字的共同性就必然显示本质的共通特性。如果我们必须同意，天主性体是通过名字来表达的，那么我们应该把那些崇高而神圣的名字应用于那个本性，而不是\Quote{\OriginalTerm{受生}{generate}}和\Quote{\OriginalTerm{非受生}{ungenerate}}这类术语，因为\Quote{善的}和\Quote{不朽的}、\Quote{正义的}和\Quote{智慧的}，以及所有这类术语只严格适用于那超越一切理解的本性，而\Quote{受生的}却显示出名称的共通性，甚至与受造物中较卑微的形式相同。因为我们称一条狗、一只青蛙、以及一切通过生殖进入世界的受造物为\Quote{受生的}。此外，\Quote{\OriginalTerm{非受生}{ungenerate}}一词不仅用于无原因存在的事物，而且还适当地应用于不存在的事物。\OriginalTerm{斯辛达普苏斯}{Scindapsus}被称为非受生的，\OriginalTerm{布利提里}{Blityri}是非受生的，\OriginalTerm{弥诺陶洛斯}{Minotaur}是非受生的，\OriginalTerm{独眼巨人}{Cyclops}、\OriginalTerm{斯库拉}{Scylla}、\OriginalTerm{喀迈拉}{Chimæra}都是非受生的，不是指无生成而存在，而是指从未生成。那么，如果更独特神圣的名字是圣子与圣父所共有的，而其他那些含糊地用于不存在物或低等动物的名字——我指的是差异的那些——就让他的\Quote{受生和非受生}如此吧：\Person{欧诺弥}对于我们有力的论证本身支持了真理的事业，它见证在本性方面没有差异，因为在名字中察觉不到差异。但如果他断言在\Quote{\OriginalTerm{受生}{generate}}和\Quote{\OriginalTerm{非受生}{ungenerate}}之间存在本质的差异，如同火与水之间的差异一样，并认为这些名字与他例子中提到的那些名字具有相同的相互关系，如同\Quote{火}和\Quote{水}一样，那么他那亵渎的可怖特征即使在我们保持沉默时也会再次显现。因为火和水具有相互毁灭的本性，如果二者相遇，强者将消灭弱者。那么，如果他提出这样的教义，即\OriginalTerm{非受生}{Ungenerate}的本性如此不同于\OriginalTerm{独生子}{Only-begotten}的本性，那么显然，他在逻辑上使这种毁灭性的对立包含在他们的本质差异中，因此，按照这个推理，它们的本性将是不相容且不可共通的，如果两者被发现相互包含或共存，一方将被另一方消灭。\par

那么，圣子如何\Quote{在圣父内}而不被消灭，圣父来到圣子内又如何持续不被消耗，如果如\Person{欧诺米}所说，火的特有属性相对于水，在\OriginalTerm{受生的}{Generate}与\OriginalTerm{不受生的}{Ungenerate}的关系中得以保持？他们的定义也不认为土与气之间存在相通，因为前者是稳定、坚固、抵抗、有下沉倾向且沉重的，而气则具有相反属性的本性。白色与黑色在颜色中是对立的，人们也同意圆与三角形不同，因为根据各自形状的定义，每一个正好是另一个所不是的。但我无法发现他在哪里看到了天主圣父与天主\OriginalTerm{独生子}{Only-begotten Son}之间的对立。同一个美善、智慧、公义、上智安排、大能、不朽——所有其他崇高意义的属性都同样赋于两者，并且一方在某种意义上在另一方内拥有其能力；因为圣父藉着圣子创造万有，而\OriginalTerm{独生子}{Only-begotten}作为圣父的能力，在自身内行事一切。那么，火与水在显示圣父与圣子的本质差异上又有何用？此外，他称我们为\Quote{鲁莽}，因为我们以\Person{伯多禄}和\Person{保禄}的本性合一与位格差异为例，并说如果我们借助物质实例将论点应用于纯粹理性对象的默观，我们便犯了极大的轻率之罪。的确，确然如此，纠正我们错误的人责备我们以物质比喻诠释天主性体是鲁莽！那么，深思熟虑而审慎的先生，你为何谈论元素？难道土是非物质的，火是纯粹理性的对象，水是无形的，气是超越感官知觉的吗？你的心智是否如此专注于目标，你在宣扬这个论证时是否如此眼光锐利，以致你的对手无法抓住，你却不看你自己身上有那些你指责他人时所责备的缺点？还是说，当你借助物质确立本质的差异时，我们应当对你让步，而当我们以我们力所能及的实例指出本性的同源特性时，我们自己却被拒绝？\par

% structure: p0015 source=h2 target=subsection reason=relative-heading-rank; base=h2

\subsection{四、他说，受造界的一切事物都由人命名，如果事实如此，每个民族对其称呼不同，正如我们赋予\Quote{不受生的}这个名称一样：但表达天主性体的神圣本质本身的专有名称，要么根本不存在，要么不为我们所知。}

但是，他说，\Person{伯多禄}和\Person{保禄}是由人取名的，因此他们的名字可以更改。然而，有什么存在的事物不是由人取名的呢？我请你为我的论证作证。因为如果你把名字的改变视为事物由人取名的标志，那么你必然也会承认，每一个名字都是我们加给事物的，因为对象的同一名称并未普遍通行。因为，就如\Person{保禄}曾经是\Person{扫禄}，\Person{伯多禄}从前是\Person{西满}，同样，地、天、空气、海以及创造的所有部分，并非所有人都以相同方式命名，而是\Place{希伯来}人以一种方式称呼，我们又以另一种方式称呼，并且每个民族都用不同的名称来指称。那么，如果\Person{欧诺米}的论证成立，他主张\Person{伯多禄}和\Person{保禄}之所以被重新取名，正是因为他们的名字是由人加给的，那么我们的教导也必然成立，因为它出发于同样的前提，即万物都是由我们命名的，因为它们的称呼因民族的不同而异。如果万物都是如此，那么\OriginalTerm{受生}{Generate}和\OriginalTerm{无始}{Ungenerate}当然也不例外，因为它们也属于名字会改变的事物。因为当我们把心中产生的对任何主体的概念，收集成一个名字的形式时，我们通过在不同时间变化的词语来宣告我们的概念，不是通过所给的名字来\Emph{制造}事物，而是\Emph{指明}事物。因为事物按其本性保持在自身内，而心灵触及存在的事物，通过可用的词语来揭示其思想。正如\Person{伯多禄}的本质没有随着他名字的改变而改变，同样，我们所思虑的任何其他事物也不会在名字变化的过程中改变。因此，我们说，\Quote{无始}这个术语是由我们用于那真实而第一位、作为万有之原因的圣父，并且如果我们承认同一个概念用另一个名字，对于指示主体来说也不会造成损害。因为，我们可以用\Quote{第一因}或\Quote{独生子之圣父}来称呼祂，而不用\Quote{无始}，或者称祂为\Quote{无因存在}，以及许多其他导向同一思想的称呼；因此，\Person{欧诺米}恰恰是在他抱怨我们的那些论证中证实了我们的教义，因为我们不知道任何能指示\OriginalTerm{天主性}{Divine Nature}的名字。我们获知祂存在的事实，同时断言，一个能够包含那不可言说和无限的天主性的称呼，要么根本不存在，要么至少不为我们所知。那么，让他放弃他惯用的神话语言，给我们展示那些指示本质的名字，然后通过名字的差异进一步区分主体。但是，只要圣经的话是真的：\Person{亚巴郎}和\Person{梅瑟}不能认识那名字，并且\Quote{从来没有人见过天主}，并且\Quote{没有人见过祂，也不能看见祂}\BibleRef{弟前 6:16}，而且祂周围的光是不可接近的\BibleRef{弟前 6:16}，并且\Quote{祂的伟大没有穷尽}——只要我们这样说并相信这些，那么，一个承诺通过名字的含义来理解和表达无限本质的论证，多么像一个人以为他可以用自己的手掌包围整个大海！因为手掌的凹陷相对于整个深渊，就如同语言的所有力量相对于那不可言说和不可理解的天主性一样。\par

% structure: p0017 source=h2 target=subsection reason=relative-heading-rank; base=h2

\subsection{五、在大量论述实际存在者、无始与善，以及天上力量的同体性之后，展示了其本质的不可测度及等级的区别，他结束了本书。}

如今我们说这些话，并非有意否认圣父\OriginalTerm{无生}{without generation}，也不打算拒绝同意\Quote{独生子天主是\OriginalTerm{受生}{generated}}这一说法——相反，后者是受生的，前者不是受生的。但是，那脱离受生而存在的祂，在本性上\Emph{是}什么，以及那被相信为受生的祂\Emph{是}什么，我们并非从\Quote{受生}和\Quote{\OriginalTerm{非受生}{ungenerate}}的涵义得知。因为当我们说\Quote{此人是受生的}（或\Quote{非受生的}）时，我们心中有两种思想：通过短语中的指示部分，我们注目于主体；通过\Quote{受生}或\Quote{非受生}这些话，我们得知在主体中所默观的东西。——因为思考那个\Quote{是者}是一回事，而思考我们在\Quote{是者}中所默观的则是另一回事。而且，用于谈论天主本性的每一个名称，无疑都隐含着\Quote{\Emph{是}}这个词——例如\Quote{公义}、\Quote{不朽}、\Quote{不死}、\Quote{非受生}，以及任何其他关于祂的说法；即使这个词并未恰好出现在短语中，说话者和听者的思想也必定使名称附着于\Quote{\Emph{是}}，因此，如果不加上这个词，称呼便徒劳无益。例如（因为最好通过举例来论述），当\Person{达味}说：\Quote{天主，公义的审判者，刚毅而含忍}时，如果短语中的每个形容语没有隐含\Quote{是}，那么这些称呼的列举似乎是无目的和不真实的，因为它们没有任何主体作为依托；但当\Quote{是}与每个名字一起被理解时，所说的话便显然有效，因为它相对于\Quote{是者}而被默观。因此，正如我们说到\Quote{祂是审判者}时，我们想到的是某种审判行动，并通过\Quote{是}将心智引向主体，从而清楚地被教导不要以为祂的\OriginalTerm{存在}{being}与行动相同；同样，由于说到\Quote{祂是受生的（或非受生的）}，我们将思想分为双重概念：通过\Quote{是}理解主体，通过\Quote{受生}或\Quote{非受生}领会属于主体的东西。因此，正如当\Person{达味}教导我们天主是\Quote{审判者}或\Quote{含忍}时，我们并非得知天主的本质，而是得知在本质中默观的某一属性；同样，在此情况下，当我们听说祂是非受生时，我们也并非藉此否定性谓述而理解\Emph{主体}，而是被引导至我们不该对主体想什么，而祂本质上是什么却始终未得到解释。同样，当圣经将其他神圣名称归于\Quote{是者}，并将无名的存在启示给\Person{梅瑟}时，那揭示该存在之本性的人，就不应列举存在的属性，而应通过言语向我们显明其实际的本性。因为你可能使用的每一个名称都是存在的一个属性，而非存在本身——\Quote{善}、\Quote{非受生}、\Quote{不朽}——但这些名称中的每一个都必然要加上\Quote{是}。因此，任何人若试图照其本来面目解释这个善的存在、这个非受生的存在，却只是列举在祂内默观的属性，而对他试图用言语解释的本质保持沉默，那便是徒劳无益。\Quote{\OriginalTerm{无生}{without generation}}是在存在内默观的属性之一，但\Quote{存在}的定义是一回事，\Quote{以某种特定方式存在}的定义是另一回事；这一点至今仍未被所引用的段落讲述或解释。那么，让他先向我们揭示本质的名称，然后再以称谓的差异来划分本性吧；——只要我们所需的东西仍未得到解释，他运用科学技巧于名称上便是徒劳的，因为这些名称没有独立的存在。\par

这就是\Person{欧诺米}对抗真理的一个更强有力的把柄，而我们却默然放过了他这部分的著作中出现的许多观点；因为我认为，那些与真理的敌人进行武装竞赛的人，应当武装自己，对抗那些被谎言的貌似可信严密防御的人，而不应让自己的论证被那些早已死亡且气味难闻的观念玷污。他假设，凡在\OriginalTerm{本质}{essence}概念上统一的东西，必然形体地存在并与\OriginalTerm{朽坏}{corruption}相连（他在其作品的这一部分中这样说道），我乐意像避开腐败气味一样绕过它，因为我认为每个有理性的人都将看出这种论证是如何腐朽与死亡。因为谁不知道，人类灵魂的数目虽不可胜数，但一个本质支撑着所有灵魂，而它们之中\OriginalTerm{同体}{consubstantial}的基质与身体的朽坏无关？因此，连孩童也能明白，身体朽坏消散，并非因为它们彼此有相同的本质，而是因为它们具有复合的\OriginalTerm{本性}{nature}。因此，复合本性的概念是一回事，它们本质共同性的概念是另一回事，因此说\Quote{可朽坏的身体是同一本质}是正确的，但相反的陈述则完全不正确，如果说\Quote{这个同体的本性也必然是朽坏的}，正如在灵魂的例子中所显示的，灵魂具有同一本质，而朽坏却并不因本质的共同性而附着于它们。而对灵魂给出的解释同样可以恰当地应用于我们在受造界中所思考的每一个理智性的存在。因为\Person{保禄}所汇集的名词，并不像\Person{欧诺米}所想要的那样，表示超世界力量彼此分歧的本性；相反，这些名称的含义清楚地表明，他在他的论述中提到的不是本性的差异，而是天上军旅\OriginalTerm{行动}{operation}的多样特性：因为他说，有\Quote{统权者}、\Quote{宝座者}、\Quote{掌权者}、\Quote{异能者}、\Quote{宰制者}。这些名称的性质使每个人立刻明白，它们的意义是按照某种\OriginalTerm{行动}{operation}安排的。因为统治、行使权力与主权、成为某人的宝座——所有这些都是任何精通论证的人都不会认为适用于\Emph{本性}的差异，因为每个名称清楚地表示的是\Emph{行动}。因此，任何人若说\Person{保禄}所列举的名字表示\Emph{本性}的差异，就是自欺，正如\Person{宗徒}所说，\BibleQuote{明白，}\BibleQuote{既不知道自己所说的，也不知道自己所主张的是甚么。}\BibleRef{弟前 1:7}因为名称的含义清楚地表明，\Person{宗徒}在可理解的权能中承认某些等级的差别，但他并不通过这些名称指示本质的各类别。\par

\SourceRecord{Translated by H.A. Wilson.；Nicene and Post-Nicene Fathers, Second Series；Vol. 5.；Edited by Philip Schaff and Henry Wace.；Buffalo, NY: Christian Literature Publishing Co.,；1893.；<http://www.newadvent.org/fathers/290107.htm>.}

% structure: p0001 source=h1 target=section reason=page-title

\section{驳斥欧诺弥（第八卷）}

% structure: p0002 source=h2 target=subsection reason=relative-heading-rank; base=h2

\subsection{第八卷有力地推翻了异端者的亵渎言论，即宣称独生子出自无有，且曾有一段时间祂不存在；并证明圣子并非新造之物，而是自永远就有，因为祂对梅瑟说\BibleQuote{我是自有者}，对玛诺亚说\BibleQuote{你为何问我的名字？它也是奇妙的}。此外，达味也对天主说\BibleQuote{你却是同一，你的年岁不会穷尽}；依撒意亚又说\BibleQuote{我是天主，首先的，以后也是我}；福音作者说\BibleQuote{他在起初，就与天主同在，就是天主}——祂既无开始也无终结。并且他证明那些说祂是新造、出自无有的人是崇拜偶像的。在此他极好地诠释了\BibleQuote{光荣的辉耀和本体的印像}。}

这些，就是欧诺弥论证的强项；我认为，当那些看似有力的论点被证明为如此脆弱和不实之后，我完全可以对其余部分保持沉默，因为其余部分实际上随着较强论点的被驳斥而同时被驳斥了；正如在军事行动中，当一支比其余更强大的部队被击败后，其余的军队在那些已经战胜了其强大部分的人眼中就不再重要了。但是，他的亵渎言论的主要部分位于他讲辞的后半段这一事实，不允许我保持沉默。因为独生子从无有进入存在这一转变，即欧诺弥那比一切不敬神更应避免的可怖而无法无天的教义，在他的论证顺序中接下来被坚持。既然每一个被这种欺骗迷惑的人，都有这样的句子挂在嘴边以维护这一教义——即\Quote{他若是存在的，就没有被生；若是被生的，就不存在}——即那从无有造出我们和万有的祂，自己也是从无有而来；并且这欺骗由此获得很大支持，因为较软弱的人被这种表面的似真性所迫，而同意这种亵渎，我们必须不可忽略这教义的\BibleQuote{苦根}，以免如宗徒所说，它\BibleQuote{长出并烦扰我们}。现在我说，我们必须首先考虑这论点本身，撇开与对手的争论，然后再着手检查和反驳他们所提出的论点。\par

真正天主性的一个标志由圣经的话指示出来，梅瑟从天上的声音学到的，当他听到那一位说：\Quote{我是自有者\BibleRef{出 3:4}。}我们认为，正确合理的是相信那在存有方面被描述为永恒和无限者，是唯一真正神圣的；在其中所观照的一切常是同一的，既不增长也不消耗；因此，若有人论及天主说，祂从前是，但现在不是，或者祂现在是，但从前不是，我们应该认为这些说法同样都是不敬的：因为两者都切断了永恒的观念，在一边或另一边被非存有所截断，无论我们观照\Quote{无}先于\Quote{存有}，还是宣称\Quote{存有}终结于\Quote{无}；而屡次重复\Quote{首先}或\Quote{最后}有关天主的非存有，也不能弥补涉及天主性的不敬概念。为此，我们宣称，坚持他们关于那真实存有者在某时非存有的教义，就是否认和拒绝祂的真正天主性；其理由是：一方面，那藉着光向梅瑟显示自己的那位，称自己为\Emph{存有者}，当祂说：\Quote{我是自有者\BibleRef{出 3:4}}；而另一方面，依撒意亚（可说成了那在他内发言者的工具）以那存有者的口吻说：\Quote{我是首先的，此后我仍然在}，如此，无论我们怎样思考，我们都在天主内领会永恒。同样，那对玛诺亚说的话也显示了天主性不能藉其名字的意义而被理解的事实，因为，当玛诺亚请求知道祂的名字，以便在应许实际成就时，他可以藉名字光荣他的恩主时，祂对他说：\Quote{你为什么问这事？它也是奇妙的}；因此我们学到，有一个名字是表明天主本性的——即那在我们心中不可言说地兴起关于祂的惊奇。同样，伟大的达味在他自省的言论中宣告同样的真理，即一切受造物都由天主造成，而祂独一永远以相同的方式存在并永久长存，他说：\Quote{但你却是同一的，你的岁月不会穷尽。}当我们从受天主感动的人那里听到这些言论及类似的话时，让我们将一切不是来自永恒的事物留给拜偶像者的敬拜，作为一件与真正天主性相异的新事物。因为那现在存在而从前不存在的，显然是新的而非永恒的，而注意任何新的敬拜对象，梅瑟称之为魔鬼的侍奉，他说：\Quote{他们祭献给魔鬼而非天主，献给他们祖先不认识的神祇；他们是新近兴起的新神。}如果敬拜中的一切都是新的就是魔鬼的侍奉，与真正天主性相异，并且如果那现在存在但不总是存在的就是新的而非永恒的，那么我们这些注意那\Emph{存有者}的人，必然将那些默想非存有附着于那存有者，并说\Quote{祂曾经不存在}的人，算在拜偶像者之列。因为我们也可以看到，伟大的若望在宣讲独生天主时，处处保护自己的陈述，使非存有的概念无法接近那存有者。因为他说祂\Quote{在起初就有}，\Quote{与天主同在}，\Quote{就是天主}，并且就是光、生命、真理，以及一切美善的事物，任何时候都从未缺少任何卓越之物，祂是一切美善的圆满，并居于圣父的怀中。如果梅瑟为我们制定了这样一条真正天主性的标志，即我们对天主一无所知，只知道祂\Emph{存有}（因为到达这一点的话是：\Quote{我是自有者\BibleRef{出 3:4}}）；而依撒意亚在他的宣讲中大声宣告那存有者的绝对无限，定义天主的存有无关于开端或终结（因为那说\Quote{我是首先的，此后我仍然在}的一位，没有在任何一个方向给祂的永恒设限，因此，如果我们看向开端，我们找不到任何标记\Emph{自}祂存在之后，并超越这个时刻祂不存在；同样，如果我们转向未来，我们不能用任何界限切断那存有者的永恒进程）——并且如果先知达味禁止我们敬拜任何新的、陌生的神（这两者都包含在异端教义中；\Quote{新}明显指向那不永恒的事物，而\Quote{陌生}是与真天主本性的疏离）——我说，如果事情是这样，我们宣称所有关于那真实存有者在某时非存有的诡辩虚构，无非是背离基督信仰，转向偶像崇拜。因为当福音作者在论及天主的本性的言论中，处处将非存有从那存有者中分离，并藉着反复使用\Quote{有}这个词，仔细清除非存有的嫌疑，并称祂为独生天主、天主的圣言、天主子、与天主相等，以及所有这样的名称时，我们心中有这个确定的判断：如果独生子是天主，我们必须相信那被认为天主的是永恒的。而且祂确实是天主，肯定永恒，绝无任何时刻被发现为非存有。因为天主，如我们常说的，如果祂现在存在，也肯定一直存在；如果祂曾经不存在，那么现在也不存在。但既然连真理的敌人都承认独生子是并继续存留为独生天主，我们说：祂在圣父内，并不是仅在某个方面，而是完全在祂内，在一切圣父被认为所是的一切方面。所以，正如祂在圣父的不朽坏性内，祂是不朽坏的；在祂的良善内是良善的；在祂的大能内是有能的；并且，在圣父被认为有的每一特殊卓越的属性内，祂就是那事物；同样，在祂的永恒性内，祂肯定是永恒的。圣父的永恒性由祂从未从非存有取得存有、也从未使存有终结于非存有所标志。因此，那拥有圣父一切所有、并在圣父的全部光荣中被观照的那一位，正如祂在圣父的无尽内没有终结，同样，祂在圣父的无始内，如宗徒所说，\Quote{没有日子的起点\BibleRef{希 7:3}，}而同时是\Quote{出自圣父}，并在圣父的永恒内被看待；在这方面，更特别地看到圣子与祂所似者相比毫无差异。在此，祂的话被证明为真，即告诉我们：\Quote{看见了我就看见了父。}此外，正是这样，宗徒的话——圣子是\Quote{祂光荣的辉映，祂本体的真像\BibleRef{希 1:3}，}——被最好地理解为具有卓越而密切的适用。因为宗徒通过感官可见的事物，向那些不能藉纯理智对象的默观将自己的思想提升到认识天主之高度的听众，传达了一种关于真理的概念。因为正如太阳的实体由它周围的整个光盘所精确映照，看太阳的人藉他所见的推断整个固体基体的存在；同样，他说，圣父的威严在圣子的大能中被精确映照，使人相信一位如同另一位的伟大；又如，光的光辉从太阳的整个光盘射出它的灿烂（因为在光盘中一部分不发光而其余部分暗），同样，圣父所是的一切光荣，从祂的全部幅度藉由从祂而来的光辉——即真光——射出灿烂；又如，光线属于太阳（因为没有太阳就没有光线），但太阳从未被想象为没有从它射出的光辉光线而单独存在；同样，宗徒将独生子从圣父所获的存有的连续性和永恒性传达给我们，称圣子为\Quote{祂光荣的辉映。}\par

% structure: p0005 source=h2 target=subsection reason=relative-heading-rank; base=h2

\subsection{二、然后他讨论圣父关于圣子诞生的\Quote{意愿}，并表明那善意的对象是永恒的，即圣子，祂存在于圣父内，且与意愿的过程紧密相连，如同光线与火焰，或视觉与眼睛。}

在我们做出这些区分之后，无人能再怀疑独生子如何同时被视为\Quote{出自圣父}，且\Emph{是}永恒地存在，即使这一短语乍看之下与另一短语似乎不一致——即宣称祂\Quote{出自圣父}的那一短语与断言祂永恒的那一短语。但若我们要以进一步的论证来证实我们的陈述，则或许可能借助我们感官能认知的事物来领悟关于此点的教义。并且，若它在现有事物中找不出一个完全足以类比和相似的方式呈现当前问题的探求对象的形象，愿无人嘲笑我们的陈述。因为我们愿意藉着一切可能的例证，说服那些声称圣父先意愿而后成为圣父，并以此为据断言圣言在存在上具有后时性的人，使他们转向正统观点。这直接的结合既不排除圣父的\Quote{意愿}——即祂并非出于本性某种必然性而毫无选择地有了一个儿子，这\Quote{意愿}也不将圣子与圣父分开，仿佛作为某种间隔介于二者之间。因此，我们既不从教义中排除生养者对圣子所怀的\Quote{意愿}——仿佛被圣子与圣父合一之结合所排斥；也绝不打破那不可分离的联系，当\Quote{意愿}被视为参与于生育之中时。因为我们的沉重迟滞之本性，其特点是愿望与拥有某物往往不同时并存；时而我们愿望没有得到的东西，时而又得到我们不愿得到的东西。然而，对于那单纯且全能的性体，一切事物都被同时且立即地领会，包括对善的意愿以及祂所意愿之物的拥有。因为那善而永恒的意志被默观为在永恒性体中运作、内住并共存，并非源于任何外在原则，也不能脱离意志的对象而被领会：因为在天主那里，既不能没有善的意志，也不能意志的对象不伴随意志的行动，因为没有原因能使相宜于圣父的事不常存，或成为拥有意志对象的任何阻碍。既然独生子天主按其本性就是善（或者说超越一切善），并且善不会不成为圣父意志的对象，由此清楚表明：圣子与圣父的结合没有任何中介，并且那常存于善性体中的意志也不因此不可分离的结合而被排斥或排除。若有人不以嘲弄的精神来听我的论证，我愿意在此基础上添加如下内容。\par

例如，若有人承认火焰具有自由选择的能力（当然，这只是假设），那么显然，火焰在其存在之时即刻就会\Emph{意愿}其光芒从自身发出，而当它意愿时，它不会无能为力（因为火焰一出现，其自然能力就立即在光芒方面实现其意志），以致毫无疑问，若承认火焰是由自由选择所驱动，我们就同时设想这一切的汇合——火焰的点燃、它对光芒的意志行动以及光芒本身；以致选择性的运动并非光芒存在之尊严的阻碍——同样，按照我们所说过的例证，你因宣认圣父内存在善的意志行动，并不会通过那意志行动将圣子与圣父分开。因为假设意愿祂存在之行动会成为祂立即存在的阻碍，是不合理的。但正如在眼睛中，看见与看见的意愿，一个是本性的运作，另一个是选择的冲动，但选择向那个特定方向的运动并未给视觉行为带来延迟——（因为这两者各自被单独看待，并非彼此存在的阻碍，而是某种程度相互并存，本性的运作与选择同时发生，而选择反过来也不缺少本性的运动）——我说，知觉自然地属于眼睛，而看见的意愿并不对实际视觉造成延迟，而是人意愿它有视觉，立即他所意愿的\Emph{是}。同样，对于那不可言说且超乎一切思想的性体，我们对一切的领悟同时汇聚——圣父的永恒存在、对圣子的意志行动以及圣子本身，祂正如若望所说，\BibleQuote{在起初}，且不被认为是在起始之后来的。如今万有的起始是圣父；但在此起始中，圣子也被宣告是存在的，祂按其本性就是起始所是的那个事物本身。因为起始是天主，而那\BibleQuote{在起初}的圣言是天主。因此，既然\BibleQuote{在起初}这一短语指向永恒，若望巧妙地将\BibleQuote{在起初的圣言}连在一起，说圣言在起初内；我相信，他断言这一事实，目的是为了使听众心中出现的首要概念不单是\BibleQuote{在起初}本身，而是在此烙印于他之前，也使他心中与起始一同呈现那\BibleQuote{在起初}内的圣言，与起始一同进入听众的理解，并与起始同时出现在他的听觉中。\par

% structure: p0008 source=h2 target=subsection reason=relative-heading-rank; base=h2

\subsection{三、随后，他略过了先前已经讨论过的有关圣子本质的事，转而处理\Quote{受生}一词所包含的意义，说受生有各种不同的方式：有由质料和技艺造成的受生，如建筑物；有由动物延续的受生；有由流出而成的受生，如太阳及其光束、灯及其光辉、香料和油膏及其散发出的气味；以及由理智所生的言语；并巧妙地讨论了从朽木的受生、从火的凝结以及无数其他原因所产生的受生。}

既然我们已经彻底审视了我们的教义，也许现在是时候陈述并考察反对的论述了，将其与我们的观点并列比较。他是这样说的：\Quote{因为虽然我们提出了两个说法，}他说，\Quote{——一个是独生子的本质在它自己的受生之前并不存在，另一个是受生之后，它在万有之前，——但他并没有证明这两个说法不真实；因为他不敢说他是在那至高的受生和形成之前存在的，因为他同时遭到圣父的本性和清醒之人的判断的反对。因为哪个清醒的人能承认圣子在至高的受生之前就存在并受生呢？那无受生者不需要受生来成为他所是的。}好吧，无论他是否说得对，当他说我们的导师反对他的对立命题是徒劳无功时，所有熟悉那位作者著作的人一定都清楚。但就我自己而言（因为我认为驳斥他在这件事上的诽谤是揭露他恶意的一小步），我将把证明我们的导师并未不加讨论地忽略这一点的工作放在一边，转而尽我所能讨论他现在提出的观点。他说他在自己的论述中谈到了两件事：一是独生子的本质在它自己的受生之前并不存在，二是它在受生之后在万有之前存在。现在我认为，通过我们已经说过的话，已经充分表明：圣父所生的没有与圣父自身内所见的本质不同的新本质，我们没有必要卷入与这种亵渎的争论中，好像这论点是首次向我们提出似的；此外，我们的论证的真正力量必须集中到一点，即他那可怕而亵渎的言论，它明确地说到天主圣言\Quote{祂不存在}。再者，因为我们前面的论述已经在一定程度上处理了他亵渎的问题，也许再用类似的考虑来证明我们已经证明的事是多余的。因为我们之所以作前面那些陈述，正是为了通过先前在听众心中留下的正统思想方式的印象，使那些断言在独生天主身上非存在先于存在的对手们的亵渎更加明显。\par

看来在此时，通过更仔细的研究，在我们的论述中探讨\Quote{受生}一词的实际意义是恰当的。这个名称向我们呈现了作为某种原因结果的存在这一事实，这是每个人都清楚的，关于这一点，我想无需争论。但由于对作为原因结果而存在的事物所做的说明是多样的，我认为通过某种科学的划分在我们的论述中澄清此事是合适的。因此，对于作为某物结果的事物，我们理解为有以下种类：有些是质料和技艺的结果，如建筑物的结构和其他作品，通过各自的质料而形成，并由某种完成所提之事的技艺指导，以达到所产生结果的适当目的；另一些是质料和自然的结果，动物的受生就是自然的建造，自然通过它们的物质身体实体进行自己的运作；另一些是物质流出的结果，在这种情况下，在先者保持其自然状态，而从它流出的东西被分别考虑，如太阳及其光束、灯及其光辉、香料和油膏及其散发的气味；因为这些东西，虽然自身保持不变不受减损，但同时各自具有与自身共存的、它们所散发的自然属性：如太阳的光束、灯的亮度、香料在空气中产生的香气。除此之外，还有另一种\Quote{受生}，其中原因是非物质和非形体的，但受生是感官的对象，通过形体方式发生——我指的是由理智所生的言语：因为理智本身是非形体的，通过感官器官产生言语。所有这些受生的种类，我们在思想上仿佛都包含在一个总的看法中。因为自然所行的一切奇事，无论是将某些动物的身体改变成不同的种类，或是从液体的变化、种子的腐败、木头的腐烂中产生某些动物，或是从火凝结的团块中，将关闭在火心中的火把发出的冷气转化产生一种称为\OriginalTerm{火蜥蜴}{salamander}的动物——这些，即使看起来超出了我们所定的界限，仍然包括在我们所提到的情况中。因为自然是借着身体塑造这些不同形态的动物的；正是身体这样那样的变化，由自然以这种或那种特殊方式安排，产生了这样那样的特定动物；这并不是一种不同于由自然和质料所完成的受生的独特种类。\par

% structure: p0011 source=h2 target=subsection reason=relative-heading-rank; base=h2

\subsection{四、祂更进一步以人的比喻表明天主的行动；因为人手、脚以及身体的其他部分如何工作，在天主身上，惟独意志代替了这些。由此也产生了受生的差异；故此祂被称为\OriginalTerm{独生子}{Only-begotten}，因为祂与受造界所见到的其他受生毫无共通之处，但祂被称为\Quote{光荣的辉映}和\Quote{香膏的馨香}时，则表明祂的本性与圣父的密切结合与同永恒。}

现在，这些\OriginalTerm{生成}{generation}的方式为人所熟知，圣神的慈爱\OriginalTerm{安排}{dispensation}在向我们传授神圣奥秘时，借着我们所能领悟的事物，传递那些超越语言的教导，如同它经常在其他地方所做的那样——当它以身体性的词语描绘神性时，在论及天主时提及祂的眼、祂的眼皮、祂的耳、祂的手指、祂的手、祂的右手、祂的臂、祂的脚、祂的鞋等等——这些事物没有一样按其本义被认为属于天主性；而是将教导转向我们容易理解的事，用人类惯用的术语描述超出一切名称的事实，同时借着论及天主所用的每一个术语，我们被类比地引向某种更高超的概念。因此，它使用多种生成形态，从受默感的教导向我们呈现\OriginalTerm{独生子}{Only-begotten}的不可言喻的存在，从每一种中只取那些可以虔敬地纳入我们对天主的观念的部分。因为当它在论及天主时提及\Quote{手指}、\Quote{手}和\Quote{臂}时，并不用这些词语描绘由骨头、筋腱、肉和韧带构成的肢体结构，而是借着这样的表达表示祂有效且运作的能力；正如它借着这类其他词语指出与它们相应的那种关于天主的观念，而不接纳这些词语的肉体意义。同样地，它确实谈到这些生成形态的方式被应用于天主性，却并不以我们惯常知识所能理解的意义来说话。因为当它谈到形成的能力时，它称那特定的动力为\Quote{生成}，因为表达神圣能力的词语必须降卑到我们的卑微；然而它并不指出与我们当中形成性生成相关的一切——既无地点、时间、材料的准备，也没有工具的协作，亦无产生事物的目的，而是将这些置之不顾，并且伟大而崇高地将存在万物的生成归于天主，如它所说：\Quote{他一发言，它们便形成；他一命令，它们便被创造。}再者，当它阐明独生子从父所具有的那不可言喻、超越的存在时，由于人的贫乏无法企及那些过于言语或思想所能及的真理，它也在此使用我们的语言，称祂为\Quote{\OriginalTerm{子}{Son}}——这个名字我们日常用法适用于由物质和自然产生的事物。但是正如那告诉我们天主创造\Quote{生成}的词语，并没有附加说它是借任何材料的帮助而生成的，而是宣告它的物质实体、它的地点、时间以及一切类似者，都存在于祂旨意的能力中；同样在这里，论及\Quote{子}时，它也排除人性在地上生成中所见的一切其他事物（我指情感、倾向、时间的协作、地点的需要，尤其是物质），没有这一切，作为自然结果的地上生成就不会发生。既然一切关于物质和间隔的观念都被排除在\Quote{子}一词的意义之外，就只剩下本性，而借此在\Quote{子}一词中，就独生子而言，宣告了祂由父而来的密切而真实的彰显特性。既然这种特殊的生成形态不足以在我们心中产生对独生子不可言喻存在的充分观念，它也采用另一种生成形态，即由\OriginalTerm{流出}{efflux}而来的形态，以表达儿子的天主性，并称祂为\BibleQuote{光荣的\OriginalTerm{光辉}{brightness}} \BibleRef{希 1:3}，\Quote{香膏的\OriginalTerm{香气}{savour}}，\BibleQuote{天主的\OriginalTerm{嘘气}{breath}} \BibleRef{智 7:25}，这些在我们习常的用法中，在已进行的科学讨论里，称为物质流出。但正如在先前的情况中，无论是受造物的制造还是\Quote{子}一词的意义，都不容许时间、物质、地点或情感；同样在这里，这句话也净化了\Quote{光辉}及其他术语的一切物质观念，只采用这种特定生成形态中适合于天主性的要素，借着这种表达方式的力量，指向祂被构想为既出于祂又同祂在一起的真理。因为\Quote{嘘气}一词并不向我们呈现从底层物质散发到空气的扩散，\Quote{芬芳}也不呈现从香膏的品质到空气的转移，\Quote{光辉}也不呈现从太阳体通过光线的流出；而唯独这样，正如我们所说，借着这个特定的生成形态被彰显出来：即祂被构想为出自祂且与祂同在，在父与出自祂的那位子之间没有中间间隔。既然圣神的恩宠在其丰厚的慈爱中，安排我们关于独生子的观念应从许多源头兴起，它也补充了生成中所思虑之事的剩余种类——我指由心神和言语而来的那一类。但崇高的若望特别谨慎，以免听者因不注意或思想软弱而陷入对\Quote{\OriginalTerm{圣言}{Word}}的普通理解，以为子就是父的声音。因此，他在最初的宣报中预备我们，视圣言为在本质内，且不在任何与那本质（即祂的存有之源）相异或分离的本质内，而是在那最初而有福的本性内。因为当他论及圣言\Quote{起初就存在}且\Quote{与天主同在}时，这正是他所教导我们的，圣言本身也是天主以及\Quote{起初}所是的一切。因为这样，他关于神性的论述才触及独生子的永恒性。既然这些生成形态（我指由原因而来的那些）在我们中间是通常已知的，并且被圣经用来教导我们关于当前的主题，正如可以预期每一样都应用于神圣观念的呈现，那么，就让我们论证的读者\Quote{按公义判断}，异端所作的任何断言是否对真理有任何效力。\par

% structure: p0013 source=h2 target=subsection reason=relative-heading-rank; base=h2

\subsection{五、然后，在指明独生者及万物的创造者的位格没有起源（如祂所造的万物那样，正如欧诺米所说），而是独生者是无始的、永恒的，与受造物既没有本质上的共性也没有名称上的共性，而是从永远就与圣父共存，正如至高的智慧所说，\Quote{时间的起始、终结与中间}，并且在就独生者的神性和永恒性、以及关于灵魂、天使、生命和死亡作了许多论述之后，他结束了本书。}

我现在再次逐字附上对手的原话。如下：— \Quote{固然有}他说，\Quote{我们提出了两个陈述：其一，独生者的本质并非在它自己的生成之前；其二，它虽是被生成的，却先于万物——}我们的教义学家向我们提出了哪一种生成？是我们能恰当地思考和论及天主的那一种？谁如此不敬神，竟在天主内预设不存在？但显然他所考虑的是我们这种物质性的生成，并且正以低下的本性作为他对独生者天主之概念的导师；既然牛、驴或骆驼都不在它自己的生成之前，他认为就连对独生者天主也该说那低下的本性在动物身上呈现给我们看的情况，而没有考虑到他这个肉体的神学家这个事实：应用于天主的\Quote{\Emph{独生}者}这个谓词，由其自身词汇指的就是不与一切生殖共通的、限于祂自己的东西。如果\Quote{独生者}这个术语与其他的生殖有意义上的共通性和同一性，它怎么可能用于这个\Quote{生成}？有一些独特而例外的关于祂的事要理解，这不存在于其他生殖中，这由\Quote{独生者}这个称谓清楚而恰当地表达出来；正如，如果低下的生殖的任何要素被设想在其中，祂在祂生成的任何属性方面被置于与受生的其他事物相等的地位，就不再是\Quote{\Emph{独生}者}了。因此，如果\Quote{独生者}的意义指明与其余受生之物没有混合和共通，我们将不承认我们在低下生殖中所看到的任何东西也被设想在圣子从圣父所拥有的存在中。但是，在生成之前不存在，是所有因生成而存在之物的属性；因此这有别于独生者的特殊性质，而\Quote{独生者}这个名称证明，没有任何属于欧诺米误解的那种常见生成方式的东西附着于祂。因此，让这位唯物主义者、感官的朋友，通过其他形式的生成来纠正他概念的错误吧。当你听到\Quote{光荣的辉耀}或\Quote{香膏的芬芳}时，你会说什么？\Quote{辉耀}不在它自己的生成之前？但如果你这样回答，你肯定也会承认\Quote{光荣}也不存在，或\Quote{香膏}也不存在；因为不能设想\Quote{光荣}本身曾存在过，暗淡无光，或\Quote{香膏}不散发出它的甜味；所以如果\Quote{辉耀}\Quote{不存在}，那么\Quote{光荣}也必然\Quote{不存在}，而\Quote{芬芳}不存在，也证明了\Quote{香膏}的不存在。但是，如果这些取自圣经的例子引起任何人的害怕，认为它们没有准确地给我们呈现独生者的威严，因为两者在本质上并不与它们的基体相同——既不是散发与香膏相同，也不是光线与太阳相同——就让真实的圣言纠正他的恐惧吧，祂在起初就已有，且是起初所是的一切，并在万有之前存在；因为若望在他的宣讲中这样宣告：\BibleQuote{圣言与天主同在，圣言就是天主。}如果圣父是天主，圣子是天主，那么关于独生者的圆满神性还有什么怀疑呢？因为通过\Quote{圣子}一词的含义承认了本性的密切关系，通过\Quote{辉耀}承认了结合与不可分离性，通过同样应用于圣父和圣子的\Quote{天主}的称谓承认了它们的绝对平等，而\Quote{本体的真像}在论及圣父的整个位格时被默观，标志着圣子本有伟大中没有任何缺陷，\Quote{天主的形象}则通过显示出所有标志天主性的特征，表明祂的完全同一。\par

让我们再次提出欧诺米的陈述。他说：\Quote{祂未曾存在，在祂自己的受生之前。} \Quote{祂未曾存在}这句话是指谁说的？让他宣示那些神性的名号，按欧诺米的说法，那位\Quote{一度未曾存在}者正是由这些名号来称呼的。我想，他会说：\Quote{光}、\Quote{福乐}、\Quote{生命}、\Quote{不朽}、\Quote{正义}、\Quote{圣化}、\Quote{德能}、\Quote{真理}，以及诸如此类。因此，说\Quote{祂在受生之前未曾存在}的人，绝对是在宣告：当祂\Quote{未曾存在}时，就没有真理、没有生命、没有光、没有德能、没有不朽，也没有其他任何在我们对祂的概念中属于卓越本质的属性；而且，更为奇妙、对不敬神者更难面对的是，那时没有\Quote{光辉}，没有\Quote{肖像}。因为说没有光辉，就必然意味着光辉所从出的发光体也不存在，正如我们从灯的比喻中所见的。因为谈及灯的光线的人，也表明灯在发光；而说光线\Quote{不存在}的人，也意味着发光之物的熄灭。所以，当圣子被说成不存在时，作为必然结果，圣父的不存在也就被主张了。因为如果两者通过连接彼此相关，按照宗徒的见证——\Quote{光辉}与\Quote{光荣}，\Quote{肖像}与\Quote{本体}，\Quote{智慧}与天主——那么，说这些连接物中的某一个\Quote{不存在}的人，必然通过废除其中一个也废除了另一个；因此，如果\Quote{光辉}\Quote{未曾存在}，那就等于承认发光之性也不存在；如果\Quote{肖像}不存在，那么被肖象的本体也不存在；如果天主的智慧和德能\Quote{未曾存在}，那么那位离开智慧和德能便无法被思想的天主也必然被承认不存在。所以，如果独生子天主，如欧诺米所说，\Quote{在祂的受生之前未曾存在}，而基督是\Quote{天主的德能和天主的智慧}\BibleRef{格前 1:24}，又是\Quote{肖像}\BibleRef{希 1:3}和\Quote{光辉}\BibleRef{希 1:3}，那么圣父也必然不存在，因为圣子正是祂的德能、智慧、肖像和光辉。因为理性无法设想一个没有肖像的本体，或没有光辉的光荣，或没有智慧的天主，或没有手的造物主，或没有圣言的元始，或没有圣子的圣父；一切这样的事物，无论为承认者还是否认者，都显然被宣示为相互结合，并且由于一个的废除，另一个也随之消失。既然他们主张圣子（即\Quote{光荣的光辉}）在被生之前\Quote{未曾存在}，而逻辑后果又意味着，随着光辉的不存在，光荣也被废除，而圣父就是那从独生子之光发出光辉的光荣，让这些过于聪明的人想想，他们显然是伊壁鸠鲁学说的支持者，假借基督宗教之名宣扬无神论。既然逻辑后果表明必然面临两种荒谬之一：要么我们必须说天主根本不存在，要么我们必须说祂的存在不是无始的，那就让他们从面前两条路中任选一条——要么被称为无神论者，要么不再说圣父的本质是无始的。我想，他们会避免被视为无神论者。因此，他们只能主张天主不是永恒的。如果已被证明的推理迫使他们走到这一步，那么他们那些花样翻新、不可逆转的名号转换又成了什么？他们那些在老妇人听来如此动听的三段论，以\Quote{\OriginalTerm{受生的}{Generated}}和\Quote{\OriginalTerm{无始的}{Ungenerate}}的对立为基础的不可抗拒的强迫力，又成了什么？\par

这些事就到此为止。但不回答他下一点，或许也不妥当；不过，让我们沉默地跳过那喜剧式的插曲，在那插曲里，我们那位聪明的演说家无论出于玩笑还是认真，都暴露了他的少年自负，以为这样就能在论证中占得上风。因为确实没有人会强迫我们加入那些目光斜视、扭曲我们视线的人，或那些患了奇症而扭曲身体的人，或加入他们身体的跳跃和扑跌。我们怜悯他们，但不会偏离我们既定的心态。然后，他转而对我们老师讲论这个题目，仿佛真的与他面对面交锋，说：\Quote{你将自投罗网。} 因为\Person{巴西略}曾说过，善常与统管万有的天主同在，而成为这样一位圣子之父是善的——因此善从未脱离祂，圣父也不愿没有圣子，祂愿意时不缺能力，既有能力又愿意以祂认为善的方式存在，祂也因常愿意那善的事而常拥有圣子（因为我们教父言论的意图正是朝此方向），\Person{欧诺米}预先将这番话拆解，并从他外来的哲学中提出这样的论证来推翻已说的话：—— \Quote{你怎么办呢？}他说道，\Quote{如果那些曾经历这类论证的人中有人这样说：\InnerQuote{如果创造是善的，且符合天主的本性，那么为何那善的且符合祂本性的事，没有以\OriginalTerm{无始的}{unoriginate}方式与祂同在，既然天主是\OriginalTerm{无始的}{unoriginate}？并且在创造之事上没有无知、软弱或年龄的阻碍，}}以及他收集并倾泻在自己身上的其余一切——因为我不能说，倾泻在天主身上。好啦，如果我们的老师能亲自回答问题，他本会以那受天主教导的口舌阐明天主的奥迹，并用他的驳斥鞭打那欺骗的冠军，从而向所有人显明，基督奥迹的仆人与一个可笑的丑角或宣扬新奇荒谬教义的人之间有多大区别。但由于他，如宗徒所说，\Quote{虽然死了，仍旧说话}向天主，而对方却提出这样的挑战，好像无人能回答他，即使我们的回答与大巴西略的话语相比可能没有那么有力，我们仍要勇敢地向提问者说：—— 你自己为推翻我们的陈述而提出的论证，正是我们在指控你们不虔敬教义时所说真话的明证。因为我们所责怪的没有比这一点更甚：你们认为创造主与受造物整体之间没有区别，而你现在所引证的正是我们所指摘的。因为如果你必须将你在受造物中所见的恰好也归于\OriginalTerm{独生子天主}{Only-begotten God}，那么我们的争论就达到了目的：你自己的话宣告了这教义的荒谬，且对所有人显而易见：我们使我们的论证保持在真理的正道上，而你对\OriginalTerm{独生子天主}{Only-begotten God}的观念如同你对其余受造物的观念一样。\par

争论是关于谁？岂不是关于\OriginalTerm{独生子天主}{Only-begotten God}，一切受造物的创造者，祂是否永远存在，还是后来才作为对圣父的附加而存在？那么我们的老师关于此事说了什么呢？他说，相信那自然为善的事不在天主内，是不虔敬的：因为他看不出任何原因可能使善不常与那善者同在，无论是因为能力不足还是意志软弱。那反对这些陈述的人说了什么？\Quote{如果你允许相信\OriginalTerm{天主圣言}{God the Word}是永恒的，你也必须允许受造物也是如此}——（他多么善于在论证中区分受造物的本性与天主的尊威！他对每件事物，什么适合于它，他可能虔诚地想什么关于天主，什么关于受造物，知道得多么清楚！）—— \Quote{如果创造主，}他说道，\Quote{从祂创造的时间开始：因为没有别的东西可以标记受造物的开始，如果时间不以它自身的间隔来限定受造物的开始和结束的话。}\par

基于此，他说时间的创造者必然从类似的起点开始存在。然而，受造物以诸世代为开端，但你能为诸世代的创造者设想出什么样的开端呢？若有人说\Quote{福音所提及的\InnerQuote{开端}}——那所指的便是圣父，并且也指出了圣子与祂同在的宣认；既然主说祂在圣父内，那祂在圣父内的存在就不可能从某一点开始。若有人谈论除此之外的另一个开端，就请他说出这个开端的名称，因为在我们所能领悟的世代之前，没有任何东西可以辨识。因此，即使老妇人也喝彩说这命题正确，这样的陈述也丝毫不会动摇我们关于独生子的正统观念。因为我们坚守从起初所确定的，我们的教义牢固地建立在真理之上：即正统教义认为我们关于独生子天主所断言的一切，与受造物毫无共同之处；而是区分万有创造者及其作品的标志相隔甚远。的确，如果圣子在任何其他方面与受造物有相通，那么我们当然应该说他甚至在存在方式上也没有与受造物分开。但如果受造物不分享我们所学到的关于圣子的一切，那么我们也必须必然地说，在这方面祂与受造物也没有相通。因为受造物不是\InnerQuote{在起初}，也不是\InnerQuote{与天主同在}，也不是天主，也不是生命、光明、复活，也不是其他神圣的名号，如真理、正义、圣化、审判者、公义者、万有创造者，存在于万世之前，直到永远；受造物不是光荣的辉耀，不是本体的肖像，不是良善的相似，不是恩宠，不是德能，不是真理，不是救恩，不是救赎；我们也找不到圣经用于独生子光荣的那些名号中任何一个属于受造物，或用于受造物——更不用说那些更崇高的言语：\Quote{我在父内，父在我内}，以及\Quote{看见我的，就是看见了父}，以及\Quote{除了父，没有人见过圣子}。如果我们的教义允许我们为受造物主张如此众多且如此伟大的事物，那么他或许有理由认为我们也应将我们在其中观察到的特性附加于我们对独生子的概念，因为转移是从同类的主题到近乎亲属。但如果所有这些概念和名号都涉及与圣父的相通，而它们又超越我们对受造物的概念，那么我们这位聪明又敏锐的朋友难道不该羞惭地溜走吗？他竟试图通过他对受造物的观察来讨论创造主的本性，却没有意识到区分受造物的标志是另一种类？一切存在物的最终区分是由\Quote{受造}与\Quote{非受造}之间的界限所作出的，其一被视为生成之物的原因，另一则由此生成。既然受造的本性与神圣本质如此划分，且在其区别属性上不容混同，我们绝不可用相似的术语来设想两者，也不可在如此分离的事物中寻求其本性观念中相同的区分标志。因此，正如最卓越的智慧之言语在某处告诉我们，存在于受造物中的本性在自身内展示出\Quote{时间的开端、终结与中间}\BibleRef{智 7:18}，并与所有时间间隔同时延伸，我们把这一特性当作该主体的一种标志，即我们从中看到其形成的某个开端，注意其中间，并将其期望延伸到终结。因为我们得知天地不是从永远就有，也不会存留到永远，因此显然它们都从某个开端开始，也必在某个终结中停止。但神圣本性不受任何限制，反而以无限性超越所有界限，远远脱离我们在受造物中发现的那些标志。因为那没有间隔、没有数量、没有限制的权能，在自身内包含所有世代及其中发生的所有受造物，并凭借其本性的无限性，在一切方面超越世代的无尽范围，要么没有任何标志表明其本性，要么有一个完全不同的标志，而不是受造物所具有的。既然受造物有开端，那么属于受造物的就与未受造的本性相异。因为若有人胆敢假定独生子天主的存在像受造物一样，从某个我们可理解的开端开始，那么他必然也要为关于圣子的陈述附带其余的结果；因为不可能只承认开端而不承认随之而来的东西。正如若承认某人是具有全部本性属性的人，他就会观察到在承认中他宣告了一个有理性的动物，以及所有关于人所设想的其他东西；同理，若我们理解受造物的任何属性存在于神圣本质中，我们便不再能避免将那纯然本性中其余被默观的属性附加其上。因为\Quote{开端}会强制并强迫随之而来的事物；如此设想的\Quote{开端}是后续之物的开端，其意义是，若后续者存在，它也存在；若与之关联的事物被除去，前项也不会存留。既然智慧篇不仅提到\Quote{开端}，还提到\Quote{中间}和\Quote{终结}，若我们按照异端教义，在独生子的本性中假定某个由时间标志界定的存在开端，智慧篇绝不让我们免于在\Quote{开端}之后加上\Quote{中间}和\Quote{终结}。若如此行，我们将发现，根据我们的论证，神圣之言向我们展示天主是会死的。因为按照智慧篇，\Quote{终结}是\Quote{开端}的必然结果，而\Quote{中间}的概念包含在两端的概念中，允许其中一个的人也在潜能上持守其他，并为无限的本性设下度量和限制的界限。若这是不虔敬且荒谬的，那么为那导致不虔敬的论证赋予开端就应受到同等甚至更大的谴责；而这荒谬教义的开端被视为假定圣子的生命被某个开端所局限。因此，他们面前有两条路：要么在前述论证的迫使下回归正确教义，默观那出自圣父的、与圣父的永恒合一的那位；要么若他们不喜欢，则必须从两方面限制圣子的永恒，并通过开端和终结将祂生命的无限性化为乌有。还有，即使灵魂和天使的本性因被造且存在开端而从某个时间点开始，却没有终结，并毫不受阻地延续到永远，因此我们的对手可以利用这一事实为基督类比，说祂不是从永远存在，却仍永远持续——任何提出这一论证的人也请考虑如下一点，即天主性在特殊属性上与受造物何等不同。因为属于天主性的是不缺少任何被视为美善的可设想之物，而受造物则通过分享比自己更好的东西而获得卓越；此外，受造物不仅有其存在的开端，而且还被发现不断处于开始拥有卓越的状态中，因为它不断进步改善，从不停止于已达到者，而其所获得的一切通过参与成为其向更大者上升的开端，并且永不停止，按\Person{保禄}的说法\Quote{向着前面的目标伸展}，\Quote{忘记背后的}\BibleRef{斐 3:13}。既然天主性就是生命本身，独生子天主就是天主、生命、真理，以及一切可设想的崇高和神圣之事，而受造物从祂那里汲取美善的供给，由此可明显看出：若受造物因分享生命而存在，那么当它停止参与时，也必停止生命。因此，若他们胆敢也对独生子天主说那些适用于受造物的真理，就让他们连这个也一起说出来：祂的存在像受造物一样有开端，并按灵魂的样式在生命中存留。但如果祂就是生命本身，不需要在自身内拥有生命（\Emph{ab extra}），而所有其他事物不是生命，仅仅参与生命，那么什么迫使我们必须根据我们在受造物中看到的东西取消圣子的永恒呢？因为那在本性上永不改变者，不允许有相反者，且不能改变到任何其他状态；而本性处于边界之上的事物则有两种倾向，任意偏向它们觉得有吸引力的东西。因此，若真正的生命是在神圣和超越的本性中被默观的，那么它的衰败似乎必然以相反状态告终。\par

现在，\Quote{生命}和\Quote{死亡}的意义是多种多样的，并不总是以相同的方式理解。就肉体而言，感官的活力与运动被称为\Quote{生命，}其消逝与分解被称为\Quote{死亡。}但对于理智本性而言，接近神圣者乃是真正的生命，远离它则被称为\Quote{死亡}；因此，最初的恶者，魔鬼，既被称为\Quote{死亡，}又被称为死亡的发明者；宗徒也说它掌握着死亡的权势。那么，正如我们由圣经所获得的，如前所述，对死亡有双重的理解，那位真正不变而不可改变者\Quote{独享不死不灭，}并居于不可企及、不可接近的光明中，远离邪恶的黑暗\BibleRef{弟前 3:16}:但所有参与死亡的受造物，因与其相反倾向而远离不死不灭，若背弃善，就会因本性的可变性接纳与更糟状况的相通，这无非就是死亡，与身体的死亡有某种对应。因为如同在那种情况下，本性的活动终止被称为死亡，同样，在理智存在中，缺乏向善的运动就是死亡和离弃生命；因此我们在无形受造物中所感知的，并不与我们驳斥异端教义的论点相冲突。因为对应于理智本性的死亡形式（即与天主分离，我们称天主为生命）甚至在潜能上并未与其本性分离；因为从无中产生显明了本性的可变性；而自身易于变化之物，因那坚固它的恩宠而得以避免参与相反状态；它并非凭本性持守于善；这样的事物不是永恒的。那么，若有人真的说实话，说不应以同一标准衡量神圣本体与受造本性，也不应以任何开端来限定天主圣子的存有，以免若承认这一点，受造的其他属性就会偕同对这一点承认而一同进来，那么，那人的教导（他利用受造的属性将独生子天主与圣父的永恒分割开来）的荒谬性就清楚地显明了。因为正如其他任何标志受造的特征都不出现在受造的造物主身上，同样，受造物是从某个开端获得存在这一事实，并不能证明圣子并非始终在圣父内——圣子是智慧、德能、光明、生命，以及一切在圣父怀中可想到的。\par

\SourceRecord{Translated by H.A. Wilson.；Nicene and Post-Nicene Fathers, Second Series；Vol. 5.；Edited by Philip Schaff and Henry Wace.；Buffalo, NY: Christian Literature Publishing Co.,；1893.；<http://www.newadvent.org/fathers/290108.htm>.}

% structure: p0001 source=h1 target=section reason=page-title

\section{驳斥欧诺弥（第九卷）}

% structure: p0002 source=h2 target=subsection reason=relative-heading-rank; base=h2

\subsection{一、第九卷宣称，欧诺弥对天主本性的叙述，在某种程度上表达良好。接着，他因着相似性，将斐洛著作中的语句混入自己的论证：\Quote{天主在一切被生成之物之前}，并加上一句：\Quote{祂主宰自己的大能}。额我略憎恶这极度的荒谬，有力地驳斥了它。}

然而他现在转向更高深的言语，自高自大，以空虚的妄自尊大而膨胀，着手说出一些配得上天主威严的话。他说：\Quote{因为天主，是万善中最崇高者，是万有中最强大者，且免于一切必然性——}这位豪迈的人何等高贵地将他的言论，像一艘没有压舱物的船，被欺骗的波涛无舵地驱赶，带进真理的港湾！\Quote{天主是万善中最崇高者。}绝妙的承认！我想他不会指控伟大的若望行为不合宪法，若望在他崇高的宣告中，宣告独生子是天主，祂与天主同在，且就是天主。那么，如果这位宣传独生子天主性的人值得信赖，且如果\Quote{天主是万善中最崇高者}，那么结论便是：圣子被祂光荣的敌人声称是\Quote{万善中最崇高者}。既然这一措辞也应用于圣父，那么\Quote{最崇高}的最高级含义便不容许通过比较而有所减少或增加。但是，既然我们已从敌人的见证中获得这些陈述来证明独生子的光荣，我们也必须加上他接下来的陈述来支持纯正的教义。他说：\Quote{天主，万善中最崇高者，既无本性的阻碍，也无原因的强制，也无需要的冲动，便按照祂自己权威的最高主权，生养并创造，以祂的意志作为足够的能力来构成所造之物。那么，如果一切善都是按照祂的意志，祂不仅决定所造之物为善，也决定它成为善的时间，也就是说，如果我们可以假设，制造自己所不愿之物乃是软弱的标志。}我们将借用到此为止，为了确认正统教义，从我们敌人的陈述中，尽管那陈述被卑劣和虚假的条款所渗透。是的，那位凭着祂权威的最高主权，在祂的意志中拥有足够构成所造之物的能力者，那位毫无本性的阻碍或原因的强制而创造万有者，确实不仅决定所造之物为善，也决定它成为善的时间。但创造万有的那一位，正如福音所宣告的，是独生子天主。祂在祂愿意的时候创造了受造界；那时，借着\OriginalTerm{环绕本质}{circumambient essence}，祂用诸天的身体环绕那个被封闭在其范围以内的整个宇宙；那时，当祂认为此事应如此时，祂使旱地显现，将水封闭在它们的凹陷之处；植物、果实、动物的生育、人的形成，都在那时出现，当时每一件事都似乎适合造物主的智慧：——而造这一切的那一位（我将再次重复我的陈述）便是创造诸世代的独生子天主。因为如果诸世代的间隔先于存在之物，那么使用时间副词是合适的，可以说\Quote{祂\Emph{那时}愿意}和\Quote{祂\Emph{那时}造了}；但是，既然世代尚未存在，既然就那不以数量或间隔来度量的神性本体而言，我们心中没有间隔的概念，那么时间表达的力量必然无效。因此，说受造界按照创造万有者的智慧的美意，被赋予了一个时间上的开始，并未超出可能性；但是，将神性本体本身视为处于一种由间隔度量的延伸之中，则只属于那些受过新智慧训练的人。这是何等的一点，埋藏在他的言辞中，我因急于触及主题而有意略过！现在我将重拾它，并读出来，以显示我们作者的巧妙。\par

\Quote{那在一切受生之物以前在天主自身内被高举到最高者，}他说道，\Quote{对自己的能力拥有统治权。}这句措辞被我们的小册子作者从希伯来人斐洛逐字逐句搬到了他自己的论证中，而欧诺米剽窃的事实，将由斐洛自身的著作向任何关心此事者证明。然而，此刻我指出这一事实，与其说是为了指责我们这位言论贩子自身论据与思想的贫乏，不如说是旨在向我的读者表明欧诺米的教义与犹太人的推理之间的密切关系。因为斐洛的这句措辞若不是两者意图之间有某种亲缘关系，就不会逐字逐句吻合他的论证。在希伯来作者那里，你可以找到如下形式的措辞：\Quote{在一切受生之物以前的天主}；而随后的话，\Quote{对自己能力拥有统治权}，是新犹太主义的添加。但这一说法的检查将清楚表明其中包含何等荒谬。他说道：\Quote{天主对自己能力拥有统治权。}请告诉我，祂是什么？祂对什么拥有统治权？祂是某种不同于自己能力的东西，且是某种不同于祂自身的能力的主吗？如此，能力就被能力的缺失所胜。因为那不同于能力的当然不是能力，因此祂被发现拥有对能力的统治权，正是在于祂不是能力。或者，天主作为能力，在祂自身内有另一种能力，并以一种能力统治另一种能力。又有什么争斗或分裂，使得天主应分裂存在于祂自身内的能力，并以一种能力推翻祂能力的另一部分？我想，祂不可能在没有某种更大、更猛烈的能力协助之下，对自己的能力拥有统治权！这就是欧诺米的天主：一个具有双重本性或组合的存在，自相分裂，一种能力与另一种不和谐，以致一种能力催促祂陷入混乱，另一种能力则约束这不和谐的运动。再者，祂统治那催促受生的能力的意图是什么？恐怕若受生不受阻碍，就会产生某种邪恶吗？或者毋让他首先解释这一点——那在本性上被统治的是什么？他的话语指向某种冲动与选择的运动，被分别、独立地思考。因为那统治的必然是同一事物，那被统治的是另一事物。现在，天主\Quote{对自己的能力拥有统治权}——而这又是什么？一种自我决定的本性？还是别的什么，迫向不安，或停留在平静状态中？好吧，如果他假定它是平静的，那安宁者不需要任何人统治它；而如果他说\Quote{祂拥有统治权}，祂显然\Quote{拥有统治权}于某种推动、运动之物：而这一点，我猜想他会说，在本性上不同于那统治它的祂。那么，让他告诉我们，他在这观念中理解了什么？是某种独立存在的、在天主之外的事物吗？另一种存在如何能在天主内？抑或是神圣本性中的某种状态，被认为具有并非自身的存在？我几乎不认为他会如此说：因为那没有自身存在的，是不存在的；那不存在之物，既不处于统治之下，也不从中解脱。那么，那在基督受生适当时间尚未来临、并将释放这能力以进行其自然运作之前，处于统治之下、在其自身活动方面受到约束的能力是什么？那导致拖延的中介原因是什么，因它天主延迟了独生子的受生，认为尚不是成为圣父的好时机？又有什么插入在圣父的生命与圣子的生命之间，既非时间、亦非空间、亦非任何延伸概念或任何类似之物？这敏锐而目光清晰的眼是何用意，它标记并观望天主生命相对于圣子生命的分离？当他走投无路时，他自己被迫承认这间隔根本不存在。\par

% structure: p0005 source=h2 target=subsection reason=relative-heading-rank; base=h2

\subsection{二、接着，他巧妙指出，圣子的受生并非如欧诺米的话所说的\Quote{圣父在他选择的时候生了他，而非在此之前}；但圣子作为一切美善与卓越的圆满，常在圣父内被默观；为此论证使用了欧诺米自己论断的支持。}

然而，尽管二者之间没有间距，他却不承认它们的共融是直接而亲密的，反而屈就于我们知识的尺度，如同我们中的一员以人的言语与我们交谈，自己默默地承认推理的无能，并诉诸于一种从未被亚里士多德及其学派教导过的论证路线。他说：\Quote{祂在祂所愿意的时候生了祂的子，这是好而适当的；在明智之人的心中，不会因此产生任何疑问，为何祂以前不这样做。} 这是什么意思，\Person{欧诺弥}？你也要像我们这些没有学问的人一样徒步而行吗？你正在离开你那艺术性的文风，实实在在地诉诸于无理性地同意吗？你，那个如此责备那些着手写作却缺乏逻辑技巧之人的人？你，对\Person{巴西略}说：\Quote{当你断言，表达神性事物的术语的定义对人是不可能的，你便显出了自己的无知，}又另在一处提出同样的指控：\Quote{当你宣称，对你而言不可能的事对所有人都是不可能时，你便使自己的无能成为大家共有的。} 你，说这些事的人，就是这样接近那询问为何圣父迟迟未成为如此一位圣子之父的人之耳吗？你认为这样说是一个充分的解释吗：\Quote{祂在祂所选择的时候生了祂：关于这一点，不要再有问题了}？你忧虑的幻想在维护你的教义时变得如此软弱吗？导致两难的前提怎么了？你那强有力的证据怎么了？你技艺中那些可怕而必然的三段论结论怎么化为虚妄和虚无？\Quote{祂在祂所选择的时候生了圣子：关于这一点，不要再有问题了！} 这就是你众多劳作、你浩大事业的成品吗？所问的问题是什么？\Quote{如果天主拥有这样一位圣子是美好而适宜的，为什么我们不应当相信这美好常与祂同在呢？} 他从他哲学的圣殿中对我们作出什么回答，用必然的必然性紧紧捆绑他的论证？\Quote{祂在祂所选择的时候造了圣子：至于为何以前不这样做，不要再有问题了。} 为什么？如果我们面前的询问是关于某个无理性的存在，它凭自然冲动行动，问它为何不更早做某事——为什么蜘蛛不早织网，蜜蜂不早酿蜜，斑鸠不早筑巢——你还能说什么呢？难道不是同样的回答已经准备好了——\Quote{它在它选择的时候做了：关于此事不要再有问题了}？不仅如此，如果是关于某个雕刻家或画家，凭其模仿性艺术从事绘画或雕刻，无论是什么（假设他行使自己的艺术而不受任何权威管辖），我想这样的回答足以应付任何想知道他为何不更早行使艺术的人——即，他并非出于必要，而是将自己的选择当作他行动的机会。因为人并不总是愿望相同的事，而且常常没有与意愿合作的能力，他们在选择倾向于工作且没有外在阻碍的时候，做看起来对他们好的事。但是那始终同一的本性，没有附加的善，在其中由错误或无知以对立方式产生的一切计划纷争毫无容身之处，没有任何因变化而来的事物（它之前没有的），也没有任何它后来选择的东西是它起初未曾视为善的——说这一本性不会始终拥有善，而是后来才选择拥有它先前未选择的东西——这属于超越我们的智慧。因为我们曾受教导，\OriginalTerm{神性本体}{Divine Nature}时时刻刻充满一切真善，或者不如说，它本身就是一切善的圆满，因为它无需附加任何东西来成就完美，而是以其本性就是善的完美。现在，完美之物同样远离增加和减少；因此，我们说我们所见神性本体中善的圆满始终如一，正如无论我们向哪个方向扩展我们的思想，我们在那里都领悟它如其本来那样。因此，神性本体从不缺乏善：但圣子是一切善的圆满：因此，祂始终被默观于那本性为一切善之完美的圣父之内。但是他说：\Quote{关于这一点，为何祂以前不这样做，不要再有问题了。} 我们将回答他——最智慧的大人，将你碰巧赞成的某个命题立为法令是一回事，通过对争议之点进行推理来使人归信则是另一回事。因此，只要你不能给出任何理由，让我们可以虔诚地说圣子是\Quote{后来}由圣父所生，你的法令在明智之人中就不会有任何效力。\par

如此，欧诺米因着他那科学的攻击，将真理向我们彰显。至于我们，将按我们惯常的做法，运用他的论证来建立真道，好使甚至通过这一段话也能清楚看见：他们在每一处，被迫违背己意，支持我们的观点。因为，如果如我们的对手所说：\Quote{祂在祂所选择的时候生了圣子}，如果祂总是选择那善的，祂的能力与祂的选择一致，那么，圣子将被视为常与圣父同在，因为圣父总是选择那卓越的，并有能力获得祂所选择的。如果我们也将其后的话化为真理，我们也不难使它们适应我们所持的教理：\Quote{在明智的人中间，不要对此有疑问，为何祂以前不这样做}——因为\Quote{以前}一词有时间含义，与\Quote{以后}和\Quote{后来}相对；但假设时间不存在，那么表达时间间隔的语词也随着它被废除。然而，主是在时间之前、在诸世代之前；对于明智之人，关于诸世代的创造者问\Quote{之前}或\Quote{之后}是无用的：因为这类语词，若不涉及时间，便毫无意义。既然主是先于时间的，那么\Quote{之前}和\Quote{之后}这两个词用于祂身上便没有位置。这话或许足以驳斥那些不需要任何人推翻、而凭自身软弱就倒下的论点。因为谁有那么多闲暇，能如此投入去听对方的论点以及我们对这些蠢话的抗辩？然而，在受不敬神之偏见影响的人心中，欺骗如同某种根深蒂固的染料，难以洗去，深深灼烧在他们的心上，让我们再花一点时间在我们论点之上，或许我们能将他们的灵魂从这邪恶的污点中洗净。在我所引的言论之后，他又仿效他的老师普吕尼库斯，加上了若干不连贯、无条理的八重侮辱和谩骂，进而到达他论点的顶点，并且搁置他那愚昧的非逻辑的陈述，再次用辩证法的武器武装他的论述，并且按他所想象的三段论方式维持他对我们的谬论。\par

% structure: p0008 source=h2 target=subsection reason=relative-heading-rank; base=h2

\subsection{三、他进一步显示：圣子在时间之前的受生，并非从普通和肉性的受生中得出的推论所谈论的对象，而是无始无终，且不依从欧诺米因对柏拉图论灵魂和希伯来人安息日安息的论述无知而虚构的捏造。}

他所说的话如下：——\Quote{既然一切受生过程并非延至无穷，而是终止于某一点；那些承认圣子有起源的人，便绝对不得不说祂那时停止了受生，并且不可怀疑那些停止受生之事的开端，因而也必有一个\Emph{开始}：因为受生的终止确立了生育和被生的开始。这些事实是不能被怀疑的，其根据既是自然本身，也是神圣的法律。}既然他企图以推论的方式确立他的论点，按照那些精通此类事物者的科学方法提出其普遍命题，并将对特殊情况的证明包含在普遍前提中；那么，让我们首先考虑他的普遍命题，然后再考察他的推论的力量。从\Quote{一切受生过程}中引出关于圣子先于时间的受生的证据，这是虔敬的做法吗？我们应当以普通本性作为关于独生者存在的导师吗？就我而言，我不会期望任何人会疯狂到这种地步，以致任何关于神圣而纯洁的受生的想法会进入他的幻想。他说：\Quote{一切受生过程并非延至无穷。}他所谓的\Quote{受生}指的是什么？他是在谈论肉体的、身体的诞生，还是无生命物体的形成？身体受生所涉及的情感是众所周知的——没有人会想到将这些情感转移到天主性体上。因此，为了使我们的论述不致因详细提及自然作品而显得多余，我们将沉默地略过这些事，因为我认为每个有理智的人自己都知晓那些使受生过程延长（无论就其开始还是终止而言）的原因：详细表达这一切会既冗长又多余——生育者的结合、被生者在子宫内的形成、产痛、出生、地点、时间，没有这些，身体的受生便不能完成——这些事都与独生者的神圣受生同样遥远：因为如果这些事中有一件被承认，其余的一切必然随之而来。因此，为使神圣的受生不受任何与情感有关的概念影响，我们将避免想到与它有关的、以间隔来测量的延展。现在，有开始和结束的东西当然被认为处于某种延展之中，而一切延展都以时间为尺度；既然时间（我们藉以标示出生的结束和开始）被排除在外，那么，对于不间断的受生，考虑结束或开始的思想便是徒然的，因为无法形成任何概念来标示这种受生开始或停止的点。另一方面，如果他指的是无生命的受造物，即使在这种情况下，同样地，地点、时间、质料、预备、工匠的能力以及许多类似的东西，共同使产品臻于完美。既然时间确实与一切产生的事物同时存在，并且既然对于一切受造物（无论有生命还是无生命），也都能设想与产品相关的构造基础，我们便能在这些情况下找到形成过程明显的开始和结束。因为即使材料的获取实际上就是织造的开始，并且是地点的标志，而且在逻辑上与时间相关。所有这些都为产品确定了它们的开始和结束；没有人能说这些事与独生天主的先于时间的受生有任何关联，以致藉助于我们眼下所考虑的事物，我们能够推算有关那受生的任何开始或结束。\par

我们既已讨论了这些事至此，让我们重新审视对手的论点。它说：\Quote{一切生成并非无限延伸，而是到达某个终结时便停止。}既然\Quote{生成}的意义已按两种含义被考虑过了——不论他用这个词是指有形之物的诞生，还是受造之物的形成（两者都与那无玷的本性毫无共同之处），这前提便显明与主题无关。因为并非如他所主张的那样绝对必要，即因为一切制造和生成都在某个界限终止，所以那些接受圣子受生的人应当用一个双重界限来限定它，就它而言推定一个开端和一个终结。因为只有以某种量化的方式被限定，事物才可以说开始或到达界限时停止，而以时间来衡量的尺度（其广延与所产生之物的量相伴）通过其间间隔区分开端与终结。但谁能衡量或视之为有广延的东西是没有量和没有广延的？他能为没有量的东西找到什么尺度，或为没有广延的东西找到什么间隔？或者谁能用\Quote{终结}和\Quote{开端}来定义无限？因为\Quote{开端}和\Quote{终结}是广延界限的名称，而哪里没有广延，哪里也没有界限。然而，天主本性是没有广延的，既没有广延，它就没有界限；那没有界限的是无限的，因此被这样称呼。所以，试图用\Quote{开端}和\Quote{终结}来限定无限是徒劳的——因为被限定的就不可能是无限的。那么，这个柏拉图派的\WorkTitle{斐德若篇}如何不连贯地把他自己的学说与\Person{柏拉图}在那篇对话中关于灵魂的思辨拼凑在一起呢？因为正如\Person{柏拉图}在那里谈到\Quote{运动的停止}，这位作者也急于谈到\Quote{受生的停止}，以便用漂亮的柏拉图式措辞来欺骗那些不了解这些事的人。他告诉我们：\Quote{这些事实不能同时基于本性自身和神圣法律而被怀疑。}但本性，从我们先前的话来看，似乎不可信赖来教导关于天主的受生——即使有人以宇宙本身作为论证的例证也不行：因为通过宇宙的创造，正如我们在\Person{梅瑟}的宇宙起源论中学习的，有时光的衡量，藉着规定的昼夜，按照一定的顺序和安排分配给每个受造之物：而这一点甚至我们的对手的声明也不承认关于独生子的存在，因为它承认主在诸时代的时间之前。\par

现在仍需考虑他藉着\Quote{神圣法律}来支持其论点，他以此来证明圣子的受生既有终结也有开端。他说：\Quote{天主愿意将创造的法律铭刻在希伯来人心中，并没有指定受生的第一天作为创造的终结，或作为创造开始的证据；因为祂赐给他们作为创造的纪念，不是受生的第一天，而是第七天，祂在那一天从祂的工程中休息了。}谁会相信这是\Person{欧诺弥}所写，而引用的这些词语并非我们插入，以歪曲他的作品，从而使他在读者面前显得可笑，因为他拉进来与问题毫无关系的事来证明他的论点？因为当前要证明的是，如他所承担的，那先前不存在的圣子进入了存在；并且祂在受生时，有了受生的开端和停止——祂的受生如同产痛一般在时间中延续。而他用来确立这一点的资源是什么？希伯来人民依照法律在第七天守安息日！这证据与当前的事情多么吻合！因为犹太人藉着懒散来尊崇他的安息日，正如他所说，证明了主既有出生的开端也停止了出生！我们的作者还忽略了多少其他证据，其分量绝不亚于他所用来确立争点的那一个！——第八天的割礼，无酵饼周，月亮的第十四天的奥秘，洁净祭，痲疯病人的观察，公绵羊、牛犊、母牛、替罪羊、公山羊。如果这些东西与论点相去甚远，就让那些对犹太奥秘如此感兴趣的人告诉我们，那个特定的事如何属于问题的范围。我们认为践踏倒下的人是卑鄙且不男子气的，我们将从他作品接下来的部分探究，是否有任何能使对手为难的东西。那么，他在下一段中所主张的一切，关于认为在圣父与圣子之间有任何中间物的不当，我将略过，因为在某种意义上它与我们的教义一致。因为如果不区分他的话中哪些是无可指摘的、哪些是该受责备的，那既是不加辨别也是不公正的，因为他一方面与自己的陈述作战，另一方面又不遵循自己的承认，谈到联系的直接性却拒绝承认其连续性，并认为圣子之前没有任何东西，同时又有些怀疑圣子\Emph{是}，却坚称祂在不存在时进入存在。我们将在这些点上花一点时间（因为论点已预先确立），然后继续处理所提出的论证。\par

任何人不可既将独生子的存有置于一切之上，又说他在出生之前并不存在，而是当圣父愿意的\Emph{那时}才被生。因为\Quote{\Emph{那时}}和\Quote{\Emph{何时}}这两个词，根据说正确话之人的普遍用法，以及根据它们在圣经中的含义，具有特别且恰当地指称时间的意义。有人可以引用\Quote{\Emph{那时}他们必在外邦人中说}，以及\BibleQuote{\Emph{当我}派遣你们时} \BibleRef{路 22:35}，和\Quote{\Emph{那时}天国好比……}，以及整部圣经中无数类似的短语，来证明这一点：圣经的通常用法用这些词来指称时间。因此，如果正如我们的对手所承认的，时间并不存在，那么指称时间的词也必然消失；而如果它不存在，那么在我们的观念中，必然要被永恒所取代。因为在\Quote{不存在}这个短语中，确实隐含了\Quote{曾}；正如，如果他谈到\Quote{不存在}而没有加上\Quote{曾}的限定，他也会否认他\Emph{现在}存在；但如果他承认他现今存在，却又反对他的永恒，那么他心中所想的肯定不是\Quote{\Emph{绝对地}不存在}，而是\Quote{\Emph{曾}不存在}。而既然这个短语若不以时间意义为基础就完全不真实，那么说在圣子之前没有任何事物，却又主张圣子并非永远存在，就是愚蠢和徒劳的。因为如果没有地方，也没有时间，也没有任何其他受造物，那\Quote{在起初就已存在的}圣言不在此处，那么说主\Quote{曾不存在}的陈述就完全从正统教义的领域中被移除。所以他与其说是与我们矛盾，不如说是与他自己矛盾，他既宣称独生子存在，又宣称他不存在。因为在承认圣子与圣父的结合不受任何事物中断时，他清楚地见证了他的永恒。但如果他说圣子不在圣父内，我们不会对这样的陈述说什么，而是会用圣经来反对它，圣经宣告圣子在圣父内，圣父在圣子内，没有在短语中加入\Quote{曾}或\Quote{何时}或\Quote{那时}，而是通过这个肯定且无条件的宣告来见证他的永恒。\par

% structure: p0013 source=h2 target=subsection reason=relative-heading-rank; base=h2

\subsection{四、然后，在表明\Person{尤诺米乌斯}对伟大的\Person{巴西略}的诽谤（即他称独生子为\Quote{\OriginalTerm{非受生者}{Ungenerate}}）是虚假的，并再次以极大的巧智讨论了独生子的永恒性、存有和无终性，以及光明与黑暗的受造之后，他结束了本书。}

至於他試圖證明我們說獨生天主是非受生的，這就好像他說我們實際上將聖父定義為受生的一樣：因為兩種說法同樣荒謬，或者說同樣褻瀆。因此，如果他決意誹謗我們，就讓他再加上另一項指控，並且毫不保留任何可能更猛烈地激怒聽眾反對我們的手段。但如果其中一項指控因明顯的誹謗性質而被保留，為何又提出另一項呢？正如我們所說，就褻瀆而言，稱聖子為非受生的和稱聖父為受生的，完全是同一回事。現在，如果在我們的著作中能找到任何這樣的詞語，稱聖子為非受生的，我們就將最終判決自己敗訴；但如果他隨意捏造虛假的指控和誹謗，為誹謗我們的教義而編造任何虛構陳述，那麼對明智之人來說，這一事實可以證明我們的正統性：當真理本身為我們作戰時，他卻用謊言來指控我們的教義，並製造與我們陳述毫無關係的異端起訴書。然而，對於這些指控，我們可以給出簡潔的回答。正如我們判定那說獨生天主是非受生的人應受詛咒，願他反過來詛咒那斷定那在起初的祂\Quote{曾不存在}的人。因為通過這種方法，將顯明誰的指控是真實的，誰是誹謗的。但如果我們否認他的指控，當我們說到聖父時，我們理解為這個詞也暗含聖子，而當我們使用\Quote{聖子}這一名稱時，我們宣稱祂真實地就是祂被稱呼的，由非受生的光通過受生而發出，那麼那些堅持說我們稱獨生天主為非受生的人，其誹謗怎能不顯明呢？然而，我們不會因為祂藉受生而存在，就因此承認祂\Quote{曾不存在}。因為每個人都知道，\Quote{存在}與\Quote{不存在}之間的矛盾是直接的，因此肯定其中一個詞就是完全否定另一個，而且正如\Quote{存在}在任何時間對任何被假定存在的事物都是相同的（因為天空、星辰、太陽以及其他\Quote{存在}的事物，現在存在的狀態並不比昨天、前天或任何更早的時間更多），同樣，\Quote{不存在}的意義在任何時間都同樣表達不存在，無論是相對於更早還是更晚。因為任何不存在的事物，其\Quote{不存在}的狀態並不比它以前不存在時更多，而是\Quote{不存在}的概念應用於在任何時間距離上\Quote{不存在}的事物。因此，在談論生物時，我們用不同的詞來表示已經存在的事物分解為\Quote{不存在}的狀態，以及從未進入存在的事物不存在的狀態，我們說要麼某物從未存在過，要麼那受生的已經死去，但無論哪種說法，我們都用語言同樣表示\Quote{不存在}。正如白天被黑夜從兩側包圍，但包圍它的黑夜部分名稱不同，我們稱一邊為\Quote{入夜後}，另一邊為\Quote{黎明前}，而兩者所指的都是黑夜，同樣，如果有人將那\Emph{不存在的}與那\Emph{存在的}相對照，他會對形成之前的狀態和形成之後解體的狀態給出不同的名稱，但對兩者所表示的狀態——形成之前的狀態和形成之後解體的狀態——卻構想為同一種狀態。因為未受生之物和已死之物的\Quote{不存在}狀態，除了名稱的差異外，是相同的——除了我們對復活希望的理解。既然我們從聖經中得知獨生天主是生命之首，就是生命、光明、真理，以及一切言語或思想中可敬的事物，我們說，在與那真實存在的祂相關聯時，構想相反的概念——無論是趨向朽壞的解體，還是形成之前的不存在——都是荒謬和褻瀆的；但當我們將思想延伸到未來或萬古之前的事物時，我們在任何地方都不在\Quote{不存在}的狀態上停止思考，認為在任何時候以不存在來中斷天主的存在同樣趨向褻瀆。因為說不朽的生命是可朽的，真理是謊言，光明是黑暗，那存在的是不存在的，都是同一回事。因此，凡拒絕允許祂在將來某一時刻不再存在的人，也將拒絕允許祂\Quote{曾不存在}，按照我們的觀點，他在兩方面避開了同樣的褻瀆：因為正如沒有死亡能中斷獨生天主生命的無窮盡，同樣，當我們回顧時，也沒有任何不存在的時期會終止祂朝向永恆的生命進程，以便那真實\Emph{存在的}可能與那真實\Emph{不存在的}完全無關。為此緣故，主渴望祂的門徒遠離這種錯誤（免得他們自己搜尋獨生天主存在之前的某物，而被他們的推理引向不存在的觀念），祂說：\BibleQuote{我在父內，父在我內}，意思是：那\Quote{不存在的}不能被構想在\Quote{存在的}之內，那\Quote{存在的}也不能被構想在\Quote{不存在的}之內。此處，詞語的順序本身就解釋了正統教義；因為聖父並非來自聖子，而是聖子來自聖父，因此祂說：\Quote{我在父內}，表明祂並非來自另一位，而是來自祂；然後反過來說：\Quote{父在我內}，表明那在好奇的推測中越過聖子的人，也越過了對聖父的理解：因為那在某某之內的事物，不可能被發現在它所在之外；因此，那人雖然不否認聖父在聖子內，卻想像自己以某種方式把握了聖子之外的聖父，是在說空話。同樣空洞的是我們對手關於\Quote{非受生}問題的虛幻爭鬥，它們在無實體的基礎上前進。然而，如果我要更充分地揭示他們論證的全部荒謬，請允許我在這個推測上再多花一點時間。因為他們說獨生天主在聖父之後\Quote{更晚}才存在，那麼他們所想像的這個\Quote{非受生的}，無論是什麼，都必然顯示出與邪惡的觀念相連。誰不知道，正如不存在的與存在的相對立，同樣，每一個善的事物或名稱都與相反的觀念相對立，例如\Quote{惡}與\Quote{善}相對，\Quote{虛假}與\Quote{真理}相對，\Quote{黑暗}與\Quote{光明}相對，以及其他所有這樣對立的事物，其中對立不容許中間狀態，兩者不可能共存，一方的出現會摧毀其對立面，而當另一方撤離時，其對立面就會出現？\par

既然这些论点已被我们接纳，那么每个人也都清楚，正如梅瑟所说，在光被造之前有黑暗，同样，就圣子而言（按照异端者的说法，圣父\Quote{在他愿意的时候造了他}），在圣父造他之前，圣子所是的光并不存在；光尚未存在，其对立面也必然存在。因为我们也从其他事例得知，出自造物主的一切都不是随意的，而是通过创造将所缺之物添加到已有之物中。因此显然，如果天主确实造了圣子，他造他是由于事物本性的缺欠。那么，正如可感觉的光尚缺时便有黑暗，而若光不曾存在，黑暗必会占据上风，同样，当圣子\Quote{尚未存在}时，那真实的光，以及圣子所是的一切，都不存在。因为即使按照异端者的明证，存在之物无需生成；所以，如果他造了他，他必定造了那不存在之物。如此，按照他们的观点，在圣子存在之前，真理尚未存在，理智之光尚未存在，生命之泉尚未存在，总之一切美好良善之物的本性尚未存在。与此同时，随着每一项的排除，其对立概念就必然存在：若光不存在，就不能否认黑暗存在；其余亦然——每一项更优越的概念之缺席，显然必然导致其对立面的存在。因此必然结论是：当圣父——如异端者所说——\Quote{尚未愿意造圣子}时，圣子所是的一切尚不存在，我们必须说，圣父被黑暗包围而非光明，被谎言包围而非真理，被死亡包围而非生命，被邪恶包围而非良善。因为造物主创造的是不存在之物；\Quote{存在之物}——如\Person{欧诺弥}所言——\Quote{无需受生}；而在被视为对立的事物中，除非更劣者存在，否则更优者不能不存在。这就是异端者的智慧献给圣父的礼物，借此贬低圣子的永恒性，并在圣子\Quote{受造}之前将这一系列邪恶归给天主圣父！\par

但愿没有人想用其他受造物的例子来反驳这个论点所导致教义荒谬的证明。或许有人会说，就如天空不存在时，并没有它的对立面，所以我们并非绝对必须承认：若那为真理的圣子不曾存在，对立面就存在了。对此我们可以回答：天空并没有相应的对立面，除非有人说它的不存在与它的存在相对立。但德行确实与恶行相对（主就是德行）；所以当天空不存在时，并不意味有任何东西存在；但当善不存在时，它的对立面就存在了；因此，说善不存在的人，即使不愿，也必然承认恶存在。\Quote{但圣父}，他说，\Quote{也是绝对的德行、生命、不可接近的光，以及一切言语思想中崇高之物：所以没有必要假设，当独生之光不存在时，与之对应的黑暗就存在了。}但这正是我所说的：黑暗从未存在；因为光从未\Quote{不存在}，因为\Quote{光}——如先知所言——\Quote{常存于光中。}然而，按照异端者的教义，\Quote{非受生的光}是一回事，\Quote{受生的光}是另一回事；前者是永恒的，后者在较晚时间才存在，那么必然的结果是：在永恒的光中，我们没有地方安置其对立面（因为若光常照，黑暗之力在其中便无位置）；而对于那据他们说后来才存在的光，除非从黑暗中发光，否则光不可能照耀；并且从永恒的光到后来出现的光之间的黑暗间隔，将在各方面被清楚地标示。因为若后来之光不是为了某种用途而被造，就没有必要造它；而光的唯一用途就是驱散笼罩的黑暗。那无创造而存在的光，凭自身本性就是其所是；但那受造的光显然是为了他物而存在。因此，它的存在之前必有黑暗，为此光必然被造，且没有任何推理能使\Quote{黑暗不先于独生之光的显现}这一观点显得合理——当然，这是假定他被认为是在后来\Quote{被造}的。这样的教义岂不是全然不敬！因此清楚证明：真理之父并没有在真理不存在的时候造了真理；相反，作为光、真理和一切美善的泉源，他从自己涌流出那独生的真理之光，他位格的荣耀被明确地映照在其中；如此，那些说圣子是后来通过创造附加给天主的人，其亵渎之言在各方面都被驳倒了。\par

\SourceRecord{Translated by H.A. Wilson.；Nicene and Post-Nicene Fathers, Second Series；Vol. 5.；Edited by Philip Schaff and Henry Wace.；Buffalo, NY: Christian Literature Publishing Co.,；1893.；<http://www.newadvent.org/fathers/290109.htm>.}

% structure: p0001 source=h1 target=section reason=page-title

\section{驳斥欧诺米（卷十）}

% structure: p0002 source=h2 target=subsection reason=relative-heading-rank; base=h2

\subsection{一、第十卷讨论对实体之探索的不可企及与不可理解性。在此，他生动地阐述了有关蚂蚁的本性与形成，以及福音中\Quote{我是门}\BibleRef{若10:9}与\Quote{道路}\BibleRef{若14:6}的段落；此外还讨论了天主之名的归属与解释，以及本雅明子孙的片段。}

然而，我们继续我们的主题。稍后，他反驳那些承认人性过于软弱，无法理解不可把捉之事的人，并以高傲的自夸大肆渲染这一话题，轻看我们在这些事上的信仰，言词如下：\Quote{因为绝不能说，若某人的心智因恶毒而蒙蔽，因而既看不见眼前之物，也看不见头顶之上之物，只是对真理的理解力略为胜任，我们就应当据此认为，对实在的发现对全人类而言是不可企及的。}然而我应对他说：凡宣称对实在的发现是可企及的人，自然已通过某种方法和逻辑过程，凭借对存在之物的认识，提升了他自身的理智，并在对相对微小且易于理解之事上受到训练之后，以此预备，将其领悟的幻想投向那超越一切概念的对象。那么，让这位夸耀已获得真实存在之知识的人，向我们解释我们眼前最微不足道的对象的真实本性，以便通过可知之物，使我们相信那隐秘之事：让他用理性解释蚂蚁的本性是什么——其生命是否由气息和呼吸维系，是否像其他动物一样由生命器官调节，其身体是否有骨架结构，骨腔内是否充满骨髓，其关节是否通过肌腱和韧带的张力连接，肌腱的位置是否由肌肉和腺体的包围维持，骨髓是否从头顶沿着椎骨延伸到尾部，是否通过肌腱膜的包围向运动的肢体赋予运动能力；这动物是否有肝脏，以及与之相连的胆囊；是否有肾脏和心脏，动脉和静脉，膜和膈；其外表是光滑还是有毛；是否分为雄雌两性；其身体哪一部分具有视觉和听觉的能力；是否享有嗅觉；其足是不分节的还是有关节的；寿命多长；它们之间如何繁衍，孕期多长；为何并非所有蚂蚁都爬行，也并非都有翅膀，有些属于地面行走的动物，有些则腾空飞翔。那么，让这位夸耀已把握真实存在知识的人，暂时向我们揭示蚂蚁的本性，然后再，并且只在那时，谈论那超越一切理解之能力的本性吧。但如果他尚未通过知识确定这微小蚂蚁的本性，他怎敢夸口通过理性的领悟把握了那亲自统御一切受造物的那一位，并说那些在自己身上承认人性软弱的人，其灵魂的感官已被蒙蔽，既不能触及眼前之物，也不能触及头顶之上的事物？\par

如今且让我们看看，那位认识存在之物者比世上众人多了什么理解。我们来听听他傲慢的言论：— \Quote{主若不自称是\InnerQuote{门}，若无人藉此门进入对圣父的理解与默观，那便是徒然；他若自称为\InnerQuote{路}，若不给愿到圣父那里去的人提供便利，那也是徒然。他如何能成为光，若不照亮人，若不照亮他们灵魂的眼目以理解他自己和那超越之光？} 好，若他在这里列举一些自己头脑中想出的论证，这些论证以其精妙逃避了听者的理解，那么或许有可能被论证的巧妙所欺骗，因为他隐含的思想常常逃过读者的注意。但既然他引用了天主的话，当然无人责备那些相信这些受默感的教导是众人共同财富的人。\Quote{因此，}他说，\Quote{主被称为\InnerQuote{门}，由此可知天主的本质是可为人理解的。}但福音并不承认这一含义。让我们听听天主的话本身。\Quote{我就是门，}基督说；\Quote{谁若经过我进去，必得救，能进能出，能找到牧场。}那么，其中哪一点是关于本质的知识？因为这里说了好几件事，每件都有其特殊含义，无法将它们全部归于本质的概念，免得天主被认为是各种不同元素组合而成的；然而要找出刚才引用的语句中哪一句最适合应用于这一主题，也非易事。主是\Quote{门}，\Quote{经过我，}他说，\Quote{谁若进去，必得救，能进能出，能找到牧场。}我们是否要说，他所说的\Quote{进入}代替了天主的本质，或进去者的\Quote{得救}，或\Quote{出去}，或\Quote{牧场}，或\Quote{找到}？——因为这些词每一个都有其独特的意义，且与其他词的含义不一致。因为进入似乎明显与\Quote{出去}相反，其他短语也是如此。因为\Quote{牧场}按其本义是一回事，\Quote{找到}则是另一回事。那么，这些当中哪一个被认为是圣父的本质？因为一个人当然不能通过说出所有这些在意义上彼此不一致的短语，用不兼容的术语来指代那单纯而不复合的本质。并且，这话如何成立：\Quote{从来没有人见过天主}，以及\Quote{没有人见过，也不能看见\BibleRef{弟前 6:16}}，以及\Quote{没有人见了上主的面而能生存\BibleRef{出 33:20}}，如果\Quote{在门内}或\Quote{门外}或\Quote{找到牧场}表示圣父的本质？因为实际上，他同时是\Quote{护围的门\OriginalTerm{护围的门}{door of encompassing}}和\Quote{防御的屋宇\OriginalTerm{防御的屋宇}{house of defence}}，正如达味称他的那样，并且他亲自接纳进入的人，在他内拯救已进入的人，又亲自领他们出来到德行的牧场，对走在救恩之路的人成为一切，以便他使自己成为每人所需的一切——既是路，又是向导，既是\Quote{护围的门}，又是\Quote{防御的屋宇}，既是\Quote{安慰的水\OriginalTerm{安慰的水}{water of comfort}}，又是\Quote{青绿的草场\OriginalTerm{青绿的草场}{green pasture}}，这在福音中他称为\Quote{牧场}；但我们的新神学家却说，主被称为\Quote{门}是因为对圣父本质的认识。那么，他为何不将\Quote{盘石}、\Quote{石头}、\Quote{泉源}、\Quote{树}等称号也强行纳入同一意义，以便借众多奇异的见证为自己的理论获得证据，因为他完全能够将他关于路、门和光所给出的同样解释应用于其中的每一个？但是，正如我受默感圣经的教导，我大胆断言，那超越一切名号者为我们有许多名号，按照他恩赐对待我们的多样性而接受这些名号，被称为光，当他驱散无知的黑暗时；被称为生命，当他赐予不朽的恩惠时；被称为路，当他引导我们从谬误走向真理时；同样，他也被称为\Quote{力量之塔\OriginalTerm{力量之塔}{tower of strength}}、\Quote{围护之城\OriginalTerm{围护之城}{city of encompassing}}、泉源、盘石、葡萄树、医生、复活，以及诸如此类，都是相对于我们而言，依照他对我们的恩惠以各种方式赐予他自己。但那些眼力超越人类能力的人，他们看见那不可理解的，却忽视那可理解的，当他们用这些称号来解释本质时，就确信他们不仅看见了，而且测量了那没有人见过也不能看见的那一位，却不用他们灵魂的眼目辨识信德——信德是我们观察范围内唯一的事，反将凭推理获得的知识置于信德之上。我听说圣经记载责备本雅明的子孙，他们不尊重法律，却能射箭毫发不差\BibleRef{民 20:16}，我想，这话表明他们热衷于追求无益之事，是远射者，精于瞄准无用虚幻之物，但对于那显然有益的事却无知且漠不关心。因为在我引用之后，历史继续讲述他们所遭遇的事，即他们怎样疯狂追逐索多玛的邪恶，以色列全民武装起来攻击他们，他们便彻底被消灭。我认为，这是一个善意的提醒，警告年轻弓箭手不要希望射中毫发之差，却对信德之门视而不见，宁可放下对不可理解之事的徒劳努力，不要失去手边可得的益处，这益处唯独因信德而得。\par

% structure: p0005 source=h2 target=subsection reason=relative-heading-rank; base=h2

\subsection{二、然后，祂奇妙地向那些不承认祂的人展示了永生，即基督，并将\Person{耶肋米亚}哀悼\Person{约雅金}的悲叹应用于他们，因他们与\Person{蒙丹}和\Person{撒伯里乌}紧密相连。}

但现在，我审视了他论著剩余的部分，我不愿再进行我的论证，因为他的话使我内心颤栗。因为，他想要证明子与永生有所不同；然而，除非永生在于子内，否则我们的信德将被证明是徒劳的，我们的宣讲也是枉然的，圣洗成为多余，殉道者的痛苦全然虚无，宗徒的劳苦对于人的生命无益且无益。因为，他们为何宣讲基督，依据\Person{欧诺米}，基督内并无永生的能力？他们为何提及那些信了基督的人，除非是藉着祂他们才能分享永生？他说：\Quote{因为那些信主的人的理智，超越了所有可感觉和可理解的存在，甚至不能停在子的受生上，而是越过了这一点，在渴望永生中飞向那原初者，急切地要遇见那第一位。}在这段话中，我最该哀悼什么呢？是这些可怜的人不认为永生在于子内，还是他们对独生子的位格抱持如此卑下和属世的观念，以致他们妄想可以在他们的推理中攀上祂的起始，从而凭藉自己理智的力量超越子的生命，并离开主的受生（在它们之下），渴望永生时就越过它飞速前行？因为我所引用的这段话的意思无非就是：人的理智，探究真实存在的知识，并提升自己超越可感觉和可理解的受造物，将原初就存在的那位天主圣言置于自身之下，正如同它把其他一切事物都置于之下一样，而自己却进到那没有天主圣言临在的祂内，藉着心智活动，踏入超越子生命的领域，在那里寻找永生，而那里没有独生子天主。他说：\Quote{因为在渴望永生时，它在思想中被带到了子之外}——显然，仿佛它没有在子内找到它所寻求的。如果永生不在于子内，那么那位说\Quote{我是生命}的，就必被证明是说谎的；或者，祂确实是生命，但不是永生。但并非永恒的东西当然有时间限制。这样的生命，无理性的动物也和人类一样共有。那么，这生命的尊荣在哪里呢？如果甚至无理性的受造物也共享它？如果圣言或神圣理性与生命是同一的，那么如果这生命只是在暂时的生命（根据他们的异端）中居于无理性的受造物内，那会怎样？因为，按照伟大的\Person{若望}，圣言就是生命，但根据他们的异端，这生命是暂时的而非永生的；而且，暂时的生命又存在于其他受造物中，那么逻辑结果是什么？要么无理性的动物是理性的，要么圣言必须被承认为非理性的。我们还需要更多的话来驳斥他们那该受诅咒的恶意亵渎吗？这样的陈述甚至假装要隐藏他们否认主的意图吗？因为，如果宗徒明白地说，非永恒的就是暂时的，而这些人只在圣父的本质中看见永生，并藉着使子远离圣父的本性，他们也把子从永生中割断了，这除了是对主的信德的明显否认和拒绝，还能是什么？而宗徒明确地说，那些\Quote{唯独在此生寄望于基督的人，是众人中最可怜的} \BibleRef{格前 15:19}。那么，如果主是生命，但不是永生，那么这生命当然就是暂时的、仅有一日的，仅在当下起作用；否则，宗徒就哀叹那些寄望于基督的人，如同他们错过了真生命。\par

然而，那些按\Person{欧诺弥}方式受光照的人越过圣子，在推理中被带到祂之外，在默观为外在于且脱离独生者的那一位身上寻求永生。对这些邪恶，除了招致哀叹和哭泣外，还能说什么呢？哀哉，我们怎能不为这悲惨可叹的一代哀叹，他们竟结出如此致命祸害的果实？在古代，热忱的\Person{耶肋米亚}哀悼以色列人民，当他们邪恶地同意带领他们走向偶像崇拜的\Person{约雅金}时，被判处亚述人的掳掠，作为他们非法崇拜的报应，从圣所流放，远离祖业。然而，在我看来，当这位\Person{约雅金}的模仿者将受他欺骗的人引向这种新型偶像崇拜，将他们从祖业——我是指信德——中放逐时，吟唱这些哀歌更加合宜。他们同样，按照圣经记载的对应方式，被掳到\Place{巴比伦}，从上方的耶路撒冷——即从天主教会——进入这有害学说的混乱，因为\Place{巴比伦}的意思是\Quote{混乱}。正如\Person{约雅金}被残害，此人自愿剥夺了真理之光，成为\Place{巴比伦}暴君的掠物，可怜的人从未学到，福音命令我们在圣父、圣子和圣神中同样看见永生，正如圣言如此论及圣父：认识祂就是永生；论及圣子：凡信祂的人有永生；论及圣神：要成为涌到永生的水泉。因此，凡渴求永生的人，当他找到了圣子——我指真正的圣子，而非虚假称为圣子的——就在祂内完全找到了他所渴望的，因为祂是生命，在祂内有生命。但此人，心思如此微妙，眼睛如此锐利，却因极其敏锐的视力在圣子内没有发现生命，反而越过祂，把祂留在身后，视之为寻找目标的障碍，在那他认为真生命不在之处寻找永生！我们能想象什么比这更可憎的亵渎，或更可悲的哀叹缘由？至于撒伯流主义和蒙丹主义的指控一再被强加于我们的教义，这几乎等同于有人将亚诺米派的亵渎归咎于我们。因为若有人仔细考察这些异端的虚假，就会发现它们与\Person{欧诺弥}的错误有很大相似性。因为它们每个都效仿犹太人的教义，既不承认独生天主，也不承认圣神分享他们称为\Quote{伟大}和\Quote{首位}的天主的神性。\Person{撒伯流}称之为三个名字的天主的那一位，\Person{欧诺弥}称之为无始：但两者都不在三个位格的天主圣三中默观神性。那么谁真正与\Person{撒伯流}相似，让阅读我们论证的人判断吧。这些事暂且至此。\par

% structure: p0008 source=h2 target=subsection reason=relative-heading-rank; base=h2

\subsection{三、接着，他展示圣子之生的永恒性，以及祂的本质与生祂者不可分离的同一性，并将\Person{欧诺弥}的愚昧比作玩沙的孩童。}

但是，由于在接下来的部分，他积极搅动他可恶尝试的恶臭，试图证明独生天主\Quote{曾经不存在}，那么，鉴于我们对此头的观点已由先前的论证阐述得相当清楚，最好不再将我们的论证同样陷入同样恶劣的境地，除非从众多论点中选出这一点添加并非不合时宜。他说（某人曾评论说：\Quote{无始的属性同样与圣父的本质相关}），\Quote{论证以与在圣子情况下得出结论时相似的步骤进行。}正统教义显然因敌人的攻击而得到加强，即：我们不应认为无始或受生与本质相同，而是这些属性固然应在主体中默观，而主体按其本有定义是除此之外的某物，既然在主体中找不到差别，因为\Quote{受生}和\Quote{无始}的差别在本质之外，并不影响本质，那么必然结论是，必须承认本质在两位中毫无变易。此外，让我们在已说过的之外探究这一点：他说\Quote{生}与圣父相异是在什么意义上——是将其构想为本质还是行动。若他将其构想为行动，则它显然同样与其结果和作者相连，因为在每一种生产中，人们都看到行动同样在产品与生产者中，出现于效果的产生中，且不与制造者分离。但若他称\Quote{生}为与圣父本质分离的本质，承认主由之而成，那么他显然将这位之于独生者置于圣父的地位，以致在圣子一事上设想出两位父：一位是徒有父名，他称之为\Quote{无始者}，与生毫无关系；另一位，他称之为\Quote{生}，对独生者履行父的角色。\par

欧诺弥自己的陈述比我们的论证更能说明这一点。因为在接下来的话中，他说：\Quote{天主是无受生的，也先于受生者}，稍后又说道：\Quote{因为祂的存在源于受生，在受生之前并不存在}。因此，如果圣父与受生毫无关系，并且圣子的存在源于受生，那么圣父对圣子的自立没有任何作为，与圣子赖以存在的受生完全无关。那么，如果圣父与圣子的受生无涉，他们要么以\Quote{受生}之名给圣子杜撰另一个父亲，要么以他们的智慧使圣子成为自生和自产的。你看到这个人在他的论证中到处向我们暴露无知的思想混乱，他的亵渎如何游荡在许多条道路上，或者更确切地说，在无路之处，没有任何可靠的引导向目标前进；正如人们可以看到婴儿的情形，当他们以幼稚的游戏用沙子模仿建造房屋时，他们所建造的并非按照任何计划或艺术规则来模仿原型，而是首先随意且愚蠢地做出某种东西，然后商量如何称呼它——我在我们的作者身上看到了这种洞察力。因为他在头脑中首先想到什么就收集亵渎的言辞，像一堆沙子，然后开始考虑他那不可理解的亵渎会走向何方，它从他所说的话中自发地增长，没有任何理性顺序。因为我不认为他最初提议将受生当作一个实际的自立体来发明，用它来代替圣父站在圣子本质上，也不认为这位修辞学家计划让圣父被视为与圣子的受生无关，也不是故意引入自生的荒谬。但是，所有这类荒谬的说法都是我们的作者不经思考地发出来的，因此，就这些错误而言，这个犯错误的人甚至不值得多加反驳，因为他——借用宗徒的话——\Quote{既不知道所说的是什么，也不知道他所肯定的是什麼\BibleRef{弟前 1:7}}。\par

\Quote{因为祂的存在源于受生}，他说，\Quote{在受生之前并不存在}。如果他在这里用\Quote{受生}一词指圣父，我就同意他，并无反对。因为人们可以用两种说法表达同一意思，或说\Person{亚巴郎}生了\Person{依撒格}，或说\Person{亚巴郎}是\Person{依撒格}的父亲。既然为父等同于已生，如果有人从一种说法转换成另一种，父性就会显示为等同于受生。因此，如果\Person{欧诺弥}所说的是：\Quote{祂的存在源于圣父，祂并不在圣父之前}，那么这话是正确的，我们赞同。但如果他在这句话中又回到我们之前谈到的那个受生，说它与圣父分离而与圣子关联，那么我认为花时间考虑这不可理解的话是浪费时间。因为无论他认为受生是一个自存的对象，还是通过这个名称他的思想被引向那没有实际存在的东西，直到今天我仍无法从他的话中找出答案。因为他那流动而无根据的论证同样适用于两种假设，随着思想者的幻想而倾向这一边或那一边。\par

% structure: p0012 source=h2 target=subsection reason=relative-heading-rank; base=h2

\subsection{四、在此之后，他证明那真正存在且在圣父怀中的圣子是单纯而不复合的；那将我们从束缚中赎回的主并不在圣父的管辖之下，也不处于奴隶状态；否则，不仅祂自己，就连那在圣子内并与祂合一的圣父也必然是奴隶；并且\Quote{存在}这个词是由\Quote{是}这个词派生而来的。在出色而卓著地讨论了这一切之后，他结束了本书。}

但他亵渎中最严重部分尚未审查，这些是紧随其后添加的。好吧，让我们逐句考虑他的话。然而我不知道如何敢让自己开口说出这位敌对基督者那可怕而无神的言语。因为我害怕，如同一些有害的药物，有害苦味的残留会沉积在话语通过的嘴唇上。宗徒说：\Quote{来到天主面前的人，必须相信祂是\BibleRef{希 11:6}}。因此，真实存在是神性的特殊区别。但\Person{欧诺弥}使那真正存在者要么根本不存在，要么不是以恰当的方式存在，这恰恰等同于完全不存在；因为那不存在于恰当方式中者，实际上根本不存在；例如，一个人在梦中被说成\Quote{跑步}，他在那个状态中幻想自己在赛跑中用力，但由于他虚假地模拟真实赛跑的样子，他幻想自己在跑步并不因此就是赛跑。但即使在不精确的意义上这样称呼，这名称仍是虚假地赋予了它。因此，凡敢断言独生天主要么不以恰当方式存在，要么根本不存在的人，显然从他的信条中抹去了对祂的一切信仰。因为谁还能相信某种不存在的东西呢？或者谁还会投奔那被真正的主的敌人证明为不恰当且无实质的祂？\par

但为使我们的陈述不被认为对对手不公，我将把不敬者的言论与之并列，其文如下：——\Quote{祂在那存在者怀裡，祂在起初就与天主同在，并非存在，或至少不是严格意义上的存在，尽管巴西略忽略了这一区别和附加，互换地使用\InnerQuote{那存在者}这一称号，违背真理——}\newline{}您说什么？祂在圣父内却不存在？祂在起初，在圣父怀裡，正因为祂在起初、在圣父内、在那存在者的怀中被辨识，所以严格地说祂不存在，因为祂在那存在者内？唉，这些空洞无理的主张！我们第一次听到这虚妄的胡言乱语——那万有藉以存在的\Person{主}严格地说并不存在。我们还未说完这可怖的陈述；后面还有更惊人的：他不仅断言祂不存在，或严格地说不存在，还断言那被构想为居于其中的本性是多样且复合的。因为他说\Quote{不存在，或不是单纯的。}那不具单纯性的东西显然是多样且复合的。那么，同一位怎能既不存在又在本质上是复合的呢？因为他们必须选择二者之一：如果他们断言祂不存在，就不能说祂是复合的；如果他们断言祂是复合的，就不能剥夺祂的存在。但为使他们的亵渎具有多种多样的形式，他们每当意欲将祂与\OriginalTerm{那存在者}{the Existent}对照时，便抓住每一个无神的观念，断言严格地说祂不存在，并在与那不复合的本性的关系中否认祂具有单纯性：——\Quote{不存在，不是单纯地存在，不是严格意义的存在。}在那些违背圣言、背弃信仰的人中，有谁曾如此慷慨地吐出否认\Person{主}的话语？他起来与\Person{若望}的神圣宣告相争。因为每当后者证实圣言\BibleQuote{是}时，他就对那一位施加一个对立的\Quote{不是}。他反对我们圣父\Person{巴西略}的神圣嘴唇，指控他\Quote{忽略了这些区别}，当他说那在圣父内、在起初、在圣父怀裡的存在时，认为\Quote{在起初}和\Quote{在圣父怀裡}的附加阻止了那一位的真实存在。虚妄的学问！欺骗的导师教的是何等东西！他们向听众介绍何等奇怪的教义！他们教导：在别的事物中的东西不存在！所以，欧诺弥，既然你的心和脑在你内，按照你的区别，它们两者都不存在。因为如果独生子天主严格地说不存在，理由就是祂在圣父怀裡，那么一切在别的事物中的东西都被剥夺了存在。但当然，你的心在你内，而非独立存在；因此，按照你的观点，你必须要么说它根本不存在，要么说它不严格地存在。然而，他语言的愚昧和亵渎是如此严重和明显，甚至在我们的论证之前就显然可见，至少对一切有理智的人而言是如此：但为使他的愚蠢和不敬更显明，我们将在前述内容上再添加一些。如果只有那圣经的言语证实其存在、脱离与任何其他事物关系的事物才可以说严格地存在，那为何他们像那些提水的人一样，在能够喝水时却渴死？即使这个人，尽管手边就有治疗他对圣子亵渎的解毒剂，却闭上眼睛跑过去，仿佛害怕得救，并指控\Person{巴西略}不公，因为他压制了修饰语，只单独引用\BibleQuote{是}来指独生子。而实际上，他完全可以看见\Person{巴西略}所看见的、每个有眼睛的人都看见的东西。在此，崇高的\Person{若望}在我看来是受到先知性的感动，为堵住那些基督斗士的口，他们基于这些附加否认基督严格意义上的存在，他们简单无修饰地说\BibleQuote{圣言是天主，是生命，是光}，并非仅仅说祂在起初、与天主同在、在圣父怀裡，以致使主绝对的存在通过其关系被取消。但他断言祂是天主，这绝对宣告脱离与任何其他事物的关系，从而割断了那些在推理中陷入不敬者的所有遁词；除此之外，还有另一件事更能证明我们对手的恶意。因为如果他们认为存在于某物中是严格意义上不存在的标志，那么他们当然承认连圣父也不是绝对存在的，因为他们在福音中学到，正如圣子住在圣父内，同样圣父也住在圣子内，按主的话。因为说圣父在圣子内，相当于说圣子在圣父怀裡。顺便让我们进一步探问：当圣子如他们所说\Quote{不存在}时，圣父的怀裡装着什么？因为他们必然要么承认它是满的，要么认为它是空的。如果怀是满的，那么圣子当然就是充满怀的那位。但如果他们想像圣父怀中有某种空虚，他们无非是在说祂通过增长而变得完美，即从空虚和欠缺的状态进入充满和完美的状态。但\BibleQuote{他们不知道，也不明白}，\Person{达味}论那些\BibleQuote{仍在黑暗中行走}的人说。因为那已成为真光仇敌的人不能使自己的灵魂留在光中。正因如此，他们没有察觉逻辑上随手可得的能纠正他们不敬的东西，如同被盲目击打，像索多玛人一样。\par

但他也说圣子的本质受圣父支配，他的原话如下：—\Quote{因为那因父而存在并生活的一位，并不将此尊位归于自己，因为那支配祂的本质将自有者这一概念引向自身。}如果这些道理为某些\Quote{外邦}智者所认可，那么福音书和圣经其他的教训就不应受到任何干扰。因为基督徒的信德与被愚弄的智慧之间有什么相干呢？但如果他倚靠圣经的支持，就让他从圣书中指出这样一个声明，我们便会沉默。我听见保禄大声呼喊：\Quote{只有一位主耶稣基督\BibleRef{格前 8:6}}。但欧诺弥却反对保禄，称基督为奴仆。因为除了受支配和被控制之外，我们不承认有别的奴仆标记。奴仆无疑是奴仆，但奴仆不可能从本性上是主，即使这个术语被不精确地用于祂。我何必要引述保禄的声明来证明主的王权呢？因为保禄的主亲自告诉祂的门徒，祂真正是主，祂接受了那些称祂为师傅和主的人的宣认。因为祂说：\Quote{你们称我为师傅和主，你们说得对，因为我就是\BibleRef{若 13:13}}。同样，祂吩咐他们称圣父为父，说：\Quote{不要称地上的人为师傅，因为你们的师傅只有一位，就是基督；也不要称地上的人为父，因为你们的父只有一位，就是天上的父\BibleRef{玛 23:8}}。我们被夹在他们之间，该听从哪一个呢？一方面，主亲自和那有基督在他内说话的人\BibleRef{格后 13:3}吩咐我们不要认为祂是奴仆，而要尊荣祂如同尊荣圣父一样；另一方面，欧诺弥对主提出控告，声称祂是奴仆，他说那肩负宇宙治理之权的祂是在管辖之下。我们的选择还有怀疑的余地吗？或者哪个更有利的决定不重要吗？欧诺弥，我该轻视保禄的劝告吗？我该认为真理的声音不如你的欺骗可靠吗？但\Quote{我若没有来对他们说话，他们就没有罪。}既然如此，祂已对他们说了话，真实地宣告自己是主，并且祂不是被虚假地称为主（因为祂说\Quote{我是}，而不是\Quote{我被称为}），他们还需要做什么呢？那不可避免的报应正因为预先警告而降临？\par

但也许为回应这一点，他会再次搬出他惯用的逻辑，并说同一存有既是奴隶又是主，既受统治权支配，又凌驾于其余之上。这些深奥的区分在十字路口被谈论，由那些迷恋谎言的人传播，他们用日常生活的实例来证实自己关于神性的空想。既然今世之事提供了此类安排的例子（例如，在一个富裕的家庭中，可以看到那位更勤快、更忠心的仆人奉主人之命被立为其他同仆之上，从而获得对同等级别中其他人的优越权），他们便将这一观念转用到关于神性的教义上，所以独生子天主，虽受其尊长的主权管辖，却在低于祂的人的管理方面丝毫未受其主权者权威的阻碍。但让我们告别这类哲学，并按照我们理智的尺度来讨论这一点。他们是否承认圣父在本性上是主？还是认为祂是通过某种拣选达到这一地位的？我认为一个稍有理智的人都不会疯狂到不承认万有的天主的宰制权是祂本性所有。因为那本性单一、纯全、不可分割者，无论它是什么，它整个地就是它自身，不是通过某种变化过程成为另一物，而是永恒地处在它所是的状态中。那么，他们关于独生子是怎样的信念呢？他们承认祂的本质是单纯的，还是认为其中有某种组合？如果他们以为祂是由许多部分组成的一种多形态之物，肯定连神性的名称都不会给祂，而会将基督的教义拖向有形质和物质的概念；但如果他们同意祂是单纯的，那么在主体的单纯性中又怎能承认相反属性的共存呢？因为就像生与死的矛盾对立不允许中间状态一样，在区分特征方面，统治与奴役也是截然对立、不可调和的。因为若单独考虑其中每一项，就无法恰当地制定一个同时适用于两者的定义；而当事物的定义不同时，它们的本性也必然不同。那么，如果主在本性上是单一而纯全的，主体中又怎能找到相反属性的结合呢？就像奴役与统治混合那样？但如果根据圣人们的教导，祂被承认为主，那么主体的单纯性就证明祂不可能有份于相反的状态；而如果他们硬说祂是奴隶，那么他们将主的称号归于祂就是徒劳的。因为本性单一的东西不会割裂成矛盾属性。但如果他们断言祂是一物，却被称作另一物，即在本性上是奴隶而仅在名义上是主，那就让他们大胆说出这一声明，使我们免于长期作答的辛劳。因为谁会如此悠闲地处理虚无之事，以致使用论证来证明显而易见、毫不含糊的东西呢？如果有人自己控告自己犯了谋杀罪，控告者就不必费心用证据来证明他血债累累。同样，当我们的对手在邪恶中走得太远时，我们不再在审查他们的案件后进行反驳。因为那宣称独生子是奴隶的人，实际上是说祂与他本人同为仆人：因而必然产生双重的丑行。因为他要么蔑视他的同仆并否认信仰，摆脱了基督宰制的轭，要么他向那奴隶下拜，背离那自主的、无主之性，从而某种程度上崇拜自己而非天主。因为如果他看到自己在奴役中，而他所崇拜的对象也在奴役中，他当然是在看自己，在他所崇拜的对象中看到自己的全部。但还有什么能计算出这种教义邪恶必然伴随的所有其他恶果呢？因为谁不知道那在本性上是奴隶、并在主人的强制下从事劳作的人，甚至不能免除恐惧的情感？受默感的宗徒为此作证，他说：\Quote{你们没有领受奴役的精神，以致再恐惧\BibleRef{罗 8:15}。} 因此，人们会发现，他们按照人的相似，也将恐惧的情感归于他们的同仆天主。\par

这就是异端的天主。但是，我们这些——用宗徒的话来说——被基督（祂已解救我们脱离了束缚）召叫到自由中的人，圣经教导我们该如何思考，我将用几句话陈述出来。我从默感的教育出发，大胆宣告：圣言不希望我们仍是奴隶，因为我们的本性现在已经变得更好；祂取了我们的一切，条件是将祂的给予我们作为回报；正如祂取了疾病、死亡、诅咒和罪恶，祂也取了我们的奴役，但不是为了自己拥有祂所取的，而是为了净化我们本性中的这些邪恶，我们的缺陷被祂无瑕的本性所吞没和消除。因此，正如在我们所盼望的生命中，既没有疾病，也没有诅咒，更没有罪恶和死亡，奴役也必将与这些一同消失。我所说的话是真实的，我呼求真理解作证，祂对门徒说：\Quote{我不再称你们为仆人，而是朋友。} 那么，如果我们的本性最终将摆脱奴役的羞辱，万有的主怎会被这些心智错乱之人的疯狂和愚昧贬为奴隶呢？他们作为逻辑推论，当然会断言祂不知道圣父的旨意，因为祂关于奴隶的声明告诉我们：\Quote{仆人不知道主人所做的事。} 但当他们这样说时，让他们听到：圣子在自己内拥有一切属于圣父的，看见圣父所做的一切，圣父的良善中没有一样是圣子所不知道的。因为祂怎么能缺少任何属于圣父的东西呢？既然祂在自己内完整地拥有圣父。因此，如果\Quote{仆人不知道主人所做的事，}而祂在自己内拥有一切属于圣父的，那么这些沉醉于烈酒的人最终应当清醒，现在，如果以前从未有过，就让他们仰望真理，看到那拥有一切属于圣父的，是万有的主，而不是奴隶。因为那没有主人管辖的位格，怎能承受奴役的烙印？万有的君王怎能没有与自己同等荣耀的形象？耻辱——因为奴役即耻辱——怎能构成真光明的光辉？君王的儿子怎能生而为奴？不，不是这样的。但是，正如祂是光之光、生命之生命、真理之真理，祂也是如此的主之主、王之王、天主的天主、至高之至高；因为祂在自己内完整地拥有圣父，圣父在自己内所有的，祂也肯定拥有；而且，圣子所有的一切都属于圣父，所以天主光荣的敌人不可避免地被迫：如果圣子是奴隶，他们也将圣父拖入奴役。因为圣子的任何属性都完全属于圣父。祂说：\Quote{因为凡是我的，都是你的；你的也是我的 \BibleRef{若 17:10}。} 那么这些可怜的人会说什么？哪个更合理：那说了\Quote{你的都是我的，并且我在他们中受到了光荣 \BibleRef{若 17:10}}的圣子，应在圣父的主权中受到光荣，还是通过圣子奴役中的贬抑而向圣父施加侮辱？因为那在自己内包含圣子所有的一切、并且自己就在圣子内的，不可能不完全处于圣子的奴役中，并在自己内有奴役。这就是欧诺米哲学所取得的成果，他通过这种方式将奴役的侮辱加给他的主，同时将同样的贬抑加给圣父无瑕的光荣。\par

然而，让我们再次回到他论文的进程。欧诺米关于独生子说了什么？他说祂\Quote{没有占取这\OriginalTerm{尊严}{dignity}，}因为他把\Quote{\OriginalTerm{存在}{being}}的称谓称为\Quote{\OriginalTerm{尊严}{dignity}}。令人震惊的哲学！在古往今来所有人类中，无论是希腊人还是外邦哲学家，无论是当代人还是历代人，有谁曾把\Quote{\OriginalTerm{存在}{being}}称为\Quote{\OriginalTerm{尊严}{dignity}}？因为按照所有使用语言的人的共同习惯，一切被视为实际存在的事物都被说成\Quote{是}；从\Quote{是}这个词形成了\Quote{\OriginalTerm{存在}{being}}这个术语。但现在，\Quote{\OriginalTerm{尊严}{dignity}}这个表达以新的方式被用于\Quote{\OriginalTerm{存在}{being}}所表达的概念。因为他说：\Quote{圣子，祂因圣父而存在并活着，并没有占取这\OriginalTerm{尊严}{dignity}。}他没有任何圣经支持他的说法，也没有通过任何逻辑推理的过程将他的陈述引向这样一个荒谬的结论，而是仿佛他肠子里进了些胀气的食物，他粗俗而杂乱地打出其亵渎的嗝，像一股难闻的气息。\Quote{他没有占取这\OriginalTerm{尊严}{dignity}。}让我们承认把\Quote{\OriginalTerm{存在}{being}}称为\Quote{\OriginalTerm{尊严}{dignity}}这一点。那么怎么样？那\Quote{是}的不是\Quote{占取}\OriginalTerm{存在}{being}？欧诺米说：\Quote{不，因为祂因圣父而存在。}那么你岂不是说那没有\Quote{占取}\OriginalTerm{存在}{being}的，是不存在的？因为\Quote{不占取}与\Quote{疏离于}具有相同的效力，并且这两个概念的相互对立是明显的。因为\Quote{本有的}不是\Quote{疏离的}，而\Quote{疏离的}不是\Quote{本有的}。因此，那没有\Quote{占取}\OriginalTerm{存在}{being}的，显然是疏离于\OriginalTerm{存在}{being}的；而疏离于\OriginalTerm{存在}{being}的，是不存在的。\par

但他的荒谬性的有力证明，他以这些话提出：\Quote{那个连祂也控制着的本质，将\OriginalTerm{存在者}{\GreekText{τὸ} \GreekText{ὄν}}的概念吸引到自己身上。}我们暂且不谈这里组合的笨拙：让我们审视他的严肃含义。有什么论证曾经证明了这一点？他多余地反复向我们陈述圣父的本质对圣子有主权。哪位福音作者是这教义的赞助者？什么样的辩证过程引导我们到达它？什么前提支持它？有什么推理路向曾经通过任何逻辑结论证明\Person{独生天主}是受管辖的？\Quote{但是，}他说，\Quote{那个对圣子有支配权的本质，将\OriginalTerm{存在者}{\GreekText{τὸ} \GreekText{ὄν}}的概念吸引到自己身上。}存在者的吸引是什么意思？而\Quote{吸引}这个说法怎么会被抛到他之前所说的话上面？当然，谁考虑词语的力量，就会自己判断。然而，关于这一点我们将什么都不说：但我们将再次抓住那个论点，即他不将本质存有赐予那位他不留下\OriginalTerm{存在者}{\GreekText{τὸ} \GreekText{ὄν}}头衔的祂。他为什么与影子进行无谓的争斗，争论不存在的东西是这样或那样？因为不存在的东西当然既不与其他东西相似，也不不相似。但是，虽然他承认祂是存在的，却禁止祂被这样称呼。唉，就一个词的声音讨价还价的徒然精确，而在更重要的事情上作出让步！但是，那位据他说对圣子有主权的祂，在什么意义上\Quote{将\OriginalTerm{存在者}{\GreekText{τὸ} \GreekText{ὄν}}的概念吸引到自己身上}？因为如果他说圣父吸引祂自己的本质，这个吸引过程是多余的：因为存在已经是祂的，无需被吸引。另一方面，如果他的意思是圣子的存在被圣父吸引，我无法理解存在怎么能从存在者那里被扭夺，并转移到那位\Quote{吸引}它的祂身上。他是不是在做梦，梦到\Person{撒柏略}的错谬，以为圣子并不自存，而是被描绘到圣父的位格存在上？这就是他在\Quote{存在者的概念被那对圣子行使支配的本质所吸引}这个表述中的意思吗？或者，他难道不否认圣子的位格存在，却仍然说祂与\Quote{\OriginalTerm{存在者}{\GreekText{τὸ} \GreekText{ὄν}}}这个术语所传达的意义分离？然而，\Quote{\OriginalTerm{存在者}{\GreekText{τὸ} \GreekText{ὄν}}}怎么能与存在的概念分离？因为只要事物是它所是的，本性并不容许它不再是它所是的。\par

\SourceRecord{Translated by H.C. Ogle and H.A. Wilson.；Nicene and Post-Nicene Fathers, Second Series；Vol. 5.；Edited by Philip Schaff and Henry Wace.；Buffalo, NY: Christian Literature Publishing Co.,；1893.；<http://www.newadvent.org/fathers/290110.htm>.}

% structure: p0001 source=h1 target=section reason=page-title

\section{\WorkTitle{驳斥欧诺弥（卷十一）}}

% structure: p0002 source=h2 target=subsection reason=relative-heading-rank; base=h2

\subsection{第十一卷显示，\Quote{善}这一称号并非如\Person{摩尼}与\Person{巴德桑}的仿效者\Person{欧诺弥}所称，唯独归于圣父，也归于圣子，祂以良善和仁爱塑造了人，并借其十字架与死亡再造了人。}

现在我们继续讨论他论证的下一步：\Quote{……独生子亲自将那唯独应归于圣父的称号归于祂。因为那位教导我们，\InnerQuote{善}这一称谓唯独属于祂自己善的及一切善的原因，且始终如此，并将所有已存在的善都归于祂的，不会轻易将万有之权柄和\OriginalTerm{自有者}{Existent}的称号据为己有。}好，只要他仍以某种面纱掩饰其亵渎，并竭力以辩证法的迷宫不知不觉地缠住受骗的听众，我就认为有必要监视他的不公和隐秘行径，并尽可能在论证中揭露潜伏的祸害。但现在他已剥去所有伪装，公开以明确的言辞宣告其渎神之言，我认为对公然显示不虔敬的人，再费无用之功去运用逻辑反驳模式已是多余。因为还有什么方法能比他们自己在著作中亲手向我们展示的更能有效证明他们的恶意呢？他说唯独圣父配得\Quote{善}的称号，这称谓唯独归于祂，理由是甚至圣子自己也同意善唯独属于祂。原告为我们作了辩护：因为在我先前的陈述中，读者可能以为我表现出某种无礼的傲慢，试图证明那些反抗基督的人将圣子与圣父的善隔绝。但我认为现在已由对手的告白证明，我们提出这样的指控并未不公。因为他说\Quote{善}的称号唯独归于圣父，这样的称呼唯独适于祂，藉此揭示其真实含义，也就公开了他先前隐藏的恶行。他说\Quote{善}的称号唯独适于圣父。他是指该称号带有其表达的意义，还是脱离其本义的称号？如果一方面，他仅以一种特殊意义将\Quote{善}的称号归于圣父，那么他因允许圣父仅有一个空洞名称的声音而值得怜悯。但如果他认为\Quote{善}这一术语所表达的概念只属于天主圣父，那么他因在自己的意见中重燃摩尼异端的瘟疫而应受憎恶。因为正如健康与疾病，善与恶亦以相互毁灭的方式存在，以致一者的缺席即为另一者的在场。因此，若他说善唯独属于圣父，他就将这些排除在除圣父之外的一切存在对象之外，以致独生子天主也连同一切被摒除于善之外。因为正如肯定唯人会笑的人，暗示其他动物无此属性，同样，断言善唯独在圣父内的人，便将万物与该属性分离。因此，如\Person{欧诺弥}所言，若圣父唯独有权享有\Quote{善}的称号，该术语将不适用于任何其他事物。但意志的每一冲动，或按善运作，或趋向其反面。因为不倾向任何一方，而保持平衡，是无生命或无感觉之受造物的特性。若圣父唯独是善，其善并非获得，而是出于本性，且若圣子（如异端所欲）不分享圣父的本性，那么不分享圣父善之本质者自然同时被排除于\Quote{善}的称号之外。但既无善之本质亦无善之名称者——他肯定不是不为人知的，即使我避免渎神之言。因为所有人都清楚，\Person{欧诺弥}如此急切的目的，是要将恶的嫌疑及其对立面引入对圣子的概念中。因为那不善良者属于何种名号，对每个有理性者都是明显的。正如不勇敢者是懦夫，不义者是不义的，不智者是愚昧的，同样，不善良者显然拥有相反的名称，而基督之敌正希望将独生子的概念推向这一点，成为教会的另一个\Person{摩尼}或\Person{巴德桑}。关于这些言论，我们认为说出与沉默同样无效。因为即使说无数的话并唤起一切可能的论证，也不能说出比对手公开无掩饰地宣告的更不利于他们的话。因为人们能发明比否认祂是善更恶毒的指控吗？\Quote{祂虽具有天主的形体，却不以自己与天主同等为应当强夺的}，却屈尊于卑贱的人性，且完全出于对人的爱。请告诉我，\Quote{你们为何这样报答上主？}\BibleRef{申 32:6}，（我要借用梅瑟的话对以色列人说）；祂难道不善吗？当你们是无生命的尘土时，祂以天主般的美装饰你们，并将你们立为祂权能的肖像，赋予灵魂。祂难道不善吗？祂为了你们取了奴仆的形体，并因那摆在面前的喜乐\BibleRef{希 12:2}，不畏惧承担因你们罪过当受的苦难，将自己作为赎价交付于你们的死亡，并为我们成了诅咒和罪。\par

% structure: p0004 source=h2 target=subsection reason=relative-heading-rank; base=h2

\subsection{二、他也巧妙地从福音书论及\Quote{善师}的段落、葡萄园的比喻、\BibleBook{依撒意亚}{先知书}和\BibleBook{保禄}{书信}中指出：在天主性内并不存在如\Person{欧诺弥}的盟友\Person{玛尔基翁}所想象的善恶二元论；并宣告圣子并不拒绝\Quote{善}或\Quote{存在者}的名号，也不承认自己与圣父分离，反而对一切受造之物同样拥有权柄。}

就连你们意见的庇护人玛尔基昂本人，在这件事上也不支持你们。诚然，他同你们一样主张天主的二元论，认为一位与另一位在本质上不同，但较为文雅的观点是将善归于福音的天主。而你们却将独生子天主与善的本质分开，为的是在教义的邪恶上甚至超越玛尔基昂。然而，他们却声称圣经支持他们，并说当他们因使用圣经的原话而被指控时，他们受到了不公正的对待。因为他们说，主亲自说过：\BibleQuote{没有良善的，唯有一个，就是天主}\BibleRef{玛 19:17}。为此，为使谬误不致胜于天主圣言，我们将简要查考福音中的这段经文。这段记述视主对之说话的那个富人为年轻人——我猜，是那种倾向于享受今生快乐的人——并且依恋其财物；因为经文说，他因放弃所有财物的劝告而忧伤，不愿以财产换取永生。这人听说有位永生导师在附近，便来见他，期望能永远享受奢华，生命无限延长，以\Quote{善}的称号奉承主——我应该说，奉承的不是我们所理解的主，而是当时以仆人形体显现的主。因为他的品性不足以使他刺透外在血肉的帷幕，看到内里的天主性圣所。于是，洞察人心的主看出了那年轻人以恳求者的身份接近祂的动机——他不是以专注于天主性的灵魂而来，而是恳求那\Quote{人}，称祂为\Quote{善师}，因为他希望从祂那里学到一些能阻挡死亡临近的学问。因此，那被如此恳求的主按照所受的称呼回答他，是合情合理的。因为恳求不是向天主圣言发出的，所以相应地，回答是由基督的人性向恳求者发出的，以此给那年轻人双重的教训。祂以同一个回答教导他两件事：一是应以生活而非奉承的言语来敬畏和尊崇天主性，即遵守诫命并以所有财产换取永生；二是人性因罪恶而沉沦于邪恶，被排除在\Quote{善}的称号之外。为此祂说：\BibleQuote{你为什么称我为善？}在回答中，祂用\Quote{我}一词暗示了那包围祂的人性，而将善归于天主性，祂就明确宣告自己是善的，因为福音宣告祂是天主。因为如果独生子被排除在天主的称号之外，那么认为祂也与\Quote{善}的名称无关，也许并非荒谬。但是，既然先知、福音作者和宗徒都大声宣告独生子的天主性，而且主自己也证实善属于天主，那么，分享天主性的祂，怎能不也分享善呢？因为先知、福音作者、门徒和宗徒都承认主是天主，这一点，凡未领受神圣奥秘者，无需特意告知。因为谁不知道，在圣咏44篇中，先知以言语肯定基督是天主，由天主所傅油？再者，凡熟悉先知预言者，谁不知道依撒意亚在其它段落中如此公开宣告圣子的天主性，他说：\BibleQuote{撒巴人，身材高大的人，要归顺你，必在你后面行走，带着锁链走过，并要在你内祈祷，因为天主在你内，除你以外没有天主；因为你是天主。}因为除了独生子之外，还有哪一位天主祂自己内有天主，而祂自己又是天主呢？让那些不听从预言的人说说吧。至于厄玛奴耳的解释，多默在认出主之后的宣认，以及若望崇高的言辞，这些连外教人也明白，我就不说了。不，我甚至认为没有必要详细引述保禄的话，因为这些话可以说是众口相传的，他不仅称主为\Quote{天主}，还称祂为\Quote{伟大的天主}和\Quote{超越万有的天主}，对罗马人说：\BibleQuote{他们的祖先也是他们的，按血统，基督也出自他们，祂是超越万有的天主，永远受赞美}\BibleRef{罗 9:5}，又写信给门徒弟铎说：\BibleQuote{随着我们的救主耶稣基督、伟大的天主的显现}，又给弟茂德写信，明确宣告：\BibleQuote{天主在肉身显现，在圣神内成义。}既然事实已从各方面证明独生子天主是天主，那么，那个说善属于天主的人，怎能竭力证明圣子的天主性与这一称号无关呢？何况主在葡萄园雇工的比喻中，实际上已为自己声称\Quote{善}的称呼。因为在那里，当先劳作的工人对所有人都得到同样报酬不满，并视最后者的好运为自己的损失时，公义的审判者对其中一个抱怨者说：\BibleQuote{朋友，我没有亏负你；我不是和你讲定一个德纳吗？拿着你的走吧！我愿意给这最后来的和给你的一样。难道我不能拿我的财物做我所愿意的吗？还是因为我善，你就眼红吗？}当然，没有人会争辩，按照功过分配报酬是审判者的特殊职能；所有福音的门徒都同意独生子天主就是审判者；因为祂说：\BibleQuote{圣父不审判任何人，却将一切审判权交给了圣子。}但他们并不反对圣经。因为他们说\Quote{独一}一词绝对是指圣父。因为祂说：\BibleQuote{没有良善的，唯有一个，就是天主。}那么，真理难道会缺乏力量为自己辩护吗？当然，有许多方法很容易就能揭露这个诡辩的欺骗性。因为那说这话论及圣父的，也对圣父说了另一句话：\BibleQuote{我的一切都是你的，你的一切也是我的，并且我在他们内得了光荣}\BibleRef{若 17:10}。现在，如果祂说圣父的一切也是圣子的，而善是圣父的属性之一，那么如果圣子没有善，祂就没有一切，他们就会说真理在说谎；或者，如果怀疑真理本身被卷入虚假是邪恶的，那么那声称圣父的一切为自己的，就以此断言祂不在善之外。因为那在自己内有圣父、并拥有圣父一切所有的，显然和\Quote{一切}一起拥有圣父的善。因此，圣子是善的。但他说：\BibleQuote{没有良善的，唯有一个，就是天主。}这是我们的对手所主张的；我自己也不拒绝这个说法。然而我并不因此否认圣子的天主性。但那些承认主是天主的人，恰恰以这个宣认也断言祂的善。因为如果善是天主的属性，而主是天主，那么根据前提，圣子被证明是天主。我们的对手说：\Quote{但是\Quote{独一}一词将圣子排除在善之外。}然而，很容易表明，就连\Quote{独一}一词也不将圣父与圣子分开。因为在所有其它情况下，\Quote{独一}一词确实带有不与任何其它事物关联的含义，但在圣父和圣子的情况下，\Quote{独一}并不意味孤立。因为祂说：\BibleQuote{我与圣父是独一。}如果善是独一的，而在圣父和圣子之内有一种特殊的合一，那么，主在将善归于\Quote{独一}时，就借着\Quote{独一}一词也为祂自己要求了\Quote{善}的称号，因为祂与圣父是独一的，并未与同一本性分离。\par

% structure: p0006 source=h2 target=subsection reason=relative-heading-rank; base=h2

\subsection{三、接着，他揭露了欧诺弥的愚昧，以及他论点的混乱与荒谬——当他把圣子称作\Quote{自有者的天使}，并声称圣子的地位低于天主性，如同圣子优于受造之物一样。在此，有一段有力而令人信服的反驳与义愤填膺的斥责，指出那向\Person{梅瑟}传授神谕的，正是那自有者本身，即独生子；祂回应\Person{梅瑟}的恳求\BibleQuote{若你不亲自与我们同行，就不要从这里领我们上去}时说\BibleQuote{我也必这样做，就是你所求的}；祂在\Person{梅瑟}和\Person{依撒意亚}书中也被称为\Quote{天使}；其中引用了经文\BibleQuote{有一个婴儿为我们诞生了}。}

但是，为了完全揭露我们这位大人物作者的学究气与文化修养，我们将逐句考察他陈述其观点的方式。他说：\Quote{圣子}，\Quote{不将自有者的尊位归给自己}，他用\InnerQuote{尊位}一词指存在本身（他何等善于恰当用词！）；并且，由于祂\InnerQuote{因圣父的缘故}，他说祂与自身分离，理由是那超越祂的本质把自有者的概念吸引到自身。这就好比有人说，一个被金钱买来的人，就他自己\Emph{是}于其存在而言，并不是被买的人，而是买主，因为他的本质人格存在被他人的本性所吸收，而那他人获得了对他的权力。这就是我们这位神学家的崇高观念：但他陈述的证明是什么呢？……他说：\Quote{独生子自己将惟独应归于圣父的称号归给圣父}，然后他提出只有圣父是善的观点。在此，圣子哪里否认了\Quote{自有者}的称号呢？然而这正是\Person{欧诺弥}所要达到的目的，他逐字逐句如下：\Quote{因为那教导我们\InnerQuote{善}的称谓只属于那自身善和一切善的原因，并且始终是善的，并将一切生成的善都归于祂的一位，祂不会轻易将一切受造之物的权柄和\InnerQuote{自有者}的称号归于自己。}\InnerQuote{权柄}与上下文有何关系？并且圣子如何因此就被剥夺了\Quote{自有者}的称号呢？但我实在不知道此时应当做什么更好——是嘲笑其缺乏教育，还是怜悯其显露出的有害愚昧。因为\Quote{祂自己的}这一表述，未按自然含义和善于使用语言者的惯常用法使用，证明了他对代词语法的广博知识（连小男孩都能毫无困难地从老师那里学会），以及他荒谬地偏离主题，谈及既与其论点无关、也与其三段论形式无关的东西，即圣子无分于\Quote{自有者}的称谓——这一论断适合他的怪诞捏造——此类的荒唐似乎结合在一起，旨在引人发笑；以至于可能较粗心的读者会感到某种这种倾向，并为其论点的支离破碎而发笑。但是，天主圣言不存在，或者至少祂不是善（这正是\Person{欧诺弥}所主张的，当他说祂不\Quote{将\InnerQuote{自有者}和\InnerQuote{善}的称号}归给自己），并且要证明一切受造之物的权柄不属于祂——这需要我们流泪，需要哀悼的哭声。\par

因为他并非只是一时冲动、轻率而偶然说出此类话，随后又竭力挽回错误；不，他久久沉溺于恶意之中，在后续陈述中力图超越先前所言。随着行文，他说圣子低于天主本性的距离，等同于天使的本性低于祂本性的距离，他固然没有明说，却试图通过所言造成这种印象。读者可自行判断他话语的含义：其言如下——\Quote{那位被称为\InnerQuote{天使}的，清楚显示了祂宣告谁的话，以及谁是存在者；同时因也被称为天主，显示了祂超越万有。因为祂是万有的天主——万有由祂所造——而祂又是统驭万有的天主的天使。}愤怒涌入我心，打断了我的论述；在这种情绪下，论点在由这类话语激起的愤怒混乱中迷失。或许我为此动情可以原谅。因为谁在想起宗徒宣告一切天使本性都隶属于主，并引用先知的崇高言论作为其教义见证时，能不被这种亵渎激起愤怒呢？——\Quote{当祂引领首生者进入世界时，祂说：天主的众天使都要崇拜祂，}以及\Quote{天主啊，你的宝座永远常存，}以及\Quote{你却一样，你的岁月没有穷尽。}既然宗徒通篇论述，以彰显独生子天主不可接近的威严，当我从基督的敌人那里听到天使之主本身不过是一位天使时——且他并非偶然说出此话，而是竭力维护这一怪异发明，以便确立他的主不优于若翰和梅瑟——我又该作何感想？因为圣经论及若翰说：\Quote{这人就是经上所记载的：\InnerQuote{看，我派遣我的天使在你面前。}}所以若翰是一位天使。但主的敌人即便承认他的主有天主之名，却使祂与梅瑟的天主地位相当，因为梅瑟也是统驭万有的天主的仆人，并被立为埃及人的天主。然而，正如先前所指出的，\Quote{统驭万有}这一短语是圣子与圣父共享的，宗徒明确将这一称号归于祂，当他说：\Quote{按肉身说，基督也是从他们来的，祂是在万有之上，永受赞美的天主。\BibleRef{罗 9:5}}但此人将天使之主贬低到天使的级别，好像他未曾听过天使是\Quote{服役的灵}，是\Quote{火焰}。因为宗徒正是通过使用这些区分性的词语，使不同主体的差异清晰无误，界定从属的本性为\Quote{灵}和\Quote{火}，并用天主之名区分至高权能。然而，虽有如此多人宣扬独生子天主的荣耀，欧诺弥却独举声音，称基督为\Quote{统驭万有的天主的天使}，通过与\Quote{统驭万有的天主}对比，将祂界定为\Quote{万有}之一，并通过赋予祂与天使相同的名字，试图确立祂在本性上丝毫不异于天使；因为他此前常言，一切同名者不可能在本性上有别。那么，这论点是否仍缺少谴责者，因为此人明说\Quote{天使}并非宣告祂自己的话，而是宣告存在者的话？他正是借此试图表明：在起初就有的圣言，那圣言就是天主，祂本身不是圣言，而是另一位圣言之圣言，是后者的仆役和\Quote{天使}。谁不知\Quote{存在者}的唯一对立面就是不存在者？因此，将圣子与存在者对立的人，显然是在扮演犹太人，剥夺基督徒道理中独生子的位格。因为当他说圣子被排除在\Quote{存在者}这一称号之外时，他无疑也在试图确立圣子处于存在范围之外：因为若他承认圣子存在，他就不会在词语发音上争执。\par

但是他试图用圣经的见证来支撑他的荒谬，并搬出\Person{梅瑟}作为他反对真理的辩护者。因为好像那是他得出论据的源头，他随意向我们摆出自己的寓言，说：\Quote{派遣\Person{梅瑟}的就是\OriginalTerm{自有者}{Existent}本身，但祂派遣和发言所藉着的那位，却是自有者的天使，也是一切其余之物的天主。}然而，他的说法并非来自圣经，这可由圣经本身确凿地证明。但如果他说这就是所写之话的含意，我们就必须查考圣经的原文。此外，让我们首先注意到，\Person{欧诺弥}在称主为祂之后万物的天主之后，并没有允许祂与天使的本性相比有任何优越性。因为\Person{梅瑟}在听说他被立为\Person{法郎}的天主时，也并未超越人性的界限，虽然在本质上他与同伴平等，但因权威的优越性而被提升到他们之上；他被称作天主并不妨碍他是人。同样，在此情况下，\Person{欧诺弥}将圣子说成是天使之一，却用天主性的名号来掩饰这种错误，按所表达的方式，以某种歧义的意义允许祂享有天主的称号。让我们再次列出并检查他口出亵渎的原话。\Quote{派遣\Person{梅瑟}的就是\OriginalTerm{自有者}{Existent}本身，但祂派遣和发言所藉着的那位，却是自有者的天使}——他给主所加的这个称号，即\Quote{天使}。那么，我们作者的荒谬被圣经本身驳斥了，其中\Person{梅瑟}恳求主不要把人民的领导权托付给一位天使，而是亲自带领他们的行程。经文是这样说的：天主在说话，\BibleQuote{去吧，下去，引导这百姓往我对你说的地方去；看，在我要巡视的那日，我的天使要在你前面走。}稍后祂又说：\BibleQuote{我要差遣我的天使在你前面。}然后，紧接着之后，祂的仆人向天主祈求如下：\BibleQuote{如果我在你眼前蒙恩宠，求主在我们中间同行，}又说：\BibleQuote{你若亲自不去，就不要从这里领我们上去。}然后天主对\Person{梅瑟}的回答是：\BibleQuote{我也要照你所说的为你做这事，因为你在我眼前蒙了恩宠，并且我认识你超过众人。}因此，如果\Person{梅瑟}请求人民不要被一位天使带领，而那位与他交谈者同意成为他的同行者和军队的向导，那么这就清楚地表明，那位以\Quote{自有者}称号启示自己的就是\OriginalTerm{独生子天主}{Only-begotten God}。\par

若有人对此持有异议，他就表明自己支持犹太人的主张，不将圣子与选民的得救联系起来。因为一方面，那与百姓同行的并非天使；另一方面，按\Person{欧诺弥}的要求，那藉\OriginalTerm{自有者}{Existent}之名显现的不是独生子，这就无异于将会堂的教义转移到天主的教会。因此，他们必须承认两者之一：要么独生子天主从未在\Person{梅瑟}前显现，要么圣子本身就是\Quote{自有者}，从祂发出话语给祂的仆人。但他反驳上述观点，引述圣经本身告知我们，有天使的声音介入，而自有者的话语就是这样传达的。然而，这并非矛盾，而是证实了我们的观点。因为我们也明白地说，先知为向人揭示关于基督的奥迹，称自有者为\Quote{\OriginalTerm{使者}{Angel}}，以免话语的意义被归于圣父——如果整篇讲论中只有\Quote{自有者}的称号，就会如此。但正如我们的话语是心灵活动的启示者和使者（或\Quote{使者}），同样，我们确信那在起初就有的真实圣言，在宣布其圣父的旨意时，被称为\Quote{使者}（或\Quote{传报者}），这称号是因祂传达信息的行动而给予祂的。如同崇高的\Person{若望}，先前称祂为\Quote{圣言}，随后又引入进一步的真理，即圣言就是天主，以免我们的思想立刻转向圣父（如果天主的称号放在前面就会如此），同样，伟大的\Person{梅瑟}先称祂为\Quote{使者}，然后在接下来的话中教导我们，祂正是自有者本身，以使关于基督的奥迹得以预示：圣经藉\Quote{使者}之名向我们保证圣言是圣父旨意的诠释者，藉\Quote{自有者}之名向我们保证圣子与圣父之间存在的亲密关系。即使他再引\Person{依撒意亚}称祂为\BibleQuote{伟大谋士的使者}，也推翻不了我们的论点。因为在那里，预言以明确不容质疑的措辞指示了祂人性的安排；因为他说：\BibleQuote{为我们有一婴孩诞生，为我们有一儿子赐下，权柄必担在祂肩上，祂的名字被称为伟大谋士的使者。} 我想，\Person{达味}正是着眼于这一点，描述祂王国的建立，并非好像祂不是君王，而是指主因安排而屈尊就卑、取了仆人的身份，被提升并吸收进祂王国威严的视角。因为他说：\BibleQuote{我被祂立为君王在祂的圣山熙雍上，宣扬上主的法令。} 因此，那藉自身揭示圣父美善的，被称为\Quote{使者}和\Quote{圣言}，\Quote{印鉴}和\Quote{肖像}，以及所有同义的类似称号。因为正如\Quote{使者}（或\Quote{传报者}）传递某人的信息，同样，圣言揭示内在的思想，印鉴藉自己的印模显示原型，肖像藉自身诠释其所肖之物的美，因此在意义上所有这些术语彼此等同。因此，\Quote{使者}的称号被置于\Quote{自有者}之前，圣子被称为\Quote{使者}乃是作为祂圣父旨意的阐释者，而被称为\Quote{自有者}则是因为没有任何名字能够给予对祂本质的认识，反而超越一切名字的表达能力。因此宗徒的著作也证明祂的名字\Quote{超乎万名之上\BibleRef{斐 2:9}}，并非好像有一个名字优于其他名字、但仍可与它们相比，而是指那真实\Emph{是}者超乎万名之上。\par

% structure: p0011 source=h2 target=subsection reason=relative-heading-rank; base=h2

\subsection{四、在此之后，他恐怕自己回应过长，便视对手的大部分陈述为已反驳而过之。但为那些认为其余部分具有很大份量的人，他简要总结并驳斥\Person{欧诺弥}的亵渎言论——他竟说主也像一切受造物中的动植物一样，在其受造之前并不存在；同样地，青蛙的产生也是如此；啊，何等亵渎！}

但我必须赶快进行，因为我看我的论文早已超出范围，我担心，若我过分延长回答，会被人视为唠叨和多话，尽管我有意省略了对手论文的许多部分，以免我的论证拖成无数的话。因为对于更勤勉的人而言，甚至不够简洁也会招来批评；但对于那些不注重效益、只在意无所事事之人的幻想的人，他们的愿望和祈求是尽可能以很少的步骤走完旅程。那么，当\Person{欧诺弥}的亵渎牵动我们时，我们该做什么？我们要追踪他的每一个转变吗？或者，反复将精力花费在类似的交锋上也许是多余而只是唠叨？因为他们接下来的所有论点都与我们已考察的一致，并没有超出之前的新点。因此，如果我们已经成功彻底推倒他之前的陈述，剩下的也一同倒塌。但为了防止好争论和固执的人认为他们论点的最强部分是我所省略的，为此，可能有必要简要触及剩余部分。\par

他说主在自己的受生以前并不存在——他无法证明主在任何方面与圣父分离。他说这话时，没有引用任何圣经作为其论断的根据，而是凭自己的论证来维持他的命题。但这一特征已被证明是受造物各部分的共同点。没有一只青蛙、一条虫、一只甲虫、一片草叶或其他任何最微小之物，在其形成之前就已存在：因此，他借助其辩证技巧，费尽辛劳和痛苦试图确立的关于圣子的情形，此前已被承认适用于受造物的任何偶然部分，而我们的作者的巨大辛劳，是要证明独生子天主藉属性的分享，与最低级的受造物处于同一水平。因此，他们关于独生子天主的意见与他们对青蛙生成方式的看法相吻合这一事实，足以表明其教义上的乖谬。接着，他力主在受生之前不存在，在事实和意义上等同于非受生。同样的论证也适用于处理这一点——一个人对狗、跳蚤、蛇或任何你喜欢的卑微生物说同样的话也不会错，因为狗在受生之前不存在，在事实和意义上等同于其非受生。但是，如果他们屡次规定的定义中，一切共享属性的也共享本性，而狗和其余各个体在受生之前不存在是它们的属性（这正是\Person{欧诺弥}认为也适用于圣子的），那么读者自会通过逻辑推理看到这一论证的结论。\par

5. \Emph{\Person{欧诺弥}再次称圣子为主、天主以及一切受造物（可理解的和可感的）的创造者，说祂从圣父领受了创造的能力和使命，被托付了创造的工作，如同一个受雇于某人的工匠，而祂的能力是外来的、附加的，是依照星辰的某种组合和位置分配给祂的能力的结果，如同命运在人出生时为他们决定命运一样。因此，他略过了\Person{欧诺弥}所写的大部分内容，驳斥了他的亵渎：即万物的创造者与地和天使一样受生，而独生子的自立体与万物的起源毫无区别；并责备他不尊重天主的奥迹，也不尊重教会的习惯，在寻求虔诚之道的努力中也不追随任何虔诚教义的导师，反而追随\Person{摩尼}、\Person{科鲁托}、\Person{亚略}、\Person{艾提乌斯}及其同类，认为基督教整体是愚昧，}\Emph{而教会的习惯和可敬的圣事是儿戏，在这方面他与外教人毫无区别，后者从我们的教义中借来了关于一位统御万有的伟大天主的观念。同样，这个新的偶像崇拜者也以同样的方式宣讲，尤其宣称圣洗是\Quote{归于一个工匠和创造者}，不怕那些对圣经加以增减者的诅咒。最后，他以指出他是假基督结束了他的著作。}\par

然而，随后他使自己的言论转向较为温和，甚至添上几分柔和；尽管就在不久前他还将圣子从\OriginalTerm{存有者}{Existent}的名号中分割出来，如今却说：\Quote{我们肯定圣子不仅是存有者，且超乎一切存有者之上，更称祂为主、天主，是每一个可感与可知存有的造物主。} 他认为这个\Quote{存有}是什么？是受造的？还是不受造的？因为如果他承认耶稣是主、天主，是全部可知存有的造物主，那么，若他说这存有是\Emph{不受造的}，他必是在说谎，将创造不受造本性的行为归于圣子；但若他相信它是\Emph{受造的}，他便使圣子成了自己的造物主。因为，若创造的行为不为了那独立而不受造者而从可知的本性中分离出来，便不再有任何区别的标志，因为可感的受造物和可知的存有将被视为隶属于同一个类别之下。但在这一点上，他提出了一个主张：\Quote{在存有者的创造中，祂已由圣父受托，负责营造一切有形与无形之物，并眷顾一切受造物，因为从高处所分予祂的权能足以产生那些已被营造之物。} 由于我们的论文已无比冗长，只得对这些主张简略带过；但可以说，亵渎之词围绕其论点，如同毒蜂般包藏大量念头。他说：\Quote{祂已受托，由圣父负责营造万物。} 但如果他是在谈论一个按其雇主旨意执行工作的工匠，他岂不会采用同样的说法？因为我们对\Person{贝匝肋耳}也这样描述，他受\Person{梅瑟}委托建造会幕，遂成了那里所提及之物的营造者，并且若非他预先因天主默感获得知识，并在\Person{梅瑟}委托他执行任务时敢于尝试，他原本不会着手进行。因此，\Quote{受托}一词暗示：祂在创造中的职务和能力是偶然得来的，意即：在祂受托负责之前，祂既无意愿亦无能力行动，但当祂获得执行工程的权威和足够的能力时，\Emph{那时}祂才成了万有的工匠，而从高处所分予祂的权能，一如\Person{欧诺米}所说，足以达成目的。那么，他是否以某种占星术的戏法，将圣子的受生置于某种命运之下——正如那些施行这种虚空骗术的人所肯定的那样，人一生的命运是在其出生之时，由星宿的某种会合或对冲决定的，上天的旋转以某种有序的进程运行，将特殊的能力分配给出世的人？或许我们这位博学之士心中想的正是这类东西；他说，那超乎一切统权、一切宰制者，超乎一切主权，超乎一切可称的名号——不但在这世界，也在那将来的世界——祂，竟如同被禁锢在某个虚空中一样，被分配了从高处来的权能，按照受造物的数量来衡量。我也将简略地越过他论述的这一部分，仅从轻微的研究开端，为较有识力的读者撒下一些种子，使它们能辨识其亵渎。此外，在接下来的部分中，好像已经写好了为我们自己辩护之词。因为当他自己亲口承认其教义之荒谬时，我们再也不能被认为误解了他论述的意图，或曲解其词句以便加以批评了。他的话说如下：\Quote{怎么？大地和天使岂不是在以前不存在的时候受造的吗？} 看，我们这位崇高的神学家多么不耻于将同一描述用于大地、天使和万有的造物主！诚然，如果他觉得适合将同一谓语用于大地与大地的主，那么他要么必须将大地奉为神，要么将造物主降格到与大地的同一水平。\par

随后，他又添加了一些言辞，使他的亵渎更加彻底地剥去了一切伪装，其荒谬之处甚至连孩童也看得出来。因为他说道：\Quote{详细列举所有理性受造物的生成模式，或那些并不都具有\InnerQuote{存在者}之本性、而是依照建造者之运作而显出差异的本质，是一项冗长的工作。}无需我们多言，此处所含的对圣子的亵渎已是昭然若揭，因为他承认，凡可归于每一种生成模式和本质的，与对独生者天主性本体的描述毫无二致。但我认为最好越过其中那些他试图坚持其亵渎的中间段落，而直奔我们必须对其教义提出的控告之首。因为他将发现，他把重生的圣事表现为无用之举，把奥秘的祭献表现为无益之事，而参与这些圣事对领受者毫无裨益。因为在那些煞费苦心的世代中，他为了贬低我们的教义，举出\Person{瓦伦提诺}、\Person{克林托}、\Person{巴西里德}、\Person{蒙丹}和\Person{马西昂}作为其拥护者，并在断言那些主张天主本性不可知、其生成方式不可知的人完全无权也不配称为基督徒之后，在把我们也算作他如此贬低的人之列后，他接着以这些话来展开自己的观点：\Quote{但我们与圣善而蒙福的人一致认为，虔敬的奥迹不在于可敬的名号，也不在于习律和圣事标记的独特特征，而在于教义的精确性。}他写这些话时，并非出于福音作者、宗徒或任何旧约作者的引导，凡稍识神圣天主圣经的人，对此都一目了然。我们自然会被引导认为，他所谓的\Quote{圣善而蒙福的人}指的是\Person{摩尼}、\Person{尼阁}、\Person{科路图}、\Person{艾迪乌}、\Person{亚略}以及同一伙的其余人，他在确立以下原则上与他们完全一致：既非宣认神圣名号，亦非教会习律，也非教会的圣事标记，能成为虔敬的确认。但我们从基督的圣言得知：\BibleQuote{人除非由水和圣神重生，不能进入天主的国}以及\BibleQuote{谁吃我的肉，喝我的血，必得永生}，我们因而确信：虔敬的奥迹因宣认神圣名号——圣父、圣子、圣神的名号——而得以确认，我们的救恩因参与圣事习律与标记而得到坚固。然而，教义常被那些与此奥迹毫无关联的人仔细探究，而我们可以听到许多这样的人将我们所持的信仰当作辩论竞争的主题，他们中的一些人甚至常常能触及真理，但尽管如此，他们仍然与信仰疏远。既然他藐视那些可敬的名号（藉此，那更神圣的诞生之能力将恩宠分赐给那些怀着信德前来领受的人），又轻视圣事习律与标记的共融（基督徒的宣信从中汲取活力），让我们稍微变通一下，对听从他谎言的众人说出先知的话：\BibleQuote{你们心迟愚钝要到几时？你们为何喜爱毁灭，追求虚谎呢？}你们为何看不见那迫害信仰的人，正在邀请那些同意他的人去违反他们的基督徒宣信呢？因为如果宣认天主圣三可敬而宝贵的名号是无用的，教会的习律是无益的，而这些习律中包括十字圣号、祈祷、圣洗、告解罪过、恪守诫命的迫切热忱、品格的正直、生活的节制、对正义的尊重、努力不为情欲所激动、不为享乐所奴役、或在道德上不缺失——如果他说这样的一切习惯都没有任何益处，而圣事标记并不如我们所相信的那样确保属灵祝福，并保护信徒免受恶者诡计的攻击，那么他还做了什么，不是在公开向人们宣告：他认为基督徒所珍视的奥迹是个寓言，嘲笑神圣名号的威严，视教会的习律为儿戏，将所有圣事运作当作空谈和愚蠢吗？除此之外，那些仍依附异教的人，在我们的信条面前还能提出什么贬低之辞？他们不是也将那些藉以确认信仰的神圣名号的威严当作笑柄吗？他们不是嘲笑那些由入门者所遵守的圣事标记和习律吗？认为虔诚只在于教义，这对谁来说比异教徒更是一种显著特征呢？因为他们也说，按照他们的看法，有比我们所宣讲的福音更有说服力的东西，他们中的一些人认为有一位超越其余诸神的伟大的神，并承认一些从属的权能，它们之间以某种常规的秩序和序列，在优越或低劣方面彼此不同，但都同样隶属于至高者。因此，这正是新偶像崇拜的导师们所宣讲的，而跟随他们的人并不畏惧那存留于违犯者身上的定罪，仿佛他们不明白：实际做一件不当之事，远比仅仅在言语上犯错更为严重。那么，那些在行为上否认信仰、轻视神圣名号的宣认、认为圣事标记所实现的圣化毫无价值、并被说服去重视那些巧设的虚构故事、幻想他们的救恩在于对受生与非受生的诡辩中的人，他们除了是救恩教义的违犯者之外，还能是什么呢？\par

但如果有人认为，我们提出这些指控是出于狭隘和不公允，那就让他自行审视我们这位作者的作品，既包括我们先前所引，也包括从我们引文逻辑推断出的结论。因为他直接违犯了主的法律——赐予我们入门方式即构成了法律——他声称：圣洗不是归于圣父、圣子及圣神，如同基督在交付奥迹时命令祂的门徒那样，而是归于一位工匠和创造者；他又说：\InnerQuote{不仅是独生子的父，也是祂的天主。}让那给邻舍喝浑浊祸水的人有祸了！他抛出泥浆，是如何搅浑和玷污了真理！面对那诅咒——凡在神圣言说上加添或妄自删减的人所承受的——他为何毫不畏惧？让我们诵读主亲口所说的宣言——祂说：\BibleQuote{你们要去，使万民成为门徒，因父、及子、及圣神之名给他们授洗。}祂在何处称圣子为受造物？圣言在何处教导说，圣父是独生子的创造者和工匠？在所引的话语中，何处教导说圣子是天主的仆人？在交付奥迹时，何处宣告了圣子的天主？你们这些被诡计拖向丧亡的人，难道不察觉、不明白，你们所托付灵魂的导师是怎样的人——他窜改圣经，歪曲神圣言说，用自己的泥浆玷污虔敬教义的纯洁，不仅用舌头攻击我们，还企图篡改真理的神圣声音，热切地要将自己的乖谬凌驾于主的教导之上？难道你们不察觉，他是在反对那令万有屈膝的名，以致终有一日主的名不再被听闻，而欧诺米将被带进教会取代基督？难道你们尚未意识到，这无神论的宣讲正是魔鬼为假基督的来临所设的预演、预备和序幕？因为那雄心勃勃要证明自己的话比基督的话更权威，并要转变信德，使之脱离神圣的名号以及圣事性的风俗和标记，而归于自己的欺骗的人——除了被称为假基督，还能恰当地称为什么？\par

\SourceRecord{Translated by H.C. Ogle and H.A. Wilson.；Nicene and Post-Nicene Fathers, Second Series；Vol. 5.；Edited by Philip Schaff and Henry Wace.；Buffalo, NY: Christian Literature Publishing Co.,；1893.；<http://www.newadvent.org/fathers/290111.htm>.}

% structure: p0001 source=h1 target=section reason=page-title

\section{\WorkTitle{驳斥欧诺弥（卷十二）}}

% structure: p0002 source=h2 target=subsection reason=relative-heading-rank; base=h2

\subsection{一、本卷对主向玛利亚所说的话作出了卓越的解释：\Quote{不要摸我，因为我还没有升到我的父那里去。}}

但是让我们看看接着这渎神言论的下一段话是什么，这段话实际上正是他们辩护其教义的关键。因为那些想要把独生子荣耀的威严贬低为奴仆般的、卑劣观念的人，认为他们在主复活后升天前向玛利亚所说的话中找到了他们论断的最强有力的证据，祂说：\Quote{不要摸我，因为我还没有升到我的父那里去：但到我的兄弟那里去，对他们说，我升到我的父和你们的父那里去，到我的天主和你们的天主那里去。}我认为，对这些话的正统解释，即我们惯常相信它们是对玛利亚所说的意义，对所有真正接受信德的人都是显而易见的。然而，对这一点我们将适时予以讨论；但与此同时，值得向那些对我们提出诸如\Quote{上升}、\Quote{被看见}、\Quote{被触摸所识别}，以及此外\Quote{因兄弟关系与人类联合}等说法的人询问，他们认为这些说法是属于神性还是属于人性。因为，如果他们看见在天主性内有可见可触、由食物和饮料支持、与人类有亲属和兄弟关系、以及所有肉体本性的属性，那么让他们也将这些以及任何他们愿意的事物，如动力和位置变化，这些是受身体限制的事物所特有的，归之于独生天主。但是，如果祂借着玛利亚与祂的兄弟交谈，而独生子没有兄弟（因为如果祂有兄弟，独生子的属性如何得以保持？），并且同一位说过\Quote{天主是神}的对祂的门徒说\Quote{摸我 \BibleRef{路 24:39}，}为要显示人性是可触摸的，而神性是不可触摸的；并且那位说\Quote{我去，}的指示位置变化，而那包含万有的，如宗徒所说，\Quote{在祂内，} \Quote{万有受造，在祂内万有也持存 \BibleRef{哥 1:16}，}在受造物中没有任何外在的东西可供祂通过任何运动移动到那里（因为运动只能通过被移动之物离开其所在之处并占据另一处来达成，而那遍及万有、在万有内、管理万有、不受任何受造物限制的，没有可移往的地方，既然没有一处缺乏天主的圆满），这些人怎能放弃相信这些说法出自那外表之物，而将它们应用于那神性、超越一切理解的本性呢？宗徒在雅典人的演讲中明确禁止我们想象任何这样的天主之事，因为天主的德能不通过触摸被发现，而是通过理性的默观和信德。再者，那位在门徒眼前吃过食物，又应许在门徒之前往加里肋亚去并在那里被他们看见的——祂所揭示的是谁？是那位无人见过或能见的天主吗？还是身体的形象，即天主在其内的仆人的形式？如果以上所述明确证明所引述的话语的含义是指那可见的、具有形状的、能运动的、与祂门徒的本性相近的，而这些属性没有一样能在那不可见、无形体、不可触摸、无形状的祂身上被察觉，那么他们如何竟然贬低那从起初就在并在圣父内的独生天主，与伯多禄、安德肋、若望及其余宗徒同等，称他们为独生子的兄弟和同仆？然而，他们所有的努力都指向这个目标，即表明在本性的尊威上，圣父与独生子的尊位、能力、本质之间的距离，与独生子和人性之间的距离一样大。他们强引这句话来支持这个含义，将天主和父的名号视为对主和祂的门徒具有共同的意义，他们认为当祂以完全相同的方式被承认为祂（主）和他们的天主和父时，本性尊位上的差异是不可想象的。\par

而他们以逻辑维系其亵渎的方式如下：要么通过所用的相对词，门徒与圣父之间也表达了本质的共融，要么我们不可藉此短语使主也共融于圣父的本性；并且，正如万有之上的天主被称为他们的天主，意味着门徒是祂的仆人，同理，由所论的话语可承认，圣子也是天主的仆人。如今，向玛利亚所说的话语并不适用于独生子的天主性，这一点可从说话的目的得知。因为那屈尊就卑、与人性渺小同列者，正是说这些话的那位。当时祂所说的话有何深意，那些藉圣神探究神圣奥迹深处的人，可完全知晓。但就我们所能及的范围，我们将追随教父的引导，简略陈述。那按本性是万有之父者，万物从祂而生，由宗徒崇高的言论被宣认为独一无二。他说：\Quote{因为只有一个天主，就是父，万物都出于祂 \BibleRef{格前 8:6}。} 因此，人性并非从别的根源进入受造界，也不是在人类始祖中自发产生，而是同样有其自身构成的作者，即万有之父。而天主性这一名称本身，无论它表示监督之权或预见之权，都暗示着与人类的某种关系。因为那赐予万有存在之力的，正是祂所产生之物的天主和监督者。但由于那在我们心中撒下悖逆莠子的仇敌之诡计，我们的本性不再保持圣父肖像的印记，反而转变为罪恶的丑陋形象，为此，藉着意志的相似，它被接枝到罪恶之父的邪恶家族中；以致良善真实的天主和圣父，不再是被自身邪恶放逐之人的天主和父亲，反而，取代那按本性是天主者，那些如宗徒所说\Quote{本来不是神的 \BibleRef{迦 4:8}} 受到了尊崇；并且取代圣父，那假称为父者被视为父亲，正如耶肋米亚先知在其隐晦话语中所说：\Quote{山鹧鸪叫了，聚集了它未曾孵化的。} 既然此乃我们灾祸的总结——人性被逐出良善的圣父，并被剥夺了天主的监督和眷顾——为此，那整个理性受造物的牧者，放弃了天上高处的无罪的超世俗羊群，因爱而动，去寻找那迷失的羊，即我们的人性。因为人性，按照比喻中的相似，唯独因恶行从理性存有的一百中走失，若与整体相比，不过是微不足道的一小部分。既然我们那疏离天主的生命无法自行返回高天之处，为此，如宗徒所说，那不知罪的，替我们成了罪，并藉承担我们的诅咒而救我们脱离诅咒，又拿起并，用宗徒的话，\Quote{杀死了} 在自己身上 \Quote{仇恨}——即那藉罪恶横亘在我们与天主之间的（事实上，罪\Emph{就是}\Quote{仇恨}）——并成为我们之所是，藉着祂自己再次将人性与天主合一。因为纯洁使那新造的人（这人是按天主造的，在他内住满了整个天主性，有形有体地）与我们本性的圣父紧密相连，祂便将所有分享祂身体并与之相似的本性，一同引入同一恩宠。这喜讯祂通过那女人宣布，不仅给那些门徒，也直到今日成为圣言之门徒的所有人——这喜讯便是：人不再被放逐，不再被逐出天主的国，而是重新成为儿子，重新处于天主为他设定的位置，因为初果的神圣也祝圣了全团。祂说：\Quote{看，我和天主所赐给我的孩子。} 那为了我们而成为血肉的分享者，祂恢复了你们，并将你们带回你们因罪恶而成为血肉时迷失之处。因此，我们从前因反叛而疏远的那位，现今成了我们的父和我们的天主。于是在上述引文中，主宣告了这恩惠的喜讯。这些话并非圣子卑微的证据，而是我们与天主和好的喜讯。因为在基督人性中所发生的事，是对全人类共同的恩赐。正如我们见祂内那本性归于地上的身体重量，上升穿过空气进入诸天，我们便相信，按宗徒的话，我们也\Quote{将被提到云中，在天空与主相遇 \BibleRef{得前 4:16}}；同样，当我们听到真实的天主和圣父成了我们初果的天主和父亲时，我们不再怀疑同一天主也成了我们的天主和父亲，因为我们得知，我们必到达同一地方，即基督作为前驱为我们所进入之处。并且，这恩宠通过一个女人启示出来，这事实本身也与我们的解释相符。因为如宗徒告诉我们：\Quote{女人被诱惑，陷于悖逆 \BibleRef{弟前 2:14}}，她藉不顺从在反叛天主中为首，为此她成了复活的首位见证人，好使她藉对复活的信仰，补偿因不顺从而导致的倾覆；并且，正如她起初为丈夫传达并倡导蛇的阴谋，将邪恶的开端及其后果引入人类生命，同样，她为门徒传递那杀死反叛毒蛇者的话，而成为人类信德的向导，藉此，死亡的首个宣告被正当地废止。确实，更勤奋的学者可能从经文发现更有益的阐释。但即使找不到，我想每位虔诚的读者都会同意，将我们对手所提出的与我们提出的相比较，前者是虚浮的。因为前者是专为毁灭独生子的光荣而捏造的，别无他意；而后者则涵盖了关于人的救恩计划的旨意。因为已表明，并非那不可触摸、不可更改、不可见的天主，而是那属于人性的可动、可见、可触的本性，命令玛利亚将这话传达给祂的门徒。\par

% structure: p0005 source=h2 target=subsection reason=relative-heading-rank; base=h2

\subsection{二、接着，他论及\Person{欧诺弥}的亵渎——这已被伟大的\Person{巴西略}驳斥过，\Person{欧诺弥}曾将独生子天主放逐到黑暗领域——以及\Person{欧诺弥}为其亵渎所作的辩护或解释；他指出，他现在的亵渎因他的辩护而变得比之前的更糟；在此他极有见地地论述了\Quote{真}光和\Quote{不可接近的}光。}

我们也应探讨这一点——在大巴西略曾指摘他错谬的那些事上，他有什么辩护可提，当他将独生天主放逐到黑暗领域，说：\Quote{就如受生者与非生者之间的差别如此之大，光与光之间的分歧也是如此之大。}因为他已表明，受生者与非生者之间的差别不仅是强度上的大小之别，而且在意义上完全对立；并且他由其前提合乎逻辑地推论出，既然父之光与子之光的差别对应于非生与受生，那么我们在子中必须设想的不只是光的减弱，而是与光的完全疏离。正如我们不能说受生是一种被修改的非生，而是术语\OriginalTerm{生成}{\GreekText{γέννησις}}与\OriginalTerm{非生}{\GreekText{ἀγεννησία}}的意义绝对矛盾且相互排斥，同样，如果要在父之光与被认为存在于子中的光之间保持同样的区分，那么逻辑上就会得出结论：子从此不再被设想为光，因为他既被排除于非生本身，也被排除于伴随该状态的光——而不同于光的那一位，显然将与其对立者相关联——我说，这种荒谬性源于他的原则，尤诺米乌斯试图以辩证的伎俩为之辩解，他这样说：\Quote{因为我们知道，我们知道真光，我们知道那在天地之后造光的那一位，我们亲耳听过生命与真理本身，即基督，向祂的门徒说：\InnerQuote{你们是世界的光}，我们从蒙福的保禄那里学到，当他称万有的天主为\InnerQuote{不可接近的光}，并以附加语界定并教导我们祂之光的超越优越性；如今我们既已得知一光与另一光之间有如此大的差别，我们将不再耐心忍受甚至连提说那两光的概念是同一个的想法。}他提出这样的论点来试图反对真理，难道是认真的吗？还是在试验那些追随他错误之人的愚蠢，看看他们能否识破如此幼稚而明显的谬误，或者是否毫无辨别力，看不出这种赤裸裸的欺骗？因为我想没有人会愚蠢到看不出尤诺米乌斯用同音异义词的戏法欺骗他自己和他的崇拜者。他说，门徒被称为光，在创造过程中产生的也叫光。但谁不知道这些只是名称相同，而所指之物在各处却完全不同？因为太阳的光使视觉辨识，而门徒的言语则在人灵中植入真理的光照。那么，如果他在那光中也认识到这种差异，以至于他认为身体的光是一回事，灵魂的光是另一回事，我们就不再与他讨论这一点，因为即使我们保持沉默，他的辩护本身也定他的罪。但如果在那光中他无法发现运作方式上的这种差异（因为他可能会说，并非眼睛的光照亮肉体，而属神的光照亮灵魂，而是这一光与另一光的运作和德能相同，在同一领域、针对同一对象运作），那么，他如何从太阳光束之光与宗徒言语之光的差异，推断出独生之光与父之光之间有类似的差异呢？\Quote{但子，}他说，\Quote{被称为\Quote{真}光，父是\Quote{不可接近的光}。}好，这些附加的区分只引进了程度上的差异，而非种类上的差异，在子之光与父之光之间。他认为\Quote{真}是一回事，\Quote{不可接近}是另一回事。我想没有人会愚蠢到看不出这两个术语在意义上的真实同一性。因为\Quote{真}和\Quote{不可接近}各自与其对立面完全同等程度地远离。因为正如\Quote{真}不容许掺杂任何假，同样\Quote{不可接近}也不容许其对立面接近。因为\Quote{不可接近}肯定不为邪恶所接近。但子之光并非邪恶；因为人怎能在邪恶中见到那真实者呢？那么，既然真理不是邪恶，没有人可以说父中的光是真理所不能接近的。因为如果它拒绝真理，它当然就会与虚假相关联。因为矛盾物的本质是这样的：较优者的缺失意味着其对立面的存在。那么，如果有人要说父之光被看作与其对立面的呈现距离遥远，他就是在与宗徒的意图相符的方式上解释\Quote{不可接近}一词。但如果他说\Quote{不可接近}意味着与善疏离，那么他无非是设想天主与自身疏离并敌对，既是善又与善对立。但这是不可能的：因为善与善是亲近的。因此，一光与另一光并不分歧。因为子是真光，父是不可接近的光。事实上我敢于说，如果有人将这两个属性互换，他也不会错。因为真是假所不能接近的，另一方面，不可接近者被发现于无玷的真理中。因此，不可接近者与真者相同，因为每个表达所意指之物同样不为邪恶所接近。那么，在与\Quote{光}一词的同音异义中欺骗自己和追随者的人，所想象的存在于这些之中的区别又是什么呢？但我们也不要略过这一点而不加注意：他是按照自己的幻想篡改了宗徒的话之后，才引用那句短语，仿佛它是出自宗徒之手。因为保禄说：\Quote{\Emph{住在}不可接近的光中\BibleRef{弟前 6:16}}。但自己是某物与在某物之中有很大区别。因为那说\Quote{住在不可接近的光中}的人，并未用\Quote{住}一词指天主本身，而是指围绕祂的事物，这在我们看来相当于福音中告诉我们父在子内的短语。因为子是真光，真理是不为虚假所接近的；所以，子就是不可接近的光，父住在其中，或者说父在祂内。\par

% structure: p0007 source=h2 target=subsection reason=relative-heading-rank; base=h2

\subsection{三、他更进一步显著地诠释了福音的话：\BibleQuote{在起初已有圣言}、以及\BibleQuote{生命}和\BibleQuote{光}，以及\BibleQuote{圣言成了血肉}，这些曾被欧诺弥误解的经文；并推翻了他的亵渎，表明主的工程是出于慈爱，而非无能，且是在圣父的合作下完成的。}

但他将力量投入空泛的争辩中说：\Quote{我从事实本身和人所相信的神谕中，为我的陈述提出证据。}这是他的承诺，但他提出的论证是否证实了他的宣称，有识别力的读者自然会考虑。他说：\Quote{蒙福的若望在说了\BibleQuote{在起初已有圣言}，并称祂为\BibleQuote{生命}，后来又给生命加上\BibleQuote{光}的称号之后，稍后说道：\BibleQuote{圣言成了血肉。}}那么，如果光是生命，圣言是生命，圣言成了血肉，由此便清楚可见：光成了血肉。那么，如何呢？因为光、生命、天主和圣言显现在肉身中，难道真的光与在父内的光有任何程度上的不同吗？不，福音证实，即使光存在于黑暗中，它仍然不为相反的元素所接近：因为他说：\BibleQuote{光照在黑暗中，黑暗却没有胜过它。}那么，如果光存在于黑暗中时被变成了相反的东西，并被黑暗压倒，那将有力地支持那些想证明这光与在父内所默观的光相比是何等低劣的人的观点。但是，如果圣言即使在肉身中仍然是圣言，如果光即使在黑暗中照耀也仍然是光，没有接纳相反者的共融，如果生命即使在死亡中仍保持自身的安全，如果天主即使屈尊取了奴仆的形体，却不成为奴仆，而是除去奴仆的从属地位，将其吸收到主宰和君权中，使那属人和卑微的成为主和基督——如果这一切都是如此，那么他又如何通过这个论证显示光变成了低劣呢？既然每一个光都同样具有不可变为恶、不可改变的特性？又，他为什么没有注意到这一点：那看见了成为血肉的圣言——祂既是光又是生命又是天主——的人，透过他所看见的光荣，认出了光荣之父，并说：\BibleQuote{我们看见了他的光荣，如同父独生者的光荣}？\par

然而，他得出了一个不可反驳的论点，我们早已察觉到它潜藏在他后续陈述的脉络中，但如今却毫无掩饰地公开宣扬了。因为他希望表明，圣子的本质是受难的、可朽坏的，与那流动不息的物质本性毫无区别，从而借此证明祂与圣父的不同。因为他说：\Quote{如果他能证明那统驭万有的天主，那不可接近的光，曾降生成人或能降生成人，服从权柄，听从命令，接受人的法律，背负十字架，就让他说光与光是平等的。}倘若这些话是我们从欧诺弥的前提出发，当作必然结论而提出的，谁会不指责我们不公，用过于精妙的辩证术将对方的陈述归结到如此荒谬的地步？但实际情况是，他们自己毫不掩饰地承认其假设自然导致的荒谬，这反而支持了我们的主张，即我们并非未经深思熟虑，而是在真理的帮助下谴责异端论的论点。因为看哪，他们对抗独生天主的斗争是多么不加掩饰、直言不讳！不，祂的敌人竟将祂的仁慈之工当作贬低和诋毁天主圣子本性的手段，仿佛祂不是出于有意的目的，而是由于本性的强迫才降生为人，并忍受十字架的苦难！正如石头本性的下落、火焰本性的上升，这些物质对象彼此不交换本性（因此石头不会向上，火焰也不会因重量而下坠），同样，他们断言受难是圣子本性的一部分，因此祂来到了与自己同类且熟悉的事物中；但圣父的本性既不受这些苦难影响，就仍然不为邪恶的接触所接近。因为他说：那统驭万有的天主，那不可接近的光，既没有降生成人，也不能降生成人。第一个说法已经足够，即圣父没有降生成人。但现在由于他的附加，产生了双重的荒谬：要么他指控圣子有邪恶，要么指控圣父无能。因为如果分享我们的肉体是邪恶的，那么他就将邪恶归于独生天主；但如果对人的慈爱是良善的，那么他就断言圣父无力行善，因为他说圣父不可能通过取得肉体来有效施予这种恩宠。然而，世界上有谁不知道生命之能在圣父和圣子中都能实际运行呢？\Quote{因为父怎样叫死人复活，使他们活着，}他说，\Quote{子也照样随自己的意愿使人活着。}——显然\Quote{死人}指的是我们这些从真生命堕落的人。如果圣子如同圣父使人活着一样，且别无二致地实现同一恩宠，那么天主的仇敌怎能同时亵渎地攻击两者，既因无能行善而侮辱圣父，又因与邪恶关联而侮辱圣子呢？但他又说：\Quote{光不与光平等，}因为他称一个为\Quote{真光，}另一个为\Quote{不可接近的光。}那么真光是否被视为不可接近之光的减损呢？为何如此？然而，他们的论点是圣父的天主性必须被认为比圣子的更伟大、更崇高，因为前者在福音中被称为\Quote{真天主，}后者被称为\Quote{天主}，没有加上\Quote{真}。那么，同一个术语，当应用于天主性时表示概念的提升，而应用于光时却表示减损，这怎么说得通呢？因为如果他们声称圣父比圣子更伟大，因为祂是真天主，那么按同样的推理，圣子就会被承认为比圣父更伟大，因为前者被称为\Quote{真光，}而后者则不是。\Quote{但这光，}欧诺弥说，\Quote{执行了仁慈的计划，而另一位在那恩慈的行动上却无所作为。}一种新奇而陌生的决定尊贵优先性的方式！他们认为对仁慈目的无效者优于有效者。但这样的观念在基督徒中既不存在，也永远不会存在——这种观念认为一切存在之善并非源自圣父。但对我们人类而言，所有善行中最大的福分被所有正直的人视为回归生命；这是通过主在人性中的经世来实现的；并非圣父袖手旁观，如同异端所认为的那样，在这经世时期无所作为、不行动。因为那位说\Quote{派遣我者与我同在}和\Quote{住在我内的父，祂作这些事}的，并不是指示这一点。那么，异端有什么权利将祂对我们的恩慈干预只归于圣子，从而排除圣父在我们对成功结局的感恩中享有任何部分或份子呢？因为自然的，感谢的回报只应归于我们的恩人，而那位不能施恩于我们的，就在我们的感恩之外。你看见了吗？他们对独生圣子的亵渎攻击偏离了目标，并出于自然结果转而指向圣父的尊威。在我看来，这是他们进路的一个必然结果。因为按神圣的宣言，\Quote{尊敬子的，就是尊敬父；}显然，另一方面，对子的攻击也打击了父。但我认为，对于那些以纯朴之心接受十字架和复活宣讲的人，同样的恩宠应该对圣子和圣父同样值得感谢；既然圣子完成了圣父的旨意（按宗徒的说法，就是\Quote{愿意所有人得救\BibleRef{弟前 2:4}}），他们就应该为这恩赐同样尊敬圣父和圣子，因为我们的救恩若不是圣父的善意通过祂自己的权能为我们付诸实施，就不会成就。我们从圣经中得知，\Quote{基督是天主的能力\BibleRef{格前 1:24}}。\par

% structure: p0010 source=h2 target=subsection reason=relative-heading-rank; base=h2

\subsection{四、然后他又指控\Person{欧诺弥}从象形文字著作、埃及神话和偶像崇拜中学到了他的术语\OriginalTerm{非受生性}{\GreekText{ἀ} \GreekText{γεννησία}}，并将\Person{阿努比斯}、\Person{奥西里斯}和\Person{伊西斯}引入基督徒的信经，并表明，若认为独生者是由于必然性而非自愿地承受苦难，则他不配从人获得任何感恩：火不因其温暖而受谢，水不因其流动性而受谢，因为它们的结果并非出自自决之权能，而是出于本性的必然。}

让我们再次注意所引用的段落。\Quote{若他能证明，}他说，\Quote{那在万有之上的天主，那不可接近的光，曾降生成人或能降生成人……那么就让他说，光与光相等。}他的话的意思从句子本身即可明了，即，他认为圣子并非凭其全能的属神性而能承受如此慈爱的形式，而是因着一种可受难的本性，祂才屈尊接受十字架的苦难。于是，当我沉思并探究\Person{欧诺弥}何以陷入如此关于天主的概念，以至于认为一方面，非受生之光因其相反者而不可接近，完全无玷、纯净，不受任何情感与感觉的影响；而另一方面，受生之光在本性上则是居间的，以至于它未能保持属神性在无受感性上的纯洁与无瑕，反而拥有一种由相反者混合与复合的本质，它既伸展以分有善，同时又消融于一种可受感的状态，既然不可能从圣经中获得支持如此荒谬理论的前提，我便产生了一个念头：他是否可能是埃及人关于属神事物之思辨的仰慕者，并将他们的幻想与他对独生者的观点混为一谈？因为据传他们说，当他们将某些无理性动物的形状与人的肢体相结合时，他们那幻想式的偶像组合方式，是那种混合本性的奥秘象征，他们称之为\Quote{神灵}，它比人的本性更微妙，在能力上远胜于我们的本性，但其中的属神因素并非不混合或不单纯，而是与灵魂的本性和身体的感知相结合，并接纳快乐与痛苦，而这两者都不能在\Quote{非受生之神}那里存在。因为他们也使用这个名称，将他们所想象的最高之神归于非受生的属性。因此，在我们看来，这位博学的神学家似乎是从埃及的神庙中引进了一个\Person{阿努比斯}、\Person{伊西斯}或\Person{奥西里斯}到基督徒的信经中，只差承认他们的名称了：但公开承认偶像名称者，与内心持有关于他们的信念却谨慎避讳其名称者，在亵渎上并无区别。那么，如果无法从圣经中为这种不敬获得任何支持，而他们的理论完全从象形文字的谜语中汲取力量，那么毫无疑问，正派之人对此应作何感想。但我们的这项指控并非侮辱性的诽谤，\Person{欧诺弥}将用自己的话为我们作证，他说非受生之光不可接近，且无力屈尊体验情感，但断言这种状态与受生者密切相关并相似：因此，人对独生天主为祂所承受的无需感恩，如果照他们所说，祂是凭其本性的自发性滑入对情感的体验，其本质天然地可被如此影响，并自然地被拖向那里，这不需要感谢。因为谁会视那必然发生的事为恩惠，即使它是有益和有利的呢？我们既不感谢火因其温暖，也不感谢水因其流动性，因为我们将这些品质归因于它们各自本性的必然，因为火不能失去其温暖之力，水也不能在斜坡上静止，因为倾斜自发地带它的运动。那么，如果他们声称圣子通过降生成人所完成的恩惠是出于本性的必然，他们当然不会感谢祂，因为他们将祂所做的归因于自然强迫，而非权威的权能。但是，如果他们一方面体验到恩赐的益处，另一方面却贬低带来恩赐的慈爱，我怕他们的不敬会转向相反的谬误，并认为那能如此受感的圣子的状态，比圣父远离此类情感的状态更值得尊敬，他们以自己的利益作为善的标准。因为如果圣子不能如此受感，如他们所说圣父的那样，那么我们的本性将仍然处于悲惨的境地，因为没有一位通过亲身经历来提升人的本性至不朽——因此，这些诡辩者以企图毁灭独生天主之尊威的同一手段，却不知不觉中将人对祂的概念提升到更宏大、更高超的高度，因为那有能力行动者，比那无力行善者更应受尊敬。\par

% structure: p0012 source=h2 target=subsection reason=relative-heading-rank; base=h2

\subsection{五、然后，再次讨论圣父和圣子的真光与不可接近的光、特殊属性、共性与本质，并表明\Quote{受生}与\Quote{非受生}的关系在意义上不包含对立，却呈现一种不容中项的对立与矛盾，他以此结束本书。}

但我感到我的论证正在失控，因为它没有保持在常规轨道上，反而像一匹热血沸腾、意气风发的小马，被对手的亵渎言辞裹挟着，在他们的体系之荒谬处驰骋。因此，当我们展示荒谬结论时，必须加以约束，以免它逾越节度的界限。然而，善意的读者无疑会原谅我们所说的一切，不会将我们探讨中出现的荒谬归咎于我们，而是归咎于那些奠定如此有害前提的人。不过，我们现在必须将注意力转向他的另一个论述。因为他声称我们的天主也是\OriginalTerm{复合的}{composite}，理由是我们虽然假定\OriginalTerm{光}{Light}是共通的，却又通过某些特殊属性与各种差异将一光与另一光分开。因为一个东西，即使被一个共同本性所统一，却被某些差异和特性的结合所分开，它仍然不亚于复合。对此，我们的回答简短且容易打发。因为他作为指控我们学说的东西，如果我们发现他并没有用自己的话确立同样的立场，我们自己就承认。让我们看看他写了什么。他称主为\Quote{真}光，又称圣父为\Quote{不可接近的}光。因此，通过这样称呼二者，他也承认它们在光方面的共通性。但由于名称适用于事物是因为它们适合，正如他经常强调的，我们不认为\Quote{光}这个名称是毫无意义地用于神圣本性，而是由于某种实在基础而被述说。因此，通过使用一个共同的名称，他们承认所指对象的同一性，因为他们已经断言，具有相同名称的事物的本性不能不同。既然\OriginalTerm{光}{Light}的意义是同一且相同的，那么根据异端言论，\Quote{不可接近的}和\Quote{真}的添加便以特殊差异分开了共同本性，以致圣父的光被设想为一种东西，圣子的光被设想为另一种，彼此被特殊属性分开。因此，要么他推翻自己的立场，以免通过自己的陈述得出神性是复合的结论；要么他就不要控告我们，因为他可以看见自己的语言中包含了同样的内容。因为我们的陈述并未由此损害神性的\OriginalTerm{单纯性}{simplicity}，因为共通性和特殊差异并非\OriginalTerm{本质}{essence}，故它们的结合不应使主体成为复合。因为一方面，本质本身保持其在自然中的样貌，它就是其所是；另一方面，每个有理性的人都会说，这些——共通性和特殊差异——属于伴随的概念和属性：因为即使在我们人类中，也能辨别出与神圣本性的一些共通性，但神性并不因此而成为人性，人性也不会成为神性。因为我们相信天主是善的，也发现圣经中这个特性被用于人。但在每种情况下的特殊含义，确立了同义词使用所产生的共通性的区别。因为那善的泉源因之而得名；而那分享了某种善的人，也分有了这名称；天主并不因为与人类共用\Quote{善}这个称号而成为复合。由这些考虑，显然必须承认共通性的概念是一回事，本质的概念是另一回事，我们并不因此更要在那与数量无关、单纯的\OriginalTerm{神性}{Deity}中主张复合或多部分，因为我们所思索的某些属性要么被视为特殊，要么具有某种共通的意义。\par

但若你乐意，让我们转向他的另一项陈述，略过中间的胡言乱语。他辛苦地重复了亚里士多德对存在物的划分，详细阐述了\Quote{属}、\Quote{种}、\Quote{差异}和\Quote{个体}，并提出了所有范畴的技术用语，以损害我们的教义。让我们略过所有这些，将我们的论述转向对付他那沉重而不可抗拒的论点。因为他以狄摩西尼式的热忱武装了他的论点，在我们眼前俨然成为奥提塞里的第二个派阿尼亚人，模仿那位演说家在和我们斗争时的严厉。我将逐字转录我们这位作者的话。他说：\Quote{是的，但如果正如受生与不受生相对立，受生之光与不受生之光同样低劣，那么一个将被发现是光，另一个是黑暗。}愿有闲暇的人从他的话语中学习他处理这种对立的方式是何等尖刻，以及它何等精确地命中目标。但我恳求这位模仿者要么说出我们所说的话，要么尽可能接近地模仿它，否则，如果他根据自己的教育和能力来处理我们的论点，就让他以自己的身份说话，而不是以我们的身份。因为我希望没有人会如此误解我们的意思，以至于认为虽然\Quote{受生}与\Quote{不受生}在意义上矛盾，但一方是另一方的减弱。因为矛盾项之间的差异不在于强度的大小，而在于它们相互排斥的意义：例如，我们说一个人睡着或没睡着，坐着或没坐着，他存在或不存在，以及其他一切遵循同样模式的例子，其中对一者的否定就是对另一矛盾者的肯定。因此，正如活着不是不活着的减弱，而是其完全对立，同样我们认为被生不是不被生成的减弱，而是对立和矛盾，不允许任何中间项，因此一方所表达的东西与另一方所表达的东西在或多或少方面毫无关系。因此，那个说两个矛盾项中一个是相对于另一个\Emph{有缺陷的}的人，让他以自己的身份说话，而不是以我们的身份。因为我们朴素的言语说，与矛盾项对应的事物彼此之间的差异与其原样相同。因此，即使欧诺弥在光中看出与受生相对于不受生相同的分歧，我将重申我的陈述：正如在一例中，矛盾的一方与其对立面毫无共同之处，同样，如果\Quote{光}被置于两个矛盾项之一的同一侧，那么图形中剩余的位置当然必须分配给\Quote{黑暗}，因为对立之必要性根据先前矛盾项\Quote{受生}与\Quote{不受生}的类比来安排光这一项相对于其对立面的位置。这就是我们——按照我们轻视的作者所说，我们未经逻辑训练就尝试写作——用我们的乡土方言对我们新的派阿尼亚人作出的笨拙回答。但要看他如何与这矛盾争辩，以狄摩西尼式的强热向我们推进那些火热而吐火的话语，让那些喜欢一笑的人去研究这位演说家的论文本身。因为我们的笔不易被激起去反驳不敬的见解，但完全不适合嘲笑未受教化者的无知。\par

\SourceRecord{Translated by H.C. Ogle and H.A. Wilson.；Nicene and Post-Nicene Fathers, Second Series；Vol. 5.；Edited by Philip Schaff and Henry Wace.；Buffalo, NY: Christian Literature Publishing Co.,；1893.；<http://www.newadvent.org/fathers/290112.htm>.}
